% !TeX program = lualatex
% Fichier LaTEI généré depuis XML-TEI Métopes / Commons-Publishing.
% Zone technique (préambule, macros, mappings) : régénérable depuis le XML source.
% NE PAS MODIFIER la zone technique — elle sera écrasée à la prochaine génération.
% Zone éditoriale réversible : l'environnement lateiDocument dans ce fichier.
% Corriger uniquement dans la zone réversible.

\documentclass[11pt,twoside,openany]{book}

\usepackage[
  paperwidth=155mm,
  paperheight=230mm,
  top=30mm,
  bottom=19mm,
  inner=20mm,
  outer=30mm,
  headheight=14pt,
  headsep=8mm,
  footskip=10mm
]{geometry}
% Référentiel PURH v0.6 §6.2 : « le bandeau généré est situé environ 3,1 mm
% plus bas que dans le PDF imprimeur » — écart de rendu, pas une remise en
% cause de la marge du haut mesurée sur le maître InDesign (profile.
% margin_top_mm, qui reste la source de vérité pour la position du corps de
% texte). \topmargin remonte le bandeau de titre courant de 3,1 mm ;
% \headsep compense d'autant pour que le corps de texte, lui, démarre
% exactement à la même position qu'avant ce correctif.
\addtolength{\topmargin}{-3.1mm}
\addtolength{\headsep}{3.1mm}
\raggedbottom

\usepackage{fontspec}
\usepackage[french]{babel}
\usepackage{csquotes}
\usepackage{amsmath}
\usepackage{microtype}
\usepackage{indentfirst}
\usepackage{emptypage}
\usepackage[normalem]{ulem}
% [table] : nécessaire pour \rowcolor sur les lignes d'entête de tableau
% (référentiel §11.3, "fond foncé noir 30 %") — charge colortbl en plus des
% commandes \color habituelles (compatible avec l'usage \color[gray]{}
% déjà en place pour le titre courant).
\usepackage[table]{xcolor}

% Référentiel PURH v0.6 §12.1 : « le texte courant du PDF imprimeur
% correspond à un noir process à 90 % [K]. La sortie actuelle utilise du
% noir plein. » CMYK explicite (0,0,0,0.9), pas une valeur de gris RVB :
% "K" désigne le noir process de la quadrichromie, un concept distinct du
% gris \color[gray]{} utilisé ailleurs (titre courant) qui n'a pas de sens
% en CMJN. Appliqué globalement dès le début du document.
\definecolor{PURHBodyBlack}{cmyk}{0,0,0,0.9}
\AtBeginDocument{\color{PURHBodyBlack}}

\IfFontExistsTF{Chaparral Pro}
  {\setmainfont{Chaparral Pro}}
  {\setmainfont{TeX Gyre Pagella}}

\IfFontExistsTF{Josefin Sans}
  {\newfontfamily\PURHTitleFont{Josefin Sans}}
  {\newfontfamily\PURHTitleFont{TeX Gyre Heros}}

% Titraille PURH (référentiel §2.3-§2.5) : partie, article et section
% observés en Josefin Sans Thin, capitales, distinct du \PURHTitleFont
% "normal weight" ci-dessus utilisé ailleurs. Chargée comme famille à part
% (et non via \bfseries sur \PURHTitleFont) car "Thin" n'est pas une série
% NFSS standard que fontspec puisse sélectionner automatiquement — "Josefin
% Sans Thin" existe en revanche comme nom de famille indépendant portant
% elle-même ses propres graisses Thin (romain) et Thin Italic.
\IfFontExistsTF{Josefin Sans Thin}
  {\newfontfamily\PURHTitreFont{Josefin Sans Thin}}
  {\newfontfamily\PURHTitreFont{TeX Gyre Heros}}

\IfFontExistsTF{Latin Modern Mono}
  {\setmonofont{Latin Modern Mono}}
  {\setmonofont{TeX Gyre Cursor}}

% Romain, pas italique (référentiel PURH §2.3 : titres courants "Josefin
% Sans Thin/Light, 10 pt, romain" — l'italique systématique précédente était
% un défaut confirmé, pas un choix).
%
% Cinq vérifications humaines successives (2026-08-04/06) ont porté sur ce
% même réglage : la première jugeait le gris trop clair (corrigé en passant
% de \PURHTitreFont, famille Thin, à \PURHTitleFont, famille standard, sans
% \bfseries) ; la seconde a trouvé ce résultat trop noir et visuellement
% gras — le PDF imprimeur, lui, n'a « pas de graisse » — d'où un retour à la
% famille Thin avec un gris explicite (\color[gray]{0.25}, approximatif) ;
% la troisième a jugé ce gris encore trop clair, d'où un premier passage au
% même système que le fond d'entête de tableau et le texte courant
% (§11.3/§12.1) — une teinte CMJN noir X % plutôt qu'un gris RVB — à 50 %
% noir ; la quatrième a de nouveau jugé ce résultat trop clair au regard du
% PDF imprimeur, où le titre courant est « presque noir » : remonté à 85 %
% noir. La cinquième vérification, sur une page réelle complète d'un livre
% du corpus (pas seulement les liminaires), a de nouveau jugé le résultat
% trop clair à 85 % — poussé au maximum possible, 100 % noir (K plein),
% sans changer de famille (\PURHTitreFont, Josefin Sans Thin) : le résultat
% restait visiblement plus clair que le PDF imprimeur, confirmant qu'à ce
% stade le levier couleur seul était épuisé — un trait Thin, même à pleine
% saturation d'encre, couvre trop peu de surface pour lire aussi sombre
% qu'un trait plus épais à la même teinte.
%
% Sixième vérification (2026-08-06), comparaison directe côte à côte avec
% le PDF imprimeur (recadrage du bandeau à haute résolution) : la graisse
% du PDF imprimeur n'est PAS aussi fine qu'un trait Thin — plutôt un trait
% "Regular" (ni fin ni gras). Changement de levier plutôt que nouvel
% incrément de couleur : \PURHTitleFont (famille Josefin Sans complète,
% graisse Regular par défaut sans \bfseries — PAS \PURHTitreFont/Thin) à
% 75 % noir, un compromis délibéré entre les deux réglages précédents (ni
% le gris à 50 % jugé trop clair, ni l'aspect "gras" que la graisse
% Regular pourrait donner à 100 %) — comparé visuellement à un jeu
% d'échantillons Thin/Light/Regular/Medium × 75 %/100 % avant d'être
% retenu comme la combinaison la plus proche du PDF imprimeur.
%
% Septième vérification (2026-08-06) : le corps (Regular) validé comme
% correct par l'utilisateur — seule la teinte, encore un peu trop noire à
% 75 %, redescendue légèrement à 65 % (un ajustement fin, pas un nouveau
% changement de levier comme aux passes précédentes).
\newcommand{\PURHHeaderFont}{\PURHTitleFont\small\color[cmyk]{0,0,0,0.65}}

% Le corps et son pas de ligne sont fixés explicitement au lieu de dépendre
% de la table de tailles du \documentclass{book} choisi : elle donne un
% pas voisin mais pas identique au pas de grille cible (référentiel PURH
% §2.4, §5.3 : corps 11 pt sur une grille de 13,5 pt).
\renewcommand{\normalsize}{\fontsize{11pt}{13.5pt}\selectfont}
\normalsize

\newcommand{\PURHBookTitle}{Héraldique et papauté. Moyen Âge-Temps modernes. II}
\newcommand{\PURHBookSubtitle}{}
\newcommand{\PURHBookAuthor}{}
\newcommand{\PURHPublisher}{PURH}
\newcommand{\PURHYear}{2025}
\newcommand{\PURHDOI}{}
\newcommand{\PURHISBN}{979-10-240-1855-3}

\title{\PURHBookTitle}
\author{\PURHBookAuthor}
\date{\PURHYear}

\setlength{\parindent}{5mm}
\setlength{\parskip}{0pt}
\linespread{1.0}
\pretolerance=100
\tolerance=500
\hyphenpenalty=500
\exhyphenpenalty=500
\emergencystretch=3em
\clubpenalty=10000
\widowpenalty=10000
\displaywidowpenalty=10000

\usepackage{enumitem}

\setlist[itemize,1]{
  label=\textendash,
  leftmargin=1.5em,
  itemsep=0.2\baselineskip,
  topsep=0.4\baselineskip
}

\setlist[enumerate,1]{
  leftmargin=1.8em,
  itemsep=0.2\baselineskip,
  topsep=0.4\baselineskip
}

\usepackage[nobottomtitles*]{titlesec}
\usepackage{titletoc}

\setcounter{secnumdepth}{0}
% Référentiel PURH v0.6 §9 ("Table des matières") : la cible exclut les
% sections internes (intertitres) de la TDM — seuls le niveau des parties
% (\part, niveau -1) et celui des ouvertures de contribution/front matter
% (\addcontentsline{toc}{chapter}, niveau 0) doivent y figurer. tocdepth=2
% incluait à tort les <div type="section1">/<div type="section2"> (niveaux
% \section=1, \subsection=2), gonflant la TDM à trois pages au lieu de deux.
\setcounter{tocdepth}{0}

% Josefin Sans Bold, 16 pt, capitales — le référentiel indiquait Josefin
% Sans Thin (§2.5, §5.3, §4.3), contredit par vérification humaine directe
% du PDF généré face au PDF imprimeur : les titres de partie y sont noirs et
% gras, alors que le Thin rendait un texte maigre et grisâtre, peu lisible.
% Observation directe suivie ici, comme pour le séparateur de note de bas de
% page défini plus bas dans ce même préambule. \PURHTitleFont (pas
% \PURHTitreFont, qui ne charge que la graisse Thin et ne peut donc pas
% produire de \bfseries réel) charge la famille Josefin Sans complète, dont
% fontspec sélectionne automatiquement la graisse Bold via NFSS.
\titleformat{\chapter}[display]
  {\PURHTitleFont\bfseries\fontsize{16pt}{19pt}\selectfont\raggedright}
  {\chaptertitlename~\thechapter}
  {10pt}
  {\MakeUppercase}

% Titre de partie : même correctif Bold que \chapter ci-dessus. Toujours
% \part* (pas de numéro affiché) : le label reste vide plutôt que d'exposer
% un numéro de partie non requis.
\titleformat{\part}[display]
  {\PURHTitleFont\bfseries\fontsize{16pt}{19pt}\selectfont\centering}
  {}
  {0pt}
  {\MakeUppercase}

% Titre de section (intertitre), 12 pt, capitales. Le référentiel §2.5
% indiquait Josefin Sans Thin ; corrigé en Bold (même correctif que la
% titraille partie/chapitre/contribution ci-dessus) après vérification
% humaine directe du PDF généré face au PDF imprimeur : les intertitres y
% apparaissent noirs et gras, pas maigres (chantier de parité v0.6,
% 2026-08-04). Alignement et espacement non chiffrés par le référentiel pour
% ce niveau : conservés tels quels (raggedright, séparations existantes).
\titleformat{\section}[block]
  {\PURHTitleFont\bfseries\fontsize{12pt}{14pt}\selectfont\raggedright}
  {}
  {0pt}
  {\MakeUppercase}

% Sous-section (référentiel v0.6 §4.3 : seul le niveau "section" — c'est-à-
% dire section1 — est chiffré ; aucune valeur propre à section2/section3
% n'est fournie). Un correctif antérieur avait délibérément retiré le gras
% à ce niveau (confirmé alors sur le livre réel — <div type="section2"> y
% porte de vrais sous-titres phrastiques, jamais des libellés courts) ; la
% vérification humaine directe du 2026-08-04 demande au contraire le même
% traitement noir et gras que les autres intertitres, suivie ici. Capitales
% non demandées à ce niveau, contrairement à \section : inchangé. Tailles
% conservées telles quelles (\large/\normalsize) faute de mesure référentiel.
\titleformat{\subsection}[block]
  {\PURHTitleFont\bfseries\large\raggedright}
  {}
  {0pt}
  {}

\titleformat{\subsubsection}[block]
  {\PURHTitleFont\bfseries\normalsize\raggedright}
  {}
  {0pt}
  {}

\titlespacing*{\part}{0pt}{0pt}{30pt}
\titlespacing*{\chapter}{0pt}{20pt}{18pt}
\titlespacing*{\section}{0pt}{18pt}{10pt}
\titlespacing*{\subsection}{0pt}{14pt}{8pt}
\titlespacing*{\subsubsection}{0pt}{12pt}{6pt}

\addto\captionsfrench{
  \renewcommand{\contentsname}{Table des matières}
}

% Entrées de contribution/front matter (référentiel PURH v0.6 §9.1,
% vérification humaine directe du 2026-08-04) : « titre de la communication
% sans graisse, bas de casse, Chaparral, calé à gauche, série de points
% puis numéro de ligne ». Chaparral Pro reste la fonte ambiante (aucune
% famille à sélectionner) : \PURHTitleFont\bfseries retiré, \hfill remplacé
% par le même filet pointillé que \section ci-dessous. \addvspace{8pt}
% inconditionnel retiré : « pas de saut de ligne entre les références sauf
% changement de section » — l'espacement avant une nouvelle partie vient
% désormais du bloc \part ci-dessous, pas d'ici.
%
% Filet pointillé + numéro de page : laissés dans ce 4e argument dédié (pas
% déplacés dans le texte de l'entrée transmis à \addcontentsline). Une
% première tentative avait déplacé \titlerule*/\contentspage dans ce texte
% pour les faire tomber sur la ligne du titre plutôt que sur celle de
% l'auteur (vérification humaine directe du 2026-08-04 : « les points de
% suite doivent être au niveau du titre, pas de l'auteur ») — abandonnée :
% le paquet bookmark, qui construit automatiquement les signets PDF depuis
% CE MÊME texte d'entrée, ne tolère pas \contentspage/\\ dans cet argument
% ("Token not allowed in a PDF string", puis désynchronisation de
% titlesec — bug réel constaté par compilation). Solution retenue : le nom
% d'auteur n'est plus concaténé DANS le texte de l'entrée de chapitre du
% tout — voir \latei_finish_contribution_toc_entry: (latei_macros.tex), qui
% l'écrit désormais comme une ligne de TDM séparée, via \addtocontents
% (jamais capturée par le mécanisme de signets de bookmark, qui n'observe
% que \addcontentsline/\contentsline).
\titlecontents{chapter}
  [0pt]
  {}
  {}
  {}
  {\titlerule*[0.5pc]{.}\contentspage}

% Entrées de partie (le référentiel dit « titres de section », mais
% désigne bien ici le niveau \part de ce document — les véritables
% intertitres/sections sont exclus de la TDM depuis §9/tocdepth=0) : Josefin
% Sans Bold, capitales (vérification humaine directe du 2026-08-04 : à la
% différence du nom d'auteur sous chaque entrée de contribution, qui doit
% rester bas de casse — voir \lateiTocAuthorLine dans latei_macros.tex —
% les titres de ce niveau doivent au contraire être en petites capitales).
% \scshape ici n'a AUCUN effet sur \PURHTitleFont (Josefin Sans, qui n'a pas
% de véritables petites capitales OpenType — confirmé par compilation
% isolée, "Font shape .../b/sc undefined", substitution silencieuse) : la
% mise en majuscules se fait donc en amont, à la source du texte, via
% \MakeUppercase{#1} dans \lateiRenderHead (latei_macros.tex) — seul le
% corps réduit ici (12 pt contre 16 pt sur la page de la partie elle-même)
% distingue ce niveau de véritables grandes capitales.
\titlecontents{part}
  [0pt]
  {\addvspace{1\baselineskip}\PURHTitleFont\bfseries\fontsize{12pt}{14pt}\selectfont\centering}
  {}
  {}
  {}
  [\addvspace{1\baselineskip}]

\titlecontents{section}
  [1.5em]
  {}
  {}
  {}
  {\titlerule*[0.5pc]{.}\contentspage}

\titlecontents{subsection}
  [3em]
  {\small}
  {}
  {}
  {\titlerule*[0.5pc]{.}\contentspage}

\usepackage{fancyhdr}

\pagestyle{fancy}
\fancyhf{}
% Folios à l'extérieur (référentiel PURH §2.3) : LE (verso) et RO (recto).
% Les titres courants intérieurs (RE, LO) et \chaptermark sont pris en
% charge par latei_macros.tex, qui distingue verso (livre/partie) et recto
% (contribution en cours) — les définir ici aussi serait dupliqué et, pire,
% écrasé silencieusement puisque ce fichier est \input après ce préambule.
\fancyhead[LE]{\PURHHeaderFont\thepage}
\fancyhead[RO]{\PURHHeaderFont\thepage}
\renewcommand{\headrulewidth}{0pt}
\renewcommand{\footrulewidth}{0pt}

\fancypagestyle{plain}{%
  \fancyhf{}%
  \renewcommand{\headrulewidth}{0pt}%
  \renewcommand{\footrulewidth}{0pt}%
}

\usepackage[flushmargin]{footmisc}

\setlength{\footnotesep}{0.6\baselineskip}
% Espace avant les notes (référentiel §5.1 : 3 mm) : \skip\footins régit
% l'espace entre la fin du corps de texte et le début de la zone de notes
% (filet inclus) — 1,2\baselineskip (~5,7 mm à 11/13,5 pt) en donnait trop.
\setlength{\skip\footins}{3mm}
% Filet de notes (référentiel §5.1 : 0,25 pt de large sur 72 pt = 25,4 mm de
% long). \footnoterule par défaut de book.cls fait 0,4 pt sur
% 0,4\columnwidth (~42 mm sur l'empagement de ce profil) — les deux
% constats "environ 0,40 pt" / "environ 42 mm" du §5.2 correspondent très
% exactement à cette valeur par défaut, jamais personnalisée jusqu'ici.
\renewcommand{\footnoterule}{%
  \kern-3pt
  \hrule width 72pt height 0.25pt
  \kern 2.6pt
}
% \footnotelayout (footmisc) n'a plus d'effet : la redéfinition plus bas du
% texte de note pour le retrait négatif de première ligne applique
% directement \fontsize{8.5pt}{10.2pt} et court-circuite
% le mécanisme normal de footmisc qui l'invoque.

\usepackage{etoolbox}

% Citations observées : 9/11 pt, retrait gauche 10 mm (pas de retrait
% droit), ~4 mm avant/après (référentiel PURH §5.3 ; état antérieur :
% 11/14 pt, retraits gauche ET droit ≈1,5em, 8pt avant/après).
\renewenvironment{quote}
  {%
    \par\begingroup
    \fontsize{9pt}{11pt}\selectfont
    \list{}{\leftmargin=10mm\rightmargin=0pt}%
    \item\relax
  }
  {%
    \endlist
    \endgroup
  }

\AtBeginEnvironment{quote}{\vspace*{4mm plus 1pt minus 1pt}}
\AtEndEnvironment{quote}{\vspace*{4mm plus 1pt minus 1pt}}

\usepackage{graphicx}
\usepackage{caption}

\graphicspath{{media/}{images/}{assets/images/}}

\captionsetup{
  font=small,
  labelfont=bf,
  labelsep=period,
  skip=11pt
}

% Référentiel PURH v0.6 §11.3 : titre de tableau 9/11 pt, centré, 10 mm
% avant, 3,5 mm après — réglage propre aux tableaux (\captionsetup[table]),
% distinct du réglage générique ci-dessus qui reste seul à s'appliquer aux
% figures (aucune cible équivalente donnée pour elles dans le référentiel).
\DeclareCaptionFont{PURHTableCaptionFont}{\fontsize{9pt}{11pt}\selectfont}
\captionsetup[table]{
  font=PURHTableCaptionFont,
  labelfont=bf,
  labelsep=period,
  justification=centering,
  aboveskip=10mm,
  belowskip=3.5mm
}

\addto\captionsfrench{
  \renewcommand{\figurename}{Figure}
  \renewcommand{\tablename}{Tableau}
}

\usepackage{array}
\usepackage{longtable}
\usepackage{tabularx}
\usepackage{booktabs}

\usepackage{verse}
\usepackage{ragged2e}
\usepackage{xurl}
\urlstyle{same}

\usepackage[
  unicode=true,
  hidelinks,
  pdfusetitle
]{hyperref}

\usepackage{bookmark}

\hypersetup{
  pdftitle={\PURHBookTitle},
  pdfauthor={\PURHBookAuthor},
  pdfsubject={\PURHPublisher},
  pdfcreator={Impressions},
  pdfproducer={LuaLaTeX}
}

% Numéro calé à gauche avec retrait négatif de première ligne (observé
% directement sur le PDF imprimeur — le référentiel indiquait un point
% après le numéro, contredit par cette observation directe, suivie ici) :
% \leftskip aligne toutes les lignes de la note sur la même marge gauche ;
% \parindent négatif ramène uniquement la première ligne (celle du numéro)
% en deçà de cette marge, jusqu'au bord. \@thefnmark seul, sans point.
% Surtout PAS de \noindent ici : \noindent annule précisément l'effet du
% \parindent négatif qu'il est censé appliquer à la première ligne — bug
% réel constaté après vérification humaine du PDF généré (la première ligne
% restait alignée sur \leftskip comme les suivantes, aucun retrait négatif
% visible) ; laisser LaTeX indenter naturellement la première ligne de
% \parindent (donc la ramener à 0, à la marge) est ce qui produit le retrait
% négatif recherché.
% \AtBeginDocument, pas une simple redéfinition de préambule : hyperref/
% bookmark redéfinissent eux-mêmes \@makefntext pour y ajouter leurs
% ancres, mais le font via leur propre \AtBeginDocument — un
% \renewcommand direct ici, même placé après leur \usepackage, se faisait
% donc encore écraser au début du document (bug réel rencontré et vérifié :
% le correctif n'avait aucun effet tant qu'il n'était pas, lui aussi,
% différé). Étant chargés plus haut, leur crochet s'exécute avant celui-ci.
% \fontsize{8.5pt}{10.2pt} explicite ici, pas laissé au
% \footnotelayout de footmisc (référentiel §5.1 : Chaparral Pro 8,5/10,2 pt) :
% cette redéfinition complète de \@makefntext remplace celle de footmisc et
% n'appelle donc plus \footnotelayout — la taille du corps du texte de note
% retombait silencieusement sur celle ambiante (~9/11 pt observés, cf. §5.2)
% tant que ce défaut n'était pas identifié.
\makeatletter
\AtBeginDocument{%
  \renewcommand{\@makefntext}[1]{%
    \fontsize{8.5pt}{10.2pt}\selectfont
    \setlength{\leftskip}{1.2em}%
    \setlength{\parindent}{-1.2em}%
    \@thefnmark\enskip#1%
  }%
}
\makeatother

% -----------------------------------------------------------------
% Macros utilitaires PURH
% -----------------------------------------------------------------
% Visibilité auteur/affiliation sur l'ouverture de contribution (référentiel
% PURH v0.6 §7.2/§17 P1 item 3) : piloté par le profil de mise en page, pas
% par une valeur fixe — « le profil doit distinguer conservation des
% métadonnées et visibilité sur la page ». Doit être déclaré ici, avant que
% les macros LaTEI ne soient chargées : elles lisent ce drapeau dans
% \lateiContributionAuthor / \lateiContributionAffiliation mais ne le
% déclarent pas elles-mêmes.
\newif\iflateiShowContributionAuthor
\lateiShowContributionAuthorfalse

\newcommand{\PURHSeparator}{%
  \par\addvspace{1.5\baselineskip}%
  \noindent\rule{5cm}{0.4pt}%
  \par\addvspace{1.5\baselineskip}%
}

% Titre principal de la page de titre (référentiel PURH v0.6 §8.1,
% vérification humaine directe du 2026-08-04) : même hauteur de départ que
% le faux-titre (\PURHTitlePage partage désormais \vspace*{0.25\textheight}
% avec \PURHFalseTitle), mais corps plus grand — Josefin Sans Bold
% capitales, comme le faux-titre, juste agrandi. Aucune mesure
% millimétrique donnée par l'utilisateur pour ce corps précis : 22/26 pt
% choisi pour rester nettement plus grand que les 12/14 pt du faux-titre.
\newcommand{\PurhTitleMain}[1]{%
  {\PURHTitleFont\bfseries\fontsize{22pt}{26pt}\selectfont\centering\MakeUppercase{#1}\par}
}

% Sous-titre : Josefin Sans Bold bas de casse (pas de \MakeUppercase, à la
% différence du titre), un peu plus petit, centré sur deux lignes — largeur
% de boîte pour forcer le retour à la ligne, même principe que
% \PURHContributionTitleWidth pour les titres d'ouverture de contribution
% (référentiel §7.3 : reproduire la coupure plutôt que l'exiger en dur).
\newcommand{\PurhSubtitle}[1]{%
  \par\vspace{0.6\baselineskip}%
  \begin{center}
  \parbox{88mm}{\PURHTitleFont\bfseries\fontsize{15pt}{18pt}\selectfont\centering #1}
  \end{center}
  \vspace{0.4\baselineskip}%
}

% Responsabilité éditoriale : Chaparral Pro (fonte principale, aucun
% changement de famille) gras bas de casse, plus petit que le sous-titre,
% sur deux lignes explicites — « sous la direction de » toujours seul sur
% sa propre ligne, puis les noms sur la suivante (référentiel §8.1,
% vérification humaine directe du 2026-08-04). #1 est déjà le texte complet
% des deux lignes, séparées par \\ côté appelant (latei_driver.py), pas
% reconstruites ici : cette macro ne fait que les mettre en forme.
\newcommand{\PurhContributors}[1]{%
  \par\vspace{0.5\baselineskip}%
  {\bfseries\fontsize{11pt}{13pt}\selectfont\centering #1\par}%
  \vspace{0.6\baselineskip}%
}

% Mention finale de la page de titre (référentiel v0.7, vérification
% humaine directe du 2026-08-04, complétée le 2026-08-05) : nom complet de
% l'éditeur, majuscules grasses, sur UNE SEULE ligne, Chaparral (fonte
% principale, aucun changement de famille) — pas \PURHTitleFont/Josefin,
% contrairement à toute la titraille (§2.4/§3 du référentiel v0.7) : cette
% mention n'est pas un niveau de titraille, c'est une mention
% institutionnelle calée en bas de page.
%
% \resizebox{0.95\linewidth}{!}{...} plutôt qu'une simple
% \fontsize{14pt}{16pt} fixe (première version, abandonnée le
% 2026-08-05) : à 14 pt, « PRESSES UNIVERSITAIRES DE ROUEN ET DU HAVRE »
% (44 caractères) ne tenait pas sur une seule ligne à la largeur
% d'empagement de ce profil (~105 mm) et retombait sur deux lignes — bug
% réel constaté par vérification humaine directe du PDF généré.
% \resizebox emballe le texte dans une boîte horizontale non coupable
% (empêchant tout retour à la ligne, contrairement à un simple changement
% de \fontsize) puis la met à l'échelle pour occuper exactement 95 % de la
% largeur de la page — garantit à la fois la ligne unique et le remplissage
% « une bonne partie de la largeur de la page » quel que soit le nom
% affiché ou le profil, sans avoir à calculer un corps de police à la main.
\newcommand{\PurhPublisherMention}[1]{%
  \noindent\resizebox{0.95\linewidth}{!}{\bfseries\MakeUppercase{#1}}\par
}

\newenvironment{PurhBlockQuote}
  {%
    \begin{quote}
  }
  {%
    \end{quote}
  }

\newenvironment{PurhBibliography}
  {%
    \par
    \setlength{\parindent}{0pt}%
    \setlength{\leftskip}{0pt}%
  }
  {%
    \par
  }

% -----------------------------------------------------------------
% Liminaires PURH (référentiel v0.6 §8.1, P0 — "construire la séquence
% complète des liminaires") : pages blanches, faux-titre, crédits, page de
% titre. Construites depuis les métadonnées du livre par latei_driver.py,
% jamais depuis le corps LaTEI réversible. Toutes en pagestyle empty (ni
% folio ni titre courant), mais comptées dans la pagination — jamais
% \begin{titlepage} (dont le mode de compatibilité LaTeX 2.09 remet parfois
% \c@page à 1, un risque inutile ici où le compte doit au contraire
% continuer sans interruption depuis la toute première page physique).
% -----------------------------------------------------------------
\newcommand{\PURHBlankPage}{%
  \clearpage
  \thispagestyle{empty}%
  \mbox{}%
  \clearpage
}

% Vérification humaine directe du 2026-08-04 (faux-titre) :
% remonté sur la page (0.35\textheight -> 0.25\textheight, approximatif —
% « un peu plus haut », aucune mesure millimétrique donnée) et repassé en
% Josefin Sans Bold capitales, même corps que les titres de section
% (\titleformat{\section}, 12/14 pt) — le référentiel ne donnait aucune
% cible pour ce niveau spécifique, seule cette vérification fait foi ici,
% comme pour les autres corrections de graisse de ce chantier.
\newcommand{\PURHFalseTitle}[1]{%
  \clearpage
  \thispagestyle{empty}%
  \begin{center}
  \vspace*{0.25\textheight}
  {\PURHTitleFont\bfseries\fontsize{12pt}{14pt}\selectfont\MakeUppercase{#1}\par}
  \end{center}
  \clearpage
}

% Chaparral Pro (fonte principale du document, aucun changement de famille
% nécessaire) 10 pt — vérification humaine directe du 2026-08-04 : \small
% ne correspond pas nécessairement à 10 pt exactement (dépend de l'échelle
% de tailles du \documentclass), remplacé par une taille explicite.
% \vspace*{\fill} (pas \vspace*{0.3\textheight}, ni un \vfill non étoilé)
% : deuxième vérification humaine directe — le colophon doit être calé en
% bas de page, pas centré verticalement. Un \vfill nu en tout début de page
% est silencieusement absorbé par l'algorithme de coupure de page de TeX
% (jamais visible, confirmé par reproduction : le contenu restait centré
% près du haut) — surtout avec \raggedbottom actif dans ce document, qui
% autorise justement les pages à ne pas s'étirer jusqu'à \textheight.
% \vspace* (étoilé) protège la glue même en tête de page.
\newcommand{\PURHCreditsPage}[1]{%
  \clearpage
  \thispagestyle{empty}%
  \begin{center}
  \vspace*{\fill}
  \fontsize{10pt}{11.5pt}\selectfont
  #1
  \end{center}
  \clearpage
}

% Vérification humaine directe du 2026-08-04 : même hauteur de départ que
% le faux-titre (0.25\textheight, pas 0.15) — seul le corps du titre change
% entre les deux pages, pas sa position verticale.
\newcommand{\PURHTitlePage}[1]{%
  \clearpage
  \thispagestyle{empty}%
  \begin{center}
  \vspace*{0.25\textheight}
  #1
  \end{center}
  \clearpage
}

% ============================================================
% Macros LaTEI — régénérables, ne pas modifier
% ============================================================

% Experimental LaTEI macros for the reversible TEI/PURH convergence path.
% These definitions are typographic only; the reversible source remains the
% controlled LaTEI body produced by purh_site.reversible.latex_writer.

\usepackage{xparse}
\usepackage{xstring}
\usepackage{newunicodechar}

% Unicode spaces remain documentary characters in the reversible body; these
% typographic mappings only prevent missing glyphs in direct LaTEI compilation.
\newunicodechar{ }{~}
\newunicodechar{ }{\nobreak\,}
\newunicodechar{ }{\,}
\newunicodechar{‑}{\nobreak\hbox{-}}
\newunicodechar{″}{\ensuremath{\prime\prime}}
% Titres courants recto/verso distincts (référentiel PURH v0.5, "Titres
% courants") : verso (pages paires, RE) porte le titre du livre ou, à
% l'intérieur d'une partie, celui de la partie ; recto (pages impaires, LO)
% porte le titre court de la contribution en cours. Les deux commandes sont
% ré-exécutées à chaque page par fancyhdr — elles doivent donc rester des
% noms de macro, jamais leur valeur expansée au moment de cet \input.
\newcommand{\lateiCurrentRunningTitle}{}
\newcommand{\lateiVersoRunningTitle}{}
\fancyhead[RE]{\PURHHeaderFont\nouppercase{\lateiVersoRunningTitle}}
\fancyhead[LO]{\PURHHeaderFont\nouppercase{\lateiCurrentRunningTitle}}
\AtBeginDocument{%
  \fancyhead[RE]{\PURHHeaderFont\nouppercase{\lateiVersoRunningTitle}}%
  \fancyhead[LO]{\PURHHeaderFont\nouppercase{\lateiCurrentRunningTitle}}%
}

% Ouverture de contribution (référentiel PURH v0.5, "Ouvertures de
% contribution") : titre, sous-titre, auteur et traducteur viennent des
% <p rend="..."> du <div type="titlePage">, seule source de vérité visible —
% ne pas les dupliquer ailleurs. Titre en Josefin Sans Bold 16 pt capitales
% centré (vérification humaine directe du PDF généré face au PDF imprimeur :
% noir et gras là où le Thin — indiqué par le référentiel §2.5/§5.3, mais
% contredit par cette observation — rendait un texte maigre et grisâtre,
% peu lisible ; même correctif que \titleformat{\part}/{\chapter} dans le
% préambule). Sous-titre d'abord laissé en Thin Italic 12 pt (référentiel
% §5.3, non signalé comme défaut lors de la première vérification), puis
% re-corrigé en Bold Italic bas de casse le 2026-08-04 : nouvelle
% vérification humaine directe montrant explicitement gras + italique sur
% le PDF imprimeur pour ce niveau — pas de majuscules forcées, déjà correct
% (aucun \MakeUppercase ici, contrairement au titre). Positions verticales
% exactes (mm) non calibrées ici — hors périmètre de cette passe.
%
% Coupures de ligne du titre (référentiel PURH v0.5/v0.6 §7.3, "coupures
% éditoriales" — P1 item 4, exécuté le 2026-08-04 sur demande explicite,
% objectif : reproduire la règle implicite du PDF imprimeur plutôt qu'exiger
% une annotation manuelle par titre) : le référentiel avertit qu'il ne faut
% pas reconstruire ces coupures « par une largeur de boîte arbitraire ».
% Vérification empirique sur 7 titres réels du PDF imprimeur de *Dissimuler
% pour mieux régner* (7, 9, 15, 27, 26, 51, 191 caractères environ) montre
% cependant que la largeur de boîte N'EST PAS arbitraire : elle correspond
% de très près (104 mm calculé contre 103,4 mm mesuré indépendamment par
% analyse de pixels sur le rendu réel) à la largeur d'empagement du profil
% 155×230 (105 mm), et un simple \parbox à cette largeur reproduit déjà
% correctement, via le même algorithme Knuth-Plass que tout paragraphe
% LaTeX centré, le nombre de lignes et la coupure exacte de plusieurs titres
% tests (ex. "Les lieux de la conjuration : société secrète et hétérotopie
% dans la littérature romantique", 4 lignes identiques au PDF imprimeur).
% Sur d'autres titres, la coupure obtenue diffère de celle du PDF imprimeur
% d'un mot (ex. "Les espaces du secret à Clarens" → "LES ESPACES DU SECRET
% À" / "CLARENS" ici, contre "LES ESPACES DU SECRET" / "À CLARENS" dans le
% PDF imprimeur) — écart probablement dû à une règle InDesign de
% conservation des mots courts (« à », « de »…) avec la ligne suivante, que
% l'algorithme de LaTeX ne reproduit pas. Valeur retenue : 104mm, meilleur
% compromis observé plutôt qu'une largeur arbitraire non vérifiée.
\newcommand{\PURHContributionTitleWidth}{104mm}
\newcommand{\lateiContributionTitle}[1]{%
  \par\vspace{0.5\baselineskip}%
  \begin{center}
  \parbox{\PURHContributionTitleWidth}{\PURHTitleFont\bfseries\fontsize{16pt}{19pt}\selectfont\centering\MakeUppercase{#1}}%
  \end{center}
  \vspace{0.3\baselineskip}%
}
\newcommand{\lateiContributionSubtitle}[1]{%
  {\PURHTitleFont\bfseries\fontsize{12pt}{14pt}\selectfont\itshape\centering #1\par}%
  \vspace{0.4\baselineskip}%
}
% \iflateiShowContributionAuthor (défini dans le préambule, valeur pilotée
% par le profil de mise en page — référentiel PURH v0.6 §7.2/§17 P1 item 3 :
% « auteur et affiliation non imprimés » sur l'ouverture de contribution,
% mais métadonnées conservées) : le contenu réversible <p rend="author-aut">
% / <p rend="authority_affiliation"> reste inchangé dans le corps LaTEI, seul
% son affichage PDF est basculé par ce drapeau, jamais son existence.
% Signature de fin d'article (référentiel PURH v0.6 §8/§17, vérification
% humaine directe du 2026-08-04) : « les articles sont signés à la fin »,
% calés à droite — un emplacement distinct de l'ouverture de contribution
% (où auteur/affiliation restent masqués, cf. \iflateiShowContributionAuthor
% ci-dessus). \lateiContributionAuthor/\lateiContributionAffiliation
% capturent donc leur contenu dans ces macros globales, INDÉPENDAMMENT du
% drapeau d'affichage en tête de page, pour que \lateiRenderContributionSignature
% (appelée après le corps de la contribution, voir \lateiRenderFrontGroup/
% \lateiRenderChapterGroup) puisse la réémettre à la fin. Réinitialisées à
% vide au début de chaque contribution (mêmes macros) pour ne jamais laisser
% fuiter la signature d'une contribution vers la suivante qui n'en a pas.
% \@empty (le sentinelle vide du noyau LaTeX) exige la catégorie de code 11
% pour "@", seulement active entre \makeatletter et \makeatother — ce
% fichier n'en a pas besoin ailleurs et n'ouvre pas ce bloc ici. \@empty y
% serait tokenisé comme "\@" (caractère isolé) suivi du texte ordinaire
% "empty", qui s'imprimerait littéralement (bug réel constaté par
% compilation : "empty empty" visible sur la page). \lateiSignatureEmpty
% propre, sans "@", sert de sentinelle à la place.
\newcommand{\lateiSignatureEmpty}{}
\newcommand{\lateiSignatureAuthor}{}
\newcommand{\lateiSignatureAffiliation}{}
\newcommand{\lateiResetContributionSignature}{%
  \global\let\lateiSignatureAuthor\lateiSignatureEmpty
  \global\let\lateiSignatureAffiliation\lateiSignatureEmpty
}
% Nom d'auteur sous une entrée de TDM (référentiel §9.1) : écrit comme une
% ligne de TDM SÉPARÉE de l'entrée de chapitre elle-même, via
% \addtocontents (voir \latei_finish_contribution_toc_entry: plus bas) —
% jamais concaténé dans le texte transmis à \addcontentsline{{toc}}{{chapter}}
% : le paquet bookmark, qui construit les signets PDF automatiquement
% depuis CE texte, ne tolère pas les macros qu'on y ajouterait pour
% positionner le filet pointillé/numéro de page au niveau du titre plutôt
% que de l'auteur (bug réel constaté par compilation, voir le commentaire
% sur \titlecontents{{chapter}} dans latei_preamble.py). \par au lieu de la
% précédente rupture "\\" à l'intérieur d'un même paragraphe : ici, il
% s'agit bien d'un paragraphe séparé, plus de retrait suspendu de
% \titlecontents à préserver.
%
% \renewcommand{\textsc}[1]{#1} ICI, dans le corps de \lateiTocAuthorLine
% (vérification humaine directe du 2026-08-05, référentiel v0.7) : le nom
% de l'auteur doit apparaître bas de casse dans la TDM, alors que
% \lateiSignatureAuthor — passé tel quel, sans capture intermédiaire — porte
% le nom de famille en petites capitales (<hi rend="small-caps">Nom</hi>,
% routé vers \teiHi[rend={small-caps}]{Nom} puis \textsc{Nom}).
%
% Une première tentative capturait une copie "texte brut" à l'avance, au
% moment de \lateiContributionAuthor, via \protected@xdef + un \textsc local
% — abandonnée : \teiHi (comme toute commande définie par \NewDocumentCommand
% de xparse) est \protected au sens eTeX, donc jamais développée par un
% \edef/\xdef quel que soit l'état de \textsc au même moment ; le fichier
% .toc obtenu contenait encore \teiHi[rend={small-caps}]{Nom} tel quel (bug
% réel constaté : petites capitales toujours visibles dans le PDF malgré la
% capture). \renewcommand fonctionne en revanche PARFAITEMENT ici : ce bloc
% s'exécute au moment où \tableofcontents \input le fichier .toc et exécute
% RÉELLEMENT \lateiTocAuthorLine{...} — une exécution normale de macro, pas
% un \edef — où \teiHi peut s'exécuter (pas seulement se développer) et
% appeler ce \textsc local, correctement redéfini à ce moment précis.
\newcommand{\lateiTocAuthorLine}[1]{%
  \par\noindent\hspace*{1em}%
  {\bfseries\renewcommand{\textsc}[1]{##1}#1\par}%
}
\newcommand{\lateiContributionAuthor}[1]{%
  \global\def\lateiSignatureAuthor{#1}%
  \iflateiShowContributionAuthor
    {\normalsize\bfseries\centering #1\par}%
    \vspace{0.3\baselineskip}%
  \fi
}
\newcommand{\lateiContributionAffiliation}[1]{%
  \global\def\lateiSignatureAffiliation{#1}%
  \iflateiShowContributionAuthor
    {\small\centering #1\par}%
    \vspace{0.3\baselineskip}%
  \fi
}
% Chaparral Pro (fonte principale, aucun changement de famille), sans
% graisse, corps un peu plus petit que le corps du texte (11 pt -> 10 pt),
% calé à droite. Le prénom en bas de casse et le nom en petites capitales
% viennent déjà tels quels du balisage source (<p rend="author-aut">Prénom
% <hi rend="small-caps">Nom</hi></p>, déjà routé vers \teiHi[rend={small-caps}]
% ailleurs) : #1 n'a besoin d'aucune reformulation ici, seulement d'un
% nouvel habillage.
\newcommand{\lateiRenderContributionSignature}{%
  \ifx\lateiSignatureAuthor\lateiSignatureEmpty\else
    \par\vspace{0.8\baselineskip}%
    {\fontsize{10pt}{12pt}\selectfont\raggedleft\lateiSignatureAuthor\par}%
    \ifx\lateiSignatureAffiliation\lateiSignatureEmpty\else
      {\fontsize{10pt}{12pt}\selectfont\raggedleft\lateiSignatureAffiliation\par}%
    \fi
  \fi
}
% Espace vertical entre le bloc de titre d'ouverture (titre, sous-titre,
% auteur/affiliation) et le début du corps : absent du généré alors que le
% PDF imprimeur montre un grand blanc, estimé à 5-6 lignes de titre (16 pt /
% 19 pt d'interligne) par vérification humaine directe (chantier de parité
% v0.6, 2026-08-04) — 5,5 × 19 pt ≈ 105 pt. Valeur volontairement fixe, pas
% pilotée par le profil (comme les autres corps de la titraille), faute de
% mesure millimétrique du référentiel à ce niveau précis (le §7.2 documente
% une position absolue du corps depuis le haut de page, pas cet écart local).
\newcommand{\PURHContributionBodyGap}{\vspace{105pt}}
\newcommand{\lateiContributionTranslator}[1]{%
  {\small\itshape #1\par}%
  \vspace{0.6\baselineskip}%
}
\NewDocumentCommand{\lateiRenderParagraph}{O{} +m}{%
  \iflateiinfootnote
    #2\unskip\space
  \else
    \IfStrEq{\lateiHeadContext}{titlePage}{%
      \IfSubStr{#1}{rend={title-main}}{\lateiContributionTitle{#2}}{%
        \IfSubStr{#1}{rend={title-sub}}{\lateiContributionSubtitle{#2}}{%
          \IfSubStr{#1}{rend={author-aut}}{\lateiContributionAuthor{#2}}{%
            \IfSubStr{#1}{rend={authority\_affiliation}}{\lateiContributionAffiliation{#2}}{%
              \IfSubStr{#1}{rend={editor-trl}}{\lateiContributionTranslator{#2}}{\par #2\par}%
            }%
          }%
        }%
      }%
    }{%
      \IfStrEq{\lateiHeadContext}{figure}{%
        \IfSubStr{#1}{rend={caption}}{\par\small #2\par}{%
          \IfSubStr{#1}{rend={credits}}{\par\footnotesize #2\par}{\par #2\par}%
        }%
      }{\par #2\par}%
    }%
  \fi
}
\NewDocumentCommand{\teiP}{O{} +m}{\lateiRenderParagraph[#1]{#2}}

% Nested footnotes are not valid LaTeX. If a teiNote appears inside another
% note, render a visible symbolic marker and inline parenthesized content.
\newif\iflateiinfootnote
\lateiinfootnotefalse
\NewDocumentCommand{\teiNote}{O{} +m}{%
  \iflateiinfootnote
    \textsuperscript{*}%
  \else
    \begingroup
      \lateiinfootnotetrue
      \footnote{#2}%
    \endgroup
  \fi
}
\NewDocumentCommand{\teiRef}{O{} +m}{\lateiRenderRef[#1]{#2}}
\NewDocumentCommand{\teiTitle}{O{} +m}{%
  \IfSubStr{#1}{level={m}}{\textit{#2}}{%
    \IfSubStr{#1}{level={j}}{\textit{#2}}{%
      \IfSubStr{#1}{level={a}}{\enquote{#2}}{#2}%
    }%
  }%
}

% <formula notation="latex"> : le contenu documentaire porte déjà ses propres
% délimiteurs de mode mathématique (\(...\) ou \[...\]) — passthrough tel
% quel, sans échappement, pour que LaTeX le compose réellement au lieu
% d'afficher le code source échappé comme texte visible.
\NewDocumentCommand{\teiFormula}{O{} +m}{#2}

% Matter switches keep book.cls's structural side effects (cleardoublepage,
% numbered vs unnumbered chapters via \@mainmatterfalse/true) but not its
% page-numbering side effect. PURH pagination is arabe continue from the
% first page to the last (référentiel PURH v0.5 §5.5/§5.6, P0 — état de
% pagination) : liminaries are never roman-numbered, and the switch to the
% body never restarts the counter at 1. \pagenumbering resets the page
% counter every time it is called, so it must be called at most once for
% the whole document, on whichever of front/main/back matter starts first.
% Every "started" flag is set with \global: each group/chapter is rendered
% from inside its own TeX group (teiElement's xparse "+b" body argument,
% IfSubStr's branch braces...), and a *local* \xxxtrue would revert the
% instant that group closes — silently re-arming the guard and making
% \lateiEnsureMainMatter re-run \cleardoublepage + reset the page counter
% on every single chapter instead of only the first one.
\newif\iflateifrontmatterstarted
\newif\iflateimainmatterstarted
\newif\iflateibackmatterstarted
\newif\iflateiinbibliography
\newif\iflateipagenumberinginitialized
\lateifrontmatterstartedfalse
\lateimainmatterstartedfalse
\lateibackmatterstartedfalse
\lateiinbibliographyfalse
\lateipagenumberinginitializedfalse
\NewDocumentCommand{\lateiEnsureContinuousArabicPagination}{}{%
  \iflateipagenumberinginitialized\else
    \pagenumbering{arabic}%
    \global\lateipagenumberinginitializedtrue
  \fi
}
\makeatletter
\NewDocumentCommand{\lateiEnsureFrontMatter}{}{%
  \iflateifrontmatterstarted\else
    \cleardoublepage
    \global\@mainmatterfalse
    \lateiEnsureContinuousArabicPagination
    \global\lateifrontmatterstartedtrue
  \fi
}
\NewDocumentCommand{\lateiEnsureMainMatter}{}{%
  \iflateimainmatterstarted\else
    \cleardoublepage
    \global\@mainmattertrue
    \lateiEnsureContinuousArabicPagination
    \global\lateimainmatterstartedtrue
  \fi
}
\NewDocumentCommand{\lateiEnsureBackMatter}{}{%
  \iflateibackmatterstarted\else
    \if@openright\cleardoublepage\else\clearpage\fi
    \global\@mainmatterfalse
    \lateiEnsureContinuousArabicPagination
    \global\lateibackmatterstartedtrue
  \fi
}
\makeatother

\ExplSyntaxOn
\tl_new:N \l_latei_option_type_tl
\tl_new:N \l_latei_option_page_title_tl
\tl_new:N \l_latei_option_target_tl
\tl_new:N \l_latei_option_url_tl
\tl_new:N \l_latei_option_n_tl
\tl_new:N \l_latei_option_rend_tl
\tl_new:N \l_latei_option_xmlid_tl
\tl_new:N \l_latei_option_internaltarget_tl
\tl_new:N \l_latei_option_numcols_tl
\tl_new:N \l_latei_graphic_path_tl
\tl_new:N \l_latei_graphic_local_path_tl
\tl_new:N \l_latei_running_title_tl
\tl_new:N \g_latei_current_running_title_tl
\tl_new:N \g_latei_verso_running_title_tl
\prop_new:N \g_latei_graphics_map_prop
\prop_new:N \g_latei_running_titles_map_prop
\cs_generate_variant:Nn \prop_get:NnNTF { NVNTF }
% Le verso par défaut est le titre du livre ; \latei_markboth_verso:n le
% remplace par le titre de la partie en cours tant qu'aucune autre partie
% ne l'écrase (référentiel : "verso : titre du livre ou de la partie").
% \PURHBookTitle est déjà défini à ce stade (préambule inséré avant ce
% fichier), mais on stocke le nom de macro, pas sa valeur expansée : la
% même robustesse que pour \lateiCurrentRunningTitle plus haut.
\tl_gset:Nn \g_latei_verso_running_title_tl { \PURHBookTitle }
\RenewDocumentCommand{\lateiCurrentRunningTitle}{}
  {
    \tl_use:N \g_latei_current_running_title_tl
  }
\RenewDocumentCommand{\lateiVersoRunningTitle}{}
  {
    \tl_use:N \g_latei_verso_running_title_tl
  }
\NewDocumentCommand{\lateiDeclareGraphic}{m m}
  {
    \prop_gput:Nnn \g_latei_graphics_map_prop { #1 } { #2 }
  }
\NewDocumentCommand{\lateiDeclareRunningTitle}{m m}
  {
    \prop_gput:Nnn \g_latei_running_titles_map_prop { #1 } { #2 }
  }
% Recto (titre court de la contribution) et verso (titre du livre ou de la
% partie) sont mis à jour indépendamment — jamais la même valeur des deux
% côtés (référentiel : "passe deux fois la même valeur à \markboth").
\cs_new_protected:Npn \latei_markboth_recto:n #1
  {
    \prop_get:NnNTF \g_latei_running_titles_map_prop { #1 } \l_latei_running_title_tl
      { \tl_gset_eq:NN \g_latei_current_running_title_tl \l_latei_running_title_tl }
      { \tl_gset:Nn \g_latei_current_running_title_tl { #1 } }
    \markboth{\tl_use:N \g_latei_verso_running_title_tl}{\tl_use:N \g_latei_current_running_title_tl}
  }
\cs_new_protected:Npn \latei_markboth_verso:n #1
  {
    \prop_get:NnNTF \g_latei_running_titles_map_prop { #1 } \l_latei_running_title_tl
      { \tl_gset_eq:NN \g_latei_verso_running_title_tl \l_latei_running_title_tl }
      { \tl_gset:Nn \g_latei_verso_running_title_tl { #1 } }
    \markboth{\tl_use:N \g_latei_verso_running_title_tl}{\tl_use:N \g_latei_current_running_title_tl}
  }
\NewDocumentCommand{\lateiMarkBoth}{m}
  {
    \latei_markboth_recto:n { #1 }
  }
\NewDocumentCommand{\lateiMarkBothVerso}{m}
  {
    \latei_markboth_verso:n { #1 }
  }
\RenewDocumentCommand{\chaptermark}{m}{\lateiMarkBoth{#1}}
\cs_new_protected:Npn \latei_chapter:n #1
  {
    \prop_get:NnNTF \g_latei_running_titles_map_prop { #1 } \l_latei_running_title_tl
      {
        \tl_gset_eq:NN \g_latei_current_running_title_tl \l_latei_running_title_tl
        \chapter{#1}
        \lateiMarkBoth{#1}
      }
      {
        \tl_gset:Nn \g_latei_current_running_title_tl { #1 }
        \chapter{#1}
        \lateiMarkBoth{#1}
      }
  }
\NewDocumentCommand{\lateiChapter}{m}
  {
    \latei_chapter:n { #1 }
  }
\keys_define:nn { latei / options }
  {
    type .tl_set:N = \l_latei_option_type_tl,
    data-page-title .tl_set:N = \l_latei_option_page_title_tl,
    target .tl_set:N = \l_latei_option_target_tl,
    url .tl_set:N = \l_latei_option_url_tl,
    n .tl_set:N = \l_latei_option_n_tl,
    rend .tl_set:N = \l_latei_option_rend_tl,
    xmlid .tl_set:N = \l_latei_option_xmlid_tl,
    internaltarget .tl_set:N = \l_latei_option_internaltarget_tl,
    numcols .tl_set:N = \l_latei_option_numcols_tl,
    unknown .code:n = {},
  }
\cs_new_protected:Npn \latei_extract_options:n #1
  {
    \tl_clear:N \l_latei_option_type_tl
    \tl_clear:N \l_latei_option_page_title_tl
    \tl_clear:N \l_latei_option_target_tl
    \tl_clear:N \l_latei_option_url_tl
    \tl_clear:N \l_latei_option_n_tl
    \tl_clear:N \l_latei_option_rend_tl
    \tl_clear:N \l_latei_option_xmlid_tl
    \tl_clear:N \l_latei_option_internaltarget_tl
    \tl_clear:N \l_latei_option_numcols_tl
    \keys_set:nn { latei / options } { #1 }
  }
% Pose une cible PDF nommée d'après le xmlid TEI, s'il est présent dans les
% options déjà extraites par \latei_extract_options:n. Utilisé pour que les
% ancres/citations restent atteignables depuis \lateiRenderRef (cf. cibles
% "#id" internes) même quand l'élément lui-même n'a pas de macro de
% sectionnement propre (qui créerait sinon sa propre destination).
\cs_new_protected:Npn \latei_hypertarget_if_xmlid:
  {
    \tl_if_empty:NF \l_latei_option_xmlid_tl
      { \hypertarget{\tl_use:N \l_latei_option_xmlid_tl}{} }
  }
% Ponts sans "_"/":" : \ExplSyntaxOff plus bas rend ces caractères inaptes
% à composer un nom de macro hors de cette zone ; \teiHead et \teiCit, définis
% après \ExplSyntaxOff, doivent donc passer par ces noms classiques.
\NewDocumentCommand{\lateiExtractOptions}{m}
  {
    \latei_extract_options:n { #1 }
  }
\NewDocumentCommand{\lateiHypertargetIfXmlid}{}
  {
    \latei_hypertarget_if_xmlid:
  }
% Nombre de colonnes du tableau en cours (calculé côté Python, cf.
% _table_column_count dans latex_writer.py — LaTeX ne peut pas le déduire
% à la volée sans avoir déjà lu toutes les lignes).
\NewDocumentCommand{\lateiTableNumCols}{}
  {
    \tl_use:N \l_latei_option_numcols_tl
  }
% Rupture structurante d'ouverture de contribution : saut de page, entrée de
% TDM et titre courant à partir de data-page-title — mais aucun titre visible
% ni numérotation de chapitre ici (référentiel PURH v0.5, "Ouvertures de
% contribution" : « sans libellé ni numérotation de chapitre »). Le titre
% réellement affiché vient du <div type="titlePage"> rendu par ailleurs
% (\lateiContributionTitle) ; le dupliquer ici romprait cette règle.
% \cleardoublepage (jamais \clearpage) : « parties et articles sur recto,
% blancs techniques si nécessaire » (référentiel §5.2) — indépendant de
% l'option de classe openany/openright, qui ne régit que le \chapter brut
% du chemin <div type="chapter">, distinct de celui-ci.
% \setcounter{footnote}{0} : référentiel PURH v0.6 §5/§17, « notes remise à 1
% par contribution » — chaque ouverture (front matter, chapitre/article, back
% matter) redémarre sa propre numérotation de notes plutôt que de poursuivre
% celle de la contribution précédente sur tout le livre.
% \thispagestyle{empty} : référentiel PURH v0.6 §7.2/§17 P1 item 2, « aucun
% en-tête ni folio » sur l'ouverture de contribution — seule cette première
% page bascule en pagestyle vide ; les pages suivantes de la contribution
% reprennent le titre courant normalement (référentiel §6).
% \g_latei_pending_toc_title_tl : l'entrée de TDM elle-même est différée
% jusqu'à \latei_finish_contribution_toc_entry: (appelée après le corps de
% la contribution, comme \lateiRenderContributionSignature) — au moment où
% cette rupture s'exécute, \lateiSignatureAuthor n'est pas encore connu (la
% page de titre, qui le capture, n'a pas encore été rendue : elle fait
% partie de #2, qui s'exécute après). Référentiel PURH v0.6 §9.1,
% vérification humaine directe du 2026-08-04 : « prénom et nom de l'auteur,
% en gras » sous chaque titre de contribution dans la TDM.
\tl_new:N \g_latei_pending_toc_title_tl
\cs_new_protected:Npn \latei_add_contribution_opening_break:
  {
    \tl_if_empty:NF \l_latei_option_page_title_tl
      {
        \cleardoublepage
        \thispagestyle{empty}
        \phantomsection
        \tl_gset_eq:NN \g_latei_pending_toc_title_tl \l_latei_option_page_title_tl
        \exp_args:NV \latei_markboth_recto:n \l_latei_option_page_title_tl
        \setcounter{footnote}{0}
      }
  }
\cs_new_protected:Npn \latei_finish_contribution_toc_entry:
  {
    \tl_if_empty:NF \g_latei_pending_toc_title_tl
      {
        % Titre seul ici : filet pointillé + numéro de page restent gérés
        % par le 4e argument de \titlecontents{chapter} (préambule), comme
        % avant — l'auteur, lui, est écrit séparément ci-dessous, jamais
        % concaténé dans CE texte (voir le commentaire sur
        % \lateiTocAuthorLine, plus haut, pour la raison : compatibilité
        % avec les signets PDF automatiques du paquet bookmark).
        \addcontentsline{toc}{chapter}{\tl_use:N \g_latei_pending_toc_title_tl}
        \ifx\lateiSignatureAuthor\lateiSignatureEmpty\else
          \addtocontents{toc}{\protect\lateiTocAuthorLine{\lateiSignatureAuthor}}
        \fi
        \tl_gclear:N \g_latei_pending_toc_title_tl
      }
  }
\NewDocumentCommand{\lateiRenderFrontGroup}{O{} +m}
  {
    \lateiEnsureFrontMatter
    \latei_extract_options:n { #1 }
    \lateiResetContributionSignature
    \latei_add_contribution_opening_break:
    #2
    \latei_finish_contribution_toc_entry:
    \lateiRenderContributionSignature
  }
\NewDocumentCommand{\lateiRenderChapterGroup}{O{} +m}
  {
    \lateiEnsureMainMatter
    \latei_extract_options:n { #1 }
    \lateiResetContributionSignature
    \latei_add_contribution_opening_break:
    #2
    \latei_finish_contribution_toc_entry:
    \lateiRenderContributionSignature
  }
\NewDocumentCommand{\lateiRenderPartGroup}{O{} +m}
  {
    \lateiEnsureMainMatter
    % Toujours forcé ici, pas seulement via \lateiEnsureMainMatter (dont le
    % \cleardoublepage ne joue qu'à la toute première bascule vers le corps) :
    % chaque partie, même la 2e ou la 3e du livre, doit ouvrir en recto.
    \cleardoublepage
    % référentiel PURH v0.6 §7.1 : « page de style vide » sur l'ouverture de
    % partie — aucun en-tête ni folio, comme pour l'ouverture de contribution.
    \thispagestyle{empty}
    \begingroup
      \lateiSetHeadContext{part}
      #2
    \endgroup
  }
\NewDocumentCommand{\lateiRenderBackGroup}{O{} +m}
  {
    \lateiEnsureBackMatter
    \latei_extract_options:n { #1 }
    \latei_add_contribution_opening_break:
    #2
  }
\NewDocumentCommand{\lateiRenderRef}{O{} +m}
  {
    \latei_extract_options:n { #1 }
    \tl_if_empty:NTF \l_latei_option_internaltarget_tl
      {
        \tl_if_empty:NTF \l_latei_option_target_tl
          { #2 }
          { \href{\tl_use:N \l_latei_option_target_tl}{#2} }
      }
      { \hyperlink{\tl_use:N \l_latei_option_internaltarget_tl}{#2} }
  }
\cs_new_protected:Npn \latei_figure_fallback:
  {
    \fbox{\parbox{0.8\linewidth}{\centering\footnotesize Image absente ou non fournie}}
  }
\cs_new_protected:Npn \latei_include_graphic:n #1
  {
    \tl_if_blank:nTF { #1 }
      { \latei_figure_fallback: }
      {
        \IfFileExists{\detokenize{#1}}
          { \includegraphics[width=0.95\linewidth,keepaspectratio]{\detokenize{#1}} }
          { \latei_figure_fallback: }
      }
  }
\NewDocumentCommand{\lateiRenderGraphic}{O{}}
  {
    \latei_extract_options:n { #1 }
    \tl_clear:N \l_latei_graphic_path_tl
    \tl_if_empty:NF \l_latei_option_url_tl
      { \tl_set_eq:NN \l_latei_graphic_path_tl \l_latei_option_url_tl }
    \tl_if_empty:NT \l_latei_graphic_path_tl
      {
        \tl_if_empty:NF \l_latei_option_target_tl
          { \tl_set_eq:NN \l_latei_graphic_path_tl \l_latei_option_target_tl }
      }
    \tl_if_empty:NT \l_latei_graphic_path_tl
      {
        \tl_if_empty:NF \l_latei_option_n_tl
          { \tl_set_eq:NN \l_latei_graphic_path_tl \l_latei_option_n_tl }
      }
    \prop_get:NVNTF \g_latei_graphics_map_prop \l_latei_graphic_path_tl \l_latei_graphic_local_path_tl
      { \exp_args:NV \latei_include_graphic:n \l_latei_graphic_local_path_tl }
      { \exp_args:NV \latei_include_graphic:n \l_latei_graphic_path_tl }
  }
\ExplSyntaxOff

% Référentiel PURH v0.6 §11.2 : Chaparral Pro 10 pt, retrait suspendu 5 mm.
% \hangindent en em dépendait de la taille de fonte ambiante (jamais fixée
% par cette macro) ; passé en mm, valeur absolue conforme à la cible.
\NewDocumentCommand{\lateiBibliographyEntry}{+m}{%
  \par\noindent\fontsize{10pt}{12pt}\selectfont\hangindent=5mm\hangafter=1 #1\par
}
\NewDocumentCommand{\lateiBibliographyBlock}{+m}{%
  \par
  \begingroup
    \lateiinbibliographytrue
    \begin{PurhBibliography}
    #1%
    \end{PurhBibliography}
  \endgroup
  \par
}

% Heads are contextual in TEI: the same <head> can name a chapter, section,
% figure, table, or list. The surrounding LaTEI environment sets this context.
\newcommand{\lateiHeadContext}{fallback}
\newcommand{\lateiSetHeadContext}[1]{\def\lateiHeadContext{#1}}
\NewDocumentCommand{\lateiRenderHead}{m}{%
  % \part* ajoute déjà lui-même son entrée de TDM sous le shape [display] de
  % \titleformat{\part} (vérifié empiriquement) : un \addcontentsline{toc}
  % {part}{...} manuel ici produirait un doublon, pas un renfort.
  %
  % \MakeUppercase{#1} ici (et non \scshape/\textsc dans
  % \titlecontents{part}, préambule) : vérification humaine directe du
  % 2026-08-05 a montré que \scshape n'a AUCUN effet visible sur
  % \PURHTitleFont (Josefin Sans) — confirmé par compilation isolée, le
  % journal LaTeX rapporte explicitement « Font shape .../b/sc undefined »,
  % silencieusement remplacée par la forme normale (bug réel, jamais une
  % erreur bloquante, donc invisible sans vérifier le journal). Cette fonte
  % n'a pas de véritables petites capitales OpenType. Palliatif : capitales
  % pleines (\MakeUppercase), à la taille réduite (12 pt) déjà fixée par
  % \titlecontents{part} — plus petites que le titre de partie affiché sur
  % sa propre page (16 pt), donc lisibles comme un niveau distinct dans la
  % TDM, sans être de véritables petites capitales typographiques.
  % Appliqué ici (à la source du texte de la partie, avant \part*) plutôt
  % que dans \titlecontents{part} : \part* réutilise ce même #1, déjà
  % majuscule, pour son propre \addcontentsline{{toc}}{{part}}{{...}} interne —
  % la mise en majuscules du titre affiché sur la page (via le \MakeUppercase
  % déjà présent dans \titleformat{{\part}}) n'est pas affectée puisqu'elle
  % s'applique une seconde fois, sans effet, sur un texte déjà majuscule.
  \IfStrEq{\lateiHeadContext}{part}{\part*{\MakeUppercase{#1}}\lateiMarkBothVerso{#1}}{%
  \IfStrEq{\lateiHeadContext}{chapter}{\lateiChapter{#1}}{%
    \IfStrEq{\lateiHeadContext}{unnumberedchapter}{\chapter*{#1}\lateiMarkBoth{#1}\addcontentsline{toc}{chapter}{#1}}{%
    \IfStrEq{\lateiHeadContext}{section}{\section{#1}}{%
      \IfStrEq{\lateiHeadContext}{subsection}{\subsection{#1}}{%
        \IfStrEq{\lateiHeadContext}{subsubsection}{\subsubsection{#1}}{%
          \IfStrEq{\lateiHeadContext}{figure}{\par\smallskip\noindent\textbf{#1}\par\smallskip}{%
            \IfStrEq{\lateiHeadContext}{table}{\par\smallskip\noindent\textbf{#1}\par\smallskip}{%
              \IfStrEq{\lateiHeadContext}{list}{\item[]\textbf{#1}}{%
                \par\smallskip\noindent\textbf{#1}\par\smallskip% fallback head
              }%
            }%
          }%
        }%
      }%
    }%
  }%
  }%
  }%
}
\NewDocumentCommand{\teiHead}{O{} +m}{%
  \lateiExtractOptions{#1}%
  \lateiHypertargetIfXmlid
  \lateiRenderHead{#2}%
}

% Apply rend values predictably, including combined values in any order. This
% mirrors the stable renderer order: small caps / position / weight / italic.
\NewDocumentCommand{\lateiIfRendSmallCaps}{m m m}{%
  \IfSubStr{\detokenize{#1}}{\detokenize{small-caps}}{#2}{%
    \IfSubStr{\detokenize{#1}}{\detokenize{small_caps}}{#2}{%
      \IfSubStr{\detokenize{#1}}{\detokenize{small caps}}{#2}{#3}%
    }%
  }%
}
\NewDocumentCommand{\lateiIfRendBold}{m m m}{%
  \IfSubStr{\detokenize{#1}}{\detokenize{bold}}{#2}{\IfSubStr{\detokenize{#1}}{\detokenize{gras}}{#2}{#3}}%
}
\NewDocumentCommand{\lateiIfRendSup}{m m m}{%
  \IfSubStr{\detokenize{#1}}{\detokenize{sup}}{#2}{\IfSubStr{\detokenize{#1}}{\detokenize{exposant}}{#2}{#3}}%
}
\NewDocumentCommand{\lateiIfRendSub}{m m m}{%
  \IfSubStr{\detokenize{#1}}{\detokenize{sub}}{#2}{\IfSubStr{\detokenize{#1}}{\detokenize{indice}}{#2}{#3}}%
}
\NewDocumentCommand{\lateiIfRendUnderline}{m m m}{%
  \IfSubStr{\detokenize{#1}}{\detokenize{underline}}{#2}{\IfSubStr{\detokenize{#1}}{\detokenize{souligne}}{#2}{#3}}%
}
\NewDocumentCommand{\lateiIfRendStrikethrough}{m m m}{%
  \IfSubStr{\detokenize{#1}}{\detokenize{strikethrough}}{#2}{\IfSubStr{\detokenize{#1}}{\detokenize{barre}}{#2}{#3}}%
}
\NewDocumentCommand{\lateiIfRendUppercase}{m m m}{%
  \IfSubStr{\detokenize{#1}}{\detokenize{uppercase}}{#2}{\IfSubStr{\detokenize{#1}}{\detokenize{majuscule}}{#2}{#3}}%
}
\NewDocumentCommand{\lateiHiSmallCaps}{m +m}{%
  \lateiIfRendSmallCaps{#1}{\textsc{#2}}{#2}%
}
\NewDocumentCommand{\lateiHiPosition}{m +m}{%
  \lateiIfRendSup{#1}{\textsuperscript{#2}}{%
    \lateiIfRendSub{#1}{$_{#2}$}{#2}%
  }%
}
\NewDocumentCommand{\lateiHiBold}{m +m}{%
  \lateiIfRendBold{#1}{\textbf{#2}}{#2}%
}
\NewDocumentCommand{\lateiHiItalic}{m +m}{%
  \IfSubStr{\detokenize{#1}}{\detokenize{italic}}{\textit{#2}}{#2}%
}
\NewDocumentCommand{\lateiHiUnderline}{m +m}{%
  \lateiIfRendUnderline{#1}{\underline{#2}}{#2}%
}
\NewDocumentCommand{\lateiHiStrikethrough}{m +m}{%
  \lateiIfRendStrikethrough{#1}{\sout{#2}}{#2}%
}
\NewDocumentCommand{\lateiHiUppercase}{m +m}{%
  \lateiIfRendUppercase{#1}{\MakeUppercase{#2}}{#2}%
}
\NewDocumentCommand{\teiHi}{O{} +m}{%
  \lateiHiItalic{#1}{\lateiHiBold{#1}{\lateiHiPosition{#1}{\lateiHiSmallCaps{#1}{\lateiHiUnderline{#1}{\lateiHiStrikethrough{#1}{\lateiHiUppercase{#1}{#2}}}}}}}%
}

\NewDocumentCommand{\teiForeign}{O{} +m}{\textit{#2}}
\NewDocumentCommand{\teiTerm}{O{} +m}{#2}
\NewDocumentCommand{\teiName}{O{} +m}{#2}
\NewDocumentCommand{\teiPersName}{O{} +m}{#2}
\NewDocumentCommand{\teiPlaceName}{O{} +m}{#2}
\NewDocumentCommand{\teiOrgName}{O{} +m}{#2}
\NewDocumentCommand{\teiDate}{O{} +m}{#2}
\NewDocumentCommand{\teiNum}{O{} +m}{#2}
\NewDocumentCommand{\teiLabel}{O{} +m}{#2}
\NewDocumentCommand{\teiQ}{O{} +m}{\enquote{#2}}
\NewDocumentCommand{\teiSaid}{O{} +m}{\enquote{#2}}
\NewDocumentCommand{\teiAuthor}{O{} +m}{#2}
\NewDocumentCommand{\teiEditor}{O{} +m}{#2}
\NewDocumentCommand{\teiPublisher}{O{} +m}{#2}
\NewDocumentCommand{\teiBiblScope}{O{} +m}{#2}
\NewDocumentCommand{\teiIdno}{O{} +m}{#2}

\NewDocumentCommand{\teiGraphic}{O{}}{\lateiRenderGraphic[#1]}
\NewDocumentCommand{\teiPtr}{O{}}{}
\NewDocumentCommand{\teiLb}{O{}}{\linebreak{}}
\NewDocumentCommand{\teiPb}{O{}}{\ignorespaces}
% <anchor xml:id="..."/> : marqueur de position TEI, sans contenu visible —
% seule sa cible PDF (pour les <ref target="#..."> qui pointent dessus) compte.
\NewDocumentCommand{\teiAnchor}{O{}}{%
  \lateiExtractOptions{#1}%
  \lateiHypertargetIfXmlid
}

\NewDocumentEnvironment{teiDiv}{O{} +b}{%
  \IfSubStr{#1}{type={titlePage}}{%
    \begingroup
    \lateiSetHeadContext{titlePage}%
    #2%
    \PURHContributionBodyGap
    \endgroup
  }{%
    \begingroup
    \lateiSetHeadContext{fallback}%
    \IfSubStr{#1}{type={chapter}}{\lateiSetHeadContext{chapter}}{}%
    \IfSubStr{#1}{type={section}}{\lateiSetHeadContext{section}}{}%
    \IfSubStr{#1}{type={section1}}{\lateiSetHeadContext{section}}{}%
    \IfSubStr{#1}{type={section2}}{\lateiSetHeadContext{subsection}}{}%
    \IfSubStr{#1}{type={subsection}}{\lateiSetHeadContext{subsection}}{}%
    \IfSubStr{#1}{type={section3}}{\lateiSetHeadContext{subsubsection}}{}%
    #2%
    \endgroup
  }%
}{}
\NewDocumentEnvironment{teiList}{O{} +b}{%
  \begingroup
  \lateiSetHeadContext{list}%
  \IfSubStr{#1}{type={ordered}}{\begin{enumerate}}{\begin{itemize}}%
  #2%
  \IfSubStr{#1}{type={ordered}}{\end{enumerate}}{\end{itemize}}%
  \endgroup
}{}
\NewDocumentCommand{\teiItem}{O{} +m}{\item #2}
\NewDocumentEnvironment{teiQuote}{O{} +b}{\begin{PurhBlockQuote}#2\end{PurhBlockQuote}}{}
% <quote> utilisée en cours de phrase (dans <p>/<item>, ou n'importe où
% dans <note> — cf. _write_quote côté Python, qui décide du choix entre
% teiQuote et teiQuoteInline) : \teiQuote/PurhBlockQuote force un
% \begin{quote}, qui coupe le paragraphe même en plein milieu d'une phrase.
% \teiQuoteInline reste une macro simple (aucun \begingroup ni saut de
% paragraphe) et ajoute les guillemets via \enquote (csquotes), comme
% \teiQ déjà.
\NewDocumentCommand{\teiQuoteInline}{O{} +m}{\enquote{#2}}
% Poésie (référentiel PURH v0.5, "Poésie") : de vrais <lg>/<l> TEI passent
% par l'environnement `verse` (aligné à gauche, une ligne par \\, jamais
% justifié) — pas le bloc de citation, qui produisait des blancs excessifs
% en forçant des \linebreak dans un paragraphe justifié. <lg> imbriqués
% (strophes) s'emboîtent naturellement, `verse` gère leur espacement.
\NewDocumentEnvironment{teiLg}{O{} +b}{%
  \begin{verse}
  #2%
  \end{verse}
}{}
\NewDocumentCommand{\teiL}{O{} +m}{#2\\}
\NewDocumentEnvironment{teiFigure}{O{} +b}{%
  \begingroup
  \lateiSetHeadContext{figure}%
  \begin{center}#2\end{center}%
  \endgroup
}{}
\NewDocumentEnvironment{teiCit}{O{} +b}{%
  \lateiExtractOptions{#1}%
  \lateiHypertargetIfXmlid
  #2%
}{}
\NewDocumentEnvironment{teiBibl}{O{} +b}{\lateiBibliographyEntry{#2}}{}

% Tableau réel (longtable/booktabs), pas une simple suite de paragraphes
% centrés : le nombre de colonnes (numcols, calculé côté Python — cf.
% _table_column_count) fixe le gabarit de colonnes. Les teiCell d'une ligne
% sont déjà séparées par un "&" littéral dans la source (écrit côté Python,
% cf. _write_row), et teiRow/teiCell sont des macros (pas des environnements) :
% \multicolumn doit être le tout premier jeton lu par le scanner d'alignement
% de TeX pour une cellule fusionnée (@cols>1) — même non protégée, l'envelopper
% dans une macro (et a fortiori un environnement, qui insère un \begingroup
% avant le corps) le rend invisible à la détection de \omit ("Misplaced
% \omit", confirmé par reproduction isolée). \multicolumn est donc écrit tel
% quel côté Python (cf. _write_cell) pour les cellules fusionnées ; teiCell ne
% gère plus que la mise en forme (gras d'entête), jamais \multicolumn. Une
% ligne d'entête (role=label/header sur la ligne, cf. TEI) reçoit
% \toprule/\midrule ; \bottomrule ferme le tableau.
% Référentiel PURH v0.6 §11.3 : texte 8/9,5 pt, marges internes 2 mm,
% filets 0,25 pt, fond d'entête noir 30 % (un gris clair : 30 % de couverture
% d'encre, pas un fond sombre). Chaparral Pro déjà la fonte principale du
% document (\setmainfont) : aucun changement de famille nécessaire, seule la
% taille change ici. \heavyrulewidth/\lightrulewidth/\cmidrulewidth
% (booktabs) valent par défaut 0,08em/0,05em/0,05em — proportionnels à la
% fonte ambiante, donc jamais 0,25 pt de façon fiable sans réglage explicite.
\newif\iflateirowisheader
\lateirowisheaderfalse
\NewDocumentEnvironment{teiTable}{O{} +b}{%
  \lateiExtractOptions{#1}%
  \begingroup
  \lateiSetHeadContext{table}%
  \fontsize{8pt}{9.5pt}\selectfont
  \setlength{\tabcolsep}{2mm}
  \setlength{\heavyrulewidth}{0.25pt}
  \setlength{\lightrulewidth}{0.25pt}
  \setlength{\cmidrulewidth}{0.25pt}
  \par\begin{center}%
  \begin{longtable}{*{\lateiTableNumCols}{l}}
  \toprule
  #2%
  \bottomrule
  \end{longtable}%
  \end{center}\par%
  \endgroup
}{}
% \rowcolor pour les lignes d'entête est écrit directement par
% purh_site/reversible/latex_writer.py (_write_row), avant même ce \teiRow,
% jamais ici : même contrainte que \multicolumn sur une cellule fusionnée
% (cf. teiCell) — un \rowcolor conditionnel à l'intérieur de cette macro
% n'est plus le premier jeton vu par le scanner d'alignement de TeX et
% produit "Misplaced \noalign" (confirmé par reproduction isolée).
\NewDocumentCommand{\teiRow}{O{} +m}{%
  \IfSubStr{#1}{role={label}}{\lateirowisheadertrue}{%
    \IfSubStr{#1}{role={header}}{\lateirowisheadertrue}{\lateirowisheaderfalse}%
  }%
  #2%
  \\%
  \iflateirowisheader\midrule\lateirowisheaderfalse\fi
}
\NewDocumentCommand{\teiCell}{O{} +m}{%
  \lateiExtractOptions{#1}%
  \iflateirowisheader\bfseries\fi
  \IfSubStr{#1}{role={label}}{\bfseries}{\IfSubStr{#1}{role={header}}{\bfseries}{}}%
  #2%
}
% teiHeader is metadata, not running text. Document wrappers pass through;
% unknown non-header elements keep their content visible rather than vanish.
\NewDocumentEnvironment{teiElement}{O{} +b}{%
  \IfSubStr{#1}{name={teiHeader}}{}{%
    \IfSubStr{#1}{name={group}}{%
      \IfSubStr{#1}{type={section1}}{\lateiRenderPartGroup[#1]{#2}}{%
      \IfSubStr{#1}{type={article}}{\lateiRenderChapterGroup[#1]{#2}}{%
      \IfSubStr{#1}{type={chapter}}{\lateiRenderChapterGroup[#1]{#2}}{%
      \IfSubStr{#1}{type={acknowledgments}}{\lateiRenderFrontGroup[#1]{#2}}{%
      \IfSubStr{#1}{type={abbreviations}}{\lateiRenderFrontGroup[#1]{#2}}{%
      \IfSubStr{#1}{type={introduction}}{\lateiRenderFrontGroup[#1]{#2}}{%
      \IfSubStr{#1}{type={foreword}}{\lateiRenderFrontGroup[#1]{#2}}{%
      \IfSubStr{#1}{type={preface}}{\lateiRenderFrontGroup[#1]{#2}}{%
      \IfSubStr{#1}{type={dedication}}{\lateiRenderFrontGroup[#1]{#2}}{%
      \IfSubStr{#1}{type={conclusion}}{\lateiRenderBackGroup[#1]{#2}}{%
      \IfSubStr{#1}{type={bibliography}}{\lateiRenderBackGroup[#1]{#2}}{#2}%
      }}}}}}}}}}%
    }{%
      \IfSubStr{#1}{name={back}}{\lateiEnsureBackMatter #2}{%
      \IfSubStr{#1}{name={listBibl}}{\lateiBibliographyBlock{#2}}{%
      \IfSubStr{#1}{name={biblStruct}}{\lateiBibliographyEntry{#2}}{#2}%
      }}%
    }%
  }%
}{}

% ---------------------------------------------------------------------------
% Commandes \latei... — corrections de mise en page PDF uniquement.
% Ces commandes ne sont pas restaurées vers XML lors du retour Métopes.
% Le parseur LaTEI parcourt leur contenu sans créer d'élément TEI.
% ---------------------------------------------------------------------------

% Macro interne — traduit les tailles normalisées en dimension PDF.
% Valeurs acceptées : small (~0,5 ligne), medium (~1 ligne), large (~2 lignes).
\NewDocumentCommand{\lateiApplyVSpace}{m}{%
  \IfStrEqCase{#1}{%
    {small}{\addvspace{0.5\baselineskip}}%
    {medium}{\addvspace{\baselineskip}}%
    {large}{\addvspace{2\baselineskip}}%
  }[\PackageError{latei}{%
    Valeur inconnue : #1. Valeurs acceptees : small, medium, large.%
  }{}]%
}

% Indentation — alinéas
% --------------------------------------------------------------------------

% \lateiNoIndent{...} — supprime le retrait de première ligne du contenu.
\NewDocumentCommand{\lateiNoIndent}{+m}{{\parindent=0pt #1}}

% \lateiIndent{...} — force le retrait de première ligne même si LaTeX le supprime.
\newlength{\lateiParindent}
\AtBeginDocument{\setlength{\lateiParindent}{\parindent}}
\NewDocumentCommand{\lateiIndent}{+m}{{\setlength{\parindent}{\lateiParindent}#1}}

% Espacement vertical autonome
% --------------------------------------------------------------------------

% \lateiVSpace{small|medium|large} — insère un blanc vertical entre deux éléments.
\NewDocumentCommand{\lateiVSpace}{m}{\par\lateiApplyVSpace{#1}}

% Espacement vertical enveloppant
% --------------------------------------------------------------------------

% \lateiSpaceBefore{small|medium|large}{...} — ajoute un blanc avant le contenu.
\NewDocumentCommand{\lateiSpaceBefore}{m +m}{\par\lateiApplyVSpace{#1}#2}

% \lateiSpaceAfter{small|medium|large}{...} — ajoute un blanc après le contenu.
\NewDocumentCommand{\lateiSpaceAfter}{m +m}{#2\par\lateiApplyVSpace{#1}}

% Sauts de page autonomes
% --------------------------------------------------------------------------

% \lateiPageBreak — force un saut de page.
\NewDocumentCommand{\lateiPageBreak}{}{\newpage}

% \lateiClearPage — force une nouvelle page en vidant les flottants.
\NewDocumentCommand{\lateiClearPage}{}{\clearpage}

% Sauts de page enveloppants
% --------------------------------------------------------------------------

% \lateiPageBreakBefore{...} — force un saut de page avant le contenu.
\NewDocumentCommand{\lateiPageBreakBefore}{+m}{\newpage#1}

% \lateiPageBreakAfter{...} — force un saut de page après le contenu.
\NewDocumentCommand{\lateiPageBreakAfter}{+m}{#1\newpage}

% \lateiClearPageBefore{...} — vide les flottants et force une nouvelle page avant.
\NewDocumentCommand{\lateiClearPageBefore}{+m}{\clearpage#1}

% \lateiClearPageAfter{...} — force une nouvelle page avec vidage des flottants après.
\NewDocumentCommand{\lateiClearPageAfter}{+m}{#1\clearpage}

% Contrôle des coupures de page
% --------------------------------------------------------------------------

% \lateiKeepWithNext{...} — interdit une coupure de page immédiatement après.
\NewDocumentCommand{\lateiKeepWithNext}{+m}{#1\nopagebreak[4]}

% \lateiKeepTogether{...} — maintient le contenu entier sur une même page.
\NewDocumentCommand{\lateiKeepTogether}{+m}{\begin{samepage}#1\end{samepage}}

% \lateiNoPageBreakBefore{...} — interdit une coupure de page immédiatement avant.
\NewDocumentCommand{\lateiNoPageBreakBefore}{+m}{\nopagebreak[4]#1}

% \lateiNoPageBreakAfter{...} — interdit une coupure de page immédiatement après.
\NewDocumentCommand{\lateiNoPageBreakAfter}{+m}{#1\nopagebreak[4]}

% lateiDocument marks the reversible editorial zone in a monofile.
% Completely transparent at compilation — only serves as a delimiter for the parser.
\newenvironment{lateiDocument}{}{}

% ============================================================
% Mappings graphiques — régénérables, ne pas modifier
% ============================================================

% Experimental LaTEI graphic mapping.
% This file is a compilation artifact, not a reversible source.
% WARNING: Image not found for LaTEI package: ../icono/br/Ch02_Doulkaridou/fig1.jpg
% WARNING: Image not found for LaTEI package: ../icono/br/Ch02_Doulkaridou/fig2.jpg
% WARNING: Image not found for LaTEI package: ../icono/br/Ch02_Doulkaridou/fig3.jpg
% WARNING: Image not found for LaTEI package: ../icono/br/Ch02_Doulkaridou/fig4%20copie.jpg
% WARNING: Image not found for LaTEI package: ../icono/br/Ch02_Doulkaridou/fig5.jpg
% WARNING: Image not found for LaTEI package: ../icono/br/Ch02_Doulkaridou/fig8a.jpg
% WARNING: Image not found for LaTEI package: ../icono/br/Ch02_Doulkaridou/fig7.jpg
% WARNING: Image not found for LaTEI package: ../icono/br/Ch03_Loskoutoff_1/fig9.jpg
% WARNING: Image not found for LaTEI package: ../icono/br/Ch03_Loskoutoff_1/fig10.jpg
% WARNING: Image not found for LaTEI package: ../icono/br/Ch03_Loskoutoff_1/fig11.jpg
% WARNING: Image not found for LaTEI package: ../icono/br/Ch03_Loskoutoff_1/fig12.jpg
% WARNING: Image not found for LaTEI package: ../icono/br/Ch03_Loskoutoff_1/fig13.jpg
% WARNING: Image not found for LaTEI package: ../icono/br/Ch03_Loskoutoff_1/fig14.jpg
% WARNING: Image not found for LaTEI package: ../icono/br/Ch03_Loskoutoff_1/fig15.jpg
% WARNING: Image not found for LaTEI package: ../icono/br/Ch03_Loskoutoff_1/fig16.jpg
% WARNING: Image not found for LaTEI package: ../icono/br/Ch03_Loskoutoff_1/fig17_9-cappsist67-7v.jpg
% WARNING: Image not found for LaTEI package: ../icono/br/Ch03_Loskoutoff_1/fig18_10-cappsist67-6v.jpg
% WARNING: Image not found for LaTEI package: ../icono/br/Ch03_Loskoutoff_1/fig19_11-cappsist213.jpg
% WARNING: Image not found for LaTEI package: ../icono/br/Ch03_Loskoutoff_1/fig20.jpg
% WARNING: Image not found for LaTEI package: ../icono/br/Ch03_Loskoutoff_1/fig21.jpg
% WARNING: Image not found for LaTEI package: ../icono/br/Ch03_Loskoutoff_1/fig22.jpg
% WARNING: Image not found for LaTEI package: ../icono/br/Ch03_Loskoutoff_1/fig23.jpg
% WARNING: Image not found for LaTEI package: ../icono/br/Ch03_Loskoutoff_1/fig24.jpg
% WARNING: Image not found for LaTEI package: ../icono/br/Ch03_Loskoutoff_1/fig25.jpg
% WARNING: Image not found for LaTEI package: ../icono/br/Ch03_Loskoutoff_1/fig26.jpg
% WARNING: Image not found for LaTEI package: ../icono/br/Ch03_Loskoutoff_1/fig27.jpg
% WARNING: Image not found for LaTEI package: ../icono/br/Ch03_Loskoutoff_1/fig28.jpg
% WARNING: Image not found for LaTEI package: ../icono/br/Ch03_Loskoutoff_1/fig29.jpg
% WARNING: Image not found for LaTEI package: ../icono/br/Ch03_Loskoutoff_1/fig30.jpg
% WARNING: Image not found for LaTEI package: ../icono/br/Ch03_Loskoutoff_1/fig31.png
% WARNING: Image not found for LaTEI package: ../icono/br/Ch03_Loskoutoff_1/fig32.jpg
% WARNING: Image not found for LaTEI package: ../icono/br/Ch03_Loskoutoff_1/fig33.jpg
% WARNING: Image not found for LaTEI package: ../icono/br/Ch03_Loskoutoff_1/fig34.jpg
% WARNING: Image not found for LaTEI package: ../icono/br/Ch03_Loskoutoff_1/fig35.jpg
% WARNING: Image not found for LaTEI package: ../icono/br/Ch03_Loskoutoff_1/fig36.jpg
% WARNING: Image not found for LaTEI package: ../icono/br/Ch03_Loskoutoff_1/fig37.jpg
% WARNING: Image not found for LaTEI package: ../icono/br/Ch03_Loskoutoff_1/fig38.jpg
% WARNING: Image not found for LaTEI package: ../icono/br/Ch03_Loskoutoff_1/fig39.png
% WARNING: Image not found for LaTEI package: ../icono/br/Ch03_Loskoutoff_1/fig40.jpg
% WARNING: Image not found for LaTEI package: ../icono/br/Ch03_Loskoutoff_1/fig41.jpg
% WARNING: Image not found for LaTEI package: ../icono/br/Ch03_Loskoutoff_1/fig42.jpg
% WARNING: Image not found for LaTEI package: ../icono/br/Ch03_Loskoutoff_1/fig43.png
% WARNING: Image not found for LaTEI package: ../icono/br/Ch04_Federici/Fig44.jpg
% WARNING: Image not found for LaTEI package: ../icono/br/Ch04_Federici/Fig45.jpg
% WARNING: Image not found for LaTEI package: ../icono/br/Ch04_Federici/Fig46.jpg
% WARNING: Image not found for LaTEI package: ../icono/br/Ch04_Federici/Fig47.jpg
% WARNING: Image not found for LaTEI package: ../icono/br/Ch04_Federici/Fig49.jpg
% WARNING: Image not found for LaTEI package: ../icono/br/Ch04_Federici/Fig50.JPG
% WARNING: Image not found for LaTEI package: ../icono/br/Ch04_Federici/Fig51.jpg
% WARNING: Image not found for LaTEI package: ../icono/br/Ch04_Federici/Fig52.jpg
% WARNING: Image not found for LaTEI package: ../icono/br/Ch04_Federici/Fig53.jpg
% WARNING: Image not found for LaTEI package: ../icono/br/Ch04_Federici/Fig54.jpg
% WARNING: Image not found for LaTEI package: ../icono/br/Ch04_Federici/Fig55.JPG
% WARNING: Image not found for LaTEI package: ../icono/br/Ch04_Federici/Fig56.jpg
% WARNING: Image not found for LaTEI package: ../icono/br/Ch04_Federici/Fig57.jpg
% WARNING: Image not found for LaTEI package: ../icono/br/Ch04_Federici/Fig58.JPG
% WARNING: Image not found for LaTEI package: ../icono/br/Ch04_Federici/Fig59.jpg
% WARNING: Image not found for LaTEI package: ../icono/br/Ch04_Federici/Fig60.jpg
% WARNING: Image not found for LaTEI package: ../icono/br/Ch05_DelPesco/fig61.jpg
% WARNING: Image not found for LaTEI package: ../icono/br/Ch05_DelPesco/fig62.jpg
% WARNING: Image not found for LaTEI package: ../icono/br/Ch05_DelPesco/fig63.jpg
% WARNING: Image not found for LaTEI package: ../icono/br/Ch05_DelPesco/fig64.jpg
% WARNING: Image not found for LaTEI package: ../icono/br/Ch05_DelPesco/fig65.jpg
% WARNING: Image not found for LaTEI package: ../icono/br/Ch05_DelPesco/fig66.jpg
% WARNING: Image not found for LaTEI package: ../icono/br/Ch05_DelPesco/fig67.jpg
% WARNING: Image not found for LaTEI package: ../icono/br/Ch05_DelPesco/fig68.jpg
% WARNING: Image not found for LaTEI package: ../icono/br/Ch05_DelPesco/fig69.jpg
% WARNING: Image not found for LaTEI package: ../icono/br/Ch05_DelPesco/fig70.jpg
% WARNING: Image not found for LaTEI package: ../icono/br/Ch05_DelPesco/fig71.jpg
% WARNING: Image not found for LaTEI package: ../icono/br/Ch05_DelPesco/fig72.jpg
% WARNING: Image not found for LaTEI package: ../icono/br/Ch05_DelPesco/fig73.jpg
% WARNING: Image not found for LaTEI package: ../icono/br/Ch05_DelPesco/fig74.jpg
% WARNING: Image not found for LaTEI package: ../icono/br/Ch05_DelPesco/fig75.jpg
% WARNING: Image not found for LaTEI package: ../icono/br/Ch05_DelPesco/fig76.jpg
% WARNING: Image not found for LaTEI package: ../icono/br/Ch05_DelPesco/fig77.jpg
% WARNING: Image not found for LaTEI package: ../icono/br/Ch05_DelPesco/fig78.jpg
% WARNING: Image not found for LaTEI package: ../icono/br/Ch05_DelPesco/fig79.jpg
% WARNING: Image not found for LaTEI package: ../icono/br/Ch05_DelPesco/fig80.png
% WARNING: Image not found for LaTEI package: ../icono/br/Ch05_DelPesco/fig81.jpg
% WARNING: Image not found for LaTEI package: ../icono/br/Ch05_DelPesco/fig82.jpg
% WARNING: Image not found for LaTEI package: ../icono/br/Ch05_DelPesco/fig83.jpg
% WARNING: Image not found for LaTEI package: ../icono/br/Ch05_DelPesco/fig84.jpg
% WARNING: Image not found for LaTEI package: ../icono/br/Ch05_DelPesco/fig85.jpg
% WARNING: Image not found for LaTEI package: ../icono/br/Ch06_Cavicchioli/Fig86.jpg
% WARNING: Image not found for LaTEI package: ../icono/br/Ch06_Cavicchioli/Fig87.JPG
% WARNING: Image not found for LaTEI package: ../icono/br/Ch06_Cavicchioli/Fig88.jpg
% WARNING: Image not found for LaTEI package: ../icono/br/Ch06_Cavicchioli/Fig89.jpg
% WARNING: Image not found for LaTEI package: ../icono/br/Ch06_Cavicchioli/Fig90.jpg
% WARNING: Image not found for LaTEI package: ../icono/br/Ch06_Cavicchioli/Fig91.jpg
% WARNING: Image not found for LaTEI package: ../icono/br/Ch06_Cavicchioli/Fig92.JPG
% WARNING: Image not found for LaTEI package: ../icono/br/Ch06_Cavicchioli/Fig93.jpg
% WARNING: Image not found for LaTEI package: ../icono/br/Ch06_Cavicchioli/Fig94.jpg
% WARNING: Image not found for LaTEI package: ../icono/br/Ch06_Cavicchioli/Fig95.jpg
% WARNING: Image not found for LaTEI package: ../icono/br/Ch06_Cavicchioli/Fig96.jpg
% WARNING: Image not found for LaTEI package: ../icono/br/Ch06_Cavicchioli/Fig97.jpg
% WARNING: Image not found for LaTEI package: ../icono/br/Ch07_Loskoutoff_2/Fig98.JPG
% WARNING: Image not found for LaTEI package: ../icono/br/Ch07_Loskoutoff_2/Fig99.JPG
% WARNING: Image not found for LaTEI package: ../icono/br/Ch07_Loskoutoff_2/Fig100.JPG
% WARNING: Image not found for LaTEI package: ../icono/br/Ch07_Loskoutoff_2/Fig101.jpg
% WARNING: Image not found for LaTEI package: ../icono/br/Ch07_Loskoutoff_2/Fig102.JPG
% WARNING: Image not found for LaTEI package: ../icono/br/Ch07_Loskoutoff_2/Fig103.JPG
% WARNING: Image not found for LaTEI package: ../icono/br/Ch07_Loskoutoff_2/Fig104.jpg
% WARNING: Image not found for LaTEI package: ../icono/br/Ch07_Loskoutoff_2/Fig105.JPG
% WARNING: Image not found for LaTEI package: ../icono/br/Ch07_Loskoutoff_2/Fig106.JPG
% WARNING: Image not found for LaTEI package: ../icono/br/Ch07_Loskoutoff_2/Fig107.JPG
% WARNING: Image not found for LaTEI package: ../icono/br/Ch07_Loskoutoff_2/Fig108.JPG
% WARNING: Image not found for LaTEI package: ../icono/br/Ch07_Loskoutoff_2/Fig109.JPG
% WARNING: Image not found for LaTEI package: ../icono/br/Ch07_Loskoutoff_2/Fig110.JPG
% WARNING: Image not found for LaTEI package: ../icono/br/Ch07_Loskoutoff_2/Fig111.JPG
% WARNING: Image not found for LaTEI package: ../icono/br/Ch07_Loskoutoff_2/Fig112.jpg
% WARNING: Image not found for LaTEI package: ../icono/br/Ch07_Loskoutoff_2/Fig113.JPG
% WARNING: Image not found for LaTEI package: ../icono/br/Ch07_Loskoutoff_2/Fig114.jpeg
% WARNING: Image not found for LaTEI package: ../icono/br/Ch07_Loskoutoff_2/Fig115.JPG
% WARNING: Image not found for LaTEI package: ../icono/br/Ch08_Hablot/Fig116.jpg
% WARNING: Image not found for LaTEI package: ../icono/br/Ch08_Hablot/Fig117.jpg
% WARNING: Image not found for LaTEI package: ../icono/br/Ch08_Hablot/Fig118.jpg
% WARNING: Image not found for LaTEI package: ../icono/br/Ch08_Hablot/Fig119.jpg
% WARNING: Image not found for LaTEI package: ../icono/br/Ch08_Hablot/Fig120.jpg
% WARNING: Image not found for LaTEI package: ../icono/br/Ch08_Hablot/Fig121.jpg
% WARNING: Image not found for LaTEI package: ../icono/br/Ch09_Moureau/Fig123.JPG
% WARNING: Image not found for LaTEI package: ../icono/br/Ch09_Moureau/Fig124.jpg
% WARNING: Image not found for LaTEI package: ../icono/br/Ch09_Moureau/Fig125.jpg
% WARNING: Image not found for LaTEI package: ../icono/br/Ch09_Moureau/Fig126.jpg
% WARNING: Image not found for LaTEI package: ../icono/br/Ch09_Moureau/Fig127.jpg
% WARNING: Image not found for LaTEI package: ../icono/br/Ch09_Moureau/Fig128.jpg
% WARNING: Image not found for LaTEI package: ../icono/br/Ch09_Moureau/Fig129.jpg
% WARNING: Image not found for LaTEI package: ../icono/br/Ch09_Moureau/Fig130.JPG
% WARNING: Image not found for LaTEI package: ../icono/br/Ch09_Moureau/Fig131.jpg
% WARNING: Image not found for LaTEI package: ../icono/br/Ch10_Larraz/Fig132.jpg
% WARNING: Image not found for LaTEI package: ../icono/br/Ch10_Larraz/Fig133.jpg
% WARNING: Image not found for LaTEI package: ../icono/br/Ch10_Larraz/Fig134.jpg
% WARNING: Image not found for LaTEI package: ../icono/br/Ch10_Larraz/Fig135.jpg
% WARNING: Image not found for LaTEI package: ../icono/br/Ch10_Larraz/Fig136.png
% WARNING: Image not found for LaTEI package: ../icono/br/Ch10_Larraz/Fig137.jpg
% WARNING: Image not found for LaTEI package: ../icono/br/Ch10_Larraz/Fig138.jpg
% WARNING: Image not found for LaTEI package: ../icono/br/Ch11_Ronzani/Fig139.jpeg
% WARNING: Image not found for LaTEI package: ../icono/br/Ch11_Ronzani/Fig140.jpg
% WARNING: Image not found for LaTEI package: ../icono/br/Ch11_Ronzani/Fig141.jpg
% WARNING: Image not found for LaTEI package: ../icono/br/Ch11_Ronzani/Fig142.jpg
% WARNING: Image not found for LaTEI package: ../icono/br/Ch11_Ronzani/Fig143.jpg
% WARNING: Image not found for LaTEI package: ../icono/br/Ch11_Ronzani/Fig144.jpg
% WARNING: Image not found for LaTEI package: ../icono/br/Ch11_Ronzani/Fig145.jpg
% WARNING: Image not found for LaTEI package: ../icono/br/Ch11_Ronzani/Fig146.jpg
% WARNING: Image not found for LaTEI package: ../icono/br/Ch11_Ronzani/Fig147.jpg
% WARNING: Image not found for LaTEI package: ../icono/br/Ch11_Ronzani/Fig148.jpg
% WARNING: Image not found for LaTEI package: ../icono/br/Ch11_Ronzani/Fig149.jpg
% WARNING: Image not found for LaTEI package: ../icono/br/Ch11_Ronzani/Fig150.jpg
% WARNING: Image not found for LaTEI package: ../icono/br/Ch11_Ronzani/Fig151.jpg
% WARNING: Image not found for LaTEI package: ../icono/br/Ch11_Ronzani/Fig152.jpg
% WARNING: Image not found for LaTEI package: ../icono/br/Ch11_Ronzani/Fig153.jpg
% WARNING: Image not found for LaTEI package: ../icono/br/Ch11_Ronzani/Fig154.jpg
% WARNING: Image not found for LaTEI package: ../icono/br/Ch11_Ronzani/Fig155.jpg
% WARNING: Image not found for LaTEI package: ../icono/br/Ch11_Ronzani/Fig156.jpg
% WARNING: Image not found for LaTEI package: ../icono/br/Ch11_Ronzani/Fig157.jpg
% WARNING: Image not found for LaTEI package: ../icono/br/Ch11_Ronzani/Fig158.jpg
% WARNING: Image not found for LaTEI package: ../icono/br/Ch12_Boiteux/Fig159.jpg
% WARNING: Image not found for LaTEI package: ../icono/br/Ch12_Boiteux/Fig160.jpg
% WARNING: Image not found for LaTEI package: ../icono/br/Ch12_Boiteux/Fig161.JPG
% WARNING: Image not found for LaTEI package: ../icono/br/Ch12_Boiteux/Fig162.jpg
% WARNING: Image not found for LaTEI package: ../icono/br/Ch12_Boiteux/Fig164.jpg
% WARNING: Image not found for LaTEI package: ../icono/br/Ch12_Boiteux/Fig165.jpg
% WARNING: Image not found for LaTEI package: ../icono/br/Ch12_Boiteux/Fig166.jpg
% WARNING: Image not found for LaTEI package: ../icono/br/Ch12_Boiteux/Fig167.jpg
% WARNING: Image not found for LaTEI package: ../icono/br/Ch13_Cousinie/Fig168.jpg
% WARNING: Image not found for LaTEI package: ../icono/br/Ch13_Cousinie/Fig169.jpg
% WARNING: Image not found for LaTEI package: ../icono/br/Ch13_Cousinie/Fig170.jpg
% WARNING: Image not found for LaTEI package: ../icono/br/Ch13_Cousinie/Fig171.jpg
% WARNING: Image not found for LaTEI package: ../icono/br/Ch13_Cousinie/Fig172.jpg
% WARNING: Image not found for LaTEI package: ../icono/br/Ch13_Cousinie/Fig173.jpg
% WARNING: Image not found for LaTEI package: ../icono/br/Ch13_Cousinie/Fig174.JPG
% WARNING: Image not found for LaTEI package: ../icono/br/Ch13_Cousinie/Fig175.JPG
% WARNING: Image not found for LaTEI package: ../icono/br/Ch13_Cousinie/Fig176.jpg
% WARNING: Image not found for LaTEI package: ../icono/br/Ch13_Cousinie/Fig177.JPG
% WARNING: Image not found for LaTEI package: ../icono/br/Ch13_Cousinie/Fig178.jpg
% WARNING: Image not found for LaTEI package: ../icono/br/Ch13_Cousinie/Fig179.JPG
% WARNING: Image not found for LaTEI package: ../icono/br/Ch13_Cousinie/Fig180.JPG
% WARNING: Image not found for LaTEI package: ../icono/br/Ch13_Cousinie/Fig181.jpg

% ============================================================
% Mappings titres courants — régénérables, ne pas modifier
% ============================================================

% Experimental LaTEI running-title mapping.
% This file is a compilation artifact, not a reversible source.
\lateiDeclareRunningTitle{Aspects ludiques dans l’appareil héraldique des manuscrits de Léon X (1513-1521)}{Aspects ludiques dans l’appareil héraldique…}
\lateiDeclareRunningTitle{L’héraldique de Jules III Ciocchi del Monte (1550-1555) dans l’ornement pour le livre}{L’héraldique de Jules III Ciocchi del Monte (1550-1555)…}
\lateiDeclareRunningTitle{Les armoiries de Jules III}{Les armoiries de Jules III}
\lateiDeclareRunningTitle{Érudits et héraldique dans la Rome d’Urbain VIII (1623-1644)}{Érudits et héraldique dans la Rome d’Urbain VIII…}
\lateiDeclareRunningTitle{Publications de sujet héraldique : Gozze et Pietrasanta}{Publications de sujet héraldique : Gozze et Pietrasanta}
\lateiDeclareRunningTitle{Francesco Gualdi et les armoiries, entre érudition et conservation}{Francesco Gualdi et les armoiries, entre érudition…}
\lateiDeclareRunningTitle{L’héraldique d’Innocent X Pamphili (1644-1655)}{L’héraldique d’Innocent X Pamphili (1644-1655)}
\lateiDeclareRunningTitle{Les lettrés du milieu romain des Pamphili et l’origine de leur emblème nobiliaire}{Les lettrés du milieu romain des Pamphili et l’origine…}
\lateiDeclareRunningTitle{La « colombe portant le rameau d’olivier au bec » à Gubbio et à Rome}{La « colombe portant le rameau d’olivier au bec »…}
\lateiDeclareRunningTitle{L’écu des Pamphili durant le pontificat d’Innocent X (1644-1655) : deux exemples en architecture}{L’écu des Pamphili durant le pontificat d’Innocent X…}
\lateiDeclareRunningTitle{Épilogue. Une nouvelle fortune pour le blason pamphili : la Raccolta Di Varie Targhe di Roma de Filippo Juvarra (1711)}{Épilogue. Une nouvelle fortune pour le blason pamphili…}
\lateiDeclareRunningTitle{« Micat in vertice »}{« Micat in vertice »}
\lateiDeclareRunningTitle{Girolamo Farnèse, légat à Bologne : «  Aeterna hic LILIA semper erunt  »}{Girolamo Farnèse, légat à Bologne : « Aeterna hic LILIA…}
\lateiDeclareRunningTitle{«  Micet in vertice  » : Pietro Vidoni et la célébration d’Alexandre VII Chigi}{« Micet in vertice » : Pietro Vidoni et la célébration…}
\lateiDeclareRunningTitle{Piranèse, l’héraldique et Clément XIII Rezzonico (1758-1769)}{Piranèse, l’héraldique et Clément XIII Rezzonico…}
\lateiDeclareRunningTitle{Clément XIII Rezzonico : un cycle de dédicaces}{Clément XIII Rezzonico : un cycle de dédicaces}
\lateiDeclareRunningTitle{La mise en signe de l’espace par les papes : pratiques et conflits}{La mise en signe de l’espace par les papes : pratiques…}
\lateiDeclareRunningTitle{Quels signes pour les papes et l’Église ?}{Quels signes pour les papes et l’Église ?}
\lateiDeclareRunningTitle{S’emparer d’une ville par les « armes », quelques exemples}{S’emparer d’une ville par les « armes », quelques exemples}
\lateiDeclareRunningTitle{La collégiale Saint-Jean-Baptiste puis Saint-Louis de Castelnau-Bretenoux : un lien étroit et discret avec la papauté}{La collégiale Saint-Jean-Baptiste puis Saint-Louis…}
\lateiDeclareRunningTitle{Avignon, une « Altera Roma »}{Avignon, une « Altera Roma »}
\lateiDeclareRunningTitle{Autres témoignages héraldiques intéressants présents dans l’église}{Autres témoignages héraldiques intéressants présents…}
\lateiDeclareRunningTitle{La basilique : un monument marqué par ses constructeurs et ses utilisateurs}{La basilique : un monument marqué par ses constructeurs…}
\lateiDeclareRunningTitle{La basilique lieu du cérémonial : cérémonies éphémères et représentations}{La basilique lieu du cérémonial : cérémonies éphémères…}
\lateiDeclareRunningTitle{L’exercice du pouvoir : quelques exemples d’utilisation des emblèmes}{L’exercice du pouvoir : quelques exemples d’utilisation…}
\lateiDeclareRunningTitle{Seul en scène : l’apothéose Ludovisi ( circa 1623)}{Seul en scène : l’apothéose Ludovisi ( circa 1623)}
\lateiDeclareRunningTitle{Dématérialisation, inchoativité, contextualisation : les Apothéoses Barberini (1632-1636) et Médici (1644-1646)}{Dématérialisation, inchoativité, contextualisation…}
\lateiDeclareRunningTitle{L’écu familial et l’ancêtre fondateur : l’apothéose Altieri ( circa 1675-1676)}{L’écu familial et l’ancêtre fondateur : l’apothéose…}
\lateiDeclareRunningTitle{Extension scénique, dissémination et transformisme : l’apothéose Aldobrandini (1596-1602)}{Extension scénique, dissémination et transformisme…}
\lateiDeclareRunningTitle{De la scène au milieu, de l’interprétation à la subjectivation}{De la scène au milieu, de l’interprétation…}

\begin{document}

\lateiEnsureContinuousArabicPagination
\PURHBlankPage
\PURHBlankPage
\PURHFalseTitle{Héraldique et papauté. Moyen Âge-Temps modernes. II}
\PURHCreditsPage{© Presses universitaires de Rouen et du Havre, 2025.\par\vspace{0.1\baselineskip}2 place Émile Blondel – 76821 Mont-Saint-Aignan Cedex\par\vspace{0.1\baselineskip}\url{http://purh.univ-rouen.fr}\par\vspace{0.1\baselineskip}ISBN : 979-10-240-1855-3\par}
\PURHTitlePage{\PurhTitleMain{Héraldique et papauté. Moyen Âge-Temps modernes. II}
\vspace*{\fill}
\PurhPublisherMention{Presses universitaires de Rouen et du Havre}}
\PURHBlankPage

% ============================================================
% Zone éditoriale réversible — corrections autorisées ici
% ============================================================

\begin{lateiDocument}
\begin{teiElement}[name={TEI},change={commons\_edition}]

  \begin{teiElement}[name={teiHeader}]

    \begin{teiElement}[name={fileDesc}]

      \begin{teiElement}[name={titleStmt}]

        \teiTitle[type={main}]{Héraldique et papauté. Moyen Âge-Temps modernes. II}
      
\end{teiElement}
      \begin{teiElement}[name={editionStmt}]

        \begin{teiElement}[name={edition}]

          \teiDate{}
        
\end{teiElement}
        \begin{teiElement}[name={respStmt}]

          \begin{teiElement}[name={resp}]

\end{teiElement}
          \teiName{}
        
\end{teiElement}
        \begin{teiElement}[name={sponsor}]

\end{teiElement}
      
\end{teiElement}
      \begin{teiElement}[name={publicationStmt}]

        \begin{teiElement}[name={ab},type={expression}]

          \begin{teiBibl}[xmlid={bibl-001}]

            \teiPublisher{PURH}
            \begin{teiElement}[name={address}]

              \begin{teiElement}[name={addrLine}]

\end{teiElement}
            
\end{teiElement}
            \begin{teiElement}[name={pubPlace}]
Mont-Saint-Aignan
\end{teiElement}
            \begin{teiElement}[name={availability}]

              \begin{teiElement}[name={licence},target={-}]

\end{teiElement}
              \teiP[xmlid={p-001}]{}
            
\end{teiElement}
          
\end{teiBibl}
        
\end{teiElement}
        \begin{teiElement}[name={ab},type={book}]

          \begin{teiBibl}[xmlid={bibl-002}]

            \teiIdno[type={ISBN-13}]{979-10-240-1855-3}
            \teiIdno[type={ISSN}]{}
            \teiDate[type={publishing},when={2025-05-22}]{}
            \teiDate[type={legal-deposit}]{}
            \teiDate[type={imprint}]{}
            \teiBiblScope[unit={page}]{}
            \begin{teiElement}[name={availability}]

              \begin{teiElement}[name={licence},target={-}]

\end{teiElement}
              \teiP[xmlid={p-002}]{}
            
\end{teiElement}
            \begin{teiElement}[name={distributor}]

\end{teiElement}
            \teiName[type={sales-agent}]{}
            \teiRef[type={site},target={-}]{}
          
\end{teiBibl}
          \teiName[type={printer}]{
            \teiName{}
            \teiNum{}
          }
          \begin{teiElement}[name={address}]

            \begin{teiElement}[name={addrLine}]

\end{teiElement}
          
\end{teiElement}
          \begin{teiElement}[name={dim},type={printingvolume}]

\end{teiElement}
          \begin{teiElement}[name={dimensions}]

            \begin{teiElement}[name={width},unit={mm}]

\end{teiElement}
            \begin{teiElement}[name={height},unit={mm}]

\end{teiElement}
            \begin{teiElement}[name={depth},unit={mm}]

\end{teiElement}
          
\end{teiElement}
          \begin{teiElement}[name={dim},type={weight},unit={gr}]

\end{teiElement}
          \begin{teiElement}[name={extent}]

\end{teiElement}
          \begin{teiElement}[name={measure},type={price},unit={EUR}]

\end{teiElement}
        
\end{teiElement}
        \begin{teiElement}[name={ab},subtype={PDF},type={digital\_download}]

          \begin{teiBibl}[xmlid={bibl-003}]

            \teiIdno[type={ISBN}]{}
            \teiIdno[subtype={\#\#},type={DOI}]{}
            \teiRef[type={DOI}]{}
            \teiRef[type={site},target={-}]{}
            \begin{teiElement}[name={distributor}]

\end{teiElement}
            \teiDate[type={publishing}]{}
            \begin{teiElement}[name={availability}]

              \begin{teiElement}[name={licence}]

\end{teiElement}
              \teiP[xmlid={p-003}]{}
            
\end{teiElement}
            \teiName[type={sales-agent}]{}
          
\end{teiBibl}
          \begin{teiElement}[name={dim},type={weight},unit={Mo}]

\end{teiElement}
          \begin{teiElement}[name={measure},type={price},unit={EUR}]

\end{teiElement}
        
\end{teiElement}
        \begin{teiElement}[name={ab},subtype={EPUB},type={digital\_download}]

          \begin{teiBibl}[xmlid={bibl-004}]

            \teiIdno[type={ISBN-13}]{}
            \teiIdno[type={DOI}]{}
            \teiRef[type={DOI}]{}
            \teiRef[type={site},target={-}]{}
            \begin{teiElement}[name={distributor}]

\end{teiElement}
            \teiDate[type={publishing}]{}
            \begin{teiElement}[name={availability}]

              \begin{teiElement}[name={licence}]

\end{teiElement}
              \teiP[xmlid={p-004}]{}
            
\end{teiElement}
            \teiName[type={sales-agent}]{}
          
\end{teiBibl}
          \begin{teiElement}[name={dim},extent={},type={weight},unit={Mo}]

\end{teiElement}
          \begin{teiElement}[name={measure},type={price},unit={EUR}]

\end{teiElement}
          \begin{teiElement}[name={desc},type={accessMode}]
textual
\end{teiElement}
          \begin{teiElement}[name={desc},type={accessibilityFeature}]
ARIA
\end{teiElement}
          \begin{teiElement}[name={desc},type={accessibilityFeature}]
highContrastDisplay
\end{teiElement}
          \begin{teiElement}[name={desc},type={accessibilityFeature}]
readingOrder
\end{teiElement}
          \begin{teiElement}[name={desc},type={accessibilityFeature}]
structuralNavigation
\end{teiElement}
          \begin{teiElement}[name={desc},type={accessibilityFeature}]
tableOfContents
\end{teiElement}
          \begin{teiElement}[name={desc},type={accessibilityFeature}]
unlocked
\end{teiElement}
          \begin{teiElement}[name={desc},subtype={soundHazard},type={accessibilityHazard}]
noSoundHazard
\end{teiElement}
          \begin{teiElement}[name={desc},subtype={flashingHazard},type={accessibilityHazard}]
noFlashingHazard
\end{teiElement}
          \begin{teiElement}[name={desc},subtype={motionSimulationHazard},type={accessibilityHazard}]
noMotionSimulationHazard
\end{teiElement}
          \begin{teiElement}[name={desc},type={accessModeSufficient}]
textual
\end{teiElement}
          \begin{teiElement}[name={desc},type={accessibilitySummary}]

\end{teiElement}
        
\end{teiElement}
        \begin{teiElement}[name={ab},subtype={HTML},type={digital\_online}]

          \begin{teiBibl}[xmlid={bibl-005}]

            \begin{teiElement}[name={distributor}]

\end{teiElement}
            \teiName[type={CMS}]{}
            \teiIdno[type={ISBN}]{}
            \teiIdno[type={DOI}]{}
            \teiRef[type={DOI}]{}
            \teiRef[type={site},internaltarget={}]{}
            \teiDate[type={publishing}]{}
            \teiDate[type={embargoend}]{}
            \begin{teiElement}[name={availability}]

              \begin{teiElement}[name={licence}]

\end{teiElement}
              \teiP[xmlid={p-005}]{}
            
\end{teiElement}
          
\end{teiBibl}
        
\end{teiElement}
      
\end{teiElement}
      \begin{teiElement}[name={sourceDesc}]

        \begin{teiBibl}[xmlid={bibl-006}]

\end{teiBibl}
        \begin{teiBibl}[type={edition1},xmlid={bibl-007}]

          \teiDate{}
          \begin{teiElement}[name={extent},rend={format}]

\end{teiElement}
          \teiPublisher{}
          \teiIdno[type={ISBN}]{}
          \begin{teiElement}[name={edition},rend={type}]

\end{teiElement}
        
\end{teiBibl}
      
\end{teiElement}
    
\end{teiElement}
    \begin{teiElement}[name={profileDesc}]

      \begin{teiElement}[name={langUsage}]

        \begin{teiElement}[name={language},ident={fr-FR}]

\end{teiElement}
      
\end{teiElement}
      \begin{teiElement}[name={textClass}]

        \begin{teiElement}[name={keywords},scheme={keywords},xmllang={fr}]

          \begin{teiList}

            \teiItem{}
            \teiItem{}
            \teiItem{}
          
\end{teiList}
        
\end{teiElement}
      
\end{teiElement}
      \begin{teiElement}[name={abstract},rend={4e-couv}]

        \teiP[xmlid={p-006}]{}
      
\end{teiElement}
      \begin{teiElement}[name={abstract},rend={resume},xmllang={fr}]

        \teiP[xmlid={p-007}]{}
      
\end{teiElement}
      \begin{teiElement}[name={abstract},rend={abstract},xmllang={fr}]

        \teiP[xmlid={p-008}]{}
      
\end{teiElement}
    
\end{teiElement}
    \begin{teiElement}[name={revisionDesc}]

      \begin{teiElement}[name={change},when={2025},who={\#\#}]

\end{teiElement}
    
\end{teiElement}
  
\end{teiElement}
  \begin{teiElement}[name={text},type={book},xmlid={book}]

    \begin{teiElement}[name={front}]

\end{teiElement}
    \begin{teiElement}[name={group},type={book},xmlid={book-001}]

      \begin{teiElement}[name={group},type={acknowledgments},data-page-title={Remerciements},data-include-href={Ch00A\_Remerciements.xml},xmlid={acknowledgments-001}]

        \begin{teiElement}[name={front}]

          \begin{teiDiv}[type={titlePage},xmlid={titlepage-001}]

            \teiP[rend={title-main},xmlid={p-009}]{Remerciements}
          
\end{teiDiv}
        
\end{teiElement}
        \begin{teiElement}[name={body},xmlid={body-001}]

          \teiP[xmlid={p1}]{Toute notre gratitude va aux institutions qui ont permis la publication de ce recueil : le Centre Saint-Louis à Rome et son directeur François-Xavier Adam, hôtes du deuxième colloque \teiHi[rend={italic}]{Héraldique et papauté, Moyen Âge-Temps modernes} (7-9 juin 2022), l’École pratique des hautes études et son directeur d’études Laurent Hablot (laboratoire SAPRAT, Savoir et pratiques du Moyen Âge à l'époque contemporaine), l’université du Havre et le laboratoire GRIC (Groupe de recherche identités et cultures).}
        
\end{teiElement}
      
\end{teiElement}
      \begin{teiElement}[name={group},type={abbreviations},data-page-title={Table des abréviations},data-include-href={Ch00B\_Table\_des\_abreviations.xml},xmlid={abbreviations-001}]

        \begin{teiElement}[name={front}]

          \begin{teiDiv}[type={titlePage},xmlid={titlepage-002}]

            \teiP[rend={title-main},xmlid={p-010}]{Table des abréviations}
          
\end{teiDiv}
        
\end{teiElement}
        \begin{teiElement}[name={body},xmlid={body-002}]

          \teiP[xmlid={p1-2}]{ADP : Archivio Doria Pamphili}
          \teiP[xmlid={p2}]{ASBO : Archivio di Stato di Bologna}
          \teiP[xmlid={p3}]{AM : Archives municipales}
          \teiP[xmlid={p4}]{BAV : Biblioteca Apostolica Vaticana}
          \teiP[xmlid={p5}]{BM : Bibliothèque municipale}
          \teiP[xmlid={p6}]{BNCR : Biblioteca nazionale centrale di Roma}
          \teiP[xmlid={p7}]{BNF : Bibliothèque nationale de France}
          \teiP[xmlid={p8}]{\teiHi[rend={italic}]{DBI} : \teiHi[rend={italic}]{Dizionario biografico degli Italiani}}
        
\end{teiElement}
      
\end{teiElement}
      \begin{teiElement}[name={group},type={introduction},data-page-title={Introduction},data-page-authors={Yvan Loskoutoff},data-include-href={Ch01\_Introduction.xml},xmlid={introduction-001}]

        \begin{teiElement}[name={front}]

          \begin{teiDiv}[type={titlePage},xmlid={titlepage-003}]

            \teiP[rend={title-main},xmlid={p-011}]{Introduction}
            \teiP[rend={author-aut},xmlid={p-012}]{Yvan Loskoutoff}
          
\end{teiDiv}
        
\end{teiElement}
        \begin{teiElement}[name={body},xmlid={body-003}]

          \teiP[xmlid={p1-3}]{Ce deuxième volume d’\teiHi[rend={italic}]{Héraldique et papauté, Moyen Âge-Temps modernes} est issu du colloque tenu au centre Saint-Louis de Rome du 7 au 9 juin 2022. L’époque moderne y domine. Les réflexions entamées dans le précédent volume s’y poursuivent, notamment sur les décors monumentaux, sur les cérémonies et sur les manuscrits. De nouveaux champs d’investigation sont ouverts, entre autres, sur des figures de pontifes en particulier ou sur le livre imprimé. Le recueil s’organise en deux parties égales de six articles chacune : l’une se consacre à des pontificats spécifiques, l’autre à la monumentalisation de l’héraldique.}
          \teiP[xmlid={p2-2}]{Six pontifes bénéficient d’une étude monographique, deux pour la Renaissance, Léon X et Jules III, trois pour le \teiHi[rend={small-caps}]{xvii}\teiHi[rend={sup}]{e} siècle, Urbain VIII, Innocent X et Alexandre VII, un pour le \teiHi[rend={small-caps}]{xviii}\teiHi[rend={sup}]{e} siècle, Clément XIII. Elli Doulkaridou-Ramantani (université de Paris I), nous entretient de l’ornementation des manuscrits de Léon X Médicis (1513-1521), s’intéressant à un pontifical décoré par Attavante degli Attavanti vers 1520 (Pierpont Morgan Library) et à trois livres de chant polyphonique, provenant du scriptorium de Petrus Alamire, offerts au pape par Marguerite d’Autriche entre 1513 et 1519 (BAV, fonds de la chapelle Sixtine). Elle étudie la « logique cumulative des symboles médicéens », mettant en lumière une tendance à la facétie dans l’esprit de la grotesque dont la mode débutait alors à Rome\teiNote[n={1},place={foot},xmlid={ftn1}]{\teiP[xmlid={p-013}]{André Chastel, \teiHi[rend={italic}]{La grottesque}, Paris, Le Promeneur, 1988.}}, volontiers inspirée du blason\teiNote[n={2},place={foot},xmlid={ftn7}]{\teiP[xmlid={p-014}]{Yvan Loskoutoff, « Le symbolisme des \teiHi[rend={italic}]{palle} médicéennes à la villa Madama », \teiHi[rend={italic}]{Journal des savants}, n\teiHi[rend={sup}]{o} 2, 2001, p. 351-391.}}. Elle analyse notamment l’image de l’ourson amateur de miel, suggérant une allusion aux supports parlants des armes de Clarice Orsini, mère de Léon X.}
          \teiP[xmlid={p3-2}]{Yvan Loskoutoff (université Le Havre Normandie, Académie des Jeux floraux) continue l’histoire du livre, manuscrit et imprimé, avec Jules III Del Monte (1550-1555). Une dizaine de manuscrits décorés par Ferdinando Ruano ou Vincent Raymond (BAV, fonds Vatican latin et fonds de la chapelle Sixtine ; Bibliothèque capitulaire de Tolède) témoignent, comme pour Léon X, de l’usage des armoiries comme des devises, à l’instar des décors monumentaux du Vatican ou de la villa Giulia. L’ornement du livre imprimé illustre le développement des vignettes de titre armoriées, qui concurrencent les marques de libraires surtout à Rome, où la pratique constante de la gravure en grisaille semble annoncer le système des hachures héraldiques habituellement daté de la fin du siècle. On remarque un seul frontispice-titre pour le \teiHi[rend={italic}]{Missarum liber primus} de Palestrina (1554), mais c’est un remploi des \teiHi[rend={italic}]{Messes} de Morales, dédiées à Cosme de Médicis (1544), dont le contenu blasonnant a été modifié.}
          \teiP[xmlid={p4-2}]{Fabrizio Federici (chercheur postdoctoral), spécialiste de l’érudit romain Francesco Gualdi\teiNote[n={3},place={foot},xmlid={ftn8}]{\teiP[xmlid={p-015}]{Voir son article, « Un partisan de la France en quête de protection : ce “bon vieux” chevalier Francesco Gualdi et le cardinal Mazarin », dans Yvan Loskoutoff et Patrick Michel (dir.), \teiHi[rend={italic}]{Mazarin, Rome et l’Italie}, Mont-Saint-Aignan, Presses universitaires de Rouen et du Havre, 2021-2022, t. 2, p. 199-238.}}, nous transporte au siècle suivant sous le pontificat d’Urbain VIII Barberini (1623-1644). Il évoque d’abord le succès des traités de blason à cette époque, notamment ceux de Gauges de’ Gozze et celui du jésuite Silvestro Pietrasanta. Dans ce contexte il analyse la présence de l’héraldique dans l’ouvrage de Francesco Gualdi consacré aux tombes médiévales, les \teiHi[rend={italic}]{Memorie sepolcrali,} ainsi que son action pour la conservation des vestiges armoriaux. Il aborde enfin la question des origines commerçantes de la noblesse romaine et des traces qu’elles ont laissées dans l’héraldique, montrant comment on s’efforça de les faire disparaître, notamment dans la famille d’Urbain VIII.}
          \teiP[xmlid={p5-2}]{Daniela Del Pesco (université de Rome III) s’attache à Innocent X Pamphili (1644-1655). Elle anticipe sur la deuxième partie de ce volume en étudiant la place de ses armoiries dans l’architecture, non sans avoir fait leur histoire depuis la ville d’origine de la famille, Gubbio en Ombrie. Elle choisit deux exemples liés au jubilé de 1650 : la fontaine des Quatre-Fleuves place Navone, due au Bernin, et la restauration de la basilique du Latran, due à son concurrent, Borromini. Elle montre comment la colombe des Pamphili s’inscrit dans le monument du point de vue formel autant que symbolique. Elle achève en évoquant sa nouvelle fortune dans le recueil de gravures de Filippo Juvarra publié en 1711, les targes (écus armoriés) de Rome.}
          \teiP[xmlid={p6-2}]{Sonia Cavicchioli (université de Bologne) nous entraîne dans le pontificat suivant, celui d’Alexandre VII Chigi (1655-1667). Elle nous montre comment, à Bologne, les armoiries du pape furent célébrées, mais aussi occultées en raison des sentiments locaux d’indépendance. En effet, le recueil manuscrit des \teiHi[rend={italic}]{Insignia }du corps municipal, conservé aux Archives d’État de la ville, témoigne par ses miniatures que pour Girolamo Farnèse, légat de 1658 à 1662, on privilégia ses lys, ignorant toute référence au pape. En revanche, l’héraldique de son successeur, Pietro Vidoni, fut combinée avec celle d’Alexandre VII, non point dans les \teiHi[rend={italic}]{Insignia,} qui restèrent a-pontificaux aussi pour lui, mais dans la galerie de son appartement dont les fresques, dues à des peintres perspectivistes locaux, Giovanni Battista Caccioli et Domenico Santi, firent honneur mythologique aux monts et à l’étoile des Chigi.}
          \teiP[xmlid={p7-2}]{Yvan Loskoutoff conclut la première partie en évoquant les rapports d’un artiste et d’un pape, Giovanni Battista Piranesi (1720-1778) et Clément XIII Rezzonico, qui régna de 1758 à 1769. C’est une idée commune que le \teiHi[rend={small-caps}]{xviii}\teiHi[rend={sup}]{e} siècle vit le déclin du blason. Piranèse nous prouve le contraire. L’artiste de la Rome antique et moderne fut aussi celui de l’héraldique. Avant de s’appliquer à la papauté, ce goût se manifeste dans la dédicace des \teiHi[rend={italic}]{Antichità romane} (1756) à Milord Charlemont, mécène trop peu généreux dont le graveur « diffama » progressivement les armoiries dans les états successifs du frontispice. Il se montre aussi dans le frontispice des \teiHi[rend={italic}]{Vasi, candelabri, cippi}… (1778) où un vase antique fut modifié de façon à évoquer la licorne des armes du dédicataire, le comte Ivan Schouvaloff. Il apparaît encore dans certaine vue de Rome où le graveur n’ignore pas le détail médiéval. Les armoiries papales se développent dans un cycle de frontispices dédiés à Clément XIII Rezzonico, vénitien comme Piranèse : \teiHi[rend={italic}]{Della magnificenza ed architettura de’ Romani} (1761), \teiHi[rend={italic}]{Lapides Capitolini }(1762), \teiHi[rend={italic}]{Antichità d’Albano e di Castel Gandolfo} (1764), ou au neveu du pape, monseigneur Giovan Battista Rezzonico, \teiHi[rend={italic}]{Diverse maniere d’adornare i cammini} (1769). Cette suite permet d’examiner l’évolution de la représentation héraldique ainsi que son adaptation à un contexte figuratif antiquisant. Giovan Battista Rezzonico fut aussi le commanditaire du décor de la villa des chevaliers de Malte sur l’Aventin, non loin de Sainte-Prisca, Santa Maria del Priorato, seule œuvre architecturale de l’artiste où il déploya tous les pouvoirs décoratifs et symboliques du blason Rezzonico.}
          \teiP[xmlid={p8-2}]{La partie précédente a centré l’intérêt sur des figures de papes en particulier, privilégiant le livre, manuscrit ou imprimé, même si dans les trois derniers articles s’annonce le sujet de la présente partie : les manifestations de l’héraldique dans les monuments, sacrés ou profanes, durables ou éphémères. Pour quelles raisons et de quelle manière s’y inscrit-elle ? Cette partie s’ouvre sur une étude médiévale due à Laurent Hablot (École pratique des hautes études). Il s’interroge sur la « mise en signe de l’espace » et sur les conflits qu’elle pouvait entraîner en Italie, en raison de la concurrence du pape, de l’empereur et des communes. Il montre, ville après ville (Bologne, Orvieto, Viterbe), par quels moyens la papauté parvint à imposer ses insignes : apposition sur les monuments, concessions d’armoiries (premier exemple connu pour la ville de Viterbe au début du pontificat de Jean XXII, pape de 1316 à 1334 ; plus tard, Salon-de-Provence qui reçoit des armes d’Innocent VIII en 1488 après avoir refusé celles de l’Empire, du roi ou de son évêque). Le cas pontifical reste une exception du fait de la spécificité de l’institution ecclésiale.}
          \teiP[xmlid={p9}]{Emmanuel Moureau (université de Toulouse, conservateur Tarn-et-Garonne) nous emmène ensuite dans la collégiale Saint-Jean-Baptiste/Saint-Louis de Castelnau-Bretenoux (Quercy). Après avoir reconstitué l’histoire de ses liens avec la famille de Castelnau¸ il en élucide le décor, rejetant certaines idées admises pour aboutir à un vitrail représentant le pape Jules II Della Rovere sous les traits de saint Jules I\teiHi[rend={sup}]{er}, qui avait confirmé l’érection de la collégiale, et le cardinal François-Guillaume de Clermont-Lodève, cousin des Castelnau, ambassadeur de Louis XII auprès du même Jules II, sous les traits de saint Jérôme. Les stalles illustrent le blason des deux personnages, feuilles de rouvre parlantes pour les La Rovere et hermine pour les Clermont-Lodève.}
          \teiP[xmlid={p10}]{Camille Larraz (université de Genève) s’intéresse pour sa part au monument éphémère, en étudiant l’entrée du cardinal Alexandre Farnèse, légat à Avignon en 1553, elle montre la place de l’héraldique dans l’image d’une « autre Rome ». Les entrées sont à la mode, la bibliographie abondante\teiNote[n={4},place={foot},xmlid={ftn9}]{\teiP[xmlid={p-016}]{Principaux ouvrages : Bernard Guenée et Françoise Lehoux, \teiHi[rend={italic}]{Les entrées royales de 1328 à 1515}, Paris, CNRS, 1968 ; Christian Desplat et Paul Mironneau (dir.), \teiHi[rend={italic}]{Les entrées : gloire et déclin d’un cérémonial}, actes du colloque de Pau, Biarritz, Société Henri IV, 1997 ; Daniel Vaillancourt et Marie-France Wagner (dir.), \teiHi[rend={italic}]{Dix-septième siècle}, numéro spécial : \teiHi[rend={italic}]{Les entrées royales}, n\teiHi[rend={sup}]{o} 3, 2001. Le GRES (Groupe de recherche sur les entrées solennelles, université Concordia, Canada) publie des \teiHi[rend={italic}]{Cahiers}.}}, mais on a surtout privilégié les entrées royales. Le mérite de Camille Larraz est de se pencher sur le cas d’une ville pontificale. Les documents d’archives éclairent précisément le rôle des peintres en blason sous l’aspect le plus matériel (fourniture en pigments, ingrédients de la peinture, l’or et ses substituts, etc.), pour non moins de 1 250 armoiries. Elle met à profit le premier livret publié pour ce genre de cérémonie à Avignon par Honoré Henry, qui lui permet de détailler le programme. Elle étend ensuite l’analyse à deux gravures pour des ouvrages du siècle suivant déployant l’héraldique des papes, des légats et de la ville dans le contexte de « l’autre Rome », la première faisant probablement référence à l’entrée du cardinal et archevêque François Bordini en 1598.}
          \teiP[xmlid={p11}]{Le père Rocco Ronzani (université du Latran), revient à Rome pour nous faire visiter son église de Sainte-Prisca. Il retrace l’histoire du monument s’arrêtant au cardinal du titre presbytéral de 1599 à 1611, Benedetto Giustiniani, qui commanda des travaux de restauration pour le jubilé de 1600. Son écu, coupé d’une tour et d’une aigle, orne l’église comme la crypte, sous sa forme complète ou ses meubles séparés. Il en est de même pour le pape Sixte Quint qui l’avait élevé à la pourpre en 1586 et dont il portait les armes en chef. On voit aussi celles de Clément VIII Aldobrandini qui l’avait créé titulaire de Sainte-Prisca en 1599. C’est donc la biographie de Benedetto Giustiniani qui se monumentalise par le blason. Les armes d’autres papes y figurent aussi, en lien avec des phases de restauration de l’église : Calixte III Borgia († 1458), Clément XII Corsini († 1740), ou encore celles de saint Jean XXIII († 1963), cardinal titulaire de 1953 à 1958.}
          \teiP[xmlid={p12}]{Martine Boiteux (École des hautes études en sciences sociales) poursuit ces visites par la basilique Saint-Pierre. Elle examine la présence de l’emblématique, blasonnante ou non, dans les tombeaux, notamment ceux de trois femmes : Mathilde de Canossa, Christine de Suède, Clémentine Sobieska. Grâce aux témoignages visuels, elle restitue l’emblématique des cérémonies en temps de vacance du trône, elle peut être héraldique (funérailles, conclave, investiture, \teiHi[rend={italic}]{possesso}) ou non (cortège funèbre). Il en est de même pour l’exercice du pouvoir, le blason est absent de la fête de la \teiHi[rend={italic}]{Croce luminosa}, mais il peut être présent aux canonisations et béatifications. Il est absent des témoignages visuels consultés pour la fête de la \teiHi[rend={italic}]{Chinea} (présentation au pape d’une jument blanche en signe d’hommage du royaume de Naples) aux \teiHi[rend={small-caps}]{xvii}\teiHi[rend={sup}]{e} et \teiHi[rend={small-caps}]{xviii}\teiHi[rend={sup}]{e} siècles\teiNote[n={5},place={foot},xmlid={ftn10}]{\teiP[xmlid={p-017}]{Certains témoignages écrits signalent pourtant sa présence au \teiHi[rend={small-caps}]{xvi}\teiHi[rend={sup}]{e} siècle sur le caparaçon de la jument, voir Yvan Loskoutoff, « Introduction », dans Yvan Loskoutoff (dir.), \teiHi[rend={italic}]{Héraldique et papauté, Moyen Âge-Temps modernes}, Mont-Saint-Aignan, PURH, 2020, p. 16.}}. Il est présent dans l’accueil des personnalités. Ce jeu de présence et d’absence s’explique par le dialogue entre la gloire mondaine et l’humilité chrétienne sur lequel repose l’exercice même de la papauté\teiNote[n={6},place={foot},xmlid={ftn11}]{\teiP[xmlid={p-018}]{Sur le blason, témoignage de la gloire mondaine confronté à l’humilité chrétienne, voir Yvan Loskoutoff, \teiHi[rend={italic}]{L’armorial de Calliope, L’œuvre du Père Le Moyne S. J. (1602-1671) : littérature, héraldique, spiritualité}, Tübingen, Gunter Narr Verlag, 2000, troisième partie.}}.}
          \teiP[xmlid={p13}]{L’héraldique n’était pas réservée aux lieux de culte. On peut même dire que son origine laïque la destinait d’abord à l’espace profane. Nous l’avons vu avec Camille Larraz, nous le voyons à nouveau avec Frédéric Cousinié (université de Rouen). Il examine les apothéoses armoriales dans les palais des grandes familles pontificales : Ludovisi (par Giovanni Luigi Valesio, artiste bolonais fort inventif en matière de blason), Barberini (par Pietro da Cortona, comparée à la salle de Mars Médicis au palais des Offices à Florence, du même), Altieri (par Carlo Maratta). Il s’attache surtout à la salle Clémentine du Vatican, exécutée par les frères Alberti pour Clément VIII Aldobrandini. Il établit une typologie de ces représentations : monumentalisation ou amplification, soutien ou élévation par des anges (\teiHi[rend={italic}]{stemma in arrivo, stemma in partenza}), extraction et dissémination des meubles, animation-activation des objets et « transformisme ». L’effet produit sur le spectateur nous convainc que l’héraldique constitue l’un des moyens d’expression privilégiés du « baroque », où finalement s’annule la distinction entre sacré et profane.}
        
\end{teiElement}
      
\end{teiElement}
      \begin{teiElement}[name={group},type={section1},xmlid={section1-001}]

        \teiHead[xmlid={premiere-partie-papes-001}]{Première partie : Papes}
        \begin{teiElement}[name={group},type={article},data-page-title={Aspects ludiques dans l’appareil héraldique des manuscrits de Léon X (1513-1521)},data-page-authors={Elli Doulkaridou-Ramantani},data-include-href={Ch02\_Aspects\_ludiques.xml},xmlid={article-001}]

          \begin{teiElement}[name={front}]

            \begin{teiDiv}[type={titlePage},xmlid={titlepage-004}]

              \teiP[rend={title-main},xmlid={p-019}]{Aspects ludiques dans l’appareil héraldique des manuscrits de Léon X (1513-1521)}
              \teiP[rend={author-aut},xmlid={p-020}]{Elli Doulkaridou-Ramantani}
            
\end{teiDiv}
          
\end{teiElement}
          \begin{teiElement}[name={body},xmlid={body-004}]

            \teiP[xmlid={p1-4}]{À l’occasion de ce deuxième volet autour de l’héraldique et de la papauté, je propose quelques réflexions sur l’appareil héraldique des manuscrits du pape Léon X. À travers cette étude, je vais montrer comment l’héraldique érudite des Médicis, et plus particulièrement celle cachée dans leurs manuscrits décorés, permet une vue plus intime et parfois nuancée du personnage de ce pontife, dont la vie coïncide avec des événements exceptionnels pour le continent européen – de Magellan jusqu’à Martin Luther – avec des ramifications pour l’Italie.}
            \teiP[xmlid={p2-3}]{J’examinerai principalement deux types de manuscrits. En premier lieu, un pontifical\teiNote[n={1},place={foot},xmlid={ftn1-2}]{\teiP[xmlid={p-021}]{C’est-à-dire un livre contenant les textes prononcés lors de la préparation vestimentaire et rituelle du pape avant d’entrer dans la chapelle pour la messe. }} daté un an avant la mort du pape et, en deuxième lieu, les livres de musique polyphonique qui lui ont été offerts par Marguerite d’Autriche (1480-1530), fille de l’empereur Maximilien I\teiHi[rend={sup}]{er} (1459-1519) et de Marie de Bourgogne\teiNote[n={2},place={foot},xmlid={ftn36}]{\teiP[xmlid={p-022}]{La culture du manuscrit enluminé dans la cour de Bourgogne est considérable. Il suffit de mentionner dans le cadre de cette étude le livre d’heures aujourd’hui conservé à Vienne. Voir Bret Rothstein, « The Rule of Metaphor and the Play of the Viewer in the Hours of Mary of Burgundy », dans Reindert Falkenburg, Walter S. Melion et Todd M. Richardson (dir\teiHi[rend={italic}]{.), Image and Imagination of the Religious Self in Late Medieval and Early Modern Europe}, Turnhout, Brepols, 2007, p. 237-275. }}, et tante du futur Charles Quint. L’examen et la juxtaposition de ces deux manuscrits dans le cadre de cette étude permettront de considérer les implications médiatiques de l’héraldique dans un périmètre européen. Plus particulièrement, cette étude ambitionne de montrer comment le langage de l’héraldique contribue à rendre compte du climat politique et révèle des équilibres délicats, tout en respectant le décorum et en amusant le spectateur. Attavante se montre ici en artiste de cour du plus haut degré.}
            \teiP[xmlid={p3-3}]{En tant que cardinal, Jean de Médicis suit les pas de ses illustres ancêtres. Une fois devenu pape sous le nom de Léon X, il se permet petit à petit des digressions ludiques dans le décor de ses manuscrits. Dans son antiphonaire pour la chapelle Sixtine, on retrouve les symboles de la famille, notamment le \teiHi[rend={italic}]{diamante} et les plumes d’autruche, la rose des Orsini et le \teiHi[rend={italic}]{broncone}, agrémentés des symboles du pape.}
            \teiP[xmlid={p4-3}]{C’est dans le feuillet ouvrant l’office des morts que l’on peut relever une première intervention, un jeu autour du lion, symbole à la fois de Florence et du pape (fig. 1). Dans la partie supérieure d’une représentation des trois morts et des trois vivants, l’artiste insère deux lions phytomorphes, à gauche et à droite du \teiHi[rend={italic}]{diamante} contenant les \teiHi[rend={italic}]{palle }médicéennes. Les lions touchent le \teiHi[rend={italic}]{diamante} d’une patte et les clefs de saint Pierre de l’autre. Au-delà de se rapprocher de l’univers des grotesques, ce qui serait une sorte de retour à leur ville d’origine, ces créations phytomorphes évoquent également le thème de la régénération perpétuelle, renvoyant de cette manière au \teiHi[rend={italic}]{motto Semper} de Laurent le Magnifique\teiNote[n={3},place={foot},xmlid={ftn37}]{\teiP[xmlid={p-023}]{Gerhart Burian Ladner, « Vegetation Symbolism and the Concept of Renaissance », dans Millard Meiss (dir.), \teiHi[rend={italic}]{De artibus opuscula XL. Essays in Honor of Erwin Panofsky}, New York, New York University Press, 1961, vol. I, p. 303-322 ; Janet Cox-Rearick,\teiHi[rend={italic}]{ Dynasty and Destiny in Medici Art. Pontormo, Leo X and the two Cosimos}, Princeton, Princeton University Press, 1984. }}. À travers la figure emblématique du lion entre les grotesques et le \teiHi[rend={italic}]{diamante}, le pape se représente comme le maillon entre la famille florentine et la papauté-Rome\teiNote[n={4},place={foot},xmlid={ftn38}]{\teiP[xmlid={p-024}]{Giovanni met en image le conseil de son père. Son cousin, Giulio de’ Medici, choisira une autre variation pour son missel personnel. Voir Elli Doulkaridou-Ramantani, « Jules de Médicis et Pompeo Colonna. Enjeux héraldiques dans les missels de deux cardinaux antagonistes », dans Yvan Loskoutoff (dir.), \teiHi[rend={italic}]{Héraldique et papauté, Moyen Âge-Temps modernes}, Mont-Saint-Aignan, PURH, 2020, p. 197-233.}}.}
            \teiP[xmlid={p5-3}]{Par ailleurs, la forme circulaire établit une continuité avec les pages de garde des \teiHi[rend={italic}]{Plutei}. Le pape florentin confère un niveau supplémentaire à cette forme livresque, devenue ex-libris de la collection médicéenne, en ajoutant un cercle extérieur où l’on trouve des croix surmontant les huit symboles médicéens – le chiffre huit renvoie quant à lui à l’éternité, autre idée chère aux Médicis – qui semblent ainsi confirmer la bénédiction de la famille par le biais de Jean, élevé à la papauté.}
            \begin{teiDiv}[type={section1},xmlid={div01}]

              \teiHead[xmlid={le-pontifical-de-1520-001}]{Le pontifical de 1520}
              \teiP[xmlid={p6-3}]{Le manuscrit illustrant le mieux la logique cumulative des symboles médicéens et la signification de leur transposition sur le sol romain est le manuel de préparation pour la messe pontificale, décoré par Attavante degli Attavanti vers 1520\teiNote[n={5},place={foot},xmlid={ftn39}]{\teiP[xmlid={p-025}]{Pierpont Morgan Library, \teiHi[rend={italic}]{Praeparatio ad Missam Pontificalem}, ms. H.6. La version numérique du manuscrit est accessible sur la plateforme Corsair de la Pierpont Morgan Library de New York : \teiRef[target={http://corsair.themorgan.org/vwebv/holdingsInfo?bibId=76897}]{\teiHi[rend={underline}]{http://corsair.themorgan.org/vwebv/holdingsInfo?bibId=76897}} (consulté le 02 janvier 2024).}}. Ce pontifical contient – en dehors des armes et des \teiHi[rend={italic}]{imprese} du pape – tous les emblèmes et \teiHi[rend={italic}]{imprese} de la maison de Médicis, en remontant jusqu’à Cosme \teiHi[rend={italic}]{pater patriae }(fig. 2). Contrairement au Pontifical recueillant les textes pour les cérémonies officiées par l’évêque ou le pape, la \teiHi[rend={italic}]{Praeparatio }est habituellement un petit livre composé de quelques cahiers et contenant les textes et instructions concernant le parement ritualisé du pape avant la messe. Plusieurs manuscrits de ce type nous sont parvenus\teiNote[n={6},place={foot},xmlid={ftn40}]{\teiP[xmlid={p-026}]{Pierpont Morgan Library, ms. H.6 ; Harvard, Houghton Library, ms. Typ 136. Voir aussi Elena de Laurentiis, « \teiHi[rend={italic}]{Praeparatio ad missam pontificalem }Vincent Raymond di Lodève (illuminator) », dans Jeffrey F. Hamburger, William P. Stoneman, Anne-Marie Eze \teiHi[rend={italic}]{et al} (dir.), \teiHi[rend={italic}]{Beyond Words: Illuminated Manuscripts in Boston Collections}, catalogue d’exposition, Chestnut Hill, McMullen Museum of Art-Boston College, 2016, cat. 221, p. 276 et 277.}}, mais le décor de celui qui appartint à Léon X permet une série d’observations concernant la dimension performative de ce livre.}
              \teiP[xmlid={p7-3}]{Dans son ouvrage récent sur la musique durant le pontificat de Léon X, John Cummings parle de « lieu de performance » et fournit une description de l’itinéraire du pape. En effet, le pontife se préparait dans la \teiHi[rend={italic}]{Sala de’ Palafrenieri} ou la \teiHi[rend={italic}]{Camera del Papagallo}. Au même moment, les cardinaux se préparaient eux aussi dans la chambre adjacente, la \teiHi[rend={italic}]{Camera de’ Paramenti}. Ils rejoignaient ensuite le pape et commençaient la procession en traversant les salles du consistoire (à savoir la \teiHi[rend={italic}]{Sala Ducale} ou \teiHi[rend={italic}]{aula tertia} et \teiHi[rend={italic}]{secunda} et la \teiHi[rend={italic}]{Sala Regia} ou \teiHi[rend={italic}]{aula prima}) pour arriver à la chapelle Sixtine\teiNote[n={7},place={foot},xmlid={ftn41}]{\teiP[xmlid={p-027}]{Anthony M. Cummings, \teiHi[rend={italic}]{The Lion’s Ear: Pope Leo X, the Renaissance Papacy, and Music}, Ann Arbor, University of Michigan Press, 2012, p. 46 et fig. 6. Voir aussi Tom Campbell, « Pope Leo X’s Consistorial “letto de paramento” and the Boughton House Cartoons », \teiHi[rend={italic}]{The Burlington Magazine}, vol. 138, n\teiHi[rend={sup}]{o} 1120, 1996, p. 436-445, ici p. 439 ; Jonathan J. G. Alexander (dir.), \teiHi[rend={italic}]{The Painted Page: Italian Renaissance Book Illumination, 1450-1550}, catalogue d’exposition, Munich-New York, Prestel, 1994, cat. 4, p. 56-60. }}. Ces manuels contiennent les prières et hymnes que le pape récitait pendant la cérémonie de son habillage. Chaque pièce de vêtement composant l’ensemble du costume ecclésiastique papal était accompagnée d’une prière particulière. Ainsi, la totalité du manuscrit était feuilletée pendant la cérémonie, rendant sa décoration régulièrement accessible. Contrairement aux autres manuscrits de la collection médicéenne, l’exhaustivité choisie pour le manuscrit servant à la préparation non seulement physique, mais également spirituelle du pape, renvoie à une sorte de rituel auquel sont conviés les illustres ancêtres et membres contemporains occupant des postes stratégiques.}
              \teiP[xmlid={p8-3}]{Plus encore, il faut supposer que le livre était placé sur un pupitre, ou bien porté par un diacre pendant la cérémonie ; par conséquent, il aurait été visible par toutes les personnalités présentes, y compris les cardinaux. L’intimité de ces lieux, combinée à l’usage régulier du manuscrit, garantissait la présence médicéenne dynastique dans l’esprit de la curie.}
              \teiP[xmlid={p9-2}]{La revendication de l’origine florentine est explicite dès les premiers feuillets du manuscrit : saint Jean Baptiste, saint Laurent, Cosme et Damien, évangélistes et docteurs de l’Église, animent les marges. Laurent le Magnifique est également présent grâce à une médaille imitant les profils impériaux romains. Les inscriptions ne laissent aucune place au doute. En passant des marges au centre de la page, nous nous trouvons dans la \teiHi[rend={italic}]{camera del papagalo} aux appartements du Vatican. Le plafond à caissons et les draps affichent les \teiHi[rend={italic}]{palle} et le \teiHi[rend={italic}]{diamante} médicéens. Attavante ajoute la date d’exécution et l’inscription \teiHi[rend={italic}]{Leo X patria florent pont max} à la base du trône papal. La jonction entre la ville natale et celle de son règne est accomplie ; dans le registre supérieur, les cieux s’ouvrent pour laisser voir le Christ et Dieu le Père bénissant le pape, entourés des chérubins.}
              \teiP[xmlid={p10-2}]{Un élément intéressant est la reprise des \teiHi[rend={italic}]{imprese} associées à Julien de Médicis (1479-1516), ainsi qu’à Pierre (1472-1503), les deux frères du pape. Le premier fut nommé gonfalonier de l’Église en 1513, lors d’une somptueuse cérémonie au Capitole\teiNote[n={8},place={foot},xmlid={ftn42}]{\teiP[xmlid={p-028}]{Voir Anthony M. Cummings, \teiHi[rend={italic}]{The Politicized Muse: Music for Medici Festivals, 1512-1537}, Princeton, Princeton University Press, 1992. Pour la présence des \teiHi[rend={italic}]{imprese }dans les manuscrits de la bibliothèque médicéenne, voir surtout Anna Lenzuni (dir.), \teiHi[rend={italic}]{All’ombra del lauro. Documenti librari della cultura in età laurenziana}, catalogue d’exposition, Cinisello Balsamo, Silvana Editoriale, 1992.}}. Les deux roues, ainsi que l’\teiHi[rend={italic}]{impresa} GLOVIS, renvoient surtout à Julien, assassiné en 1516 ; les insignes de Pierre, le soleil et la pluie, le luth et le monogramme IR servent quant à eux à consolider l’importance de la famille, mais surtout à rappeler que la disparition de ce dernier a permis le retour des Médicis à Florence en 1512, événement qui marqua justement le tournant de leur fortune dynastique\teiNote[n={9},place={foot},xmlid={ftn43}]{\teiP[xmlid={p-029}]{Erkinger Schwarzenberg, « Glovis, impresa di Giuliano de’ Medici », \teiHi[rend={italic}]{Mitteilungen des Kunsthistorischen Institutes in Florenz}, n\teiHi[rend={sup}]{o} 39, 1995, p. 140-166, ici p. 155: «\teiHi[rend={italic}]{ Con la morte di Piero, nel 1503, scomparve l’ostacolo maggiore al ritorno dei Medici a Firenze e sotto il papato di Giulio II non mancarono le occasioni che diedero a Giuliano la certezza che la ruota della Fortuna, rimasta incagliata per tanti anni, si sarebbe rimessa a girare.} »}}.}
              \teiP[xmlid={p11-2}]{La représentation d’un ourson se régalant du miel d’une ruche, jusqu’à présent passée sous silence, est plus difficile à interpréter (fig. 3). La scène est accompagnée de l’inscription « \teiHi[rend={italic}]{Uno orsacchio ghiotto sono sio mimpugno egle pur buono}\teiNote[n={10},place={foot},xmlid={ftn44}]{\teiP[xmlid={p-030}]{« Je suis un ourson rassasié, quand je m’attrape à quelque chose je ne lâche pas prise. » Transcription et interprétation de l’auteur. Aucune mention dans la bibliographie existante n’a été trouvée. Dans le catalogue de la Pierpont Morgan Library, l’inscription n’est pas accompagnée d’interprétation ou de commentaire. }}\teiHi[rend={italic}]{. }» Il est certes tentant d’y voir une allusion au patronyme de la mère du pape, Clarice Orsini (1453-1488). En effet, les décors après le retour des Médicis soulignent souvent l’alliance avec les Orsini, le jeu de mots Orsini/Orsa/Ursa étant monnaie courante\teiNote[n={11},place={foot},xmlid={ftn45}]{\teiP[xmlid={p-031}]{Voir surtout Marica Tacconi, « Appropriating the Instruments of Worship: The 1512 Medici Restoration and the Florentine Cathedral », \teiHi[rend={italic}]{Renaissance Quarterly}, vol. 56, n\teiHi[rend={sup}]{o} 2, 2003, p. 333-376, ici p. 348.}}. Dans notre cas, l’ourson accompagné du \teiHi[rend={italic}]{motto} pourrait constituer un jeu sur l’\teiHi[rend={italic}]{impresa} aux pots de miel et abeilles de son père, en faisant allusion au petit ourson Jean et à sa prédilection pour les plaisirs de la table\teiNote[n={12},place={foot},xmlid={ftn46}]{\teiP[xmlid={p-032}]{Anthony M. Cummings, \teiHi[rend={italic}]{The Lion’s Ear}, \teiHi[rend={italic}]{op. cit.}, p. 114 ; Jeremy Kruse, « Hunting, Magnificence and the Court of Leo X », \teiHi[rend={italic}]{Renaissance Studies}, vol. 7, n\teiHi[rend={sup}]{o} 3, 1993, p. 243‑257.}}.}
              \teiP[xmlid={p12-2}]{Mais toutes ces observations ne suffisent pas pour expliquer cette image à la fois originale et incongrue eu égard à la date d’exécution – 1520 –, au contexte d’usage du manuscrit et au caractère autodérisoire d’un ourson qui se présente comme étant glouton. Ce que nous avons ici est à mon avis une double ruse.}
              \teiP[xmlid={p13-2}]{L’artiste et son commanditaire adoptent la stratégie médicéenne ancestrale d’appropriation d’éléments communs et bien connus. D’où puisent-ils les formes archétypales ? Une première réponse évidente serait l’univers héraldique des Orsini ; non seulement autour de la mère Clarice, à laquelle renvoient également et avec plus d’élégance les roses, mais de la famille en général, maintenant que Jean est à Rome et doit valoriser son ascendance romaine.}
              \teiP[xmlid={p14}]{Toutefois, ce n’est pas un simple ours héraldique comme c’est le cas dans les manuscrits des Orsini ou même dans un livre de chœur de Santa Maria del Fiore, où Clarice prend la forme d’une ourse parmi les lions symbolisant les Médicis\teiNote[n={13},place={foot},xmlid={ftn47}]{\teiP[xmlid={p-033}]{Marica Tacconi, art. cité.}}. Ni une référence à la constellation de la Grande Ourse comme c’était le cas dans les festivités pour l’investiture de Julien en 1513\teiNote[n={14},place={foot},xmlid={ftn48}]{\teiP[xmlid={p-034}]{Paolo Palliolo, \teiHi[rend={italic}]{Le feste pel conferimento del patriziato romano a Giuliano e Lorenzo de’ Medici}, Bologne, G. Romagnoli, 1885, p. 114 et suiv.}}. Ce que l’ourson souligne est le lien de parenté. Dans la littérature médiévale et contemporaine – de Dante aux \teiHi[rend={italic}]{Vaticinia des Summis pontificibus} et de poèmes d’Antonio Nerli ou de Teofilo Folengo, la locution « \teiHi[rend={italic}]{figlio de l’orso/orsa }» est amplement utilisée pour signifier un membre de la famille des Orsini\teiNote[n={15},place={foot},xmlid={ftn49}]{\teiP[xmlid={p-035}]{Agostino Paravicini Bagliani, \teiHi[rend={italic}]{Il bestiario del papa}, Turin, Einaudi, 2016, p. 249 et suiv. Voir aussi Michel Pastoureau, \teiHi[rend={italic}]{L’ours. L’histoire d’un roi déchu}, Paris, Seuil, 2016, p. 113.}}. Quant au terme « \teiHi[rend={italic}]{ghiotto }», le dictionnaire de l’Accademia della Crusca atteste de son association répandue à l’ours friand du miel\teiNote[n={16},place={foot},xmlid={ftn50}]{\teiP[xmlid={p-036}]{\teiHi[rend={italic}]{Grande Dizionario della lingua italiana}, Turin, UTET, 1972, t. VI, p. 742 ; pour les textes de la période voir par exemple Sabba da Castiglione, \teiHi[rend={italic}]{Ricordi di monsignor Sabba da Castiglione caualier Gierosolimitano, di nuouo corretti, et ristampati. Con vna tauola copiosissima nuouamente aggiunta. Et appresso breuemente e descritta la vita dell’auttore}, In Venetia, Per Paolo Gerardo, 1560, p. 133 : « \teiHi[rend={italic}]{più ghiotto che l’orso del miele} ». Sur le poème « \teiHi[rend={italic}]{Le api} » de Giovanni Rucellai et le poème \teiHi[rend={italic}]{Caos del Triperuno de Teofilo Folengo}, voir William Roscoe, \teiHi[rend={italic}]{Life and Pontificate of Leo X}, Londres, Henry G. Bohn, 1846, 2 vol., t. II, p. 163 et 171.}}. Nous avons donc une première base littéraire qui couvre plusieurs siècles et traditions.}
              \teiP[xmlid={p15}]{Tournons-nous maintenant vers l’univers matériel médicéen lors des années de formation des fils du Magnifique : l’Ésope des Médicis affirme l’emploi de livres illustrés pour l’éducation des enfants\teiNote[n={17},place={foot},xmlid={ftn51}]{\teiP[xmlid={p-037}]{New York Public Library, ms Spencer 50, \teiRef[target={https://digitalcollections.nypl.org/collections/fables\#/?tab=about}]{https://digitalcollections.nypl.org/collections/fables\#/?tab=about} (consulté le 6 avril 2023).}}. En ce qui concerne plus particulièrement la fable de l’ours et du miel associée au vice de la colère, il fut connu à travers le \teiHi[rend={italic}]{Fior di virtù}, un traité populaire d’histoires moralisatrices dont les exemplaires circulaient dans les bibliothèques de l’élite européenne\teiNote[n={18},place={foot},xmlid={ftn52}]{\teiP[xmlid={p-038}]{Curt F. Bühler, « Studies in the Early Editions of the “Fiore Di Virtù” », \teiHi[rend={italic}]{The Papers of the Bibliographical Society of America, }vol. 49, n\teiHi[rend={sup}]{o} 4, 1955, p. 315-339 ; Federico Botana, \teiHi[rend={italic}]{Learning through Images in the Italian Renaissance: Illustrated Manuscripts and Education in Quattrocento Florence}, Cambridge, Cambridge University Press, 2020, p. 58‑79, https://doi.org/10.1017/9781108867313.004.}}.}
              \teiP[xmlid={p16}]{La source textuelle apparaît dans les collections médicéennes en deux exemplaires, qui sont pourtant dépourvus d’illustrations\teiNote[n={19},place={foot},xmlid={ftn53}]{\teiP[xmlid={p-039}]{Biblioteca Medicea Laurenziana, Plutei 41.3, 76.63. }}. Toutefois, un exemplaire contemporain illustré par Mariano del Buono, artiste employé par Laurent le Magnifique, conservé aujourd’hui à la Biblioteca Riccardiana de Florence\teiNote[n={20},place={foot},xmlid={ftn54}]{\teiP[xmlid={p-040}]{Biblioteca Riccardiana, ms. 1774.}} pourrait constituer une source d’inspiration pour Attavante degli Attavanti (fig. 4). Les deux exemplaires de la Riccardiana présentent des illustrations où la ruche est similaire, tandis qu’une autre copie contemporaine du texte contenant des illustrations moins élaborées pourrait servir de modèle pour l’ourson\teiNote[n={21},place={foot},xmlid={ftn55}]{\teiP[xmlid={p-041}]{Biblioteca Riccardiana, ms. 1763, f. 13r.}} (fig. 5). La source iconographique étant associée aux manuels destinés à l’éducation des enfants, elle enrichit notre interprétation en lui ajoutant une dimension proprement liée à l’univers de l’enfance.}
              \teiP[xmlid={p17}]{Nous avons donc une image qui serait reconnue et reliée, dans un premier temps grâce à la culture du spectateur, à une représentation moralisatrice du vice de la colère. Sauf que Léon X n’était point colérique et, qui plus est, l’image nous présente un ourson tout à fait content et indifférent aux abeilles.}
              \teiP[xmlid={p18}]{Disons d’emblée que l’esprit ludique se trouve au cœur d’une image qui, au sein de la page décorée, prétend être une devise comme ses voisines. Les éléments formels qui participent à cette opération de tromperie du regard sont les couleurs, le cadre et surtout le ruban donnant l’impression de contenir le \teiHi[rend={italic}]{motto} de la devise. L’iconographie aide aussi à fondre l’image dans le tissu décoratif de l’ensemble du manuscrit où l’on trouve d’autres animaux sans rapport héraldique avec le pape. Mais contrairement aux adverbes et phrases à la troisième personne, neutres comme le demande le caractère d’un proverbe, l’opération est ici différente : c’est l’ourson qui parle à la première personne et informe le lecteur qu’il est un ourson bien glouton et que s’il s’attrape il ne laissera pas prise. Il s’agit d’une fausse \teiHi[rend={italic}]{impresa} de caractère autodérisoire qui, de plus, renvoie à une fable dont le sujet n’est pas à propos : nous avons à faire à une double ruse, picturale/héraldique d’un côté, textuelle/littéraire de l’autre.}
              \teiP[xmlid={p19}]{Comment pourrait-on expliquer cette ruse qui exploite le langage héraldique ? Afin de répondre à cette question, je dois ouvrir l’analyse au deuxième groupe de manuscrits, à savoir les chants polyphoniques provenant de l’atelier de Petrus Alamire.}
            
\end{teiDiv}
            \begin{teiDiv}[type={section1},xmlid={div02}]

              \teiHead[xmlid={le-cadeau-de-lempire-001}]{Le cadeau de l’Empire}
              \teiP[xmlid={p20}]{La riche personnalité de Léon X et la conjoncture politique chargée de la deuxième décennie du \teiHi[rend={small-caps}]{xvi}\teiHi[rend={sup}]{e} siècle l’ont amené à nouer des relations avec les cours royales qui, à l’époque, se disputaient le contrôle du continent. Parmi les cadeaux diplomatiques qu’il a reçus, on trouve trois livres de chants polyphoniques provenant du scriptorium de Petrus Alamire, scribe, courtisan et espion d’Henri VIII d’Angleterre\teiNote[n={22},place={foot},xmlid={ftn56}]{\teiP[xmlid={p-042}]{BAV, Capp.Sist.160, \teiRef[target={https://digi.vatlib.it/view/MSS\_Capp.Sist.160}]{https://digi.vatlib.it/view/MSS\_Capp.Sist.160} (consulté le 03 mars 2025).}}. Les manuscrits ont été offerts à Léon X par Marguerite d’Autriche, tante de Charles Quint et à l’époque régente des Pays-Bas. La recherche n’a pas encore tranché sur les dates précises de la remise, possible entre 1513 et 1519. Leur présence dans le fonds de la chapelle Sixtine permet au moins deux conclusions : le cadeau fut accepté et les manuscrits ont été utilisés par les chantres\teiNote[n={23},place={foot},xmlid={ftn57}]{\teiP[xmlid={p-043}]{Sur le rôle de Marguerite d’Autriche dans la conception des livres de chœur de l’atelier d’Alamire, voir Bonnie Blackburn, « Messages in Miniature: Pictorial Programme and Theological Implications in the Alamire Choirbooks », dans \teiHi[rend={italic}]{The Burgundian-Habsburg Court Complex of Music Manuscripts (1500-1535) and the Workshop of Petrus Alamire: Colloquium Proceedings, Leuven, 25-28 November 1999}, Louvain-Neerpelt, Alamire Foundation, 2003, p. 161-184. }}.}
              \teiP[xmlid={p21}]{Les manuscrits d’Alamire, une cinquantaine à ce jour, étaient souvent le joyau des collections royales et réunissaient à un moment donné les acteurs de toute l’Europe. De l’Angleterre jusqu’au Vatican, en passant par Maximilien I\teiHi[rend={sup}]{er} et Fréderic III de Saxe, le regard de tous a croisé les chants et leur décor exceptionnel\teiNote[n={24},place={foot},xmlid={ftn58}]{\teiP[xmlid={p-044}]{Anne Margreet W. As-Vijvers, « The Grotesque Initials in the Alamire Choirbooks », \teiHi[rend={italic}]{Journal of the Alamire Foundation}, vol. 11, n\teiHi[rend={sup}]{o} 1-2, 2019, p. 13‑46 ; Klaas van der Heide, « How Many Paths Must a Choirbook Tread Before it Reaches the Pope? », \teiHi[rend={italic}]{ibid.}, p. 47‑70. }}. Quel est leur trait caractéristique ? La présence d’initiales grotesques au début de chaque voix (\teiHi[rend={italic}]{bassus, contrabassus, tenor, contratenor}). Il s’agit de conceptions hybrides et humoristiques de style flamand. Ces têtes humaines phytomorphes très colorées, habillées et même déguisées avec des coiffures extravagantes, mettent en place une cour virtuelle composée de bouffons. Leur présence ouvre sur un espace d’expression plus large au sein duquel on a le droit de se moquer des autres sans pénalité. Par ailleurs, l’association du bouffon à la mécanique pneumatique de son corps établit un lien entre les chantres et leur activité\teiNote[n={25},place={foot},xmlid={ftn59}]{\teiP[xmlid={p-045}]{Benjamin Dang, « Le bouffon, un héraut du vide », \teiHi[rend={italic}]{Les chantiers de la création. Revue pluridisciplinaire en lettres, langues, arts et civilisations}, n\teiHi[rend={sup}]{o} 13, 12 juillet 2021, https://doi.org/10.4000/lcc.4119.}}. Enfin, le bouffon comme contre-exemple peut encourager les associations théologiques. D’une allusion généraliste aux écrits de saint Paul\teiNote[n={26},place={foot},xmlid={ftn60}]{\teiP[xmlid={p-046}]{Surtout 1 Corinthiens 3:18. }}, on peut facilement arriver à Martin Luther, qui en 1520 se présente comme un bouffon dans son manifeste \teiHi[rend={italic}]{À la noblesse chrétienne de la nation allemande}\teiNote[n={27},place={foot},xmlid={ftn61}]{\teiP[xmlid={p-047}]{Hub Zwart, \teiHi[rend={italic}]{Ethical Consensus and the Truth of Laughter: The Structure of Moral Transformations}, Kampen, Kok Pharos, 1996, p. 156.}}.}
              \teiP[xmlid={p22}]{Après une étude comparative des initiales du corpus d’Alamire, et malgré les motifs répétitifs et donc partagés entre les différents possesseurs, il semblerait que le décor ait été repensé pour Léon X. Les initiales sont plus colorées et les figures sont encore plus déguisées et même maquillées\teiNote[n={28},place={foot},xmlid={ftn62}]{\teiP[xmlid={p-048}]{BAV, Capp.Sist.160, \teiRef[target={https://digi.vatlib.it/view/MSS\_Capp.Sist.160/0199}]{https://digi.vatlib.it/view/MSS\_Capp.Sist.160/0199} (consulté le 18 décembre 2024).}}. Si le trait commun de ces figures est l’incapacité sensorielle, les têtes de Léon X sont encore plus handicapées. Elles sont aveuglées par les grands chapeaux – souvent des chapeaux de cardinal. Elles sont aussi bâillonnées. Quand on compare ces initiales avec les initiales du manuscrit ayant appartenu à Maximilien I\teiHi[rend={sup}]{er}, le contraste est parlant. Ce qui permet de penser que l’humour et son goût pour la frivolité étaient non seulement connus, mais dans ce cas-là exploités.}
              \teiP[xmlid={p23}]{Mais pour revenir à l’héraldique, ces manuscrits mettent en lumière le revers de cette pratique de prestige et de pouvoir, à savoir les obligations qui en découlent. Dans le cas des Médicis, une famille des banquiers qui cultiva avec assiduité et persévérance son image, complètement factice en fin de compte, l’insertion de leurs armes au sein du décor d’un cadeau royal constitue en quelque sorte leur légitimation en dehors de l’Italie en même temps que leur obligation à se plier aux conjonctures, ou du moins un rappel diplomatique de leur appartenance à un plus vaste réseau d’intérêts\teiNote[n={29},place={foot},xmlid={ftn63}]{\teiP[xmlid={p-049}]{Caroline Elam, « Art and Diplomacy in Renaissance Florence », \teiHi[rend={italic}]{RSA Journal}, vol. 136, n\teiHi[rend={sup}]{o} 5387, 1988, p. 813‑826, ici p. 814.}}.}
              \teiP[xmlid={p24}]{Les précédents travaux d’historiens de la musique et des musicologues ont mis en lumière l’agenda théologique lié au dogme de l’Immaculée Conception et à la réforme de l’Église, insinuée surtout à travers le motet intitulé \teiHi[rend={italic}]{L’Homme armé} et renvoyant clairement à l’empereur. Pour ma part, je vais combiner quelques éléments et proposer de nouvelles observations en ce qui concerne l’emploi et la signification de l’héraldique de Léon X.}
              \teiP[xmlid={p25}]{Le cadeau comporte trois volumes\teiNote[n={30},place={foot},xmlid={ftn64}]{\teiP[xmlid={p-050}]{BAV, Capp.Sist.160, \teiRef[target={https://digi.vatlib.it/mss/detail/223173}]{https://digi.vatlib.it/mss/detail/223173} (consulté le 18 décembre 2024) ; Capp.Sist.34 ; Capp.Sist.36. }}. Premièrement, l’un d’entre eux n’était à l’origine pas destiné au pape, mais à un conseiller de Marguerite d’Autriche, Jean III de Glymes, seigneur de Bergen op Zoom. Son prénom identique à celui du pape permettait ce recyclage. Cela veut dire que les armes ont été peintes par-dessus celles du conseiller, ce qui est encore visible aujourd’hui\teiNote[n={31},place={foot},xmlid={ftn65}]{\teiP[xmlid={p-051}]{Klaas van der Heide, « How Many Paths Must a Choirbook Tread Before it Reaches the Pope? », art. cité.}}.}
              \teiP[xmlid={p26}]{Deuxièmement : la taille des armes papales. Quand on les compare aux autres manuscrits d’Alamire, on se rend vite compte que leur échelle ne correspond pas au statut papal. En effet, dans les manuscrits d’Alamire, les familles royales des Habsbourg, mais aussi d’Angleterre, sont mises en valeur par le biais de portraits et l’insertion d’un appareil héraldique de proportions impressionnantes (fig. 6 et 7). Nous pourrions même nous demander si les dimensions des armes médicéennes ne cachent pas un certain mépris que l’on pourrait expliquer en utilisant des arguments politiques et surtout la mentalité césaropapiste traversant la politique de Maximilien I\teiHi[rend={sup}]{er}. C’est un manuscrit conservé aujourd’hui à Mechelen qui permet de prendre la mesure des ambitions de Maximilien eu égard à la réforme de l’Église : il se présente en tant que chef temporel suprême, dont le couronnement est assisté par le pape\teiNote[n={32},place={foot},xmlid={ftn66}]{\teiP[xmlid={p-052}]{Vincenzo Borghetti, « Music and the Representation of Princely Power in the Fifteenth and Sixteenth Century », \teiHi[rend={italic}]{Acta Musicologica, }vol. 80, n\teiHi[rend={sup}]{o} 2, 2008, p. 179‑214.}}.}
              \teiP[xmlid={p27}]{Comme le tableau de Holbein le met clairement en relief, la musique polyphonique appartenait avec la cartographie et la cosmographie à l’arsenal des cadeaux diplomatiques de l’époque. La musique polyphonique flamande et le décor profondément héraldique et propagandiste devaient infiltrer le Vatican avec ou sans l’empereur. D’aucuns ont proposé d’interpréter le choix comme une suggestion au pape de tempérer sa préférence pour les compositeurs français. Pour ma part, je pense que l’enjeu est plus directement lié au pouvoir. Si les manuscrits partagés entre les possesseurs visaient à consolider leurs liens à travers des références culturelles et théologiques communes, la mise en place d’une cour royale virtuelle, par le biais de bouffons ornant les initiales des chants, permettait aussi la présence de l’Empire au sein de la chapelle Sixtine. En faisant appel au goût de Léon X pour la musique et les bouffons, tout en lui présentant des répertoires inédits et innovants, le langage codé et hiérarchique de l’héraldique, où l’échelle et les proportions sont significatives, indique que derrière la flatterie résidaient les objectifs et les principes impériaux qui, quelques années plus tard, se matérialiseront de manière dramatique avec le sac de Rome et le couronnement de Charles Quint par Clément VII à Bologne.}
            
\end{teiDiv}
            \begin{teiDiv}[type={section1},xmlid={div03}]

              \teiHead[xmlid={conclusions-001}]{Conclusions}
              \teiP[xmlid={p28}]{Ce serait un euphémisme de dire que le pape était sous pression en 1520. Deux ans plus tôt, il avait survécu à une tentative d’empoisonnement\teiNote[n={33},place={foot},xmlid={ftn67}]{\teiP[xmlid={p-053}]{Ludwig von Pastor, \teiHi[rend={italic}]{The History of the Popes: From the Close of the Middle Ages, Drawn From the Secret Archives of the Vatican and Other Original Sources}, R. F. Kerr (éd.), Londres, Kegan Paul, Trench, Trübner \& co, 1910, vol. 8, p. 184. }}. En dehors de la disparition d’Alfonsina Orsini, facteur de stabilité à Florence et proche de Léon X, le pape devait également agir contre Martin Luther avec la bulle \teiHi[rend={italic}]{Exsurge Domine}. C’était au même moment, pour le continent, une période de grands changements dans la répartition du pouvoir. À l’est, Suleyman le Magnifique accédait au trône de l’empire ottoman en même temps que Charles Quint prenait les rênes à l’ouest, alors que François I\teiHi[rend={sup}]{er} et Henry VIII d’Angleterre s’efforçaient de dissiper les tensions avec la rencontre du camp du Drap d’or.}
              \teiP[xmlid={p29}]{La confection de la \teiHi[rend={italic}]{Praeparatio}, un livre contenant les paroles capables de le transformer en \teiHi[rend={italic}]{princeps apostolorum}, peut être comprise comme un geste de réaffirmation de la part du pape. Il choisit l’enlumineur médicéen par excellence, Attavante degli Attavanti, désormais célèbre, et lui demande d’accentuer la présence médicéenne en tête de l’Église catholique en agrémentant les marges avec tous les symboles des membres de sa famille. Au sein de ce réseau figuratif chargé, Léon dépose son amour pour la musique et les facéties – le livre contient plusieurs vignettes où l’on voit des singes jouant de la musique ou imitant les gestes de parement du pape. En dernier lieu, il incorpore une image entièrement affective : l’ourson.}
              \teiP[xmlid={p30}]{Ce dernier passe facilement inaperçu grâce à un véritable bestiaire réparti dans les marges du manuscrit par Attavante. Il arrive ainsi à le cacher en pleine vue parmi les animaux héraldiques – le lion, la licorne – et d’autres d’une symbolique plus généralement liée aux marges des manuscrits enluminés, comme les singes, les satyres ou encore les hippocampes. En parlant à la première personne, l’ourson personnifie le cardinal qui arriva malade au conclave de 1513 et dont l’élection fut favorisée par des spéculations tenant compte de son mauvais état de santé. Sept ans plus tard, le \teiHi[rend={italic}]{figlio de l’Orsa }rappelle aux cardinaux qui l’entourent dans ses appartements privés – y compris ceux des autres pays de l’Europe et par extension à leurs princes, symbolisés par les abeilles\teiNote[n={34},place={foot},xmlid={ftn68}]{\teiP[xmlid={p-054}]{Waldemar Deonna, « L’abeille et le roi », \teiHi[rend={italic}]{Revue belge d’archéologie et d’histoire de l’art}, vol. 25, 1956, p. 105‑131. }} – que c’est bien lui le chef solide de l’Église catholique. Malheureusement les conjonctures le démentiront un an plus tard.}
              \teiP[xmlid={p31}]{Bien que ludiques, ces constructions dévoilent des intentions sérieuses, leur fonctionnement anxiolytique pointe justement l’état d’âme d’un pontife qui subit d’énormes pressions et cherche à se rassurer en faisant appel à ces ancêtres et plus précisément à la présence maternelle. Mais plus encore, il opte pour l’humour curatif, thérapeutique que l’on rencontre également dans un format monumental sous les traits du nain exposant son testicule au spectateur dans la \teiHi[rend={italic}]{Sala di Costantino}, sans parler des décors de la Villa médicéenne à Poggio a Caiano\teiNote[n={35},place={foot},xmlid={ftn69}]{\teiP[xmlid={p-055}]{Voir surtout Robin O’Bryan, « The Medici Pope, Curative Puns, and the Panacean Dwarf in the Sala di Costantino », \teiHi[rend={italic}]{Secac Review}, vol. 16, n\teiHi[rend={sup}]{o} 5, p. 590-607.}}. L’artiste et le commanditaire exploitent la pompe figurative du langage héraldique pour dissimuler un élément affectif et personnel qui rappelle la caresse tendre que le pape bibliophile accorde à la Bible dans son portrait de Raphaël.}
              \teiP[xmlid={p32}]{Cette mise en perspective des manuscrits polyphoniques d’Alamire et de la commande papale successive d’un livre tout autant intime que d’apparat met en relief la relativité et la polyvalence du signe héraldique. Entre instrument de présentation et moyen de projection d’un soi politique, il peut aussi revêtir des fonctions psychologiques et apotropaïques tout en respectant les codes de communication des cours et des réseaux diplomatiques.}
              \teiP[xmlid={p33}]{De ce point de vue, étudier l’héraldique dans les livres enluminés ayant appartenu aux papes s’avère un instrument précieux pour obtenir des résultats plus intimes, plus précis, sur les équilibres politiques des réseaux de pouvoir de la période.}
              \begin{teiFigure}[xmlid={figure01}]

                \teiGraphic[url={../icono/br/Ch02\_Doulkaridou/fig1.jpg}]
                \teiHead[xmlid={figure-1-antiphonaire-de-leon-x-rome-bav-ms-capp-sist-10-amico-aspertini-fol-66v-detail-001}]{Figure 1. \teiHi[rend={italic}]{Antiphonaire de Léon X}, Rome, BAV, ms. Capp.Sist.10 : Amico Aspertini, fol. 66v (détail).}
              
\end{teiFigure}
              \begin{teiFigure}[xmlid={figure02}]

                \teiGraphic[url={../icono/br/Ch02\_Doulkaridou/fig2.jpg}]
                \teiHead[xmlid={figure-2-praeparatio-ad-missam-pontificalem-new-york-the-pierpont-morgan-library-ms-h-6-attavante-degli-attavanti-fol-iv-iir-001}]{Figure 2. \teiHi[rend={italic}]{Praeparatio ad missam pontificalem}, New York, The Pierpont Morgan Library, ms. H.6 : Attavante degli Attavanti, fol. Iv-IIr.}
              
\end{teiFigure}
              \begin{teiFigure}[xmlid={figure03}]

                \teiGraphic[url={../icono/br/Ch02\_Doulkaridou/fig3.jpg}]
                \teiHead[xmlid={figure-3-praeparatio-ad-missam-pontificalem-new-york-the-pierpont-morgan-library-ms-h-6-attavante-degli-attavanti-fol-7v-001}]{Figure 3.\teiHi[rend={italic}]{ Praeparatio ad missam pontificalem}, New York, The Pierpont Morgan Library, ms. H.6 : Attavante degli Attavanti, fol. 7v.}
              
\end{teiFigure}
              \begin{teiFigure}[xmlid={figure04}]

                \teiGraphic[url={../icono/br/Ch02\_Doulkaridou/fig4\%20copie.jpg}]
                \teiHead[xmlid={figure-4-fiore-di-virtu-florence-biblioteca-riccardiana-ms-ricc-1774-mariano-del-buono-fol-18r-001}]{Figure 4. \teiHi[rend={italic}]{Fiore di Virtù}, Florence, Biblioteca Riccardiana, ms. Ricc. 1774 : Mariano del Buono, fol. 18r.}
              
\end{teiFigure}
              \begin{teiFigure}[xmlid={figure05}]

                \teiGraphic[url={../icono/br/Ch02\_Doulkaridou/fig5.jpg}]
                \teiHead[xmlid={figure-5-fior-di-virtu-florence-biblioteca-riccardiana-ms-ricc-1763-giovanni-di-gregorio-dantonio-ghinghi-copiste-fol-13r-001}]{Figure 5. \teiHi[rend={italic}]{Fior di Virtù} Florence, Biblioteca Riccardiana, ms. Ricc.1763 : Giovanni di Gregorio d’Antonio Ghinghi (copiste), fol. 13r.}
              
\end{teiFigure}
              \begin{teiFigure}[xmlid={figure06}]

                \teiGraphic[url={../icono/br/Ch02\_Doulkaridou/fig8a.jpg}]
                \teiHead[xmlid={figure-6-motets-jena-thuringer-universitats-und-landesbibliothek-ms-2-fol-2r-detail-001}]{Figure 6. \teiHi[rend={italic}]{Motets}, Jena, Thüringer Universitäts- und Landesbibliothek, ms. 2, fol. 2r (détail).}
              
\end{teiFigure}
              \begin{teiFigure}[xmlid={figure07}]

                \teiGraphic[url={../icono/br/Ch02\_Doulkaridou/fig7.jpg}]
                \teiHead[xmlid={figure-7-motets-atelier-de-petrus-alamire-rome-bav-ms-capp-sist-36-atelier-de-petrus-alamire-fol-66v-detail-001}]{Figure 7 : \teiHi[rend={italic}]{Motets}, atelier de Petrus Alamire, Rome, BAV, ms. Capp.Sist.36 : Atelier de Petrus Alamire, fol. 66v (détail).}
              
\end{teiFigure}
            
\end{teiDiv}
          
\end{teiElement}
        
\end{teiElement}
        \begin{teiElement}[name={group},type={article},data-page-title={L’héraldique de Jules III Ciocchi del Monte (1550-1555) dans l’ornement pour le livre},data-page-authors={Yvan Loskoutoff},data-include-href={Ch03\_Heraldique\_JulesIII.xml},xmlid={article-002}]

          \begin{teiElement}[name={front}]

            \begin{teiDiv}[type={titlePage},xmlid={titlepage-005}]

              \teiP[rend={title-main},xmlid={p-056}]{L’héraldique de Jules III Ciocchi del Monte (1550-1555) dans l’ornement pour le livre}
              \teiP[rend={author-aut},xmlid={p-057}]{Yvan Loskoutoff}
            
\end{teiDiv}
          
\end{teiElement}
          \begin{teiElement}[name={body},xmlid={body-005}]

            \teiP[xmlid={p1-5}]{Le livre constitue un support privilégié du blason dans sa reliure comme dans son décor intérieur peint pour les manuscrits et gravé pour les imprimés\teiNote[n={1},place={foot},xmlid={ftn1-3}]{\teiP[xmlid={p-058}]{Voir par exemple, Matthieu Desachy (dir.), \teiHi[rend={italic}]{L’héraldique et le livre}, catalogue d’exposition, Montauban-Albi-Toulouse, Paris, Somogy, 2002.}}. Les reliures aux armes du pape Jules III (1550-1555) ont été étudiées par Tammaro de Marinis, qui en a recensé une dizaine d’exemplaires\teiNote[n={2},place={foot},xmlid={ftn36-2}]{\teiP[xmlid={p-059}]{\teiHi[rend={italic}]{La legatura artistica in Italia nei secoli XV e XVI, Notizie ed elenchi}, Florence, Fratelli Alinari, 1960, 3 vol., t. 1, n\teiHi[rend={sup}]{o} 871, 877, 879-884, 906. Voir aussi \teiHi[rend={italic}]{Legature papali da Eugenio IV a Paolo VI}, catalogue d’exposition, Cité du Vatican, BAV, 1977, n\teiHi[rend={sup}]{o} 95-100.}}. Nous nous intéresserons ici au décor intérieur des manuscrits puis des imprimés afin d’examiner comment l’héraldique pontificale y est traitée.}
            \begin{teiDiv}[type={section1},xmlid={div01-2}]

              \teiHead[xmlid={les-armoiries-de-jules-iii-001}]{Les armoiries de Jules III}
              \teiP[xmlid={p2-4}]{Originaire de Monte San Savino près d’Arezzo, Jules III portait des armes parlantes, faisant allusion à son nom et à sa ville. Elles sont ainsi décrites dans le diplôme par lequel le pape autorisait les comtes Spada de Terni à les porter en chef : « d’azur à la bande de gueules bordée d’or chargée de trois monts de trois coupeaux de même, et accostée de deux couronnes de laurier aussi d’or\teiNote[n={3},place={foot},xmlid={ftn37-2}]{\teiP[xmlid={p-060}]{Ferruccio Pasini Frassoni, \teiHi[rend={italic}]{Essai d’armorial des papes}, Rome, Collège héraldique, 1906, p. 37. Cet auteur ne donne pas la localisation du document. Aux archives du Vatican, nous n’avons trouvé que deux concessions de ses propres armoiries par Jules III, elles accompagnent des créations de comtes palatins, l’une à Polidamas Maynus de Pavie (14 mai 1550), l’autre à Florianus Aurelius de Bologne (5 novembre 1550). Dans les deux cas, il est précisé que ces armes s’ajouteront à celles du récipiendaire, et pour le deuxième en chef (Archivio Apostolico Vaticano, Armadio XXXIX.59, fol. 267 et 444).}} ». Son oncle, le cardinal Antonio Maria del Monte, qui s’était élevé par ses talents de juriste et qui avait favorisé la carrière de son neveu, portait les mêmes, comme on peut le voir sur les pages de titre des recueils publiés à l’occasion des fêtes romaines de Pasquin dont il était le protecteur\teiNote[n={4},place={foot},xmlid={ftn38-2}]{\teiP[xmlid={p-061}]{Valerio Marucci, Antonio Mazzo, Angela Romano \teiHi[rend={italic}]{et al}. (éd.), \teiHi[rend={italic}]{Pasquinate romane del Cinquecento}, Rome, Salerno Editrice, 1983, 2 vol., t. 2, pl. XXXVI à XL. Voir aussi Yvan Loskoutoff, « Regrets et pasquinades. Le cycle de Jules III dans le recueil de Du Bellay », \teiHi[rend={italic}]{Revue d’histoire littéraire de la France}, vol. 123, n\teiHi[rend={sup}]{o} 1, 2023, p. 5-24.}}.}
              \teiP[xmlid={p3-4}]{Les Archives d’État de Rome conservent un témoignage sur l’histoire de ces armes, selon lequel Jules III serait intervenu pour les modifier :}
              \begin{teiCit}[xmlid={cit01}]

                \begin{teiQuote}
« Del Monte San Savino : d’azur à une bande de gueules bordée d’or, incluant trois monts de trois coupeaux chacun. Tel fut l’antique écu de cette maison qui prit son nom du château dit Monte S. Savino près d’Arezzo, appelé d’abord des Ciocchi. Furent ensuite ajoutés dans l’azur aux côtés de la bande deux rameaux de laurier recourbés en forme de couronne obtenus sous forme de concession par Fabiano del Monte fameux juriste. Ces armes furent finalement changées par le cardinal Gian Maria del Monte après qu’il fut élevé au pontificat sous le nom de Jules III, lequel, ayant entendu d’un certain astrologue que ces armes présageaient je ne sais quoi de funeste, changea la disposition des monts, de tombant qu’ils étaient il les mit sur pied et enlevant les deux rameaux de laurier il mit en leur place deux couronnes d’olivier en guirlandes\teiNote[n={5},place={foot},xmlid={ftn39-2}]{\teiP[xmlid={p-062}]{Archivio di Stato di Roma, Cartari Febei, reg. 164, f. 122, cité par Maria Giulia Aurigemma, \teiHi[rend={italic}]{Palazzo Firenze in Campo Marzio}, Rome, Istituto poligrafico e zecca dello Stato-Libreria dello Stato, 2007, p. 94, n. 70 : « \teiHi[rend={italic}]{Del Monte San Savino : D’azzurro con una banda di rosso profilata d’oro, cornice di tre mucchi di monti ciascuno. Questa fu l’arma antica di questa casa che prese il cognome dal Castello detto il Monte di S. Savino presso ad Arezzo, chiamato prima de’ Ciocchi. Furono poi aggiunti nell’azzurro à lati della banda due rami d’alloro ritorti in giro a modo di corone, ottenuti per concessione da Fabiano del Monte chiaro Giurista. Mutossi finalmente quest’arme dal Cardinal Gian Maria del Monte dopo ch’ei fu assunto al Pontificato col nome di Giulio terzo, il quale avendo inteso da un certo Astrologo, che questa insegna presagiva non sò che d’infausto, mutò la disposizione de’monti e da cadenti che erano gli pose tutti in piè, e togliendo i due rami d’alloro vi pose in lor vece due corone à ghirlande di olivo.} » }}. »
\end{teiQuote}
              
\end{teiCit}
              \teiP[xmlid={p4-4}]{La modification porte sur deux points : les monts de mauvais présage, car tombants qui sont redressés, illustrant la croyance aux pouvoirs auguraux du blason pontifical que nous avons étudiée ailleurs\teiNote[n={6},place={foot},xmlid={ftn40-2}]{\teiP[xmlid={p-063}]{Yvan Loskoutoff, « Augures héraldiques de la papauté », dans Andreas Gottsmann, Pierantonio Piatti et Andreas E. Rehberg (dir.), \teiHi[rend={italic}]{Incorrupta monumenta Ecclesiam defendunt, Studi offerti a Sergio Pagano, Prefetto dell’Archivio Segreto Vaticano}, Cité du Vatican, Archivio Segreto Vaticano, 2018, 3 vol., t. 2, p. 976-994.}} et les couronnes de laurier qui deviennent d’olivier, symbole de paix. Un témoignage imprimé légèrement postérieur à Jules III confirme que les couronnes furent effectivement modifiées : « le champ d’azur contient une bande de gueules chargée de trois monts d’or qui est accostée de deux couronnes de poètes ou impériales d’or dans ledit champ d’azur (changées ensuite par le pape Jules III en deux simples rameaux de laurier\teiNote[n={7},place={foot},xmlid={ftn41-2}]{\teiP[xmlid={p-064}]{Agostino Fortunio, \teiHi[rend={italic}]{Cronichetta del Monte San Savino di Toscana}, In Firenze, Nella Stamperia di Bartolomeo Sermartelli, 1583, p. 38 : « \teiHi[rend={italic}]{il campo azurro contiene per compartimento vna sbarra rossa con tre monti d’oro, che rimane attrauersata da due Corone poetali, o imperiali doro in detto campo azurro (divertite poi da Papa Giulio III. in due semplici Ramoscelli di lauro).} »}}). »}
              \teiP[xmlid={p5-4}]{Les deux textes s’accordent sur le fait qu’une modification fut bien voulue par Jules III, mais ils divergent sur sa nature. Quoi qu’il en soit, cela n’eut pas d’incidence notable sur la représentation des armes papales. Il est généralement difficile de discerner si le feuillage des couronnes est constitué de laurier ou d’olivier. Concernant les triples monts « remis sur pied » alors qu’ils étaient placés sur le bord gauche de la bande prenant un air penché : cela veut-il dire mis dans le sens de la bande ? Ou encore à l’horizontale par rapport à l’ensemble de l’écu ? On rencontre également ces trois sortes de dispositions.}
            
\end{teiDiv}
            \begin{teiDiv}[type={section1},xmlid={div02-2}]

              \teiHead[xmlid={lornement-des-manuscrits-001}]{L’ornement des manuscrits}
              \teiP[xmlid={p6-4}]{Nous avons identifié une dizaine de manuscrits présentant un décor intérieur de nature héraldique pour Jules III. Ils se répartissent en trois fonds : à la bibliothèque Vaticane, fonds latin et fonds de la chapelle Sixtine, et en provenance de ce dernier, à la Bibliothèque capitulaire de Tolède. Dans tous les cas, il sera intéressant d’examiner non seulement la présence des armoiries pontificales, mais aussi comment les meubles en furent extraits et employés pour constituer un décor.}
              \begin{teiDiv}[type={section2},xmlid={div03-2}]

                \teiHead[xmlid={le-fonds-vatican-latin-001}]{Le fonds Vatican latin}
                \teiP[xmlid={p7-4}]{Dans le fonds Vatican latin, quatre manuscrits sont l’œuvre du \teiHi[rend={italic}]{scriptor} Ferdinando Ruano, calligraphe et miniaturiste\teiNote[n={8},place={foot},xmlid={ftn42-2}]{\teiP[xmlid={p-065}]{Pierre Petitmengin, avec la collaboration de Jeannine Fohlen, « I manoscritti latini della Vaticana. Uso, acquisizioni, classificazioni », dans Massimo Ceresa (éd.), \teiHi[rend={italic}]{Storia della Biblioteca apostolica vaticana}, vol. II :\teiHi[rend={italic}]{ La Biblioteca vaticana tra riforma cattolica, crescita delle collezioni e nuovo edificio (1535-1590)}, Cité du Vatican, BAV, 2012, p. 43-90, ici p. 52.}}. Deux d’entre eux exploitent la même formule : ils montrent les armoiries entourées d’une branche de laurier. Un meuble a donc été emprunté à l’écu pour lui servir d’encadrement. Certes, il n’était pas rare avant Jules III de placer l’écu dans une couronne de feuillage, les exemples abondent. Mais dans ce cas la couronne est un meuble armorial, il ne s’agit d’ailleurs pas à proprement parler d’une couronne, mais d’une branche de laurier recourbée sur elle-même comme celles qui se trouvent dans les armes de part et d’autre de la bande. Ainsi, l’intérieur et l’extérieur de l’écu se répondent.}
                \teiP[xmlid={p8-4}]{Un exemple en est offert par le \teiHi[rend={italic}]{Registre des dépenses} de la bibliothèque Vaticane en trois volumes (Vat. lat. 3967-3969). Il fut commencé le 28 octobre 1548 et terminé le 4 avril 1555. Le premier volume montre à la page de titre les armes de Paul III Farnèse, ayant été exécuté à la fin de son pontificat. Les deux volumes suivants sont ornés des armes de Jules III (fig. 8). L’écu ovale, forme habituelle pour les papes, s’inscrit dans un rectangle à enroulements, la couronne de laurier servant d’encadrement est formée d’une seule branche recourbée sur elle-même comme le meuble armorial. Peut-être pour souligner ce parallèle, les deux couronnes dans l’écu ne sont pas d’or, mais de sinople. De part et d’autre, l’inscription en lettres d’or indique le nom et la qualité du pontife.}
                \teiP[xmlid={p9-3}]{Le volume des \teiHi[rend={italic}]{Lettres} de Léon IX daté de 1551 offre un exemple similaire (Vat. lat. 3790). Sur la page de garde, nous retrouvons sans surprise la même formule appréciée de ce \teiHi[rend={italic}]{scriptor} (fig. 9). Mais cette fois le laurier est d’or dans l’écu et hors de l’écu ovale, placé dans un cartouche à enroulements. À l’incipit, la lettrine reprend les trois émaux des armes de Jules : l’azur, le gueules et l’or, discrète référence héraldique.}
                \teiP[xmlid={p10-3}]{Les \teiHi[rend={italic}]{Lettres} d’Hildebert, évêque du Mans puis archevêque de Tours (1056-1135), aussi datées de 1551\teiNote[n={9},place={foot},xmlid={ftn43-2}]{\teiP[xmlid={p-066}]{Mandat de paiement de la part du cardinal bibliothécaire Marcello Cervini dans BAV, Vat. lat. 3965, fol. 29r-v.}}, donnent un plein développement à la formule sur le premier folio du manuscrit (Vat. lat. 3841). Cette fois il ne s’agit pas seulement d’extraire un meuble pour en faire un support de l’écu, mais de multiplier les deux meubles afin de créer le décor d’encadrement de la page entière (fig. 10). L’écu a pris la forme d’un bouclier à pointe arrondie. Dans les marges alternent la branche de laurier et les trois monts (en italien \teiHi[rend={italic}]{trimonzi}). Sur la droite, les branches de laurier se referment gracieusement sur elles-mêmes de façon à mieux rappeler le meuble armorial. Dans le reste du manuscrit alternent de petites initiales d’or au fond, tantôt d’azur, tantôt de gueules, rappelant les couleurs des armoiries.}
                \teiP[xmlid={p11-3}]{Enfin, la même manière d’utiliser des meubles pour constituer le décor est reprise au titre de l’\teiHi[rend={italic}]{Exposition sur les Psaumes} de saint Jérôme datée de 1554, mais avec de considérables variations (Vat. lat. 317). Dans un cartouche pourpre à enroulements, le titre en capitales est entouré de trois \teiHi[rend={italic}]{trimonzi} d’or directement extraits de la bande des armoiries de Jules III (fig. 11). Deux épis couronnent le tout. On est d’abord tenté de les interpréter simplement comme des symboles d’abondance, mais il s’agit en fait d’une allusion héraldique à Marcello Cervini, premier cardinal bibliothécaire nommé par Jules III en 1550. Ses armes parlantes représentaient un cervidé, mais aussi des épis de blé, comme on peut le voir, par exemple, dans un manuscrit du fonds Barberiniani graeci (Barb. gr. 265). Le principe d’emprunter des meubles au blason pour décorer la page reste le même, mais cette fois l’allusion héraldique se dédouble.}
                \teiP[xmlid={p12-3}]{Dans les quatre manuscrits examinés, outre la représentation des armoiries complètes, Ferdinando Ruano est resté fidèle à un même procédé qui consistait à en extraire les meubles pour constituer le décor de la page. Dans un cas il est même allé jusqu’à combiner les allusions héraldiques à ses deux protecteurs, le pape Jules III et le cardinal bibliothécaire Marcello Cervini.}
              
\end{teiDiv}
              \begin{teiDiv}[type={section2},xmlid={div04}]

                \teiHead[xmlid={le-fonds-vatican-de-la-chapelle-sixtine-001}]{Le fonds Vatican de la chapelle Sixtine}
                \teiP[xmlid={p13-3}]{La bibliothèque Vaticane conserve trois manuscrits à décor héraldique de Jules III dans le fonds de la chapelle Sixtine. Juan Maria Llorens attribue ce décor en totalité à Vincent Raymond\teiNote[n={10},place={foot},xmlid={ftn44-2}]{\teiP[xmlid={p-067}]{Juan Maria Llorens, « Miniaturas de Vincent Raymond en los manuscritos musicales de la Capilla Sixtina », dans \teiHi[rend={italic}]{Miscelànea Homenaje a Mons. Higinio Anglès}, Barcelone, Consejo Superior de Investigaciones científicas, 1958-1961, 2 vol., t. 1, p. 475-494.}}, Emilia Anna Talamo n’y reconnaît pas sa main partout et l’associe à son atelier\teiNote[n={11},place={foot},xmlid={ftn45-2}]{\teiP[xmlid={p-068}]{Emilia Anna Talamo, \teiHi[rend={italic}]{Codices cantorum, miniature e disegni nei codici della Cappella Sistina}, Florence, Officine del Novecento, 1997, p. 132.}}. Alors qu’il était jeune, sous Léon X ou Clément VII, cet artiste avait été amené à Rome par le cardinal de Clermont-Lodève qui avait été ambassadeur de Louis XII auprès de Jules II\teiNote[n={12},place={foot},xmlid={ftn46-2}]{\teiP[xmlid={p-069}]{Personnage qui apparaît dans l’article d’Emmanuel Moureau du présent volume : « La collégiale Saint-Jean-Baptiste puis Saint-Louis de Castelnau-Bretenoux : un lien étroit et discret avec la papauté ».}}. En 1549, un \teiHi[rend={italic}]{motu proprio} de Paul III l’avait créé enlumineur officiel de la chapelle Sixtine, titre qu’il avait conservé sous Jules III\teiNote[n={13},place={foot},xmlid={ftn47-2}]{\teiP[xmlid={p-070}]{Maria Francesca Saffiotti Dale, « Raymond de Lodève, Vincent », \teiHi[rend={italic}]{Dizionario biografico dei miniatori italiani}, Milan, Sylvestre Bonnard, 2004, p. 899-902.}}.}
                \teiP[xmlid={p14-2}]{À la différence de Ferdinando Ruano qui se contente de déployer les meubles héraldiques tels quels, Vincent Raymond recourt volontiers au genre de la devise. Inventée au Moyen Âge comme une image personnelle distincte du blason, souvent accompagnée d’une inscription, celle-ci connut ses premiers traités imprimés dans l’orbite de la papauté à l’époque dont nous traitons. Alors, en théorie comme en pratique, plutôt que de se différencier complètement du blason, elle tendait à lui emprunter ses meubles\teiNote[n={14},place={foot},xmlid={ftn48-2}]{\teiP[xmlid={p-071}]{Pour un développement sur ce point, voir notre introduction aux actes du colloque de l’École française de Rome (juin 2016), Yvan Loskoutoff (dir.), \teiHi[rend={italic}]{Héraldique et papauté. Moyen Âge-Temps modernes}, Mont-Saint-Aignan, PURH, 2020.}}. Paolo Giovio, protégé des papes Médicis, Léon X et Clément VII, puis de Paul III Farnèse, publie en 1555, année de la mort de Jules III, le premier traité imprimé du genre, le \teiHi[rend={italic}]{Dialogo dell’imprese militari et amorose}. Il signale les exemples de princes qui tirèrent leur devise de leurs armoiries\teiNote[n={15},place={foot},xmlid={ftn49-2}]{\teiP[xmlid={p-072}]{Paolo Giovio,\teiHi[rend={italic}]{ Dialogo dell’imprese militari et amorose di Monsignor Paolo Giouio Vescouo di Nucera}, In Roma, Appresso Antonio Barre, 1555, p. 130-132.}}. En 1556, Girolamo Ruscelli, originaire de Viterbe, fondateur à Rome, en 1541, de l’\teiHi[rend={italic}]{Accademia dei Sdegnati }alors qu’il est attaché au cardinal Grimani, publie à son tour un traité. Il y évoque les liens du genre avec le blason\teiNote[n={16},place={foot},xmlid={ftn50-2}]{\teiP[xmlid={p-073}]{Paolo Giovio, \teiHi[rend={italic}]{Ragionamento di Mons. Paolo Giouio sopra i motti, \& disegni d’arme, \& d’amore, che communemente chiamano imprese, con vn discorso di Girolamo Ruscelli intorno allo stesso soggetto}, In Venetia, Appresso Giordano Ziletti, 1556, p. 230.}}. Dix ans plus tard, il développe cette conception dans un recueil de devises gravées sous lesquelles figurent les armes de leur possesseur :}
                \begin{teiCit}[xmlid={cit02}]

                  \begin{teiQuote}
En ce qui regarde les figures, se trouvent très belles les devises qui se tirent ou se forment des armes, ou des insignes propres à la maison ou à celui-là même par qui elles sont faites, ajoutant ou retranchant, et les modifiant selon les besoins de son intention, assortissant les mots avec règle et beauté\teiNote[n={17},place={foot},xmlid={ftn51-2}]{\teiP[xmlid={p-074}]{Girolamo Ruscelli,\teiHi[rend={italic}]{ Le imprese illustri con espositioni, et discorsi del S.or Ieronimo Ruscelli}, Venise, \lbrack{}Damiano Zenaro\rbrack{}, 1566, p. 14 : « \teiHi[rend={italic}]{CHE inquanto alle figure, riescono bellissime quelle Imprese, che si traggono, ò si formano dall’Arme, ò dall’Insegne proprie della casa, ò di colui stesso, da chi si fanno, aggiungendoui, ò togliendone, \& mutandole secondo il bisogno dell’intention sua, accomodandoui le parole regolatamente, \& con leggiadria} ».}}.
\end{teiQuote}
                
\end{teiCit}
                \teiP[xmlid={p15-2}]{C’est effectivement ce qui s’était pratiqué pour Jules III.}
                \teiP[xmlid={p16-2}]{Commençons par le plus simple, le missel coté Cappella Sistina 154. Au début (fol. 3v), figurent les armes du pape (fig. 12). Comme avait procédé Ferdinando Ruano, les branches de laurier ont été extraites de l’écu pour lui servir de support, cependant, elles ne sont pas recourbées en couronne, mais disposées en sautoir. Par ailleurs elles ont été traitées en sinople, différemment de celles de l’écu qui sont d’or. Ruano avait, lui, choisi la même couleur. Le folio suivant montre une devise dans un cartouche d’or ovale à enroulements (fig. 13). Sur un champ de gueules se détache le \teiHi[rend={italic}]{trimonzio} d’où surgissent des fleurs et des fruits suggérant l’abondance du pontificat, motif présent de manière accessoire sur deux médailles, celle de l’\teiHi[rend={italic}]{Hilaritas publica }et celle de Bologne\teiNote[n={18},place={foot},xmlid={ftn52-2}]{\teiP[xmlid={p-075}]{Adolfo Modesti, \teiHi[rend={italic}]{Corpus numismatum omnium romanorum pontificum}, Rome, Adolfo Modesti, 2003, vol. 2, n\teiHi[rend={sup}]{o} 398, p. 290 et n\teiHi[rend={sup}]{o} 414, p. 322.}}. Le mot dit que ses fondements sont sur les saints monts : \teiHi[rend={italic}]{Fundamenta eius in montibus sanctis} (psaume 86, v. 1). Alessandro Nova a montré combien le thème du « mont », en référence au nom et au blason du pontife, dominait sa célébration\teiNote[n={19},place={foot},xmlid={ftn53-2}]{\teiP[xmlid={p-076}]{\teiHi[rend={italic}]{The Artistic Patronage of Pope Julius III (1550-1555), Profane Imagery and Buildings for the De Monte Family in Rome}, New York-Londres, Garland Publishing, 1988, p. 28.}}. Au centre de la villa Giulia comme des appartements de Jules III au Vatican, décorés pour l’essentiel par l’atelier de Prospero Fontana, se trouve la salle des sept monts de Rome, chacun étant représenté dans la frise. Par ailleurs, le verset des Psaumes que nous venons de citer constitue un topos de la littérature encomiastique. Le docteur juif converti Andrea de Monte\teiNote[n={20},place={foot},xmlid={ftn54-2}]{\teiP[xmlid={p-077}]{Les Juifs convertis prenaient parfois le nom et les armes du pape qui avait patronné leur conversion. Sous Grégoire XIII Boncompagni, voir le cas des financiers Corcos dans Yvan Loskoutoff, « Concessions héraldiques à la fin de la Renaissance d’après les recueils de brefs de saint Pie V (1566-1572), Grégoire XIII (1572-1585) et Sixte-Quint (1585-1590) conservés aux Archives du Vatican », dans Yvan Loskoutoff (dir.), \teiHi[rend={italic}]{Héraldique et papauté. Moyen Âge-Temps modernes}, \teiHi[rend={italic}]{op. cit.}, p. 179-196. Sur les Corcos voir Isabelle Poutrin, \teiHi[rend={italic}]{Les convertis du pape. Une famille de banquiers juifs à Rome au }\teiHi[rend={small-caps italic}]{xvi}\teiHi[rend={italic sup}]{e}\teiHi[rend={italic}]{ siècle}, Paris, Seuil, 2023.}}, dans un discours prononcé devant le pape où il entend montrer que les armoiries de celui-ci étaient annoncées par la sainte Écriture, en commente divers passages relatifs à des monts. Dans celui du psaume 86, il considère qu’ils sont le fondement de l’Église\teiNote[n={21},place={foot},xmlid={ftn55-2}]{\teiP[xmlid={p-078}]{Vat. lat. 3561, sans titre, sans date, fol. 8. Le texte a été publié par Antonio Blado, sans date, sous le titre : \teiHi[rend={italic}]{Oratio habita apud S. D. N. Iulium diuina prouidentia Papam Tertium per Andream de Monte doctorem, quondam Haebreum iam vero sacro baptismate initiatum, ex hebreo in latinum versa, in qua suadetur summum Pontificem ex illustri familia montis \& illius insignia in sacris literis fuisse praenunciata}. }}. Mariano Cavense, dans son poème de cent folios en six livres, \teiHi[rend={italic}]{Triumphus montium,} daté de mai 1551, s’inspire plusieurs fois du premier verset du psaume 86 qu’il relie aux armoiries du pape. Pour lui c’est aussi l’Église qui est fondée sur les monts : « \teiHi[rend={italic}]{Fundamenta eius sanctis in montibus aedis} » \lbrack{}Les fondements du sanctuaire sont sur les saints monts\teiNote[n={22},place={foot},xmlid={ftn56-2}]{\teiP[xmlid={p-079}]{Biblioteca Angelica, Rome, ms. 928, fol. 32v, voir aussi fol. 18v, 20v, 21, 83 et 83v.}}\rbrack{}. C’était un lieu commun de la célébration pontificale. Une médaille lui fut consacrée en 1553 montrant une église imaginaire sur les trois monts avec le verset du psaume pour légende\teiNote[n={23},place={foot},xmlid={ftn57-2}]{\teiP[xmlid={p-080}]{Adolfo Modesti, \teiHi[rend={italic}]{op. cit.}, n\teiHi[rend={sup}]{o} 435, p. 364.}}.}
                \teiP[xmlid={p17-2}]{L’\teiHi[rend={italic}]{Officium ad Tertiam }coté Cappella Sistina 67, écrit par Galeazzo Herculano, se révèle intéressant pour ses lettrines. Au début du volume, quatre initiales O de \teiHi[rend={italic}]{Os lingua} contiennent une devise. Au folio 4v nous retrouvons les trois monts porteurs de fleurs et de fruits, cette fois munis de l’inscription : \teiHi[rend={italic}]{Fructiferi} (fig. 14). Au-dessous, un autre O contient une targe suspendue où se lit le nom du pape : \teiHi[rend={italic}]{Iulius III Pont. M}. Au folio 5, le \teiHi[rend={italic}]{trimonzio} reparaît supportant cette fois une triple lampe, gage non plus d’abondance, mais d’enseignement, la légende énonce qu’elle luit pour tous : \teiHi[rend={italic}]{Omnibus lucet}. Plus bas dans le folio, le même \teiHi[rend={italic}]{trimonzio} disposé sur trois marches supporte une église sous laquelle se lit l’inscription rencontrée dans le précédent manuscrit pour l’image des fleurs et des fruits. Elle paraît plus appropriée ici : \teiHi[rend={italic}]{Fundamenta eius in montibus sanctis} (fig. 15). Emilia Anna Talamo suppose que les trois portes font « avec bonne probabilité » allusion au jubilé de 1550 et en tire argument pour dater le manuscrit de cette année-là\teiNote[n={24},place={foot},xmlid={ftn58-2}]{\teiP[xmlid={p-081}]{Emilia Anna Talamo,\teiHi[rend={italic}]{ op. cit}., p. 132.}}. Au folio 7, reparaît dans une lettrine « A » la devise de la triple lampe sur les monts avec la même légende, un phylactère ajoutant le mot : \teiHi[rend={italic}]{Supra}. Enfin, comme pour clore cette suite symbolique au folio 7v, une grande lettrine « L » inclut les armes du pape dont le nom est inscrit sur le jambage vertical (fig. 16). Outre cette suite de devises, on constate aussi que certaines initiales contiennent le \teiHi[rend={italic}]{trimonzio} en creux ou en relief dans leur corps même (fig. 17). Parfois il est remplacé par le lys de Paul III Farnèse. C’était le pape qui avait créé Giovanni Maria del Monte cardinal en 1536, c’était aussi celui auquel Vincent Raymond devait son office d’enlumineur de la chapelle Sixtine. Un lien évident existe entre les manuscrits Cappella Sistina 154 et 67, le second utilisant une devise présente dans le premier avec une variante pour l’inscription qui passe à une autre devise.}
                \teiP[xmlid={p18-2}]{Le lectionnaire coté Cappella Sistina 213 présente le même jeu de lettrines que Cappella Sistina 67, incluant dans les jambages tantôt le \teiHi[rend={italic}]{trimonzio} de Jules III tantôt le lys de Paul III. Les couleurs sont celles du blason de Jules : or, azur et gueules. Le premier folio est entièrement décoré (fig. 18). Emilia Anna Talamo voit là le fruit d’une collaboration entre Vincent Raymond et son adjoint Apollonio de’ Bonfratelli\teiNote[n={25},place={foot},xmlid={ftn59-2}]{\teiP[xmlid={p-082}]{\teiHi[rend={italic}]{Ibid}., p. 135.}}. Les armoiries sont placées en bas sur un linge rappelant celui de sainte Véronique, le texte étant consacré à la Passion selon saint Matthieu. De part et d’autre un \teiHi[rend={italic}]{putto} en vol couronne de laurier le \teiHi[rend={italic}]{trimonzio}, motif repris en haut de la page, au-dessus du titre. Même sans inscription, c’était une devise présente dans les décors des résidences de Jules III attribués à Prospero Fontana. Dans les appartements du Vatican elle orne les angles de la salle des sept monts de Rome. À la villa Giulia, aux angles de la salle correspondante, elle fait l’objet d’une variante, non pas couronnée, mais entourée de laurier, motif multiplié dans le décor à grotesques du portique (fig. 19). Dans la salle des sept monts, la vue de l’Esquilin inclut en bas à gauche une fontaine montrant le \teiHi[rend={italic}]{trimonzio} couronné de lauriers d’où l’eau jaillit. Il faut l’interpréter comme une signature ingénieuse de Fontana. La parenté entre décor architectural et décor de manuscrit n’est pas surprenante, un même motif symbolique pouvant servir à différents supports.}
              
\end{teiDiv}
              \begin{teiDiv}[type={section2},xmlid={div05}]

                \teiHead[xmlid={la-bibliotheque-capitulaire-de-tolede-001}]{La Bibliothèque capitulaire de Tolède}
                \teiP[xmlid={p19-2}]{En 1798, l’occupation des troupes françaises causa la dispersion des manuscrits de la chapelle Sixtine, dont les miniatures ou les folios découpés alimentèrent diverses collections. Le cardinal Antonio de Lorenzana y Buitròn parvint à sauver quelques volumes, aujourd’hui à la Bibliothèque capitulaire de Tolède ou à la Bibliothèque nationale d’Espagne\teiNote[n={26},place={foot},xmlid={ftn60-2}]{\teiP[xmlid={p-083}]{Elena De Laurentiis et Emilia Anna Talamo, \teiHi[rend={italic}]{The Lost Manuscripts from the Sistine Chapel, An Epic Journey from Rome to Toledo}, catalogue d’exposition, Dallas, Meadows Museum-Southern Methodist University, 2010.}}. Le dernier folio ajouté dans un pontifical du cardinal Pietro Barbo (1400-1450) est le premier folio d’un évangéliaire de Jules III attribué à Vincent Raymond. On peut y voir les armoiries de ce pape disposées dans un cercle au bas de l’encadrement (fig. 20).}
                \teiP[xmlid={p20-2}]{La Bibliothèque capitulaire de Tolède conserve un autre manuscrit intéressant : le missel du cardinal Girolamo Dandini, élevé à la pourpre par Jules III en 1551. Son écu figure au bas du folio d’incipit parti avec celui de ce pape (fig. 21). Il était effectivement fréquent qu’un cardinal associât ses armes à celles du pontife qui l’avait promu. Cela pouvait se faire de plusieurs manières sur l’écu : en chef, coupé, parti ou écartelé. Il semble que la disposition partie ait joui d’un certain succès. Outre le présent exemple, elle fut aussi choisie par le cardinal Ricci de Montepulciano, de la même promotion de 1551, pour le décor de ses appartements du Vatican\teiNote[n={27},place={foot},xmlid={ftn61-2}]{\teiP[xmlid={p-084}]{Maria Giulia Aurigemma, \teiHi[rend={italic}]{Palazzo Firenze in Campo Marzio},\teiHi[rend={italic}]{ op. cit.}, illustration p. 111.}} comme pour sa chapelle à San Pietro in Montorio, en face de celle de la famille del Monte\teiNote[n={28},place={foot},xmlid={ftn62-2}]{\teiP[xmlid={p-085}]{Alessandro Zuccari (dir.),\teiHi[rend={italic}]{ La Spagna sul Gianicolo}, Rome, Eurografica, 2004, t. 1, \teiHi[rend={italic}]{San Pietro in Montorio}, fig. 35, p. 137 et fig. 42, p. 145.}}. Par ailleurs, la lettrine G de \teiHi[rend={italic}]{Gaudete} associe dans une disposition écartelée la couronne de laurier de Jules III avec les trois étoiles et les trois lys de Dandini.}
                \teiP[xmlid={p21-2}]{Que ce soit chez Ferdinando Ruano ou chez Vincent Raymond le décor héraldique du manuscrit obéit aux mêmes principes : figuration des armes complètes au bas de la première page et extraction des meubles, soit pour créer des supports, soit pour remplir les encadrements, soit pour illustrer les lettrines. Vincent Raymond se distingue par l’emploi d’une suite de devises, elles aussi inspirées du blason, comme c’était devenu la règle au temps de Jules III. On constate des rapports avec les décors monumentaux qui obéissaient aux mêmes principes : armoiries au centre et diffusion des meubles, fréquemment sous la forme des devises.}
              
\end{teiDiv}
            
\end{teiDiv}
            \begin{teiDiv}[type={section1},xmlid={div06}]

              \teiHead[xmlid={lornement-du-livre-imprime-001}]{L’ornement du livre imprimé}
              \teiP[xmlid={p22-2}]{Au temps de Jules III, les grands centres d’édition sont Venise et Florence. L’héraldique pontificale n’y apparaît guère dans le décor des livres où dominent les marques de libraires. C’est à Rome, centre d’édition d’importance moindre, qu’elle se développe, mais sans excès. On n’y trouve ni frontispices, ni bandeaux, ni culs-de-lampe. Nous avons rencontré un titre gravé et trois lettrines. Elle se résume presque entièrement aux vignettes de titre. Comment les armoiries y sont-elles représentées ? Nous commencerons par Rome, puis traiterons de Venise, Florence et Bologne, petit centre d’édition dans les États pontificaux, ainsi que de Rimini.}
              \begin{teiDiv}[type={section2},xmlid={div07}]

                \teiHead[xmlid={rome-001}]{Rome}
                \teiP[xmlid={p23-2}]{À Rome domine Antonio Blado, l’éditeur des textes officiels de la papauté, qui crée des types de vignettes héraldiques largement utilisés. Les petits éditeurs créent aussi parfois les leurs qui sont souvent des hapax. Nous commencerons par les armoiries simples, puis représentées dans une couronne ou plus rarement munies de tenants, pour terminer par le seul cas de titre gravé.}
                \begin{teiDiv}[type={section3},xmlid={div08}]

                  \teiHead[xmlid={les-armoiries-simples-001}]{Les armoiries simples}
                  \teiP[xmlid={p24-2}]{Les cas d’armoiries représentées simplement dans un cartouche ou même sans cartouche sont peu fréquents, on préfère les orner. Chez les frères Valerio et Luigi Dorico on voit paraître une élégante vignette ovale datable de 1550 pour le \teiHi[rend={italic}]{De exercitatione iurisperitorum} d’Antonio Massa (fig. 22). Elle est reprise en 1551 pour la \teiHi[rend={italic}]{Genealogia d’imperatori romani et constantinopolitani} d’Andrea Angelo Commeno, en 1552 pour les \teiHi[rend={italic}]{Commentariorum de bello Aphrodisiensi libri quinque} d’Orazio Nucula et en 1554 pour la \teiHi[rend={italic}]{Copia delle lettere del Serenissimo Re d’Inghilterra, \& del Reuerendissimo Card. Polo Legato della S. Sede Apostolica alla Santità di N. S. Iulio Papa III. sopra la reduttione di quel Regno alla vnione della Santa Madre Chiesa ; \& obedienza della Sede Apostolica}. C’est la vignette propre de la maison. En 1553, une version plus élaborée est créée dans un format plus grand pour la \teiHi[rend={italic}]{Erudita atque luculenta disputatio} de Giovanni Girolamo Albani. Un rang de perles a été ajouté au bord et des étoiles au fond, un cartouche bordé de masques entoure l’écu en forme de bouclier à pointe arrondie (fig. 23).}
                  \teiP[xmlid={p25-2}]{Une autre formule simple représente les armoiries dans un carré, nous la trouvons chez le même éditeur en 1550 pour le \teiHi[rend={italic}]{De iobileo et indulgentiis} de Giovanni Battista Poiana (fig. 24) et la même année dans une version légèrement différente chez Antonio Blado pour les \teiHi[rend={italic}]{Regulae tres S. D. N. D. Iulii diuina prouidentia Papae. III }(fig. 25). On assiste dans ce dernier cas à une tentative de créer l’illusion des couleurs dans le noir et blanc de la gravure : la bande a été hachurée. C’est aussi le cas pour le \teiHi[rend={italic}]{De concordia servanda} de Silvestro Sigonio publié la même année par Antonio Blado avec la sous-traitance de Stefano Nicolini, dont nous ne connaissons que ce seul exemple en hapax (fig. 26). Les armes ovales sont figurées dans un cartouche ajouré. Toujours chez Antonio Blado, un autre hapax, présentant un cartouche ajouré, cherche aussi à rendre sensible la couleur de la bande pour un édit de grand format contre les Juifs détenant des livres hostiles à la foi catholique, datable de 1554 (vignette : 16 × 17 cm) (fig. 27). L’effort pour exprimer la différence de couleur de la bande par des hachures est fréquent dans les autres cas étudiés, il domine même très largement. On le remarque dans quinze figures sur vingt concernées (notamment fig. 25 à 33 sans interruption). Les \teiHi[rend={italic}]{trimonzi} sont aussi souvent hachurés, mais partiellement pour rendre le volume autant que la couleur (fig. 23, 26-28, 31, 34, 38-40 et 42), il en est de même pour le champ de l’écu (fig. 38 à 41). La critique situe l’apparition des hachures héraldiques dans les cartes géographiques anversoises à « l’extrême fin du \teiHi[rend={small-caps}]{xvi}\teiHi[rend={sup}]{e} siècle\teiNote[n={29},place={foot},xmlid={ftn63-2}]{\teiP[xmlid={p-086}]{Voir Michel Pastoureau, « Aux origines des hachures héraldiques (\teiHi[rend={small-caps}]{xv}\teiHi[rend={sup}]{e}-\teiHi[rend={small-caps}]{xvii}\teiHi[rend={sup}]{e} siècle) », \teiHi[rend={italic}]{Revue française d’héraldique et de sigillographie}, t. 65, 1995, p. 21-31, ici p. 26, « Les premières hachures ». }} ». Une cinquantaine d’années plus tôt, il semble que le livre romain annonce cette systématisation. Certes, les hachures ne désignent pas encore des couleurs précises, mais on a manifestement cherché à matérialiser leurs différences tout autant que les effets de volume. Dès lors, il n’est pas surprenant que cette invention soit attribuée à un jésuite romain, le père Pietrasanta, dans la première moitié du \teiHi[rend={small-caps}]{xvii}\teiHi[rend={sup}]{e} siècle.}
                
\end{teiDiv}
                \begin{teiDiv}[type={section3},xmlid={div09}]

                  \teiHead[xmlid={les-armoiries-dans-une-couronne-001}]{Les armoiries dans une couronne}
                  \teiP[xmlid={p26-2}]{Nous avons déjà rencontré les armoiries ceintes d’une couronne dans les manuscrits, nous les retrouvons dans les imprimés. Un premier exemple est offert par l’édition grecque du commentaire d’Eustathius sur l’\teiHi[rend={italic}]{Iliade}. L’entreprise de vaste envergure en quatre volumes in-folio avait commencé sous Paul III chez Antonio Blado. Le quatrième volume fut publié au début du pontificat de Jules III en 1550. Le titre porte la marque de l’éditeur, c’est le privilège qui est orné des armes du pontife (fig. 29). Néanmoins, la couronne qui sert d’encadrement n’est pas de laurier ou d’olivier, mais de fruits divers en signe d’abondance. Cela s’explique, car il s’agit de la reprise d’une pièce largement utilisée sous Paul III chez le même éditeur, par exemple en 1544 pour le \teiHi[rend={italic}]{De nova christiani orbis pace oratio} de Filippo Archinto. La même année 1550, Antonio Blado pratiqua aussi la réutilisation pour la \teiHi[rend={italic}]{Bulla confirmationis priuilegiorum militum sancti Pauli} dont le même modèle de couronne fruitière avait souvent servi sous Paul III, par exemple pour la \teiHi[rend={italic}]{Bulla indulgentiarum et facultatum pro erectione Ecclesiae, \& Monasterii S. Tholemaei in Vrbae Nepesina}.}
                  \teiP[xmlid={p27-2}]{Cependant, Antonio Blado ne se limita pas à la récupération pour Jules III. Il créa une autre couronne fruitière employée pour le discours dédié en 1553 au pontife par Uberto Foglietta à l’occasion de la Toussaint, \teiHi[rend={italic}]{In festo omnium Beatorum }mais aussi pour diverses pièces officielles. Elle servit assez souvent et elle est reconnaissable aux fissures du bois en haut et en bas (fig. 29).}
                  \teiP[xmlid={p28-2}]{Il existe en revanche une couronne, non de fruits, mais de laurier, créée sous le pontificat de Jules III par Antonio Blado. Elle est attachée par un ruban où se déroule une inscription de vœu de longue prospérité : « \teiHi[rend={italic}]{Diu felix} ». On la trouve pour les \teiHi[rend={italic}]{Enarrationes in quinque priora capita libri geneseos} d’Ambrosio Catharino publiées sous la direction de Lancelotto Politi en 1552 non sans luxe, la vignette étant dans ce cas encadrée d’une bordure. On peut la considérer comme la pièce classique par excellence du pontificat de Jules III. Elle fut très largement utilisée. Donnons l’exemple, la même année, des \teiHi[rend={italic}]{Novae regvlae cancellariae apostolicae} (fig. 30). Antonio Blado la reprit avec changement des armoiries pour les papes qui suivirent : Pie IV (1559-1565), saint Pie V (1566-1572) et jusqu’à Grégoire XIII (1572-1585\teiNote[n={30},place={foot},xmlid={ftn64-2}]{\teiP[xmlid={p-087}]{Pour les références, voir Yvan Loskoutoff, \teiHi[rend={italic}]{Un art de la Réforme catholique}, vol. 2, \teiHi[rend={italic}]{La symbolique du pape Grégoire XIII (1572-1585) et des Boncompagni}, Paris, Honoré Champion, 2018, p. 157.}}).}
                  \teiP[xmlid={p29-2}]{À la première page des \teiHi[rend={italic}]{Enarrationes}, au début de la dédicace à Jules III, on remarque une lettrine pour le « L » de \teiHi[rend={italic}]{Legitur} (fig. 31). Elle aussi fut remployée. On la trouve la même année au début du \teiHi[rend={italic}]{Libellus vere Christiana lectione dignus diuersas res Turcharum} de Bartolomej Georgijevic et du \teiHi[rend={italic}]{Rimedio delle podagre} d’Andrès Laguna, ou encore en 1553 au \teiHi[rend={italic}]{proemium} du \teiHi[rend={italic}]{Trattato di Scientia d’Arme} de Camillo Agrippa. La lettre « L » s’inscrit devant les armoiries placées sur le fond hachuré de la vignette, créant un effet de profondeur à trois plans successifs. Nous connaissons une autre lettrine datant du début du pontificat, mais en une seule occurrence dans le recueil des \teiHi[rend={italic}]{Bullae diuersorum pontificum} publié chez une petite éditrice, Girolama Cartolari, en 1550. Chaque pape bénéficie d’une pièce d’incipit montrant ses armes, dont Jules III, qui conclut l’ouvrage (fig. 32).}
                
\end{teiDiv}
                \begin{teiDiv}[type={section3},xmlid={div10}]

                  \teiHead[xmlid={les-armoiries-a-tenants-001}]{Les armoiries à tenants}
                  \teiP[xmlid={p30-2}]{Les tenants des armoiries papales sont traditionnellement des \teiHi[rend={italic}]{putti} ou des anges. Nous n’en avons trouvé que deux exemples parmi les vignettes de titre romaines, l’une largement employée, l’autre une seule fois. La première revient à Antonio Blado. Elle orne plusieurs textes officiels. On y voit deux \teiHi[rend={italic}]{putti} dans des positions asymétriques, l’un accroupi de face, l’autre debout de dos, tenir la couronne de feuillage qui entoure les armoiries. Son emploi se révèle intéressant à la fin du \teiHi[rend={italic}]{Libellus} de Bartolomej Georgijevic (fig. 33), déjà rencontré pour la lettrine « L ». Elle sert en effet d’illustration à une épigramme latine commentant les armes de Jules : « Regarde comme l’olivier entoure les trois monts verdoyants / Et comme un seul sommet s’élève en trois. / De très grands mystères viennent à l’esprit, Jules entoure / De la même paix l’Europe, l’Asie et la Lybie. » La pointe finale fait référence à Jules César pour évoquer Jules qui concluait la paix avec Parme l’année de la publication de l’ouvrage en 1552. Il s’agit bien ici de rameaux d’olivier. La vignette de titre est devenue l’illustration d’un texte poétique, formule qui ne resta pas réservée à Jules III. Nous en connaissons au moins un autre exemple pour Grégoire XIII dans un recueil de vers sur la victoire contre les Turcs publié en 1572\teiNote[n={31},place={foot},xmlid={ftn65-2}]{\teiP[xmlid={p-088}]{\teiHi[rend={italic}]{Ibid}., p. 136, fig. 34.}}.}
                  \teiP[xmlid={p31-2}]{L’occurrence unique revient à un petit éditeur, Giovanni Maria Viotti, qui dirigeait les presses de Sainte Brigitte, fondées par deux Suédois, les frères Magnus, pour exalter l’histoire catholique de leur nation\teiNote[n={32},place={foot},xmlid={ftn66-2}]{\teiP[xmlid={p-089}]{Valentino Romani, « Per la storia dell’editoria italiana del Cinquecento: le edizioni romane “In aedibus sanctae Brigidae (1553-1557)” », \teiHi[rend={italic}]{Rara Volumina}, n\teiHi[rend={sup}]{o} 1, 1998, p. 23-36.}}. L’un d’eux, Iohannes, archevêque d’Uppsala, était l’auteur de cette \teiHi[rend={italic}]{Historia de omnibus Gothorum Sueonumque regibus} in-folio, datée de 1554. Comme dans l’exemple précédent, la vignette est mise en rapport avec un texte, une inscription qui la domine et qui développe aussi le thème de la paix : « \teiHi[rend={italic}]{Suscipiant montes pacem populo} » \lbrack{}Les monts produisent la paix pour le peuple\rbrack{} (fig. 34). Les armes sont tenues par deux \teiHi[rend={italic}]{putti}, cette fois symétriques, brandissant chacun une couronne qui rappelle celles de l’écu. Le \teiHi[rend={italic}]{putto} porte couronne n’est pas absent des décors de Jules III. Dans les appartements du Vatican il suffit de penser à la salle des \teiHi[rend={italic}]{putti}, où plusieurs de ceux qui jouent sur la frise autour du nom du pape sont munis de cet accessoire. Mais nous ne saurions dire s’il s’agit de laurier ou d’olivier.}
                
\end{teiDiv}
                \begin{teiDiv}[type={section3},xmlid={div11}]

                  \teiHead[xmlid={le-missarum-liber-primus-de-palestrina-001}]{Le \teiHi[rend={italic}]{Missarum liber primus} de Palestrina}
                  \teiP[xmlid={p32-2}]{L’ornement héraldique le plus complet pour le livre imprimé sous Jules III se trouve dans le grand in-folio du \teiHi[rend={italic}]{Missarum liber primus} de Palestrina, publié par les frères Dorico en 1554. Il offre en effet le seul exemple de titre gravé blasonnant sous ce pontificat (fig. 35). Palestrina jouissait de la faveur du pape, ce que laisse transparaître le luxe du volume. Le cadre est composé d’instruments de musique. On voit en bas la marque du libraire, Pégase faisant jaillir la source Hippocrène. Cette marque est entourée d’Orphée charmant les créatures sauvages à gauche et du châtiment de Marsyas à droite. Au centre de la page, le compositeur présente son ouvrage à son bienfaiteur assis dans un fauteuil dont le dossier et l’accoudoir portent le \teiHi[rend={italic}]{trimonzio}. Ce genre de fauteuil héraldique était devenu de rigueur dans le portrait pontifical au moins depuis celui de Léon X, où Raphaël avait introduit la \teiHi[rend={italic}]{palla} des Médicis. Mais ce titre, aussi magnifique soit-il, n’est que l’adaptation de celui qui avait été créé par les mêmes éditeurs pour le deuxième volume des \teiHi[rend={italic}]{Messes} de Cristobal de Morales, dédié à Paul III Farnèse en 1544. Les visages ont été changés et les \teiHi[rend={italic}]{trimonzi} sont venus remplacer sur le fauteuil des boules portant les lys des Farnèse. En bas se trouvait une composition aux armes de Paul III. Celle que nous voyons pour Jules III provient du premier volume des \teiHi[rend={italic}]{Messes} de Morales dédié en 1544 à Cosme de Médicis.}
                  \teiP[xmlid={p33-2}]{Là ne s’arrête pas le décor blasonnant. Le volume de Palestrina contient cinq messes. Chacune est pourvue d’une grande initiale ornée d’une scène religieuse, sauf la première qui inclut le portrait du pape et ses armes pour la messe \teiHi[rend={italic}]{Ecce sacerdos magnus} (fig. 36). Le lion n’est pas tout à fait dans la symbolique de Jules III, mais l’écu qu’il supporte est bien le sien, ceint d’une couronne de feuillage. Il s’agit du \teiHi[rend={italic}]{marzocco} florentin et sa présence s’explique, car nous avons là le remaniement d’une initiale du volume dédié à Cosme de Médicis qui portait ses armes. Enfin, seulement dans cette première messe, les Dorico ont ponctué le début de certaines portées de leur élégante vignette ovale (fig. 22). Ce volume paraît une tentative de transposer l’art du manuscrit dans l’imprimé.}
                
\end{teiDiv}
              
\end{teiDiv}
              \begin{teiDiv}[type={section2},xmlid={div12}]

                \teiHead[xmlid={hors-de-rome-001}]{Hors de Rome}
                \teiP[xmlid={p34}]{À Venise, Florence, Bologne et Rimini, les vignettes héraldiques de Jules III se révèlent très exceptionnelles, les éditeurs utilisant plutôt leurs marques.}
                \teiP[xmlid={p35}]{À Venise, dès l’élection du pontife en 1550, l’éditeur Prospero Danza publie un livret décrivant les cérémonies : \teiHi[rend={italic}]{La elettion del beatissimo \& summo pontifice Iulio tertio} (fig. 37). La vignette carrée spécialement conçue pour l’occasion est curieuse : elle montre les armes Farnèse coupées de celles de del Monte (avec une erreur sur la bande aux trois monts qui est devenue une barre), le futur Jules III avait en effet été élevé à la pourpre par Paul III Farnèse dans la promotion du 22 décembre 1536. Les cardinaux portaient fréquemment leurs armes coupées de celles du pape qui les avait créés. Néanmoins, les pièces extérieures sont celles de la papauté : la tiare et les clefs. Il semble que l’on ait voulu montrer le moment intermédiaire où le cardinal accède au pontificat. Nous ne connaissons pas d’autre exemple où Jules III ait brisé ses armes de celles de Paul III. Cette pièce présente donc une iconographie héraldique d’une certaine rareté.}
                \teiP[xmlid={p36}]{Deux autres pièces vénitiennes datent de 1553. L’une illustre l’œuvre d’un ecclésiastique, Giacomo Nacchiante, les \teiHi[rend={italic}]{Enarrationes piae, doctae, et catholicae : in epistolam Pauli ad Ephesios} (fig. 38). Les armes sont présentées dans un ovale grisé, incluses dans un cartouche à tête de lion au sommet et bustes d’angelots sur les côtés. L’autre ouvrage est celui d’un philologue de Ravenne, Tommaso Giannotti Rangoni, et explique comment prolonger la vie de l’homme au-delà de 120 ans, \teiHi[rend={italic}]{De vita hominis ultra CXX. annos protrahenda}. Les armes ovales sont présentées sur une tête d’angelot dans un cartouche à enroulements (fig. 39).}
                \teiP[xmlid={p37}]{À Florence, en 1550, ce sont deux ouvrages d’un même auteur, le père Lorenzo Davidico\teiNote[n={33},place={foot},xmlid={ftn67-2}]{\teiP[xmlid={p-090}]{Voir Carlo von Flüe, « Davidico, Lorenzo », \teiHi[rend={italic}]{DBI}, Rome, Istituto della Enciclopedia italiana, vol. 33, 1987, p. 157-160.}}, qui portent la même vignette héraldique : \teiHi[rend={italic}]{Il vittorioso trionfo di Maria V. Contra Luterani, }dédié au Sacré Collège et l’\teiHi[rend={italic}]{Anotomia delli vitii, }dédiée au pape (fig. 40). Les armes sont présentées dans un cartouche à enroulement sur fond de tenture. Le choix de la vignette s’est évidemment fait en lien avec la dédicace, peut-être sur intervention de l’auteur, car nous n’avons pas trouvé d’autre exemple de vignette papale dans cette ville.}
                \teiP[xmlid={p38}]{À Bologne, on remarque chez l’éditeur Anselmo Giaccarello une vignette à supports de \teiHi[rend={italic}]{putti} utilisée pour deux ouvrages : en 1550, le \teiHi[rend={italic}]{De nothis spuriisque filiis liber }de Gabriele Paleotti, et en 1551, le \teiHi[rend={italic}]{De concilio liber I} de Lorenzo Dal Gambaro Sclarici, dédié au cardinal Innocenzo del Monte, neveu adoptif de Jules III (fig. 41). Les armes de celui-ci reparaissent à la toute fin du pontificat en 1554 sur deux brefs dans une composition dominée par les armes du pontife. Enfin, dans cette même ville, Bartolomeo Maggi publia en 1552 chez Bartholomeo Bonardi un traité \teiHi[rend={italic}]{De vulnerum sclopetorum et bombardarum curatione}, dédié au frère aîné de Jules III, Baldovino del Monte. En lien avec cette dédicace une vaste pièce de titre fut conçue montrant les armes de Baldovino, une femme en cimier tenant un phylactère inscrit du mot \teiHi[rend={italic}]{Merito} (fig. 42). En 1497, Fabiano del Monte, oncle de Jules III, fut accueilli dans la noblesse d’Arezzo en récompense de ses talents de juriste. Il se choisit alors la devise de Pallas avec le mot \teiHi[rend={italic}]{Pro merito} qui aurait été abrégé par Jules III en \teiHi[rend={italic}]{Merito}\teiNote[n={34},place={foot},xmlid={ftn68-2}]{\teiP[xmlid={p-091}]{Agostino Fortunio, \teiHi[rend={italic}]{Cronichetta del Monte San Savino di Toscana},\teiHi[rend={italic}]{ op. cit.}, p. 38, cité par Alessandro Nova, « Bartolomeo Ammannati e Prospero Fontana a Palazzo Firenze. Architettura e emblemi per Giulio III Del Monte », \teiHi[rend={italic}]{Ricerche di Storia dell’arte}, n\teiHi[rend={sup}]{o} 21, 1983, p. 53-76, ici p. 74, n. 59.}}. On est tenté d’en déduire que la femme représentée en cimier pourrait être la déesse de la sagesse.}
                \teiP[xmlid={p39}]{À Rimini, en 1552, chez Erasmo Virginio, le chanoine Publio Francesco Modesti consacra sa \teiHi[rend={italic}]{Christiana pietas} à la louange de Jules III. La pièce de titre est une copie fidèle de celle qui avait été utilisée à Bologne par Anselmo Giaccarello en 1550 pour le \teiHi[rend={italic}]{De nothis spuriisque filiis liber }de Gabriele Paleotti et en 1551 pour le \teiHi[rend={italic}]{De concilio liber I }de Lorenzo Dal Gambaro Sclarici (fig. 41). Seule a changé la disposition des \teiHi[rend={italic}]{trimonzii}, dans le sens de la bande à Bologne, sur le bord gauche de la bande à Rimini. Cette similitude n’est pas surprenante, Rimini étant voisine de Bologne. Il pourrait s’agir du même instrument dont un détail a été modifié. Nous avons peut-être ici la trace du changement de disposition des monts qui aurait été voulu par le pontife selon le témoignage cité au début de cette étude. Un texte latin a été ajouté autour de la vignette disant : « Les choses que tu vois dépeintes ici vivent dans mon cœur. Si ce sont des signes, à qui appartiennent ces signes importe encore plus ». L’exemplaire conservé à la bibliothèque Vaticane présente l’une des rares reliures connues aux armes de Jules III\teiNote[n={35},place={foot},xmlid={ftn69-2}]{\teiP[xmlid={p-092}]{BAV, Stamp. De Marinis 82.}}.}
                \teiP[xmlid={p40}]{À Rome, en ce qui regarde l’art du manuscrit, deux artistes se sont distingués dans l’ornement héraldique en l’honneur de Jules III : Ferdinando Ruano et Vincent Raymond. Le premier a pratiqué l’extraction et la diffusion des meubles, le second a exploité la devise suivant des modes similaires à ceux du décor monumental. En ce qui regarde l’art de l’imprimé, deux éditeurs romains se remarquent : Antonio Blado et les frères Dorico. Eux seuls ont conçu des vignettes qui ont servi pour plusieurs ouvrages au point de devenir des classiques, l’ovale des Dorico et la couronne de laurier « \teiHi[rend={italic}]{Diu felix} » de Blado qui durera plus d’un siècle, jusqu’à Grégoire XIII. La présence de l’héraldique s’explique chez Blado par le caractère officiel de la plupart des pièces et chez les Dorico par une intention dédicatoire. Pour le reste, à Rome et hors de Rome, des vignettes ont servi pour une, parfois deux éditions généralement en lien avec des stratégies de dédicace au pontife. Dans la plupart, on remarque une tentative de rendre la différence des couleurs par les hachures de la gravure, procédé qui se systématisera plus tard. Après le manuscrit et son héritage médiéval, le livre imprimé demeure un support de l’image héraldique même si elle se trouve en concurrence avec les marques de libraires.}
                \begin{teiFigure}[xmlid={figure01-2}]

                  \teiGraphic[url={../icono/br/Ch03\_Loskoutoff\_1/fig9.jpg}]
                  \teiHead[xmlid={figure-8-ferdinando-ruano-registre-des-depenses-bav-vat-lat-3968-fol-1r-001}]{Figure 8. Ferdinando Ruano, \teiHi[rend={italic}]{Registre des dépenses}, BAV, Vat. lat. 3968, fol. 1r.}
                
\end{teiFigure}
                \begin{teiFigure}[xmlid={figure02-2}]

                  \teiGraphic[url={../icono/br/Ch03\_Loskoutoff\_1/fig10.jpg}]
                  \teiHead[xmlid={figure-9-ferdinando-ruano-lettres-de-leon-ix-1551-bav-vat-lat-3790-fol-3v-001}]{Figure 9. Ferdinando Ruano, \teiHi[rend={italic}]{Lettres de Léon IX}, 1551, BAV, Vat. lat. 3790, fol. 3v.}
                
\end{teiFigure}
                \begin{teiFigure}[xmlid={figure03-2}]

                  \teiGraphic[url={../icono/br/Ch03\_Loskoutoff\_1/fig11.jpg}]
                  \teiHead[xmlid={figure-10-ferdinando-ruano-lettres-dhildebert-eveque-du-mans-1551-bav-vat-lat-3841-fol-1r-001}]{Figure 10. Ferdinando Ruano, \teiHi[rend={italic}]{Lettres d’Hildebert, évêque du Mans}, 1551, BAV, Vat. lat. 3841, fol. 1r.}
                
\end{teiFigure}
                \begin{teiFigure}[xmlid={figure04-2}]

                  \teiGraphic[url={../icono/br/Ch03\_Loskoutoff\_1/fig12.jpg}]
                  \teiHead[xmlid={figure-11-ferdinando-ruano-saint-jerome-exposition-sur-les-psaumes-1554-bav-vat-lat-317-fol-2v-001}]{Figure 11. Ferdinando Ruano, saint Jérôme, \teiHi[rend={italic}]{Exposition sur les Psaumes}, 1554, BAV, Vat. lat. 317, fol. 2v.}
                
\end{teiFigure}
                \begin{teiFigure}[xmlid={figure05-2}]

                  \teiGraphic[url={../icono/br/Ch03\_Loskoutoff\_1/fig13.jpg}]
                  \teiHead[xmlid={figure-12-missel-bav-ms-capp-sist-154-fol-3v-detail-001}]{Figure 12. Missel, BAV, ms. Capp.Sist.154, fol. 3v (détail).}
                
\end{teiFigure}
                \begin{teiFigure}[xmlid={figure06-2}]

                  \teiGraphic[url={../icono/br/Ch03\_Loskoutoff\_1/fig14.jpg}]
                  \teiHead[xmlid={figure-13-missel-bav-ms-capp-sist-154-fol-4r-detail-001}]{Figure 13. Missel, BAV, ms. Capp.Sist.154, fol. 4r (détail).}
                
\end{teiFigure}
                \begin{teiFigure}[xmlid={figure07-2}]

                  \teiGraphic[url={../icono/br/Ch03\_Loskoutoff\_1/fig15.jpg}]
                  \teiHead[xmlid={figure-14-officium-ad-tertiam-bav-ms-capp-sist-67-fol-4v-detail-001}]{Figure 14. \teiHi[rend={italic}]{Officium ad Tertiam}, BAV, ms. Capp.Sist.67, fol. 4v (détail).}
                
\end{teiFigure}
                \begin{teiFigure}[xmlid={figure08}]

                  \teiGraphic[url={../icono/br/Ch03\_Loskoutoff\_1/fig16.jpg}]
                  \teiHead[xmlid={figure-15-officium-ad-tertiam-bav-ms-capp-sist-67-fol-5r-detail-001}]{Figure 15. \teiHi[rend={italic}]{Officium ad Tertiam}, BAV, ms. Capp.Sist.67, fol. 5r (détail).}
                
\end{teiFigure}
                \begin{teiFigure}[xmlid={figure09}]

                  \teiGraphic[url={../icono/br/Ch03\_Loskoutoff\_1/fig17\_9-cappsist67-7v.jpg}]
                  \teiHead[xmlid={figure-16-officium-ad-tertiam-bav-ms-capp-sist-67-fol-7v-detail-001}]{Figure 16. \teiHi[rend={italic}]{Officium ad Tertiam}, BAV, ms. Capp.Sist.67, fol. 7v (détail).}
                
\end{teiFigure}
                \begin{teiFigure}[xmlid={figure10}]

                  \teiGraphic[url={../icono/br/Ch03\_Loskoutoff\_1/fig18\_10-cappsist67-6v.jpg}]
                  \teiHead[xmlid={figure-17-officium-ad-tertiam-bav-ms-capp-sist-67-fol-6v-detail-001}]{Figure 17. \teiHi[rend={italic}]{Officium ad Tertiam}, BAV, ms. Capp.Sist.67, fol. 6v (détail).}
                
\end{teiFigure}
                \begin{teiFigure}[xmlid={figure11}]

                  \teiGraphic[url={../icono/br/Ch03\_Loskoutoff\_1/fig19\_11-cappsist213.jpg}]
                  \teiHead[xmlid={figure-18-vincent-raymond-et-apollonio-de-bonfratelli-lectionnaire-bav-ms-capp-sist-213-folio-de-titre-001}]{Figure 18. Vincent Raymond et Apollonio de’ Bonfratelli, Lectionnaire, BAV, ms. Capp.Sist.213, folio de titre.}
                
\end{teiFigure}
                \begin{teiFigure}[xmlid={figure12}]

                  \teiGraphic[url={../icono/br/Ch03\_Loskoutoff\_1/fig20.jpg}]
                  \teiHead[xmlid={figure-19-prospero-fontana-et-atelier-devise-de-jules-iii-portique-de-la-villa-giulia-001}]{Figure 19. Prospero Fontana et atelier, devise de Jules III, portique de la villa Giulia.}
                
\end{teiFigure}
                \begin{teiFigure}[xmlid={figure13}]

                  \teiGraphic[url={../icono/br/Ch03\_Loskoutoff\_1/fig21.jpg}]
                  \teiHead[xmlid={figure-20-vincent-raymond-attribue-a-evangeliaire-de-jules-iii-bibliotheque-capitulaire-de-tolede-premier-folio-001}]{Figure 20. Vincent Raymond (attribué à), \teiHi[rend={italic}]{Évangéliaire de Jules III}, Bibliothèque capitulaire de Tolède, premier folio.}
                
\end{teiFigure}
                \begin{teiFigure}[xmlid={figure14}]

                  \teiGraphic[url={../icono/br/Ch03\_Loskoutoff\_1/fig22.jpg}]
                  \teiHead[xmlid={figure-21-missel-du-cardinal-girolamo-dandini-bibliotheque-capitulaire-de-tolede-premier-folio-001}]{Figure 21. Missel du cardinal Girolamo Dandini, Bibliothèque capitulaire de Tolède, premier folio.}
                
\end{teiFigure}
                \begin{teiFigure}[xmlid={figure15}]

                  \teiGraphic[url={../icono/br/Ch03\_Loskoutoff\_1/fig23.jpg}]
                  \teiHead[xmlid={figure-22-vignette-de-titre-pour-antonio-massa-de-exercitatione-iurisperitorum-rome-1550-bnf-001}]{Figure 22. Vignette de titre pour Antonio Massa,\teiHi[rend={italic}]{ De exercitatione iurisperitorum}, Rome, 1550, BNF.}
                
\end{teiFigure}
                \begin{teiFigure}[xmlid={figure16}]

                  \teiGraphic[url={../icono/br/Ch03\_Loskoutoff\_1/fig24.jpg}]
                  \teiHead[xmlid={figure-23-vignette-de-titre-pour-giovanni-girolamo-albani-erudita-atque-luculenta-disputatio-rome-1553-biblioteca-general-historica-de-la-universidad-de-salamanca-espagne-001}]{Figure 23. Vignette de titre pour Giovanni Girolamo Albani, \teiHi[rend={italic}]{Erudita atque luculenta disputatio, }Rome, 1553, Biblioteca General Histórica de la Universidad de Salamanca (Espagne).}
                
\end{teiFigure}
                \begin{teiFigure}[xmlid={figure17}]

                  \teiGraphic[url={../icono/br/Ch03\_Loskoutoff\_1/fig25.jpg}]
                  \teiHead[xmlid={figure-24-vignette-de-titre-pour-giovanni-battista-poiana-de-iobileo-et-indulgentiis-rome-1550-bibliotheque-casanatense-001}]{Figure 24. Vignette de titre pour Giovanni Battista Poiana, \teiHi[rend={italic}]{De iobileo et indulgentiis}, Rome, 1550, bibliothèque Casanatense.}
                
\end{teiFigure}
                \begin{teiFigure}[xmlid={figure18}]

                  \teiGraphic[url={../icono/br/Ch03\_Loskoutoff\_1/fig26.jpg}]
                  \teiHead[xmlid={figure-25-vignette-de-titre-pour-les-regulae-tres-s-d-n-d-iulii-diuina-prouidentia-papae-iii-rome-1550-bncr-001}]{Figure 25. Vignette de titre pour les \teiHi[rend={italic}]{Regulae tres S. D. N. D. Iulii diuina prouidentia Papae. III, }Rome, 1550, BNCR.}
                
\end{teiFigure}
                \begin{teiFigure}[xmlid={figure19}]

                  \teiGraphic[url={../icono/br/Ch03\_Loskoutoff\_1/fig27.jpg}]
                  \teiHead[xmlid={figure-26-vignette-de-titre-pour-silvestro-sigonio-de-concordia-servanda-rome-1550-bncr-001}]{Figure 26. Vignette de titre pour Silvestro Sigonio,\teiHi[rend={italic}]{ De concordia servanda,} Rome, 1550, BNCR.}
                
\end{teiFigure}
                \begin{teiFigure}[xmlid={figure20}]

                  \teiGraphic[url={../icono/br/Ch03\_Loskoutoff\_1/fig28.jpg}]
                  \teiHead[xmlid={figure-27-vignette-de-titre-pour-contra-hebreos-rome-1554-bibliotheque-casanatense-001}]{Figure 27. Vignette de titre pour \teiHi[rend={italic}]{Contra hebreos}…, Rome, 1554, bibliothèque Casanatense.}
                
\end{teiFigure}
                \begin{teiFigure}[xmlid={figure21}]

                  \teiGraphic[url={../icono/br/Ch03\_Loskoutoff\_1/fig29.jpg}]
                  \teiHead[xmlid={figure-28-vignette-pour-le-privilege-du-commentaire-deustathius-sur-l-iliade-vol-iv-rome-1550-bibliotheque-sainte-genevieve-fol-y-7-inv-11-res-001}]{Figure 28. Vignette pour le privilège du commentaire d’Eustathius sur l’\teiHi[rend={italic}]{Iliade, }vol. IV, Rome, 1550, bibliothèque Sainte-Geneviève, FOL Y 7 INV 11 RES.}
                
\end{teiFigure}
                \begin{teiFigure}[xmlid={figure22}]

                  \teiGraphic[url={../icono/br/Ch03\_Loskoutoff\_1/fig30.jpg}]
                  \teiHead[xmlid={figure-29-vignette-de-titre-pour-uberto-foglietta-in-festo-omnium-beatorum-rome-1553-universite-de-turin-001}]{Figure 29. Vignette de titre pour Uberto Foglietta, \teiHi[rend={italic}]{In festo omnium Beatorum, }Rome, 1553, Université de Turin.}
                
\end{teiFigure}
                \begin{teiFigure}[xmlid={figure23}]

                  \teiGraphic[url={../icono/br/Ch03\_Loskoutoff\_1/fig31.png}]
                  \teiHead[xmlid={figure-30-vignette-de-titre-pour-les-novae-regvlae-cancellariae-apostolicae-rome-1552-bncr-001}]{Figure 30. Vignette de titre pour les \teiHi[rend={italic}]{Novae regvlae cancellariae apostolicae, }Rome, 1552, BNCR.}
                
\end{teiFigure}
                \begin{teiFigure}[xmlid={figure24}]

                  \teiGraphic[url={../icono/br/Ch03\_Loskoutoff\_1/fig32.jpg}]
                  \teiHead[xmlid={figure-31-lettrine-pour-ambrosio-catharino-enarrationes-in-quinque-priora-capita-libri-geneseos-rome-1552-national-central-library-of-rome-001}]{Figure 31. Lettrine pour Ambrosio Catharino, \teiHi[rend={italic}]{Enarrationes in quinque priora capita libri geneseos}, Rome, 1552, National Central Library of Rome.}
                
\end{teiFigure}
                \begin{teiFigure}[xmlid={figure25}]

                  \teiGraphic[url={../icono/br/Ch03\_Loskoutoff\_1/fig33.jpg}]
                  \teiHead[xmlid={figure-32-vignette-pour-bullae-diuersorum-pontificum-rome-1550-bncr-001}]{Figure 32. Vignette pour \teiHi[rend={italic}]{Bullae diuersorum pontificum}, Rome, 1550, BNCR.}
                
\end{teiFigure}
                \begin{teiFigure}[xmlid={figure26}]

                  \teiGraphic[url={../icono/br/Ch03\_Loskoutoff\_1/fig34.jpg}]
                  \teiHead[xmlid={figure-33-vignette-pour-bartolomej-georgijevic-libellus-vere-christiana-lectione-rome-1552-bncr-001}]{Figure 33. Vignette pour Bartolomej Georgijevic, \teiHi[rend={italic}]{Libellus vere Christiana lectione…, }Rome, 1552, BNCR.}
                
\end{teiFigure}
                \begin{teiFigure}[xmlid={figure27}]

                  \teiGraphic[url={../icono/br/Ch03\_Loskoutoff\_1/fig35.jpg}]
                  \teiHead[xmlid={figure-34-ioannes-magnus-historia-de-omnibus-gothorum-sueonumque-regibus-rome-1554-bnf-001}]{Figure 34. Ioannes Magnus, \teiHi[rend={italic}]{Historia de omnibus Gothorum Sueonumque regibus,} Rome, 1554, BNF.}
                
\end{teiFigure}
                \begin{teiFigure}[xmlid={figure28}]

                  \teiGraphic[url={../icono/br/Ch03\_Loskoutoff\_1/fig36.jpg}]
                  \teiHead[xmlid={figure-35-page-de-titre-pour-palestrina-missarum-liber-primus-rome-1554-museo-internazionale-e-biblioteca-della-musica-bologne-001}]{Figure 35. Page de titre pour Palestrina, \teiHi[rend={italic}]{Missarum liber primus}, Rome, 1554, Museo internazionale e biblioteca della musica, Bologne.}
                
\end{teiFigure}
                \begin{teiFigure}[xmlid={figure29}]

                  \teiGraphic[url={../icono/br/Ch03\_Loskoutoff\_1/fig37.jpg}]
                  \teiHead[xmlid={figure-36-lettrine-pour-palestrina-missarum-liber-primus-rome-1554-bibliotheque-du-conservatoire-de-sainte-cecile-001}]{Figure 36. Lettrine pour Palestrina, \teiHi[rend={italic}]{Missarum liber primus}, Rome, 1554, Bibliothèque du conservatoire de Sainte-Cécile.}
                
\end{teiFigure}
                \begin{teiFigure}[xmlid={figure30}]

                  \teiGraphic[url={../icono/br/Ch03\_Loskoutoff\_1/fig38.jpg}]
                  \teiHead[xmlid={figure-37-vignette-de-titre-pour-prospero-danza-la-elettion-del-beatissimo-summo-pontifice-iulio-tertio-venise-1550-bav-001}]{Figure 37. Vignette de titre pour Prospero Danza, \teiHi[rend={italic}]{La elettion del beatissimo \& summo pontifice Iulio tertio}, Venise, 1550, BAV.}
                
\end{teiFigure}
                \begin{teiFigure}[xmlid={figure31}]

                  \teiGraphic[url={../icono/br/Ch03\_Loskoutoff\_1/fig39.png}]
                  \teiHead[xmlid={figure-38-vignette-de-titre-pour-giacomo-naccchiante-enarrationes-piae-doctae-et-catholicae-in-epistolam-pauli-ad-ephesios-venise-1553-biblioteca-civica-bertoliana-001}]{Figure 38. Vignette de titre pour Giacomo Naccchiante, \teiHi[rend={italic}]{Enarrationes piae, doctae, et catholicae : in epistolam Pauli ad Ephesios, }Venise, 1553, Biblioteca civica Bertoliana.}
                
\end{teiFigure}
                \begin{teiFigure}[xmlid={figure32}]

                  \teiGraphic[url={../icono/br/Ch03\_Loskoutoff\_1/fig40.jpg}]
                  \teiHead[xmlid={figure-39-vignette-de-titre-pour-tommaso-giannoti-rangoni-de-vita-hominis-ultra-cxx-annos-protrahenda-venise-1553-bnf-001}]{Figure 39. Vignette de titre pour Tommaso Giannoti Rangoni, \teiHi[rend={italic}]{De vita hominis ultra CXX. annos protrahenda, }Venise, 1553, BNF.}
                
\end{teiFigure}
                \begin{teiFigure}[xmlid={figure33}]

                  \teiGraphic[url={../icono/br/Ch03\_Loskoutoff\_1/fig41.jpg}]
                  \teiHead[xmlid={figure-40-vignette-de-titre-pour-lorenzo-davidico-anotomia-delli-vitii-florence-1550-bibliotheque-casanatense-001}]{Figure 40. Vignette de titre pour Lorenzo Davidico, \teiHi[rend={italic}]{Anotomia delli vitii, }Florence, 1550, bibliothèque Casanatense.}
                
\end{teiFigure}
                \begin{teiFigure}[xmlid={figure34}]

                  \teiGraphic[url={../icono/br/Ch03\_Loskoutoff\_1/fig42.jpg}]
                  \teiHead[xmlid={figure-41-vignette-de-titre-pour-lorenzo-dal-gambaro-sclarici-de-concilio-liber-i-bologne-1551-bibliotheque-casanatense-001}]{Figure 41. Vignette de titre pour Lorenzo Dal Gambaro Sclarici, \teiHi[rend={italic}]{De concilio liber I, }Bologne, 1551, bibliothèque Casanatense.}
                
\end{teiFigure}
                \begin{teiFigure}[xmlid={figure35}]

                  \teiGraphic[url={../icono/br/Ch03\_Loskoutoff\_1/fig43.png}]
                  \teiHead[xmlid={figure-42-vignette-de-titre-pour-bartolomeo-maggi-de-vulnerum-sclopetorum-et-bombardarum-curatione-bologne-1552-bncr-001}]{Figure 42. Vignette de titre pour Bartolomeo Maggi, \teiHi[rend={italic}]{De vulnerum sclopetorum et bombardarum curatione, }Bologne, 1552, BNCR.}
                
\end{teiFigure}
              
\end{teiDiv}
            
\end{teiDiv}
          
\end{teiElement}
        
\end{teiElement}
        \begin{teiElement}[name={group},type={chapter},data-page-title={Érudits et héraldique dans la Rome d’Urbain VIII (1623-1644)},data-page-authors={Fabrizio Federici},data-include-href={Ch04\_Erudits\_et\_heraldique\_Rome.xml},xmlid={chapter-001}]

          \begin{teiElement}[name={front}]

            \begin{teiDiv}[type={titlePage},xmlid={titlepage-006}]

              \teiP[rend={title-main},xmlid={p-093}]{Érudits et héraldique dans la Rome d’Urbain VIII (1623-1644)}
              \teiP[rend={author-aut},xmlid={p-094}]{Fabrizio Federici}
              \teiP[rend={editor-trl},xmlid={p-095}]{Traduit de l’italien par Yvan Loskoutoff.}
            
\end{teiDiv}
          
\end{teiElement}
          \begin{teiElement}[name={body},xmlid={body-006}]

            \teiP[xmlid={p1-6}]{L’intérêt pour l’héraldique dans la Rome des Barberini fut, comme on pouvait s’y attendre, celui d’une ville envahie par l’essaim vivace, varié et multiforme des abeilles armoriales du pontife régnant Urbain VIII (1623-1644\teiNote[n={1},place={foot},xmlid={ftn1-4}]{\teiP[xmlid={p-096}]{Sur les recherches généalogiques et héraldiques à Rome, en particulier sur celles de la deuxième moitié du \teiHi[rend={small-caps}]{xvi}\teiHi[rend={sup}]{e} siècle, objet de la présente contribution, voir Andreas Rehberg\teiHi[rend={italic}]{, }« Collecting and Drawing against Oblivion. Panvinio, Ceccarelli and Chacón and their Search for the Genealogical-Heraldic Identity of the Families of Rome », dans Noëlle-Laetitia Perret et Hans-Joachim Schmidt (dir.), \teiHi[rend={italic}]{Memories Lost in the Middle Ages. Collective Forgetting as an Alternative Procedure of Social Cohesion}, Turnhout, Brepols, 2023, p. 295-324. Sur les abeilles des Barberini, voir Louise Rice, « Branding, Barberini-style », dans Maurizia Cicconi, Flaminia Gennari Santori et Sebastian Schütze (dir.), \teiHi[rend={italic}]{L’immagine sovrana. Urbano VIII e i Barberini}, catalogue d’exposition, Rome, Officina libraria, 2023, p. 245-250.}}). Cet intérêt se manifesta plus largement dans la production bibliographique relative au blason que nous évoquerons. Il fut notamment le fait de l’antiquaire Francesco Gualdi qui consacra un traité aux pierres tombales mais qui œuvra aussi pour la conservation des vestiges armoriaux. Enfin, à Rome, la question se pose des origines commerçantes de la noblesse, ayant des répercussions dans la symbolique armoriale.}
            \begin{teiDiv}[type={section1},xmlid={div01-3}]

              \teiHead[xmlid={publications-de-sujet-heraldique-gozze-et-pietrasanta-001}]{Publications de sujet héraldique : Gozze et Pietrasanta}
              \teiP[xmlid={p2-5}]{Sur le marché éditorial on dénombre des publications à sujet héraldique de types différents, depuis les petits volumes jusqu’aux tomes copieux. Parmi les premiers, figure l’opuscule \teiHi[rend={italic}]{Se dalle armi, o insegne, che parlano, overo se da’ corpi delle armi, che rappresentano i cognomi, si possi argomentare ignobiltà in quella fameglia, che le usa \lbrack{}...\rbrack{}}, dû à l’antiquaire de Pesaro Gauges de’ Gozze et publié chez l’éditeur Vitale Mascardi en 1637\teiNote[n={2},place={foot},xmlid={ftn48-3}]{\teiP[xmlid={p-097}]{Gauges de’ Gozze, \teiHi[rend={italic}]{Se dalle armi, o insegne, che parlano, overo se da’ corpi delle armi, che rappresentano i cognomi, si possi argomentare ignobiltà in quella fameglia, che le usa \lbrack{}...\rbrack{}} \lbrack{}\teiHi[rend={italic}]{Si des armes ou insignes parlants, ou si du corps des armes qui représentent les noms, on peut déduire la noblesse de la famille qui en use..}.\rbrack{}, In Roma, Per Vitale Mascardi, 1637. Gozze fut un antiquaire de vastes intérêts, qui écrivit et publia sur l’héraldique mais aussi sur l’épigraphie, sur l’histoire antique, la linguistique, la philosophie. Il fut un des collaborateurs à la rédaction du traité des \teiHi[rend={italic}]{Memorie sepolcrali}, entrepris par son protecteur Francesco Gualdi (voir \teiHi[rend={italic}]{infra}). Sur lui voir, outre l’introduction à l’édition critique des \teiHi[rend={italic}]{Memorie} préparée par nos soins et en cours de publication, Ingo Herklotz, \teiHi[rend={italic}]{Cassiano Dal Pozzo und die Archäologie des 17. Jahrhunderts}, Munich, Hirmer, 1999, p. 289 et Anna Polo, \teiHi[rend={italic}]{Las }Annotationi \teiHi[rend={italic}]{de Gauges de‘ Gozze (1631) en relación con la obra didáctica de Franciosini}, Padoue, CLEUP, 2021.}}. Gozze voulut dédier cette œuvre au collectionneur Francesco Gualdi, son protecteur, non seulement pour « la profonde affection que par votre bonté vous me portez et parce que vous avez daigné plusieurs fois me signifier par des effets que vous me tenez en une certaine estime, moi et tout ce qui me regarde\teiNote[n={3},place={foot},xmlid={ftn49-3}]{\teiP[xmlid={p-098}]{Gauges de’ Gozze, \teiHi[rend={italic}]{Se dalle armi, o insegne},\teiHi[rend={italic}]{ op. cit.}, p. 3 (lettre dédicatoire, datée du 22 juillet 1637) : « \teiHi[rend={italic}]{molto affetto, che per sua bontà mi porta, e che più volte si è con effetti degnat}\lbrack{}\teiHi[rend={italic}]{o}\rbrack{}\teiHi[rend={italic}]{ significarmi di tener in qualche stima me, e le cose mie. }»}} », mais surtout pour le fait que c’était ce gentilhomme qui avait convaincu le jeune homme de composer cette dissertation, au cours d’une « grave maladie » qui l’affligea comme il le rapporte dans l’avis au lecteur :}
              \begin{teiCit}[xmlid={cit01-2}]

                \begin{teiQuote}
\lbrack{}...\rbrack{} Étant né un désaccord entre certains gentilshommes de solide et forte littérature et moi au sujet des armes parlantes, eux affirmant qu’elles sont des indices de familles non nobles et moi le niant constamment, nous restâmes alors séparés de part et d’autre par les raisons que nous avions adoptées \lbrack{}...\rbrack{}. Il arriva ensuite que je fus réduit à l’extrémité de la vie par la véhémence de très cruelles douleurs \lbrack{}...\rbrack{} lors de cette grave maladie, je fus plusieurs fois visité avec humanité par monsieur le chevalier Gualdi et entre plusieurs discours il me revint en mémoire le raisonnement que nous avions eu sur les armes et je lui donnai ma parole que si je me relevais de mon mal je développerai mon point de vue sur une telle matière, ce que j’exécutais ensuite en l’espace de deux jours et c’est cela qui t’est maintenant présenté\teiNote[n={4},place={foot},xmlid={ftn50-3}]{\teiP[xmlid={p-099}]{\teiHi[rend={italic}]{Ibid}., p. 6 (\teiHi[rend={italic}]{A chi legge }\lbrack{}À qui lit\rbrack{}) : « \lbrack{}…\rbrack{} \teiHi[rend={italic}]{essendo nato disparere fra certi gentilhuomini di soda e granita litteratura, e me intorno alle armi, che parlano, ed affermando essi esser elleno inditij di fameglie non nobili, ed io ciò costantemente negando, rimanemmo per allhora dell’una, e dell’altra parte per le ragioni addotte }\lbrack{}...\rbrack{}\teiHi[rend={italic}]{. Succedé poi, ch’io mi riducessi all’estremo della vita, per la veemenza di crudelissimi dolori }\lbrack{}...\rbrack{}\teiHi[rend={italic}]{ in questa grave infermità fui più volte humanamente visitato dal signor cavalier Gualdi, e fra diversi discorsi egli mi tornò a memoria l’havuto ragionamento delle armi, ed io li diedi parola, che rihavuto che mi fossi dal male, harrei diste\lbrack{}s\rbrack{}o il mio parere intorno a sì fatta materia, il che poi nello spatio di due giorni eseguij, ed è questo, ch’hora ti si presenta. }»}}.
\end{teiQuote}
              
\end{teiCit}
              \teiP[xmlid={p3-5}]{Le raisonnement de Gozze pour défendre les armes « parlantes », où le nom de la maison qui les possède se traduit en image, est le suivant : comme tous les écus, excepté ceux qui sont nés de la fantaisie d’un caprice personnel, peuvent se dire parlants parce qu’ils nous communiquent des informations sur les familles qui en usent, « que les corps des armes qui expriment les patronymes soient arguments de noblesse dans la famille qui l’utilise, je ne sais comment le persuader facilement et par une simple affirmation : bon nombre d’exemples nous forcent à suivre l’avis contraire \lbrack{}...\rbrack{}. Dans de nombreux cas on utilisa des choses viles et abjectes pour démontrer qu’on s’était élevé de cette bassesse par sa propre vertu\teiNote[n={5},place={foot},xmlid={ftn51-3}]{\teiP[xmlid={p-100}]{\teiHi[rend={italic}]{Ibid.}, p. 24 : « \teiHi[rend={italic}]{che i corpi delle armi, che esprimono i cognomi, sieno argomento d’ignobiltà in quella fameglia, che l’usa, io non so come mi si possa agevolmente, e con un semplice attestato persuadere: havemo esempij in buon numero che ’l contrario parere ci sforzano a seguitare }\lbrack{}...\rbrack{}\teiHi[rend={italic}]{. In molti usarono cose vili, ed abbiette in dimostramento che la propia virtù da cotali bassezze gli haveva sollevati }\lbrack{}...\rbrack{}\teiHi[rend={italic}]{ }».}} \lbrack{}...\rbrack{} ». Dans la lettre dédicatoire \teiHi[rend={italic}]{À Monsieur Francesco Gualdi chevalier de Saint Stéphane de Rimini}, il est rappelé que le tombeau de la famille Gualdi dans le cloître du Temple de Malatesta était surmonté d’une fresque aujourd’hui perdue « de Giotto », c’est-à-dire d’un suiveur de l’artiste à Rimini : « qu’elle \lbrack{}la famille Gualdi\rbrack{} était originaire d’Allemagne, un indice évident en est l’antique motto qui se voit sur la pierre tombale des Gualdi dans le cloître de saint François à Rimini sous la peinture de Giotto, en caractères allemands \teiHi[rend={italic}]{Anaupe poper} : « et elle se glorifie que dans sa patrie, depuis 400 ans et plus, elle a toujours été une famille de renom\teiNote[n={6},place={foot},xmlid={ftn52-3}]{\teiP[xmlid={p-101}]{\teiHi[rend={italic}]{Ibid.}, p. 4 : « \teiHi[rend={italic}]{la qual \lbrack{}scil. famiglia Gualdi\rbrack{} habbia origine da Germania, è grand’inditio l’antichissimo motto, che si vede nella lapide sepolcrale de’ Gualdi nel claustro di San Francesco di Rimini a piè la pittura di Giotto, in carattere todesco Anaupe poper: ed ella si gloria che nella patria pel corso di 400 e più anni sia stata sempre fameglia graduata }\lbrack{}...\rbrack{}\teiHi[rend={italic}]{. }» Les Gualdi figurent aussi dans la liste des « très nobles familles titrées comme privées, dont les corps des armes représentent vivement leurs noms \lbrack{}\teiHi[rend={italic}]{fameglie nobilissime sì titolate, come private, i corpi delle armi delle quali rappresentano vivamente i cognomi di esse}\rbrack{} ». Cette famille porte « un lion qui parle également car en langue d’Allemagne (dont la famille est originaire) Guualt ou Geuualt \lbrack{}Gewalt\rbrack{} signifie pouvoir, ce qui est représenté par cet animal \lbrack{}\teiHi[rend={italic}]{un leone, che parimente parla, percioché in lingua todesca (donde la fameglia è originaria) Guualt, o Geuualt \lbrack{}scil. Gewalt\rbrack{} significa potenza, il che viene rappresentato con quell’animale}\rbrack{} » (\teiHi[rend={italic}]{ibid}., p. 25).}} ». À la page de titre de l’opuscule sont reproduites, conjointement aux armes du chevalier Francesco, celles qui figurent sur la tombe des Gualdi à Rimini, surmontées du \teiHi[rend={italic}]{motto} mystérieux inséré dans le losange quadrilobé qui devait encadrer l’écu sur la pierre tombale (fig. 43). Les deux écus sont cernés de la bordure indentée ou, comme on disait fréquemment au \teiHi[rend={small-caps}]{xvii}\teiHi[rend={sup}]{e} siècle, de la « ceinture militaire \lbrack{}\teiHi[rend={italic}]{cingolo militare}\rbrack{} », souvent concédée par les Malatesta aux familles qui leur étaient fidèles.}
              \teiP[xmlid={p4-5}]{À peu près contemporain du traité sur les « armes parlantes », il existe un autre écrit de Gozze à sujet héraldique, de nettement plus grande ampleur, le traité \teiHi[rend={italic}]{Dell’antichità delle armi, o insegne delle famiglie }\lbrack{}\teiHi[rend={italic}]{De l’antiquité des armes ou insignes des familles}\rbrack{}, dédié au cardinal Francesco Barberini et resté manuscrit. Cette œuvre de la main de Gozze constitue le manuscrit It. 364 de la Bibliothèque nationale de France. L’antiquaire de Pesaro la composa en 1635-1636 : le 16 janvier 1636 il en envoyait à Peiresc une « esquisse \lbrack{}\teiHi[rend={italic}]{schizzo}\rbrack{} » (à identifier avec le manuscrit parisien ?), et quelques mois plus tard il sollicitait l’avis du Provençal, qui ne s’était pas encore prononcé sur ce traité\teiNote[n={7},place={foot},xmlid={ftn53-3}]{\teiP[xmlid={p-102}]{BNF, manuscrit Fr. 9542, fol. 100r-101r, lettre de Gozze à Peiresc du 3 avril 1636, fol. 100v : « \teiHi[rend={italic}]{A’ 16 gennaro passato mandai a Vostra Signoria Illustrissima uno schizzo, per così dire, del mio trattato dell’antichità delle armi per attenderne da lei il suo giudicio, non tengo però finhora alcun avviso del ricapito. Non entro a trattar delle armi delle fameglie moderne, ovvero viventi per non tirarmi adosso l’odio di altre che pretendono d’esser eguali alle più illustri} ». Sur les rapports entre le savant de Pesaro et le grand érudit provençal, voir Fabrizio Federici, « Alla ricerca dell’esattezza: Peiresc, Francesco Gualdi e l’antico », Marc Bayard (dir.), \teiHi[rend={italic}]{Rome-Paris 1640. Transferts culturels et renaissance d’un centre artistique}, actes du colloque, Rome, villa Médicis, 17-19 avril 2008, Paris, Somogy, 2010, p. 229-273, \teiHi[rend={italic}]{passim}.}}. Avec ampleur d’argumentation et abondance de citations d’auteurs antiques, Gozze y affirme que l’origine des armoiries remonte à l’antiquité classique, et donc à une époque très antérieure à ce que soutenaient divers auteurs : « Disons donc que l’usage des armes, c’est-à-dire des insignes des maisons et familles, est non seulement plus ancien que le temps de l’empereur Charlemagne et du roi Artus de Bretagne, comme le pensent certains auteurs modernes, mais ancien au-delà de toute mémoire et même s’il ne s’en trouve pas fait spécifiquement et expressément mention chez les auteurs de l’ancienne langue latine \lbrack{}...\rbrack{} il s’y trouve au moins en général et tacitement mais avec des signes clairs et évidents\teiNote[n={8},place={foot},xmlid={ftn54-3}]{\teiP[xmlid={p-103}]{BNF, manuscrit It. 364, fol. 44r-v : « \teiHi[rend={italic}]{Diciamo dunque, che l'uso delle armi, cioè delle insegne delle Case, e famiglie è non solo più antico del tempo dell'imperador Carlo Magno, e del re Artù di Brettagna, come sono di parere alcuni scrittori moderni, ma antichissimo affatto oltre ogni memoria: e quantunque non si truovi specificamente, ed espressamente fatta mentione di esse appo gli autori della lingua latina antica }\lbrack{}...\rbrack{} \teiHi[rend={italic}]{sì vi si truova almeno in generale, e tacitamente, ma però con chiari, ed evidenti segnali }\lbrack{}...\rbrack{}. »}} \lbrack{}...\rbrack{} ».}
              \teiP[xmlid={p5-5}]{La valeur du traité doit néanmoins être relativisée car il s’agit en bonne partie du plagiat d’un écrit de sujet héraldique dû à un autre antiquaire, en certains endroits sensiblement amplifié. L’œuvre que Gozze a pillée, reprenant des passages entiers pratiquement tels quels et s’en servant comme canevas pour son argumentation, est le traité \teiHi[rend={italic}]{Delle antichità delle armi gentilizie }\lbrack{}\teiHi[rend={italic}]{De l’antiquité des armes nobiliaires}\rbrack{}, composé à la fin du \teiHi[rend={small-caps}]{xiv}\teiHi[rend={sup}]{e} siècle par le grand philologue et historien de la langue Celso Cittadini. Longtemps resté inédit, il fut publié seulement en 1741 à Lucques, par les soins de Giovan Girolamo Carli\teiNote[n={9},place={foot},xmlid={ftn55-3}]{\teiP[xmlid={p-104}]{\teiHi[rend={italic}]{Delle antichità delle armi gentilizie, trattato di Celso Cittadini, colle annotazioni di Giovan Girolamo Carli}, Lucca, Per Salvatore e Giandomenico Marescandoli, 1741. Le passage copié par Gozze, dont la référence est à peine donnée dans le texte, se trouve aux p. 20-22 de cette édition.}}. Gozze dut connaître l’œuvre manuscrite entre 1631 et 1634, au cours de son séjour à Sienne, ville où avait vécu et travaillé Cittadini, et il s’en servit sans aucun scrupule.}
              \teiP[xmlid={p6-5}]{Parmi les passages qui ne proviennent pas de Cittadini figure celui dans lequel Gozze, se servant de certaines figurations sculptées sur la Colonne Trajane, fait remonter à une habitude des soldats romains l’usage de mettre des cimiers au-dessus des écus : « Sur cette colonne j’observe pour la première fois que l’usage de placer des morions, des salades ou des heaumes sur les écus de nos armoiries tire son origine de la coutume que les soldats avaient d’appuyer l’écu à un pal et d’y poser ensuite le morion, alors qu’ils s’occupaient à faire des abris, ou qu’ils se reposaient, restant néanmoins armés, comme cela s’y voit exprimé\teiNote[n={10},place={foot},xmlid={ftn56-3}]{\teiP[xmlid={p-105}]{BNF, manuscrit It. 364, fol. 64r : « \teiHi[rend={italic}]{Da questa colonna osservo il primo fra tutti l'usanza del porre i morioni, o celate, od elmi sopra gli scudi delle nostre armi, trarre origine dal costume, che i soldati havevano di appoggiare ad un palo lo scudo, e mettervi poi sopra il morione, mentre essi s'affadigavano intorno a' ripari, ovvero riposavano, stando del rimanente armati, come ivi si vede espresso }\lbrack{}...\rbrack{} ». Une référence à ce passage se trouve dans l’un des écrits que Gozze composa pour les \teiHi[rend={italic}]{Memorie sepolcrali }(cf. \teiHi[rend={italic}]{infra}) : « Gauges de’ Gozze de Pesaro dans le livre qu’il fait sur les armoiries des familles porte l’opinion que l’usage du heaume au-dessus de l’écu provient de l’antiquité, se voyant dans la Colonne Trajane, et dans d’autres vestiges précédemment par lui observés, les soldats se reposant ou travaillant non loin à des réparations, ils appuyaient l’écu à un pal et sur lui plaçaient le heaume, d’où provient manifestement l’usage des armoiries avec le heaume ou cimier à notre époque \lbrack{}\teiHi[rend={italic}]{Gauges de’ Gozze da Pesaro nel suo libro che fa delle armi delle fameglie porta oppinione che l’uso dell’elmo sopra lo scudo sia tolto dall’antico, vedendosi nella Colonna Traiana, ed in altre memorie da lui prima osservate, che riposandosi i soldati, o pur lavorando intorno a’ ripari, ch’essi appoggiavan lo scudo ad un palo, e sopra di esso vi collocavan l’elmo, onde se ne cava manifestamente l’uso delle armi coll’elmo, o cimiero del tempo nostro}\rbrack{} » (BAV, manuscrit Vat. Lat. 8254, II, fol. 347r).}} \lbrack{}...\rbrack{} ».}
              \teiP[xmlid={p7-5}]{Une œuvre de plus grande ampleur, l’une des principales du siècle et pas seulement en contexte romain, est le traité du jésuite Silvestro Pietrasanta \teiHi[rend={italic}]{Tesserae gentilitiae }\lbrack{}\teiHi[rend={italic}]{Les armoiries nobiliaires}\rbrack{}, volume fondamental autant par la richesse de ses informations que par la mise au point de conventions graphiques pour la reproduction imprimée des écus\teiNote[n={11},place={foot},xmlid={ftn57-3}]{\teiP[xmlid={p-106}]{Silvestro Pietrasanta, \teiHi[rend={italic}]{Tesserae gentilitiae}, Romae, Typis haeredum Francisci Corbelletti, 1638.}} (fig. 44). L’ouvrage, où sont examinées diverses typologies d’écus, est ponctué de centaines de petits écus calcographiques des familles nobles, souvent romaines (fig. 45). Parmi les sources utilisées par Pietrasanta figurent des pierres tombales, des écus sur des édifices, des recueils manuscrits, mais pas seulement : on reproduit et décrit soigneusement le sceau du \teiHi[rend={small-caps}]{xiv}\teiHi[rend={sup}]{e} siècle de Giovanni di Vico, préfet de Rome\teiNote[n={12},place={foot},xmlid={ftn58-3}]{\teiP[xmlid={p-107}]{\teiHi[rend={italic}]{Ibid.}, p. 656 et 657.}}, et un grand relief est accordé en conclusion à une importante peinture perdue. Pietrasanta réserve les ultimes pages du traité à un argument déjà abordé au début de l’œuvre, soit l’origine supposée des écus dans les vêtements utilisés par les souverains et la noblesse. Comme illustration du lien étroit entre les écus et les ornements vestimentaires, le jésuite mentionne la pierre tombale disparue d’Angelotto Normanni († 1378), sur laquelle le défunt portait une armure ornée de ses propres symboles héraldiques (« \teiHi[rend={italic}]{partim liliato panno, et partim schemate vestium undulato} \lbrack{}de vêtements partie d’une bande fleurdelysée et partie d’un dessin ondulé\teiNote[n={13},place={foot},xmlid={ftn59-3}]{\teiP[xmlid={p-108}]{\teiHi[rend={italic}]{Ibid.}, p. 677, où la pierre tombale est mentionnée et décrite mais non reproduite, pour ne pas conclure l’œuvre, – est-il affirmé dans le texte – par la triste image d’une sépulture : « \teiHi[rend={italic}]{Demum ut in exordio \lbrack{}}scil. \teiHi[rend={italic}]{in exitu\rbrack{} etiam operis confirmem id, quod in illius exordio dixi: nempe gentilitijs tesseris plurimis et vetustissimis esse originem datam, ex vestium cultu domi forisque; possem subijcere oculis tumulum marmoreum, in quo spectatur miles indutus habitu quadripartito; nempe partim liliato panno, et partim schemate vestium undulato. Inscribitur vero semi-barbare:} Hic iacet Angeloctus \lbrack{}...\rbrack{}. \teiHi[rend={italic}]{Adijcitur vero ei hoc stemma. Sed ut finiam aspectu laetiore, atque ab imagine sepulchri funesta luctuosaque hoc loco abstineam} ». Une note marginale indique la localisation suivante de l’objet, originellement situé dans l’église de Sant’Antonio à Camaldoli, ou de San Niccolò de Forbitoribus, démolie en 1631 pour construire le sanctuaire jésuitique de Saint-Ignace : « La pierre tombale est au Collège romain \lbrack{}\teiHi[rend={italic}]{Lapis sepulchralis in Collegio Romano}\rbrack{} ». Sur cette belle pierre tombale, qui fut reproduite dans une des xylographies commandées par Francesco Gualdi pour le traité des \teiHi[rend={italic}]{Memorie sepolcrali }(fig. 6), voir Jörg Garms, Roswitha Juffinger et Bryan Ward-Perkins (dir.), \teiHi[rend={italic}]{Die mittelalterlichen Grabmäler in Rom und Latium vom 13. bis zum 15. Jahrhundert}: I, \teiHi[rend={italic}]{Die Grabplatten und Tafeln}, Rome-Vienne, Österreichische Akademie der Wissenschaften, 1981-1994, 2 vol., t. I, p. 235 et 236.}}\rbrack{} »), il mentionne aussi les représentations contenues « \teiHi[rend={italic}]{in perveteri tabula, quae extat hodieque, cum icone regis ac reginae Britannorum} \lbrack{}dans un très ancien tableau qu’on peut voir aujourd’hui portant l’image du roi et de la reine des Britanniques\rbrack{} ». Une note marginale révèle le lieu de conservation de cette œuvre perdue : à Rome, au collège des Anglais (\teiHi[rend={italic}]{Romae in Collegio Anglorum}). Pietrasanta identifie dans les deux personnages agenouillés figurant sur le panneau, qui portent des habits couverts de leurs emblèmes héraldiques, le roi d’Angleterre Richard II et la reine Anne de Bohême, et il publie deux petites eaux-fortes qui les montrent (fig. 46). Le roi est présenté à la Vierge Marie par son saint patron, Jean-Baptiste. Le sens de l’image est précisé par une inscription latine : « Ceci est ta dot, pieuse Vierge, afin que tu puisses la dominer \lbrack{}\teiHi[rend={italic}]{Dos tua Virgo pia haec est : quare rege Maria}\teiNote[n={14},place={foot},xmlid={ftn60-3}]{\teiP[xmlid={p-109}]{Silvestro Pietrasanta, \teiHi[rend={italic}]{Tesserae gentilitiae, op. cit.}, p. 677 et 678 : « \teiHi[rend={italic}]{Species alia offert se augusta et plena regiae maiestatis, in perveteri tabula, quae extat hodieque, cum icone regis ac reginae Britannorum, hoc vestium ornatu. }\lbrack{}effigie du roi agenouillé\rbrack{}\teiHi[rend={italic}]{ Amiciuntur ambo, tum liliato panno, tum veste leopardis picta. }\lbrack{}image de la reine\rbrack{}\teiHi[rend={italic}]{ Et regina insuper cyclade inornatur aquilis et auro intexta} \lbrack{}écu du roi\rbrack{} \teiHi[rend={italic}]{utrique vero consentanei indumentis }\lbrack{}écu de la reine\rbrack{}\teiHi[rend={italic}]{ clypei tesserarij, hoc est ex eis nati, adpinguntur. Videntur ij esse Riccardus II Rex: qui proximus regnavit ab Edoardo rege; a quo coepta sunt Francica lilia iungi leopardis Britannis; atque Anna Bohema uxor eius; quae cum aut filia aut soror fuerit Wenceslai Bohemi regis Romanorum, ac dein etiam imperatoris, aquila ideo insuper vestiebatur. Haerent vero nixi genibus, dum interim, interprete Divo Ioanne, regiam insulam Britanniae offerunt Deiparae Virgini, cum adscripta ea epigraphe, }Dos tua Virgo pia haec est: quare rege Maria\teiHi[rend={italic}]{. Ex avita nimirum regni olim religiosissimi sanctitate: ex qua heu quantum distat iam ibi pravitas haereseos Calvinianae. }»}}\rbrack{} ».}
              \teiP[xmlid={p8-5}]{La ressemblance du panneau avec le diptyque Wilton\teiNote[n={15},place={foot},xmlid={ftn61-3}]{\teiP[xmlid={p-110}]{Visible sur : \teiRef[target={https://commons.wikimedia.org/wiki/File:Wilton\_diptych.jpg}]{https://commons.wikimedia.org/wiki/File:Wilton\_diptych.jpg} (consulté le 03 mars 2025).}}, pour la forme comme pour le contenu religieux, est claire. Il est vrai que dans la vaste littérature sur le chef-d’œuvre de la National Gallery de Londres, on rencontre de nombreuses références à la peinture romaine\teiNote[n={16},place={foot},xmlid={ftn62-3}]{\teiP[xmlid={p-111}]{Dillian Gordon, Lisa Monnas, Caroline Elam (dir.), \teiHi[rend={italic}]{The regal image of Richard II and the Wilton Diptych}, Londres, Harvey Miller, 1997 ; Dillian Gordon, \teiHi[rend={italic}]{The Wilton Diptych}, Londres, The National Gallery, 2015.}}, depuis un article de Charles Coupe, publié en 1895, où pour la première fois l’attention a été attirée sur le texte de Pietrasanta, présentant pour l’occasion l’unique autre description ancienne de la peinture, contenue dans le manuscrit Harley 360 de la British Library\teiNote[n={17},place={foot},xmlid={ftn63-3}]{\teiP[xmlid={p-112}]{C. Coupe, « An Old Picture », \teiHi[rend={italic}]{The Month}, n\teiHi[rend={sup}]{o} 84, 1895, p. 229-242.}}. La figuration et l’inscription du panneau romain ont été en particulier utilisées pour renforcer la lecture du diptyque londonien comme représentation du roi Richard faisant don de son royaume à la Vierge : une idée de l’Angleterre comme « dot de la Madonne » qui était répandue à l’époque, symbolisée dans le diptyque par l’étendard à la croix de saint Georges surmonté d’un petit globe portant une carte de l’Angleterre tenu par un ange et présenté à la Vierge à l’Enfant. On pourrait vraiment faire l’hypothèse que le diptyque et le panneau romain furent deux moments d’un même programme artistique, conçu par le roi Richard peu après la mort d’Anne de Bohême en 1394 et peu avant son remariage\teiNote[n={18},place={foot},xmlid={ftn64-3}]{\teiP[xmlid={p-113}]{J. H. Harvey, « The Wilton Diptych: A Re-examination », \teiHi[rend={italic}]{Archaeologia}, n\teiHi[rend={sup}]{o} 98, 1961, p. 1-28, ici p. 20 : « \lbrack{}…\rbrack{} \teiHi[rend={italic}]{it seems highly probable that the Diptych and the Roman polyptych were two related parts of a single artistic programme in Richard's mind, and that they were of approximately the same date, soon after the death of Anne of Bohemia and before the king's remarriage} \lbrack{}il semble hautement probable que le diptyque et le polyptyque romain formaient deux parties liées d’un même programme artistique dans l’esprit de Richard et que leurs dates étaient proches : peu après la mort d’Anne de Bohême et avant le remariage du roi\rbrack{}. »}}.}
              \teiP[xmlid={p9-4}]{Malgré l’importance du panneau jadis conservé au collège des Anglais jouxtant l’église Saint-Thomas-de-Cantorbéry, on sait peu de chose de cette œuvre : seulement ce que nous disent les descriptions précédemment mentionnées, on ignore par ailleurs tout de la commande de la peinture et de sa probable destruction. Le paradoxe est évident : si cette œuvre est bien connue dans les études sur la peinture gothique anglaise, la présence d’un exemple notoire de peinture médiévale provenant d’Angleterre passe totalement inaperçue dans la littérature sur l’art au Moyen Âge à Rome et sur la réception de l’art médiéval romain à l’époque moderne. Nous devons aux intérêts héraldiques de Pietrasanta et au rôle du panneau comme source de la connaissance des usages héraldiques du Moyen Âge tardif le fait que le jésuite en parle avec tant d’attention et qu’il nous en ait laissé une reproduction partielle.}
            
\end{teiDiv}
            \begin{teiDiv}[type={section1},xmlid={div02-3}]

              \teiHead[xmlid={francesco-gualdi-et-les-armoiries-entre-erudition-et-conservation-001}]{Francesco Gualdi et les armoiries, entre érudition et conservation}
              \teiP[xmlid={p10-4}]{Francesco Gualdi, évoqué au début de cet article en tant que dédicataire de l’opuscule de Gauges de’ Gozze \teiHi[rend={italic}]{Se dalle armi, o insegne, che parlano}…, fut le concepteur des \teiHi[rend={italic}]{Memorie sepolcrali }\lbrack{}Mémoires sépulcraux\rbrack{}, un vaste traité sur les pierres tombales médiévales des églises romaines écrit au début des années 1640 avec l’aide d’une équipe de collaborateurs qui comprenait l’historien Costantino Gigli\teiNote[n={19},place={foot},xmlid={ftn65-3}]{\teiP[xmlid={p-114}]{Gigli fut le plus proche collaborateur de Gualdi. Son intérêt pour l’héraldique est surtout documenté par deux tomes manuscrits qui constituent le codex de la BAV, Vat. Lat. 8253, dans lequel des centaines de transcriptions d’épitaphes des églises de Rome sont accompagnées de descriptions et de rapides esquisses des écus présents sur les pierres tombales et sur les monuments. À Gigli est lié l’album de dessins de sujet sépulcral RCIN 970334 de la Royal Library de Windsor : voir Fabrizio Federici et Jörg Garms, « “Tombs of illustrious Italians at Rome”. L’album di disegni RCIN 970334 della Royal Library di Windsor », \teiHi[rend={italic}]{Bollettino d’Arte}, volume spécial, 2011.}} et le poète Ottavio Tronsarelli\teiNote[n={20},place={foot},xmlid={ftn66-3}]{\teiP[xmlid={p-115}]{Le traité des \teiHi[rend={italic}]{Memorie sepolcrali}, dont la composition remonte aux années 1640-1644, donc aux derniers temps du pontificat d’Urbain VIII, resta inédit jusqu’à la mort de Gualdi en 1657. Nous en avons préparé une édition critique de prochaine publication aux éditions de la BAV. Les citations d’extraits de cette œuvre effectuées dans la présente contribution font référence à la division en chapitres et paragraphes adoptée pour cette édition. Sur Gualdi et le traité voir Claudio Franzoni et Alessandra Tempesta, « Il museo di Francesco Gualdi nella Roma del Seicento tra raccolta privata ed esibizione pubblica », \teiHi[rend={italic}]{Bollettino d’Arte}, vol. LXXVII, n\teiHi[rend={sup}]{o} 73, 1992, p. 1-42 et Fabrizio Federici, « Il trattato \teiHi[rend={italic}]{Delle memorie sepolcrali} del cavalier Francesco Gualdi: un collezionista del Seicento e le testimonianze figurative medievali », \teiHi[rend={italic}]{Prospettiva}, n\teiHi[rend={sup}]{o} 110-111, avril-juin 2003, p. 149-159.}}. Dans les \teiHi[rend={italic}]{Memorie, }un certain intérêt pour l’héraldique se révèle : les armoiries sculptées ou gravées sur les pierres tombales sont soigneusement décrites et souvent reproduites dans les xylographies destinées à constituer l’appareil illustratif de l’œuvre dont elles constituent la majeure partie (fig. 47). Même des armoiries non sépulcrales sont évoquées à plusieurs occasions en tant que témoignages permettant d’attribuer à certaine famille une commande architecturale\teiNote[n={21},place={foot},xmlid={ftn67-3}]{\teiP[xmlid={p-116}]{Chap. II, § 11-12 : « \teiHi[rend={italic}]{La famiglia Alberina fabricò non solo la cappella antica di San Domenico, in questo medesimo sito, nel quale si fabrica questa che vediamo, di San Domenico di Soriano; ma fece ancora tutto ’l muro, che si vede in questa testa della traversa, come si prova dalle due porte antiche, delle campane, e della sagrestia, che mettono in mezzo la cappella; nell’architrave delle quali è scolpita l’arme della suddetta famiglia, con queste parole sotto: }\lbrack{}…\rbrack{} \lbrack{}La famille Alberina fabriqua non seulement l’ancienne chapelle de Saint-Dominique sur ce même site dans lequel se fabrique celle que nous voyons de Saint Dominique de Soriano ; mais elle fit encore tout le mur qui se voit au début du côté comme le montrent les deux anciennes portes, les cloches et la sacristie qui distinguent la chapelle ; sur leur architrave sont sculptées les armoiries de ladite famille sous lesquelles sont ces mots :…\rbrack{} » ; \teiHi[rend={italic}]{ibid}., § 16 : « \teiHi[rend={italic}]{Vedesi ancora nella traversa di questa chiesa, similmente a man sinistra, vicino all’altare della Nunziata, un’altra arme della Casa Alberina, fatta di musaico, la quale si conosce esser frammento di qualche altra lapida antica, oppure, più probabilmente, della fronte d’un altare, quivi fatto dall’istessa famiglia} \lbrack{}On voit encore sur le côté de cette église similairement à main gauche, près de l’autel de l’Annonciation, une autre armoirie de la maison Alberina faite en mosaïque et connue pour être un fragment d’une autre ancienne pierre ou plus probablement du fronton d’un autel fait là par la même famille\rbrack{} ».}} ou comme preuve du patronage d’une chapelle\teiNote[n={22},place={foot},xmlid={ftn68-3}]{\teiP[xmlid={p-117}]{Chap. XXVIII, § 10-11 : « \teiHi[rend={italic}]{Né altri marmi sepolcrali di questa famiglia a nostra notitia pervenuti sono, che li sopradetti, e la cappella loro antica nel detto tempio d’Araceli del Santissimo Sacramento }\lbrack{}…\rbrack{}\teiHi[rend={italic}]{, e dove continuamente sepelite si sono la gente de’ Conti, senza esservi alcuna particolar memoria, né epitaffij dei nomi de’ loro defonti, se non quanto su la volta mirasi un’arma picciola del loro casato antica, come anche occorre alla gente della Casa Cesarina, quale ha per antico uso di sepelirsi nella loro cappella, la qual è medesimamente in Araceli, et in essa non appare di fuori alcuna inscrittione, né memoria, se non l’antica loro arma} \lbrack{}D’autres marbres sépulcraux de cette famille ne sont pas venus à notre connaissance sinon ceux que nous avons mentionnés et leur antique chapelle dans le sanctuaire d’Aracoeli du Saint-Sacrement \lbrack{}...\rbrack{}, où l’on a constamment enseveli les membres de la famille Conti sans mémoire particulière ni épitaphe aux noms de leurs défunts sinon quand on voit à la voûte une petite armoirie de leur antique maison, comme cela arrive aussi pour ceux de la maison Cesarina qui ont pour antique usage de se faire ensevelir dans leur chapelle laquelle est aussi à l’Aracoeli et en celle-ci n’apparaît ni inscription, ni mémoire, sinon leurs antiques armoiries\rbrack{} ».}}.}
              \teiP[xmlid={p11-4}]{Dans le chapitre portant sur la famille Anguillara, un ample passage est consacré à un épisode dont Gualdi lui-même fut le protagoniste, s’employant à sauver deux écus de cette antique maison qui risquaient d’être perdus ou détruits (fig. 48). Deux écus des Anguillara datant du \teiHi[rend={small-caps}]{xv}\teiHi[rend={sup}]{e} siècle, placés sur la façade de l’hôpital de Saint-Jean-de-Latran pour témoigner qu’un membre de cette famille, Everso II (vers 1398-1464), avait agrandi l’édifice, furent détachés lors des travaux de réfection du bâtiment ordonnés par le pape Urbain VIII, datés de 1636 dans l’épigraphe célébrative fixée à l’angle de la façade vers la rue de Saint-Jean-de-Latran\teiNote[n={23},place={foot},xmlid={ftn69-3}]{\teiP[xmlid={p-118}]{Sur les travaux entrepris par Urbain VIII à l’Hôpital du Latran voir Martin Raspe, « La Corsia Nuova dell’Ospedale del SS. Salvatore ad Sancta Sanctorum e i suoi architetti Giacomo e Giovanni Battista Mola », dans Philine Helas et Patrizia Tosini (dir.), \teiHi[rend={italic}]{Tra Campidoglio e curia: l’Ospedale del SS. Salvatore ad Sancta Sanctorum tra Medioevo ed Età Moderna}, Cinisello Balsamo, Silvana, « Studi della Bibliotheca Hertziana, 11 », 2017, p. 137-147.}}. Le chevalier Gualdi adressa au cardinal Barberini et aux gardiens de l’hôpital un discours aujourd’hui perdu, dont un ample résumé figure dans les \teiHi[rend={italic}]{Memorie sepolcrali}. Dans son écrit le gentilhomme plaidait pour la conservation des deux écus, rappelant une série d’exemples de ce type en rapport avec les témoignages du passé :}
              \begin{teiCit}[xmlid={cit02-2}]

                \begin{teiQuote}
Les deux armoiries mentionnées plus haut, ayant été enlevées de leur lieu les années passées lors de la construction du nouveau bâtiment de l’hôpital, me conduisirent, poussé par mon propre génie de la conservation des antiques vestiges, sacrés comme profanes, à représenter au seigneur cardinal Francesco Barberini et aux seigneurs gardiens de cet hôpital, combien il était inconvenant dans la ville de Rome de priver les familles vivantes ou éteintes de la gloire justement méritée, laquelle ne se peut longuement conserver sinon dans le bronze, dans le marbre ou dans les écritures publiques. Je mentionnai l’exemple de Lautrec, général en Italie de Charles VIII roi de France, lequel voulant prendre la ville de Rimini ma patrie, alors défendue par son seigneur Sigismond Malatesta, le convainquit de renoncer à détruire le célèbre pont d’Octave Auguste, construit sur la Marechia (fleuve formant avec la mer Adriatique le port de la ville rendu célèbre par saint Antoine de Padoue qui en ce lieu où l’on voit une petite chapelle attira les poissons pour entendre sa prédication), Malatesta avait déjà décidé de démolir ce pont pour couper la route aux ennemis de ce côté, mais le grand capitaine Lautrec plutôt que de détruire ce remarquable monument choisit d’essayer de pénétrer dans la ville par des passages plus malaisés et exposés à toutes sortes de dangers. Je mentionnai aussi l’exemple de la noble cité de Florence où une ancienne loi ordonne que si l’on bâtit un nouvel édifice dans un site où se trouvaient des armes ou autres antiques vestiges des familles, ces mêmes témoignages soient remis au même lieu où ils se trouvaient déjà, dans toute leur ancienneté et dans toute leur simplicité\teiNote[n={24},place={foot},xmlid={ftn70}]{\teiP[xmlid={p-119}]{Cela fait référence à la « \teiHi[rend={italic}]{Legge contro chi rimovesse o violasse armi, inscrittioni o memorie existenti apparentemente nelli edifitii così publici come privati} \lbrack{}Loi contre ceux qui enlèveraient ou violeraient armoiries, inscriptions ou vestiges existant visiblement dans les édifices publics ou privés\rbrack{} » promulguée par le grand-duc de Toscane le 30 mai 1571, dans laquelle est évoqué « \teiHi[rend={italic}]{l’uso et inveterata consuetudine della città \lbrack{}scil. Firenze\rbrack{} (della quale non è memoria alcuna in contrario)} \lbrack{}l’usage et habitude enracinée de la ville \lbrack{}Florence\rbrack{} (dont il n’est aucun témoignage contraire)\rbrack{} », selon lequel « \teiHi[rend={italic}]{non era lecito a chi comperava, o per qual si voglia modo acquistava alcuno edifitio, rimuovere, estinguere o violare le dette armi et memorie de’ construttori o fondatori d’esso} \lbrack{}il n’était pas licite à qui achetait, où d’une quelconque manière acquérait un édifice, d’enlever, effacer ou violer lesdites armes ou vestiges de leurs fondateurs ou constructeurs\rbrack{} » (la mesure est rapportée dans Andrea Emiliani, \teiHi[rend={italic}]{Leggi, bandi e provvedimenti per la tutela dei beni artistici e culturali negli antichi stati italiani 1571-1860 }\lbrack{}Alfa, 1978\rbrack{}, Bologne, Nuova Alfa, 1996, p. 23-24). Teodoro Amayden est aussi un admirateur enthousiaste de la tradition florentine, voir \teiHi[rend={italic}]{La Storia delle famiglie romane, con note ed aggiunte del Comm. Carlo Augusto Bertini} \lbrack{}Collegio Araldico, 1910\rbrack{}, 2 vol., Rome, Edizioni Romane Colosseum, 1987, chap. I, p. 395 : « \teiHi[rend={italic}]{Siano per mille volte benedetti i Fiorentini, che non cancellano mai né inscrizione, né arme }\lbrack{}Soient mille fois bénis les Florentins qui n’effaçaient jamais ni inscriptions ni armoiries\rbrack{} ».}}. J’alléguai enfin l’exemple de Léon XI, d’éternelle mémoire, dont je fus effectivement serviteur 14 ans durant son cardinalat, et 27 jours qu’il plut à Dieu de le donner ou plutôt montrer au monde comme pontife, celui-ci s’étant vu attribuer le titre de Sainte-Praxède et voulant par conséquent mettre cette église en meilleur état, dans les restaurations qu’il fit effectuer, l’ornant de peintures, il voulut qu’avec ses armes on mit encore celles de saint Charles Borromée qui avait été titulaire de la même église, afin qu’avec sa mémoire se conservât aussi celle de son prédécesseur\teiNote[n={25},place={foot},xmlid={ftn71}]{\teiP[xmlid={p-120}]{Chap. VI, § 15-17 : « \teiHi[rend={italic}]{Queste due armi dette di sopra, essendo state gli anni addietro levate dal lor luogo, coll’occasione della nuova fabrica dello spedale, mossero me (il cavalier Gualdi), spinto a ciò dal proprio genio verso la conservatione delle antiche memorie, sì sacre, come profane, a rappresentare al signor cardinale Francesco Barberini, et a’ signori guardiani dello spedale suddetto, quanto sia sconvenevole nella città di Roma il privar le famiglie, o viventi, o vero estinte, della gloria giustamente meritata, la quale non si può lungamente conservare, se non ne’ bronzi, e ne’ marmi, o nelle publiche scritture. Apportai l’esempio di Lotrecco, generale in Italia di Carlo VIII re di Francia, il quale volendo pigliar la città di Rimini mia patria, difesa allhora da Sigismondo Malatesta, signor di quella, fece intendergli, che lasciasse di rovinar il ponte celebre d’Ottaviano Augusto, fabricato sulla Marechia (forma questo fiume, insieme col Mare Adriatico, il porto della città, reso celebre da Sant’Antonio da Padova, il quale nel luogo, ove si vede una cappelletta, fece comparire ad udir la sua predica i pesci del mare), il qual ponte era già destinato dal Malatesta ad esser demolito, per tagliar da quella parte il passo al nemico; eleggendo quel magnanimo capitano, più tosto che la rovina di quell’eccelsa fabrica, il tentar l’entrata nella città per passi più malagevoli, et esposti a qualsivoglia pericolo. Apportai anche l’esempio della nobile città di Fiorenza, nella quale per antica legge vien ordinato, che rifacendosi nuova fabrica in qualche sito, dove fusser poste o armi, o altra memoria antica di famiglie, quelle stesse memorie, così vecchie, e così semplici, come anticamente si solevan fare, si ripongano nell’istesso luogo, dov’eran già collocate. Addussi finalmente l’esempio di Leone XI, di memoria eterna, di cui fui attual servitore per lo spazio di quattordici anni nel cardinalato, e per 27 giorni, che piacque a Dio di darlo, o pure mostrarlo al mondo, pontefice; al quale essendo toccato il titolo di Santa Prassede, e volendo perciò ridurre in migliore stato quella chiesa, ne’ risarcimenti che vi fece, ornandola di pitture, volle che colle armi sue si ponessero quelle ancora del cardinale San Carlo Borromeo, già benefico titolare della medesima chiesa, accioché colla presente sua memoria si conservasse ancora quella del suo antecessore }».}}.
\end{teiQuote}
              
\end{teiCit}
              \teiP[xmlid={p12-4}]{Grâce à l’intervention de Gualdi, les deux écus ne furent pas perdus et furent replacés exactement au même endroit qu’ils occupaient avant le lancement des travaux\teiNote[n={26},place={foot},xmlid={ftn72}]{\teiP[xmlid={p-121}]{Une fois replacés, les écus furent accompagnés d’une inscription rappelant qu’ils avaient retrouvé leur position d’origine sur la façade de l’hôpital du Latran : « \teiHi[rend={italic}]{Hoc insigne repertum / affixum muro veteri / d(omini) custodes / in muro novo eodem / in loco affigi mandarunt }\lbrack{}Cet insigne ayant été trouvé fixé dans un vieux mur, les seigneurs gardiens ordonnèrent qu’il fût fixé au même endroit dans le mur nouveau\rbrack{} » (Vincenzo Forcella, \teiHi[rend={italic}]{Iscrizioni delle chiese e d’altri edificii di Roma dal secolo XI fino ai nostri giorni}, Rome, Tipografia delle scienze matematiche e fisiche, 1869-1884, 14 vol., t. VIII, p. 133, n\teiHi[rend={sup}]{o} 347). Un des deux écus des Anguillara est encore visible sur la façade de l’hôpital tournée vers la basilique du Latran, accompagné de l’inscription du \teiHi[rend={small-caps}]{xvii}\teiHi[rend={sup}]{e} siècle ; l’autre, jadis inséré dans la façade arrière, vers San Clemente, se trouve maintenant dans une salle de l’hôpital, encastré dans une cloison (les deux côtés de la plaque sont visibles, et l’on peut ainsi constater que, pour réaliser l’écu, on réutilisa un parapet du haut Moyen Âge, décoré de croix et de motifs géometriques).}}. La “bataille” de Gualdi pour la conservation de ces deux armoiries est à replacer dans le cadre d’une infatigable lutte pour la défense des monuments antiques et médiévaux en péril conduite par le gentilhomme\teiNote[n={27},place={foot},xmlid={ftn73}]{\teiP[xmlid={p-122}]{Voir Fabrizio Federici, « Battaglie per la tutela nella Roma barocca: Francesco Gualdi e la difesa delle “memorie antiche” », \teiHi[rend={italic}]{Studi Romani}, vol. LXII, n\teiHi[rend={sup}]{o} 1-4, 2014, p. 149-172 (en particulier, pour l’affaire des armoiries des Anguillara, p. 153-155).}}, qui eut même des retombées légales : ce fut probablement sous la pression du chevalier Gualdi que le cardinal Antonio Barberini sénior promulgua le 2 octobre 1640 un édit où il était intimé « à tous les supérieurs des églises séculiers ou réguliers \lbrack{}\teiHi[rend={italic}]{a tutti i superiori delle chiese, tanto secolari, quanto regolari}\rbrack{} » de ne pas faire enlever « témoignages, inscriptions et pierres \lbrack{}\teiHi[rend={italic}]{memorie, inscrittioni, e lapidi}\rbrack{} » – expression qui faisait évidemment aussi référence aux armoiries – « sans le consentement des parties et notre autorisation écrite \lbrack{}\teiHi[rend={italic}]{senza il consenso delle parti, e nostra licenza in scritto}\rbrack{} », dans lequel il était prescrit de sévères punitions pour les contrevenants\teiNote[n={28},place={foot},xmlid={ftn74}]{\teiP[xmlid={p-123}]{\teiHi[rend={italic}]{Ibid.}, p. 158-159.}}.}
              \teiP[xmlid={p13-4}]{À l’empressement montré par Gualdi pour sauvegarder les anciens vestiges, on peut ajouter d’autres exemples d’une sensibilité similaire pour la conservation, précédant ou suivant l’action du gentilhomme, issus de la volonté des familles de préserver les traces de leur antique noblesse, ou déterminés par la volonté de transmettre à la postérité les témoignages historiques sans le motif des liens du sang. On peut donner en exemple les vestiges de la famille Casali rassemblés dans l’atrium du palais ancestral de Campo Marzio, au 23 de la rue de la Stelletta : parmi ceux-ci ressortent la pierre tombale de Francesca Casali († 1463), autrefois à Sant’Agostino\teiNote[n={29},place={foot},xmlid={ftn75}]{\teiP[xmlid={p-124}]{Sur la partie haute de la pierre tombale est sculpté l’écu des Casali, dans un losange quadrilobé ; sous celui-ci est gravée l’inscription sépulcrale, en capitales. Dans l’angle inférieur droit est gravé un chandelier. Pour le texte de l’épitaphe, voir Forcella, \teiHi[rend={italic}]{Iscrizioni delle chiese},\teiHi[rend={italic}]{ op. cit.}, t. V, p. 12, n\teiHi[rend={sup}]{o} 23, qui considérait l’objet comme perdu (de même dans Jan L. de Jong, \teiHi[rend={italic}]{Tombs in Early Modern Rome (1400-1600). Monuments of Mourning, Memory and Meditation}, Leiden, Brill, 2023, p. 228). La pierre tombale est localisée à Sant’Agostino dans les \teiHi[rend={italic}]{Memorie sepolcrali} de Gualdi (chap. XXIII, § 3) et aussi par Pier Luigi Galletti, \teiHi[rend={italic}]{Inscriptiones Romanae infimi aevi Romae exstantes}, Romae, Typis Generosi Salomoni Bibliopolae, 1760, p. LXIII, n. 7. L’extraction de la pierre tombale du pavement de l’église et son transfert sous le porche du palais durent se produire dans les années qui succédèrent immédiatement à la publication du volume de Galletti, quand la chapelle Casali fut démolie (voir à ce propos l’inscription de 1764 rapportée dans Forcella, \teiHi[rend={italic}]{Iscrizioni delle chiese, op. cit.}, t. V, p. 106, n\teiHi[rend={sup}]{o} 316 et Lothar Sickel, « L’antica cappella Casali in Sant’Agostino a Roma in un disegno di Taddeo Kuntz: una testimonianza per un affresco scomparso del Pintoricchio? », \teiHi[rend={italic}]{Bollettino d’Arte}, n\teiHi[rend={sup}]{o} 43, 2020, p. 59-78).}}, et un écu des Casali accompagné d’une inscription vulgaire en majuscules gothiques datée de 1429, où est évoquée la construction d’un puits\teiNote[n={30},place={foot},xmlid={ftn76}]{\teiP[xmlid={p-125}]{L’inscription inédite dit : « \teiHi[rend={italic}]{Q(ue)sto poçço a(n)no fac(t)o / li figli d(e) Ia(n)ni d(e) Romano / Casal(e) allo(r) spese et / no alle spe(se) de la contrada c(onstruct)o }\lbrack{}?\rbrack{} \teiHi[rend={italic}]{/ p(er) lo(r) case a(nn)o D(omini) M}\teiHi[rend={italic sup}]{o}\teiHi[rend={italic}]{IIII}\teiHi[rend={italic sup}]{C}\teiHi[rend={italic}]{XXIX }\lbrack{}Ce puits, les fils de Jean de Romain Casali l’ont construit à leurs dépens et non à ceux du district pour leurs maisons, l’an du Seigneur 1429\rbrack{} ». Au milieu du \teiHi[rend={small-caps}]{xvii}\teiHi[rend={sup}]{e} siècle, Teodoro Amayden situait l’objet dans le jardin des Casali à Marmorata : « Ils ont dans ledit jardin une pierre tombale de celles qu’on appelle antiques-modernes, avec les armes en haut et au-dessous l’inscription en lettres mi-gothiques \lbrack{}\teiHi[rend={italic}]{Hanno nel sopradetto giardino una lapide di quelle che si chiamano antiche moderne, col arme in cima, e la scrittura sotto in lettera mezza gotica}\rbrack{} ». Suit la trancription de l’épigraphe (Biblioteca Casanatense, ms. 1335, fol. 254r ; il s’agit d’un passage du chapitre de la \teiHi[rend={italic}]{Storia delle famiglie romane} consacré aux Casali non transcrit dans la discutable édition de Carlo Augusto Bertini, \teiHi[rend={italic}]{La Storia delle famiglie romane...}, Rome, 1910, 2 vol.). Sur l’adjectif « \teiHi[rend={italic}]{anticomoderno }» employé par Amayden, utilisé à l’époque moderne pour désigner les derniers siècles du Moyen Âge, sinon toute l’époque médiévale, voir Fabrizio Federici, «\teiHi[rend={italic}]{ }“Anticomoderno”: significati ed usi del termine nella letteratura artistica tra Cinque e Settecento\teiHi[rend={italic}]{ }», dans Maria Monica Donato et Massimo Ferretti (dir.), \teiHi[rend={italic}]{« Conosco un ottimo storico dell’arte… ». Per Enrico Castelnuovo. Scritti di allievi e amici pisani}, Pise, Edizioni della Normale, 2012, p. 291-296.}} (fig. 49). Ce dernier témoignage offre un intérêt particulier : Giovanni Battista Casali, un membre de la famille qui dans la première moitié du \teiHi[rend={small-caps}]{xvii}\teiHi[rend={sup}]{e} siècle collectionnait les antiques et publia plusieurs ouvrages sur l’ancienne Rome, les rituels religieux des Égyptiens, des Romains ou des chrétiens, fit tirer de cet écu une petite xylographie reproduite à la fin de ses doctes volumes en guise de signature\teiNote[n={31},place={foot},xmlid={ftn77}]{\teiP[xmlid={p-126}]{Casali inséra cette gravure à la fin de ses travaux d’antiquaire sur les rites et ensuite en conclusion du volume qui réunit la série, voir Giovanni Battista Casali, \teiHi[rend={italic}]{De profanis, et sacris veteribus ritibus}, Romae, Ex typographia Andreae Phaei, 1645.}} (fig. 50). La gravure, accompagnée d’une brève didascalie (« \teiHi[rend={italic}]{Ex antiquo marmore apud auctorem} \lbrack{}d’après un marbre antique, chez l’auteur\rbrack{} »), représente l’écu avec exactitude mais elle omet complètement l’inscription sous-jacente, à l’exception de la date conclusive. En choisissant de sceller les volumes par une antique représentation de son propre écu, Casali entendait souligner son statut d’érudit issu d’une noble lignée, membre d’une illustre maison romaine. À cette xylographie on peut en joindre une autre, figurant parmi celles que Gualdi fit réaliser pour les \teiHi[rend={italic}]{Memorie sepolcrali} : il s’agit d’une pièce elle aussi de format réduit (65 × 50 mm) qui reproduit l’écu des Gualdi inséré sur le côté nord du Palais Sénatorial, au Capitole (fig. 51 et 52). L’écu de marbre fut placé là dans la première moitié du \teiHi[rend={small-caps}]{xvi}\teiHi[rend={sup}]{e} siècle par un des deux ancêtres du chevalier Gualdi qui assumèrent la charge prestigieuse de sénateur de Rome, Galeotto ou Francesco\teiNote[n={32},place={foot},xmlid={ftn78}]{\teiP[xmlid={p-127}]{L’écu reproduit est celui autour duquel Gualdi disposa trois inscriptions en mémoire des sénateurs Galeotto et Francesco, peut-être la même année, 1655, où il érigea non loin de là le monument à Scipion l’Africain (on ne peut cependant exclure que les inscriptions aient été mises à une date antérieure). Sur le monument à Scipion voir Franzoni, Tempesta, \teiHi[rend={italic}]{Il museo di Francesco Gualdi},\teiHi[rend={italic}]{ op. cit.} et Fabrizio Federici, « Baroque Returns: the Donations and Reuses of Francesco Gualdi », dans Eva-Maria Troelenberg, Felicity Bodenstein et Damiana Otoiu (dir.), \teiHi[rend={italic}]{Contested Holdings: Museum Collections in Political, Epistemic and Artistic Processes of Return}, actes du colloque \teiHi[rend={italic}]{What do Contentious Objects Want? Political, Epistemic and Artistic Cultures of Return}, Florence, Kunsthistorisches Institut in Florenz, 21-22 ottobre 2016, New York-Oxford, Berghahn, 2022, p. 197-219. La xylographie commandée par Gualdi est publiée dans Ferdinando Ughelli, \teiHi[rend={italic}]{Italia Sacra}, Rome, Apud Bernardinum Tanum, 1647, t. II, col. 432, à la fin d’un éloge de cette famille de Rimini.}}. Nous ne savons pas à quel endroit des \teiHi[rend={italic}]{Memorie }cette xylographie était destinée : on peut penser que comme l’autre réalisée dans les mêmes années elle était prévue pour clore le volume, pour souligner le prestige de la famille du collectionneur de Rimini et le rang sénatorial de bien deux de ses ancêtres.}
              \teiP[xmlid={p14-3}]{Si la décision de préserver les vestiges et les armoiries des Casali dans l’atrium de leur palais fut déterminée par la volonté de sauvegarder la mémoire familiale, le déplacement de la pierre tombale de Giuliano Porcari († 1282) à San Giovanni della Pigna (fig. 53) fut en revanche motivé par une sensibilité pour la conservation indépendante de l’orgueil nobiliaire. En 1562, le recteur de l’église, Nicola Martinelli, fit extraire du pavement du sanctuaire la belle pierre tombale qui montre une croix et un écu nobiliaire en mosaïque dans un cadre de même, et la fit insérer dans un mur, expliquant les raisons de son initiative dans un inscription malheureusement disparue mais recueillie par Forcella : « La pierre que tu vois se trouvait devant l’entrée du sanctuaire ; afin que le signe de la croix ne fût pas foulé aux pieds de ceux qui entrent ni que cette mémoire de l’antiquité de la famille Porcari ne tombât dans l’oubli, Nicola Martinelli, recteur de cette église, l’a fait poser dans ce lieu plus honorable et plus commode, l’an du salut 1562\teiNote[n={33},place={foot},xmlid={ftn79}]{\teiP[xmlid={p-128}]{Vincenzo Forcella, \teiHi[rend={italic}]{op. cit.}, t. IX, p. 487, n\teiHi[rend={sup}]{o} 976. Sur la pierre tombale du \teiHi[rend={small-caps}]{xiii}\teiHi[rend={sup}]{e} siècle, voir Jörg Garms, Roswitha Juffinger et Bryan Ward-Perkins (dir.),\teiHi[rend={italic}]{ op. cit.}, t. I, p. 89-90 : « \teiHi[rend={italic}]{Lapis quem cernis ante ostium templi / iacebat et ne ab introentibus signu(m) / crucis pedib(us) conculcaretur Porciaeq(ue) / familiae vetustatis memoria haec / oblivioni traderetur Nicolaus / Martinellus huius ecclesiae / rector in hunc et honestiorem / et commodiorem locu(m) posuit / anno Salutis MDLXII }».}} ». Le déplacement fut donc motivé non seulement par une raison religieuse (empêcher que la croix fût piétinée par ceux qui entraient dans l’église), mais aussi par la volonté d’éviter que l’antique mémoire des Porcari « tombât dans l’oubli ».}
              \teiP[xmlid={p15-3}]{Pour en revenir aux \teiHi[rend={italic}]{Memorie sepolcrali} et à l’attention que les armoiries nobiliaires reçoivent dans cette œuvre, un aspect qui reste à mettre en évidence est la faveur avec laquelle est généralement considérée la simplicité des représentations héraldiques sur les pierres tombales du passé. Dans le contexte plus large de l’éloge de la pureté et de la modestie des derniers siècles du Moyen Âge, concentré principalement sur les habits, en particulier féminins, même les armoiries sont exploitées comme preuve d’une intégrité morale supérieure désormais perdue, supplantée par la vaine « pompe » des temps présents. En prenant l’exemple de la pierre tombale de Pietro Muti (fig. 54), il est affirmé :}
              \begin{teiCit}[xmlid={cit03}]

                \begin{teiQuote}
La simplicité qui au cours des siècles passés s’observait dans différentes choses a été tellement corrompue qu’enfin dans les armoiries ou insignes elle se trouve tout à fait altérée, d’où on doit bien observer dans les armoiries de nos pierres sépulcrales comment pour être de ces temps elles se trouvent toutes mises et formées sans aucune sorte de cimier ou couronne ou autre ornement extérieur quoiqu’elles soient mémoires de familles romaines ayant dominé des villes et ayant gouverné des États. Cela étant, ils représentaient purement leurs armoiries dans le seul écu ou targe et l’on voit que Bartolomeo Cassaneo, très fameux jurisconsulte qui vivait vers l’an 1529, dans son catalogue \teiHi[rend={italic}]{Gloriae mundi} quand il traite des insignes et des armes nous les représente sans aucun ornement quoiqu’elles appartinssent à de grands seigneurs comme on peut l’y voir\teiNote[n={34},place={foot},xmlid={ftn80}]{\teiP[xmlid={p-129}]{Barthélémy de Chasseneuz, \teiHi[rend={italic}]{Catalogus gloriae mundi}..., Lugduni, Per honoratum virum Dionysium de Harsy, 1529, partie I, p. 26-34.}}, signe évident que dans son temps cela n’était pas encore en usage et c’est un peu plus d’un siècle après que s’introduisit l’usage des cimiers et des couronnes. De là on voit que la maison Malatesta parmi beaucoup d’autres, originaire de ma patrie Rimini, laquelle a dominé à plusieurs reprises jusqu’à vingt-deux villes ne s’autorisa pas à déployer ses armes sinon simples, comme le firent beaucoup d’autres très nobles familles de Rome et d’Italie, mais il peut bien être arrivé par concession de potentats, ou par autre raison, qu’on ait commencé à introduire les heaumes, les couronnes et autres ornements\teiNote[n={35},place={foot},xmlid={ftn81}]{\teiP[xmlid={p-130}]{Chap. XLIX, § 14-16 : « \teiHi[rend={italic}]{La semplicità, che negl’andati secoli s’osservava in diverse cose, talmente è stata corrotta, che infin nelle arme, od insegne vedesi affatto alterata, onde ben devesi osservare nelle arme di queste nostre lapidi sepolcrali come, per esser di quei tempi, si trovino tutte collocate, e formate senza sorte alcuna di cimiero, overo corone, od altri ornamenti all’intorno, ancorché sieno memorie di famiglie romane ch’hebbero dominato città, et havuto giurisdittione di Stati. Nondimeno puramente rappresentavano le loro imprese col solo scudo, o targa; e vedesi che Bartolomeo Cassaneo famosissimo giureconsulto, il qual viveva circa l’anno 1529, nel suo catalogo Gloriae mundi, mentre tratta delle insegne, e dell’arme, egli ce le rappresenta spogliate d’ogni ornamento ancorché sieno di signori grandi, come ivi si può vedere, segno evidente che a’ suoi tempi ancor ciò non era in uso; et è poco più del corso d’un secolo in qua, che si è introdotto l’uso del cimiero, e delle corone. Laonde vedesi la Casa Malatesta tra molte altre, originaria di Rimini mia patria, la quale ha dominato in più volte nell’Italia da ventidue città, e nondimeno non si arrogò mai di spiegare le sue arme se non semplici, sì come molte altre nobilissime famiglie di questa città di Roma, e dell’Italia; ma può ben essere accaduto, che poi forse per concessione de’ potentati, o per altra cagione si sieno cominciati ad introdurre gli elmi, e le corone, e gl’altri ornamenti }». Suivent aux § 17-19 quelques considérations sur les origines des armoiries : « \teiHi[rend={italic}]{Diedero alcuni l’origine delle arme agl’Egittij ne’ loro simboli, altri a’ Thebani nelle lor guerre, chi ad Artù et a’ cavalieri della tavola rotonda, e chi a’ paladini di Carlo Magno, liberatore del mondo. Non ha dubio però che tra gli antichi il valore de’ soldati che molte cose in guerra generosamente acquistarono, o perché erano dagl’imperadori di doni militari, e di varie corone adorni, diede campo di poter da quelle ritrarne un principio delle nostre arme, le quali dopo il secolo di Barbarossa pigliarono qualche forma, e sempre si sono andate avanzando, et ad imitatione di quei doni de’ capitani antichi hanno con bei simboli divisate le loro insegne; ed i regi dagli antichi imperadori presero le corone radiate, et i principi dagl’Augusti inferiori trassero l’uso delle diademe, o fascie gioiellate per coronare le arme loro, ed il nobile da’ soldati di gran fattione l’uso di alzare il cimiero sopra le loro imprese, e finalmente ciascuno nella vanità degli scudi ripone la gloria della pace} \lbrack{}Certains donnèrent l’origine des armes aux Égyptiens dans leurs symboles, d’autres aux Thébains dans leurs guerres, qui à Artus et aux chevaliers de la Table ronde, et qui aux paladins de Charlemagne, libérateur du monde. Il n’y a pourtant aucun doute que, dans l’Antiquité, les soldats ayant généreusement acquis de nombreux gains pendant la guerre ou s’étant vu décerner par les empereurs des dons militaires et orner de couronnes variées, cela donna lieu de pouvoir en tirer l’origine de nos armoiries, lesquelles prirent quelque forme après le siècle de Barbarossa et allèrent toujours en s’avançant, et à l’imitation des dons des capitaines antiques on orna les insignes de beaux symboles ; et les rois prirent les couronnes radiées des antiques empereurs, et les princes prirent l’usage des diadèmes ou des bandeaux ornés de gemmes des Augustes inférieurs pour couronner leurs armoiries, et le noble des soldats méritants l’usage d’élever un cimier au-dessus de leurs armes, et finalement chacun place la gloire de la paix dans la vanité des écus\rbrack{} ». Sur la question des origines des armoiries (que la critique actuelle tend à situer dans la seconde moitié du \teiHi[rend={small-caps}]{xii}\teiHi[rend={sup}]{e} siècle), voir Michel Pastoureau, \teiHi[rend={italic}]{Typologie des sources du Moyen Âge occidental}, Turnhout, Brepols, 1976, fasc. 20 : \teiHi[rend={italic}]{Les armoiries}, p. 24-27.}}.
\end{teiQuote}
              
\end{teiCit}
              \teiP[xmlid={p16-3}]{On trouve des considérations de la même teneur dans un des discours que Gauges de’ Gozze prépara pour son traité et qui ne trouvèrent pas leur place dans la version finale de l’œuvre :}
              \begin{teiCit}[xmlid={cit04}]

                \begin{teiQuote}
Peu d’armoiries ou insignes de familles se trouvent sculptés sur les pierres tombales avec le heaume ou le cimier au-dessus de l’écu dans les siècles passés où l’on aimait une pure et franche simplicité éloignée de toute fumée de grandeur et d’ambition. Après ces bons temps vraiment heureux, étant, outre les autres abus, crû parmi les hommes un désir immodéré de gloire, ils valorisèrent progressivement les titres lesquels étant initialement de sens honorable et appropriés à la qualité de chacun atteignirent ensuite un tel degré de grandeur que les plus élevés se trouvant usurpés par des personnes médiocres sinon viles furent la cause que les cardinaux de notre temps laissèrent l’« illustrissime » devenu trop commun et s’arrogèrent pour eux seuls celui d’« éminentissime ». Ainsi, à la place des simples armes ou écu, les princes dans le simple entourage de la couronne ont ajouté de longs lambrequins et autres ornements, et les nobles (suivis par les plébéiens, quoique leurs singes très vils, ce qui mérite qu’on y remédie) ont placé sur l’écu le heaume et divers emblèmes ou devises\teiNote[n={36},place={foot},xmlid={ftn82}]{\teiP[xmlid={p-131}]{Le passage précède immédiatement celui sur l’origine antique de la coutume de poser des cimiers sur les écus que nous avons déjà cité : « \teiHi[rend={italic}]{Poche armi, od insegne di fameglie si trovano scolpite in lapidi sepolcrali coll’elmo, o cimiero sopra lo scudo negli andati secoli, i quali amavano una pura, e schietta simplicità lontana da ogni fumo di grandezza, ed ambitione: doppo que’ buon tempi veramente felici oltre agli altri abusi, essendo negli huomini cresciuto uno smoderato disiderio di gloria, cominciarono a campeggiar a poco a poco i titoli, i quali da principio essendo di honorevol senso, ed appropriati alla qualità di ciascuno crebbero poi a tal grado di grandezza, che venendo da persone mediocri, nonché vili, usurpati i titoli maggiori, cagionarono che i cardinali a’ tempi nostri lasciassero l’illustrissimo reso troppo commune, e pigliassero quel d’eminentissimo e l’appropriassero a sé soli: così in luogo delle semplici armi, o scudi, i principi al semplice giro della corona hanno aggionte longhissime punte, ed altri ornamenti, ed i nobili (seguiti da plebei etiandio vilissimi loro simie, cosa degna di rimedio) hanno posto sopra lo scudo l’elmo con diversi emblemi, o imprese }».}}.
\end{teiQuote}
              
\end{teiCit}
            
\end{teiDiv}
            \begin{teiDiv}[type={section1},xmlid={div03-3}]

              \teiHead[xmlid={commerce-et-noblesse-001}]{Commerce et noblesse}
              \teiP[xmlid={p17-3}]{Outre les armoiries, un autre élément d’identification des défunts reçoit une grande attention dans les \teiHi[rend={italic}]{Memorie sepolcrali} : les marques de marchands que l’on pouvait et peut encore voir aujourd’hui gravées ou sculptées sur les pierres tombales du Moyen Âge tardif et de la Renaissance\teiNote[n={37},place={foot},xmlid={ftn83}]{\teiP[xmlid={p-132}]{Sur les marques de marchands, voir Umberto Gnoli, « Nobiltà romana e marchi mercantili nel Medioevo\teiHi[rend={italic}]{ }», \teiHi[rend={italic}]{L’Urbe}, vol. VIII, n\teiHi[rend={sup}]{o} 11-12, 1943, p. 15-32 ; Jérôme Hayez, « Un segno fra altri segni. Forme, significati e usi della marca mercantile verso il 1400 », dans Elena Cecchi Aste (dir.), \teiHi[rend={italic}]{Di mio nome e segno. « Marche » di mercanti nel carteggio Datini (secc. XIV-XV)}, Prato, Istituto di studi postali Aldo Cecchi, « Quaderni di storia postale, 30 », 2010, p. IX-XLVI.}}. Deux marques étaient visibles sur la pierre tombale de Lello Maddaleni dans l’église de la Minerve (fig. 55 et 56), qui suscitent un intéressant débat sur ce genre de signes :}
              \begin{teiCit}[xmlid={cit05}]

                \begin{teiQuote}
Nous voyons que depuis environ cinq à six siècles les familles de particuliers ont inventé leurs propres armes ou insignes, et avant on utilisait à leur place les lettres ou caractères voulons-nous dire, ce qui se voit clairement chez Ciaccone, dans les vies des pontifes et des cardinaux de ces temps, ainsi que chez d’autres auteurs. Donc, à propos de ces mêmes lettres ou caractères, les hommes qui faisaient commerce dans les pays lointains, pour marquer leurs marchandises afin que leur propre patron fût reconnu, prirent l’habitude d’y imprimer une lettre ou autre chiffre similaire qui déclarerait le nom propre du marchand et patron de ces marchandises par sa première lettre ou plusieurs ensemble, laquelle se nomme marque ou merque ; ils se servaient d’elle non seulement pour sceau et pour marquer ou merquer lesdites marchandises mais ils l’imprimaient aussi sur toutes leurs affaires, sur leur bétail ou sur toute autre chose, et non contents de cela ils la faisaient enfin graver sur les portes comme sur les meubles de leurs maisons. Plus encore, ils les sculptaient jusque sur leurs sépultures et voulaient les emporter dans la mort comme cela se voit pour notre Madaleni et pour d’autres dans ce livre. Et aujourd’hui tout cela est de coutume pour les marchands puissants et riches\teiNote[n={38},place={foot},xmlid={ftn84}]{\teiP[xmlid={p-133}]{Chap. XVIII, § 3-5 : « \teiHi[rend={italic}]{Noi vediamo che d’intorno a cinque, o sei secoli passati le famiglie de’ particolari hanno inventate le proprie arme, od insegne; e per prima da alcuni usavansi in vece d’esse le lettere, o caratteri che dir vogliamo, il che nel Ciaccone nelle vite di que’ pontefici, e cardinali di que’ tempi, et in altri auttori chiaramente vedesi. Quindi è, che dalle medesime lettere, o caratteri gli huomini che mercatantie in lontane parti facevano, per segnare le loro merci, acciò il proprio padrone riconosciuto fosse, usarono d’imprimere sopra esse una lettera, o altra simile zifra con la quale la prima, o più lettere assieme dichiarassero il nome del proprio mercadante e padrone di quelle merci, la quale marchio, o merchio nominossi; ove con essa non solo per suggello, et per marcare, o merchiare le dette mercadantie servivansi, ma anco vi segnavano tutte le robbe, o bestiami, o altro che fosse, né di ciò contenti infin sopra le porte e suppellettili delle loro case intagliar facevano. Anzi infino nelle sepolture le scolpivano, e seco in morte portar le volevano, come il nostro Madaleni et altri in questo libro vedesi. Et hoggidì il tutto da’ facultosi, e nobili mercadanti usarsi suole} ».}}.
\end{teiQuote}
              
\end{teiCit}
              \teiP[xmlid={p18-3}]{La présence conjointe de marques de marchands et d’armoiries nobiliaires obligeait à réfléchir sur un thème, celui du rapport entre commerce et noblesse, qui était particulièrement ressenti à Rome où la noblesse autochtone était souvent accusée de ne pouvoir se vanter d’origines très anciennes et de s’être élevée d’humbles conditions seulement dans des temps récents : « la majeure partie des familles aujourd’hui estimées à Rome sont issues de basses origines, comme notaire, apothicaire, ce qui serait supportable, mais aussi de l’art malodorant du tannage des cuirs ; c’est une chose que remarque très souvent Metellino \lbrack{}Castalio Metallino\rbrack{} : V. et N. gentilhomme dans l’art du tannage, chose en vérité ridicule, composant sa noblesse avec tant de souillure et cette origine qu’ils eurent est très fraîche », lit-on dans une relation sur les familles romaines du milieu du \teiHi[rend={small-caps}]{xvii}\teiHi[rend={sup}]{e} siècle\teiNote[n={39},place={foot},xmlid={ftn85}]{\teiP[xmlid={p-134}]{Archivio Apostolico Vaticano, Misc. Arm. II 150, fol. 646r-716r, ici fol. 675v-676r : « \teiHi[rend={italic}]{la maggior parte delle famiglie hoggi stimate a Roma nobili vengono da bassi principij, come di notaro, spetiale, che pur sarebbe da soportare, ma dall’arte puzzolente della concia de’ corami; è cosa che nota più, e più volte il Metellino }\lbrack{}Castalio Metallino\rbrack{}\teiHi[rend={italic}]{: V. et N. nobilhuomo dell’arte della concia, cosa invero ridicola, componendo la nobiltà con tanta lordura; e questo principio ch’hebbero è molto fresco} ». Un peu plus loin, il est ajouté : « C’est une chose certaine qu’aucune famille à Rome ne descend des anciens Romains car Constantin transporta le Sénat à Byzance et avec lui tout ce qu’il y avait de bon à Rome, chacun suivant le prince \lbrack{}\teiHi[rend={italic}]{Certissima cosa è, che non è famiglia in Roma, che discenda dagl’antichi Romani, poiché Costantino condusse il senato in Bisanzio, e seco ciò che di buono era in Roma, seguendo ogn’uno il prencipe}\rbrack{} ».}}. De semblables mots sont prononcés dans ces mêmes années par l’historien français Jean Le Laboureur : « Par cette table l’on peut iuger de la noblesse romaine d’auiourd’huy ; dont la pluspart est de naissance estrangere, et si l’on la remenoit en sa source, l’on ne trouveroit que de la bourbe en plusieurs, qui se sont enrichis par les charges, par la banque, et par le trafic, ou par les fortunes de la cour ; particulierement les romaines, ou l’on pourroit à peine en choisir une douzaine de maisons vrayement nobles\teiNote[n={40},place={foot},xmlid={ftn86}]{\teiP[xmlid={p-135}]{Jean Le Laboureur, \teiHi[rend={italic}]{Relation du voyage de la Royne de Pologne et du retour de Mme la mareschalle de Guébriant, ambassadrice extraordinaire...}, Paris, Chez la veuve Jean Camusat, et Pierre Le Petit, 1647, p. 239.}} ».}
              \teiP[xmlid={p19-3}]{À Rome, face à ces accusations, plusieurs auteurs revendiquèrent le fait que la pratique du commerce et plus généralement des activités professionnelles n’est pas contraire à la noblesse. C’est ce qui se lit dans divers passages des \teiHi[rend={italic}]{Memorie sepolcrali}\teiNote[n={41},place={foot},xmlid={ftn87}]{\teiP[xmlid={p-136}]{Chap. X, § 9 (« \teiHi[rend={italic}]{In questo sepolcro del Bellapani, che è di maniera gotica, vediamo scolpito negl’estremi delle colonnette il marchio di mercadante, ma ciò punto non noce all’esser ben nato, e tenuto per nobile }\lbrack{}…\rbrack{} ») ; chap. XVIII, § 6-13 ; chap. XXXIX, § 4-12.}}, comme dans certains extraits d’un texte à peu près contemporain, la\teiHi[rend={italic}]{ Storia delle famiglie romane} de Teodoro Amayden : « Cette famille – écrit le polygraphe hollandais à propos de la famille Cavalieri – a sans aucun doute parfois exercé le négoce parce que l’on voit sur une ancienne porte de marbre d’une de leurs maisons à l’Argentina les armes de la famille flanquées de marques de marchands, ce qui ne déroge aucunement à la noblesse de la famille, parce que lesdites armes portent la bordure, signe de noblesse passée et présente, cela ne doit pas paraître étrange car dans l’église de Santa Maria dell’Anima, dans la chapelle faite par Focari \lbrack{}Fugger\rbrack{} après qu’ils furent devenus princes on voit les armes de la famille d’un côté et de l’autre la marque du négoce\teiNote[n={42},place={foot},xmlid={ftn88}]{\teiP[xmlid={p-137}]{Teodoro Amayden, \teiHi[rend={italic}]{op. cit.}, t. I, p. 286-287 : « \teiHi[rend={italic}]{Questa famiglia – scrive il poligrafo olandese a proposito della famiglia Cavalieri – senza fallo ha esercitato alcune volte il negozio, perché si vede sopra una porta antica di marmo d’una sua casa sotto l’Argentina l’arme della famiglia immezzo la porta e le note mercantili da’ fianchi, il che non deroga punto alla nobiltà della famiglia, perché detta arme porta la risega, segno di nobiltà precedente e concomitante, il che non deve parer strano, poiché nella chiesa di Santa Maria dell’Anima nella cappella fatta da’ Focari }\lbrack{}Fugger\rbrack{}\teiHi[rend={italic}]{ dopo che sono prencipi, si vede l’arme della famiglia e dall’altro lato il merco del negozio }\lbrack{}…\rbrack{} ».}} \lbrack{}...\rbrack{} ». La référence est selon toute probabilité à un écu des Cavalieri inséré dans le mur à côté de l’entrée de la tour du Papito, qui s’élève sur l’aire archéologique du Largo de Torre Argentina\teiNote[n={43},place={foot},xmlid={ftn89}]{\teiP[xmlid={p-138}]{L’écu, de facture assez ordinaire, n’est pas flanqué des marques de marchands évoquées par Amayden, probablement gravées sur des portions de travertin qui ne nous sont pas parvenues. Comme d’habitude, l’écu Cavalieri porte la « \teiHi[rend={italic}]{risega} », ou bordure, ici indentée.}} (fig. 57). Amayden encore souligne de manière polémique dans un autre passage : « J’observe bien qu’au cours de pas tellement plus de deux cents ans la noblesse romaine a beaucoup changé, les nobles ayant exercé communément la marchandise ce pourquoi ils étaient riches, alors qu’au jour d’aujourd’hui ils en ont honte et ils sont pauvres\teiNote[n={44},place={foot},xmlid={ftn90}]{\teiP[xmlid={p-139}]{Teodoro Amayden, \teiHi[rend={italic}]{op. cit.}, t. II, p. 47 : « \teiHi[rend={italic}]{Ben osservo che, in corso di duecento e poco più anni, si è mutata molto la nobiltà romana, esercitando comunemente la mercanzia li nobili, e perciò erano ricchi, ove alli dì d’oggi l’hanno per vergogna e sono poveri }».}}. » En accord avec les positions exprimées par Gualdi et Amayden est ce qu’affirme Giovanni Pietro Crescenzi, soutenant que le petit commerce n’est pas noble, alors que la « marchandise opulente \lbrack{}\teiHi[rend={italic}]{mercatura opulenta}\rbrack{} », les grands traffics ne sont pas en fait incompatibles avec la noblesse\teiNote[n={45},place={foot},xmlid={ftn91}]{\teiP[xmlid={p-140}]{Giovanni Pietro Crescenzi, \teiHi[rend={italic}]{Il nobile romano, o sia trattato di nobiltà}, Bologna, Per gli eredi d’Antonio Pisarri, 1693, p. 193-207.}}.}
              \teiP[xmlid={p20-3}]{Aux yeux des nobles qui ne pouvaient se vanter d’une illustre lignée, toutefois, ces arguments ne devaient pas sembler suffisants, et les racines mercantiles de leurs familles continuaient à constituer une source d’embarras. À tel point que certains tentèrent d’occulter leurs propres origines, intervenant matériellement sur les vestiges familiaux. C’est sans doute le cas de la pierre tombale déjà mentionnée de Lello Maddaleni : si la marque mercantile sur le coussin, à droite de la tête du défunt, semble presque entièrement disparue à la suite de l’usure généralement caractéristique des pierres tombales insérées dans le pavement (comme est désormais peu visible le petit écu aussi placé sur le coussin à gauche de la tête de Maddaleni), on ne peut exclure que la marque bien plus évidente entre les jambes du défunt, attestée par la xylographie et dont il ne reste nulle trace, ait été effacée intentionnellement (fig. 56).}
              \teiP[xmlid={p21-3}]{Il est certain que l’altération de certains souvenirs familiaux gênants fut pratiquée dans les hauts rangs de la curie pontificale : les Barberini firent effacer sur quelques anciens écus de famille les ciseaux qu’un de leurs ancêtres, Taddeo di Cecco, y avait fait figurer pour signaler le métier de tisserand exercé par lui et par sa famille. Nous savons en particulier qu’en 1636 le cardinal Francesco, par l’intermédiaire du grand érudit et historien Carlo di Tommaso Strozzi, ordonna à un certain Zanobi Radicchi de supprimer les ciseaux dans certains écus de Taddeo présents à Barberino di Val d’Elsa. La lettre où Radicchi annonce au cardinal qu’il a accompli sa mission nous restitue pour ainsi dire une scène de comédie à l’italienne, avec l’émissaire du cardinal, le curé et le maçon qui « expédient l’affaire \lbrack{}\teiHi[rend={italic}]{fanno il negotio}\rbrack{} » en pleine nuit, pour n’être surpris par personne :}
              \begin{teiCit}[xmlid={cit06}]

                \begin{teiQuote}
Éminentissime Seigneur et Patron très respectable, / Ces derniers jours on m’a fait entendre par le seigneur Carlo Strozzi que Votre Seigneurie Éminentissime aurait aimé que les ciseaux fussent enlevés des trois armoiries qui sont ici à Barberino ; n’ayant d’autre désir que de satisfaire à l’obligation que j’ai de la servir, professant toujours que ses moindres signes me sont des commandements très exprès, dès que j’entendis cela je vins à Barberino, et parce que deux de sesdites armes sont dans l’église de ce château, il me fut nécessaire d’en conférer avec le curé de cette église, le sachant prêtre à qui l’on pourrait confier toute plus importante affaire et ayant pris rendez-vous avec lui pour le soir où cela devait se faire j’arrivai sur les lieux à la deuxième heure de la nuit avec le maçon du monastère et jusqu’à minuit ledit prêtre, le maçon et moi nous expédiâmes l’affaire avec tant de dextérité qu’il n’a pas été su que personne s’en soit aperçu et qu’il ne semble pas qu’on y soit jamais allé, et tant que je serai jugé digne de ses commandements je ferai en sorte qu’elle soit servie avec toute diligence. Pour ne pas l’ennuyer plus, je terminerai demandant pour elle au Seigneur avec toute ma famille le comble de toute félicité. De Barberino le 9 novembre 1636 / de Votre Seigneurie Éminentissime / très humble et très obligé serviteur / Zanobi Radicchi\teiNote[n={46},place={foot},xmlid={ftn92}]{\teiP[xmlid={p-141}]{La lettre a été publiée par Pio Pecchiai, \teiHi[rend={italic}]{Un assassinio politico a Roma nel Cinquecento}, Rome, Biblioteca d’Arte Editrice, 1956, p. 91 : « \teiHi[rend={italic}]{Eminentissimo Signore e Padron Colendissimo, / Mi fu accennato alli giorni passati dal signor Carlo Strozzi come Sua Signoria Eminentissima haverebbe hauto gusto che si fosser levate le forbice dalle tre arme che sono qui in Barberino; onde non havendo altro desiderio che di satisfare all'obligo che tengo di servirla, professando sempre che li di lor minimi cenni mi sieno espressissimi comandamenti, subito che ciò hebbi inteso me ne venni a Barberino; e perché dua delle dette lor armi sono nella chiesa di questo castello, mi fu necessario conferirlo con il curato di detta chiesa, conoscendolo per sacerdote da poterli confidare qual si voglia più importante negotio, e rimanendo seco in appuntamento della sera che si doveva fare, arrivai in quella alle due hore di notte con il muratore del monasterio, et in su la meza notte il detto prete, il muratore et io facemmo il negotio con tanta destrezza che non s'è sentito che alcuno se ne sia accorto, né si pare che vi sieno mai state, e sempre che sarò fatto degno de' suoi comandamenti cercherò con ogni diligentia resti servita. Per non l'infastidir maggiormente farò fine, pregandoli sempre con tutta la mia famiglia dal Signore il colmo d'ogni felicità. Di Barberino li 9 novembre 1636 / di Vostra Signoria Eminentissima / humilissimo et obbligatissimo servitore / Zanobi Radicchi.} » Sur cet épisode et en général sur les rapports des Barberini avec leur passé, voir Michele Di Monte, « La cultura antica e lo specchio della storia », dans Maurizia Cicconi, Flaminia Gennari Santori et Sebastian Schütze (dir.), \teiHi[rend={italic}]{L’immagine sovrana. Urbano VIII e i Barberini}, catalogue d’exposition, Rome, Officina libraria, 2023, p. 265-269. Les écus de Barberino de Val d’Elsa ne furent probablement pas les seuls à être intéressés par la « censure » : deux écus de Taddeo di Cecco à Santa Croce à Florence, l’un à l’intérieur de la basilique et l’autre dans le grand cloître, semblent porter des traces d’une intervention destinée à supprimer ces encombrants ciseaux.}}.
\end{teiQuote}
              
\end{teiCit}
              \teiP[xmlid={p22-3}]{Le travail fut si bien exécuté que, aux dires de Radicchi, il ne semblait pas qu’il y ait jamais eu de ciseaux sur ces armoiries. Qui contrôlait le présent était en mesure, comme chez Orwell, de contrôler le passé, modifiant ses traces de façon à pouvoir le réécrire. Mais parfois le passé n’est pas aussi docile ni aussi facile à domestiquer : un écu des Barberini doté des ciseaux échappa à l’intervention de Radicchi, il peut encore être contemplé aujourd’hui dans la petite ville toscane, dans la sacristie de la chapelle de San Bartolo (fig. 58), nous rappelant les origines humbles et laborieuses de la famille du pape Urbain VIII\teiNote[n={47},place={foot},xmlid={ftn93}]{\teiP[xmlid={p-142}]{La localisation de l’écu n’est pas celle d’origine. À Barberino, dans la chapelle de San Bartolo, et à Tavarnelle, dans l’église paroissiale de Santa Lucia al Borghetto, se trouvent aussi deux écus des Barberini « traités ». Dans la lettre de Radicchi, il est d’abord question de modifier « trois armoiries », mais ensuite on ne parle que d’une intervention sur deux écus, conservés « dans l’église de ce château », ou dans l’église paroissiale de San Bartolomeo : il semble toutefois improbable que l’envoyé du cardinal ait oublié un écu ; plutôt, dans la lettre, il parle de ces deux-là pour expliquer l’implication du prêtre, alors qu’un troisième écu qui fut privé de ses ciseaux devait se trouver ailleurs, peut-être à l’extérieur. Mais il y en eut un quatrième, dont on ignorait l’existence et qui pour cette raison échappa à la censure. Par ailleurs, comme on le sait, la présence des ciseaux n’était pas le seul élément gênant dans l’héraldique des Barberini : nous faisons naturellement référence à la présence originelle dans les armes de la famille non de gracieuses abeilles mais de taons \lbrack{}\teiHi[rend={italic}]{tafani}\rbrack{}, dont il semble que le cardinal Maffeo, bien avant de monter sur le trône de Pierre, voulut ennoblir l’aspect. À l’évidence, il aurait été beaucoup plus difficile, sinon impossible, d’intervenir sur des taons, car leur modification en abeilles aurait entraîné en pratique la destruction des écus d’origine. On doit écarter l’idée, qui continue à jouir d’une certaine fortune, selon laquelle le nom même de la famille fut à l’origine Tafani : comme l’a déjà fait remarquer Pio Pecchiai, « le nom de Tafani ne se trouve jamais mentionné dans les documents notariés », on le fit probablement dériver à posteriori « des armes adoptées par la famille », originaire d’une zone du fief de Barberino dénommée Tafania (Pio Pecchiai, \teiHi[rend={italic}]{I Barberini}, Rome, Biblioteca d’Arte Editrice, 1959, p. 33). }}.}
              \teiP[xmlid={p23-3}]{Dans les écrits comme dans les événements que nous avons évoqués, on remarque avec quelle attention la Rome du \teiHi[rend={small-caps}]{xvii}\teiHi[rend={sup}]{e} siècle considérait l’héraldique : on s’interrogeait sur ses origines, on en reconstituait l’histoire, on veillait à sa conservation ou bien au contraire on y intervenait pour effacer certains aspects gênants. Cette attention pour le passé était toujours reliée au présent, cela ne surprend pas pour une société d’Ancien Régime qui, malgré de notables transformations, demeurait dans une certaine continuité avec le Moyen Âge.}
              \begin{teiFigure}[xmlid={figure01-3}]

                \teiGraphic[url={../icono/br/Ch04\_Federici/Fig44.jpg}]
                \teiHead[xmlid={figure-43-gauges-de-gozze-se-dalle-armi-o-insegne-che-parlano-roma-vitale-mascardi-1637-page-de-titre-bncr-001}]{Figure 43. Gauges de’ Gozze, \teiHi[rend={italic}]{Se dalle armi, o insegne, che parlano }\lbrack{}...\rbrack{}, Roma, Vitale Mascardi, 1637, page de titre, BNCR.}
              
\end{teiFigure}
              \begin{teiFigure}[xmlid={figure02-3}]

                \teiGraphic[url={../icono/br/Ch04\_Federici/Fig45.jpg}]
                \teiHead[xmlid={figure-44-silvestro-pietrasanta-tesserae-gentilitiae-romae-typis-haeredum-francisci-corbelletti-1638-page-de-titre-gravee-par-johann-friedrich-greuter-sur-un-dessin-de-giovan-francesco-romanelli-001}]{Figure 44. Silvestro Pietrasanta, \teiHi[rend={italic}]{Tesserae gentilitiae}, Romae, Typis haeredum Francisci Corbelletti, 1638, page de titre gravée par Johann Friedrich Greuter sur un dessin de Giovan Francesco Romanelli.}
              
\end{teiFigure}
              \begin{teiFigure}[xmlid={figure03-3}]

                \teiGraphic[url={../icono/br/Ch04\_Federici/Fig46.jpg}]
                \teiHead[xmlid={figure-45-silvestro-pietrasanta-tesserae-gentilitiae-p-661-001}]{Figure 45. Silvestro Pietrasanta, \teiHi[rend={italic}]{Tesserae gentilitiae}, p. 661.}
              
\end{teiFigure}
              \begin{teiFigure}[xmlid={figure04-3}]

                \teiGraphic[url={../icono/br/Ch04\_Federici/Fig47.jpg}]
                \teiHead[xmlid={figure-46-silvestro-pietrasanta-tesserae-gentilitiae-p-677-detail-des-effigies-de-richard-ii-et-de-la-reine-anne-sur-le-tableau-perdu-du-college-anglais-001}]{Figure 46. Silvestro Pietrasanta, \teiHi[rend={italic}]{Tesserae gentilitiae}, p. 677, détail des effigies de Richard II et de la reine Anne sur le tableau perdu du Collège anglais.}
              
\end{teiFigure}
              \begin{teiFigure}[xmlid={figure05-3}]

                \teiGraphic[url={../icono/br/Ch04\_Federici/Fig49.jpg}]
                \teiHead[xmlid={figure-47-xylographie-pour-les-memorie-sepolcrali-figurant-la-pierre-tombale-perdue-dangelotto-normanni-1378-autrefois-dans-leglise-de-santantonio-a-camaldoli-collection-particuliere-fabrizio-federici-001}]{Figure 47. Xylographie pour les \teiHi[rend={italic}]{Memorie sepolcrali} figurant la pierre tombale perdue d’Angelotto Normanni († 1378), autrefois dans l’église de Sant’Antonio à Camaldoli, collection particulière, Fabrizio Federici.}
              
\end{teiFigure}
              \begin{teiFigure}[xmlid={figure06-3}]

                \teiGraphic[url={../icono/br/Ch04\_Federici/Fig50.JPG}]
                \teiHead[xmlid={figure-48-ecu-deverso-ii-dellanguillara-premiere-moitie-du-xv-e-siecle-rome-hopital-du-latran-001}]{Figure 48. Écu d’Everso II dell’Anguillara, première moitié du \teiHi[rend={small-caps}]{xv}\teiHi[rend={sup}]{e} siècle, Rome, Hôpital du Latran.}
              
\end{teiFigure}
              \begin{teiFigure}[xmlid={figure07-3}]

                \teiGraphic[url={../icono/br/Ch04\_Federici/Fig51.jpg}]
                \teiHead[xmlid={figure-49-armes-de-la-famille-casali-1429-rome-palais-casali-001}]{Figure 49. Armes de la famille Casali, 1429, Rome, Palais Casali.}
              
\end{teiFigure}
              \begin{teiFigure}[xmlid={figure08-2}]

                \teiGraphic[url={../icono/br/Ch04\_Federici/Fig52.jpg}]
                \teiHead[xmlid={figure-50-xylographie-representant-les-armes-de-la-famille-casali-dapres-g-b-casali-de-profanis-et-sacris-veteribus-ritibus-romae-ex-typographia-andreae-phaei-1645-bibliotheque-municipale-de-grenoble-c-1969-001}]{Figure 50. Xylographie représentant les armes de la famille Casali, d’après G. B. Casali, \teiHi[rend={italic}]{De profanis, et sacris veteribus ritibus}, Romae, Ex typographia Andreae Phaei, 1645, bibliothèque municipale de Grenoble, C.1969.}
              
\end{teiFigure}
              \begin{teiFigure}[xmlid={figure09-2}]

                \teiGraphic[url={../icono/br/Ch04\_Federici/Fig53.jpg}]
                \teiHead[xmlid={figure-51-xylographie-pour-le-traite-des-memorie-sepolcrali-figurant-lecu-gualdi-sur-le-flanc-nord-du-palais-senatorial-collection-particuliere-fabrizio-federici-001}]{Figure 51. Xylographie pour le traité des \teiHi[rend={italic}]{Memorie sepolcrali} figurant l’écu Gualdi sur le flanc nord du Palais sénatorial, collection particulière, Fabrizio Federici.}
              
\end{teiFigure}
              \begin{teiFigure}[xmlid={figure10-2}]

                \teiGraphic[url={../icono/br/Ch04\_Federici/Fig54.jpg}]
                \teiHead[xmlid={figure-52-armes-dun-senateur-de-la-famille-gualdi-galeotto-ou-francesco-premiere-moitie-du-xvi-e-siecle-rome-palais-senatorial-001}]{Figure 52. Armes d’un sénateur de la famille Gualdi, Galeotto ou Francesco, première moitié du \teiHi[rend={small-caps}]{xvi}\teiHi[rend={sup}]{e} siècle, Rome, Palais sénatorial.}
              
\end{teiFigure}
              \begin{teiFigure}[xmlid={figure11-2}]

                \teiGraphic[url={../icono/br/Ch04\_Federici/Fig55.JPG}]
                \teiHead[xmlid={figure-53-pierre-tombale-de-giuliano-porcari-1282-rome-san-giovanni-della-pigna-001}]{Figure 53. Pierre tombale de Giuliano Porcari († 1282), Rome, San Giovanni della Pigna.}
              
\end{teiFigure}
              \begin{teiFigure}[xmlid={figure12-2}]

                \teiGraphic[url={../icono/br/Ch04\_Federici/Fig56.jpg}]
                \teiHead[xmlid={figure-54-xylographie-pour-le-traite-des-memorie-sepolcrali-figurant-la-pierre-tombale-de-pietro-muti-1390-dans-leglise-de-san-pantaleo-collection-particuliere-fabrizio-federici-001}]{Figure 54. Xylographie pour le traité des \teiHi[rend={italic}]{Memorie sepolcrali} figurant la pierre tombale de Pietro Muti († 1390) dans l’église de San Pantaleo, collection particulière, Fabrizio Federici.}
              
\end{teiFigure}
              \begin{teiFigure}[xmlid={figure13-2}]

                \teiGraphic[url={../icono/br/Ch04\_Federici/Fig57.jpg}]
                \teiHead[xmlid={figure-55-xylographie-pour-le-traite-des-memorie-sepolcrali-figurant-la-pierre-tombale-de-lello-maddaleni-1390-dans-leglise-de-santa-maria-sopra-minerva-collection-particuliere-fabrizio-federici-001}]{Figure 55. Xylographie pour le traité des \teiHi[rend={italic}]{Memorie sepolcrali} figurant la pierre tombale de Lello Maddaleni († 1390) dans l’église de Santa Maria sopra Minerva, collection particulière, Fabrizio Federici.}
              
\end{teiFigure}
              \begin{teiFigure}[xmlid={figure14-2}]

                \teiGraphic[url={../icono/br/Ch04\_Federici/Fig58.JPG}]
                \teiHead[xmlid={figure-56-pierre-tombale-de-lello-maddaleni-rome-santa-maria-sopra-minerva-001}]{Figure 56. Pierre tombale de Lello Maddaleni, Rome, Santa Maria sopra Minerva.}
              
\end{teiFigure}
              \begin{teiFigure}[xmlid={figure15-2}]

                \teiGraphic[url={../icono/br/Ch04\_Federici/Fig59.jpg}]
                \teiHead[xmlid={figure-57-armes-de-la-famille-cavalieri-a-lentree-de-la-tour-du-papito-rome-largo-di-torre-argentina-001}]{Figure 57. Armes de la famille Cavalieri à l’entrée de la Tour du Papito, Rome, Largo di Torre Argentina.}
              
\end{teiFigure}
              \begin{teiFigure}[xmlid={figure16-2}]

                \teiGraphic[url={../icono/br/Ch04\_Federici/Fig60.jpg}]
                \teiHead[xmlid={figure-58-armes-de-taddeo-di-cecco-da-barberino-2-e-moitie-du-xiv-e-siecle-barberino-di-val-delsa-chapelle-de-san-bartolo-sacristie-001}]{Figure 58. Armes de Taddeo di Cecco da Barberino, 2\teiHi[rend={sup}]{e} moitié du \teiHi[rend={small-caps}]{xiv}\teiHi[rend={sup}]{e} siècle, Barberino di Val d’Elsa, chapelle de San Bartolo, sacristie.}
              
\end{teiFigure}
            
\end{teiDiv}
          
\end{teiElement}
        
\end{teiElement}
        \begin{teiElement}[name={group},type={chapter},data-page-title={L’héraldique d’Innocent X Pamphili (1644-1655)},data-page-subtitle={Origines, architecture, décoration},data-page-authors={Daniela del Pesco},data-include-href={Ch05\_Heraldique\_Innocent.xml},xmlid={chapter-002}]

          \begin{teiElement}[name={front}]

            \begin{teiDiv}[type={titlePage},xmlid={titlepage-007}]

              \teiP[rend={title-main},xmlid={p-143}]{L’héraldique d’Innocent X Pamphili (1644-1655)}
              \teiP[rend={title-sub},xmlid={p-144}]{Origines, architecture, décoration}
              \teiP[rend={author-aut},xmlid={p-145}]{Daniela del Pesco}
              \teiP[rend={editor-trl},xmlid={p-146}]{Traduit de l’italien par Yvan Loskoutoff.}
            
\end{teiDiv}
          
\end{teiElement}
          \begin{teiElement}[name={body},xmlid={body-007}]

            \teiP[xmlid={p1-7}]{Dirk Amayden, juriste et érudit originaire du Brabant espagnol (1586-1656), dans son \teiHi[rend={italic}]{Histoire des familles romaines}, écrit : « Ma bonne fortune a voulu que dès l’enfance j’aie été familier de cette maison \lbrack{}les Pamphili\rbrack{} favorisée par la fortune et favorisant la fortune, anciens et nobles de la première classe à Gubbio depuis des centaines d’années\teiNote[n={1},place={foot},xmlid={ftn1-5}]{\teiP[xmlid={p-147}]{\teiHi[rend={italic}]{La Storia delle famiglie romane, con note ed aggiunte del Comm. Carlo Augusto Bertini} \lbrack{}Rome, Collegio araldico, 1910\rbrack{}\teiHi[rend={italic}]{,} Bologne, Forni, 1967, vol. 2, p. 124 et 125 (entrée « Pamphilj », p. 124-127). Amayden arrive à Rome pour l’année jubilaire 1600, en 1601-02 il s’y installe pour y mener ses recherches.}}. »}
            \teiP[xmlid={p2-6}]{Le progrès économique et social de cette maison se déroula de fait sur la longue durée dans la cité ombrienne. Le 10 novembre 1460, Sixte IV nomme Antonio Pamphili docteur en droit civil et ecclésiastique, procurateur fiscal de la Chambre Apostolique : il se transfère à Rome, initiant l’ascension de sa famille dans le milieu de la Curie pontificale. En 1604, Girolamo Pamphili, neveu d’Antonio, est élu cardinal\teiNote[n={2},place={foot},xmlid={ftn80-2}]{\teiP[xmlid={p-148}]{Girolamo Pamphili (1545-1610), docteur \teiHi[rend={italic}]{in utroque iure}, décan de la Sainte Rote (à partir de 1602), cardinal prêtre de San Biagio dell’Anello (à partir de 1604), vicaire général de Sa Sainteté pour le diocèse de Rome (de 1605 à 1610). Très lié aux Oratoriens, il résida au palais de la place Navone, le partageant avec son frère Camillo. Vers 1606 furent entrepris des travaux d’adaptation fonctionnelle à la dignité cardinalice de Girolamo. Benedetta Borello, « Pamphili, Girolamo », dans \teiHi[rend={italic}]{DBI}, Rome, Istituto della Enciclopedia italiana, 2014, vol. 80, 2014, p. 670 et 671.}}, comme en 1630 son neveu Giovan Battista, qui va être élu pape en 1644. Une gravure, conçue par Giovanni Luigi Valesio en 1606, représente l’écu de Girolamo Pamphili entouré des attributs cardinalices, inséré entre des figures allégoriques, sur le fond d’un ciel peuplé de colombes\teiNote[n={3},place={foot},xmlid={ftn81-2}]{\teiP[xmlid={p-149}]{Babette Bohn, \teiHi[rend={italic}]{Ludovico Carracci and the art of drawing}, Turnhout, Brepols, 2004, p. 377, R 59. Un dessin, signé et daté de 1606, a servi de modèle pour la gravure (fig. 1) incluse dans le recueil intitulé : \teiHi[rend={italic}]{Stampe bolognese di Giovanni Valesio, inventore e intagliatore, }qui présente de nombreux écus de cardinaux de l’époque. La gravure est conservée à Bologne, Pinacoteca Nazionale, Gabinetto Disegni e Stampe (290 × 380 mm). Voir Giovanna Gaeta Bertelà (éd.), avec la collaboration de Stefano Ferrara, \teiHi[rend={italic}]{Incisori bolognesi ed emiliani del sec. XVII}, Bologne, Associazione per le arti Francesco Francia, \lbrack{}1973\rbrack{}, n\teiHi[rend={sup}]{o} 1009. L’image des armes Pamphili avec les attributs cardinalices se rencontre dans diverses estampes représentant le portrait de Girolamo Pamphili.}} (fig. 59). Dans l’église romaine de Santa Maria in Vallicella, la pierre tombale du cardinal récapitule son \teiHi[rend={italic}]{cursus honorum}\teiNote[n={4},place={foot},xmlid={ftn82-2}]{\teiP[xmlid={p-150}]{L’épigraphe se conclut sur les mots : \teiHi[rend={italic}]{IOANNES BAPTISTA PAMPHILIUS AB VRBANO VIII CARD CREATUS ET PAMPHILIUS FRATRIS FILII MONUMENTUM POSUERE}, c’est-à-dire qu’il est affirmé que « posèrent la pierre tombale » les fils de Camillo, frère de Girolamo (« \teiHi[rend={italic}]{Fratris filii} »), soit Pamphilio (1563-1639, mari d’Olimpia Maidalchini) et Giovan Battista (1574-1655), futur Innocent X.}} (fig. 60). Dessinée par Paolo Maruscelli en 1634, elle montre, sous l’épigraphe, gravée sur une élégante dalle de marbre noir, l’image d’une colombe blanche appuyée sur deux rameaux d’olivier verts et feuillus\teiNote[n={5},place={foot},xmlid={ftn83-2}]{\teiP[xmlid={p-151}]{Lionello Neppi, \teiHi[rend={italic}]{Schede del catalogo di opere d’arte mobili (Santa Maria in Vallicella)}, Rome, Ministero della Pubblica Istruzione, 1957, n. 90, l. 389. Girolamo est le premier des Pamphili à ne pas avoir été enterré à San Lorenzo in Damaso, où la famille était titulaire de l’autel de San Domenico (deuxième dans la nef de gauche), par ailleurs détruit lors du réaménagement de l’église (Adam Poos, \teiHi[rend={italic}]{San Lorenzo in Damaso a Roma}, \teiHi[rend={italic}]{Guida alla basilica}, Rome, GB Editoria, 2015, p. 8).}}.}
            \teiP[xmlid={p3-6}]{Du point de vue de l’héraldique et de la compréhension du message confié à l’écu des Pamphili, nous verrons comment cette image, qui se diffuse à Rome comme conséquence de l’élévation à la pourpre de Girolamo Pamphili, était déjà présente dans les phases précédentes de l’histoire de la famille. Le meuble de la « colombe au rameau d’olivier dans le bec », qui remplit l’écu, se révèle très malléable pour exprimer des significations propres à différentes circonstances. Nous chercherons à les décrypter : en fait, comme on le lit dans le \teiHi[rend={italic}]{Miroir symbolique des armes nobiliaires} d’Andrea Cellonese (1667), la création d’un écu répond à la nécessité, propre aux maisons en ascension, de se forger une « devise symbolique qui montre une chose à la vue et une autre à l’intellect, faisant allusion aux vertus de la famille\teiNote[n={6},place={foot},xmlid={ftn84-2}]{\teiP[xmlid={p-152}]{Andrea Cellonese, \teiHi[rend={italic}]{Specchio simbolico delle armi gentilizie, }Naples, Giovanni Francesco Paci, 1667, p. 46 : « \teiHi[rend={italic}]{impresa simbolica che altro mostra alla vista et altro all’intelletto significa e alle virtù della famiglia allude} ». }} ».}
            \begin{teiDiv}[type={section1},xmlid={div01-4}]

              \teiHead[xmlid={les-lettres-du-milieu-romain-des-pamphili-et-lorigine-de-leur-embleme-nobiliaire-001}]{Les lettrés du milieu romain des Pamphili et l’origine de leur emblème nobiliaire}
              \teiP[xmlid={p4-6}]{Amayden, dans son ouvrage consacré aux familles romaines, se pose la question de l’origine du nom des Pamphili, « vocable grec », sans trouver de réponse satisfaisante. Il affirme utiliser les écrits de Giovan Battista Cantelmaggio « qui a beaucoup travaillé à chercher et à trouver l’origine et les progrès de cette maison ». Il rapporte en conséquence les données qu’il retient comme documentées : en particulier la présence des Pamphili à Gubbio depuis 1150\teiNote[n={7},place={foot},xmlid={ftn85-2}]{\teiP[xmlid={p-153}]{\teiHi[rend={italic}]{La Storia delle famiglie romane...},\teiHi[rend={italic}]{ op. cit.}, p. 125 : « \teiHi[rend={italic}]{che molto ha faticato in cercare e trovare l’origine e i progressi di questa Casa} ». Giovan Battista Cantelmaggio était auditeur de Rote, c’est-à-dire membre du tribunal ecclésiastique de la Sainte Rote romaine. Dirk Amayden affirme (p. 127) que l’inventeur du beau nom de Pamphilio voulut lui donner une devise correspondante et rapporte qu’Andrea de’ Mantici, descendant de Pamphilio, fut un fameux médecin employé par Louis, roi de Hongrie en 1345 et que l’empereur Frédéric III concéda un très noble privilège à Antonio Pamphili en 1461, monseigneur du ministère public \lbrack{}\teiHi[rend={italic}]{fiscale}\rbrack{} sous Sixte IV en 1473, donnant, comme on l’a dit, origine à l’histoire romaine de la famille. Amayden exclut que l’origine du nom « Pam-filio », « ami de tous », terme grec, ait été déterminé par la reconnaissance envers la bénéfique activité médicale de membres de la maison.}}.}
              \teiP[xmlid={p5-6}]{Au milieu du \teiHi[rend={small-caps}]{xvii}\teiHi[rend={sup}]{e} siècle, l’écu Pamphili est ainsi décrit dans un document de l’archive Doria Pamphili, attribuable à Amayden : « Une colombe, le rameau d’olivier au bec, double symbole de paix, symbolisant \lbrack{}…\rbrack{} l’amour et l’amitié compagnons éternels de la quiétude, éloignés du fracas des guerres. Dans la partie supérieure de l’écu, on voit trois lys d’or en chef d’azur ordinaire mais erronément divisé par des pals de gueules qui ont été interprétés comme le rabat ou retombée d’un baldaquin du roi de France, je ne sais sur quel fondement\teiNote[n={8},place={foot},xmlid={ftn86-2}]{\teiP[xmlid={p-154}]{« \teiHi[rend={italic}]{Una colomba con l’uliva in bocca, dupplicato simbolo di Pace, simbolizzante }\lbrack{}…\rbrack{}\teiHi[rend={italic}]{ l’amore ed amicizia compagni eterni della quiete, e lontani da strepiti di guerra. Nella parte superiore dello scudo si vedono tre gigli d’oro in campo azzurro traversato commune, ma erroneamente con asse rossi, i quali vengono interpretati da alcuno, che rappresentino il lembo, o vero ricaduta di un Baldacchino del Re di Francia. Non so con che fondamento }», ADP, Fonds Archiviolo, n\teiHi[rend={sup}]{o} 145, transcrit dans Giuseppe Simonetta, Laura Gigli et Gabriella Marchetti, \teiHi[rend={italic}]{Sant’Agnese in Agone a Piazza Navona. Immagine, Luce, Ordine, Suono nelle fabbriche Pamphili, }Rome, Gangemi, 2005, appendice 2, p. 208, et aussi, avec de petites variantes dans le texte, Dirk Amayden, \teiHi[rend={italic}]{La Storia delle famiglie romane…}, \teiHi[rend={italic}]{op. cit.}, p. 127. Sur l’écu, voir aussi : Giovanni Battista Di Crollalanza\teiHi[rend={italic}]{, Dizionario storico blasonico delle famiglie nobili e notabili italiane estinte e fiorenti}, Pise, Presso la direzione del Giornale araldico, 1888, vol. II, p. 266. Donald Lindsay Galbreath, \teiHi[rend={italic}]{Papal Heraldry}, Cambridge, W. Heffer, 1930, part. I, p. 98, cite un exemple de l’écu Pamphili vu par lui au « musée de Pérouse », qu’il date du \teiHi[rend={small-caps}]{xv}\teiHi[rend={sup}]{e} siècle et mentionne aussi celui qui figure sur la pierre tombale à San Pietro in Montorio à Rome, avec la date de 1600, date non rencontrée sur la sépulture qui porte l’inscription « DOM, Pamphiliis/ad Diem novissimam ». }} » (fig. 61). Le même document affirme que, dans cet écu, « les lys ne sont ni français, ni espagnols, mais signes et marques des factions comme en usait jadis la noblesse dans les cités italiennes \lbrack{}...\rbrack{}, ou que les trois lys faisant partie des armes de la ville de Gubbio, les Pamphili y étant la principale famille noble, ils ont été honorés desdits lys. \lbrack{}...\rbrack{} Dans les armes Pamphili le champ de gueules n’est qu’un mur aux créneaux à l’ancienne dans lesquels furent plantés des lys selon l’usage des maisons fortes\teiNote[n={9},place={foot},xmlid={ftn87-2}]{\teiP[xmlid={p-155}]{« \teiHi[rend={italic}]{i gigli non sono né francesi, né spagnuoli, ma segni, e note de’fazzioni come anticamente usava la nobiltà nelle città d’Italia }\lbrack{}…\rbrack{} \teiHi[rend={italic}]{overo che, sendo i tre gigli parte dell’arme della Città d’Augubio ed i Pamphili, più principale, e più nobile, siano stati onorati di detti gigli.}\lbrack{}…\rbrack{}\teiHi[rend={italic}]{ Nell’arme pamphilia, il campo rosso non è che un muro colli merli all’antica, in mezzo alli quali fossero piantati gigli a usanza di casa forte} ».}} ». Il n’est pas surprenant qu’Amayden se fasse le défenseur d’une origine « autochtone » de l’écu des Pamphili, exempt de l’influence de puissances étrangères, telles l’Espagne ou la France, la référence à cette dernière s’exprimant par la présence de pals de gueules pour séparer les meubles des lys, alors qu’Amayden les voudrait au contraire en chef d’azur uni, annulant la référence française.}
              \teiP[xmlid={p6-6}]{Le polygraphe de Gubbio Vincenzo Armanni (1608-1684), dans ses \teiHi[rend={italic}]{Réflexions sur la célébrité de l’Excellentissime maison Pamphili de l’an 850 à 1300}\teiNote[n={10},place={foot},xmlid={ftn88-2}]{\teiP[xmlid={p-156}]{\teiHi[rend={italic}]{Riflessioni intorno alla chiarezza dell’Eccellentissima Casa Panphilia dall’anno 850 fino al 1300, }Archive Doria Pamphili, scaff. 93, busta 54, interno 2, transcrit dans Giuseppe Simonetta, Laura Gigli et Gabriella Marchetti, \teiHi[rend={italic}]{op. cit.}, appendice 2, p. 209 et 210. Voir aussi : Vincenzo Armanni, \teiHi[rend={italic}]{Delle Lettere del Signor V. A., colla vita dell'autore per Carlo Cartari}, Rome, Iacomo Dragondelli, t. I, 1663 ; \teiHi[rend={italic}]{ibid}., t.\teiHi[rend={italic}]{ }II et III, Macerata, Giuseppe Piccini, 1674. Dans une lettre, Armanni, souvent accusé de manque de sérieux, défend la rigueur de sa méthode de constitution des généalogies (vol. II, p. 88). Sur Armanni et la généalogie des Pamphili, Gaetano Moroni, \teiHi[rend={italic}]{Dizionario di erudizione storico ecclesiastica}, In Venezia, dalla Tipografia Emiliana, 1840-1878, vol. XXXIII, 1845, « Gubbio », p. 162 ; Umberto Coldagelli, « Armanni, Vincenzo »,\teiHi[rend={italic}]{ DBI}, \teiHi[rend={italic}]{op. cit}., 1962, vol. 4, p. 222-224. L’activité d’Armanni fut certainement encouragée par Camillo Pamphili, bibliophile passionné et collectionneur d’art, fortement engagé dans l’élaboration de l’image de prestige de la famille.}}\teiHi[rend={italic}]{, }examine les « armes » de saint Lodolfo Pamphili, évêque de Gubbio en 967, « dit Colombino, et sa congrégation dite de la Colombe, étant très probable que la cause principale de cette dénomination fut la colombe, armes ou insigne de ce saint\teiNote[n={11},place={foot},xmlid={ftn89-2}]{\teiP[xmlid={p-157}]{« \teiHi[rend={italic}]{detto Colombino, e la sua congregazione detta della Colomba ond’è probabilissimo che la cagion principale di così fatta denominazione fosse la colomba, Arme o insegna di esso Beato}… », Archive Doria Pamphili, scaff. 93, busta 54, interno 2, transcrit dans Giuseppe Simonetta, Laura Gigli et Gabriella Marchetti, \teiHi[rend={italic}]{op. cit.}, p. 209 et 210.}}… ». En examinant les conjectures disponibles, Armanni revendique le rôle de Lodolfo dans la fondation de l’ordre des camaldules et affirme que « nous ne pouvons douter que dès ce temps-là cette maison ait eu pour armes la colombe\teiNote[n={12},place={foot},xmlid={ftn90-2}]{\teiP[xmlid={p-158}]{\teiHi[rend={italic}]{Ibid}., « \teiHi[rend={italic}]{non possiamo dubitare che questa casa havesse anche hallora per arme sua la Colomba }»\teiHi[rend={italic}]{.}}} ». Il ajoute qu’au palais des Pamphili dans le quartier San Lorenzo à Gubbio « on voit peintes depuis plusieurs centaines d’années les mêmes armes et insigne d’une colombe d’argent en champ de gueules qui porte au bec un rameau d’olivier, surmontée de trois lys d’or entre les dents d’un râteau de gueules en champ d’azur. En certains endroits pourtant, particulièrement au plafond de ce palais, on voit au-dessus de la colombe pendre trois drapeaux d’azur, séparés par un espace de gueules, chaque drapeau ayant en son milieu un lys d’or : le pontife régnant et les seigneurs de sa très excellente maison ont composé leurs armes de cette façon\teiNote[n={13},place={foot},xmlid={ftn91-2}]{\teiP[xmlid={p-159}]{\teiHi[rend={italic}]{Ibid., }« \teiHi[rend={italic}]{vedesi dipinta di più centinaia d’anni la medesima arme, et insegna di una colomba bianca, che porta nella bocca un ramoscello verde d’oliva in campo rosso, havendo di sopra tre gigli d’oro tra li denti d’un Rastrello rosso in Campo azzurro. In alcune parti però, spetialmente nel soffitto di esso Palazzo, veggonsi al di sopra della Colomba pendenti tre Drapelloni turchini, in ciascheduno de’quali lascia tra sè tanto spatio rosso, quanto basta a dividergli con havere ogni Drapellone in mezzo un giglio d’oro: et in tal maniera han seguitato a far l’arme il pontefice regnante, e i signori della sua eccellentissima Casa. }»}} ». L’allusion renvoie à Camillo Pamphili, neveu d’Innocent X, auquel Armanni dédie ses écrits.}
              \teiP[xmlid={p7-6}]{L’hypothèse du rôle du camaldule Lodolfo pour la constitution de l’écu des Pamphili est reprise par Gaetano Moroni dans son \teiHi[rend={italic}]{Dictionnaire d’érudition historico-ecclésiastique }(1851). Il affirme qu’un lien entre l’emblème de la colombe et cette famille se serait vérifié dans le milieu de l’ermitage de Fonte Avellana, fondé en 1009, alors que Lodolfo Pamphili participe à la création de l’ermitage consacré à la dévotion à la Sainte Croix, en consonance avec l’œuvre de Romuald de Ravenne (951/953-1027), créateur de l’ordre des camaldules et des monastères voisins de Fonte Avellana\teiNote[n={14},place={foot},xmlid={ftn92-2}]{\teiP[xmlid={p-160}]{Gaetano Moroni, \teiHi[rend={italic}]{op. cit.}, 1851, vol. LI, p. 85 (entrée « Pamphili »), écrit que Lodolfo, évêque de Gubbio fut fondateur de la congrégation Avellana, « qui fut dite aussi de la colombe d’après les armoiries de son fondateur ». Au vol. XXXIII, p. 171 (entrée « Gubbio »), on lit que l’ermitage de l’Avellana /Serra Sant’Abbondio (Pesaro et Urbino), affilié à la règle de saint Benoît, « devint vite un exemple de doctrine ». Voir \teiHi[rend={italic}]{Vita del b. Lodolfo Pamphilij vescovo di Gubbio, e abate fondatore della Congreg. Dell’Avellana, detta della Colomba descritta dal p.d. Agostino Romano Fiori monaco camaldolese}, Rome, Per Antonio de’ Rossi nella strada del Seminario romano, vicino alla Rotonda, 1722. Par la bulle de Pie V \teiHi[rend={italic}]{Quantum animus noster} du 10 décembre 1569, la congrégation de l’Avellana fut supprimée et, en 1570, les biens du monastère dévolus à l’ordre des camaldules. En 1578, par les dispositions de Grégoire XIII, ces mêmes biens passèrent à l’administration du Collège germanique. Au \teiHi[rend={small-caps}]{xvii}\teiHi[rend={sup}]{e} siècle, le calice présent dans l’écu fut surmonté d’une étoile ou d’une comète d’or et l’écu, d’un chapeau de prélat avec mitre et pastoral. }}. Toutefois, l’écu de l’ordre est actuellement documenté en 1183, à l’époque de la reconstruction de l’abbaye, par le sceau du prieur des camaldules\teiNote[n={15},place={foot},xmlid={ftn93-2}]{\teiP[xmlid={p-161}]{L’abbaye de l’Avellana fut reconstruite en 1171 et consacrée le 31 août 1197. Italo Moretti et Renato Stopani, \teiHi[rend={italic}]{Architettura romanica religiosa a Gubbio, }Florence, Salimbeni, 1973, p. 33-37.}}. Il fut ensuite adopté dans presque tous les monastères et congrégations camaldules : il montre, en champ d’azur, un calice d’or accosté de deux colombes d’argent s’abreuvant. La devise de la congrégation est : « \teiHi[rend={italic}]{Ego vobis, vos mihi} » \lbrack{}Je suis pour vous, et vous pour moi\rbrack{}. Les deux colombes représentent les cénobites et les ermites qui s’abreuvent au même calice, symbole du Christ\teiNote[n={16},place={foot},xmlid={ftn94}]{\teiP[xmlid={p-162}]{Giuseppe Cacciamani, \teiHi[rend={italic}]{Dizionario degli Istituti di Perfezione}, Milan, Paoline, 1974, vol. 1, col. 1724.}}.}
              \teiP[xmlid={p8-6}]{À l’entourage de Camillo Pamphili (1622-1666) appartient aussi le jésuite Nicolangelo Caferri, auteur d’une vie de Bartolomeo Sacchi, dit Plàtina (1421-1481), de biographies des pontifes et de nombreux écrits portant sur la généalogie des Pamphili. Dans son \teiHi[rend={italic}]{Syntema Vetustatis} (\teiHi[rend={italic}]{Abrégé de l’Antiquité,} Rome, 1670), compendium de l’histoire du monde des origines à 1665, Caferri traite de la geste antique de la maison, rendue illustre par l’aïeul Amanzio « héros germanique, dans l’entourage de Charles le Grand \lbrack{}Charlemagne\teiNote[n={17},place={foot},xmlid={ftn95}]{\teiP[xmlid={p-163}]{Angelo Caferri, \teiHi[rend={italic}]{Synthema Vetustatis}, Romae, Ex tipographia Iacobi Dragondelli, 1670, p. 84-88. Sur ses écrits, voir Giuseppe Simonetta, Laura Gigli et Gabriella Marchetti,\teiHi[rend={italic}]{ op. cit.}, p. 33, 43 et 209.}}\rbrack{} ». Cette hypothèse est aussi acceptée par Moroni, qui affirme que « Amanzio Pamphilj, de l’ordre équestre, \lbrack{}…\rbrack{} suivant Charlemagne en Italie, fut récompensé par celui-ci avec quelques châteaux et avec les armes des Français dont la famille fait encore usage\teiNote[n={18},place={foot},xmlid={ftn96}]{\teiP[xmlid={p-164}]{Gaetano Moroni, \teiHi[rend={italic}]{op. cit.}, vol. LI, p. 85 (entrée « Pamphili ») ; dans le même Moroni, \teiHi[rend={italic}]{ibid}., vol. XXXVI, p. 12 et 13, figure l’entrée biographique consacrée à Innocent X.}} ». Il rapporte en outre qu’en 917, Gubbio fut reconstruite grâce à l’œuvre de Pietro Pamphili, fils d’Amanzio\teiNote[n={19},place={foot},xmlid={ftn97}]{\teiP[xmlid={p-165}]{\teiHi[rend={italic}]{Ibid}., vol. XXXIII, p. 167 (entrée « Gubbio », p. 153-178).}}.}
              \teiP[xmlid={p9-5}]{L’attitude d’adulation propre à la littérature encomiastique et l’opinion que la noblesse d’une famille se mesure par l’antiquité de ses origines, portent les lettrés de l’entourage des Pamphili qui écrivent dans les années du pontificat d’Innocent X et ensuite dans le milieu de son neveu Camillo, à insister sur les origines médiévales de l’écu aux lys et à la colombe\teiNote[n={20},place={foot},xmlid={ftn98}]{\teiP[xmlid={p-166}]{Pour le processus de construction d’une noblesse « antique » comme élément destiné à conférer plus de prestige à une maison, particulièrement en Toscane, voir Vanna Arrighi et Elisabetta Insabato, « Tra storia e mito, la ricostruzione del passato familiare nella nobiltà toscana dei secoli XVI-XVIII »\teiHi[rend={italic}]{, }dans \teiHi[rend={italic}]{L’identità genealogica e araldica. Fonti, metodologie, interdisciplinarietà, }actes du colloque, Turin, Archive d’État, 21-26 septembre 1988, Rome, Ministero per i beni e le attività culturali, 2000, p. 1099-1121.}}.}
            
\end{teiDiv}
            \begin{teiDiv}[type={section1},xmlid={div02-4}]

              \teiHead[xmlid={la-colombe-portant-le-rameau-dolivier-au-bec-a-gubbio-et-a-rome-001}]{La « colombe portant le rameau d’olivier au bec » à Gubbio et à Rome}
              \teiP[xmlid={p10-5}]{Des représentations de l’emblème de la colombe tenant le rameau d’olivier dans son bec sont encore visibles à Gubbio sur des édifices liés à la présence des Pamphili.}
              \teiP[xmlid={p11-5}]{Leur demeure révèle un processus d’agrégation complexe initié vers « 1260 alors que Mantia, fils de Pamphilio \lbrack{}...\rbrack{} acquit une maison et un terrain dans le \teiHi[rend={italic}]{rione} San Martino, quartier dans lequel les Pamphili de Gubbio résidèrent au moins jusqu’à la fin du \teiHi[rend={small-caps}]{xvii}\teiHi[rend={sup}]{e} siècle\teiNote[n={21},place={foot},xmlid={ftn99}]{\teiP[xmlid={p-167}]{Benedetta Borello, « Alleanze matrimoniali e mobilità sociale e geografica. Il caso dei Pamphilj (XV-XVII secolo) », \teiHi[rend={italic}]{Mélanges de l'École française de Rome}, vol. 115, n\teiHi[rend={sup}]{o} 1, 2003, p. 345-366, ici p. 352 : « \teiHi[rend={italic}]{1260 quando Mantia, figlio di Pamphilio \lbrack{}…\rbrack{} acquistò una casa e un terreno nel rione San Martino, quartiere nel quale i Pamphili eugubini risiedettero almeno fino alla fine del Seicento.} »}} ». La construction et les transformations s’y succédèrent sans parvenir à une configuration homogène, laissant voir dans les parements muraux des modifications partielles et des restaurations effectuées au moyen de pierres et de briques de types différents\teiNote[n={22},place={foot},xmlid={ftn100}]{\teiP[xmlid={p-168}]{Au palais Pamphili, l’assemblage de corps de fabrique est évident : ils s’articulent autour d’une cour rectangulaire. Dans cette cour, à l'entrée du palais vers la place adjacente, s’inscrit un portique de douze colonnes à section octagonale et arcs en plein cintre. À côté, il y a aussi une petite cour à portique aux arcs sur quatre colonnes en \teiHi[rend={italic}]{pietra serena }(grès). }} (fig. 62).}
              \teiP[xmlid={p12-5}]{L’appartement de l’étage noble du palais conserve des décorations à emblèmes héraldiques qui apparaissent bien plus fragmentaires que celles décrites par Vincenzo Armanni. Une colombe au rameau d’olivier est sculptée dans le linteau en grès de la cheminée d’une salle proche de la salle principale (fig. 63). Là, des fresques assez endommagées représentent, au-dessus d’un lambris décoré de lys, une suite de vases enrichis d’éléments végétaux parmi lesquels on peut observer au moins une colombe (fig. 64).}
              \teiP[xmlid={p13-5}]{Dans la chapelle Pamphili de l’ex-cimetière voisin du presbytère de San Secondo, des fresques signées par Jacopo Bedi (« \teiHi[rend={italic}]{Iacopus pinxit} »), datées du 7 septembre 1458, illustrent la légende de saint Sébastien :\teiHi[rend={italic}]{ Flagellation, Martyre }et \teiHi[rend={italic}]{Ensevelissement}, évoqués pour conjurer les périls de l’épidémie de peste qui frappa Gubbio précisément cette année-là\teiNote[n={23},place={foot},xmlid={ftn101}]{\teiP[xmlid={p-169}]{Sur Jacopo Bedi on est renseigné de 1432 à 1478. À la voûte de la chapelle Pamphili figurent les effigies des Docteurs de l’Église. Dans deux \teiHi[rend={italic}]{clipei} monochromes sont les portraits de saint Jean et de saint Jean l’évangéliste en train de recevoir la vision de l’Apocalypse. Le même Bedi réalisa au palais des Pamphili une \teiHi[rend={italic}]{Annonciation. }La fresque fragmentaire fut redécouverte en 1920-1922 (Umberto Gnoli, \teiHi[rend={italic}]{Pittori e miniatori nell'Umbria}, Spolète, Argentieri, 1923, p. 131 et 132. Voir aussi Filippo Todini, \teiHi[rend={italic}]{La pittura umbra dal Duecento al primo Cinquecento}, Milan, Longanesi, 1989, p. 31.}}. Le martyre est actualisé : y assistent les habitants de la ville figurée au fond ; le personnage représenté sur la gauche de la scène ornant la paroi centrale pourrait bien être, comme c’est l’usage, le commanditaire (fig. 65). Pour la dévotion de Gubbio, le site de San Secondo offre un intérêt particulier en tant que lieu de formation de saint Ubaldo, évêque et patron vénéré de la ville (1084-1160). Dans le pavement de la chapelle, l’accès à la cavité sépulcrale est clos d’une pierre tombale sur laquelle on remarque l’écu Pamphili aux lys et à la colombe munie du rameau d’olivier, finement encadré d’une bordure au motif du cordon enroulé (fig. 66).}
              \teiP[xmlid={p14-4}]{Les symboles des lys inclus dans la partie supérieure de l’écu Pamphili concordent avec les événements historiques qui, entre le \teiHi[rend={small-caps}]{xiii}\teiHi[rend={sup}]{e} et le \teiHi[rend={small-caps}]{xiv}\teiHi[rend={sup}]{e} siècle, impliquèrent l’Italie centrale et cette famille. En fait, le développement du quartier de San Martino à Gubbio se situe précisément après la compilation des listes de prescription de gibelins en 1315 et avant la crise du milieu du \teiHi[rend={small-caps}]{xiv}\teiHi[rend={sup}]{e }siècle, arc de temps qui voit la construction, le long de l’axe de la rue Dogana (aujourd’hui rue Cavour), des maisons des familles de fidélité guelfe éprouvée, comme les Andreoli, les Beccoli et les Pamphili, qui, par leur activité médicale, avant même celle de judicature, acquirent une position de relief dans la communauté de Gubbio\teiNote[n={24},place={foot},xmlid={ftn102}]{\teiP[xmlid={p-170}]{Paolo Micalizzi, \teiHi[rend={italic}]{Gubbio. Storia dell’architettura e della città} \lbrack{}2\teiHi[rend={sup}]{e} éd.\rbrack{}, Gubbio, L’arte grafica, 2009, p. 185-187. Il s’agit de maisons plutôt que de palais ; leur réalisation, entre la fin du \teiHi[rend={small-caps}]{xiv}\teiHi[rend={sup}]{e} et les premières décennies du siècle suivant, fut effectuée par la fusion progressive de propriétés, suivant une procédure qui se répètera plus tard pour la demeure de la branche romaine de la maison. Sur les activités médicales et juridiques comme sur l’ascension sociale des Pamphili à Gubbio, voir aussi Benedetta Borello, \teiHi[rend={italic}]{L’apprentissage de Rome à la Renaissance, officiers à l’ombre de la Curie }\teiHi[rend={small-caps italic}]{xv}\teiHi[rend={italic}]{-}\teiHi[rend={small-caps italic}]{xvii}\teiHi[rend={italic sup}]{e}\teiHi[rend={italic}]{ siècles,} Turnhout, Brepols, 2021, chap. I en particulier ; Alberto Luongo, \teiHi[rend={italic}]{Gubbio nel Trecento}, Rome, Viella, 2016, p. 293-294 et 555 (sur la situation des Pamphili à Gubbio au début de la seigneurie des Montefeltro). Dirk Amayden, \teiHi[rend={italic}]{op. cit.}, p. 125, affirme que « les Pamphili sont antiques et nobles de première classe dans la région de Gubbio depuis des centaines d’années, antiquité dont témoigne la maison située dans le quartier de San Martino, près du couvent de San Domenico, place de San Lorenzo ».}}. En 1342, deux membres de la famille Pamphili sont nommés « surintendants des corps de garde du quartier San Martino, institués pour défendre la ville des agressions des gibelins », ceux-ci étant par ailleurs réintégrés à l’avènement de la seigneurie des Montefeltro en 1384\teiNote[n={25},place={foot},xmlid={ftn103}]{\teiP[xmlid={p-171}]{Oderigi Lucarelli, \teiHi[rend={italic}]{Memorie e guida storica di Gubbio, }Città di Castello, S. Lappi, 1888, p. 309. L’avènement de la seigneurie des Montefeltro (1384-1632) détermine le développement de Gubbio qui implique une nouvelle définition formelle et planimétrique du tissu urbain. On constate le retour des familles qui avaient été bannies et une sorte de réappropriation gibeline d’une partie du quartier de San Martino. Dans les premières décennies de la seigneurie des Montefeltro on assiste au développement de ce quartier, dans l’espace entourant l’église homonyme (ensuite consacrée à saint Dominique), qui devient le centre de gravité de la zone, contrebalançant les implantations plus anciennes et donnant à Gubbio la forme d’un habitat constitué de deux parties équivalentes, disposées symétriquement (Paolo Micalizzi, \teiHi[rend={italic}]{op. cit.}, p. 182-185).}}.}
              \teiP[xmlid={p15-4}]{Ce que nous observons dans le tiers supérieur de l’écu des Pamphili est un signe de leur appartenance au parti guelfe, dans la mesure où c’est une version de ce qui est appelé « chef d’Anjou », une figuration de tradition française, adoptée par Charles I en 1247, composée précisément de deux pals de gueules entre trois lys se détachant sur un champ d’azur\teiNote[n={26},place={foot},xmlid={ftn104}]{\teiP[xmlid={p-172}]{Le chef est une « pièce honorable » qui occupe le tiers supérieur de l’écu ; il est délimité par une ligne horizontale. Alessandro Savorelli, « Segni e simboli araldici nell’arte fiorentina dal Medioevo al Rinascimento », dans Maria Monica Donato et Daniela Parenti (dir.), \teiHi[rend={italic}]{Dal giglio al David. Arte civica a Firenze fra Medioevo e Rinascimento}, Florence-Milan, Giunti, 2013, p. 73-77, ici p. 74. En août 1380, Charles III d’Anjou entra à Gubbio, ville des Pamphili. }}.}
              \teiP[xmlid={p16-4}]{Après la défaite de Manfred de Souabe à Bénévent en 1266 et de son neveu Conradin, à Tagliacozzo, en 1268, l’influence de Charles I s’étend sur certaines zones du Piémont, sur la Toscane et sur la région de Montefeltro, se consolidant dans le royaume de l’Italie méridionale. Dans le contexte des affrontements entre factions guelfe et gibeline, l’azur devient la couleur du parti guelfe, soutien des Angevins qui concèderont longtemps et largement à leurs partisans l’usage des lys, insérés entre les pals du haut de l’écu, comme symbole de fidélité, noblesse, force, richesse et gloire. Souvent les vassaux adoptèrent le chef de leurs patrons d’Anjou, les lys, le lambel, les pals et les couleurs de France, comme expression d’hommage et requête de protection\teiNote[n={27},place={foot},xmlid={ftn105}]{\teiP[xmlid={p-173}]{Goffredo Di Crollalanza, \teiHi[rend={italic}]{Gli emblemi dei guelfi e dei ghibellini} \lbrack{}San Casciano, F. Cappelli, 1878\rbrack{}, Milan, Orsini De Marzo, 2010, chap. III, p. 57 ; Giacomo Carlo Bascapé et Marcello Dal Piazzo, \teiHi[rend={italic}]{Insegne e simboli. Araldica pubblica e privata medievale e moderna, }Rome, Ministero per i beni culturali e ambientali, 1983 ; Alessandro Savorelli, « “Les fleurs de lys y sont partout”. Origine et signification du chef d’Anjou en Italie (\teiHi[rend={small-caps}]{xiii}\teiHi[rend={sup}]{e}-\teiHi[rend={small-caps}]{xiv}\teiHi[rend={sup}]{e} siècles », \teiHi[rend={italic}]{Armas e troféus}, vol. 20, 2018, p. 93-107.}}. Des villes comme Gubbio, Prato, San Gimignano, Florence et des familles guelfes des territoires voisins témoigneront par ce moyen de leur fidélité au parti de l’Église et à la dynastie française de Naples.}
              \teiP[xmlid={p17-4}]{Des symboles et des figurations angevines envahirent les espaces publics et les palais municipaux, représentant les aspects positifs du pouvoir royal pour les villes : tutelle, adhésion politique, appartenance à un réseau. Souvent, les insignes pontificaux et ceux des Angevins entrèrent dans les armes des institutions locales, comme cela s’observe à Gubbio, ville des Pamphili, sur l’architrave de la porte du palais des Consuls qui porte des dates de la construction (1332-octobre 1336) et inclut une suite de cinq lys séparés par les pendants d’un lambel\teiNote[n={28},place={foot},xmlid={ftn106}]{\teiP[xmlid={p-174}]{Matteo Ferrari, Riccardo Rao et Pierluigi Terenzi, « Rappresentazioni del potere angioino nell’Italia comunale: sovrani, ufficiali, città », dans Thierry Pécout (dir.), \teiHi[rend={italic}]{Les officiers et la chose publique dans les territoires angevins (}\teiHi[rend={small-caps italic}]{xiii}\teiHi[rend={italic sup}]{e}\teiHi[rend={italic}]{-}\teiHi[rend={small-caps italic}]{xvi}\teiHi[rend={italic sup}]{e}\teiHi[rend={italic}]{ siècle) : vers une culture politique }, Rome, École française de Rome, « École française de Rome, 518, 4 », 2020. Voir aussi Francesca Fumi Cambi Gado, « Stemmi ed emblemi nella decorazione degli edifici\teiHi[rend={italic}]{ }», dans Amerigo Restucci (dir.),\teiHi[rend={italic}]{ L’architettura civile in Toscana. Il Medioevo, }Sienne, Monte dei Paschi di Siena, 1995, p. 401-441, ici p. 423.}} (fig. 67).}
              \teiP[xmlid={p18-4}]{Bien que la présence des Anjou fût assez contestée, le lys fut conservé dans les répertoires iconographiques locaux, même après leur départ d’Italie\teiNote[n={29},place={foot},xmlid={ftn107}]{\teiP[xmlid={p-175}]{Vieri Favini et Alessandro Savorelli, \teiHi[rend={italic}]{Segni di Toscana. Identità e territorio attraverso l’araldica dei comuni: storia e invenzione grafica (secoli XIII-XVII),} Florence, Le Lettere, 2006 ; Alessandro Savorelli, « L’araldica per la storia: una fonte ausiliaria ? », dans Maria Pia Paoli (dir.), \teiHi[rend={italic}]{Nel laboratorio della storia: una guida alle fonti dell'età moderna}, Rome, Carocci, 2013, p. 289-315.}}. On peut ainsi comprendre comment la maison des Pamphili, lorsqu’elle définit son propre écu nobiliaire, inséra l’allusion à un moment clef de son histoire, évoquant sa tenace fidélité à l’Église, exprimée dans l’appartenance au parti guelfe\teiNote[n={30},place={foot},xmlid={ftn108}]{\teiP[xmlid={p-176}]{Giacomo Carlo Bascapé et Marcello Dal Piazzo, \teiHi[rend={italic}]{op. cit.}, 1983, p. 41, 57, 76 et 86.}}.}
              \teiP[xmlid={p19-4}]{Pour ce qui regarde la branche romaine de la famille, le groupe des trois lys fut inclus dans l’écu pour évoquer le passé ombrien, référence à une antiquité contribuant au prestige de la maison, engagée à Rome dans la consolidation de sa dignité curiale, pourvoyeuse de droits et de prérogatives toujours plus exclusifs.}
              \teiP[xmlid={p20-4}]{Comme on l’a dit, en 1460, avec le transfert d’Antonio Pamphili, commença le destin de la branche romaine de la maison qui se détacha progressivement de l’Ombrie\teiNote[n={31},place={foot},xmlid={ftn109}]{\teiP[xmlid={p-177}]{Un autre moment de la séparation de la branche romaine des Pamphili d’avec Gubbio, est attesté en 1501, quand les trois héritiers d’Angelo Benedetto Pamphili vendent quelques propriétés. Benedetta Borello, , « Alleanze matrimoniali… », art. cité, p. 362-363 ; Benedetta Borello, \teiHi[rend={italic}]{op. cit}., 2021, p. 124-128 et 131-134.}} : toutefois, encore en 1663, son lien avec Gubbio est mis en évidence dans la vue axonométrique de la \teiHi[rend={italic}]{Cité royale très antique }\lbrack{}\teiHi[rend={italic}]{Città Regia Antichiss}\lbrack{}\teiHi[rend={italic}]{ima}\rbrack{}\teiHi[rend={italic}]{,} dessinée par le citoyen de Gubbio, Ignazio Cassetta, et éditée à Amsterdam par Johannes Blaeu, avec dédicace à Vincenzo Armanni\teiNote[n={32},place={foot},xmlid={ftn110}]{\teiP[xmlid={p-178}]{\teiHi[rend={italic}]{Theatrum Civitatum et admirandorum Italiae}, vol. 1, \teiHi[rend={italic}]{Stato della Chiesa}, Amsterdam, s. n., 1663, 430 × 575 mm. Dans cette édition en haut à droite, se lit la dédicace de la carte de la part du « Prieur Ignatio Cassetta » (1592-1667), à Vincenzo Armanni. Ce dernier avait accompagné d’une description historique de Gubbio en latin, utile pour élaborer la légende, l’envoi de la carte à Rome à Carlo Cartari, avocat concistorial, qui, à son tour, l’envoya à Blaeu. Dans une autre légende, placée dans la partie inférieure gauche, sont contenues les indications relatives aux lieux et institutions « civils » de la ville ; le palais des Pamphili est indiqué par la lettre « Z ». La vue fut rééditée avec quelques variantes et sans la dédicace à Armanni, par Pierre Mortier dans \teiHi[rend={italic}]{Nouveau théâtre de l’Italie,} publié à Amsterdam en 1704-1705 (vol. II) et réédité par Rutger Christoffel Alberts en 1724. Giuseppe Maria Nardelli, \teiHi[rend={italic}]{La pianta di Gubbio di Ignazio Cassetta stampata da Joannis Blaeu e Pierre Mortier: rappresentazione e simbolismi}, Gubbio, Rotary Club, 2001. Nardelli étudie en détail les éléments naturels et artificiels présents dans la représentation. }}. De fait, aux limites nord-ouest de l’habitat, émerge l’imposant édifice du « Palais des Pamphili », unique demeure privée indiquée par une légende dans la carte (fig. 68).}
              \teiP[xmlid={p21-4}]{Les années 1630, avec l’élévation au cardinalat de Giovan Battista Pamphili, furent un moment clef pour la fortune romaine de l’écu Pamphili. Nous avons déjà évoqué le monument commémoratif en l’honneur de Girolamo Pamphili à l’église de la Vallicella, réalisée en 1638. La même année, l’écu à la colombe et aux lys, indiqué comme « à Rome, des Pamphili \lbrack{}\teiHi[rend={italic}]{Romae Pamphiliorum}\rbrack{} », est publié dans les \teiHi[rend={italic}]{Armoiries nobiliaires} du jésuite Silvestro Pietrasanta, qui ne tardera pas à devenir un influent collaborateur du pape Pamphili\teiNote[n={33},place={foot},xmlid={ftn111}]{\teiP[xmlid={p-179}]{Silvestro Pietrasanta, \teiHi[rend={italic}]{Tesserae gentilitiae,} Romae, Typis haeredum Francisci Corbelletti, 1638, p. 438. Pietrasanta décrit ainsi l’écu des Pamphili : « \teiHi[rend={italic}]{Argentea rursus in scuti valvulo muricato, ac rostro attinens oleae ramulum prasinum, cum quadrifida punicea laciniola superne, inumbrante tria aurata lilia est Romae Pamphiliorum} ». Marc Vulson de La Colombière accusa Pietrasanta d’avoir copié dans les \teiHi[rend={italic}]{Tesserae} la méthode qu’il avait élaborée pour rendre graphiquement les images héraldiques. Le 8 août 1642, l’Inquisition emprisonna le fondateur de l’ordre piariste, Giuseppe Calasanzio (1557-1648). Le 9 mai 1643, Pietrasanta fut nommé visiteur apostolique de cet ordre qui depuis 1612 avait son siège avec son école à San Pantaleo, près du palais Pamphili de la place Navone. À la suite des critiques de Pietrasanta, Innocent X, par un bref du 16 mars 1646, réduisit l’ordre à une congrégation séculière subordonnée aux évêques locaux. L’ordre fut réhabilité en 1656 et rénové en 1669. }} (fig. 69).}
              \teiP[xmlid={p22-4}]{Particulièrement significatives pour la fortune romaine de l’écu furent les circonstances de la réalisation du palais de la famille. Le problème qu’Antonio Pamphili avait affronté à son arrivée à Rome en 1460 était de réussir à s’intégrer à la classe dirigeante de la ville, objectif poursuivi grâce à une habile politique de mariages avec des membres de grandes familles romaines\teiNote[n={34},place={foot},xmlid={ftn112}]{\teiP[xmlid={p-180}]{Benedetta Borello, « Alleanze matrimoniali... », art. cité, p. 345-366 ; Christof Friedrich Weber, « Papstgeschichte und Genealogie », \teiHi[rend={italic}]{Römische Quartalschrift}, n\teiHi[rend={sup}]{o} 84, 1989, p. 331-400, en particulier p. 394 et 395. }} et à un choix judicieux du lieu de résidence\teiNote[n={35},place={foot},xmlid={ftn113}]{\teiP[xmlid={p-181}]{Une communauté originaire de Gubbio est présente dans la zone de Parione, mais elle n’est pas assez structurée pour garantir à Antonio Pamphili un rôle social correspondant au rang acquis : Benedetta Borello, « Alleanze matrimoniali… », art. cité, p. 345-366 ; Christof Weber, art. cité, en particulier p. 394 et 395. }}.}
              \teiP[xmlid={p23-4}]{Entre 1471 et 1479, il acquiert quelques propriétés place Pasquino et rue de l’Anima, dans le quartier de Parione, une zone en expansion, stratégiquement située sur la voie papale, la \teiHi[rend={italic}]{via Papalis}, parcours solennel pendant la cérémonie de l’intronisation du Pontife entre le Vatican et Saint-Jean-de-Latran\teiNote[n={36},place={foot},xmlid={ftn114}]{\teiP[xmlid={p-182}]{Luigi De Gregori, « Piazza Navona prima d’Innocenzo X », \teiHi[rend={italic}]{Roma}, n\teiHi[rend={sup}]{o} 4, 1926, p. 14-25 et 97-116, ici p. 111 et 112 ; Mirka Beneš, dans \teiHi[rend={italic}]{The Villa Pamphilj (1630,1670): Family, Gardens, and Land in Papal Rome}, thèse de doctorat, Université de Yale, 1989, p. 33, 36, 37, 59 et 60, publie les documents d’archive relatifs aux propriétés acquises par Antonio Pamphili à différentes époques. Anna Esposito, « Il rione Parione durante il Pontificato Sistino: analisi di un'area campione. Osservazione sulla popolazione rionale », dans Massimo Miglio, Francesca Niutta, Diego Quaglioni \teiHi[rend={italic}]{et al}. (dir.), \teiHi[rend={italic}]{Un pontificato e una città. Sisto IV (1471-1484)}, actes de colloque, Rome, 3-7 décembre 1984, Cité du Vatican, Scuola Vaticana di Paleografia, Diplomatica e Archivistica, « Littera antiqua, 5 », 1986, p. 651-666.}}. Finalement, en 1497, le fils d’Antonio, Angelo Benedetto Pamphili, est enregistré dans les documents comme citoyen romain et scripteur des lettres apostoliques \lbrack{}\teiHi[rend={italic}]{cives romanus ac litterarum apostolicarum scriptor}\teiNote[n={37},place={foot},xmlid={ftn115}]{\teiP[xmlid={p-183}]{Benedetta Borello, « Strategie di insediamento in città: i Pamphilj a Roma nel primo Cinquecento », dans Maria Antonietta Visceglia (dir.), \teiHi[rend={italic}]{La nobiltà romana in età moderna}, Rome, Carocci, 2001, p. 31-61, en particulier p. 39 et 54. Les conditions sociales de l’implantation d’Antonio Pamphili et de ses descendants à Rome sont approfondies par Benedetta Borello, \teiHi[rend={italic}]{op. cit.}, en particulier au chap. II, p. 69-121 et chap. IV, p. 182-184.}}\rbrack{}\teiHi[rend={italic}]{. }Les propriétés de ce quartier de Parione s’étendirent progressivement grâce à l’acquisition d’édifices adjacents au noyau initial : en 1615, les documents enregistrent un projet pour intégrer les parties les plus anciennes dans d’autres donnant sur la place Navone, enregistrées à partir de 1592\teiNote[n={38},place={foot},xmlid={ftn116}]{\teiP[xmlid={p-184}]{Stéphanie Leone, \teiHi[rend={italic}]{The Palazzo Pamphilj in Piazza Navona: Constructing Identity in Early Modern Rome}, Londres, H. Miller, 2008. Voir aussi « Cardinal Pamphilj Builds a Palace: Self-Representation and Familial Ambition in Seventeenth-Century Rome », \teiHi[rend={italic}]{Journal of the Society of Architectural Historians}, vol. 63, n\teiHi[rend={sup}]{o}4, 2004, p. 440-471 (la référence à la décoration des arches de la cour se trouve à la p. 450) ; « La costruzione di palazzo Pamphilj », dans \teiHi[rend={italic}]{Palazzo Pamphilj: ambasciata del Brasile a Roma,} Turin, U. Allemandi, 2016, p. 15-67. Stéphanie Leone (\teiHi[rend={italic}]{op. cit}., p. 88, fig. 39) observe, sur la base d’un dessin de Stefano Pignatelli de 1615 (ADP 86.2, int.2), que, en réalité, les travées du palais vers la place de Pasquin étaient au nombre de cinq et non de sept comme on le voit dans la gravure de Giovanni Maggi de 1625. Du vestibule donnant sur la place, on accède à un monumental escalier de cérémonie qui conduit aux appartements de représentation du cardinal.}}. En 1634, Giovan Battista Pamphili acquiert d’autres biens immobiliers de la famille Teofili, doublant l’ampleur de sa propriété ; il demande ensuite à l’architecte Francesco Peperelli (vers 1585-1641) d’unir en une structure unique les corps de fabrique existants et d’installer une demeure correspondant à sa dignité de cardinal\teiNote[n={39},place={foot},xmlid={ftn117}]{\teiP[xmlid={p-185}]{Stéphanie Leone, \teiHi[rend={italic}]{op. cit.}, p. 111-126 ; Mirka Beneš, \teiHi[rend={italic}]{op. cit.}, p. 75 et 76, confirme que dans cette phase du premier \teiHi[rend={small-caps}]{xvii}\teiHi[rend={sup}]{e} siècle, les structures préexistantes furent réutilisées.}}.}
              \teiP[xmlid={p24-3}]{L’aménagement du palais prévoit la création d’appartements destinés à Giovan Battista Pamphili, à son frère et à sa femme Olimpia Maidalchini. Dans le corps sud-ouest, vers la place Pasquino, qui pousse vers les parties les plus anciennement acquises par les Pamphili, est inséré un escalier à plan carré autour d’une cage centrale. Dans l’aménagement voulu par Peperelli, il permet l’accès aux appartements de Donna Olimpia, et, aux étages supérieurs, aux membres de sa suite. Sur le portail d’entrée de l’appartement de représentation\teiNote[n={40},place={foot},xmlid={ftn118}]{\teiP[xmlid={p-186}]{Stéphanie Leone, \teiHi[rend={italic}]{op. cit}., p. 120-126.}}, on observe un écu de pierre à la colombe et aux lys disposé sur une plaque ornée de rubans déployés, de facture plus ancienne que ceux qui sont visibles ailleurs dans le palais (fig. 70).}
              \teiP[xmlid={p25-3}]{Dans la configuration qui suivit, les appartements s’articulent autour de deux cours. Deux vestibules en portiques, accessibles, respectivement, depuis la rue de l’Anima et la place Navone, servent d’entrée pour la cour principale. Dans la \teiHi[rend={italic}]{Mesure et estimation des travaux de stuc faits au palais situé place Navone par l’Emminentissime Seigneur Cardinal Pamphili} \lbrack{}\teiHi[rend={italic}]{Misura e stima de’ lavori di stucco fatti nel palazzo posto in Navona dell’…Em.m Sig. Card.le Panfili}\rbrack{}, datée du 9 juin 1636, le paragraphe relatif à la « façade de la loggia vers la cour » \lbrack{}\teiHi[rend={italic}]{facciata della loggia verso il cortile}\rbrack{} enregistre un paiement « pour avoir dessiné et stuqué les six palombes dans les triangles au-dessus des trois arcs avec leurs rameaux d’olivier dans le bec\teiNote[n={41},place={foot},xmlid={ftn119}]{\teiP[xmlid={p-187}]{« \teiHi[rend={italic}]{per aver abozzato e stuccato le 6 palombe nelli triangoli sop.a li 3 archi con suo rame d’oliva in bocca }»,\teiHi[rend={italic}]{ Misura e stima de’lavori di stucco fatti di tutta robba eccetto la calce bianca da M.ro Ludovico Bossi Capomastro Muratore nel palazzo posto in Navona dell’…Em.m Sig. Card.le Panfili}, ADP, 88.34.2, fol. 215, transcrit partiellement dans Stéphanie Leone, \teiHi[rend={italic}]{op. cit.,} Appendix, doc. 3, p. 319-321, ici p. 320.}} » (fig. 71 et 72). Dans le même document, un autre paiement indique que dans la « loggia de la cour, à l’étage noble » \lbrack{}\teiHi[rend={italic}]{loggia del cortile, al piano nobile}\rbrack{}, sont insérées, dans les volutes de la console, « deux nacelles et en leur milieu \lbrack{}...\rbrack{} une palombe en bas-relief \lbrack{}...\rbrack{} avec son rameau d’olivier au bec\teiNote[n={42},place={foot},xmlid={ftn120}]{\teiP[xmlid={p-188}]{\teiHi[rend={italic}]{Ibid.}, fol. 212, p. 320 : « \teiHi[rend={italic}]{doi baccelli... et nel suo mezzo fatto una palomba dì mezzo rilievo…con suo ramo d’olive in bocca }». À l’est, la loggia était ouverte, alors que les arches vers l’ouest étaient aveugles. Stéphanie Leone, \teiHi[rend={italic}]{op. cit}., p. 112. La loggia manque sur le côté sud, qui coïncide avec les parties les plus anciennes du palais. À l’étage noble, la nouvelle loggia fut de dimentions plus modestes que la précédente : elle incluait quatre colonnes de « \teiHi[rend={italic}]{bigio} » et deux de « \teiHi[rend={italic}]{cipollino} ». Ces colonnes furent réutilisées, après que leur hauteur eut été réglée, adossées à des pilastres. Elles sont perdues comme les décorations de stuc à la base des arches de la voûte, mentionnées dans le document (Stéphanie Leone, \teiHi[rend={italic}]{op. cit}., p. 114). Des colombes furent aussi incluses dans la décoration de la tour existant aux limites du palais Teofili, comme il résulte du paiement pour « avoir fait le relief du corps des douze palombes dessinées et stuquées avec des rameaux d’olivier au bec dans lesdits triangles sur les arches » \lbrack{}\teiHi[rend={italic}]{Aver fatto l’aggetto del corpo delle n. 12 palombe ne’detti triangoli sopra l’archi abozz.e e stucc.te con rame d’olivo in becco}...\rbrack{}, ADP, 88.34.2, fol. 207, dans Stéphanie Leone, \teiHi[rend={italic}]{op. cit}., p. 320. Dans un autre document, présent dans le même dossier de l’archive Pamphili, daté du 21 mars 1638, fol. 226-325, ici fol. 247 (Stéphanie Leone, \teiHi[rend={italic}]{op. cit.}, chap. \teiHi[rend={italic}]{Lavori fatti nella facciata di d.o Cortile dove è l’Entrone verso la Strada dell’Anima}, p. 321-325, ici p. 323) nous avons l’indication que la loggia du rez-de-chaussée est limitée par « deux pilastres isolés ». Vers 1653, un autre document (ADP, 88.34.1 ; Stéphanie Leone, \teiHi[rend={italic}]{op. cit}., Appendix, doc. 4, p. 325-350), inclus dans le chapitre relatif aux « travaux de la façade extérieure sud. Rehaussement vers la place » \lbrack{}\teiHi[rend={italic}]{lavori nella facciata di fuori del sud. Rialzamento verso la piazza}\rbrack{}, relève l’insertion du motif décoratif des lys (fol. 18) et de la « devise de palombes sur les reliefs des susdits termes faits, la coquille avec des canaux et une cascade de festons au dessous » \lbrack{}\teiHi[rend={italic}]{impresa di Palombe dentro sotto sopra li aggetto delli suddetti termini fattici la conchiglia con canali e cascata di festoni sotto}\rbrack{} (fol. 21, Stéphanie Leone, \teiHi[rend={italic}]{op. cit}., p. 326). L’emblème de la colombe ne semble pas figurer dans les peintures des salles décorées durant le cardinalat de Giovan Battista Pamphili, en particulier dans les\teiHi[rend={italic}]{ Histoires de Joseph }(1630-1632), les \teiHi[rend={italic}]{Histoires de Moïse} et dans les \teiHi[rend={italic}]{Marines }d’Agostino Tassi (1635) ; Patrizia Cavazzini, « Towards a Chronology of Agostino Tassi », \teiHi[rend={italic}]{The Burlington Magazine}, n\teiHi[rend={sup}]{o} 144, 2002, p. 396-408\teiHi[rend={italic}]{. }Vers 1647, sont documentées les recherches promues par donna Olimpia Maidalchini sur l’héraldique de sa famille paternelle ; voir aussi Tommaso Tomasi, \teiHi[rend={italic}]{Della esaltazione di papa Innocentio X. Lettera panegirica}, In Roma, Nella stamperia di Bernardino Tani, 1644.}} ». L’uniformité qui s’observe dans la forme des colombes sur les façades de la cour centrale et de celle du nord, réalisée après 1645, indique que ces décorations de stuc ont été modifiées après la restructuration documentée de 1636-1638, ou qu’elles ont été exécutées avec une attention particulière pour obtenir un effet homogène\teiNote[n={43},place={foot},xmlid={ftn121}]{\teiP[xmlid={p-189}]{Vers 1653, le document ADP, 88.34.1 (Stéphanie Leone, \teiHi[rend={italic}]{op. cit.}, Appendix, doc. 4, p. 325-350), fol. 79, au chapitre « Façade de la galerie dans ladite loggia » \lbrack{}\teiHi[rend={italic}]{Facciata della galleria in detta loggia}\rbrack{} (Stéphanie Leone, \teiHi[rend={italic}]{op. cit}., p. 333), il y a encore un paiement pour le relevé des arcades de la loggia au premier étage de la cour : « pour le relief stuqué des bandes qui encadrent les triangles , pour six palombes dans lesdits triangles avec leurs rameaux d’olivier et l’enduit de stuc sur le fond de chacun » \lbrack{}\teiHi[rend={italic}]{per l’aggetto stuccato delle fascie, che fanno riquadramento nelli triangoli; per 6 Palombe in d.ti triangoli, con suoi rami d’Olivo et incollato di stucco il campo di ciascuno}\rbrack{}.}}.}
              \teiP[xmlid={p26-3}]{Après l’élection de Giovan Battista Pamphili au pontificat, en 1644, et avant 1650, date du jubilé, le palais subit des modifications substantielles, confiées aux architectes Girolamo Rainaldi et Francesco Borromini, assumant un aspect encore visible aujourd’hui, où l’emblème de la colombe au rameau d’olivier dans le bec est largement présent\teiNote[n={44},place={foot},xmlid={ftn122}]{\teiP[xmlid={p-190}]{Girolamo Rainaldi, impliqué dans le réaménagement du palais à partir de 1644, créa une troisième cour, reliée par un portique à celle du centre. Les projets de Rainaldi pour l’agrandissement de l’édifice figurent parmi les papiers de Virgilio Spada dans le manuscrit de la BAV, Vat. lat. 11258, fol. 170, 144 et 167, qui contient aussi le projet de Borromini pour la distribution intérieure (fol. 168, 169, 181 et 192). }}.}
              \teiP[xmlid={p27-3}]{L’image symbolique s’enrichit d’autres références : c’est un phénomène récurrent dans le contexte de l’ascension des familles aristocratiques à l’époque moderne que la récupération de références antiques se radicalise au point de devenir légende. Le recours à la figure de l’Énée virgilien caractérise le récit des origines aussi bien chez les Barberini que chez les Pamphili, en lien avec la conception de Rome comme siège mythique d’une autorité politique et spirituelle universelle, diffusée au moins depuis le \teiHi[rend={small-caps}]{xiii}\teiHi[rend={sup}]{e} siècle\teiNote[n={45},place={foot},xmlid={ftn123}]{\teiP[xmlid={p-191}]{Maffeo Barberini, avant d’être élu pape, s’était forgé la devise « \teiHi[rend={italic}]{hic domus} » \lbrack{}ici est ma demeure\rbrack{}, citation des mots que le héros antique aurait prononcés à son arrivée dans le Latium, auspice de la fondation de Rome et d’un futur âge d’or.}}.}
              \teiP[xmlid={p28-3}]{Dans la décoration du palais de la place Navone, dans la \teiHi[rend={italic}]{Salle de Bacchus}, qui prend son nom de la frise peinte par Andrea Camassei en 1646-1648, aux embrasures des fenêtres sont figurées Vénus et Junon. Un panégyrique en l’honneur d’Innocent X indique que la colombe et le lys sont leurs attributs. L’assimilation de la colombe pamphili à Vénus, mère d’Énée, considérée comme fondatrice de la maison, fait partie de l’invention tardive d’un passé aussi glorieux qu’imaginaire, promu par les Pamphili\teiNote[n={46},place={foot},xmlid={ftn124}]{\teiP[xmlid={p-192}]{Giuseppe Simonetta, Laura Gigli et Gabriella Marchetti, \teiHi[rend={italic}]{Sant’Angese in Agone a Piazza Navona. Bellezza, Proporzione, Armonia nelle fabbriche Pamphili}, Rome, Gangemi, 2003, p. 55 et 56. Vénus apparaît aussi dans la frise d’Andrea Camassei. Au cours de la démolition de l’ancienne église de Sainte-Agnès, le 17 avril 1652, fut retrouvée une pierre portant une épigraphe avec le nom : « A. Volturgio Pamphilio », qui, pour le jésuite Girolamo Petrucci (1585-1669), collaborateur d’Athanasius Kircher, constituait une nouvelle preuve de l’antiquité de la maison Pamphili. Les affirmations de Petrucci furent contredites sur des bases philologiques par Denis Lambin et par André Schott, qui, à leur tour, furent durement critiqués par le jésuite (\teiHi[rend={italic}]{ibid.}, p. 77-80).}}.}
              \teiP[xmlid={p29-3}]{La référence à Énée triomphe dans les \teiHi[rend={italic}]{Histoires }de la galerie du palais, œuvre de Pietro da Cortona (1651-1654), connaisseur averti de l’\teiHi[rend={italic}]{Énéide}\teiNote[n={47},place={foot},xmlid={ftn125}]{\teiP[xmlid={p-193}]{Rudolf Preimesberger, « “Pontifex Romanus per Aeneam Praesignatus”, Die Galleria Pamphilj und ihre Fresken », \teiHi[rend={italic}]{Römisches Jahrbuch für Kunstgeschichte}, n\teiHi[rend={sup}]{o} 16, 1976, p. 221-287. }}, comme chez Angelo Caferri déjà cité, auteur du programme iconographique des peintures qui montrent un flagrant exemple de ce que l’on a appelé « généalogies incroyables\teiNote[n={48},place={foot},xmlid={ftn126}]{\teiP[xmlid={p-194}]{Roberto Bizzocchi, \teiHi[rend={italic}]{Genealogie incredibili. Scritti di storia nell’Europa moderna}, Bologne, Il Mulino, 1995.}} ». Il faut rappeler que Caferri affirme dans un \teiHi[rend={italic}]{Discours }que Numa Pompilius, qu’il appelle « Pamphilius », serait un descendant de la « famille Pamphilia de Sparte dans cette cité fondée par Pamphilio, roi grec des Doriens, 350 ans avant Rome\teiNote[n={49},place={foot},xmlid={ftn127}]{\teiP[xmlid={p-195}]{« \teiHi[rend={italic}]{famiglia Panfilia di Sparta, in quella città fondata da Panfilio re greco de’Dorici, 350 anni prima della fondazione di Roma }», Gaetano Moroni, \teiHi[rend={italic}]{op. cit.}, 1851, vol. LI, entrée « Pamphili », p. 85, rapporte cette hypothèse, avancée par Nicolò Angelo Caferri, et affirme qu’elle fut dictée par « l’adulation ». L’affirmation de Caferri sur l’origine spartiate des Pamphili et sur leur descendance de Numa Pompilius est déjà avancée par Girolamo Brusoni dans son ode intitulée \teiHi[rend={italic}]{Le glorie Pamphilie, }publiée en 1656, et, en 1662, par Camillo Pamphili dans ses compositions poétiques incluses dans le recueil dirigé par Girolamo Brusoni sous le titre\teiHi[rend={italic}]{ Degli Allori d’Eurota, Poesie diverse dell’Eccellentiss. Sig. Principe D. Camillo Pamphilo, vol. 2: Raccolte dal Cavalier Girolamo Brusoni, e Dedicate all'Eccellentissima Signora Principessa Donna Olimpia Aldobrandina Pamphilii}, parte prima, In Venetia, Per il Valuasense, 1662. Ces hypothèses suscitèrent aussi les critiques de Denis Lambin et d’André Schott (Giuseppe Simonetta, Laura Gigli et Gabriella Marchetti, \teiHi[rend={italic}]{Sant’Angese in Agone a Piazza Navona. Bellezza, Proporzione, Armonia nelle fabbriche Pamphili}, \teiHi[rend={italic}]{op. cit}, p. 79).}} ». L’affirmation de l’origine spartiate des Pamphili et de leur descendance de Numa Pompilius sont avancées et répétées dans plusieurs écrits par Girolamo Brusoni (1614-1686), érudit et historiographe, originaire du Polesine, chartreux ondoyant, qui, après avoir renié son adhésion aux idées libertines, adopta un radical conformisme, témoignant son respect à l’Église et aux puissants de l’époque\teiNote[n={50},place={foot},xmlid={ftn128}]{\teiP[xmlid={p-196}]{Gaspare De Caro, « Brusoni, Girolamo »,\teiHi[rend={italic}]{ DBI}, \teiHi[rend={italic}]{op. cit}., vol. 14, 1972, p. 712-720.}}.}
            
\end{teiDiv}
            \begin{teiDiv}[type={section1},xmlid={div03-4}]

              \teiHead[xmlid={lecu-des-pamphili-durant-le-pontificat-dinnocent-x-1644-1655-deux-exemples-en-architecture-001}]{L’écu des Pamphili durant le pontificat d’Innocent X (1644-1655) : deux exemples en architecture}
              \teiP[xmlid={p30-3}]{Durant le pontificat d’Innocent X (1644-1655), l’écu, centré sur la colombe au rameau d’olivier, se diffuse dans les œuvres architecturales commanditées par le pape et par les membres de sa famille, en particulier son neveu Camillo et sa belle-sœur Olimpia Maidalchini, aussi bien à Rome que dans les fiefs de San Martino al Cimino et de Valmontone.}
              \teiP[xmlid={p31-3}]{Un exemple, choisi parmi les œuvres promues par le pontife à Rome, permettra d’éclairer les objectifs et les modalités d’usage de l’emblème dans cette période. Les significations qui lui sont attribuées deviennent parfaitement adéquates à la situation historico-politique et démontrent comment deux caractéristiques apparemment contradictoires peuvent être attribuées à ces images symboliques sur la durée de plusieurs siècles : efficacité communicative et malléabilité.}
              \teiP[xmlid={p32-3}]{Pour le premier aspect, on peut relever que l’élaboration d’un écu prévoit souvent l’adoption de symboles disponibles pour plusieurs maisons. Le meuble de la colombe (sans le rameau d’olivier toutefois) est présent dès 1470 dans la chapelle Palombara, dédiée à saint François dans l’église romaine de San Silvestro in Capite\teiNote[n={51},place={foot},xmlid={ftn129}]{\teiP[xmlid={p-197}]{Dirk Amayden, \teiHi[rend={italic}]{op. cit}., p. 123-124. Une colombe « au naturel » figure dans l’écu de la famille Rosolini (\teiHi[rend={italic}]{ibid.}, p. 171). Pietrasanta aussi (\teiHi[rend={italic}]{op. cit.}, p. 438) illustre les écus de plusieurs familles où sont incluses des colombes.}}. L’écu des Ciccolini de Macerata est « d’azur au mont à six coupeaux d’or, surmonté d’une colombe d’argent tenant en son bec un épi d’or\teiNote[n={52},place={foot},xmlid={ftn130}]{\teiP[xmlid={p-198}]{\teiHi[rend={italic}]{Ibid}., p. 359 : « \teiHi[rend={italic}]{d’azzurro al monte di sei cime d’oro, cimato da una colomba d’argento tenente nel becco una spiga di grano d’oro} ».}} ». L’écu du Pérugin Vermigliolo Vermiglioli, conservateur à Rome en 1634, présente, dans sa partie inférieure en champ de gueules « une colombe d’argent béquetant un rameau d’olivier sur lequel elle est posée\teiNote[n={53},place={foot},xmlid={ftn131}]{\teiP[xmlid={p-199}]{\teiHi[rend={italic}]{Ibid}., p. 227 : « \teiHi[rend={italic}]{una colomba d’argento in atto di beccare un ramo d’ulivo sul quale è posata }». La tombe de Giovanni Maria Vermiglioli, au cloître de Santa Maria della Pace, est datée de 1644.}} ».}
              \teiP[xmlid={p33-3}]{Pour ce qui regarde l’adéquation de l’emblème aux objectifs politiques et idéologiques du pontificat d’Innocent X, une lettre de Fabio Chigi à Virgilio Malvezzi, politologue philo-espagnol et son ami proche, offre une clef de lecture précise. Le cardinal Chigi écrit : « L’État ecclésiastique, l’Italie, le monde entier espèrent beaucoup d’être sauvés des maux présents et que la colombe portant l’olive soit un signe certain que nous avons passé le déluge\teiNote[n={54},place={foot},xmlid={ftn132}]{\teiP[xmlid={p-200}]{« \teiHi[rend={italic}]{Lo Stato Ecclesiastico, l’Italia, il mondo tutto spera il gran ristoro de correnti mali, e che la colomba, e l’oliva che porta, siano certo segno del passato diluvio} », Münster, le 5 novembre 1644, de Chigi à Malvezzi, BAV, ms. Chigiani, A.II.29, fol. 45, transcrit dans Vlastimil Kybal et Giovanni Incisa della Rocchetta (éd.), \teiHi[rend={italic}]{La nunziatura di Fabio Chigi}, Rome, Tip. Ist. Grafico Tiberino, 1943, vol. 1, p. 536. Je remercie Julien Régibeau pour m’avoir signalé cette lettre.}}. » La lettre est datée du 5 novembre 1644 ; Giovan Battista Pamphili avait été élu pape le 15 septembre. Sur son initiative, le 23 décembre 1643, Fabio Chigi avait été envoyé comme nonce extraordinaire au congrès de Münster pour négocier la paix qui devait mettre fin à la guerre de Trente Ans.}
              \teiP[xmlid={p34-2}]{Dans le domaine religieux, depuis les premiers temps chrétiens, la colombe, animal pur et doux, était le symbole de la Sagesse divine exprimée par le Saint-Esprit\teiNote[n={55},place={foot},xmlid={ftn133}]{\teiP[xmlid={p-201}]{Dans le \teiHi[rend={italic}]{Cantique des Cantiques}, 2:14, « Ma colombe » est un qualificatif affectueux réservé à la Sulamite par le pasteur amoureux. Dans la langue hébraïque, Giona (Yohnàh, « colombe ») est un nom masculin commun. L’identification de la colombe à l’Esprit-Saint est dans Tertullien, \teiHi[rend={italic}]{De Baptismo}, chap. VIII. Cette identification est illustrée par la médaille d’Innocent X à la colombe dans une couronne d’olivier, \teiHi[rend={italic}]{Replevit orbem terrarum}. Voir Yvan Loskoutoff, « Hercule et Atlas à Rome puis à Modène, Un présent diplomatique du cardinal Francesco degli Albizzi au cardinal Mazarin », \teiHi[rend={italic}]{Dix-septième siècle,} n\teiHi[rend={sup}]{o} 241, octobre-décembre 2008, p. 677-708.}}. À partir du \teiHi[rend={small-caps}]{iv}\teiHi[rend={sup}]{e} siècle, et dans les miniatures des manuscrits des \teiHi[rend={small-caps}]{v}\teiHi[rend={sup}]{e} et \teiHi[rend={small-caps}]{vi}\teiHi[rend={sup}]{e} siècles, dans les représentations de l’Annonciation à Marie et de la Trinité, la colombe devient l’icône du Saint-Esprit. Elle est aussi un emblème de renaissance et de salut. Dans la Genèse, une colombe porte à Noé le rameau d’olivier qui annonce la fin du déluge universel et le début d’une nouvelle ère de paix entre Dieu et les hommes\teiNote[n={56},place={foot},xmlid={ftn134}]{\teiP[xmlid={p-202}]{\teiHi[rend={italic}]{Genèse}, 8:11.}}. De manière analogue, Pâques, résurrection du Christ, a comme symbole une colombe blanche qui tient un rameau d’olivier dans le bec. L’association entre colombe et renaissance est aussi propre au sacrement du baptême : dans l’évangile de Matthieu il est dit qu’une colombe descendant du ciel apparut à Jean-Baptiste lors du baptême du Christ\teiNote[n={57},place={foot},xmlid={ftn135}]{\teiP[xmlid={p-203}]{Matthieu, \teiHi[rend={italic}]{Évangile}, 3:16.}}.}
              \teiP[xmlid={p35-2}]{Les références à saint Jean-Baptiste conviennent à la figure du pape Pamphili, dont le nom était justement Jean-Baptiste, de même que celles à la paix et à la renaissance l’étaient à la situation européenne. Les récits des événements qui précédèrent l’élection d’Innocent X sont parsemés d’apparitions de colombes au rameau d’olivier, de même que les fêtes pour l’intronisation du pontife, la \teiHi[rend={italic}]{possessio. }À cette occasion, place Navone, fut construite une machine éphémère représentant Noé qui accueille dans ses bras une colombe en vol ; elle couronnait un parcours allégorique complexe imaginé par Girolamo et Carlo Rainaldi\teiNote[n={58},place={foot},xmlid={ftn136}]{\teiP[xmlid={p-204}]{Francesco Cancellieri, \teiHi[rend={italic}]{Storia de’solenni possessi de’Sommi Pontefici, detti anticamente processi}, Rome, Presso Luigi Lazzarini stampatore della R. C. A., 1802, p. 253. À la p. 244, Cancellieri rapporte, citant le \teiHi[rend={italic}]{Journal }de Gigli, que le 15 septembre 1644, durant le conclave, une colombe était entrée dans la salle pour aller se poser sur la loge des bénédictions à Saint Pierre (p. 251) ; une colombe aurait aussi pénétré dans la chambre du futur pontife (p. 208). Cancellieri rapporte aussi que le 23 novembre 1644, date de la cérémonie de la \teiHi[rend={italic}]{possessio}, on annonçait une tempête, mais une colombe dissipa les nuages et « apporta la sérénité désirée » et le soleil « apparut riant de joie avec plus de splendeur \lbrack{}…\rbrack{} comme dans les premiers siècles, l’innocente colombe apporta la tranquillité au bon Noé » (p. 244 et 245). Sur ces manifestations surnaturelles, voir Yvan Loskoutoff, « Augures héraldiques de la papauté », dans Andreas Gottsmann, Pierantonio Piatti et Andreas E. Rehberg (dir.), \teiHi[rend={italic}]{Incorrupta monumenta Ecclesiam defendunt, Studi offerti a Sergio Pagano, Prefetto dell’Archivio Segreto Vaticano}, Cité du Vatican, Archivio Segreto Vaticano, 2018, 3 vol., t. 2, p. 976-994.}}.}
              \teiP[xmlid={p36-2}]{La guerre qui avait agité l’Europe centrale se conclut en 1648 par la paix de Westphalie. Le pontife fut contraint de reconnaître la liberté de culte dans les États des princes allemands. Dès lors, tous ses efforts consistèrent à minimiser cette défaite et à réaffirmer le rôle politique de la sainte Église romaine sur la scène internationale.}
              \teiP[xmlid={p37-2}]{Le jubilé de 1650 attira à Rome des foules nombreuses et fut donc une occasion parfaite pour la stratégie du pontife. Les œuvres qu’il commanda montrent largement l’image de la colombe au rameau d’olivier, appropriée pour un pape se proposant comme l’arbitre de la paix et de la renaissance universelle après le dramatique conflit. Elle paraît dans les édifices publics, civils et religieux, réalisés par la volonté d’Innocent X : le nouveau palais du Capitole, la rénovation de la basilique du Latran et des nefs de Saint-Pierre, les nouvelles prisons à la via Giulia, la coupole de Saint-Yves-de-la-Sapience, où l’image de la colombe triomphe à l’intérieur, au sommet de la lanterne, exploitant la coïncidence entre l’emblème papal et celui de la Sagesse divine. Dans cette phase prévaut le message de paix et de renaissance : la colombe est souvent représentée sans les lys qui complètent l’écu des Pamphili, mais toujours surmontée de la tiare et des clefs papales.}
              \teiP[xmlid={p38-2}]{Pour Giovan Battista Pamphili, l’ascension au trône pontifical fut aussi l’occasion de manifester le prestige de sa propre maison. Sur ce point, l’initiative prioritaire se développa sur le vaste chantier de la place Navone où, entre 1644 et 1650, reprirent les travaux du palais familial. Ils s’accompagnèrent du réaménagement de l’espace public de l’antique stade de Domitien, célébré depuis le \teiHi[rend={small-caps}]{iv}\teiHi[rend={sup}]{e} siècle comme lieu de martyre. L’exaltation du pouvoir personnel du pontife se développa par la construction d’une nouvelle façade pour la demeure, par le projet de transformer l’église de Sainte-Agnès-in-Agone en mausolée familial et de créer le Collège Innocentien, édifices qui furent tous semés d’une myriade de colombes au rameau d’olivier.}
              \teiP[xmlid={p39-2}]{La sacralité du lieu et la continuité entre monde païen et monde chrétien furent rehaussées par la monumentale présence de l’obélisque antique, élevé en 1649 au centre de la fontaine des Quatre Fleuves, conçue par le Bernin. L’œuvre de récupération fut soutenue par les recherches du jésuite Athanasius Kircher qui considérait les hiéroglyphes comme des symboles unissant les mystères antiques à ceux de la foi des chrétiens, éléments de connexion entre la réalité terrestre et le monde surnaturel originaire. Pour Kircher l’antiquité représentait l’époque la plus proche de la « Vérité Première », son étude se révélait importante en vue de la restauration d’une harmonie que les hommes avaient perdue par des siècles de péché et d’idolâtrie\teiNote[n={59},place={foot},xmlid={ftn137}]{\teiP[xmlid={p-205}]{Sur Kircher : Valerio Rivosecchi, \teiHi[rend={italic}]{Esotismo in Roma Barocca. Studi sul Padre Kircher}, Rome, Bulzoni, 1982 ; Carlos Ziller Camenietzki, \teiHi[rend={italic}]{L’harmonie du monde au }\teiHi[rend={small-caps italic}]{xvii}\teiHi[rend={italic sup}]{e}\teiHi[rend={italic}]{ siècle. Essai sur la pensée scientifique d’Athanasius Kircher}, thèse de doctorat, Maurice Clavelin (dir.), Université Paris IV, 1995, en particulier p. 466. On ignorait alors que l’obélisque, retrouvé dans le cirque de Maxence, où il gisait brisé en quatre morceaux, avait été fabriqué en Égypte sous Domitien. }}. Dans le volume monumental consacré par Kircher à l’\teiHi[rend={italic}]{Obeliscus Pamphilius}, publié en 1650, « \teiHi[rend={italic}]{Anno Jubilaei} », le portrait du pontife est accompagné des mots qui illustrent son emblème : « \lbrack{}Innocent X\rbrack{} lie les lys, la palme et l’olivier, il aime le temps prospère de la paix victorieuse » (fig. 73). Dans la dédicace du volume au pontife « Pam-philio », c’est-à-dire « Amour universel », Kircher affirme : « La colombe visible sur le vert rameau d’olivier s’élève de son vol pour apaiser la divinité irritée et pour obtenir la tranquillité de la paix\teiNote[n={60},place={foot},xmlid={ftn138}]{\teiP[xmlid={p-206}]{Rudolf Preimesberger, « Obeliscus Pamphilius. Beitrage zur Vorgeschichte und Ikonographie des Vierstromebrunnes auf Piazza Navona », \teiHi[rend={italic}]{Münchner Jahrbuch der Bildenden Kunst}, n\teiHi[rend={sup}]{o} 25, 1974. p. 77-162 : « \teiHi[rend={italic}]{i gigli, la palma e l’ulivo lega insieme, ama il tempo prospero della pace vittoriosa} » ; « \teiHi[rend={italic}]{La colomba, visibile sul verde ramo d'ulivo, si alza in volo per placare la divinità arrabbiata, e per \lbrack{}ottenere\rbrack{} la tranquillità della pace} ».}} ».}
              \teiP[xmlid={p40-2}]{L’érection de l’obélisque place Navone peut être considérée comme un épisode du repositionnement des obélisque antiques comme pierres milliaires de la sainte Rome moderne, initié par Sixte V. Dans un document de 1647 on lit que par l’érection de l’obélisque de la place Navone « furent imités les vestiges de Sixte V\teiNote[n={61},place={foot},xmlid={ftn139}]{\teiP[xmlid={p-207}]{Maria Grazia D’Amelio, Tod A. Marder, « La fontana dei Quattro Fiumi a piazza Navona, iconologia e costruzione », dans Jean-François Bernard (dir.), \teiHi[rend={italic}]{Piazza Navona, ou Place Navone, la plus belle \& \lbrack{}et\rbrack{} la plus grande}, Rome, École française de Rome, « Collection de l’École Française de Rome, 493 », 2014, p. 399-419, voir aussi Stephanie Leone (dir.), \teiHi[rend={italic}]{The Pamphilj and the Arts: Patronage and Consumption in Baroque Rome, }Chestnut Hill-Chicago, McMullen Museum of Art-Boston College-University of Chicago Press, 2011, p. 23-36.}} », qui en avait fait les symboles de la puissance du catholicisme et de la réorganisation de la ville comme centre de l’Église universelle et destination de pèlerinage salvateur.}
              \teiP[xmlid={p41}]{Dans le cas de la place Navone, la présentation de l’obélisque assume une signification symbolique actualisée et spécifique, qui s’exprime en usant de formes différentes de celles, sobres, du \teiHi[rend={small-caps}]{xvi}\teiHi[rend={sup}]{e} siècle, par l’introduction de l’invention berninienne révolutionnaire du « rocher » comme soutien des statues des quatre fleuves. L’insertion du « rocher » permettait de compenser la hauteur limitée du monument antique ; d’autre part, sa structure en X soulignait sa signification de symbole solaire orienté vers les quatre points cardinaux tout en évoquant le chiffre antique du nom du pontife, Innocent « X ».}
              \teiP[xmlid={p42}]{L’iconographie des fleuves comme allusion aux quatre parties du monde et la polyvalence symbolique de la colombe sommitale entendaient célébrer sur un mode visuel efficace le triomphe de l’Église romaine sur les quatre continents et manifester un auspice de paix universelle, objectifs perçus comme réalisables seulement par l’œuvre du pontife, ministre du Saint-Esprit\teiNote[n={62},place={foot},xmlid={ftn140}]{\teiP[xmlid={p-208}]{L’insertion du Rio della Plata, exploré en 1516, évoque l’évangélisation des territoires américains par les pères jésuites réformés du Saint Nom de Jésus. }}. La fontaine des Quatre Fleuves mettait en spectacle \lbrack{}\teiHi[rend={italic}]{mostra}\rbrack{} le prolongement de l’\teiHi[rend={italic}]{acqua Vergine} ; d’un point de vue symbolique, la présence de l’eau permettait une reconsécration de l’espace, scellée par la colombe, symbole baptismal de la purification offerte aux pèlerins pour l’année jubilaire.}
              \teiP[xmlid={p43}]{Une série d’écrits, publiés en 1651, parallèlement au Jubilé et à l’inauguration de la fontaine, tels le \teiHi[rend={italic}]{Très copieux discours de la fontaine} d’Antonio Bernal di Gioya\teiNote[n={63},place={foot},xmlid={ftn141}]{\teiP[xmlid={p-209}]{Antonio Bernal di Gioya\teiHi[rend={italic}]{, Copiosissimo discorso della fontana, e guglia eretta in Piazza Navona, per ordine della Santità di Nostro Signore Innocentio X dal Signor Cavalier Bernini: con una abondante dichiaratione delli quattro fiumi...}, Roma, Tiberius, 1651. On y lit qu’Innocent X « \teiHi[rend={italic}]{imposuit niloticis enigmatibus exaratum lapidem ut salubrem spatiantibus amenitatem, sitientibus potum, meditantibus escam, magnifice largiretur} ».}} ou la \teiHi[rend={italic}]{Description de la fontaine Pamphilia} de Michelangelo Lualdi, explicitent ces significations\teiNote[n={64},place={foot},xmlid={ftn142}]{\teiP[xmlid={p-210}]{Michelangelo Lualdi, \teiHi[rend={italic}]{Descrittione della fontana Pamphilia, dove fu già il cerchio Agonale}, In Roma, Nella stamperia di Francesco Moneta, 1651.}}.}
              \teiP[xmlid={p44}]{À de tels éléments s’ajoute sur la fontaine une composante héraldique explicite, spécifiquement allusive à la maison des Pamphili, constituée par les deux écus monumentaux qui ornent les versants nord et sud, couronnés de la tiare et des clefs de saint Pierre. Un relief particulier leur est accordé par l’usage de précieux marbre de Carrare se détachant sur le travertin des parties environnantes. L’écu du côté nord (fig. 74), réalisé par Jacopo Antonio Fancelli, est soutenu par le Nil. Il est placé au sein d’une coquille, encadré d’épis, de palmes et de pivoines, symboles de paix et d’immortalité. L’écu du côté sud (fig. 75), œuvre d’Antonio Raggi, est porté par le Danube, symbole de l’aire contestée de la nouvelle orthodoxie protestante. Il est flanqué de cornes d’abondance et surmonté d’une coquille jouxtée de dauphins. La majeure partie des éléments présents autour de l’écu du nord se réfère donc à la paix, alors que pour celui du sud, les dauphins font allusion au prince « clément et vigilant » et les cornes d’abondance à la prospérité, dans un temps qui subissait en réalité les famines causées par la guerre de Trente Ans.}
              \teiP[xmlid={p45}]{Cette fontaine, assez connue et étudiée, doit toutefois être replacée dans le contexte d’une recherche dédiée à l’héraldique et la papauté, car elle représente sur ce point une œuvre emblématique exemplaire du pontificat Pamphili, une remarquable image de pierre, unitaire et complexe, à une époque où le langage de l’allégorie utilise les instruments de communication de l’héraldique.}
              \teiP[xmlid={p46}]{Le message de paix proposé par la fontaine des Quatre Fleuves fut ultérieurement repris et explicité dans l’iconographie du défilé de Pâques jubilaires de 1650, organisé sur la place Navone par la Confraternité espagnole de la Résurrection et dans l’église Saint-Jacques-des-Espagnols, qui répétait de différentes manières le message des machines éphémères installées lors de la \teiHi[rend={italic}]{possessio} d’Innocent X\teiNote[n={65},place={foot},xmlid={ftn143}]{\teiP[xmlid={p-211}]{Colombes et lys constellent aussi les apparats dessinés par Girolamo et Carlo Rainaldi pour la cérémonie de la \teiHi[rend={italic}]{possessio} (Francesco Cancellieri, \teiHi[rend={italic}]{op. cit.}, p. 248, 249 et 251). }}. L’installation, dessinée par Carlo Rainaldi, artiste de confiance des ambassadeurs d’Espagne autant que des Pamphili\teiNote[n={66},place={foot},xmlid={ftn144}]{\teiP[xmlid={p-212}]{Maurizio Fagiolo Dell’Arco, \teiHi[rend={italic}]{Corpus delle feste a Roma, }vol.\teiHi[rend={italic}]{ }1\teiHi[rend={italic}]{, La festa barocca, }Rome, De Luca, 1997, p. 349-353. Rainaldi aurait aussi construit le catafalque pour les funérailles en effigie de Philippe IV à Sainte-Marie-Majeure (1665).}}, prévoyait un théâtre fortifié et des arches de bois, avec des obélisques faisant allusion à la mission de défense de la foi catholique confiée à cette nation\teiNote[n={67},place={foot},xmlid={ftn145}]{\teiP[xmlid={p-213}]{Maria Antonietta Visceglia, « Giubilei tra pace e guerre (1625-1650) », \teiHi[rend={italic}]{Roma moderna e contemporanea}, vol. V, n\teiHi[rend={sup}]{o} 2-3, 1997, p. 431-474. }}.}
              \teiP[xmlid={p47}]{Les événements prirent en réalité un tour bien différent de ce qu’on attendait du jubilé à Rome : en ville on dut déplorer des rixes, des hostilités entre Français et Espagnols, une pénurie de denrées alimentaires, qui conduisirent les objectifs poursuivis par le pape Pamphili à un certain échec, faisant de l’\teiHi[rend={italic}]{Urbe} l’arène privilégiée des conflits européens et non le siège de l’accomplissement de la paix.}
              \teiP[xmlid={p48}]{Aux initiatives prévues pour le jubilé de 1650 appartient aussi la restauration de la basilique Saint-Jean-de-Latran\teiNote[n={68},place={foot},xmlid={ftn146}]{\teiP[xmlid={p-214}]{Augusto Roca De Amicis, \teiHi[rend={italic}]{L’opera di Borromini in San Giovanni in Laterano: gli anni della fabbrica (1646,1650)}, Rome, Dedalo, 1996.}}. Son réaménagement aurait été suggéré de manière miraculeuse par la découverte d’une colombe peinte sur les murs antiques\teiNote[n={69},place={foot},xmlid={ftn147}]{\teiP[xmlid={p-215}]{Alessandro Zuccari, « Borromini tra religiosità borromaica e cultura scientifica », dans Alessandro Zuccari et Stefania Macioce (dir.), \teiHi[rend={italic}]{Innocenzo X Pamphili. Arte e potere a Roma nell’età barocca }\lbrack{}2\teiHi[rend={sup}]{e} éd.\rbrack{}, actes de colloque, Rome, novembre 1990, Rome, Logart Press, 2001, p. 35-74.}}. En valorisant la basilique, on mettait en œuvre une sorte de triangulation du pouvoir papal répartie entre les pôles du Vatican, du Quirinal, demeure du pontife, et du Latran, son siège d’évêque de Rome.}
              \teiP[xmlid={p49}]{Dans les interventions de Borromini à Saint-Jean, le blason des Pamphili est exhibé avec variété, acquérant une valeur de source pour l’invention des formes. Innocent X et Borromini, aidés de Virgilio Spada, \teiHi[rend={italic}]{moderator consiliorum omnium atque arbiter} \lbrack{}modérateur et arbitre de tous les conseils\rbrack{}, confirmèrent leur intérêt pour les vestiges des origines chrétiennes, promues par Baronio et par les oratoriens, auxquels le pontife était aussi très lié, suivant les orientations de son oncle le cardinal Girolamo. L’objectif fut de « conserver les antiquités », en suivant une direction opposée à celle qu’on avait adoptée pour les travaux de la basilique Saint-Pierre, démolie au \teiHi[rend={small-caps}]{xvi}\teiHi[rend={sup}]{e} siècle et complètement reconstruite. En vue du jubilé, sous la direction de Gian Lorenzo Bernini, les pilastres de la nef de Saint-Pierre s’enrichirent d’une décoration de marbres polychromes, exaltant l’histoire de l’Église en même temps que les pontifes du passé et du présent : non moins de 104 emblèmes de la colombe Pamphili « plus grands que nature » y furent placés à côté des portraits des papes.}
              \teiP[xmlid={p50}]{Au Latran, en revanche, Borromini réussit à conserver le plan à cinq nefs et les parois de l’édifice constantinien qui, quoiqu’assez détériorées, furent habilement consolidées et restèrent visibles à travers douze oculi ouverts sur la partie haute de la nef centrale. La basilique se configura comme une extraordinaire « architecture parlante »\teiHi[rend={italic}]{,} dans laquelle un nouvel espace se trouva plus que jamais conjugué avec des significations ingénieuses \lbrack{}\teiHi[rend={italic}]{concetti}\rbrack{}, convergeant dans un message clair et cohérent, amplifié par la blancheur de la nef\teiNote[n={70},place={foot},xmlid={ftn148}]{\teiP[xmlid={p-216}]{Marcello Fagiolo, « Borromini in Laterano. Il Nuovo Tempio per il Concilio universale », \teiHi[rend={italic}]{L’arte}, n\teiHi[rend={sup}]{o} 13, 1971, p. 1-46 ; « I geroglifici, gli emblemi e l’araldica: note sul ragionamento simbolico di Borromini », dans Richard Boesel et Christoph L. Frommel (dir.), \teiHi[rend={italic}]{Borromini e l’universo barocco}, Milan, Electa, 1999, p. 95-105.}} (fig. 76). Le choix de l’architecte de faire confluer formes et motifs symboliques dans un tout intimement lié se manifeste par la copieuse élaboration de figures héraldiques répandues dans toute la basilique. Même le pavement propose la combinaison de marbres antiques et modernes qui montrent encore une fois la colombe et des lys. L’usage du symbole du \teiHi[rend={italic}]{chrismon }est répandu\teiHi[rend={italic}]{, }faisant allusion au Christ, mais aussi à Innocent « X\teiNote[n={71},place={foot},xmlid={ftn149}]{\teiP[xmlid={p-217}]{Alessandro Zuccari, art. cité, p. 35-74. Dans le même volume, voir Paolo Portoghesi, « Borromini e Innocenzo X: architettura e politica pontificia », p. 27-34.}} ». Dans la décoration en stuc surmontant la double porte d’accès à la basilique depuis le palais du Latran trône le blason du pape Pamphili, inséré sur une sphère symbolisant le monde dominé par son autorité ; des emblèmes du pontife donnent force plastique aux portes correspondantes vers l’intérieur de l’église (fig. 77 et 78). L’emblème reparaît aussi dans les stucs monumentaux des arches qui séparent le transept des nefs latérales, et dans les plus amples arcades centrales, inscrites au milieu de la séquence qui divise la nef principale des nefs secondaires.}
              \teiP[xmlid={p51}]{De cette façon, l’espace de la basilique évoque l’image de la Jérusalem céleste, décrite dans l’\teiHi[rend={italic}]{Apocalypse} de saint Jean : les tabernacles de marbre foncé de la nef centrale, destinés à accueillir les statues des apôtres, font référence aux douze portes de la cité biblique. Même les chérubins qui se succèdent dans les ailes latérales rappellent ceux qui décoraient le temple de Jérusalem, fondé par Salomon en tant que temple de la paix \lbrack{}\teiHi[rend={italic}]{templum Pacis}\rbrack{}, comme les palmes, les grenades et les images des colonnes salomoniques, sculptées dans des stucs clairs qui décorent la trabéation de la contre-façade\teiNote[n={72},place={foot},xmlid={ftn150}]{\teiP[xmlid={p-218}]{Marcello Fagiolo, art\teiHi[rend={italic}]{. }cité, p. 99 et 100. Pour la corrélation entre Saint-Yves-de-la-Sapience et le temple/palais de Salomon : Alessandro Zuccari, art. cité, p. 55, 56 et 59-61.}}.}
              \teiP[xmlid={p52}]{Un autre aspect de ce projet met en évidence de manière exemplaire l’utilisation de l’emblème papal comme source de solutions architecturales et décoratives visuellement suggestives. De nombreux dessins conservés à l’Albertina de Vienne, documentent la recherche passionnée et tenace de Borromini pour inscrire l’image de la colombe au centre de la composition des verrières\teiNote[n={73},place={foot},xmlid={ftn151}]{\teiP[xmlid={p-219}]{Les didascalies autographes de Borromini pour les verrières, inscrites dans les dessins de l’Albertina, précisent l’emplacement exact prévu pour les armoiries papales complètes, avec les clefs et la tiare : les lunettes des « nefs extrêmes » (Az Rom 440 et 441) ; « les verrières de la nef principale » (Az Rom 442) ; les « fenêtres ovales entre une chapelle et l’autre » (27 avril 1650, Az Rom 443) ; la « fenêtre principale de la façade » (Az Rom 444) ; les verrières de la « nef du milieu », c’est-à-dire de « la seconde nef, entre la majeure et la mineure » (Az Rom 445) ; les « fenêtres supérieures de la nef centrale » (Az Rom 451). Augusto Roca de Amicis, \teiHi[rend={italic}]{op. cit.}, 1996, p. 108-111 et 184-186.}}. Celles-ci forment une part essentielle du projet d’élaborer un ingénieux système de gradation de la lumière des nefs de la basilique, jouant habilement des différents niveaux des toits\teiNote[n={74},place={foot},xmlid={ftn152}]{\teiP[xmlid={p-220}]{Az Rom 454 : Augusto Roca De Amicis, \teiHi[rend={italic}]{op. cit.}, fig. 100, p. 186.}}(fig. 79). Aux dessins d’étude s’ajoutent les dessins de « présentation », réalisés en collaboration avec le forgeron ; dans ceux-ci, le motif héraldique est mis en évidence par la simple structure du fer au carré, précisément définie dans les mesures. Ces verrières ont disparu ; toutefois les dessins de Borromini montrent comment la fusion de la forme et du symbole, de la transparence et des images, systématiquement poursuivie, était parvenue à de remarquables résultats.}
              \teiP[xmlid={p53}]{Le blason présenté en peintures et sculptures, dans des édifices, des stucs et des plafonds de bois, mais aussi reproduit dans des gravures, des reliures de livres, des monnaies et des médailles, fut un instrument de la divulgation de l’image héraldique identitaire du pontificat d’Innocent X. Une propagande bien diffusée qui se heurtait toutefois à une situation internationale difficile. Le pape lui-même dans son bref \teiHi[rend={italic}]{Zelo domus Dei}, promulgué le 20 août 1650, la considère compromise par l’« infâme paix de Münster ». En effet, sur la scène européenne, le projet universel d’Innocent X se trouva affaibli par la diffusion du culte protestant et par l’émergence des identités nationales.}
            
\end{teiDiv}
            \begin{teiDiv}[type={section1},xmlid={div04-2}]

              \teiHead[xmlid={epilogue-une-nouvelle-fortune-pour-le-blason-pamphili-la-raccolta-di-varie-targhe-di-roma-de-filippo-juvarra-1711-001}]{Épilogue. Une nouvelle fortune pour le blason pamphili : la \teiHi[rend={italic}]{Raccolta Di Varie Targhe di Roma }de Filippo Juvarra (1711)}
              \teiP[xmlid={p54}]{L’emblème du pape Pamphili fut traité au siècle suivant dans un contexte nouveau et différent. Le \teiHi[rend={italic}]{Recueil de diverses targes de Rome faites par les premiers professeurs à Rome, dessinées et gravées} de Filippo Juvarra, fut publié par Antonio De Rossi, à Rome, en 1711\teiNote[n={75},place={foot},xmlid={ftn153}]{\teiP[xmlid={p-221}]{Voir Raffaella Zega, « La “Raccolta di varie targhe di Roma...” nell'attività di Filippo Juvarra », \teiHi[rend={italic}]{Archivum Arcis}, n\teiHi[rend={sup}]{o} 2, 1989, p. 10-51 ; Maria Vittoria Cattaneo et Elena Gianasso, « Il Corpus Juvarrianum della Biblioteca Nazionale Universitaria di Torino », dans Franca Porticelli, Costanza Roggero, Chiara Devoti \teiHi[rend={italic}]{et al}. (dir.), \teiHi[rend={italic}]{Filippo Juvarra regista di corti e capitali dalla Sicilia al Piemonte all'Europa}, Turin, Centro Studi Piemontesi, 2020, p. 197-210.}}. Il s’agit d’une anthologie de 50 gravures représentant des armoiries de papes ou de cardinaux, figurées dans des œuvres d’architectes célèbres, de Bramante à l’époque contemporaine. Une bonne vingtaine d’images proviennent du Bernin, six de Borromini. L’écu d’Innocent X est présent dans les images qui reproduisent les targes insérées dans la façade latérale du Capitole (n\teiHi[rend={sup}]{o} 45), dans la fontaine de la cour du Vatican (n\teiHi[rend={sup}]{o} 35), sur les deux côtés de la fontaine des Quatre Fleuves, vers le nord (n\teiHi[rend={sup}]{o} 15, fig. 80) et vers le sud (n\teiHi[rend={sup}]{o} 18, fig. 81). Trois planches illustrent les écus de Saint-Jean-de-Latran (n\teiHi[rend={sup}]{o} 27 ; n\teiHi[rend={sup}]{o} 30, fig. 82 ; n\teiHi[rend={sup}]{o} 31\teiNote[n={76},place={foot},xmlid={ftn154}]{\teiP[xmlid={p-222}]{Deux images sont consacrées aux écus de Camillo Pamphili, figurant à San Nicola di Tolentino (pl. 34) et à Sant’Andrea al Quirinale (pl. 13).}}). En conclusion de l’œuvre, une gravure propose quelques graphiques indiquant les normes pour les proportions et la composition géométrique de divers types d’écus, différenciés suivant le rang du personnage célébré.}
              \teiP[xmlid={p55}]{Au cours de son activité, Juvarra se révéla très intéressé par la décoration héraldique, comme le prouve la grande quantité d’esquisses et de dessins préparatoires au \teiHi[rend={italic}]{Recueil des targes.} Nous disposons de près de 235 dessins réunis dans des albums conservés à New York, au Metropolitan Museum et à Turin, au Musée municipal, à la Bibliothèque nationale et dans la collection Tournon\teiNote[n={77},place={foot},xmlid={ftn155}]{\teiP[xmlid={p-223}]{L’album du Metropolitan Museum (MM 165) contient 12 dessins d’écus préparatoires au recueil. Mary Myers (\teiHi[rend={italic}]{Architectural and ornament drawings}: \teiHi[rend={italic}]{Juvarra, Vanvitelli, The Bibiena family \& other Italian draughtsmen, }New York, The Metropolitan Museum of Art, 1975) démontre que les dessins du musée municipal de Turin, préparatoires pour le recueil, furent exécutés par Juvarra sur la base des dessins du Metropolitan et de ceux de la réserve 59.4 de la Bibliothèque nationale de Turin, preuve d’un processus d’étude soigneux et persévérant de la part de l’artiste. Sarah Mc Phee, « A new sketch-book by Filippo Juvarra », \teiHi[rend={italic}]{The Burlington Magazine}, n\teiHi[rend={sup}]{o} 135, 1993, p. 346-350 ; « Filippo Juvarra: a new volume of drawings from the Roman years », \teiHi[rend={italic}]{Studi piemontesi}, n\teiHi[rend={sup}]{o} 22, 1993, p. 377-383 ; Henry A. Millon, \teiHi[rend={italic}]{Filippo Juvarra. Drawings from the roman period 1704-1714, }Rome, Edizioni dell’Elefante, 1984, vol. 1, p. XXII, 181, 188-200, 248 et 336-337. Une liste comparative des dessins d’étude de Juvarra et des planches contenues dans le recueil se trouve dans Tommaso Manfredi, \teiHi[rend={italic}]{Filippo Juvarra. Gli anni giovanili, }Rome, Argos, 2010, p. 426 ; là, les notes 228, 229, 231, 235 sont consacrées aux écus d’Innocent X.}}. Quelques études, consacrées à des écus du pape Innocent X, sont conservées à la bibliothèque Vaticane dans le manuscrit Vat. lat. 13295, d’attribution controversée\teiNote[n={78},place={foot},xmlid={ftn156}]{\teiP[xmlid={p-224}]{Le ms. BAV Vat. lat. 13295 contient dix-huit dessins pour la décoration héraldique des piliers de la basilique Saint-Pierre. L’édition critique du recueil vatican, dans lequel au moins huit dessins sont préparatoires pour le recueil des \teiHi[rend={italic}]{Targhe,} se trouve dans Henry A. Millon, Andreina Griseri, Sarah Mc Phee \teiHi[rend={italic}]{et al.} (dir.), \teiHi[rend={italic}]{Filippo Juvarra: drawings from the Roman period 1704-1714}, Rome, Edizioni dell’Elefante, 1999, vol. 2. Pour l’attribution du manuscrit vatican : Tommaso Manfredi, \teiHi[rend={italic}]{op. cit.}, p. 427.}}. Les dessins ont souvent été réalisés en réinterprétant de façon personnelle les modèles et enrichis d’ombres à l’aquarelle qui en augmentent les effets picturaux (fig. 83). Pour Juvarra, ce travail revêtait une fonction de formation et d’apprentissage, conduit tantôt sur la base d’une observation directe des emblèmes présents sur des constructions, tantôt en analysant les estampes et les dessins qui les reproduisent. La dimension des feuilles fait penser que, au moins depuis 1706, il poursuivait le projet de traduire ses études dans un recueil de gravures, afin de réaliser une sorte d’encyclopédie de la décoration baroque.}
              \teiP[xmlid={p56}]{Les expressions héraldiques innovantes et spectaculaires des maîtres, diffusées grâce à la sérialité des estampes, devinrent des modèles pour la mise au point d’un nouveau langage décoratif. Dédiant son œuvre au cardinal Ottoboni, Juvarra écrit qu’elle est destinée « au bénéfice de ceux qui \lbrack{}…\rbrack{} se proposent de tirer profit de l’Académie du dessin\teiNote[n={79},place={foot},xmlid={ftn157}]{\teiP[xmlid={p-225}]{« \teiHi[rend={italic}]{a beneficio \lbrack{}di coloro\rbrack{} i quali…nell’Accademia del disegno procurano di approfittarsi} », Filippo Juvarra,\teiHi[rend={italic}]{ Raccolta Di Varie Targhe di Roma Fatte Da Professori Primarj In Roma, Disegnate, ed intagliate}, Rome, per Antonio de’ Rossi, 1711, lettre dédicatoire.}} » : au fin de fournir un répertoire satisfaisant les exigences d’une clientèle désireuse d’autoreprésentation, s’ajoute la fonction d’instrument utile aux enseignements de l’Académie romaine de saint Luc, lieu institutionnel de la formation des artistes.}
              \begin{teiFigure}[xmlid={figure01-4}]

                \teiGraphic[url={../icono/br/Ch05\_DelPesco/fig61.jpg}]
                \teiHead[xmlid={figure-59-giovanni-luigi-valesio-armes-du-cardinal-girolamo-pamphili-stemma-del-cardinale-innocenzo-x-panfili-e-figure-allegoriche-gravure-bologne-pinacoteca-nazionale-gabinetto-disegni-e-stampe-83-picc-coll-ii-001}]{Figure 59. Giovanni Luigi Valesio, armes du cardinal Girolamo Pamphili, \teiHi[rend={italic}]{Stemma del cardinale Innocenzo X Panfili e figure allegoriche, }gravure. Bologne, Pinacoteca Nazionale, Gabinetto Disegni e Stampe, 83 Picc. Coll. II.}
              
\end{teiFigure}
              \begin{teiFigure}[xmlid={figure02-4}]

                \teiGraphic[url={../icono/br/Ch05\_DelPesco/fig62.jpg}]
                \teiHead[xmlid={figure-60-rome-santa-maria-in-vallicella-paolo-maruscelli-et-collaborateurs-memorial-du-cardinal-girolamo-pamphili-1634-001}]{Figure 60. Rome, Santa Maria in Vallicella, Paolo Maruscelli et collaborateurs, mémorial du cardinal Girolamo Pamphili, 1634.}
              
\end{teiFigure}
              \begin{teiFigure}[xmlid={figure03-4}]

                \teiGraphic[url={../icono/br/Ch05\_DelPesco/fig63.jpg}]
                \teiHead[xmlid={figure-61-rome-palais-pamphili-salle-des-paysages-armes-papales-dinnocent-x-pamphili-001}]{Figure 61. Rome, palais Pamphili, salle des Paysages, armes papales d’Innocent X Pamphili.}
              
\end{teiFigure}
              \begin{teiFigure}[xmlid={figure04-4}]

                \teiGraphic[url={../icono/br/Ch05\_DelPesco/fig64.jpg}]
                \teiHead[xmlid={figure-62-gubbio-palais-pamphili-facade-sur-la-place-001}]{Figure 62. Gubbio, palais Pamphili, façade sur la place.}
              
\end{teiFigure}
              \begin{teiFigure}[xmlid={figure05-4}]

                \teiGraphic[url={../icono/br/Ch05\_DelPesco/fig65.jpg}]
                \teiHead[xmlid={figure-63-gubbio-palais-pamphili-salle-a-letage-noble-embleme-de-la-colombe-pamphili-inscrit-dans-la-cheminee-001}]{Figure 63. Gubbio, palais Pamphili, salle à l’étage noble, emblème de la colombe Pamphili inscrit dans la cheminée.}
              
\end{teiFigure}
              \begin{teiFigure}[xmlid={figure06-4}]

                \teiGraphic[url={../icono/br/Ch05\_DelPesco/fig66.jpg}]
                \teiHead[xmlid={figure-64-gubbio-palais-pamphili-salon-a-letage-noble-detail-de-la-decoration-a-fresque-avec-lembleme-de-la-colombe-pamphili-001}]{Figure 64. Gubbio, palais Pamphili, salon à l’étage noble, détail de la décoration à fresque avec l’emblème de la colombe Pamphili.}
              
\end{teiFigure}
              \begin{teiFigure}[xmlid={figure07-4}]

                \teiGraphic[url={../icono/br/Ch05\_DelPesco/fig67.jpg}]
                \teiHead[xmlid={figure-65-gubbio-ancien-cimetiere-pres-de-leglise-san-secondo-chapelle-pamphili-jacopo-bedi-martyre-de-saint-sebastien-1458-fresque-detail-001}]{Figure 65. Gubbio, ancien cimetière près de l’église San Secondo, chapelle Pamphili, Jacopo Bedi, \teiHi[rend={italic}]{Martyre de saint Sébastien}, 1458, fresque (détail).}
              
\end{teiFigure}
              \begin{teiFigure}[xmlid={figure08-3}]

                \teiGraphic[url={../icono/br/Ch05\_DelPesco/fig68.jpg}]
                \teiHead[xmlid={figure-66-gubbio-ancien-cimetiere-pres-de-leglise-san-secondo-chapelle-pamphili-pierre-tombale-avec-lecu-des-pamphili-001}]{Figure 66. Gubbio, ancien cimetière près de l’église San Secondo, chapelle Pamphili, pierre tombale avec l’écu des Pamphili.}
              
\end{teiFigure}
              \begin{teiFigure}[xmlid={figure09-3}]

                \teiGraphic[url={../icono/br/Ch05\_DelPesco/fig69.jpg}]
                \teiHead[xmlid={figure-67-gubbio-palais-des-consuls-architrave-du-portail-dentree-avec-emblemes-de-la-ville-pontificaux-des-angevins-et-les-dates-de-la-construction-1332-1336-001}]{Figure 67. Gubbio, palais des consuls, architrave du portail d’entrée avec emblèmes de la ville, pontificaux des Angevins et les dates de la construction (1332-1336).}
              
\end{teiFigure}
              \begin{teiFigure}[xmlid={figure10-3}]

                \teiGraphic[url={../icono/br/Ch05\_DelPesco/fig70.jpg}]
                \teiHead[xmlid={figure-68-gubbio-citta-regia-antichiss-ima-dell-umbria-gravure-dapres-un-dessin-dignazio-cassetta-johannes-blaeu-amsterdam-1663-reproduit-dans-johannes-blaeu-nouveau-theatre-ditalie-ou-description-exacte-de-ses-villes-palais-eglises-etc-tome-quatrieme-a-amsterdam-par-le-soin-de-pierre-mortier-libraire-1704-detail-montrant-la-maison-des-pamphili-indiquee-par-la-lettre-z-001}]{Figure 68.\teiHi[rend={italic}]{ Gubbio Città Regia Antichiss\lbrack{}ima\rbrack{} dell'Umbria}, gravure, d’après un dessin d’Ignazio Cassetta, Johannes Blaeu, Amsterdam, 1663, reproduit dans Johannes Blaeu, \teiHi[rend={italic}]{Nouveau théâtre d’Italie ou description exacte de ses villes, palais, églises, etc. Tome Quatrième}, À Amsterdam, Par le soin de Pierre Mortier Libraire, 1704. Détail montrant la maison des Pamphili, indiquée par la lettre Z.}
              
\end{teiFigure}
              \begin{teiFigure}[xmlid={figure11-3}]

                \teiGraphic[url={../icono/br/Ch05\_DelPesco/fig71.jpg}]
                \teiHead[xmlid={figure-69-armoiries-romae-panfiliorum-dans-silvestro-pietrasanta-tesserae-gentilitiae-romae-typis-haered-francisci-corbelletti-1638-p-438-001}]{Figure 69. Armoiries « \teiHi[rend={italic}]{Romae Panfiliorum} », dans Silvestro Pietrasanta, \teiHi[rend={italic}]{Tesserae gentilitiae,} Romae, Typis haered. Francisci Corbelletti, 1638, p. 438.}
              
\end{teiFigure}
              \begin{teiFigure}[xmlid={figure12-3}]

                \teiGraphic[url={../icono/br/Ch05\_DelPesco/fig72.jpg}]
                \teiHead[xmlid={figure-70-armoiries-de-la-famille-pamphili-sans-attributs-cardinalices-ou-papaux-rome-palais-pamphili-escalier-de-laile-vers-la-place-pasquino-001}]{Figure 70. Armoiries de la famille Pamphili (sans attributs cardinalices ou papaux), Rome, palais Pamphili, escalier de l’aile vers la place Pasquino.}
              
\end{teiFigure}
              \begin{teiFigure}[xmlid={figure13-3}]

                \teiGraphic[url={../icono/br/Ch05\_DelPesco/fig73.jpg}]
                \teiHead[xmlid={figure-71-rome-palais-pamphili-loggia-de-la-cour-centrale-angle-nord-est-001}]{Figure 71. Rome, palais Pamphili, loggia de la cour centrale, angle nord-est.}
              
\end{teiFigure}
              \begin{teiFigure}[xmlid={figure14-3}]

                \teiGraphic[url={../icono/br/Ch05\_DelPesco/fig74.jpg}]
                \teiHead[xmlid={figure-72-rome-palais-pamphili-loggia-de-la-cour-centrale-detail-001}]{Figure 72. Rome, palais Pamphili, loggia de la cour centrale (détail).}
              
\end{teiFigure}
              \begin{teiFigure}[xmlid={figure15-3}]

                \teiGraphic[url={../icono/br/Ch05\_DelPesco/fig75.jpg}]
                \teiHead[xmlid={figure-73-portrait-dinnocent-x-dans-athanasius-kircher-obeliscus-pamphilius-gravure-emblem-collection-of-the-university-of-illinois-001}]{Figure 73. Portrait d’Innocent X, dans Athanasius Kircher, \teiHi[rend={italic}]{Obeliscus Pamphilius}, gravure, Emblem Collection of the University of Illinois.}
              
\end{teiFigure}
              \begin{teiFigure}[xmlid={figure16-3}]

                \teiGraphic[url={../icono/br/Ch05\_DelPesco/fig76.jpg}]
                \teiHead[xmlid={figure-74-jacopo-antonio-fancelli-dapres-un-dessin-de-gian-lorenzo-bernini-armes-dinnocent-x-pamphili-marbre-fontaine-des-quatre-fleuve-armoiries-cote-nord-001}]{Figure 74. Jacopo Antonio Fancelli d’après un dessin de Gian Lorenzo Bernini, armes d’Innocent X Pamphili, marbre, Fontaine des Quatre Fleuve (armoiries côté nord).}
              
\end{teiFigure}
              \begin{teiFigure}[xmlid={figure17-2}]

                \teiGraphic[url={../icono/br/Ch05\_DelPesco/fig77.jpg}]
                \teiHead[xmlid={figure-75-antonio-raggi-dapres-un-dessin-de-gian-lorenzo-bernini-armes-dinnocenzo-x-pamphili-marbre-fontaine-des-quatre-fleuves-armoiries-cote-sud-001}]{Figure 75. Antonio Raggi, d’après un dessin de Gian Lorenzo Bernini, armes d’Innocenzo X Pamphili, marbre, Fontaine des Quatre Fleuves (armoiries côté sud).}
              
\end{teiFigure}
              \begin{teiFigure}[xmlid={figure18-2}]

                \teiGraphic[url={../icono/br/Ch05\_DelPesco/fig78.jpg}]
                \teiHead[xmlid={figure-76-basilique-de-saint-jean-de-latran-interieur-001}]{Figure 76. Basilique de Saint-Jean-de-Latran, intérieur.}
              
\end{teiFigure}
              \begin{teiFigure}[xmlid={figure19-2}]

                \teiGraphic[url={../icono/br/Ch05\_DelPesco/fig79.jpg}]
                \teiHead[xmlid={figure-77-basilique-de-saint-jean-de-latran-decor-en-stuc-avec-lecu-du-pape-pamphili-double-porte-entre-la-basilique-et-le-palais-du-latran-dapres-un-dessin-de-francesco-borromini-001}]{Figure 77. Basilique de Saint-Jean-de-Latran, décor en stuc avec l’écu du pape Pamphili, double porte entre la basilique et le palais du Latran, d’après un dessin de Francesco Borromini.}
              
\end{teiFigure}
              \begin{teiFigure}[xmlid={figure20-2}]

                \teiGraphic[url={../icono/br/Ch05\_DelPesco/fig80.png}]
                \teiHead[xmlid={figure-78-francesco-borromini-etude-pour-la-porte-double-entre-le-palais-du-latran-et-la-basilique-dessin-a-la-mine-de-plomb-vienne-albertina-az-rom-416-detail-on-lit-un-globo-che-rappresenta-il-mondo-sara-congiunto-insieme-il-spirituale-e-il-temporale-001}]{Figure 78. Francesco Borromini, étude pour la porte double entre le palais du Latran et la basilique, dessin à la mine de plomb, Vienne, Albertina, Az Rom 416 (détail). On lit : « \teiHi[rend={italic}]{un globo che rappresenta il mondo sarà congiunto insieme il spirituale e il temporale }».}
              
\end{teiFigure}
              \begin{teiFigure}[xmlid={figure21-2}]

                \teiGraphic[url={../icono/br/Ch05\_DelPesco/fig81.jpg}]
                \teiHead[xmlid={figure-79-francesco-borromini-etude-pour-la-verriere-de-la-fenetre-de-la-facade-orientale-de-saint-jean-de-latran-dessin-a-la-mine-de-plomb-vienne-albertina-az-rom-454-on-lit-finestra-della-facciata-alta-palmi-40-larga-20-portello-dellarma-del-papa-alto-10-34-largo-8-001}]{Figure 79. Francesco Borromini, étude pour la verrière de la fenêtre de la façade orientale de Saint-Jean-de-Latran, dessin à la mine de plomb, Vienne, Albertina, Az Rom 454. On lit : « \teiHi[rend={italic}]{finestra della facciata alta palmi 40 larga 20, portello dell’arma del papa alto 10 ¾ largo 8 }».}
              
\end{teiFigure}
              \begin{teiFigure}[xmlid={figure22-2}]

                \teiGraphic[url={../icono/br/Ch05\_DelPesco/fig82.jpg}]
                \teiHead[xmlid={figure-80-armes-dinnocent-x-pamphili-a-la-fontaine-des-quatre-fleuves-cote-nord-dans-filippo-juvarra-raccolta-di-varie-targhe-di-roma-fatte-da-professori-primarj-in-roma-disegnate-ed-intagliate-antonio-de-rossi-ed-a-roma-1711-pl-15-gravure-001}]{Figure 80. Armes d’Innocent X Pamphili à la fontaine des Quatre Fleuves, côté nord, dans Filippo Juvarra, \teiHi[rend={italic}]{Raccolta Di Varie Targhe di Roma Fatte Da Professori Primarj In Roma, Disegnate, ed intagliate,} Antonio De Rossi ed., a Roma, 1711, pl. 15, gravure.}
              
\end{teiFigure}
              \begin{teiFigure}[xmlid={figure23-2}]

                \teiGraphic[url={../icono/br/Ch05\_DelPesco/fig83.jpg}]
                \teiHead[xmlid={figure-81-armes-dinnocent-x-pamphili-a-la-fontaine-des-quatre-fleuves-cote-sud-dans-filippo-juvarra-raccolta-di-varie-targhe-di-roma-fatte-da-professori-primarj-in-roma-disegnate-ed-intagliate-antonio-de-rossi-ed-a-roma-1711-pl-18-gravure-001}]{Figure 81. Armes d’Innocent X Pamphili, à la fontaine des Quatre Fleuves, côté sud, dans Filippo Juvarra\teiHi[rend={italic}]{, Raccolta Di Varie Targhe di Roma Fatte Da Professori Primarj In Roma, Disegnate, ed intagliate}, Antonio De Rossi ed., a Roma, 1711, pl. 18, gravure.}
              
\end{teiFigure}
              \begin{teiFigure}[xmlid={figure24-2}]

                \teiGraphic[url={../icono/br/Ch05\_DelPesco/fig84.jpg}]
                \teiHead[xmlid={figure-82-armes-dinnocent-x-pamphili-a-saint-jean-de-latran-arche-centrale-de-la-nef-principale-in-filippo-juvarra-raccolta-di-varie-targhe-di-roma-fatte-da-professori-primarj-in-roma-disegnate-ed-intagliate-antonio-de-rossi-ed-a-roma-1711-pl-30-gravure-001}]{Figure 82. Armes d’Innocent X Pamphili à Saint-Jean-de-Latran, arche centrale de la nef principale, in Filippo Juvarra, \teiHi[rend={italic}]{Raccolta Di Varie Targhe di Roma Fatte Da Professori Primarj In Roma, Disegnate, ed intagliate}, Antonio De Rossi ed., a Roma, 1711, pl. 30, gravure.}
              
\end{teiFigure}
              \begin{teiFigure}[xmlid={figure25-2}]

                \teiGraphic[url={../icono/br/Ch05\_DelPesco/fig85.jpg}]
                \teiHead[xmlid={figure-83-filippo-juvarra-etude-pour-la-pl-30-de-la-raccolta-di-varie-targhe-ecu-pamphili-saint-jean-de-latran-arche-centrale-de-la-nef-principale-dessin-mm-117-plume-crayon-encre-marron-aquarelle-new-york-the-metropolitan-museum-of-art-001}]{Figure 83. Filippo Juvarra, étude pour la pl. 30 de la \teiHi[rend={italic}]{Raccolta Di Varie Targhe }(Écu Pamphili, Saint-Jean-de-Latran, arche centrale de la nef principale), dessin, MM 117, plume, crayon, encre marron, aquarelle. New York, The Metropolitan Museum of Art.}
              
\end{teiFigure}
            
\end{teiDiv}
          
\end{teiElement}
        
\end{teiElement}
        \begin{teiElement}[name={group},type={chapter},data-page-title={« Micat in vertice »},data-page-subtitle={L’héraldique d’Alexandre VII Chigi à Bologne, entre occultation et célébration (1658-1665)},data-page-authors={Sonia Cavicchioli},data-include-href={Ch06\_Micat\_in\_vertice.xml},xmlid={chapter-003}]

          \begin{teiElement}[name={front}]

            \begin{teiDiv}[type={titlePage},xmlid={titlepage-008}]

              \teiP[rend={title-main},xmlid={p-226}]{« \teiHi[rend={italic}]{Micat in vertice} »}
              \teiP[rend={title-sub},xmlid={p-227}]{L’héraldique d’Alexandre VII Chigi à Bologne, entre occultation et célébration (1658-1665)}
              \teiP[rend={author-aut},xmlid={p-228}]{Sonia Cavicchioli}
              \teiP[rend={editor-trl},xmlid={p-229}]{Traduit de l’italien par Yvan Loskoutoff.}
            
\end{teiDiv}
          
\end{teiElement}
          \begin{teiElement}[name={body},xmlid={body-008}]

            \teiP[xmlid={p1-8}]{Dans les siècles qui vont de l’intronisation de Jules II della Rovere, en novembre 1506, à l’arrivée des troupes françaises en 1796, Bologne appartint sans interruption aux États pontificaux. La cité continua toutefois à cultiver un sentiment municipal qui trouva son expression dans la volonté ou velléité des principales familles de faire valoir leur propre rôle. Bien que le gouvernement municipal fût de fait exercé par le cardinal légat représentant le pontife, une telle ambition fut satisfaite par le Saint-Siège en laissant persister, au moins pour la forme, les anciennes magistratures, particulièrement le Sénat qui, à la fin du \teiHi[rend={small-caps}]{xvi}\teiHi[rend={sup}]{e} siècle, comptait cinquante membres appartenant à autant de familles nobles autorisées à siéger par le pontife. À leur tête se trouvait le gonfalonier de justice qui restait en charge deux mois, assisté de huit Anciens, les consuls\teiNote[n={1},place={foot},xmlid={ftn1-6}]{\teiP[xmlid={p-230}]{Sur la particularité politique de Bologne, qui se répercute sur les commandes des multiples acteurs de la vie publique voir Angela De Benedictis, « Tra “affratellamento comunitario” e “commensualità rituale”. La “condotta di vita” dei cittadini di governo tra ’400 e ’500 », dans Maurizio Ricci (dir.), \teiHi[rend={italic}]{L’architettura a Bologna nel Rinascimento (1460-1550): centro o periferia?}, Bologne, Minerva, 2001, p. 17-27, et pour un vision plus articulée, \teiHi[rend={italic}]{Repubblica per contratto. Bologna, una città europea nello Stato della Chiesa, }Bologne, Il Mulino, 1995.}}.}
            \teiP[xmlid={p2-7}]{Le siège du gouvernement bolonais se concentra pendant des siècles au palais connu aujourd’hui sous le nom d’Accursio ou Public, qui constitue le côté occidental de la grand-place, où la coexistence des pouvoirs s’exprimait pour ainsi dire plastiquement. Sur la portion de gauche se trouvaient les bureaux et la résidence des Anciens, alors que sur celle de droite cohabitaient le Sénat municipal et le légat. Au premier étage se réunissait le Sénat, au second se trouvait l’appartement du légat, qui constituait notamment la vraie et propre résidence pontificale hors de Rome\teiNote[n={2},place={foot},xmlid={ftn43-3}]{\teiP[xmlid={p-231}]{Camilla Bottino (dir.), \teiHi[rend={italic}]{Il palazzo comunale di Bologna. Storia, architettura e restauri}, Bologne, Compositori, 1999, \teiHi[rend={italic}]{passim} et en particulier dans ce volume, Hans W. Hubert, « La nascita e lo sviluppo architettoncio del Palazzo del Comune di Bologna fra potere comunale e potere papale », p. 64-87 (synthèse de Hans W. Hubert, \teiHi[rend={italic}]{Der Palazzo Comunale von Bologna. Vom Palazzo della Biada zum Palatium Apostolicum}, Cologne, Weimar, 1993).}}. Une telle multiplicité de pouvoirs trouvait naturellement un écho dans d’autres édifices publics et à l’extérieur de la cité. Il en reste un témoignage visuel dans l’inscription et dans la série d’écus disposés sur la façade tournée vers la ville d’une des portes urbaines, la porte Galliera, reconstruite en 1661 suite à un écroulement dû à des infiltrations. Si l’épigraphe signalant la reconstruction in « \teiHi[rend={italic}]{augustiorem formam} » fut ordonnée par le Sénat de la cité (SPQB) au cours du pontificat d’Alexandre VII Chigi et dans la quatrième année de légation du cardinal Girolamo Farnèse, la suite d’écus placée au-dessus de l’arche est encore plus révélatrice du \teiHi[rend={italic}]{status} de la ville (fig. 84) : au sommet apparaît celui du pape, et sur les côtés, plus bas, figurent ceux du légat et du cardinal neveu Flavio Chigi (peut-être ici représenté comme préfet de la congrégation des confins et surintendant général de l’État ecclésiastique), pour finir par ceux de la famille sénatoriale Bonfioli et de la commune de Bologne. Particulièrement significative est la présence de l’écu des Bonfioli, qui semble renvoyer à la charge de gonfalonier assumée par Lelio Bonfioli dans le cinquième bimestre de cette année 1661\teiNote[n={3},place={foot},xmlid={ftn44-3}]{\teiP[xmlid={p-232}]{On remarquera que le dernier quartier des armes papales oblitère un quartier Della Rovere (concédés, comme on sait, par Jules II à Agostino Chigi), une croix lui a été substituée. Une occurrence singulière que je n’ai pas le loisir d’élucider ici. }}.}
            \teiP[xmlid={p3-7}]{La cohabitation des pouvoirs impliqués dans le gouvernement municipal est ici représentée dans un monument visible de tous. Mais dans les intérieurs bolonais et dans des documents destinés à une circulation restreinte, elle s’exprime de manières plus complexes, qui proviennent de la volonté d’autopromotion de l’une ou de l’autre des composantes politiques et laissent parfois dans l’ombre justement l’écu du pontife régnant. Les commandes artistiques du \teiHi[rend={small-caps}]{xvi}\teiHi[rend={sup}]{e} et du \teiHi[rend={small-caps}]{xvii}\teiHi[rend={sup}]{e} siècle constituent un témoignage particulièrement intéressant de la volonté d’autoreprésentation des trois pouvoirs entremêlés mais distincts : les Anciens, les familles sénatoriales et le légat. Laissant de côté les grandes familles, sur lesquelles les études se sont concentrées depuis longtemps, cette contribution se propose de mettre en lumière le caractère composite de la situation bolonaise prenant en considération la magistrature des Anciens et le légat, par une étude sur un nombre limité d’années et sur certaines commandes spécifiques dans lesquelles, d’après moi, se révèle une utilisation de l’héraldique dynamique et sémantiquement assez riche. Il s’agit de deux entreprises décoratives importantes dues à l’initiative des cardinaux légats successifs Girolamo Farnèse et Pietro Vidoni, et de certaines miniatures appartenant à la série connue sous le nom d’\teiHi[rend={italic}]{Insignia degli Anziani consoli} \lbrack{}Insignes des Anciens, les consuls\rbrack{}. Les années considérées vont de 1658 à 1665 (1658-1662, la légation Farnèse, puis celle de Vidoni), elles correspondent à une bonne partie du pontificat d’Alexandre VII Chigi, qui, comme on sait, régna de 1655 à 1667.}
            \begin{teiDiv}[type={section1},xmlid={div01-5}]

              \teiHead[xmlid={girolamo-farnese-legat-a-bologne-aeterna-hic-lilia-semper-erunt-001}]{Girolamo Farnèse, légat à Bologne : «\teiHi[rend={italic}]{ Aeterna hic LILIA semper erunt }»}
              \teiP[xmlid={p4-7}]{Les \teiHi[rend={italic}]{Insignia degli Anziani consoli} se composent de centaines de miniatures sur parchemin commandées tous les deux mois de 1530 à 1796 pour célébrer l’intronisation d’un nouveau gonfalonier de justice ou d’un Ancien, aujourd’hui conservées en seize volumes dans l’archive d’État de Bologne\teiNote[n={4},place={foot},xmlid={ftn45-3}]{\teiP[xmlid={p-233}]{Giuseppe Plessi (dir.), \teiHi[rend={italic}]{Le Insignia degli Anziani del comune dal 1530 al 1796: catalogo-inventario}, Bologne, ASBO, 1954. Les miniatures relatives aux années du pontificat d’Alexandre VII sont conservées dans le volume VIII. En réalité tous les bimestres ne sont pas représentés : cela provient probablement de l’absence de commande de la part du gonfalonier.}}. L’héraldique y domine, car les nouveaux magistrats se qualifient et se présentent dans ces images principalement à travers leurs armoiries familiales, ici accompagnées de la liste de leurs noms. La miniature relative au premier bimestre de 1658, année où Girolamo Farnèse fut envoyé comme légat à Bologne, illustre bien cette typologie particulière\teiNote[n={5},place={foot},xmlid={ftn46-3}]{\teiP[xmlid={p-234}]{ASBO, \teiHi[rend={italic}]{Insignia}, vol. VIII, p. 28.}} (fig. 85). Sur un ciel de pur azur traversé par l’écliptique, à gauche est figurée Felsina, personnification de Bologne, en train de montrer la ville au pied des collines, alors qu’à droite paraît une sorte de cippe monumental couronné de l’ombelle papale et orné de neuf écus des membres du conseil dont les noms sont reportés dans l’inscription sous-jacente. En recourant à l’astrologie et au mythe, l’image célèbre la venue, en 1658, d’une sorte de nouvel âge d’or, comme le déclare avec emphase l’inscription :}
              \begin{teiCit}[xmlid={cit01-3}]

                \begin{teiQuote}
\teiHi[rend={italic}]{Anno Domini 1658, culminantia Saturni astra civilem preduxere foelicitatem. Pan choreas duxit Ganimedes vina profudit cum summa innocentia et sobrietate. Primum hoc bimestre de aureo saeculo, reservatum nobis videtur. }\lbrack{}…\rbrack{}\teiHi[rend={italic}]{ Rempublicam curavere tunc temporis cum Em.}\teiHi[rend={italic sup}]{mo}\teiHi[rend={italic}]{ \& R.}\teiHi[rend={italic sup}]{mo}\teiHi[rend={italic}]{ Ioanne Hyeronimo Card.}\teiHi[rend={italic sup}]{li}\teiHi[rend={italic}]{ Lomellino, Ill.}\teiHi[rend={italic sup}]{mi}\teiHi[rend={italic}]{ DD.Berlingeus Gypsisius Vexillifer Iustitiae … }
\end{teiQuote}
              
\end{teiCit}
              \teiP[xmlid={p5-7}]{L’épigraphe explicite la signification politique de l’image, faisant référence à une ère de félicité publique, sous la protection de Saturne, avec la faveur de Pan (qui conduit les danses) et de Ganymède qui verse le vin « avec une innocente sobriété ». Le deux figures mythologiques, représentées sur la bande céleste, font allusion à la félicité reconquise (du corps comme de l’esprit), que la suite de l’inscription attribue au mérite des excellents gouvernants. Et si l’inscription définit Bologne « \teiHi[rend={italic}]{Rempublicam} » (république), outre les magistrats municipaux, soit le gonfalonier Berlingero Gessi et les huit Anciens, celle-ci implique aussi le cardinal légat Giovanni Gerolamo Lomellini et le vice-légat Gasparo Lascaris. Mais elle ne fait pas mention du pontife. Si nous passons à l’image, dont l’impact est bien plus fort et immédiat que celui du texte, nous observons que les armoiries des trois autorités ecclésiastiques n’y paraissent pas. Nous avons là un des nombreux indices de ce qui semble une stratégie précise de communication des \teiHi[rend={italic}]{Insignia}, tout du moins au \teiHi[rend={small-caps}]{xvii}\teiHi[rend={sup}]{e }siècle pris en considération par la recherche : la tendance systématique à occulter le pouvoir ecclésiastique, probablement à mettre sur le compte de la fonction de propagande des miniatures, orchestrée par les gonfaloniers en charge qui en étaient les commanditaires. Celle-ci représente un élément de continuité au cours du siècle et rend ainsi particulièrement pertinent le changement qui s’observe à partir de l’année 1659 et pour quelque temps, en coïncidence avec la légation de Girolamo Farnèse.}
              \teiP[xmlid={p6-7}]{Le changement commence avec la miniature du cinquième bimestre, Bartolomeo Lambertini étant gonfalonier (fig. 86) : une image allégorique rigidement symétrique, sur les côtés de laquelle des rameaux d’olivier et des palmes évoquent les valeurs habituelles de ces plantes, paix et gloire. Le long de l’axe central, parmi des anges en vol, sont figurés la tiare papale et le chapeau cardinalice, entourés des symboles du pouvoir ecclésiastique d’un côté (crosse pastorale et mitre épiscopale), et politique de l’autre (sceptres et couronnes). L’image dans son ensemble et son centre est pourtant dominée par une pluie de lys de couleur indigo, formant aussi un cadre et une protection pour les écus du gonfalonier lui-même et des Anciens, qui ornent une sorte d’autel sur lequel on lit :}
              \begin{teiCit}[xmlid={cit02-3}]

                \begin{teiQuote}
Que personne dorénavant n’ose définir les fleurs comme caduques ; ici les LYS seront toujours éternels\teiNote[n={6},place={foot},xmlid={ftn47-3}]{\teiP[xmlid={p-235}]{« \teiHi[rend={italic}]{Ne quisquam posthac flores dixisse caducos / Audeat Aeterna hic LILIA semper erunt} ».}}.
\end{teiQuote}
              
\end{teiCit}
              \teiP[xmlid={p7-7}]{L’image et l’inscription célèbrent les lys de l’écu Farnèse, donc le cardinal légat, et elles suggèrent la volonté du gonfalonier ainsi que des Anciens de se confier à sa protection\teiNote[n={7},place={foot},xmlid={ftn48-4}]{\teiP[xmlid={p-236}]{La miniature est signée en marge en bas à droite : « \teiHi[rend={italic}]{D. Gio. And. Abanti scr(isse) e min(iò)} ». }}. À partir de ce moment, la référence au cardinal devient fréquente dans les \teiHi[rend={italic}]{Insignia}, subvertissant l’habituelle absence ou marginalité des écus du pontife régnant et du légat apostolique notées précédemment.}
              \teiP[xmlid={p8-7}]{Entre le cinquième bimestre de 1659 que nous venons de considérer et la fin de l’année 1661, alors que le cardinal Farnèse s’apprête à terminer son mandat bolonais, les miniatures exécutées à l’occasion des intronisations du gonfalonier et des Anciens lors de bimestres successifs sont au nombre de sept\teiNote[n={8},place={foot},xmlid={ftn49-4}]{\teiP[xmlid={p-237}]{Comme on l’a dit plus haut, tous les bimestres ne sont pas représentés. Dans l’intervalle de temps considéré manquent le sixième bimestre de 1659, le premier et deuxième de 1660, le premier, deuxième et cinquième de 1661.}} : non moins de cinq d’entre elles sont dédiées au cardinal. Du reste, élevé à la pourpre par Alexandre VII le 29 mai 1658 et immédiatement envoyé à Bologne, le cardinal s’était révélé un homme très actif et un administrateur appliqué\teiNote[n={9},place={foot},xmlid={ftn50-4}]{\teiP[xmlid={p-238}]{Stefano Andretta, « Farnese, Girolamo »,dans \teiHi[rend={italic}]{DBI}, Rome, Istituto della Enciclopedia italiana, 1995, vol. 45, p. 95-98, \teiRef[target={https://www.treccani.it/enciclopedia/girolamo-farnese\_Dizionario-Biografico/}]{https://www.treccani.it/enciclopedia/girolamo-farnese\_Dizionario-Biografico/} ; Alfeo Giacomelli, « La storia di Bologna dal 1650 al 1796: un racconto e una cronologia », dans Adriano Prosperi (dir.), \teiHi[rend={italic}]{Bologna nell’età moderna (secoli XVI-XVIII). }Vol. I.\teiHi[rend={italic}]{ Istituzioni, forme del potere, economia e società}, Bologne, Bologna University Press, 2009, p. 61-198, ici p. 68.}}. Il devait laisser dans la ville des traces importantes et un excellent souvenir. Sa magnificence et sa profonde implication dans le gouvernement municipal semblent avoir conduit les divers gonfaloniers à lui témoigner une reconnaissance particulière\teiNote[n={10},place={foot},xmlid={ftn51-4}]{\teiP[xmlid={p-239}]{Ce furent Enrico Hercolani (troisième bimestre 1660) ; Marc’Antonio Ranuzzi (quatrième bimestre 1660) ; Tommaso Campeggi (sixième bimestre 1660) ; Carlo Francesco Caprara (troisième bimestre 1661) ; et Nicola Calderini (quatrième bimestre 1661). Bien que ce ne soit pas l’objectif de cette étude, il vaut la peine de noter que chaque miniature adopte des solutions figuratives et une structure narrative autonome, et cette variété fait du recueil un objet d’étude passionnant et de grand potentiel.}}. Parmi les miniatures, la plus intéressante pour notre recherche est celle du sixième bimestre 1660, qui eut comme gonfalonier de justice le marquis Tommaso Campeggi\teiNote[n={11},place={foot},xmlid={ftn52-4}]{\teiP[xmlid={p-240}]{ASBO, \teiHi[rend={italic}]{Insignia}, vol. VIII, p. 46. Signé en marge en bas à droite, « D. Gio. And. Abanti ».}} (fig. 87). On y voit dans une stricte symétrie la statue de métal d’Alexandre VII qui fut offerte au pontife par le cardinal, comme on le lit sur le socle de l’effigie\teiNote[n={12},place={foot},xmlid={ftn53-4}]{\teiP[xmlid={p-241}]{« \teiHi[rend={italic}]{ALEX. VII. / PONT. MAX. / BENEFACTORI / HIERONIMUS /CARD. FARNES. / AN. MDCLIX}. ». Le nom du cardinal est réitéré par deux plus petites inscriptions sur des cheminées à droite et à gauche : « \teiHi[rend={italic}]{Farnesius Legatus} », « \teiHi[rend={italic}]{Hieronymus Cardin(alis)} ».}}. L’image célèbre donc les deux personnages et suggère que Tommaso Campeggi leur était particulièrement proche, ou entendait obtenir leur faveur : les écus des huit Anciens sont disposés comme s’ils constituaient les métopes de la frise dorique qui se déploie sur le mur derrière la statue, alors que celui de Campeggi, de grande dimension, paraît au premier plan. Il est dévotement présenté à la statue du pape par deux figures féminines que leurs attributs permettent d’identifier comme la Justice et la Paix. De l’autre côté Felsina offre une maquette de Bologne à la statue, accompagnée d’une jeune femme couronnée de fleurs, portant caducée et cornucopie, qui correspond à la Félicité publique\teiNote[n={13},place={foot},xmlid={ftn54-4}]{\teiP[xmlid={p-242}]{Comme une bonne partie des personnifications utilisées dans la décoration de l’époque, qui constituent un répertoire de connaissances communes, et souvent tiré de Cesare Ripa, qui s’inspire des monnaies romaines. Elle se trouve là en image dès la première édition illustrée : Cesare Ripa, \teiHi[rend={italic}]{Iconologia}, In Roma, Appresso Lepido Facji, 1603, p. 154-155.}}. Ces quatre personnifications, que les contemporains devaient facilement reconnaître, font référence aux vertus inspiratrices de toute action de gouvernement, comme au but fondamental du bon gouvernement, c’est-à-dire la félicité du peuple\teiNote[n={14},place={foot},xmlid={ftn55-4}]{\teiP[xmlid={p-243}]{Nous aurons l’occasion de revenir sur la pertinence politique des personnifications de la Justice et de la Paix.}}. La composition rhétorique de l’image, où l’écu Campeggi est offert en leur présence à la statue d’Alexandre VII, don du cardinal Farnèse, exprime l’hommage du gonfalonier aux deux prélats, et souligne sa volonté d’en partager l’action politique. Un fait peu fréquent dans la Bologne de l’époque moderne, où les \teiHi[rend={italic}]{Insignia}, omettant la plupart du temps toute référence au pape et au légat en fonction, revendiquent l’autonomie du gonfalonier et des Anciens.}
              \teiP[xmlid={p9-6}]{Comme on le voit, ces miniatures offrent un témoignage visuel particulièrement éloquent si on les étudie de manière comparative, elles révèlent comment l’héraldique procurait à la vie politique un instrument de communication assez souple, grâce auquel il est possible de percevoir des aspects importants de ses dynamiques.}
              \teiP[xmlid={p10-6}]{La statue figurée dans la miniature, qui, heureusement, a été conservée, présente elle aussi son intérêt. Cette image souligne avec une certaine subtilité et sur un mode original la dévotion du cardinal envers le pontife, que l’inscription du socle définit comme « bienfaiteur\teiNote[n={15},place={foot},xmlid={ftn56-4}]{\teiP[xmlid={p-244}]{Quoique le socle actuel ne soit pas l’original, il en reste le souvenir et de son inscription dans la miniature et surtout dans la maquette préparatoire de la statue conservé au Victoria \& Albert Museum de Londres (inv. A.60 : 1, 2-1921), pour laquelle on renvoie à la notice en ligne, avec bibliographie et fiche technique (\teiRef[target={http://collections.vam.ac.uk}]{http://collections.vam.ac.uk}.). Il faut aussi noter les inscriptions latérales présentes sur le socle de la maquette, qui fournissent des informations importantes, qu’il n’est pas possible de commenter ici. Sur le côté droit on lit : « \teiHi[rend={italic}]{CURATORE / ILL.mo DOMno / SENATre /MARCHne / BAYULIVO / FERDINANDO / DE COSPIIS / EX MANDATis/ EMmi /LEGATI} » ; sur le côté gauche : « \teiHi[rend={italic}]{PROTOTYPUM / SIMULACRI / STIVATI IN / AULA / HELVETICORUM / A / DORASTANTE / OSIO FLORENTINO CONSTRUCTUM} ».}} ». Œuvre du Florentin Dorastante Maria d’Osio, ou Osio\teiNote[n={16},place={foot},xmlid={ftn57-4}]{\teiP[xmlid={p-245}]{Le sculpteur est précisément florentin, non siennois, comme on le lit curieusement dans Carla Bernardini, « Bologna 1660. Una committenza Farnese per papa Alessandro VII », dans Andrea Bacchi et Luca M. Barbero (dir.), \teiHi[rend={italic}]{Studi in onore di Stefano Tumidei}, Venise-Bologne, Fondazione Cini-Fondazione Zeri, 2016, p. 279-287. D’autant plus que la signature même de l’artiste, gravée sous le pied gauche de la statue, dit sans aucun doute possible : « \teiHi[rend={italic}]{DORASTANTI M.}\teiHi[rend={italic sup}]{I}\teiHi[rend={italic}]{ DEOSIIS FLO(RENTINI) AD MDCXXXXXX FECIT} » (voir la fig. 4 de l’étude précédemment citée). }}, c’est une effigie plutôt singulière, en feuille de cuivre sur âme de bois, qui reprend une typologie assez ancienne de portrait pontifical, dont on conserve à Bologne un exemple important avec le \teiHi[rend={italic}]{Boniface VIII} daté de 1301\teiNote[n={17},place={foot},xmlid={ftn58-4}]{\teiP[xmlid={p-246}]{Sur la statue et l’artiste voir Simona Moretti, « Manno di Bandino », dans \teiHi[rend={italic}]{DBI}, \teiHi[rend={italic}]{op. cit}., 2007, vol. 69, p. 109-110, \teiRef[target={https://www.treccani.it/enciclopedia/manno-di-bandino\_Dizionario-Biografico/}]{https://www.treccani.it/enciclopedia/manno-di-bandino\_Dizionario-Biografico/}, avec bibliographie, et Raffaella Pini, « La statua di Bonifacio VIII, Manno da Siena e gli orefici a Bologna », dans Maria Andaloro (dir.), \teiHi[rend={italic}]{Le culture di Bonifacio VIII}, Rome, Istituto Storico Italiano per il Medioevo, 2006, p. 231-240. On peut se demander si cette commande originale de Farnèse, très sophistiquée dans sa rareté, est une reprise consciente d’un modèle antique (dans ce cas médiéval) destinée à satisfaire aussi bien les attentes du destinataire, le très cultivé pontife Chigi, que l’intérêt diffus à partir de la Contre-Réforme pour remettre en usage les pratiques de l’Église des origines ou médiévale. }}. Le cardinal fit placer la statue dans la salle des Suisses, atrium de l’appartement du légat\teiNote[n={18},place={foot},xmlid={ftn59-4}]{\teiP[xmlid={p-247}]{Cela est rappelé par une des inscriptions du socle, voir n. 16. }}. À cette même époque il avait fait décorer de fresques la vaste pièce qui sert de vestibule solennel à l’appartement, conçue comme l’équivalent bolonais de la Sala Regia au palais du Vatican, ainsi que les études de Hans Hubert l’ont démontré\teiNote[n={19},place={foot},xmlid={ftn60-4}]{\teiP[xmlid={p-248}]{Voir n. 3. Cette salle, aujourd’hui connue comme la salle Farnèse, est aussi le vestibule de la chapelle du palais, décorée par Prospero Fontana avec des fresques mariales lors de la légation de Charles Borromée, en 1562.}}. À la lumière de ce décor on peut aussi et surtout considérer la statue en l’honneur d’Alexandre VII comme une sorte de compensation de la part du cardinal, les fresques étant entièrement consacrées à l’exaltation des Farnèse et de leur pape Paul III. L’œuvre fut confiée à Carlo Cignani, le meilleur héritier vivant de la tradition carrachesque, après la mort de son maître Francesco Albani en 1660. L’équipe de Cignani conçut une somptueuse structure décorative, mettant en œuvre des solutions caractéristiques du perspectivisme bolonais \lbrack{}\teiHi[rend={italic}]{quadratura bolognese}\rbrack{} destinées à exalter la monumentalité de l’ensemble\teiNote[n={20},place={foot},xmlid={ftn61-4}]{\teiP[xmlid={p-249}]{Les peintres actifs dans la salle furent Giovanni Maria Bibiena, Lorenzo Pasinelli, Emilio Taruffi, Luigi Scaramuccia, Girolamo Bonini, Antonio Catalani, Bartolomeo Morelli et Francesco Quaini.}} : elle encadre quelques scènes soigneusement sélectionnées pour synthétiser l’histoire de Bologne en lien avec l’Église romaine. Dans les fresques se mêlent en particulier histoire municipale, exaltation du rôle des légats, hommage à la famille du commanditaire et de ses membres illustres\teiNote[n={21},place={foot},xmlid={ftn62-4}]{\teiP[xmlid={p-250}]{Sur la décoration de la salle, voir Beatrice Buscaroli Fabbri, « Il Cardinal Farnese e la sua sala. Un ciclo di affreschi per la famiglia e la città », dans Camilla Bottino (dir.), \teiHi[rend={italic}]{Il palazzo comunale di Bologna}, \teiHi[rend={italic}]{op. cit}., p. 99-110. Sur les autres pièces décorées et l’appartement dans son ensemble, voir en particulier Carla Bernardini, « Per una storia dell’Appartamento del Legato, sede delle Collezioni Comunali d’Arte di Palazzo d’Accursio. Una ricerca in prospettiva museografica », \teiHi[rend={italic}]{Schede Umanistiche}, n\teiHi[rend={sup}]{o} 3, 1989, p. 81-92 ; et aussi Davide Righini, « Iconografie celebrative. Gli emblemi e i miti raffigurati nelle sale di rappresentanza del Palazzo Pubblico di Bologna (secoli XV-XVIII) », \teiHi[rend={italic}]{Arte a Bologna. Bollettino dei Musei Civici d’Arte Antica}, n\teiHi[rend={sup}]{o} 7-8, 2010-2011, p. 128-140. }}. Parmi les événements cruciaux dont la ville avait été le théâtre comme le \teiHi[rend={italic}]{Couronnement de Charles V de la part de Clément VII }(1530) et la \teiHi[rend={italic}]{Rencontre de François I}\teiHi[rend={italic sup}]{er }\teiHi[rend={italic}]{de France avec le pape Léon X }(1515), où le souverain avait accepté de pratiquer l’attouchement miraculeux des malades, aussi représenté, on trouve l’\teiHi[rend={italic}]{Entrée de Paul III Farnèse à Bologne }(1543) (fig. 88). Le pontife est bien reconnaissable au premier plan, chevauchant devant le Palais public à côté d’un imaginaire monument équestre à Pietro Farnèse, ancêtre qui, en 1361, à la tête des troupes pontificales, avait vaincu les Visconti, reconquérant Bologne. Ce passage de Paul III dans la ville était somme toute un épisode historique assez marginal que le programme décoratif élève au rang de moment emblématique, l’associant au souvenir du condottière Pietro par le moyen d’un autre tour de force, aucune statue ne lui ayant jamais été consacrée. Cette iconographie dynastique se trouve conçue comme appartenant à un programme visuel destiné à encadrer la statue en marbre, aujourd’hui perdue, de Paul III Farnèse que Giovanni Zacchi da Volterra avait exécutée pour la salle justement en 1543\teiNote[n={22},place={foot},xmlid={ftn63-4}]{\teiP[xmlid={p-251}]{Elle était placée sur le côté court entre les fenêtres, sous un baldaquin peint qui encadre maintenant l’\teiHi[rend={italic}]{Alexandre VII Chigi }de Dosio.}}. Cela est rappelé et souligné par l’inscription peinte sur la porte d’entrée de l’appartement du légat, où le cardinal Girolamo explique avoir ennobli la salle à la suite des Bolonais qui avaient placé là au siècle précédent la statue de son ancêtre le pontife\teiNote[n={23},place={foot},xmlid={ftn64-4}]{\teiP[xmlid={p-252}]{« \teiHi[rend={italic}]{AULAM HANC UBI BONONIENSIUM INCLYTA FIDES PAULO III STATUAM OLIM POSUIT IN AUGUSTIOREM FORMAM EXORNANDAM CURAVIT HYERONIMUS CARDINAL FARNES . LEG. A.D. MDCLX} ».}}.}
              \teiP[xmlid={p11-6}]{Mais ce n’est pas tout, car l’héraldique aussi, avec son langage immédiat, est sollicitée pour exalter la lignée, sur les murs comme sur le plafond. Les lys Farnèse sont figurés dans des écus au-dessus des grandes scènes historiques qui illustrent les parois (fig. 88), ils paraissaient aussi aux caissons d’origine du plafond, aujourd’hui disparu mais décrit dans une source contemporaine qui dit :}
              \begin{teiCit}[xmlid={cit03-2}]

                \begin{teiQuote}
en son milieu \lbrack{}du plafond de bois\rbrack{}, dans un écu bien grand et majestueux, doré avec profusion, se voient les armes révérées des Chigi, de notre Seigneur le pape Alexandre VII et dans les embrasures qui sont toutes en partie couvertes d’or se trouvent les lys d’azur qui font un beau concert\teiNote[n={24},place={foot},xmlid={ftn65-4}]{\teiP[xmlid={p-253}]{« \teiHi[rend={italic}]{nel suo mezzo} \lbrack{}du plafond de bois\rbrack{}, \teiHi[rend={italic}]{in ben grande e maestoso scudo, tutto profusamente indorato, vi si vede la riverita arma Ghisia, di Nostro Signore papa Alessandro VII, e nei vani, che sono partitamente tutti posti a oro vi sono gigli azzurri che fanno un bellissimo concerto} ». L’extrait fait partie d’une relation anonyme relative aux embellissements dus à la munificence du cardinal Girolamo Farnèse, adressée en 1661 à Monseigneur Girolamo Boncompagni, archevêque de la ville. Le texte entier figure dans Gaetano Giordani, \teiHi[rend={italic}]{Pitture della sala Farnese}, Bologne, Tipografia Guidi all’Ancora, 1845, p. 13 et 14. La lecture de cette source est intéressante aussi pour la description des fresques et du placement de la statue de d’Osio figurant \teiHi[rend={italic}]{Alexandre VII }dans la salle des Suisses.}}.
\end{teiQuote}
              
\end{teiCit}
              \teiP[xmlid={p12-6}]{Quoique placées en position éminente, les armes des Chigi étaient entourées de caissons ornés des lys Farnèse. On comprend comment le poids visuel fut accordé à l’héraldique du légat au détriment de celle du pontife.}
            
\end{teiDiv}
            \begin{teiDiv}[type={section1},xmlid={div02-5}]

              \teiHead[xmlid={micet-in-vertice-pietro-vidoni-et-la-celebration-dalexandre-vii-chigi-001}]{« \teiHi[rend={italic}]{Micet in vertice }» : Pietro Vidoni et la célébration d’Alexandre VII Chigi}
              \teiP[xmlid={p13-6}]{La légation de Girolamo Farnèse, qui fut en poste dans la ville de 1658 à 1662, se conclut par une apothéose d’hommages. Suivit le mandat du Crémonais Pietro Vidoni (1610-1681), qui avait longtemps été apprécié en tant qu’évêque de Lodi. Comme son prédécesseur, Vidoni avait été élevé à la pourpre par Alexandre VII, quoique plus tard, dans le concistoire du 5 avril 1660\teiNote[n={25},place={foot},xmlid={ftn66-4}]{\teiP[xmlid={p-254}]{Le diocèse lui avait été confié en 1644, mais de fait il s’était absenté de la cité lombarde depuis 1652, année où il avait été appelé à la nonciature près le roi de Pologne Jean Casimir. L’habileté diplomatique dont Vidoni fit preuve en tant que nonce lui valut sa nomination comme protecteur du royaume de Pologne près le Saint-Siège et le fort soutien du roi Jean pour son élévation à la pourpre. Voir à ce propos Gaetano Moroni, \teiHi[rend={italic}]{Dizionario di erudizione storico-ecclesiastica da S. Pietro sino ai nostri giorni}, in Venezia, dalla Tipografia Emiliana, 1840-1861, vol. XCIX, p. 246-247 ; Patrice Gauchat (dir.), \teiHi[rend={italic}]{Hierarchia Catholica Medii et Recentioris Aevi}, Regensburg, Sumptibus et typis Librariae Monasterii, 1935, vol. V., \teiHi[rend={italic}]{1592-1667}, p. 34. Giampiero Brunelli, « Vidoni, Pietro senior », \teiHi[rend={italic}]{DBI}, \teiHi[rend={italic}]{op. cit}, 2020, vol. 99, p. 204-207, \teiRef[target={https://www.treccani.it/enciclopedia/pietro-senior-vidoni\_Dizionario-Biografico/}]{https://www.treccani.it/enciclopedia/pietro-senior-vidoni\_Dizionario-Biografico/}.}}. Cela est rappelé, comme d’habitude dans le recueil, grâce à un portrait gravé par Albert Clouwet sur un dessin de Pietro Martire Neri dans la série des \teiHi[rend={italic}]{Effigies cardinalium nunc viventium, }publiée à Rome par Giovanni Giacomo de Rossi. Il porte les armes du pontife auteur de l’élévation à la pourpre en haut à gauche, celles du cardinal à droite\teiNote[n={26},place={foot},xmlid={ftn67-4}]{\teiP[xmlid={p-255}]{L’inscription dit : « \teiHi[rend={italic}]{Petrus Tit. S. Callisti S. R. E. Presb. Card. / Vidonus. Episcopus Laudensis / Cremonensis V. Aprilis. MDCLX} ».}} (fig. 89). Il vaut la peine de rappeler le blasonnement de l’écu des Vidoni, qui est assez peu connu :}
              \begin{teiCit}[xmlid={cit04-2}]

                \begin{teiQuote}
D’argent, à une tour de gueules ajourée de sable. À une vigne de sinople feuillée du même et fruitée au naturel, naissant de la porte et accolant la tour, entrant et sortant par les fenêtres et couronnant la cime\teiNote[n={27},place={foot},xmlid={ftn68-4}]{\teiP[xmlid={p-256}]{« \teiHi[rend={italic}]{D’argento, la torre di rosso aperta e finestrata di nero. La vite sradicata, pampinosa di verde coll’uva, nascente dalla porta ed accollante la torre, entrando e uscendo per le finestre coronando la cima} ». Je remercie Gian Luca Tusini pour son aide.}}.
\end{teiQuote}
              
\end{teiCit}
              \teiP[xmlid={p14-5}]{Pour en rester aux \teiHi[rend={italic}]{Insignia}, l’avènement de Vidoni ne passa pas inaperçu. La miniature du troisième bimestre de l’année, qui correspond au moment de son arrivée dans la ville, est très éclairante\teiNote[n={28},place={foot},xmlid={ftn69-4}]{\teiP[xmlid={p-257}]{ASBO, \teiHi[rend={italic}]{Insignia degli Anziani consoli}, vol. VIII, p. 56. Signée, comme une bonne partie des précédentes, « \teiHi[rend={italic}]{D. Gio. And. Abanti scr. e min. }».}} (fig. 90). Dans un pré lumineux ceint d’une balustrade faisant face à un ample paysage, est campée la tour de l’écu Vidoni, surmontée du chapeau cardinalice et de l’ombelle, sur laquelle est fixé l’écu du gonfalonier de justice, Giovan Francesco Isolani, alors que les armes des huit Anciens sont prêtes à y prendre leur place après avoir été consacrées par Apollon. Cette miniature aussi montre un usage dynamique de l’héraldique : on peut dire qu’elle illustre les armes Vidoni, complétées par les insignes propres du légat (ombelle, clefs en sautoir, chapeau), et en même temps elle les anime, pour ainsi dire, rendant la tour géante, élément iconique, sorti de l’usuel écu et utilisé comme une architecture symbolique autant que physique, sur laquelle peuvent s’appuyer les écus héraldiques des Anciens. La Justice garde la tour, à côté de la vigne vidonienne surgie de terre, alors qu’à gauche la Paix a déposé le rameau d’olivier afin d’offrir à Apollon l’écu Bentivoglio pour qu’il le couronne, et l’Abondance (qui semble reconnaissable grâce à la corne d’abondance et à la couronne de fleurs) s’apprête à offrir ceux des sept autres Anciens. En marge, en bas à droite, une divinité fluviale, probablement le Reno (fleuve de Bologne, identifiant traditionnellement cette ville), fait jaillir l’eau d’une amphore sur laquelle leurs noms sont inscrits. Un cartel proclame l’intention de la ville de se placer sous la protection de Vidoni par une métaphore : « \teiHi[rend={italic}]{sub umbra alarum tuarum }», à l’ombre de tes ailes, qui pourraient être les houppes du chapeau ou les rameaux de la vigne.}
              \teiP[xmlid={p15-5}]{Une image aussi emphatique conduirait à penser que Vidoni avait montré depuis son entrée en fonction une attitude de magnificence, et peut-être instauré un lien de protection avec les Anciens et le Sénat. Mais la sensation qui ressort en parcourant les \teiHi[rend={italic}]{Insignia }suivants est autre. On a plutôt l’impression qu’au début de sa légation les gonfaloniers successifs ont exprimé à travers les miniatures un augure de continuation du mandat Farnèse, et qu’ensuite cette attente a été démentie par les faits : dans les miniatures suivantes la présence de l’écu Vidoni se réduit peu à peu pour disparaître complètement\teiNote[n={29},place={foot},xmlid={ftn70-2}]{\teiP[xmlid={p-258}]{On fait ici référence aux miniatures du quatrième bimestre 1662, p. 58 (gonfalonier Achille Volta) et du cinquième bimestre 1662, p. 60 (gonfalonier Gerolamo Capacelli Albergati). Il disparaît à partir du sixième et dernier bimestre de l’année (gonfalonier Odoardo Pepoli). L’écu réapparaît, mais comme élément du décor, inséré dans une architecture éphémère, dans la miniature du quatrième bimestre 1663 (gonfalonier Angelo Michele Guastavillani), p. 68.}}.}
              \teiP[xmlid={p16-5}]{Et pourtant Vidoni, comme Girolamo Farnèse avant lui, laissa durant les quelques années de sa légation un témoignage important de son rôle de commanditaire dans l’appartement du légat. Il s’agit de la décoration de la voûte de la galerie qui porte son nom, la Vidoniana, mais qui mériterait plutôt d’être rebaptisée galerie Chigiana, si l’on tient compte de la signification de ses fresques (fig. 91). Contrairement à celles de la sala Regia voulues par son prédécesseur, celles-ci sont entièrement consacrées à Alexandre VII : on peut en déduire un lien plus étroit de Vidoni avec le pontife, probablement la gratitude pour son élévation à la pourpre et peut-être aussi pour sa nomination comme légat dans la ville la plus riche des États pontificaux après Rome.}
              \teiP[xmlid={p17-5}]{Les travaux datent des derniers mois de la légation, vers 1665, et la documentation d’archive montre que Vidoni ne fut pas seulement responsable de la décoration mais aussi de la création de la galerie elle-même\teiNote[n={30},place={foot},xmlid={ftn71-2}]{\teiP[xmlid={p-259}]{À propos du choix de cette typologie, qui eut un grand succès à cette époque, et pour une étude plus complète de la décoration, voir Sonia Cavicchioli, « Il cardinale Pietro Vidoni e l’omaggio ad Alessandro VII Chigi nell’appartamento del legato a Bologna (1665). Una decorazione politica », \teiHi[rend={italic}]{Studi di Storia dell’Arte}, n\teiHi[rend={sup}]{o} 29, 2018, p. 167-178.}}, située dans une partie de l’appartement dont l’espace et le choix de la disposition étaient limités, la résidence ayant déjà été largement structurée dans les siècles et décennies précédents. Les fresques furent confiées à deux maîtres de l’école de peinture locale, qui au \teiHi[rend={small-caps}]{xvii}\teiHi[rend={sup}]{e} siècle était peut-être la plus appréciée en Europe, spécialement pour la grande décoration : il s’agit des perspectivistes (\teiHi[rend={italic}]{quadraturisti}) Giovanni Battista Caccioli (1623-1675) et Domenico Santi, dit le Mengazzino (1621-1694), respectivement l’un peintre de figures et l’autre spécialiste d’architecture peinte, élèves le premier de Domenico Canuti, le second d’Agostino Mitelli\teiNote[n={31},place={foot},xmlid={ftn72-2}]{\teiP[xmlid={p-260}]{La bibliographie sur ces artistes est mince. Sur Caccioli, voir Luigi Crespi, \teiHi[rend={italic}]{Felsina pittrice}, In Roma, Nella Stamperia di Marco Pagliarini, 1769, p. 119-120 (en particulier p. 120 pour la référence à la galerie) et p. 57-59 (en particulier p. 59) qui la date de 1665. Ebria Feinblatt, « Giovan Battista Cavalcaselle: a Drawing identified », \teiHi[rend={italic}]{Master Drawings}, vol. V, n\teiHi[rend={sup}]{o} 2, 1967, p. 569-576 ; Clara Roli Guidetti, « Caccioli, Giovanni Battista », dans \teiHi[rend={italic}]{DBI}, \teiHi[rend={italic}]{op. cit}., 1973, vol. 16, p. 39-40, \teiRef[target={https://www.treccani.it/enciclopedia/giovanni-battista-caccioli\_Dizionario-Biografico}]{https://www.treccani.it/enciclopedia/giovanni-battista-caccioli\_Dizionario-Biografico}. Sur Santi, voir les références dans Giovan Pietro Zanotti, \teiHi[rend={italic}]{Storia dell’Accademia Clementina di Bologna}, \teiHi[rend={italic}]{Giuseppe Rolli}, In Bologna, Per Lelio Dalla Volpe, 1739, vol. I, p. 29, 270, 323 et 373. Et aussi Simonetta Stagni, \teiHi[rend={italic}]{Domenico Maria Canuti}, Rimini, Luisè, 1988, p. 136-139 et 149-151. En particulier sur leur travail dans l’« \teiHi[rend={italic}]{aula Vidonia} », Ebria Feinblatt,\teiHi[rend={italic}]{ Seventeenth-century Bolognese ceiling decorators}, Santa Barbara, Fithian Press, 1992, p. 161-162 et p. 163, n. 31. }}. La décoration rend hommage au pontife de manière sophistiquée, mêlant références héraldiques et emblématiques. La voûte est rythmée par quatre arcs transversaux imaginaires qui délimitent trois travées présentées comme des voûtes de pierre fortement abaissées, au centre desquelles est peinte une scène céleste, ceinte d’un cadre en feint bronze doré. L’entière superficie est peuplée de satyres et satyresses couronnés de pampres, et partout sont peints des sarments et des grappes, en alternance avec de plus communs festons de fruits et bouquets de fleurs. La présence de la vigne, qui sert discrètement de fond aux fresques, suggère une référence héraldique vu son rôle dans le blason du commanditaire.}
              \teiP[xmlid={p18-5}]{L’impression est confirmée par un petit poème en latin publié à Bologne en 1665, en parallèle à l’achèvement des fresques : \teiHi[rend={italic}]{Aula Vidonia}. Constitué de 186 hexamètres, il offre une sorte d’\teiHi[rend={italic}]{ekphrasis} de la décoration de la galerie, consacrée dès le titre au nom et au « \teiHi[rend={italic}]{numen} » du pape\teiNote[n={32},place={foot},xmlid={ftn73-2}]{\teiP[xmlid={p-261}]{La page de titre donne le titre complet : \teiHi[rend={italic}]{Aulam Vidoniam picturis exornatam Alexandri VII Pont. Opt. Max. Nomini et Numini consecratam Vincentius Maria Marescalchius hisce carminibus adumbrabat}, Bologna, Typis Ferronianis, 1665. « \teiHi[rend={italic}]{Numen} » peut être compris au sens de majesté, autorité religieuse. La référence à la vigne se lit aux vers 62-65.}}. C’est l’œuvre de Vincenzo Maria Marescalchi, personnalité d’un certain relief dans la Bologne d’alors : plusieurs fois Ancien, et membre de l’importante Académie des Gelati, dont il fut prince en 1652. En 1661, encore au temps de la légation Farnèse, il participa à un recueil collectif d’\teiHi[rend={italic}]{Ossequii poetici nell’ammirarsi in Bologna l’Heroica risoluzione della sereniss. Principessa Caterina Farnese di monacarsi}\teiNote[n={33},place={foot},xmlid={ftn74-2}]{\teiP[xmlid={p-262}]{\teiHi[rend={italic}]{Ossequii poetici nell’ammirarsi in Bologna l’Heroica risoluzione della sereniss. Principessa Caterina Farnese di monacarsi nel Ven. Monastero delle Madri carmelitane scalze di Santa Teresa in Parma}, In Bologna, Per Giacomo Monti, 1661.}}\teiHi[rend={italic}]{,} entreprise liée à la prise de voile de la fille du duc de Parme Odoardo I et de Marguerite de Médicis, advenue dans cette ville au monastère des carmélites, dont l’écho peut être compris en relation avec le séjour à Bologne du cardinal de la même famille. Le sonnet composé par Marescalchi pour l’occasion développe la métaphore du Christ « éternel jardinier » et souligne le renoncement de Catherine au monde, qui la portera à cultiver ses fleurs, soit les lys des armes Farnèse, tendant à la vertu (« Va, grande CATHERINE, quitte les parures de pourpre / Pour l’austère voile et les vastes salles pour la cellule monastique, / Distille les rayons de tes yeux en pures fontaines / Et tu verras chacune de tes fleurs \lbrack{}les lys des Farnèse\rbrack{} se changer en étoile », « \lbrack{}…\rbrack{}aujourd’hui transplantées\lbrack{}…\rbrack{}/ Dans le terrain de Thérèse l’AUGUSTE LYS\teiNote[n={34},place={foot},xmlid={ftn75-2}]{\teiP[xmlid={p-263}]{\teiHi[rend={italic}]{Ibid.}, p. 40 : « \teiHi[rend={italic}]{Vanne gran CATERINA, e gli ostri in lame / Ruvide muta, e le gran Sale in Cella, / Stilla da i raggi tuoi pure fontane, / E ogni tuo fior vedrai cangiare in Stella} »; « \lbrack{}…\rbrack{}\teiHi[rend={italic}]{oggi trapianti }\lbrack{}\teiHi[rend={italic}]{…}\rbrack{}\teiHi[rend={italic}]{/ Nel terren di Teresa il GIGLIO AUGUSTO} ».}} »). La gravure destinée à constituer le frontispice du recueil, due au Parmesan Domenico Fontana, se concentre aussi sur l’héraldique Farnèse. Dans un cartouche surmonté d’une couronne, l’écu Farnèse se trouve pour ainsi dire fixé par deux anges sur le plus grand écu carmélitain\teiNote[n={35},place={foot},xmlid={ftn76-2}]{\teiP[xmlid={p-264}]{Aux Farnèse revenait le titre, donc la couronne, d’Altesse Sérénissime. Comme on le sait, l’écu de l’ordre est constitué du mont Carmel surmonté de la croix, propre à l’ordre des carmélites déchaux, d’une étoile d’argent au pied du mont et de deux autres étoiles de gueules symétriques sur les côtés de la croix.}} (fig. 92). Pour souligner l’engagement total de la jeune Farnèse, qui avait choisi d’entrer en religion sous le nom de Thérèse, fondatrice de l’ordre, deux anges tiennent respectivement une flèche et un clou, allusion aux extases de la sainte espagnole, alors que la couronne ducale propre à l’écu Farnèse se trouve remplacée par celle d’épines. L’inscription rend la signification de l’image explicite :}
              \begin{teiCit}[xmlid={cit05-2}]

                \begin{teiQuote}
Thérèse, de lignée princière, une couronne d’épines te joint au Christ avec une flèche et un clou, et te place en ce lieu ; avec le dard et le clou le Christ se fait maître des parfums. Ainsi et pour toujours fleurissent les lys mêlés aux ronces\teiNote[n={36},place={foot},xmlid={ftn77-2}]{\teiP[xmlid={p-265}]{« \teiHi[rend={italic}]{Te iaculo iungit clavoq(ue) Teresia Christo, / Regalisq(ue) loco spinea serta locat. / Et iaculo, et clavo Christus fit cultor odoru(m). / SIC SEMPER FLORENT LILIA mixta rubis} ». La flèche est une référence particulière à la transverbération de Thérèse d’Avila (à propos de laquelle la référence obligatoire est la sculpture de Gian Lorenzo Bernini dans la chapelle Cornaro, à Santa Maria della Vittoria à Rome), moment d’extase mystique et de communion suprême avec l’amour divin décrit par la sainte dans sa propre autobiographie, le clou faisant allusion aux stigmates imprimés dans son cœur, toujours selon son témoignage.}}.
\end{teiQuote}
              
\end{teiCit}
              \teiP[xmlid={p19-5}]{Cette conjonction de l’image et du mot, qui repose sur le langage de l’héraldique et de l’emblématique, est caractéristique de la culture du \teiHi[rend={small-caps}]{xvii}\teiHi[rend={sup}]{e} siècle, et très ancrée dans celle de Bologne. Elle est aussi à la base des fresques de la galerie commandée par le cardinal Vidoni, dont la signification se révèle grâce au petit poème de Vincenzo Maria Marescalchi déjà cité. Grâce à lui, on peut en comprendre certains détails autrement difficiles à interpréter. Il en ressort l’impression que ce fut Marescalchi qui composa le programme iconographique de la décoration, et qu’il dut, lui ou un autre des érudits qui participèrent aux \teiHi[rend={italic}]{Ossequii poetici per Caterina Farnese,} diriger le burin de Domenico Fontana.}
              \teiP[xmlid={p20-5}]{L’\teiHi[rend={italic}]{Aula Vidonia} décrit la décoration dans un langage dense en métaphores visuelles, commençant par la célébration de la splendeur de l’étoile placée au-dessus des monts de l’écu Chigi (« \teiHi[rend={italic}]{de vertice Montis Chisia Stella micet} » \lbrack{}l’étoile Chigi resplendira du sommet du mont\rbrack{}, v. 2-3). Marescalchi énonce tout de suite le motto chigien « \teiHi[rend={italic}]{Micat in vertice} » \lbrack{}Elle brille au sommet\rbrack{} en le déclinant au futur, et cet élément se révèle crucial dans le poème et dans la décoration de la voûte, dont les scènes principales sont disposées avec intelligence autour de cette image.}
              \teiP[xmlid={p21-5}]{Le médaillon central montre Mercure déplaçant la voûte du ciel étoilé comme s’il s’agissait d’un manteau et découvrant la vision lumineuse de Jupiter flanqué de son aigle sur un nuage, qui pose une étoile sur une tiare papale ceinte de laurier, présentée par deux figures féminines agenouillées et embrassées (fig. 93). Il n’y a pas de doute sur la signification politique de l’image : la tiare, symbole du pouvoir papal, y compris temporel, se trouve justement couronnée par l’étoile à huit rayons, meuble des armes Chigi. Le concept visuel, qui joue indubitablement sur l’identification des monts chigiens à la tiare, est enrichi de deux figures agenouillées, identifiables comme Justice et Paix par leurs attributs, l’épée et le rameau d’olivier, dont l’union est enrichie par la couronne de laurier autour de la tiare, rappel de la vertu dont cette plante est le symbole, et peut-être de l’universalité du pontificat, successeur des anciens empereurs qui étaient effectivement couronnés de laurier. Les deux personnifications, déjà vues sur la miniature du sixième bimestre 1660 liée au gonfalonier Campeggi, renvoient évidemment au célèbre verset 11-12 du psaume\teiHi[rend={italic}]{ }85 (« Amour et vérité se rencontreront, Justice et Paix s’embrasseront »), et constituent une allégorie politique de la fortune, constante et répandue\teiNote[n={37},place={foot},xmlid={ftn78-2}]{\teiP[xmlid={p-266}]{Sur l’histoire iconographique de ce motif à forte valeur politique, en particulier sur son usage à Bologne entre \teiHi[rend={small-caps}]{xvi}\teiHi[rend={sup}]{e} et \teiHi[rend={small-caps}]{xvii}\teiHi[rend={sup}]{e} siècle, je renvoie à Sonia Cavicchioli, art. cité, p. 171 et 172. Voir aussi son usage dans la propagande, y compris posthume, du cardinal Mazarin : Yvan Loskoutoff, \teiHi[rend={italic}]{Rome des Césars Rome des Papes. La propagande du cardinal Mazarin}, Paris, Honoré Champion, 2007, p. 184-197.}}. La fresque vise à exalter le pontificat Chigi et son bon gouvernement, soulignant l’importance d’Alexandre VII sur le trône de Pierre. On pourrait l’intituler \teiHi[rend={italic}]{Jupiter couronne de l’étoile chigienne la tiare présentée par la Justice et par la Paix}. Grâce à la description donnée par Marescalchi, nous pouvons aussi interpréter les quatre cartouches (« \teiHi[rend={italic}]{bullae} ») splendidement peints en monochrome, insérés dans l’encadrement feint du médaillon. Ils figurent les quatre éléments par les figures de \teiHi[rend={italic}]{Cibèle }(la terre), \teiHi[rend={italic}]{Vulcain }(le feu), \teiHi[rend={italic}]{Junon }(l’air) et une \teiHi[rend={italic}]{Divinité fluviale }(l’eau). Naturellement, leur présence souligne l’universalité du pouvoir d’Alexandre VII célébré au centre de la fresque.}
              \teiP[xmlid={p22-5}]{Pour en venir aux fresques des travées latérales (fig. 94 et 95), il est intéressant de remarquer que, différemment du médaillon central, elles sont enrichies d’un phylactère\teiNote[n={38},place={foot},xmlid={ftn79-2}]{\teiP[xmlid={p-267}]{Le poème de Marescalchi suggère toutefois qu’une inscription était aussi prévue dans la peinture centrale, énonçant le \teiHi[rend={italic}]{motto} : « \teiHi[rend={italic}]{Ambiat hunc apicem nostrum, sine cuspide, fulgur} », traduisible en : « que la foudre ceigne notre tiare dépourvue de pointe », le terme « \teiHi[rend={italic}]{fulgur} » étant volontairement ambigu, pouvant se référer à la foudre, attribut de Jupiter, comme à l’étoile des Chigi.}}. Son inscription rend explicite leur nature d’emblème politique, outre celle d’image allégorique. L’étoile chigienne joue son rôle dans les deux cas, et là aussi l’élément héraldique agit comme un symbole dynamique. Dans l’une des deux fresques, Neptune est assis sur un nuage devant un ciel tempétueux, montrant de sa dextre la Prudence, parfaitement identifiable grâce à ses attributs usuels du serpent enroulé autour du bras et du miroir : plutôt que de se regarder dans celui-ci, elle intercepte un des huit rayons de l’étoile chigienne, rayon qui resplendit parmi les nuées (détail actuellement peu visible d’en bas) et s’imprime sur le cristal, alors que de son autre main elle tient la grande voile d’une figure identifiable comme la Fortune, qui paraît vouloir s’enfuir. La fresque a une évidente signification morale, outre que politique\teiNote[n={39},place={foot},xmlid={ftn80-3}]{\teiP[xmlid={p-268}]{L’idée est renforcée par les vers qui proposent la comparaison entre la Prudence dirigeant la voile de la Fortune et l’image connue de l’aurige qui tient la bride à un cheval rebelle (« \teiHi[rend={italic}]{indomiti \lbrack{}…\rbrack{} rector equi} », v. 132). Il faut remarquer que la présence de Neptune peut renvoyer au rôle reconnu à cette divinité antique depuis le \teiHi[rend={small-caps}]{xvi}\teiHi[rend={sup}]{e} siècle, documenté par les commentaires du passage de l’\teiHi[rend={italic}]{Énéide }virgilienne qui commence avec le célèbre « \teiHi[rend={italic}]{Quos ego}… » (I. 135) et par la fontaine réalisée justement sous les fenêtres de l’appartement du légat (voir Irving Lavin, « Giambologna’s Neptune at the Crossroads », dans \teiHi[rend={italic}]{Past-Present Essays on Historicism in Art from Donatello to Picasso}, Berkeley, University of California Press, 1993, p. 63-83).}}, et le \teiHi[rend={italic}]{motto} peint sur le phylactère, identiquement présent dans le poème de Marescalchi, en révèle le sens : « \teiHi[rend={italic}]{Sidus Prudentem ducit, Prudentia Sortem} » \lbrack{}L’étoile conduit celui qui est prudent, la Prudence fait de même avec la Fortune\rbrack{}, v. 135. La scène représente donc \teiHi[rend={italic}]{La Prudence, inspirée par l’étoile }(chigienne, évidemment)\teiHi[rend={italic}]{, qui tient en lisière la Fortune}. L’image et le \teiHi[rend={italic}]{motto} expriment une idée politique, non seulement parce que l’emblème des Chigi est là pour rappeler le rôle et le pouvoir du pape régnant, mais aussi parce que la Prudence elle-même est la vertu du gouvernement et de la politique par excellence. Ce n’est pas un hasard si elle est aussi représentée dans les \teiHi[rend={italic}]{specula principis}\teiNote[n={40},place={foot},xmlid={ftn81-3}]{\teiP[xmlid={p-269}]{Voir Giorgio Forni, « Nicolò Machiavelli », dans Ezio Raimondi (dir.), \teiHi[rend={italic}]{La letteratura italiana. Dalle origini al Cinquecento}, Milan, Bruno Mondadori, 2007, p. 346. }} \lbrack{}miroirs des princes\rbrack{}\teiHi[rend={italic}]{. }L’image de la Fortune à la voile, issue des monnaies antiques et reprise à l’époque moderne, est tout à fait cohérente avec l’idée que la vie humaine et les responsabilités du gouvernement peuvent s’identifier à la navigation et à ses périls, d’autant que dans son poème, Marescalchi identifie le pontife au pilote d’un navire occupé à conduire à bon port la « \teiHi[rend={italic}]{fidei cymba} », le navire de la foi.}
              \teiP[xmlid={p23-5}]{La signification de la troisième travée se trouve aussi éclairée par le poème. Il rappelle que la figure féminine « \teiHi[rend={italic}]{clavigera} » (elle tient les clefs de la main droite) écrase du talon le rebelle Pluton, le précipitant récemment enchaîné dans l’Adès avec les Furies, que la fresque dépeint minutieusement, avec leurs attributs mythologiques : la torche, les serpents et le cyprès, une des plantes sacrées propres à ces figures. Le poème précise que Pluton et les Furies représentent les puissances infernales, rebelles au ciel, et que la figure féminine – définie comme « Force » \lbrack{}\teiHi[rend={italic}]{Vis}\rbrack{} ou « Suprême Pouvoir » \lbrack{}\teiHi[rend={italic}]{Suprema Potestas}\rbrack{} – est en mesure de les écraser de son regard, de son pied et de sa dextre : « \teiHi[rend={italic}]{Ore, pede, et dextra, premit ima, suprema Potestas} » (v. 161), lisons-nous justement aussi bien dans le poème que dans le phylactère porté par le putto en vol dans la fresque. La fresque a donc pour sujet \teiHi[rend={italic}]{Le Pouvoir domine les puissances infernales }\lbrack{}\teiHi[rend={italic}]{La Potestà domina le potenze infere}\rbrack{}. Cette personnification semble tirée de l’\teiHi[rend={italic}]{Iconologia }de Cesare Ripa, qui identifie l’« Autorité ou Pouvoir » comme « une matrone \lbrack{}…\rbrack{} vêtue d’un habit riche et somptueux \lbrack{}qui\rbrack{} dans sa main droite levée tienne les deux clefs élevées\teiNote[n={41},place={foot},xmlid={ftn82-3}]{\teiP[xmlid={p-270}]{Cesare Ripa, \teiHi[rend={italic}]{op. cit}., p. 35-36 : « \teiHi[rend={italic}]{matrona (…) vestita d’abito ricco e sontuoso \lbrack{}che\rbrack{} con la destra mano alzata tenghi due chiavi elevate} ».}} ». Un détail s’ajoute à ce que décrit Ripa : la tête de la figure rayonne de lumière telle un astre, et je crois qu’elle doit être identifiée encore une fois comme un pouvoir non générique mais spécifique, celui d’Alexandre VII. L’article que lui consacre Ripa suggère que : « \lbrack{}l\rbrack{}es clefs dénotent l’autorité et le pouvoir spirituel, comme le démontre fort bien le Christ Notre Seigneur et Rédempteur, quand par leur moyen il donna la suprême autorité à saint Pierre \lbrack{}...\rbrack{}. Elle tient ces clefs dans la main droite parce que le pouvoir spirituel est le principal et le plus noble de tous \lbrack{}...\rbrack{} et il n’y a personne qui ne soit sujet à celui du souverain pontife vicaire du Christ sur la terre\teiNote[n={42},place={foot},xmlid={ftn83-3}]{\teiP[xmlid={p-271}]{\teiHi[rend={italic}]{Ibid.} : « \teiHi[rend={italic}]{Le chiavi denotano l’auttorità e potestà spirituale, come benissimo lo dimostra Cristo Nostro Signore e Redentore, quando per mezzo d’esse diede quella suprema auttorità a San Pietro (…). Tiene dette chiavi nella destra, perché la potestà spirituale è la principale; è più nobile di tutte l’altre (…) e non è alcuno che non sia suddito a quella del sommo Pontefice Vicario di Cristo in terra} ».}} ». Un hommage qui ne pouvait qu’être agréable au pontife dédicataire, exprimé par le jeu séculaire de l’héraldique au cœur de la ville qui s’obstinait à se proclamer république quoique ne l’étant pas.}
              \begin{teiFigure}[xmlid={figure01-5}]

                \teiGraphic[url={../icono/br/Ch06\_Cavicchioli/Fig86.jpg}]
                \teiHead[xmlid={figure-84-ecus-bologne-porte-galliera-facade-vers-la-ville-001}]{Figure 84.\teiHi[rend={italic}]{ Écus}, Bologne, Porte Galliera, façade vers la ville.}
              
\end{teiFigure}
              \begin{teiFigure}[xmlid={figure02-5}]

                \teiGraphic[url={../icono/br/Ch06\_Cavicchioli/Fig87.JPG}]
                \teiHead[xmlid={figure-85-anonyme-miniature-du-primier-bimestre-1658-bologne-archive-detat-anziani-consoli-insignia-vol-viii-p-28-001}]{Figure 85. Anonyme, \teiHi[rend={italic}]{Miniature du primier bimestre 1658}, Bologne, Archive d’État, \teiHi[rend={italic}]{Anziani Consoli, Insignia}, vol. VIII, p. 28.}
              
\end{teiFigure}
              \begin{teiFigure}[xmlid={figure03-5}]

                \teiGraphic[url={../icono/br/Ch06\_Cavicchioli/Fig88.jpg}]
                \teiHead[xmlid={figure-86-giovanni-andrea-abanti-miniature-du-cinquieme-bimestre-1658-bologne-archive-detat-anziani-consoli-insignia-vol-viii-p-38-001}]{Figure 86. Giovanni Andrea Abanti, \teiHi[rend={italic}]{Miniature du cinquième bimestre 1658}, Bologne, Archive d’État, \teiHi[rend={italic}]{Anziani Consoli, Insignia}, vol. VIII, p. 38.}
              
\end{teiFigure}
              \begin{teiFigure}[xmlid={figure04-5}]

                \teiGraphic[url={../icono/br/Ch06\_Cavicchioli/Fig89.jpg}]
                \teiHead[xmlid={figure-87-giovanni-andrea-abanti-miniature-du-sixieme-bimestre-1660-bologne-archive-detat-anziani-consoli-insignia-vol-viii-p-46-001}]{Figure 87. Giovanni Andrea Abanti, \teiHi[rend={italic}]{Miniature du sixième bimestre 1660}, Bologne, Archive d’État, \teiHi[rend={italic}]{Anziani Consoli, Insignia}, vol. VIII, p. 46.}
              
\end{teiFigure}
              \begin{teiFigure}[xmlid={figure05-5}]

                \teiGraphic[url={../icono/br/Ch06\_Cavicchioli/Fig90.jpg}]
                \teiHead[xmlid={figure-88-carlo-cignani-e-emilio-taruffi-entree-de-paul-iii-farnese-a-bologne-en-1543-bologne-palais-public-sala-regia-001}]{Figure 88. Carlo Cignani e Emilio Taruffi, \teiHi[rend={italic}]{Entrée de Paul III Farnèse à Bologne en 1543}, Bologne, Palais Public, sala Regia.}
              
\end{teiFigure}
              \begin{teiFigure}[xmlid={figure06-5}]

                \teiGraphic[url={../icono/br/Ch06\_Cavicchioli/Fig91.jpg}]
                \teiHead[xmlid={figure-89-albert-clouwet-sur-un-dessin-de-pietro-martire-neri-le-cardinal-pietro-vidoni-in-effigies-cardinalium-nunc-viventium-rome-giovanni-giacomo-de-rossi-1660-1676-the-trustee-of-the-british-museum-001}]{Figure 89. Albert Clouwet, sur un dessin de Pietro Martire Neri, \teiHi[rend={italic}]{Le cardinal Pietro Vidoni}, in \teiHi[rend={italic}]{Effigies cardinalium nunc viventium, }Rome, Giovanni Giacomo de Rossi, 1660-1676, The Trustee of the British Museum.}
              
\end{teiFigure}
              \begin{teiFigure}[xmlid={figure07-5}]

                \teiGraphic[url={../icono/br/Ch06\_Cavicchioli/Fig92.JPG}]
                \teiHead[xmlid={figure-90-giovanni-andrea-abanti-miniature-du-troisieme-bimestre-1662-bologne-archive-detat-anziani-consoli-insignia-vol-viii-p-56-001}]{Figure 90. Giovanni Andrea Abanti, \teiHi[rend={italic}]{Miniature du troisième bimestre 1662}, Bologne, Archive d’État, \teiHi[rend={italic}]{Anziani Consoli, Insignia}, vol. VIII, p. 56.}
              
\end{teiFigure}
              \begin{teiFigure}[xmlid={figure08-4}]

                \teiGraphic[url={../icono/br/Ch06\_Cavicchioli/Fig93.jpg}]
                \teiHead[xmlid={figure-91-domenico-santi-dit-le-mengazzino-et-giovan-battista-caccioli-galleria-vidoniana-appartement-du-legat-bologne-palais-communal-vue-densemble-001}]{Figure 91. Domenico Santi, dit le Mengazzino et Giovan Battista Caccioli, \teiHi[rend={italic}]{Galleria Vidoniana}, Appartement du légat, Bologne, Palais communal, vue d’ensemble.}
              
\end{teiFigure}
              \begin{teiFigure}[xmlid={figure09-4}]

                \teiGraphic[url={../icono/br/Ch06\_Cavicchioli/Fig94.jpg}]
                \teiHead[xmlid={figure-92-domenico-fontana-frontispice-des-ossequii-poetici-nellammirarsi-in-bologna-lheroica-risoluzione-della-sereniss-principessa-caterina-farnese-di-monacarsi-nel-ven-monastero-delle-madri-carmelitane-scalze-di-santa-teresa-in-parma-in-bologna-per-giacomo-monti-1661-bologne-bibliotheque-communale-de-larchiginnasio-001}]{Figure 92. Domenico Fontana, frontispice des \teiHi[rend={italic}]{Ossequii poetici nell’ammirarsi in Bologna l’Heroica risoluzione della sereniss. Principessa Caterina Farnese di monacarsi nel Ven. Monastero delle Madri carmelitane scalze di Santa Teresa in Parma}, In Bologna, Per Giacomo Monti, 1661, Bologne, Bibliothèque communale de l’Archiginnasio.}
              
\end{teiFigure}
              \begin{teiFigure}[xmlid={figure10-4}]

                \teiGraphic[url={../icono/br/Ch06\_Cavicchioli/Fig95.jpg}]
                \teiHead[xmlid={figure-93-giovan-battista-caccioli-et-domenico-santi-dit-le-mengazzino-jupiter-place-letoile-chigienne-sur-la-tiare-presentee-par-la-justice-et-la-paix-bologne-palais-communal-galleria-vidoniana-detail-de-la-voute-001}]{Figure 93. Giovan Battista Caccioli et Domenico Santi, dit le Mengazzino, \teiHi[rend={italic}]{Jupiter place l’étoile chigienne sur la tiare présentée par la Justice et la Paix}, Bologne, Palais communal, Galleria Vidoniana, détail de la voûte.}
              
\end{teiFigure}
              \begin{teiFigure}[xmlid={figure11-4}]

                \teiGraphic[url={../icono/br/Ch06\_Cavicchioli/Fig96.jpg}]
                \teiHead[xmlid={figure-94-giovan-battista-caccioli-et-domenico-santi-dit-le-mengazzino-la-prudence-inspiree-par-letoile-chigienne-tient-la-fortune-en-lisiere-bologne-palais-communal-galleria-vidoniana-detail-de-la-voute-001}]{Figure 94. Giovan Battista Caccioli et Domenico Santi, dit le Mengazzino, \teiHi[rend={italic}]{La Prudence, inspirée par l’étoile chigienne tient la Fortune en lisière}, Bologne, Palais communal, Galleria Vidoniana, détail de la voûte.}
              
\end{teiFigure}
              \begin{teiFigure}[xmlid={figure12-4}]

                \teiGraphic[url={../icono/br/Ch06\_Cavicchioli/Fig97.jpg}]
                \teiHead[xmlid={figure-95-giovan-battista-caccioli-et-domenico-santi-dit-le-mengazzino-le-pouvoir-domine-les-puissances-infernales-bologne-palais-communal-galleria-vidoniana-detail-de-la-voute-001}]{Figure 95. Giovan Battista Caccioli et Domenico Santi, dit le Mengazzino, \teiHi[rend={italic}]{Le Pouvoir domine les puissances infernales}, Bologne, Palais communal, Galleria Vidoniana, détail de la voûte.}
              
\end{teiFigure}
            
\end{teiDiv}
          
\end{teiElement}
        
\end{teiElement}
        \begin{teiElement}[name={group},type={chapter},data-page-title={Piranèse, l’héraldique et Clément XIII Rezzonico (1758-1769)},data-page-authors={Yvan Loskoutoff},data-include-href={Ch07\_Piranese\_heraldique.xml},xmlid={chapter-004}]

          \begin{teiElement}[name={front}]

            \begin{teiDiv}[type={titlePage},xmlid={titlepage-009}]

              \teiP[rend={title-main},xmlid={p-272}]{Piranèse, l’héraldique et Clément XIII Rezzonico (1758-1769)}
              \teiP[rend={author-aut},xmlid={p-273}]{Yvan Loskoutoff}
            
\end{teiDiv}
          
\end{teiElement}
          \begin{teiElement}[name={body},xmlid={body-009}]

            \teiP[xmlid={p1-9}]{C’est une idée commune que le \teiHi[rend={small-caps}]{xviii}\teiHi[rend={sup}]{e} siècle vit le déclin de l’héraldique. Piranèse va nous prouver le contraire. Cela peut surprendre, mais l’artiste des ruines romaines fut aussi celui du blason. En publiant ses carnets de dessins conservés à Modène, Mario Bevilacqua écrit : « Outre qu’il a inventé des écus nobiliaires originaux, Piranèse en a souvent copié\teiNote[n={1},place={foot},xmlid={ftn1-7}]{\teiP[xmlid={p-274}]{\teiHi[rend={italic}]{Piranesi, Taccuini di Modena}, Rome, Artemide, 2008, 2 vol., t. 1, c. 64v, « \teiHi[rend={italic}]{Oltre ad inventarne di originali, Piranesi ha ripetutamente copiato stemmi nobiliari }».}} ». Dans un autre écrit, le même auteur donne comme exemple d’invention de ce genre le dessin du Louvre intitulé \teiHi[rend={italic}]{Scène d’apothéose} ou \teiHi[rend={italic}]{Apothéose d’Isaac Newton} : « Au centre apparaît l’image ironique et profanatoire d’un écu fictif\teiNote[n={2},place={foot},xmlid={ftn34}]{\teiP[xmlid={p-275}]{« The Young Piranesi: the Itineraries of his formation », dans Mario Bevilacqua, Heather Hyde Minor et Fabio Barri (dir.), \teiHi[rend={italic}]{The Serpent and the Stylus, Essays on G. B. Piranesi}, Ann Arbor, The University of Michigan Press, « Memoirs of the American Academy in Rome », 2006, p. 13-53, ici p. 47 : « \teiHi[rend={italic}]{In the middle of it all appears the ironic, desecrating image of a mock escutcheon.} »}} ». On perçoit une approche complexe : Piranèse copie le blason mais il s’en joue aussi. Mario Bevilacqua traite des dessins, ce goût de l’artiste se manifeste néanmoins d’abord dans la gravure et dans l’architecture. Nous essaierons de montrer comment il s’appliqua à la papauté dans ces deux domaines.}
            \begin{teiDiv}[type={section1},xmlid={div01-6}]

              \teiHead[xmlid={piranese-et-lheraldique-001}]{Piranèse et l’héraldique}
              \teiP[xmlid={p2-8}]{Avant d’aborder la papauté, nous voudrions montrer l’intime familiarité de l’artiste avec le blason et même sa créativité en la matière. Nous avons choisi pour cela trois exemples pris dans l’ensemble de sa production.}
              \begin{teiDiv}[type={section2},xmlid={div02-6}]

                \teiHead[xmlid={le-frontispice-charlemont-001}]{Le frontispice Charlemont}
                \teiP[xmlid={p3-8}]{En 1756, Piranèse publie à Rome \teiHi[rend={italic}]{Le antichità romane}, dont les quatre volumes in-folio sont dédiés à son mécène, Milord Charlemont, un seigneur irlandais qui avait séjourné à Rome. Malheureusement, le soutien financier nécessaire à cette coûteuse entreprise ne correspondit pas aux attentes et l’artiste décida de supprimer la quadruple dédicace de l’édition en cours. Il rendit compte de cette décision l’année suivante dans un opuscule in-4° où il publiait la correspondance échangée pour l’occasion : \teiHi[rend={italic}]{Lettere di giustificazione scritte a Milord Charlemont e a’ di lui agenti di Roma, dal signor Piranesi socio della real società degli antiquari di Londra intorno la dedica della sua opera delle antichità rom. fatta allo stesso signore ed vltimamente soppressa}. \lbrack{}\teiHi[rend={italic}]{Lettres de justification écrites à Milord Charlemont et à ses agents de Rome par Monsieur Piranèse, membre de la Société royale des antiquaires de Londres, sur la dédicace de son œuvre des Antiquités de Rome faite audit seigneur et récemment supprimée}\rbrack{}. L’affaire a été étudiée mais sans aborder son caractère héraldique\teiNote[n={3},place={foot},xmlid={ftn35}]{\teiP[xmlid={p-276}]{Heather Hyde Minor, « Engraved in porphyry, printed on paper: Piranesi and Lord Charlemont », dans \teiHi[rend={italic}]{ibid.}, p. 123-147 ; Pierluigi Panza, « L’araldica e Giovan Battista Piranesi », \teiHi[rend={italic}]{Rivista del Collegio araldico}, 2017, p. 142-154.}}.}
                \teiP[xmlid={p4-8}]{Avant la rupture, le frontispice du premier volume des \teiHi[rend={italic}]{Antichità romane} situe le blason du dédicataire dans le contexte glorieux de la Rome antique (fig. 96). Piranèse s’explique dans la première de ses \teiHi[rend={italic}]{Lettere di giustificazione} :}
                \begin{teiCit}[xmlid={cit01-4}]

                  \begin{teiQuote}
La première planche, qui se trouve en tête du premier volume représente un marbre découvert parmi les ruines et montre l’inscription désignant le sujet auquel l’œuvre était consacrée. Sur les trophées antiques qui l’accompagnent on voit un écu portant les armes et les attributs de la maison de Milord, pour en montrer l’antiquité et le lustre\teiNote[n={4},place={foot},xmlid={ftn36-3}]{\teiP[xmlid={p-277}]{\teiHi[rend={italic}]{Lettere di giustificazione...}, In Roma, s. n., 1757, p. II, n. 3 : « \teiHi[rend={italic}]{La prima Tavola, che sta alla testa del primo Volume, rappresenta un marmo scoperto fra le rovine, e mostra l’Iscrizione col nome del Suggetto, a cui era dedicata l’Opera. Su i trofei antichi, che l’accompagnano, sono in uno Scudo le armi, e gli attributi della Casa di Milord, per mostrarne l’antichità ed il lustro.} »}}.
\end{teiQuote}
                
\end{teiCit}
                \teiP[xmlid={p5-8}]{Dans une intention de célébration, Piranèse avait donc associé deux imaginaires différents, celui du blason d’origine médiévale et celui de l’antiquité romaine. Le frontispice montre en effet sur un fond d’arc de triomphe une stèle portant l’inscription dédicatoire, et à gauche, parmi des trophées, les armoiries du Milord dont ses agents avaient fourni le modèle. Elles se décrivent ainsi et elles ont été fidèlement représentées :}
                \begin{teiCit}[xmlid={cit02-4}]

                  \begin{teiQuote}
Burelé d’argent et de gueules ; au canton de gueules chargé d’un léopard d’or. Cimier : une tête et col de dragon de gueules, colleté d’une jumelle d’argent. Supports : deux dragons de gueules ailés de sable, colleté chacun d’une jumelle d’argent. Devise : DEO DUCE, FERRO COMITANTE\teiNote[n={5},place={foot},xmlid={ftn37-3}]{\teiP[xmlid={p-278}]{Jean-Baptiste Rietstap, \teiHi[rend={italic}]{Armorial général}, Londres, Heraldry Today, 1988, 2 vol., t. 1, p. 388, « Caulfeild, comte de Charlemont ».}}.
\end{teiQuote}
                
\end{teiCit}
                \teiP[xmlid={p6-8}]{Placé sur un bouclier ovale l’écu s’inscrit parfaitement dans le trophée qui l’entoure. Néanmoins la célébration ne paraît pas univoque. En haut à droite, un fragment sculpté empiète sur la tête du dragon montrant une danseuse nue. On sait que Milord Charlemont menait galante vie et qu’il dut quitter Rome pour une affaire de jalousie. Dès avant la rupture, Piranèse nuançait donc son œuvre encomiastique d’un soupçon de satire.}
                \teiP[xmlid={p7-8}]{Dans la deuxième des \teiHi[rend={italic}]{Lettere di giustificazione}, Piranèse précise l’exécution de sa vengeance :}
                \begin{teiCit}[xmlid={cit03-3}]

                  \begin{teiQuote}
\lbrack{}…\rbrack{} pour ne pas gâter mes planches, j’imiterai ce qui se voit sur l’arc de Septime Sévère dans la ligne où fut enlevé le nom de Geta sur ordre de son frère Caracalla, et je substituerai un compliment au public qui sera autant juge que témoin de ce qui advint de part et d’autre relativement à cette dédicace\teiNote[n={6},place={foot},xmlid={ftn38-3}]{\teiP[xmlid={p-279}]{\teiHi[rend={italic}]{Lettere di giustificazione}..., \teiHi[rend={italic}]{op. cit}. : « \lbrack{}…\rbrack{} \teiHi[rend={italic}]{per non guastar le mie Tavole, imiterò quel che si vede sull’Arco di Settimio Severo nella linea, da cui fu tolto il nome di Geta per ordine di Caracalla suo fratello, e vi sostituirò un complimento al Pubblico, che sarà insieme e giudice e testimone delle cose passate dall’una e l’altra parte intorno a quella dedica.} »}}.
\end{teiQuote}
                
\end{teiCit}
                \teiP[xmlid={p8-8}]{Il y eut en effet deux états successifs de la gravure remaniée. Certains exemplaires, comme celui du fonds Barberini à la Bibliothèque vaticane\teiNote[n={7},place={foot},xmlid={ftn39-3}]{\teiP[xmlid={p-280}]{BAV, Stamp. Barb. X.I.20-23.}}, contiennent les trois frontispices. Dans le premier état modifié, l’inscription est effacée comme le voulait la \teiHi[rend={italic}]{damnatio memoriae} antique à laquelle Piranèse se réfère (fig. 97). Mais la modification s’applique aussi au blason qui se trouve « diffamé » suivant une pratique médiévale\teiNote[n={8},place={foot},xmlid={ftn40-3}]{\teiP[xmlid={p-281}]{Voir Laurent Hablot, « “\teiHi[rend={italic}]{Sens dessoubz dessus}”. Le blason de la trahison », dans Myriam Soria et Maïté Billoré (dir.), \teiHi[rend={italic}]{La trahison au Moyen Âge. De la monstruosité au crime politique (}\teiHi[rend={small-caps italic}]{v}\teiHi[rend={italic sup}]{e}\teiHi[rend={italic}]{-}\teiHi[rend={small-caps italic}]{xv}\teiHi[rend={italic sup}]{e}\teiHi[rend={italic}]{ siècle)}, Rennes, PUR, 2010, p. 331-347.}} dont la lettre ne dit rien. Dans la gloire comme dans la honte, l’antique et le médiéval cheminaient de concert. Comme prévu, dans le dernier état une nouvelle inscription remplaça la précédente, cette fois vouée non à un décevant mécène mais à l’utilité publique, \teiHi[rend={italic}]{VTILITATI PVBLICAE}. Piranèse en profita pour accentuer la diffamation des armoiries, qui fut appliquée non plus seulement à l’écu mais aussi à ses pièces extérieures (fig. 98).}
              
\end{teiDiv}
              \begin{teiDiv}[type={section2},xmlid={div03-5}]

                \teiHead[xmlid={le-frontispice-schouvaloff-001}]{Le frontispice Schouvaloff}
                \teiP[xmlid={p9-7}]{Favori de l’impératrice Élisabeth, le comte Ivan Schouvaloff (1727-1797), à l’avènement de l’impératrice Catherine, voyagea dans les capitales européennes pour s’installer à Rome où il résida de 1765 à 1774. Là, il multiplia les achats d’antiques pour les institutions d’État russes et pour son compte personnel. Il se lia avec Piranèse qui lui dédia le second volume des \teiHi[rend={italic}]{Vasi, candelabri, cippi, sarcofagi, tripodi, lucerne ed ornamenti antichi} publiés en 1778, l’année de sa mort. Le frontispice (fig. 99) a été étudié par Andrea Busiri Vici\teiNote[n={9},place={foot},xmlid={ftn41-3}]{\teiP[xmlid={p-282}]{« L’erudito della corte russa del Settecento, Ivan Ivanovitch Schouvaloff, ed i suoi rapporti con Roma », \teiHi[rend={italic}]{L’Urbe}, n\teiHi[rend={sup}]{o} 39, 1976, p. 39-52.}}, Giulia Fusconi\teiNote[n={10},place={foot},xmlid={ftn42-3}]{\teiP[xmlid={p-283}]{« Da Bartoli a Piranesi: spigolature dai Codici Ottoboniani Latini della Raccolta Ghezzi », \teiHi[rend={italic}]{Xenia Antiqua}, n\teiHi[rend={sup}]{o} 3, 1994, p. 145-172.}} et Olga Medvedkova\teiNote[n={11},place={foot},xmlid={ftn43-4}]{\teiP[xmlid={p-284}]{«\teiRef[target={https://fr.wikipedia.org/wiki/Š}]{ Š}uvalov à Rome (1765-1774). Histoire d’une dédicace », \teiHi[rend={italic}]{Cahiers du monde russe et soviétique}, vol. 52, n\teiHi[rend={sup}]{o} 1, 2011, p. 45-73.}} qui n’ont pas remarqué son contenu héraldique.}
                \teiP[xmlid={p10-7}]{Sur le haut de la gravure on lit : « \teiHi[rend={italic}]{Bassorilievo antico che si vede nel Portico della Chiesa de SS. Apostoli. Le vasi sono estratti dalla raccolta del Cavalier Ghezzi esistente nella Biblioteca Vaticana} » \lbrack{}Bas-relief antique que l’on voit sous le portique de l’église des Saints-Apôtres. Les vases sont extraits de la collection du chevalier Ghezzi existant à la Bibliothèque Vaticane\rbrack{}. Giulia Fusconi a en effet retrouvé deux de ces vases dans les recueils de dessins d’antiques réunis par Pier Leone Ghezzi et conservés à la Bibliothèque vaticane : celui du milieu, dont l’inscription ajoutée se rapporte au prince Théodore Golitsine, neveu du comte Schouvaloff qui l’accompagnait à Rome\teiNote[n={12},place={foot},xmlid={ftn44-4}]{\teiP[xmlid={p-285}]{Olga Medvedkova, art. cité.}}, et le vase à boire ou rhyton qui surmonte la dédicace (fig. 100). À propos de ce dernier, Giulia Fusconi écrit :}
                \begin{teiCit}[xmlid={cit04-3}]

                  \begin{teiQuote}
Quoique fidèle aux modèles, Piranèse enrichit ses copies de forts contrastes en clair-obscur et insère des éléments iconographiques nouveaux, comme la corne sur la tête du protome caprin ou la couronne de pampres au col du vase qui confèrent à l’objet des caractéristiques décoratives surabondantes et fantaisistes\teiNote[n={13},place={foot},xmlid={ftn45-4}]{\teiP[xmlid={p-286}]{Art. cité, p. 161 : « \teiHi[rend={italic}]{Pur fedele ai modelli, il Piranesi arrichisce le sue copie di forti contrasti chiaroscurali e inserisce elementi iconografici nuovi, come il corno sulla testa della protome caprina, o il serto di pampini attorno al collo del vaso, che conferisce all’oggetto caratteristiche di decorativismo sovrabbondante e fantasioso.} »}}.
\end{teiQuote}
                
\end{teiCit}
                \teiP[xmlid={p11-7}]{Ce qui peut paraître, de prime abord, un ajout de fantaisie décorative sur la source antique est en réalité un hommage héraldique au dédicataire opportunément situé au-dessus de la dédicace. Les Schouvaloff portent en effet : « Coupé : au 1 de sable au chevron d’or accosté de trois étoiles (5) d’argent et chargé de trois grenades du même, allumées de gueules ; au 2 de gueules à une licorne courante d’argent\teiNote[n={14},place={foot},xmlid={ftn46-4}]{\teiP[xmlid={p-287}]{Jean-Baptiste Rietstap, \teiHi[rend={italic}]{Armorial général},\teiHi[rend={italic}]{ op. cit.}, t. 2, p. 740, « Schuwalow ».}} ». Comme pour Milord Charlemont, l’héraldique se greffe sur l’antique, mais ici elle se fait allusive, si allusive qu’elle a échappé à la critique. On peut également se demander si l’aigle antique conservée à l’église des Saints-Apôtres, aimée de Piranèse qui la représente souvent et que nous retrouverons plus tard dans cette étude, n’évoque pas l’aigle impériale russe qui figure d’ailleurs en cimier des armes des Schouvaloff. L’allusion au blason serait ainsi plus complète, se référant non seulement au meuble de l’écu mais aussi à la pièce extérieure qui domine le tout.}
                \teiP[xmlid={p12-7}]{Giulia Fusconi écrit encore : « Une recherche plus étendue dans les manuscrits vaticans de la collection Ghezzi permettrait peut-être d’identifier les autres modèles utilisés\teiNote[n={15},place={foot},xmlid={ftn47-4}]{\teiP[xmlid={p-288}]{Art. cité, p. 161 : « \teiHi[rend={italic}]{Una più capillare ricerca fra i codici vaticani della raccolta Ghezzi permetterebe forse di individuare anche gli altri modeli utilizzati.} »}} ». Dans le même recueil qui contient le prototype du rhyton Schouvaloff, on trouve en effet une lampe à tête de taureau qui a inspiré le vase voisin du vase Golitsine\teiNote[n={16},place={foot},xmlid={ftn48-5}]{\teiP[xmlid={p-289}]{BAV, ms. Ottoboniani latini 3105.}}.}
              
\end{teiDiv}
              \begin{teiDiv}[type={section2},xmlid={div04-3}]

                \teiHead[xmlid={le-tombeau-de-cecilia-metella-001}]{Le tombeau de Cecilia Metella}
                \teiP[xmlid={p13-7}]{Dans ses \teiHi[rend={italic}]{Vues de Rome} exécutées tout au long de sa carrière, Piranèse représente les armoiries quand le monument qu’il traite en comporte. Le palais Farnèse, par exemple, est montré avec de grandes armes sur sa façade tout simplement parce qu’elles sont là. Cela n’attire aucune remarque particulière. Elles ne jouissent pas d’un traitement privilégié. Une de ces vues de Rome suscite pourtant le commentaire : celle du tombeau de Cecilia Metella sur la voie Appienne (fig. 101). L’artiste avait déjà traité le sujet dans ses \teiHi[rend={italic}]{Antichità Romane}, purement à l’antique. Néanmoins, le même sujet dans les \teiHi[rend={italic}]{Vedute di Roma} introduit au premier plan sur la droite une portion de mur médiéval portant les armoiries des Caetani, grands barons qui avaient dominé cette région. On s’interroge sur les intentions de l’artiste qui aurait pu se dispenser de ce détail. Il est vrai que cet ajout confère une épaisseur historique à la scène, le regard franchissant les époques, d’abord depuis le Moyen Âge pour ensuite arriver à l’Antiquité. Comme dans les exemples précédents, l’héraldique se trouve combinée à l’antique, cette fois non en synthèse mais par plans successifs. Dans la tour éventrée se dresse une silhouette phallique. La scène reste énigmatique.}
                \teiP[xmlid={p14-6}]{Ces trois exemples préliminaires nous ont offert un échantillon du sentiment héraldique chez Piranèse, toujours mis en relation avec l’antique, mais de manières variées, depuis la \teiHi[rend={italic}]{damnatio memoriae} satirique de Charlemont jusqu’à la célébration allusive de Schouvaloff. Avec la papauté nous allons retrouver la célébration, cryptée ou explicite.}
              
\end{teiDiv}
            
\end{teiDiv}
            \begin{teiDiv}[type={section1},xmlid={div05-2}]

              \teiHead[xmlid={clement-xiii-rezzonico-un-cycle-de-dedicaces-001}]{Clément XIII Rezzonico : un cycle de dédicaces}
              \teiP[xmlid={p15-6}]{Mario Bevilacqua signale dans les carnets de Modène un dessin des armes du pape Clément X Altieri (1670-1676) : il s’agit de celles qui se trouvent au-dessus de l’entrée du palais familial à Rome, place du Gesù, sans autre intention que la reproduction de la réalité\teiNote[n={17},place={foot},xmlid={ftn49-5}]{\teiP[xmlid={p-290}]{\teiHi[rend={italic}]{Piranesi, Taccuini di Modena, op. cit., }t. 1, c. 64v, « \teiHi[rend={italic}]{Stemma di papa Clemente X Altieri sulla facciata di palazzo Altieri al Gesù} ».}}. C’est sous le règne de Clément XIII Rezzonico (1758-1769), vénitien comme Piranèse, que l’inspiration héraldique prit chez celui-ci son plein développement encomiastique. Il composa, dans la première moitié du pontificat, une suite de trois frontispices pour trois in-folios dédiés au pape, auxquels il faut en ajouter un à la fin du pontificat à l’intention de son neveu Monseigneur Giovan Battista Rezzonico. Ces pièces ont été présentées dans des catalogues d’exposition mais sans étude approfondie du contenu armorial\teiNote[n={18},place={foot},xmlid={ftn50-5}]{\teiP[xmlid={p-291}]{Barbara Jatta (dir.), \teiHi[rend={italic}]{Piranesi e l’Aventino}, Milan, Electa, 1998, chap. IV : « Piranesi negli anni del Priorato » ; \teiHi[rend={italic}]{Clemente XIII Rezzonico, Un papa veneto nella Roma di metà Settecento}, catalogue d’exposition, Padoue, Museo Diocesiano, 12 décembre 2008-15 mars 2009, Milan, Silvana Editoriale, 2008, n\teiHi[rend={sup}]{o} 126-145, p. 188-205.}}.}
              \teiP[xmlid={p16-6}]{Il faut d’abord signaler un dessin de Piranèse au Metropolitan Museum dont la destination d’usage n’est pas connue\teiNote[n={19},place={foot},xmlid={ftn51-5}]{\teiP[xmlid={p-292}]{The Metropolitan Museum of Art, \teiHi[rend={italic}]{Coat of arms of Clement XIII (Rezzonico)}, Accession Number 52.570.155.}} (fig. 102). Il représente les armoiries de Clément XIII avec mention des couleurs, ce qui laisse supposer une étude préparatoire pour une peinture. Les armes se décrivent ainsi :}
              \begin{teiCit}[xmlid={cit05-3}]

                \begin{teiQuote}
Écartelé au 1\teiHi[rend={sup}]{er} de gueules à la croix d’argent, au 2 et 3 d’azur à la tour d’argent, au 4\teiHi[rend={sup}]{e} barré de gueules et d’argent. Sur le tout d’or à l’aigle à deux têtes de sable. Cet écusson couronné d’une couronne à l’antique d’or\teiNote[n={20},place={foot},xmlid={ftn52-5}]{\teiP[xmlid={p-293}]{Ferruccio Pasini Frassoni, \teiHi[rend={italic}]{Essai d’armorial des papes}, Rome, Collège héraldique, 1906, p. 48, « Clément XIII (1758-69) ».}}.
\end{teiQuote}
              
\end{teiCit}
              \teiP[xmlid={p17-6}]{Quintiliano Rezzonico, riche marchand, s’était vu conférer la noblesse par le Sénat de Venise en 1687 pour avoir généreusement financé la guerre de Morée. La croix du premier quartier renverrait à la ville de Côme où la famille prit son essor. Elle portait le titre de comtes\teiHi[rend={italic}]{ della Torre,} ce qui explique la présence de la tour aux quartiers 2 et 3, armes parlantes. L’aigle à deux têtes sur le tout indique le titre de barons du Saint-Empire.}
              \teiP[xmlid={p18-6}]{Les quatre frontispices que nous allons étudier forment une série par leur inspiration commune : tous représentent une pièce d’architecture à l’antique où les armoiries ont été sculptées.}
              \begin{teiDiv}[type={section2},xmlid={div06-2}]

                \teiHead[xmlid={della-magnificenza-ed-architettura-de-romani-1761-001}]{\teiHi[rend={italic}]{Della magnificenza ed architettura de’ Romani} (1761)}
                \teiP[xmlid={p19-6}]{L’ouvrage étant bilingue, il comporte deux frontispices-titres. Dans l’exemplaire que nous avons consulté, le titre italien précède le titre latin\teiNote[n={21},place={foot},xmlid={ftn53-5}]{\teiP[xmlid={p-294}]{BAV, Cicognara.X.3853.}}. L’un et l’autre sont traités comme des stèles montrant une inscription gravée dans le marbre, sous laquelle vient s’insérer un motif sculpté. Dans la version italienne, nous voyons l’aigle des Saints-Apôtres qui reparaîtra au frontispice Schouvaloff. Au titre latin sont réservées les armoiries, ce qui peut sembler paradoxal : leur origine médiévale les eût plutôt associées à l’italien, mais le latin est pourvoyeur de gloire (fig. 103). Sans aucun souci de vraisemblance historique, Piranèse a voulu leur conférer la grandeur antique et il l’a fait en imitant un relief célèbre de la colonne Trajanne : la Victoire qui écrit sur un bouclier. Installée dans un médaillon, ici elle n’inscrit pas le nom des batailles gagnées ou des villes prises mais elle achève de sculpter les armes du pontife, mettant une ultime touche à la tiare. Elle est devenue une artiste au travail, telle un autre Piranèse ; une règle est appuyée sur le bouclier, lui-même adossé à une sphinge sur un sarcophage. Suit un portrait de Clément XIII dont le fauteuil porte les armes, ce qui était de tradition dans ce genre de représentation même si elles sont à demi dissimulées ici. Enfin, le texte de la dédicace, sur le modèle des manuscrits, commence par une initiale, V, ornée de la tiare et des clefs, celle-ci étant conçue comme un chapiteau, en accord avec les frontispices architecturaux. En trois étapes, le titre, le portrait et l’initiale, se développe l’appareil de la gloire héraldique du dédicataire au seuil de l’ouvrage.}
                \teiP[xmlid={p20-6}]{La même année, Piranèse publia \teiHi[rend={italic}]{Le rovine del castello dell’acqua Giulia}. S’il n’y a pas de grand appareil dédicatoire au pape, on remarque quand même un discret hommage héraldique dans deux culs-de-lampe. Le premier se compose d’un fragment d’architecture accompagné de livres et de médailles. Du sommet crénelé, à l’arrière-plan, émergent les deux têtes de l’aigle Rezzonico (fig. 104). Le second est semblablement conçu. Autour d’un bouclier sont disposés les instruments des arts et des sciences. À l’arrière-plan sur la droite, se profilent l’aigle à deux têtes et la tour. Ces deux vignettes suggèrent l’idée du mécénat pontifical.}
              
\end{teiDiv}
              \begin{teiDiv}[type={section2},xmlid={div07-2}]

                \teiHead[xmlid={lapides-capitolini-1762-001}]{\teiHi[rend={italic}]{Lapides Capitolini} (1762)}
                \teiP[xmlid={p21-6}]{L’année suivante Piranèse dédia au pontife ses \teiHi[rend={italic}]{Lapides Capitolini sive fasti consularesque Romanorum}. Cette fois, les armoiries sont réservées non point au titre, mais à la dédicace (fig. 105). Celle-ci a aussi été conçue comme une stèle, moins ruiniste cependant : on constate seulement une ébréchure en haut à droite. L’héraldique forme un encadrement. En haut, les armoiries, disposées entre deux \teiHi[rend={italic}]{putti} phytomorphes, composent une frise. Comme précédemment, il s’agit d’une citation de l’antique, renvoyant à un fragment de frise du forum de Trajan, alors conservé à Rome au palais Aldobrandini de Magnanapoli. Ce fragment sera reproduit dans le second volume des \teiHi[rend={italic}]{Vasi, candelabri, cippi}. Comme précédemment, les \teiHi[rend={italic}]{putti} formant supports sont représentés en train de sculpter les armoiries. À un an de distance, les deux compositions obéissent à des principes similaires. La dédicace, inscrite sur une feuille de papier, est encadrée d’une guirlande, supportée en bas par deux meubles issus du blason du pape : la tour et l’aigle. Le portrait n’a pas été traité séparément mais figuré comme l’avers d’une médaille du pontife en haut de la feuille dédicatoire.}
              
\end{teiDiv}
              \begin{teiDiv}[type={section2},xmlid={div08-2}]

                \teiHead[xmlid={antichita-dalbano-e-di-castel-gandolfo-1764-001}]{\teiHi[rend={italic}]{Antichità d’Albano e di Castel Gandolfo} (1764)}
                \teiP[xmlid={p22-6}]{Deux ans plus tard, paraissent les \teiHi[rend={italic}]{Antichità d’Albano e di Castel Gandolfo. }Après le frontispice-titre, composé d’une accumulation de fragments antiques, vient le frontispice-dédicatoire en format double dépliant (fig. 106). Il représente une stèle où la dédicace s’inscrit dans le cercle du zodiaque. Tout autour sont répandus des fragments d’inscriptions et des médailles. À gauche de la stèle, sur un fragment de petit format, plongé dans une demi-pénombre, sont sculptées les armoiries papales. Évidemment, cela surprend, après les deux précédents frontispices où les armes du pape Rezzonico occupaient une position centrale. Ici, elles semblent réduites à l’accessoire. Peut-être l’artiste a-t-il voulu varier la présentation. En tout cas, c’est le dernier des trois frontispices aux armes de Clément XIII. Il n’y en aura plus d’autres jusqu’à sa mort cinq ans plus tard en 1769. On peut se demander si la diminution des armoiries n’annonce pas ce silence à venir. Le texte de la dédicace commence par un S orné de l’aigle des Saints-Apôtres.}
              
\end{teiDiv}
              \begin{teiDiv}[type={section2},xmlid={div09-2}]

                \teiHead[xmlid={diverse-maniere-dadornare-i-cammini-1769-001}]{\teiHi[rend={italic}]{Diverse maniere d’adornare i cammini} (1769)}
                \teiP[xmlid={p23-6}]{Si Piranèse ne dédia plus rien à Clément XIII, cela ne signifie pas qu’il rompit avec la cour pontificale. En 1769, il dédia à Monseigneur Giovan Battista Rezzonico, neveu du pape, les \teiHi[rend={italic}]{Diverse maniere d’adornare i cammini. }Il explique dans l’épître dédicatoire avoir travaillé à la décoration des appartements du prélat ainsi que de ceux de son frère, prenant la défense des modèles égyptiens ou étrusques. Le frontispice dépliant de grand format inclut le titre et la dédicace, conçus comme un monument sur fond nocturne (fig. 107). Le décor n’est pas ruiniste mais égyptianisant. De part et d’autre de l’inscription, deux Diane polymastes entourent une momie couverte de hiéroglyphes. L’héraldique, qui s’était associée à la gloire des ruines de Rome, ne s’accorda pas moins aux mystères de l’Égypte ancienne. Certains théoriciens ne faisaient-ils pas, en effet, remonter les origines du blason aux hiéroglyphes ? Sous l’inscription, figurent les armes du prélat dans un écu en « tête de cheval », supporté par une guirlande que tiennent deux des Diane polymastes. L’ordre des quartiers a été modifié : au 1 et 4 la tour, au 2 la croix, au 3 les barres. Les autres exemples que nous connaissons des armes du prélat respectent pourtant le modèle de celles du pape. Il s’agit probablement d’un effet d’inversion dû à la gravure, car les barres sont devenues des bandes. Par ailleurs, on peut se demander si le choix de l’écu en « tête de cheval » n’est pas à mettre en rapport avec les couples de chevaux affrontés qui se trouvent juste au-dessus.}
                \teiP[xmlid={p24-4}]{Dans cette suite de quatre frontispices, Piranèse s’est efforcé d’inclure l’héraldique dans l’antique, qu’il fût romain ou égyptien. C’était une manière d’inscrire la célébration du dédicataire dans la longue durée et d’accroître ainsi sa gloire.}
              
\end{teiDiv}
            
\end{teiDiv}
            \begin{teiDiv}[type={section1},xmlid={div10-2}]

              \teiHead[xmlid={loeuvre-architecturale-001}]{L’œuvre architecturale}
              \teiP[xmlid={p25-4}]{Piranèse fut l’auteur de deux projets architecturaux, celui de Saint-Jean de Latran et celui de l’église Santa Maria du prieuré de Malte. Seul le second fut réalisé, mais l’un et l’autre présentent un caractère héraldique.}
              \begin{teiDiv}[type={section2},xmlid={div11-2}]

                \teiHead[xmlid={saint-jean-de-latran-001}]{Saint-Jean de Latran}
                \teiP[xmlid={p26-4}]{En 1763, Piranèse reçut commande de Clément XIII pour la réfection du chœur de Saint-Jean de Latran. Le projet fut abandonné en 1767, mais la bibliothèque Avery de l’université de Columbia conserve les dessins préparatoires de l’artiste\teiNote[n={22},place={foot},xmlid={ftn54-5}]{\teiP[xmlid={p-295}]{Voir \teiHi[rend={italic}]{Piranesi architetto}, catalogue d’exposition, American Academy in Rome, 15 mai-5 juillet 1992, Rome, Edizioni dell’Elefante, 1992. }}. Parmi eux, quatre projets de baldaquin d’esprit borrominien. L’un s’orne de l’aigle des Saints-Apôtres, qui équivalait dans l’imaginaire piranésien à celui des Rezzonico. L’autre présente la tour surmontée de l’aigle.}
              
\end{teiDiv}
              \begin{teiDiv}[type={section2},xmlid={div12-2}]

                \teiHead[xmlid={leglise-santa-maria-del-priorato-sur-laventin-001}]{L’église Santa Maria del Priorato sur l’Aventin}
                \teiP[xmlid={p27-4}]{En 1764, alors qu’il travaillait au projet de décor de la basilique du Latran, Piranèse reçut commande de la part de Monseigneur Giovan Battista Rezzonico, neveu de Clément XIII et grand prieur de l’ordre de Malte, pour la rénovation de l’église Santa Maria del Priorato sur l’Aventin. Les travaux se déroulèrent de novembre 1764 à octobre 1766. Le détail en a été recueilli dans un livre de comptes aujourd’hui conservé à la bibliothèque Avery de l’université de Columbia\teiNote[n={23},place={foot},xmlid={ftn55-5}]{\teiP[xmlid={p-296}]{Ce livre a été publié par Pierluigi Panza en annexe de \teiHi[rend={italic}]{Piranesi architetto}, Milan, Guerini Studio, 1998.}}. Concernant le décor sculpté, Piranèse retint le principe de répandre partout les meubles du blason Rezzonico. On peut examiner ce décor en s’aidant d’une description manuscrite situable entre 1766 et 1769, conservée dans les archives de l’ordre et publiée à la fin du siècle dernier. On procédera en trois étapes : le portail d’entrée sur la place adjacente, l’extérieur de l’église et enfin son intérieur.}
                \begin{teiDiv}[type={section3},xmlid={div13}]

                  \teiHead[xmlid={la-place-des-chevaliers-et-le-portail-dentree-001}]{La place des chevaliers et le portail d’entrée}
                  \teiP[xmlid={p28-4}]{Le manuscrit conservé dans les archives de l’ordre précise que l’on a voulu donner au portail d’entrée l’apparence d’un arc triomphal. Il est composé de trois reliefs, l’un au-dessus de la porte, les deux autres de part et d’autre. Pour le relief central (fig. 108), le document nous renseigne sur la double interprétation de l’un des meubles du blason Rezzonico, la croix grecque ou croix de saint Jean-Baptiste présente au premier quartier de l’écu :}
                  \begin{teiCit}[xmlid={cit06-2}]

                    \begin{teiQuote}
\lbrack{}…\rbrack{} à la clef de l’arche de la porte on a mis en bas-relief un type de la puissance maltaise, exprimé par un grand heaume à l’héroïque soutenu par l’emblème de deux navires et entouré d’écus, drapeaux et hallebardes, tant des leurs que de ceux pris aux barbares. Ce heaume couronne la croix, enseigne de Malte et appartenant aussi aux armoiries de l’écu nobiliaire de la maison Rezzonico, dont les attributs se voient pour cette raison représentés en certain des écus précédemment mentionnés\teiNote[n={24},place={foot},xmlid={ftn56-5}]{\teiP[xmlid={p-297}]{John Edward Critien, « Un manoscritto dell’Archivio Magistrale del Sovrano Militare Ordine di Malta di Roma », dans Barbara Jatta (dir.), \teiHi[rend={italic}]{op. cit.}, p. 79-85, ici p. 85 : « \lbrack{}…\rbrack{}\teiHi[rend={italic}]{ nella proteride della porta si è posto in bassorilievo un Tipo della potenza Maltese, espresso con un grand’elmo all’eroica, sostenuto dall’Emblemma di due Navilj, e circondato da scudi, e banderie e labarde, si sue che prese a barbari. Quest’elmo incorona la Croce, o si Insegna tanto di Malta, quanto appartenente allo stemma gentilizio della Casa Rezzonico, i cui attributi si veggono perciò effiggiati in qualcuno de’detti scudi. }» Piranèse aurait participé à la rédaction de ce texte : voir Fabio Barry, « “Onward Christian Soldier”: Piranesi at Santa Maria del Priorato », dans Mario Bevilacqua et Daniela Gallavotti Cavallero (dir.), \teiHi[rend={italic}]{L’Aventino dal Rinascimento a oggi},\teiHi[rend={italic}]{ Arte e architettura}, Rome, Editoriale Artemide, 2010, p. 140-175.}}.
\end{teiQuote}
                  
\end{teiCit}
                  \teiP[xmlid={p29-4}]{Par cette croix de saint Jean, Piranèse avait trouvé le moyen d’associer deux gloires en un même symbole, celle de l’ordre de Malte et celle de la famille de son grand prieur, qui était aussi celle du pape. Ce double potentiel symbolique explique le choix d’une position centrale dans le dispositif décoratif. Comme le précise la description, les autres meubles des armoiries Rezzonico sont présents parmi les trophées : on remarque à gauche la tour des comtes della Torre sur un écu et à droite une tête d’aigle en proue de navire devant les drapeaux.}
                  \teiP[xmlid={p30-4}]{Les panneaux latéraux, quant à eux, comportent en leur centre l’aigle à deux têtes. Le manuscrit décrit correctement chaque panneau mais en intervertissant celui de droite et celui de gauche :}
                  \begin{teiCit}[xmlid={cit07}]

                    \begin{teiQuote}
Au centre des deux bas-reliefs se trouve l’aigle bicéphale, meuble nobiliaire de la même maison Rezzonico. On voit dans le bas-relief à main droite un développement d’écus, l’arc, la pharètre, le heaume, la lance, l’ancre, le timon, la tour et autres symboles sur lesquels glisse le serpent de l’Éternité pour en perpétuer la gloire\teiNote[n={25},place={foot},xmlid={ftn57-5}]{\teiP[xmlid={p-298}]{John Edward Critien, art. cité. : « \teiHi[rend={italic}]{Siede in mezzo ai due bassorilievi laterali la doppia Aquila pur gentilizia della stessa Casa Rezzonico. Vedesi nel bassorilievo di man diritta una espansione di scudi, l’arco, la faretra, l’elmo, l’aste, l’ancora, il timone, la torre ed altri simboli, su quali scorre a perpetuarne la gloria il serpente dell’Eternità} ».}}.
\end{teiQuote}
                  
\end{teiCit}
                  \teiP[xmlid={p31-4}]{La tour se trouve en effet à l’arrière-plan sur la droite du relief (fig. 109). On a remarqué à juste titre que l’aigle bicéphale avait été représentée au sein d’une couronne de lauriers sur le modèle de l’aigle des Saints-Apôtres, certainement l’une des images les plus récurrentes dans l’œuvre de Piranèse. Le panneau de droite est pareillement conçu :}
                  \begin{teiCit}[xmlid={cit08}]

                    \begin{teiQuote}
Dans le bas-relief de droite se voit une aile déployée symbole de la Réputation unie au volume de la geste et de la force de la Sacrée Religion signifiée par la tour dessinée dans le volume même et faisant aussi allusion à la maison Rezzonico, par la colonne et les écus que couvre la même aile. Il y a la trompe de cette Réputation et ses mains qui la tiennent et des acrostoles ou extrémités de navires, indiquant la principale sorte d’armée par laquelle la Sacrée Religion a remporté ses victoires\teiNote[n={26},place={foot},xmlid={ftn58-5}]{\teiP[xmlid={p-299}]{\teiHi[rend={italic}]{Ibid. }: « \teiHi[rend={italic}]{Nel bassorilievo poi da sinistra vedesi un’ala spansa, simbolo della Fama unita al volume delle gesta e della fortezza della Sacra Religione Significata per la Torre delineata nel volume stesso ed insieme allusiva alla casa Rezzonico, per la colonna e per gli scudi coperti dalla stessa ala. Sonovi le trombe della stessa Fama, e le mani di essa tenenti questa trombe, e degli acrostolij, o siano estremità de’navilj, indicanti la principale specie di milizia con cui la Sac: Religione ha riportato le sue vittorie.} »}}.
\end{teiQuote}
                  
\end{teiCit}
                  \teiP[xmlid={p32-4}]{La tour apparaît effectivement sur le volume de la geste des chevaliers et, de même que la croix de saint Jean, un double sens lui est attribué en relation avec l’ordre de Malte, la force, et avec la famille Rezzonico dont elle est le symbole.}
                  \teiP[xmlid={p33-4}]{La place est ceinte de quatre stèles, une forme architecturale aimée de Piranèse, particulièrement pour ses frontispices. Le manuscrit nous apprend que la stèle centrale, celle qui fait face au portail, est ornée d’un trophée d’armes turques placé derrière l’écu de Malte. Dans la partie inférieure, une tête de Gorgone signifie la terreur des ennemis, et le manuscrit ajoute à propos des trophées qu’ils : « reçoivent en leur milieu ici aussi la tour nobiliaire de la maison Rezzonico\teiNote[n={27},place={foot},xmlid={ftn59-5}]{\teiP[xmlid={p-300}]{\teiHi[rend={italic}]{Ibid. }: « \teiHi[rend={italic}]{ricevono in mezzo anco qui la torre gentilizia della Casa Rezzonico} ».}} ». Mais nous savons très bien que cette tour est aussi devenue, dans ce contexte, le symbole de la force des chevaliers. Elle figure également sur un croissant turc au sommet des deux stèles latérales. Pénétrons maintenant dans les jardins de la villa et dirigeons-nous vers l’église.}
                
\end{teiDiv}
                \begin{teiDiv}[type={section3},xmlid={div14}]

                  \teiHead[xmlid={lexterieur-de-leglise-001}]{L’extérieur de l’église}
                  \teiP[xmlid={p34-3}]{Piranèse avait ajouté à la façade un fronton. On peut s’en faire une idée grâce à une gravure de Cassini datant de 1779. Les deux tourelles latérales sont ornées de l’aigle bicéphale. Au tympan, les armes sont celles de saint Pie V préalablement présentes. Un dessin anonyme du musée Soane à Londres offre une image plus détaillée et différente sur certains points. Les aigles bicéphales sont surmontées par les deux tours et la croix de saint Jean. Sur les montants latéraux des boucliers montrent les barres. Au tympan, se dessinent les armoiries complètes de Clément XIII dont le livre de compte dit qu’elles furent dorées. En 1849, un bombardement des troupes françaises détruisit le fronton et au tympan les armoiries furent remplacées par la croix grecque de Malte et des Rezzonico\teiNote[n={28},place={foot},xmlid={ftn60-5}]{\teiP[xmlid={p-301}]{Sur l’histoire de la façade d’après les livres de compte, voir Pierluigi Panza, \teiHi[rend={italic}]{Piranesi architetto},\teiHi[rend={italic}]{ op. cit.}, p. 79-81.}}.}
                  \teiP[xmlid={p35-3}]{Les quatre pilastres se révèlent intéressants (fig. 110). Leurs chapiteaux montrent la tour entre deux sphinges, d’après un motif antique alors conservé à la villa Borghèse et représenté à la planche XIII du \teiHi[rend={italic}]{Della magnificenza}. Comme dans ses frontispices gravés, l’artiste inclut des sources antiques. Dans les années où il travaillait à cette église, en 1765, il publia son \teiHi[rend={italic}]{Parere sul’architettura}. Il y est fait référence à un ordre architectural britannique que James Adam, le frère de Robert Adam, lui aurait montré en 1762, exemple de l’application de l’antiquité aux besoins de l’époque moderne\teiNote[n={29},place={foot},xmlid={ftn61-5}]{\teiP[xmlid={p-302}]{John Wilton-Ely, \teiHi[rend={italic}]{Piranesi as Architect and Designer}, New York-New Haven-Londres, The Pierpont Morgan Library-Yale University Press, 1993, p. 51.}}. Or, cet ordre britannique inclut dans un chapiteau corinthien les deux supports des armoiries du Royaume-Uni, la licorne et le lion. On peut donc établir un parallèle entre cet exemple illustrant le \teiHi[rend={italic}]{Parere su l’architettura} et l’ordre Rezzonico inventé à la même époque pour Santa Maria del Priorato. Cela dit, la pratique d’introduire des meubles armoriaux dans les chapiteaux n’était pas nouvelle. On peut penser à ce propos aux colombes d’Innocent X Doria Pamphili sur la façade de Sant’Agnese place Navone, dues à Borromini, artiste admiré de Piranèse (1653-1657). Pour la même famille, on peut mentionner les lis et les rameaux d’olivier des colonnes de leur palais sur le Corso, œuvre de Gabriele Valvassori (1734). Piranèse, tout en flattant sa passion armoriale, se pliait à une pratique déjà bien admise.}
                  \teiP[xmlid={p36-3}]{Sur les pilastres eux-mêmes, sont insérés des reliefs représentant des épées dont certaines sont porteuses de la symbolique Rezzonico. Celle de l’extrême gauche montre un fourreau décoré de la tour surmontée de l’aigle bicéphale et une poignée en tête d’aigle. À droite, celle qui est proche de la porte laisse voir les deux têtes de l’aigle couronnée en guise de poignée.}
                  \teiP[xmlid={p37-3}]{À l’arrière de l’église, derrière le chœur, sous l’aigle et la croix grecque, on remarque deux médailles dédiées au commanditaire Monseigneur Giovan Battista Rezzonico (fig. 111). Elles sont inspirées de ses armoiries. Celle de gauche montre la tour accompagnée du mot : \teiHi[rend={italic}]{SACRA AEDE RESTITVT\lbrack{}A\rbrack{}} \lbrack{}Le sanctuaire ayant été restauré\rbrack{}. Celle de droite, les barres accompagnées du nom : \teiHi[rend={italic}]{IOAN. BAPTISTA. REZZONICO. M. P.} Piranèse avait un goût marqué pour la numismatique\teiNote[n={30},place={foot},xmlid={ftn62-5}]{\teiP[xmlid={p-303}]{Carolyn Yorke et Heather Hyde Minor, \teiHi[rend={italic}]{Piranese unbound}, Princeton, Princeton University Press, 2020, chapitres « Archaeology, inscriptions, iconography » et « Minting for Mr Menzies ».}}, tout comme Clément XIII Rezzonico\teiNote[n={31},place={foot},xmlid={ftn63-5}]{\teiP[xmlid={p-304}]{Giancarlo Alteri, « Monete e medaglie del pontificato di Clemente XIII (1758-1769) », \teiHi[rend={italic}]{Piranesi e l’Aventino},\teiHi[rend={italic}]{ op. cit}., p. 204.}} (le frontispice des \teiHi[rend={italic}]{Lapides Capitolini} est orné de ses médailles, voir fig. 105). Nous ne connaissons pas d’exemplaire frappé de ces médailles (ou deux faces d’une même médaille) qui demeurent donc à l’état d’ornement architectural.}
                
\end{teiDiv}
                \begin{teiDiv}[type={section3},xmlid={div15}]

                  \teiHead[xmlid={linterieur-de-leglise-001}]{L’intérieur de l’église}
                  \teiP[xmlid={p38-3}]{À l’intérieur, se poursuit le procédé des chapiteaux héraldiques. Cette fois, les sphynges ont été remplacées par un simple ornement ionique où s’inscrivent en alterance non seulement la tour mais aussi les autres meubles du blason, croix et aigle bicéphale couronnée (fig. 112). La pièce la plus importante se situe dans l’abside du chœur : on y voit les armes de Giovan Battista Rezzonico disposées dans une coquille (fig. 112). Au-dessus de l’aigle centrale, a été ajoutée la « basilique », c’est-à-dire l’ombrelle et les clefs croisées, privilège réservé aux membres de familles papales.}
                  \teiP[xmlid={p39-3}]{Sur la voûte de la nef, le relief central montre la croix de saint Jean dans une gloire. Elle figure aussi sur une tunique de chevalier surmontée de la tiare papale. Un des deux reliefs latéraux, celui du côté du chœur (fig. 113), est ainsi décrit dans le manuscrit précédemment cité :}
                  \begin{teiCit}[xmlid={cit09}]

                    \begin{teiQuote}
Puis dans le second se voit l’écu emblématique de la maison Rezzonico flanqué de deux heaumes où sont figurés des tritons, qui, mi-hommes, mi-poissons, font allusion au chevaliers du Sacré Ordre puissants sur mer comme sur terre\teiNote[n={32},place={foot},xmlid={ftn64-5}]{\teiP[xmlid={p-305}]{John Edward Critien, art. cité : « \teiHi[rend={italic}]{Vedesi poi nel secondo lo stemma emblematico della Casa Rezzonico, fiancheggiato da due elmi, ove son figurati alcuni Tritoni, che mezz’uomini, e mezzo pesci, alludono a Cavalieri del Sacro Ordine potenti in Mare ed in terra.} »}}.
\end{teiQuote}
                  
\end{teiCit}
                  \teiP[xmlid={p40-3}]{À vrai dire, il ne s’agit pas exactement des armes Rezzonico mais d’une composition qui en est dérivée. Sous le chapeau de prélat, sur un bouclier avec les barres (librement traitées en chevrons), ont été disposés l’un au-dessus de l’autre les trois meubles restants : la croix grecque, la tour, l’aigle bicéphale couronnée.}
                  \teiP[xmlid={p41-2}]{On a parfois qualifié Piranèse de « gothique » au sens large de non classique ou même d’étrangement inquiétant\teiNote[n={33},place={foot},xmlid={ftn65-5}]{\teiP[xmlid={p-306}]{Exemples dans Pierluigi Panza, \teiHi[rend={italic}]{Piranesi architetto},\teiHi[rend={italic}]{ op. cit.}, p. 68 ; ou dans Henri Focillon, \teiHi[rend={italic}]{G. B. Piranesi}, Paris, Librairie Renouard-Henri Laurens, 1928, p. 21 et p. 302.}}. L’héraldique peut entrer dans cette catégorie. Si cet artiste fut un maître de l’antique il fut aussi un maître de l’héraldique. Dans son œuvre de graveur ou d’architecte, ces deux imaginaires à priori divergents s’unissent avec une inventivité toujours renouvelée pour la célébration du pape Rezzonico et de son neveu Giovan Battista, le grand prieur de Malte. L’idée répandue que le \teiHi[rend={small-caps}]{xviii}\teiHi[rend={sup}]{e} siècle a sonné le glas de l’héraldique est une idée à revoir.}
                  \begin{teiFigure}[xmlid={figure01-6}]

                    \teiGraphic[url={../icono/br/Ch07\_Loskoutoff\_2/Fig98.JPG}]
                    \teiHead[xmlid={figure-96-giovanni-battista-piranesi-le-antichita-romane-1756-detail-du-frontispice-1-er-etat-accademia-di-belle-arti-di-firenze-001}]{Figure 96. Giovanni Battista Piranesi, \teiHi[rend={italic}]{Le antichità romane}, 1756, détail du frontispice, 1\teiHi[rend={sup}]{er} état, Accademia di Belle Arti di Firenze.}
                  
\end{teiFigure}
                  \begin{teiFigure}[xmlid={figure02-6}]

                    \teiGraphic[url={../icono/br/Ch07\_Loskoutoff\_2/Fig99.JPG}]
                    \teiHead[xmlid={figure-97-giovanni-battista-piranesi-le-antichita-romane-1756-detail-du-frontispice-2-e-etat-bav-001}]{Figure 97. Giovanni Battista Piranesi, \teiHi[rend={italic}]{Le antichità romane}, 1756, détail du frontispice, 2\teiHi[rend={sup}]{e} état, BAV.}
                  
\end{teiFigure}
                  \begin{teiFigure}[xmlid={figure03-6}]

                    \teiGraphic[url={../icono/br/Ch07\_Loskoutoff\_2/Fig100.JPG}]
                    \teiHead[xmlid={figure-98-giovanni-battista-piranesi-le-antichita-romane-1756-detail-du-frontispice-3-e-etat-collection-museum-boijmans-van-beuningen-rotterdam-001}]{Figure 98. Giovanni Battista Piranesi, \teiHi[rend={italic}]{Le antichità romane}, 1756, détail du frontispice, 3\teiHi[rend={sup}]{e} état, collection Museum Boijmans Van Beuningen, Rotterdam.}
                  
\end{teiFigure}
                  \begin{teiFigure}[xmlid={figure04-6}]

                    \teiGraphic[url={../icono/br/Ch07\_Loskoutoff\_2/Fig101.jpg}]
                    \teiHead[xmlid={figure-99-giovanni-battista-piranesi-vasi-candelabri-1778-frontispice-du-vol-2-001}]{Figure 99. Giovanni Battista Piranesi, \teiHi[rend={italic}]{Vasi, candelabri, …}, 1778, frontispice du vol. 2.}
                  
\end{teiFigure}
                  \begin{teiFigure}[xmlid={figure05-6}]

                    \teiGraphic[url={../icono/br/Ch07\_Loskoutoff\_2/Fig102.JPG}]
                    \teiHead[xmlid={figure-100-dessin-dun-rhyton-album-ghezzi-bav-ott-lat-3105-001}]{Figure 100. Dessin d’un rhyton, album Ghezzi, BAV, Ott. lat. 3105.}
                  
\end{teiFigure}
                  \begin{teiFigure}[xmlid={figure06-6}]

                    \teiGraphic[url={../icono/br/Ch07\_Loskoutoff\_2/Fig103.JPG}]
                    \teiHead[xmlid={figure-101-giovanni-battista-piranesi-detail-du-tombeau-de-cecilia-metella-vedute-di-roma-collection-particuliere-yvan-loskoutoff-001}]{Figure 101. Giovanni Battista Piranesi, détail du tombeau de Cecilia Metella, \teiHi[rend={italic}]{Vedute di Roma}, collection particulière, Yvan Loskoutoff.}
                  
\end{teiFigure}
                  \begin{teiFigure}[xmlid={figure07-6}]

                    \teiGraphic[url={../icono/br/Ch07\_Loskoutoff\_2/Fig104.jpg}]
                    \teiHead[xmlid={figure-102-giovanni-battista-piranesi-dessin-des-armes-de-clement-xiii-rezzonico-the-elisha-whittelsey-collection-the-elisha-whittelsey-fund-1952-001}]{Figure 102. Giovanni Battista Piranesi, dessin des armes de Clément XIII Rezzonico. @ The Elisha Whittelsey Collection, The Elisha Whittelsey Fund, 1952.}
                  
\end{teiFigure}
                  \begin{teiFigure}[xmlid={figure08-5}]

                    \teiGraphic[url={../icono/br/Ch07\_Loskoutoff\_2/Fig105.JPG}]
                    \teiHead[xmlid={figure-103-giovanni-battista-piranesi-de-romanorum-magnificentia-et-architectura-1761-frontispice-001}]{Figure 103. Giovanni Battista Piranesi, \teiHi[rend={italic}]{De Romanorum magnificentia et architectura}, 1761, frontispice.}
                  
\end{teiFigure}
                  \begin{teiFigure}[xmlid={figure09-5}]

                    \teiGraphic[url={../icono/br/Ch07\_Loskoutoff\_2/Fig106.JPG}]
                    \teiHead[xmlid={figure-104-giovanni-battista-piranesi-le-rovine-del-castello-dellacqua-giulia-1761-cul-de-lampe-yale-university-art-gallery-001}]{Figure 104. Giovanni Battista Piranesi, \teiHi[rend={italic}]{Le rovine del castello dell’acqua Giulia, }1761, cul-de-lampe, Yale University Art Gallery.}
                  
\end{teiFigure}
                  \begin{teiFigure}[xmlid={figure10-5}]

                    \teiGraphic[url={../icono/br/Ch07\_Loskoutoff\_2/Fig107.JPG}]
                    \teiHead[xmlid={figure-105-giovanni-battista-piranesi-lapides-capitolini-sive-fasti-consularesque-romanorum-1762-frontispice-bibliotheques-mcgill-001}]{Figure 105. Giovanni Battista Piranesi, \teiHi[rend={italic}]{Lapides Capitolini sive fasti consularesque Romanorum, }1762, frontispice, bibliothèques McGill.}
                  
\end{teiFigure}
                  \begin{teiFigure}[xmlid={figure11-5}]

                    \teiGraphic[url={../icono/br/Ch07\_Loskoutoff\_2/Fig108.JPG}]
                    \teiHead[xmlid={figure-106-giovanni-battista-piranesi-antichita-dalbano-e-di-castel-gandolfo-1764-frontispice-001}]{Figure 106. Giovanni Battista Piranesi, \teiHi[rend={italic}]{Antichità d’Albano e di Castel Gandolfo}, 1764, frontispice.}
                  
\end{teiFigure}
                  \begin{teiFigure}[xmlid={figure12-5}]

                    \teiGraphic[url={../icono/br/Ch07\_Loskoutoff\_2/Fig109.JPG}]
                    \teiHead[xmlid={figure-107-giovanni-battista-piranesi-diverse-maniere-dadornare-i-cammini-1769-frontispice-bnf-001}]{Figure 107. Giovanni Battista Piranesi, \teiHi[rend={italic}]{Diverse maniere d’adornare i cammini,} 1769, frontispice, BNF.}
                  
\end{teiFigure}
                  \begin{teiFigure}[xmlid={figure13-4}]

                    \teiGraphic[url={../icono/br/Ch07\_Loskoutoff\_2/Fig110.JPG}]
                    \teiHead[xmlid={figure-108-giovanni-battista-piranesi-portail-du-prieure-de-malte-rome-detail-001}]{Figure 108. Giovanni Battista Piranesi, portail du prieuré de Malte, Rome (détail).}
                  
\end{teiFigure}
                  \begin{teiFigure}[xmlid={figure14-4}]

                    \teiGraphic[url={../icono/br/Ch07\_Loskoutoff\_2/Fig111.JPG}]
                    \teiHead[xmlid={figure-109-giovanni-battista-piranesi-portail-du-prieure-de-malte-rome-detail-001}]{Figure 109. Giovanni Battista Piranesi, portail du prieuré de Malte, Rome (détail).}
                  
\end{teiFigure}
                  \begin{teiFigure}[xmlid={figure15-4}]

                    \teiGraphic[url={../icono/br/Ch07\_Loskoutoff\_2/Fig112.jpg}]
                    \teiHead[xmlid={figure-110-giovanni-battista-piranesi-facade-de-santa-maria-del-priorato-rome-detail-dipartimento-di-architettura-progetto-sapienza-universita-di-roma-paolo-marcoaldi-001}]{Figure 110. Giovanni Battista Piranesi, façade de Santa Maria del Priorato, Rome (détail), Dipartimento di Architettura. Progetto Sapienza Università di Roma, Paolo Marcoaldi.}
                  
\end{teiFigure}
                  \begin{teiFigure}[xmlid={figure16-4}]

                    \teiGraphic[url={../icono/br/Ch07\_Loskoutoff\_2/Fig113.JPG}]
                    \teiHead[xmlid={figure-111-giovanni-battista-piranesi-exterieur-de-labside-du-choeur-de-santa-maria-del-priorato-rome-detail-001}]{Figure 111. Giovanni Battista Piranesi, extérieur de l’abside du chœur de Santa Maria del Priorato, Rome (détail).}
                  
\end{teiFigure}
                  \begin{teiFigure}[xmlid={figure17-3}]

                    \teiGraphic[url={../icono/br/Ch07\_Loskoutoff\_2/Fig114.jpeg}]
                    \teiHead[xmlid={figure-112-giovanni-battista-piranesi-armes-de-giovan-battista-rezzonico-choeur-de-santa-maria-del-priorato-rome-001}]{Figure 112. Giovanni Battista Piranesi, armes de Giovan Battista Rezzonico, chœur de Santa Maria del Priorato, Rome.}
                  
\end{teiFigure}
                  \begin{teiFigure}[xmlid={figure18-3}]

                    \teiGraphic[url={../icono/br/Ch07\_Loskoutoff\_2/Fig115.JPG}]
                    \teiHead[xmlid={figure-113-giovanni-battista-piranesi-composition-heraldique-en-lhonneur-de-giovan-battista-rezzonico-santa-maria-del-priorato-rome-001}]{Figure 113. Giovanni Battista Piranesi, composition héraldique en l’honneur de Giovan Battista Rezzonico, Santa Maria del Priorato, Rome.}
                  
\end{teiFigure}
                
\end{teiDiv}
              
\end{teiDiv}
            
\end{teiDiv}
          
\end{teiElement}
        
\end{teiElement}
      
\end{teiElement}
      \begin{teiElement}[name={group},type={section1},xmlid={section1-002}]

        \teiHead[xmlid={deuxieme-partie-monuments-001}]{Deuxième partie : Monuments}
        \begin{teiElement}[name={group},type={chapter},data-page-title={La mise en signe de l’espace par les papes : pratiques et conflits},data-page-authors={Laurent Hablot},data-include-href={Ch08\_Mise\_en\_signe.xml},xmlid={chapter-005}]

          \begin{teiElement}[name={front}]

            \begin{teiDiv}[type={titlePage},xmlid={titlepage-010}]

              \teiP[rend={title-main},xmlid={p-307}]{La mise en signe de l’espace par les papes : pratiques et conflits}
              \teiP[rend={author-aut},xmlid={p-308}]{Laurent Hablot}
            
\end{teiDiv}
          
\end{teiElement}
          \begin{teiElement}[name={body},xmlid={body-010}]

            \teiP[xmlid={p1-10}]{Dès la fin du \teiHi[rend={small-caps}]{xii}\teiHi[rend={sup}]{e} siècle, et dans le prolongement probable de pratiques impériales plus anciennes, l’apposition de signes emblématiques dans des lieux symboliques de l’espace public permet de faire connaître et d’imposer l’autorité qui s’y exerce. Cet usage s’est appliqué à travers toute l’Europe et jusqu’à la fin des anciens régimes, comme en témoignent toujours d’innombrables traces, visibles ou documentées. Il se poursuit d’ailleurs aujourd’hui encore notamment à travers le déploiement dans l’espace public des drapeaux et des enseignes de l’État. Au Moyen Âge, ce marquage emblématique de l’espace politique correspond néanmoins à des droits précis très disputés et énonce des pouvoirs spécifiques. Il a donc été bien souvent l’occasion de conflits, joués à grande ou petite échelle, depuis les incessantes querelles de prééminences dans le décor héraldiques des églises paroissiales, si bien documentées dans la Bretagne des \teiHi[rend={small-caps}]{xv}\teiHi[rend={sup}]{e}-\teiHi[rend={small-caps}]{xvii}\teiHi[rend={sup}]{e} siècle\teiNote[n={1},place={foot},xmlid={ftn1-8}]{\teiP[xmlid={p-309}]{Les prééminences correspondent aux privilèges de représentation détenus par une autorité seigneuriale et comprenant notamment le droit d’apposer ses armoiries sur tels ou tels lieux et supports, en particulier dans les églises (vitraux, clefs de voûtes, enfeux, bancs, litres, etc.). Voir sur cette question Michel Nassiet, « Signes de parenté signes de seigneurie : un système idéologique », \teiHi[rend={italic}]{Mémoires de la société d’Histoire et d’Archéologie de Bretagne}, n\teiHi[rend={sup}]{o} 67, 1991, p. 175-232 ; Paul-François Broucke, \teiHi[rend={italic}]{Les prééminences héraldiques de la cathédrale de Quimper}, mémoire de master I, Université de Bretagne Occidentale, 2010 et \teiHi[rend={italic}]{La cathédrale de Saint-Pol de Léon : le chantier des parties orientales du }\teiHi[rend={small-caps italic}]{xi}\teiHi[rend={italic sup}]{e}\teiHi[rend={italic}]{ au }\teiHi[rend={small-caps italic}]{xvi}\teiHi[rend={italic sup}]{e}\teiHi[rend={italic}]{ siècle}, mémoire de master II, Université de Bretagne Occidentale, 2012.}}, jusqu’au niveau des espaces urbains, des territoires ou des États. Les sources médiévales et modernes témoignent amplement de ces pratiques de mises en signe et des moyens humains, financiers, techniques et artistiques qui leurs sont consacrés. Mais elles révèlent également les querelles qu’elles induisent et leurs conséquences, entraînant la dégradation, la diffamation, le bûchage, la destruction et parfois le rétablissement de signes héraldiques apposés sur des monuments et des objets symboliques de toutes sortes. D’ailleurs, les grands épisodes de l’histoire, révoltes, révolutions, guerres civiles ou de religions, se traduisent presque toujours par ces « guerres des signes » qui, par les traces qui en sont conservées, marquent encore très souvent notre environnement monumental historique\teiNote[n={2},place={foot},xmlid={ftn37-4}]{\teiP[xmlid={p-310}]{Voir en dernier lieu sur le sujet Laurent Hablot, « Le bris des armes. L’iconoclasme héraldique dans la société médiévale », dans Marc Gil, Pascale Charron et Ambre Vilain (dir.), \teiHi[rend={italic}]{La pensée du regard. Études d’histoire de l’art du Moyen Âge offertes à Christian Heck}, Turnhout, Brepols, 2016, p. 181-191 ; et Steven Thiry, \teiHi[rend={italic}]{Matter(s) of State. Heraldic Display and Discourse in the early Modern Monarchy (c. 1480-1650)}, Ostfildern, Thorbecke Verlag, « Heraldic Studies, 2 », 2018.}}.}
            \teiP[xmlid={p2-9}]{Dès lors qu’ils ont fait usage de signes héraldiques, les chefs de l’Église, princes temporels autant que spirituels, n’ont pas dérogé à cette règle pour baliser leurs possessions territoriales et affirmer leur pouvoir et leurs droits à travers le monde chrétien\teiNote[n={3},place={foot},xmlid={ftn38-4}]{\teiP[xmlid={p-311}]{Je remercie vivement Matteo Ferrari pour m’avoir indiqué de nombreuses et précieuses références et pour les riches échanges que nous avons pu avoir autour de cette question.}}. Dans ce cas toutefois, la mise en signe des biens de l’Église confond souvent, de façon presque inextricable, les intérêts du pontife et de sa famille avec ceux de l’Institution, dans la mesure où les emblèmes personnels des papes accompagnent, complètent ou remplacent les insignes de l’Église. De très nombreuses attestations documentées ou toujours visibles témoignent encore de cette emblématisation du territoire, spécialement en Italie. Les particularités historiques de cet espace politique, notamment la problématique de l’expansion concurrente de l’Empire et des États de l’Église, ainsi que du système des communes, autant que le statut singulier des papes, n’ont pas manqué d’entretenir d’actives politiques d’apposition d’emblèmes et de soulever de nombreuses contestations et d’importants conflits de marquage héraldique.}
            \teiP[xmlid={p3-9}]{Pour ouvrir ce dossier, qui demande à être bien plus largement et plus précisément considéré que dans ce bref article, il convient de revenir, à travers quelques cas documentés, sur les supports de cette mise en signe, sur l’emblématique développée par les papes et sur la législation qui encadre ces pratiques d’apposition d’armoiries.}
            \begin{teiDiv}[type={section1},xmlid={div01-7}]

              \teiHead[xmlid={les-supports-de-la-mise-en-signe-001}]{Les supports de la mise en signe}
              \teiP[xmlid={p4-9}]{Dans le cas qui nous intéresse ici, l’emblématisation de l’espace s’applique prioritairement aux cités. Elle se focalise sur quelques supports monumentaux symboliques, fréquemment sollicités, sur lesquels se concentre habituellement l’information héraldique et connus par les hommes de ce temps pour être les habituels relais de l’information politique\teiNote[n={4},place={foot},xmlid={ftn39-4}]{\teiP[xmlid={p-312}]{Vues de Sienne et de Florence. Florence, Biblioteca Medicea Laurenziana, Tempiano, n\teiHi[rend={sup}]{o} 3, planche 1, vers 1330.}}. Il s’agit bien évidemment des lieux de passage tels que les portes d’entrée des villes et les ponts, des résidences de pouvoir telles que les palais et avant tout le palais communal\teiNote[n={5},place={foot},xmlid={ftn40-4}]{\teiP[xmlid={p-313}]{Voir sur le sujet les travaux Matteo Ferrari, en dernier lieu \teiHi[rend={italic}]{La politica in figure. Temi, funzioni, attori della communicazione visiva nei Comuni lombardi (XII-XIV secolo)}, Rome, Viella, 2022 ; Maria Monica Donato, « “Ogni cosa è pieno d’arme”. Uno sguardo dall’esterno », dans Matteo Ferrari (dir.), \teiHi[rend={italic}]{L’Arme segreta. Araldica e storia dell’arte nel Medioevo (secoli XIII-XIV)}, Florence, Le Lettere, 2015, p. 267-283 ; Christoph Friedrich Weber, \teiHi[rend={italic}]{Zeichen der Ordnung und des Aufruhrs. Heraldische Symbolik in italienischen Stadtkommunen des Mittelalters}, Cologne, Böhlau, 2011. }}, des monuments symboliques comme le beffroi ou les principaux édifices religieux séculiers ou conventuels, des constructions évergétiques et symboliques telles que les fontaines, les halles et les marchés, et des constructions seigneuriales comme les moulins et les pressoirs. La mise en signe de ces supports est matérialisée, de façon pérenne, par des bas-reliefs sculptés, le plus souvent rehaussés de couleurs, ou par de très amples représentations héraldiques peintes. Elle est complétée, de manière plus éphémère, par tout un déploiement de signes mobiles, bannières et enseignes de toile ou de tôle, le plus souvent apposées sur les faîtages. Ce marquage est relayé par tout un ensemble d’objets mobiliers du pouvoir tels que le sceau ou la cloche, qui diffusent l’information visuelle ou sonore, le char de guerre, conservé dans l’église majeure de la cité et porteur de l’étendard et de la cloche, divers instruments symboliques du pouvoir tels que les masses ou les uniformes des agents publics, sans oublier les nombreux décors textiles comme la bannière ou le dais, déployés dans toutes les occasions de manifestations du pouvoir : entrées, fêtes patronales et liturgiques, funérailles, Te Deum, etc.}
              \teiP[xmlid={p5-9}]{Identifiés à leur tour par des signes héraldiques, les souverains pontifes s’insèrent progressivement dans cette politique de prise de possession symbolique de l’espace public.}
            
\end{teiDiv}
            \begin{teiDiv}[type={section1},xmlid={div02-7}]

              \teiHead[xmlid={quels-signes-pour-les-papes-et-leglise-001}]{Quels signes pour les papes et l’Église ?}
              \teiP[xmlid={p6-9}]{Évoquer la mise en signe de l’espace par les papes impose de s’interroger sur les emblèmes qu’ils apposent et notamment la concurrence qui s’exerce entre des armories institutionnelles, en l’occurrence celles de l’Église, des insignes de fonction, et les armoiries personnelles ou familiales des pontifes.}
              \teiP[xmlid={p7-9}]{Comme l’ont précisé les travaux d’Édouard Bouyé, les papes ne commencent à faire un usage assuré d’armoiries qu’à partir de Nicolas III (Giovanni Gaetano Orsini, pape de 1277-1280). Même si quelques occurrences laissent penser que ses prédécesseurs Clément IV (Gui Foucois, pape de 1265-1268) et Innocent IV (Sinibaldo de Fleschi, pape de 1245-1253) aient pu faire usage d’armoiries personnelles ou familiales\teiNote[n={6},place={foot},xmlid={ftn41-4}]{\teiP[xmlid={p-314}]{Édouard Bouyé, « Les armoiries pontificales à la fin du \teiHi[rend={small-caps}]{xiii}\teiHi[rend={sup}]{e} siècle : construction d’une campagne de communication », \teiHi[rend={italic}]{Médiévales}, n\teiHi[rend={sup}]{o} 44, 2003, p. 177-178, \teiRef[target={http://medievales.revues.org/938}]{http://medievales.revues.org/938} (consulté le 24 septembre 2024).}}.}
              \teiP[xmlid={p8-9}]{Passé le pontificat de Boniface VIII (Benedetto Caetani, pape de 1294-1303) qui opère, comme le souligne Édouard Bouyé, un « tournant dans le développement de l’héraldique pontificale » mais contribue aussi à la diffusion du portrait pontifical\teiNote[n={7},place={foot},xmlid={ftn42-4}]{\teiP[xmlid={p-315}]{Boniface VIII est aussi bien connu pour s’être exposé par portraits interposés. Voir à ce sujet Agostino Paravicini Bagliani, « Boniface VIII en images. Vision d'église et mémoire de soi », dans Dominic Olariu (dir.), \teiHi[rend={italic}]{Le portrait individuel. Réflexions autour d’une forme de représentation, }\teiHi[rend={small-caps italic}]{xiii}\teiHi[rend={italic sup}]{e}\teiHi[rend={italic}]{-}\teiHi[rend={small-caps italic}]{xv}\teiHi[rend={italic sup}]{e}\teiHi[rend={italic}]{ siècles}, Berne, Peter Lang, 2009, p. 65-81.}}, tous les pontifes du Moyen Âge, à quelques exceptions près, font usage de leurs armoiries personnelles et/ou familiales. La célèbre Loggia du Latran (fig. 114), réalisée à l’extrême fin du \teiHi[rend={small-caps}]{xiii}\teiHi[rend={sup}]{e} siècle, livre un des premiers témoignages documentés du marquage héraldique de l’espace par les papes. Connu par un relevé produit au \teiHi[rend={small-caps}]{xvi}\teiHi[rend={sup}]{e }siècle des fresques qui étaient figurées dans ce même monument\teiNote[n={8},place={foot},xmlid={ftn43-5}]{\teiP[xmlid={p-316}]{Giacomo Grimaldi, \teiHi[rend={italic}]{Instrumenta translationum}, Milan, Bibliothèque ambrosienne, ms. F 227 inf., copie de la scène de la bénédiction du cycle de fresques de la loggia du Latran.}}, ce décor se déploie sur les corniches et le balcon de la loggia depuis laquelle le pape bénit le peuple de Rome. Les célèbres armes des Caetani (d’or à deux jumelles ondées d’azur posées en bande) y alternent avec plusieurs insignes pontificaux : l’\teiHi[rend={italic}]{ombrellino} rayé jaune et rouge et les clefs entrecroisées timbrées d’une tiare.}
              \teiP[xmlid={p9-8}]{La cohabitation sur cet ensemble des insignes familiaux avec ceux de la fonction souligne la diffusion, dès la fin du \teiHi[rend={small-caps}]{xiii}\teiHi[rend={sup}]{e} siècle et sans doute antérieurement à l’usage des armoiries familiales, d’autres signes liés à l’institution ou à la charge pontificale. Parmi ceux-ci se retrouvent les insignes pontificaux (clefs, \teiHi[rend={italic}]{ombrellino} et tiare), les armoiries de l’Église (de gueules à deux clefs d’argents posées en sautoir) et la bannière ou gonfanon pontifical (de gueules à la croix d’argent cantonnée de quatre clefs en pal du même alias huit clefs posées en sautoir).}
              \teiP[xmlid={p10-8}]{Comme en témoignent divers exemples, l’ensemble de ces emblèmes, distribués selon une logique difficile à saisir, semble fournir un panel de signes dont la hiérarchie et les associations, entre eux et avec d’autres emblèmes, sont peut-être signifiantes et relatives, selon les lieux et les supports, à la nature du pouvoir ou des droits exercés, aux fonctions déléguées, etc.}
              \teiP[xmlid={p11-8}]{Cette panoplie permet donc aux papes d’occuper symboliquement l’espace visuel en tenant compte de l’expression et des droits respectifs des autres pouvoirs en exercice.}
            
\end{teiDiv}
            \begin{teiDiv}[type={section1},xmlid={div03-6}]

              \teiHead[xmlid={le-droit-dapposer-ses-armoiries-001}]{Le droit d’apposer ses armoiries}
              \teiP[xmlid={p12-8}]{Rappelons en effet que l’apposition des armoiries dans l’espace public correspond à un droit spécifique strictement réglementé. La performativité de ce « geste », symbole de prise de possession, de protection, d’exercice de droits, de juridiction et de \teiHi[rend={italic}]{memoria}, est bien établie, tant par les exemples de marquage liés à la domination d’un espace que par les contestations que ces appositions suscitent. De ce fait, l’exposition de signes héraldiques dans l’espace sacré et, dans l’espace public, le droit de « faire blason fenêtre », d’orner sa demeure ou des édifices publics d’armoiries « de particuliers », sont encadrés par un ensemble de règles coutumières régulièrement couchées par écrit dans les sources normatives ou institutionnelles des communautés, grands fiefs, communes ou cités. De nombreux exemples peuvent être cités à travers l’Europe des \teiHi[rend={small-caps}]{xiii}\teiHi[rend={sup}]{e}-\teiHi[rend={small-caps}]{xvi}\teiHi[rend={sup}]{e} siècles, à l’instar des villes de Poitiers ou de Lyon\teiNote[n={9},place={foot},xmlid={ftn44-5}]{\teiP[xmlid={p-317}]{À Poitiers, l’exposition des armoiries du maire sur des bâtiments publics nécessite également l’accord du corps de ville. En 1471, ce dernier autorise son maire à apposer ses armes sur la fontaine du pont Joubert qu’il a fait restaurer. Robert Favreau (éd.), \teiHi[rend={italic}]{Poitiers sous le règne de Louis XI, de 1471 à 1482. Registre de délibérations du corps de ville}, n\teiHi[rend={sup}]{o} 7, Poitiers, Société des antiquaires de l’Ouest, 2015, p. 186. Sur le sujet, voir l’article de Matteo Ferrari, « Héraldique et “mise en signe” de l’espace urbain en Poitou au Moyen Âge », \teiHi[rend={italic}]{L’héraldique dans la ville}, Stuttgart, Thorbecke, « Heraldic Studies, 5 », 2025.}}.}
              \teiP[xmlid={p13-8}]{La péninsule italienne constitue cependant un contexte singulier qui, en raison des concurrences politiques du pape, de l’empereur et de la commune et des multiples pouvoirs qui s’y exercent (\teiHi[rend={italic}]{magnati}, podestats et capitaines, partis, corporations, confréries, etc.), exacerbe les enjeux de ce marquage emblématique de l’espace. Il convient également d’y distinguer d’une part les villes des États pontificaux, théoriquement soumises à l’autorité du pape et, à ce titre, marquées de ses signes et de ceux de l’Église et, d’autre part, les communes d’obédience guelfe, qui par libre choix de leurs dirigeants et en signe d’alliance, affichent les armes de l’Église.}
              \teiP[xmlid={p14-7}]{Le cas spécifique des communes guelfes du nord de l’Italie a été envisagé dans un récent article dans lequel Gian Paolo Ermini rappelle les interdictions faites par Pérouse, (1297), Brescia (1290) et Sienne (1310) de figurer des armes de particuliers sur des lieux publics\teiNote[n={10},place={foot},xmlid={ftn45-5}]{\teiP[xmlid={p-318}]{Gian Paolo Ermini, « La décoration picturale de la grande salle du palais de la commune d’Orvieto », dans Torsten Hiltman et Miguel de Seixas (dir.), \teiHi[rend={italic}]{Heraldry in Medieval and Early Modern State Rooms}, Ostfildern, Thorbecke Verlag, « Heraldic Studies, 3 », 2020, p. 165-184, ici p. 167-171. Voir aussi d’autres exemples dans Matteo Ferrari, \teiHi[rend={italic}]{La politica in figure}, \teiHi[rend={italic}]{op. cit.}, p. 121.}}. Plusieurs de ces réglementations somptuaires concernent les armoiries pontificales. Ainsi, à Florence, les gouverneurs de la commune, cherchant à annihiler les pouvoirs des partis et factions, après avoir ordonné en 1312 la destruction de toutes les aigles impériales figurées dans la cité, décident, suite à l’épisode du gouvernement de Charles de Calabre en 1328, d’interdire la figuration et d’exiger la destruction de toute sculpture ou peinture figurant des armoiries personnelles. N’autorisant que les figurations du Christ, de la Vierge et de saints « \teiHi[rend={italic}]{vel armorum Sancte Romane Ecclesie vel domini Summi Pontificis vel Regis Karoli seu descendentium vel Regis Francie} », la commune spécifie encore que les armes d’Anjou ne doivent être associées avec aucune armoirie personnelle, de la commune, du peuple, d’aucune société ou parti guelfe (voir fig. 115, la prééminence des insignes pontificaux sur le Palazzo Vecchio). En 1390 encore, des lois somptuaires relatives à l’exposition publique des armoiries à Florence sont prises, confirmant la difficulté des autorités à faire appliquer ces lois\teiNote[n={11},place={foot},xmlid={ftn46-5}]{\teiP[xmlid={p-319}]{Voir sur le sujet Maria Monica Donato, « “Ogni cosa è pieno d’arme”… », art. cité. Voir aussi l’article de Matteo Ferrari, « Stemmi esposti. Presenze araldiche nei broletti lombardi », dans Matteo Ferrari (dir.), \teiHi[rend={italic}]{L’Arme segreta}, \teiHi[rend={italic}]{op. cit}., p. 91-107.}}.}
              \teiP[xmlid={p15-7}]{Les mêmes mesures restrictives sont appliquées à Prato en 1329, visant notamment l’exposition d’armoiries d’autres communes (Florence principalement). Ce qui n’empêche pas la cité de faire réaliser en 1336, sur le palais du Peuple, deux séries d’écus sculptés aux armes du pape (en fait de l’Église), du roi (c’est-à-dire d’Anjou), de la commune et du vicaire.}
              \teiP[xmlid={p16-7}]{En 1338, la commune de Gubbio suit l’exemple florentin et, après avoir imposé la destruction des emblèmes sur les édifices publics, n’autorise que les « \teiHi[rend={italic}]{arma sive insignia Sancte Romane Ecclesie ac etiam Summi Pontificis et domini Regis Roberti et domini Regis Francie et communis Eugubii }». C’est dans ce contexte qu’est sculptée en 1336 l’architrave du palais du Peuple toujours en place (fig. 116) figurant les armes de la ville sous le \teiHi[rend={italic}]{Capo d’Angiò}\teiNote[n={12},place={foot},xmlid={ftn47-5}]{\teiP[xmlid={p-320}]{Le \teiHi[rend={italic}]{capo} ou chef de l’écu est une pièce ajoutée à une armoirie initiale et qui intègre les armoiries d’une autorité de tutelle, plaçant le bénéficiaire dans une forme de patronage ou de sujétion honorable. En Italie, les personnes et les institutions affichent ainsi leur obédience à l’empereur ou au pape en retenant le \teiHi[rend={italic}]{capo dell’Impero} (d’or à l’aigle de sable) ou le \teiHi[rend={italic}]{capo d’Angio} (aux armes des rois de Naples, défenseurs de la papauté). Voir sur le sujet Michel Popoff, « Le “capo dello scudo” dans l’héraldique florentine, \teiHi[rend={small-caps}]{xiii-xvi}\teiHi[rend={sup}]{e} siècles », dans \teiHi[rend={italic}]{Brisures, augmentations et changements d'armoiries}, Bruxelles, Académie internationale d’héraldique, 1988, p. 251-255 et mon mémoire inédit d’HDR, \teiHi[rend={italic}]{Affinités héraldiques}, Michel Pastoureau (dir.), École pratique des hautes études, 2015, « Capo d’Angio Capo d’Impero, la mise en signe de la péninsule », p. 239 et suiv. Voir aussi Luisa Clotilde Gentile, « La diffusion du capo dell'Impero dans l’Italie du nord », \teiHi[rend={italic}]{Armas e troféus}, n\teiHi[rend={sup}]{o} 20, 2018, p. 109-138.}} (de gueules au mont à cinq coupeaux d’argent au chef d’azur semée de fleurs de lis d’or au lambel de gueules brochant) et celles du roi de Naples encadrant, en prééminence, les armes de l’Église dans la version de l’écu « de gueules à la croix d’argent cantonnée de huit clefs en sautoir et liées d’argent ». Mais d’autres formules semblent encore possibles. Ainsi, un relief de pierre de la cité d’Ascoli Piceno dans les Marches (fig. 117), expose, en prééminence, un parti aux armes de l’Église (une croix cantonnée de quatre sautoirs de clefs\teiNote[n={13},place={foot},xmlid={ftn48-6}]{\teiP[xmlid={p-321}]{Sur ces armes à la croix et aux clefs, voir Édouard Bouyé, « Les clefs de saint Pierre, sur la terre comme au ciel », dans Catalina Girbea, Christine Manigand, Laurent Hablot et \teiHi[rend={italic}]{al.} (dir.), \teiHi[rend={italic}]{Signes et couleurs des identités politiques}, Rennes, PUR, 2008, p. 275-312.}}) et du pape Boniface IX (1389-1404) accostées par celles de la cité et celles du capitaine du peuple.}
              \teiP[xmlid={p17-7}]{Dans les États de l’Église, en 1330, Amelia d’Ombrie, ville pontificale agitée par une forte opposition entre les partis guelfe et gibelin, est l’objet d’une violente révolte antipapale suscitée par la récente descente dans la Péninsule de l’empereur Louis IV en 1328 et qui s’incarne assurément par la destruction des signes pontificaux. Cette action motive en tous cas la rédaction de statuts interdisant la figuration de tous signes héraldiques dans la ville\teiNote[n={14},place={foot},xmlid={ftn49-6}]{\teiP[xmlid={p-322}]{Gian Paolo Ermini, « La décoration picturale… », art. cité, p. 170 et171.}}. Malgré cela, en 1357, les règlements pris à ce sujet par les Constitutions égidiennes rétablissent de nouveau la primauté des armes du pape et de l’Église sur celles des seigneurs locaux.}
              \teiP[xmlid={p18-7}]{Ces luttes politiques par images interposées ont marqué les contemporains tels Franco Sacchetti ou saint Bernardin de Sienne\teiNote[n={15},place={foot},xmlid={ftn50-6}]{\teiP[xmlid={p-323}]{Maria Monica Donato, « “Ogni cosa e pieno d’arme”… », art. cité, p. 21 et 23}}. Elles font d’ailleurs écho au célèbre attentat d’Anagni de 1303 durant lequel la prise de possession emblématique de l’espace avait joué un rôle essentiel. En effet, lorsque la troupe de Guillaume de Nogaret, qui avait symboliquement occupé le palais pontifical sans aucun doute marqué aux armes des Caetani et de l’Église, en balisant l’espace de bannières fleurdelisées, est chassée de la ville, la population s’en prend aussitôt aux emblèmes royaux\teiNote[n={16},place={foot},xmlid={ftn51-6}]{\teiP[xmlid={p-324}]{\teiHi[rend={italic}]{Chronique métrique de Godefroy Paris}, J.-A. Buchon (éd.), Paris, Verdière, 1827, v. 1995 à 2042 et 2124 et suiv. : « Ainsi ques l’ont (les conseillers du pape) reconforté / La bannière le roy osté / ont, qui sus le schastel estoit mise. / Chascun sa partie en a prise, / Et l’ont despecie trestoute / pièce ni ot qui ne fut route. »}}.}
            
\end{teiDiv}
            \begin{teiDiv}[type={section1},xmlid={div04-4}]

              \teiHead[xmlid={semparer-dune-ville-par-les-armes-quelques-exemples-001}]{S’emparer d’une ville par les « armes », quelques exemples}
              \teiP[xmlid={p19-7}]{Favorisés par les lois somptuaires des communes de leurs États dans le marquage de l’espace public, les papes doivent néanmoins, y compris dans leurs domaines, composer avec les autres pouvoirs en place et, selon les contextes, avoir recours aux différentes formules emblématiques disponibles : armoiries familiales, emblèmes de l’Église et concessions d’armoiries\teiNote[n={17},place={foot},xmlid={ftn52-6}]{\teiP[xmlid={p-325}]{Voir Pierre-Yves Le Pogam, \teiHi[rend={italic}]{De la « cité de Dieu » au « palais du pape » : les résidences pontificales dans la seconde moitié du }\teiHi[rend={small-caps italic}]{xiii}\teiHi[rend={italic sup}]{e}\teiHi[rend={italic}]{ siècle (1254-1304)}, Rome, École française de Rome, 2005, en particulier, sur la lutte entre Boniface VIII et la famille Colonna par armoiries interposées, p. 550-553.}}.}
              \teiP[xmlid={p20-7}]{L’historiographie de l’héraldique pontificale rapporte un exemple très précoce et assez singulier – mais très probablement imaginaire – qui introduit la pratique du partage ou de la concession de leurs armoiries par les papes à des communautés, solution qui offre un moyen plus subtil mais non moins efficace de présence héraldique dans l’espace politique. La \teiHi[rend={italic}]{Chronique} de Giovanni Villani mentionne ainsi les armes que Clément IV aurait concédées vers 1265 à ses partisans guelfes florentins\teiNote[n={18},place={foot},xmlid={ftn53-6}]{\teiP[xmlid={p-326}]{Sur ces armoiries, voir les notices 24 (\teiHi[rend={italic}]{Stemma della Parte Guelfa}), 25 (\teiHi[rend={italic}]{Sigillo della Parte Guelfa}) et 26 (\teiHi[rend={italic}]{Sigillo della Parte Guelfa}) du catalogue de l’exposition \teiHi[rend={italic}]{Dal Giglio al David}, \teiHi[rend={italic}]{Arte civica a Firenze fra Medioevo e Rinascimento}, Maria Monica Donato et Daniela Parenti (dir.), Florence, Giunti, 2013, p. 157-159.}} : « d’argent à l’aigle de gueules terrassant un dragon de sinople\teiNote[n={19},place={foot},xmlid={ftn54-6}]{\teiP[xmlid={p-327}]{« \teiHi[rend={italic}]{In questi tempi \lbrack{}1265\rbrack{} i Guelfi usciti di Firenze e dell’altre terre di Toscana }\lbrack{}...\rbrack{}\teiHi[rend={italic}]{ mandarono loro ambasciadori a papa Chimento, accio che gli raccomandasse al conte Carlo eletto re di Cicilia, e profferendosi al servigio di santa Chiesa; i quali dal detto papa furono ricevuti graziosamente, e proveduti di moneta e d’altri benifici; e volle il detto papa che per suo amore la parte guelfa di Firenze portasse sempre la sua arme propia in bandiera e in suggello, la quale era, ed è, il campo bianco con una aguglia vermiglia in su uno serpente verde, la quale portarono e tennero poi, e fanno insino a’ nostri presenti tempi} », dans \teiHi[rend={italic}]{Nuova Cronica}, Giuseppe Porta (éd.), Parma, Guanda, 1990-1991, livre VII, vol. 2.}} ». Aucune autre attestation ne vient confirmer l’usage de ces armes par le pape Clément IV mais une matrice de sceau du parti guelfe datée des années 1270 les porte bien et confirme leur usage par ce groupe politique très puissant. À partir de la fin du XIII\teiHi[rend={sup}]{e} siècle, ce signe du parti guelfe est régulièrement cité dans la quadrilogie des armes liées à la Commune : le lis de gueules de la commune, la croix de gueules du Popolo, les armes parties d’argent et de gueules de l’armée communale. Dans les décors, cet ensemble est d’ailleurs régulièrement distribué autour des armes de l’Église placées en prééminence, conformément aux statuts communaux cités plus haut. C’est par exemple le cas dans les décors héraldiques du Palazzo della Lana, peints vers 1350, ou dans l’allégorie héraldique de la ville peinte vers 1366 par Jacoeo di Cione dans le Palazzo dell’Arte dei guidici e dei notari, thème en grande partie repris vers 1360-1370 dans le Palazzo dell’Arte della Seta (voir en fig. 118 l’allégorie héraldique de la ville de Florence). Le célèbre monument funéraire de Piero da Farnese, réalisé vers 1376, offre, lui aussi, un bel exemple de ce déploiement héraldique de la cité distribué autour des armes de l’Église associées à celles du roi de Naples.}
              \teiP[xmlid={p21-7}]{Un riche exemple de ces concurrences de mise en signe est fourni par la ville de Bologne. En 1349, Galéas Visconti s’empare de la ville au nom de son oncle Jean Visconti, archevêque et seigneur de Milan, et, après une entrée militaire où se déploient les bannières aux armes de Visconti et d’Empire, il fait apposer ses emblèmes à travers toute la cité après y avoir fait effacer les lis, les lambels (armes d’Anjou), les clefs (armes de l’Église) et les lions. Le chroniqueur rapporte même la frayeur que ces monstres héraldiques jamais observés encore, guivre et aigle, suscite auprès des petits enfants\teiNote[n={20},place={foot},xmlid={ftn55-6}]{\teiP[xmlid={p-328}]{\teiHi[rend={italic}]{Annales mediolanenses ab anno 1230 usque ad annum 1402}, dans Lodovico Antonio Muratori (éd.), \teiHi[rend={italic}]{Rerum Italicarum Scriptores}, Mediolani, Typographia societatis palatinae in regiae curia, 1730, t. 16., coll. 721-722, CXVI.}}. En 1360, le pape reprend possession de la cité. À lire les chroniques, ce changement d’obédience est formalisé par une ample campagne d’apposition des armes de l’Église. Un char à la bannière de la ville est élevé à Saint-François tandis que des enseignes aux clefs et aux armes de la commune sont apposées sur le palais communal. Une véritable chasse aux figurations de l’aigle impériale est organisée à travers toute la ville\teiNote[n={21},place={foot},xmlid={ftn56-6}]{\teiP[xmlid={p-329}]{Hyeronimus de Bursellis, \teiHi[rend={italic}]{Cronica}, A. Sorbelli (éd.), Città di Castello, S. Lapi, 1912, p. 48 et Flavia Gramellini, \teiHi[rend={italic}]{Le Antichità di Bologna di Bartolomeo della Pugliola}, thèse de doctorat, Luciano Formisano (dir.), Bologne, Université de Bologne, 2008, p. 38.}}. La célèbre statue de Boniface VIII, apposée sur la façade du palais de la commune depuis 1301, est rénovée et repeinte aux armes de l’Église (« \teiHi[rend={italic}]{depicta cum armis ecclesie} », « \teiHi[rend={italic}]{si s’orno et depinsesi con l’armadura della Santa Chiesa}\teiNote[n={22},place={foot},xmlid={ftn57-6}]{\teiP[xmlid={p-330}]{Respectivement, Hyeronimus de Bursellis, \teiHi[rend={italic}]{Cronica}, \teiHi[rend={italic}]{op. cit}., p. 49 et Flavia Gramellini, \teiHi[rend={italic}]{op. cit}., p. 43.}} »). En 1410 encore, les comptes de la ville signalent le salaire de six peintres chargés de figurer les armes du pape sur le palais et sur les galons (\teiHi[rend={italic}]{guazalones}) du baldaquin aux armes du pape et de l’Église\teiNote[n={23},place={foot},xmlid={ftn58-6}]{\teiP[xmlid={p-331}]{Guido Filippini et Francesco Zucchini,\teiHi[rend={italic}]{ Miniatori e pittori a Bologna. Documenti del secolo XV}, Rome, Accademia nazionale dei Lincei, 1968, p. 172 et Hans W. Hubert, \teiHi[rend={italic}]{Der Palazzo Comunale von Bologna. Vom Palazzo della Biada zum Palatium Apostolicum}, Cologne, 1993, p. 181. Notons d’ailleurs qu’à cette date, le trône pontifical est sans doute vacant car Alexandre V meurt en juin 1409 et Jean XXIII n’est élu qu’en mai 1410. Les armes du pape et de l’Église seraient donc une forme redondante classique dans les textes médiévaux. Je remercie Matteo Ferrari pour cette précieuse remarque.}}. En 1425, le peintre Francesco Lola est payé pour les travaux de peinture et d’armoiries réalisés aux quatre principales portes de la ville et sur le palais communal, aux armes du pape et de l’Église et des seigneurs du gouvernement et de la commune de Bologne\teiNote[n={24},place={foot},xmlid={ftn59-6}]{\teiP[xmlid={p-332}]{Davide Ravaioli, « Pitture murali di Francesco Lola sulla facciata del Palazzo Comunale », \teiHi[rend={italic}]{Arte a Bologna}, n\teiHi[rend={sup}]{o} 6, 2007, p. 133-136 et principalement p. 132 et 133.}}.}
              \teiP[xmlid={p22-7}]{À Orvieto enfin, d’après les chroniqueurs dont Luca Manenti, en 1287, le pape Honorius IV Savelli et le podestat Bertoldo Orsini apposent leurs armes sur le pont de Rigochiaro tandis que le pape fait peindre le palais de la commune, sans doute à ses armes\teiNote[n={25},place={foot},xmlid={ftn60-6}]{\teiP[xmlid={p-333}]{« \teiHi[rend={italic}]{feci pingere il palazzo nuovo de signori VII} (palais de la commune) », Gian Paolo Ermini, « La décoration picturale… », art. cité, p. 180.}}. Le même chroniqueur rapporte que, sous Boniface VIII, sont ajoutées aux portes de la ville deux statues du pontife ornées de ses armes et de celles de l’Église\teiNote[n={26},place={foot},xmlid={ftn61-6}]{\teiP[xmlid={p-334}]{Voir Gian Paolo Ermini, « La campana del palazzo del popolo di Orvieto (1316) », dans Matteo Ferrari (dir.), \teiHi[rend={italic}]{L’Arme segreta, op. cit.}, p. 109-125, ici p. 115.}}. Cet auteur nous apprend encore qu’en 1304, après la mort de Boniface VIII, « \teiHi[rend={italic}]{fu ordinato che in tutto lo stato de Orvieto sopre le porte et palazzi si debia pingere l’arme de la communità de Orvieto} », ce qui implique l’effacement des armes pontificales et révèle en creux le déplaisir de la commune face à cette emprise emblématique du pape.}
              \teiP[xmlid={p23-7}]{À Viterbe, dans un style très proche du décor révélé par les fresques du Latran, les restes de la loggia pontificale, édifiée en 1267, offrent encore un magnifique exemple de la mise en signe des édifices pontificaux (fig. 119). Au-dessus des arcades, quatre registres sculptés alternent de haut en bas, une frise portant une succession de clefs en sautoir et de tiares, puis les armes fascées du capitaine du peuple Raniero Gatti, issu d’une des plus puissantes familles de la cité, et celles à l’aigle du parti guelfe ou de l’Empire, une frise constituée d’une succession de tiares, une frise alternant les armes des Gatti et celles de la ville au lion passant regardant brochant sur une palme\teiNote[n={27},place={foot},xmlid={ftn62-6}]{\teiP[xmlid={p-335}]{Un sceau de 1198 lui donne pour armes un lion tenant un arbre déraciné associé à la légende \teiHi[rend={italic}]{NON METUENS VERBUM LEO SUM QUI SIGNO VITERBUM}\teiHi[rend={small-caps}]{. }Cité dans Martin Roland et Andreas Zajic, « Les chartes médiévales enluminées dans les pays d’Europe centrale », \teiHi[rend={italic}]{Bibliothèque de l’école des Chartes}, n\teiHi[rend={sup}]{o} 169, 2011, p. 187 et n. 172.}}. Véritable synthèse des pouvoirs en exercice dans la cité, ce programme héraldique donne néanmoins la prééminence à l’Église de Clément IV qui ne fait pas usage d’armes personnelles. La construction d’une autre partie du palais sous le pontificat de Nicolas III Orsini (1277-1280) et sous l’autorité de son neveu Orso, podestat de Viterbe entre 1277 et 1280, leur permet en revanche de s’imposer face aux Gatti en apposant notamment leurs armes sur la célèbre fontaine du Sepale accompagnée de la mention « \teiHi[rend={italic}]{Papa Nicolaus regnabat in orbe beatus // Orso dominus prudens, clarus et potens} ». Leurs armes seront d’ailleurs effacées après la mort de Nicolas III et la révolte de la ville contre Orso\teiNote[n={28},place={foot},xmlid={ftn63-6}]{\teiP[xmlid={p-336}]{Gian Paolo Ermini, « La décoration picturale… », art. cité, p. 181 et 182.}}.}
              \teiP[xmlid={p24-5}]{Mais la présence pontificale s’impose à Viterbe par un moyen plus efficace encore. La première augmentation d’armoirie réellement accordée par un pape est précisément documentée pour Viterbe par une lettre d’armoiries datant du début du pontificat de Jean XXII (1316-1334) et constituant un des premiers exemples du genre si ce n’est le premier. En 1316, en effet, la ville reçoit du vicaire général du pape en Tuscie, Bernard de Coucy, une charte lui attribuant le titre de « gonfalonier de l’armée pontificale », la possession de la seigneurie de Montefiascone pour dix ans et de nouvelles armoiries : l’ajout au lion passant regardant sur la palme de ses armes d’un gonfanon pontifical dans la patte dextre du lion (de gueules à la croix d’argent cantonnée de quatre clefs posées en pal du même\teiNote[n={29},place={foot},xmlid={ftn64-6}]{\teiP[xmlid={p-337}]{Sur cette concession, sa bibliographie et celle sur les armes de Viterbe, voir Martin Roland et Andreas Zajic, « Les chartes médiévales enluminées… », art. cité, p. 183-186.}}). Cette composition, blasonnée et illustrée dans l’acte\teiNote[n={30},place={foot},xmlid={ftn65-6}]{\teiP[xmlid={p-338}]{Visible sur \teiRef[target={https://www.monasterium.net/mom/IlluminierteUrkunden/1316-03-11\_Viterbo/charter}]{https://www.monasterium.net/mom/IlluminierteUrkunden/1316-03-11\_Viterbo/charter} (consulté le 12/03/2025).}}, rappelle l’action de la cité en faveur du pape et de son représentant dans la reconquête des territoires du patrimoine de Saint-Pierre ayant fait sécession après la mort de Clément V. Il n’est d’ailleurs pas impossible que la bannière pontificale tenue par le lion soit venue remplacer dans ces armes une bannière impériale (à la croix) accordée par les empereurs au milieu du \teiHi[rend={small-caps}]{xii}\teiHi[rend={sup}]{e} siècle. Cette charte fut même peinte, à une date inconnue, sur les murs du palais communal de Viterbe où elle se voyait encore au milieu du \teiHi[rend={small-caps}]{xviii}\teiHi[rend={sup}]{e} siècle, établissant publiquement la présence du pape dans le paysage visuel de la ville\teiNote[n={31},place={foot},xmlid={ftn66-5}]{\teiP[xmlid={p-339}]{\teiHi[rend={italic}]{Ibid}., p. 185 et notes, rappellent que les sources du \teiHi[rend={small-caps}]{xv}\teiHi[rend={sup}]{e} siècle précisent que ce gonfanon pontifical remplace une bannière impériale concédée à Viterbe par Frédéric I\teiHi[rend={sup}]{er} Barberousse en 1167 ou 1170 et confirmée le 19 mars 1172 à Sienne par le légat impérial. Mais les auteurs ont établi que l’acte en question de 1172 (Viterbe, Archivio comunale, Pergamene n\teiHi[rend={sup}]{o} 11) porte seulement au dos la confirmation de la possession des biens dont l’empereur avait gratifié la ville au moyen du rituel de l’étendard impérial « \teiHi[rend={italic}]{per vexillum imperiale} ». Voir aussi Christoph Friedrich Weber, \teiHi[rend={italic}]{op. cit.}, \lbrack{}n. 170\rbrack{}, p. 70-73, 173 et \teiHi[rend={italic}]{passim} (résumé p. 507-521) qui indique que Crémone, investie de ses biens par l’empereur Henri VI en 1195 par une bannière (de gueules à la croix d’argent), a en effet ajouté cet emblème dans ses armes. }}.}
              \teiP[xmlid={p25-5}]{Cette pratique de concession de l’armoirie de l’Église permet à l’évidence une voie médiane moins problématique que l’apposition des armes sur les monuments publics. Un exemple plus tardif et attestant de la diffusion de la pratique au-delà des États de l’Église est fourni par la ville de Salon-de-Provence. La cité avait primitivement reçu pour armes celles de l’Empire, d’or à l’aigle bicéphale de sable, modifiées par l’évêque Eustache de Lévis, devenu seigneur de Salon lors de son rattachement au royaume en 1486. Après avoir fait arracher les armes impériales et les avoir remplacées par celles du roi, l’évêque avait alors imposé ses propres armes à la ville : d’or à trois chevrons de sable au lambel de gueules. À la demande des habitants, se plaignant de cette soumission héraldique de la ville, le pape Innocent VIII (Giovanni Battista Cybo, 1481 à 1492) donne de nouvelles armes à Salon par une bulle \teiHi[rend={italic}]{motu proprio} datée du 12 décembre 1488 : coupé de gueules à deux clefs d’argent et d’or à l’aigle de sable\teiNote[n={32},place={foot},xmlid={ftn67-5}]{\teiP[xmlid={p-340}]{Arthur de Boislile, « Rapport de M. de Boislile sur une communication de M. Destandau », \teiHi[rend={italic}]{Bulletin historique et philologique du Comité des travaux historiques et scientifiques}, 1905, p. 356-359, https://gallica.bnf.fr/ark:/12148/bpt6k5721727m. L’acte était conservé aux archives municipales sous les cotes A.A 3, n\teiHi[rend={sup}]{o} 1 et D.D. 7, fol. 31 et 51.}}. Par lettres d’avril 1537, François I\teiHi[rend={sup}]{er} modifiera de nouveau ces armes\teiNote[n={33},place={foot},xmlid={ftn68-5}]{\teiP[xmlid={p-341}]{Voir dans mon mémoire inédit d’HDR, \teiHi[rend={italic}]{Affinités héraldiques}, « Les concessions dans le royaume de France ».}}.}
              \teiP[xmlid={p26-5}]{Le marquage emblématique de l’espace par les papes s’inscrit bien dans le cadre de pratiques héraldiques observables dans l’ensemble de l’Occident entre la fin du XIII\teiHi[rend={sup}]{e} siècle et celle du XVI\teiHi[rend={sup}]{e} siècle\teiNote[n={34},place={foot},xmlid={ftn69-5}]{\teiP[xmlid={p-342}]{Voir la base de données ARMMA (Armorial monumental du Moyen Âge) : \teiRef[target={https://armma.saprat.fr}]{https://armma.saprat.fr} (consultée le 5 novembre 2024).}}. En France comme dans l’Empire et les autres royaumes d’Europe, qui restent tous la somme de pouvoirs partagés et négociés, ces mises en signes construisent et exposent de très subtils équilibres entre les pouvoirs en place. Elles sont finalement, lorsqu’elles résultent d’un commun accord et du respect des directives juridiques locales, une preuve éclatante de la concorde qui règne dans l’espace concerné. Mais les enjeux pratiques et honorifiques que ces figurations recouvrent font également de ces manifestations héraldiques un \teiHi[rend={italic}]{casus belli} exemplaire, première cible des conflits et des contestations.}
              \teiP[xmlid={p27-5}]{Cet intérêt ou ce rejet pour les expositions publiques d’armoiries attestent d’ailleurs de la culture héraldique commune des hommes de ce temps, capables à l’évidence de reconnaître ces signes, de repérer et d’interpréter les informations qui leur sont ainsi communiquées et de les contester et de les modifier si nécessaire.}
              \teiP[xmlid={p28-5}]{Le cas pontifical reste néanmoins une forme d’exception, tant à cause des particularités du système politique de l’Italie médiévale qu’en raison de la spécificité de la fonction papale et du statut très singulier de l’institution ecclésiale. La cohabitation de signes de différents statuts dans les manifestations emblématiques des papes en témoigne. Ces programmes permettent notamment de suivre les évolutions de l’image emblématique du pape, depuis la bannière de l’Église et les insignes de la fonction jusqu’aux armoiries des dynasties pontificales. Une analyse plus fine de l’utilisation de ces différents registres de signes reste à faire, en multipliant les études de cas finement contextualisées pour tenter de déterminer les motivations et les fonctions de tel ou tel signe et la valeur de leur association. L’analyse de ces pratiques dans le contexte particulier des États pontificaux comparés aux villes guelfes, dans le cadre de l’exil avignonnais ou celui du Grand Schisme et des conflits entre deux papautés apporterait également d’importantes informations.}
              \teiP[xmlid={p29-5}]{Compte-tenu de l’inflexion des discours relatifs aux armoiries souveraines, que des légendes héraldiques et des pratiques juridiques contribuent à sacraliser et à protéger sous le régime de la lèse-majesté\teiNote[n={35},place={foot},xmlid={ftn70-3}]{\teiP[xmlid={p-343}]{Voir Laurent Hablot, « Sacralisation of the royal coats of arms in Europe in the Middle Ages », dans Monserrat Herrero, Jaume Aurell, Angela Concetta Micelli Stout (dir.), \teiHi[rend={italic}]{Political theology in Medieval and Early Modern Europe. Discourses, Rites and Representations}, Turnhout, Brepols, 2017, p. 313-336.}}, il serait également pertinent de s’interroger plus avant sur la nature sacrée des armes de l’Église et du pape\teiNote[n={36},place={foot},xmlid={ftn71-3}]{\teiP[xmlid={p-344}]{En 1408, encore, Monstrelet évoque les envoyés de Pedro de Luna, anti-pape en Avignon sous le nom de Benoît XIII et qui refuse d’admettre la soustraction d’obédience décidée par Charles VI, qui sont publiquement diffamés à Paris. Déguisés en évêques, portant les armes renversées de leur maître autour du cou, ils sont promenés sur un tombereau, puis mis au pilori dans la cour du palais de la Cité, chapeautés de mitres inscrites : « eulx cy sont desloyaux à l’Eglise et au Roy ». Enguerrand de Monstrelet, \teiHi[rend={italic}]{Chronique}, Louis Douët d’Arcq (éd.), Paris, V\teiHi[rend={sup}]{e} Jules Renouard, 1857-1862, 6 vol., t. I, chap. 43, p. 264 et 265.}}.}
              \begin{teiFigure}[xmlid={figure01-7}]

                \teiGraphic[url={../icono/br/Ch08\_Hablot/Fig116.jpg}]
                \teiHead[xmlid={figure-114-giacomo-grimaldi-instrumenta-translationum-milan-bibliotheque-ambrosienne-ms-f-227-inf-copie-de-la-scene-de-la-benediction-du-cycle-de-fresques-de-la-loggia-du-latran-001}]{Figure 114. Giacomo Grimaldi, \teiHi[rend={italic}]{Instrumenta translationum}, Milan, bibliothèque Ambrosienne, ms. F 227 inf., copie de la scène de la bénédiction du cycle de fresques de la loggia du Latran.}
              
\end{teiFigure}
              \begin{teiFigure}[xmlid={figure02-7}]

                \teiGraphic[url={../icono/br/Ch08\_Hablot/Fig117.jpg}]
                \teiHead[xmlid={figure-115-la-preeminence-des-insignes-pontificaux-sur-le-palazzo-vechio-001}]{Figure 115. La prééminence des insignes pontificaux sur le Palazzo Vechio.}
              
\end{teiFigure}
              \begin{teiFigure}[xmlid={figure03-7}]

                \teiGraphic[url={../icono/br/Ch08\_Hablot/Fig118.jpg}]
                \teiHead[xmlid={figure-116-architrave-du-palais-du-peuple-de-gubbio-001}]{Figure 116. Architrave du palais du Peuple de Gubbio.}
              
\end{teiFigure}
              \begin{teiFigure}[xmlid={figure04-7}]

                \teiGraphic[url={../icono/br/Ch08\_Hablot/Fig119.jpg}]
                \teiHead[xmlid={figure-117-relief-de-pierre-de-la-cite-dascoli-piceno-001}]{Figure 117. Relief de pierre de la cité d’Ascoli Piceno.}
              
\end{teiFigure}
              \begin{teiFigure}[xmlid={figure05-7}]

                \teiGraphic[url={../icono/br/Ch08\_Hablot/Fig120.jpg}]
                \teiHead[xmlid={figure-118-allegorie-heraldique-de-la-ville-de-florence-vers-1360-1370-par-chini-galileo-fin-du-xix-e-siecle-palazzo-dellarte-della-seta-florence-001}]{Figure 118. Allégorie héraldique de la ville de Florence, vers 1360-1370, par Chini Galileo (fin du \teiHi[rend={small-caps}]{xix}\teiHi[rend={sup}]{e} siècle), Palazzo dell’Arte della Seta, Florence.}
              
\end{teiFigure}
              \begin{teiFigure}[xmlid={figure06-7}]

                \teiGraphic[url={../icono/br/Ch08\_Hablot/Fig121.jpg}]
                \teiHead[xmlid={figure-119-loggia-de-viterbe-001}]{Figure 119. Loggia de Viterbe.}
              
\end{teiFigure}
            
\end{teiDiv}
          
\end{teiElement}
        
\end{teiElement}
        \begin{teiElement}[name={group},type={chapter},data-page-title={La collégiale Saint-Jean-Baptiste puis Saint-Louis de Castelnau-Bretenoux : un lien étroit et discret avec la papauté},data-page-authors={Emmanuel Moureau},data-include-href={Ch09\_Collegiale\_Saint\_Jean\_Baptiste.xml},xmlid={chapter-006}]

          \begin{teiElement}[name={front}]

            \begin{teiDiv}[type={titlePage},xmlid={titlepage-011}]

              \teiP[rend={title-main},xmlid={p-345}]{La collégiale Saint-Jean-Baptiste puis Saint-Louis de Castelnau-Bretenoux : un lien étroit et discret avec la papauté}
              \teiP[rend={author-aut},xmlid={p-346}]{Emmanuel Moureau}
            
\end{teiDiv}
          
\end{teiElement}
          \begin{teiElement}[name={body},xmlid={body-011}]

            \teiP[xmlid={p1-11}]{Si le château féodal de Castelnau-Bretenoux marque le paysage du nord du Lot\teiNote[n={1},place={foot},xmlid={ftn1-9}]{\teiP[xmlid={p-347}]{Commune de Prudhomat.}}, imposant vaisseau de pierres rouges qui rappelle encore l’important pouvoir de ses seigneurs, l’humble collégiale Saint-Louis\teiNote[n={2},place={foot},xmlid={ftn33}]{\teiP[xmlid={p-348}]{Le vocable de Saint-Louis s’est substitué à celui primitif de Saint-Jean-Baptiste après le don à la fin du \teiHi[rend={small-caps}]{xvi}\teiHi[rend={sup}]{e} siècle d’un bras reliquaire de saint Louis par Louise de Bretagne, épouse de Guy de Castelnau et dame d’atour de Catherine de Médicis.}}, blottie aux pieds de ses remparts, passe inaperçue. Elle conserve pourtant un riche décor héraldique supporté par des clés de voûte, des vitraux et surtout les imposantes stalles sculptées. Sur ces dernières se sont glissées de discrètes allusions à deux bienfaiteurs du chapitre : le pape Jules II et l’un de ses cardinaux, François Guillaume de Clermont-Lodève.}
            \begin{teiDiv}[type={section1},xmlid={div01-8}]

              \teiHead[xmlid={une-famille-et-une-collegiale-001}]{Une famille et une collégiale}
              \teiP[xmlid={p2-10}]{La famille des Castelnau est connue en Quercy depuis la donation à l’abbaye de Cluny du monastère de Beaulieu-sur-Dordogne par Hugues de Castelnau en 1076\teiNote[n={3},place={foot},xmlid={ftn34-2}]{\teiP[xmlid={p-349}]{Archives départementales du Lot, Cahors : F 365. Pour le texte, voir Maximin Deloche, \teiHi[rend={italic}]{Cartulaire de l’abbaye de Beaulieu (en Limousin)}, Paris, Imprimerie impériale, 1859, p. XXIV. }}. Les barons de Castelnau ont bâti vers 1100 un premier château, remplacé au \teiHi[rend={small-caps}]{xiii}\teiHi[rend={sup}]{e} siècle par l’imposante forteresse qui marque toujours le paysage aujourd’hui. En 1280, après d’âpres luttes, les Castelnau se sont émancipés de la tutelle des vicomtes de Turenne en obtenant le privilège de dépendre directement du roi de France. En 1297, Raymond de Calmont d’Olt, évêque de Rodez et héritier de la seigneurie familiale depuis le décès de son frère Bégon, institue comme héritiers universels ses petits-neveux Hugues de Castelnau et Pierre Pelet\teiNote[n={4},place={foot},xmlid={ftn35-2}]{\teiP[xmlid={p-350}]{Hugues est le fils d’Alaisie de Calmont (fille de Bégon) et de Matfred de Castelnau ; Pierre le fils d’Alix de Calmont (fille de Bégon) et de Raymond Pellet. Voir Hippolyte de Barrau, \teiHi[rend={italic}]{Documents historiques et généalogiques sur les familles et les hommes remarquables du Rouergue dans les temps anciens et modernes}, Rodez, impr. de N. Ratery, 1853-1860, 4 vol., t. 1, p. 385.}}. Ce dernier est mort avant son père Raymond, qui lègue en 1315 par son testament tous ses biens à Hugues de Castelnau. En 1395, Jean I\teiHi[rend={sup}]{er} de Castelnau institue comme légataire universel, aux préjudices des enfants de sa fille, son neveu Pons de Caylus\teiNote[n={5},place={foot},xmlid={ftn36-4}]{\teiP[xmlid={p-351}]{Il est le fils de Déodat de Caylus et d’Hélène de Castelnau. Voir Hyppolite de Barrau, \teiHi[rend={italic}]{op. cit.}, p. 386.}}, issu d’une famille du Rouergue, à condition que celui-ci substitue à ses armes celles des Castelnau-Calmont. Pons épouse le 28 janvier 1401 Bourguine de Clermont-Lodève. Le couple a six enfants, dont Pons, qui relève le nom et les armes des Clermont-Lodève, Jean, évêque de Cahors (1438-1460) et Antoine, qui poursuit le lignage paternel. Il épouse en 1436 Catherine de Chauvigny qui lui donne notamment Guy, évêque de Cahors puis de Périgueux et Jean II de Castelnau. Du mariage de ce dernier avec Marie de Culant\teiNote[n={6},place={foot},xmlid={ftn37-5}]{\teiP[xmlid={p-352}]{L’épouse de Jean II de Castelnau est souvent prénommée Anne. Toutefois, elle est figurée sur un vitrail de la collégiale de Prudhomat avec sa sainte patronne Marie Jacobé. }} naissent Jacques et Jean de Castelnau-Calmont.}
              \teiP[xmlid={p3-10}]{Dans son testament, rédigé le 5 février 1505, Jean II de Castelnau exprime son désir et celui de son épouse de fonder une collégiale avec huit chanoines dirigés par un doyen dont les prières quotidiennes soutiendront les âmes de ses parents, de son épouse et de lui-même ainsi que de l’ensemble de leur lignage. La bulle de fondation du pape Jules II précise que le baron de Castelnau souhaite ainsi accumuler des trésors dans le ciel et rendre pérenne dans les cieux et pour l’éternité son passage sur la terre\teiNote[n={7},place={foot},xmlid={ftn38-5}]{\teiP[xmlid={p-353}]{Archives nationales, ms. Fr. 23154, fol. 273v.}}. La nouvelle fondation est placée sous le vocable et la protection de saint Jean-Baptiste, patron de Jean II. Cette nouvelle communauté religieuse remplacera les six chapellenies perpétuelles, destinées à six prêtres, que le baron a auparavant instituées au sein de la chapelle de son château. Les chanoines doivent prendre place dans un premier temps au sein de la chapelle castrale\teiNote[n={8},place={foot},xmlid={ftn39-5}]{\teiP[xmlid={p-354}]{La chapelle du château de Castelnau-Bretenoux existe encore. Des stalles sont disposées de part et d’autre de l’autel, témoins de la présence des chanoines de la collégiale avant la construction de cette dernière.}}, en attendant que soit construit un édifice spécifique et plus digne contre les murailles de ce dernier. Ce type de fondation ecclésiastique est assez surprenant en Quercy. Il n’existe en effet que très peu de collégiales séculières dans cette ancienne province qui n’a pas été – contrairement à la plupart des autres territoires du royaume –, entre le \teiHi[rend={small-caps}]{xi}\teiHi[rend={sup}]{e} et le début du \teiHi[rend={small-caps}]{xiv}\teiHi[rend={sup}]{e} siècle, un territoire propice à ce type d’établissement religieux, et il n’en existe aucune fondée par de simples laïcs. Plusieurs raisons expliquent ce phénomène : la forte présence des moines bénédictins, aux anciennes et puissantes maisons, assortie à un contexte politique agité sur plusieurs décennies et à une crise spirituelle qui conduit à un développement d’hérésies et d’un sentiment de méfiance vis-à-vis de l’Église. Il faut attendre la papauté d’Avignon pour que soit fondé le premier chapitre de chanoines de ce type, Saint-Étienne du Tescou à Montauban, en 1318. Le cardinal des Prés fonde par la suite le chapitre de Montpezat entre 1323 et 1343, Bernard Carit, évêque d’Évreux, la modeste collégiale Notre-Dame de Grâce de Puylaroque en 1366 et enfin en 1425 Bernard de Maumont, évêque de Tulle, érige en collégiale les prêtres desservants Saint-Sauveur de Rocamadour\teiNote[n={9},place={foot},xmlid={ftn40-5}]{\teiP[xmlid={p-355}]{Emmanuel Moureau, « Les collégiales du Quercy et la Curie pontificale du \teiHi[rend={small-caps}]{xiv}\teiHi[rend={sup}]{e} au \teiHi[rend={small-caps}]{xvi}\teiHi[rend={sup}]{e} siècle », \teiHi[rend={italic}]{Bulletin de la société archéologique et historique de Tarn-et-Garonne}, vol. 145, 2020, p. 19-48.}}.}
              \teiP[xmlid={p4-10}]{Pourquoi Jean II et sa femme portent-ils leur choix sur un type d’établissement ecclésiastique qui n’a jamais été prisé en Quercy ? Une hypothèse peut être suggérée en examinant la vie non pas du fondateur mais de son épouse. Marie de Culan est issue d’une ancienne famille du Berry, dont l’aïeul Louis a combattu aux côtés de Jeanne d’Arc. Elle est la fille de Philippe, seigneur de Jaloignes, maréchal de France en 1441 et d’Anne de Beaujeu. Un de ses oncles, Charles de Culan, a épousé en secondes noces Catherine de Castelnau, fille d’Antoine et sœur de son mari Jean II de Castelnau. Or, cette même Catherine de Castelnau avait épousé en premières noces Louis de Culan, cousin de Marie. Ce dernier, conseiller et chambellan du roi, a fondé en 1480 dans son château de Culan (au diocèse de Bourges) une collégiale castrale, réservée à son propre usage et desservie par six chanoines. François de Lignières, époux de Jeanne de Beaujeu, tante de Marie de Culan, a lui aussi fondé près de sa demeure du Berry un chapitre de chanoines similaire\teiNote[n={10},place={foot},xmlid={ftn41-5}]{\teiP[xmlid={p-356}]{Julien Noblet, \teiHi[rend={italic}]{En perpétuelle mémoire. Collégiales castrales et Saintes-Chapelles à vocation funéraire en France (1450-1560)}, Rennes, PUR, 2009.}}. Le choix de Jean II de Castelnau de privilégier un chapitre canonial séculier comme fondation familiale et mémorielle n’est donc pas incongru et lui a très certainement été inspiré par son épouse Marie, dont la famille avait dans un passé très proche fondé des collégiales.}
              \teiP[xmlid={p5-10}]{Cette nouvelle église ne relève toutefois que du seul seigneur : il s’agit là d’une véritable chapelle castrale à usage privé. Le droit de patronage du chapitre, qui inclut la nomination des chanoines, est réservé à la famille du fondateur. Les seigneurs de Castelnau, tout comme à Puylaroque, ne sont pas inhumés dans l’église Saint-Jean mais dans le cimetière avoisinant, au sein d’une modeste chapelle funéraire. La collégiale de Castelnau est là pour perpétuer le souvenir de ses deux fondateurs, mais les chanoines se retrouvent en quelque sorte les gardiens d’un tombeau vide.}
              \teiP[xmlid={p6-10}]{Jean II de Castelnau a obtenu de l’évêque de Cahors, Antoine de Luzech (1501-1510), la reconnaissance de sa collégiale, mais il souhaite toutefois que l’approbation soit également donnée par le pape lui-même, pour renforcer le prestige de la nouvelle fondation. Ce type de demande est assez courante dans les étapes de création de chapitres collégiaux séculiers, notamment dans le centre de la France, comme l’a exposé Julien Noblet dans sa thèse\teiNote[n={11},place={foot},xmlid={ftn42-5}]{\teiP[xmlid={p-357}]{\teiHi[rend={italic}]{Ibid}.}}. Pour la collégiale Saint-Jean-Baptiste, la confirmation pontificale est accordée en 1506, soit très rapidement après le vœu de Jean II, ce qui suggère une intervention directe auprès du souverain pontife. La bulle de Jules II précise que Guy de Castelnau, alors protonotaire apostolique, a accepté de se démettre de son bénéfice de curé de l’église prieurale Saint-Michel de Loubéjou en faveur de la nouvelle fondation\teiNote[n={12},place={foot},xmlid={ftn43-6}]{\teiP[xmlid={p-358}]{Ms. Fr. 23154, fol. 273r.}}. Il semble toutefois qu’un autre personnage plus influent que lui soit intervenu auprès de la Curie. Par comparaison, Louis II de la Trémouille a obtenu en 1516 des privilèges du pape Léon X pour sa collégiale de Thouars, suite à une lettre envoyée par Louis XII. Jean Stuart intervient directement auprès du même pontife en 1520 pour sa Sainte-Chapelle de Vic-le-Comte\teiNote[n={13},place={foot},xmlid={ftn44-6}]{\teiP[xmlid={p-359}]{Julien Noblet, \teiHi[rend={italic}]{op. cit}., p. 78.}}. Cette hypothèse est confirmée par la statuaire et les meubles héraldiques qui ornent les stalles de la collégiale de Prudhomat.}
            
\end{teiDiv}
            \begin{teiDiv}[type={section1},xmlid={div02-8}]

              \teiHead[xmlid={un-decor-heraldique-omnipresent-et-un-paradis-familial-001}]{Un décor héraldique omniprésent et un paradis familial}
              \teiP[xmlid={p7-10}]{Suite au testament de Jean II de Castelnau, décédé cette même année 1505, le chantier de la nouvelle collégiale Saint-Jean-Baptiste est lancé suivant son souhait. Jean Lartigaut a découvert dans les archives notariales deux prix-faits qui prouvent que les travaux ont débuté en 1507\teiNote[n={14},place={foot},xmlid={ftn45-6}]{\teiP[xmlid={p-360}]{Jean Lartigaut, « Sur la date de construction de la collégiale de Castelnau-Bretenoux », \teiHi[rend={italic}]{Bulletin de la société des études du Lot}, vol. XCVIII, 1977, p. 167-168.}}. Les historiens et historiens de l’art qui se sont penchés sur la collégiale – les chanoines Poulbrière\teiNote[n={15},place={foot},xmlid={ftn46-6}]{\teiP[xmlid={p-361}]{Jean-Baptiste Poulbrière, \teiHi[rend={italic}]{Notice historique et archéologique sur Castelnau-de-Bretenoux, Lot}, Tulle, Imprimerie E. Crauffon, 1873.}} et Gouzou\teiNote[n={16},place={foot},xmlid={ftn47-6}]{\teiP[xmlid={p-362}]{Joseph-Simon Gouzou, \teiHi[rend={italic}]{Bretenoux en Haut-Quercy}, Villefranche-de-Rouergue, Impr. Salingardes, 1955.}}, Jacques Juillet\teiNote[n={17},place={foot},xmlid={ftn48-7}]{\teiP[xmlid={p-363}]{Jacques Juillet, \teiHi[rend={italic}]{Les 38 barons de Castelnau : seconds barons chrétiens du royaume et leur baronnie, leur château, leurs églises}, s. l., Imprimerie Fabrègue, 1971.}} ou dernièrement Françoise Gatouillat\teiNote[n={18},place={foot},xmlid={ftn49-7}]{\teiP[xmlid={p-364}]{Françoise Gatouillat, « Prudhomat. Église Saint-Louis, ancienne collégiale Saint-Jean-Baptiste », dans Michel Hérold (dir.), \teiHi[rend={italic}]{Corpus vitrearum. Les vitraux du midi de la France}, Rennes, PUR, 2020, p. 238-240.}} – pensaient que c’était Jacques de Castelnau, fils de Jean II et époux de Françoise de La Tour, fille d’Agnet IV de La Tour, vicomte de Turenne et d’Agnès de Beaufort\teiNote[n={19},place={foot},xmlid={ftn50-7}]{\teiP[xmlid={p-365}]{Il l’a épousée le 1\teiHi[rend={sup}]{er} février 1500. Voir Alexandre Bruel, \teiHi[rend={italic}]{Inventaire d'une partie des titres de famille et documents historiques de la maison de La Tour d'Auvergne conservés dans les papiers Bouillon pour faire suite aux inventaires rédigés par Baluze (2e pt.)}, Paris, Archives nationales, 1905, p. 33. }} qui avait initié et suivi le chantier de l’église. Toutefois, l’examen du décor héraldique des clés de voûtes de l’édifice permet d’avancer une autre hypothèse. En effet, la clé de voûte du chœur de l’église porte les armes des Castelnau-Calmont – écartelé en 1 et 3 de gueules au château d’or et en 2 et 4 d’azur au lion d’argent – timbrées d’une crosse et d’une mitre. Deux griffons ailés supportent l’écu et s’appuient sur un phylactère qui porte l’inscription suivante en onciales gothiques : « LAN MIL CINC C XIII » (fig. 120). La date ainsi donnée est très certainement celle de l’achèvement des travaux de l’église, qui peut même être affinée avant le 26 mai 1514, date du mariage dans cette collégiale de Pierre Guillaume de Clermont-Lodève et de Marguerite de La Tour-Turenne. Le fait que cette date symbolique soit associée aux armes de Guy de Castelnau-Calmont, évêque de Périgueux depuis 1511, tend à prouver que c’est bien ce prélat qui a exécuté les dernières volontés de son frère Jean II et mené à bien le chantier de construction de la collégiale Saint-Jean-Baptiste.}
              \teiP[xmlid={p8-10}]{Les autres clés de voûtes de la nef de l’église sont toutes sculptées d’armoiries posées sur des réseaux flamboyants. Une première porte le blason familial des Castelnau-Calmont puis chacune des autres reprend un écu mi-parti qui rappelle les différents seigneurs du lieu et leurs alliances\teiNote[n={20},place={foot},xmlid={ftn51-7}]{\teiP[xmlid={p-366}]{La polychromie actuelle des clés de voûtes, qui paraît dater du \teiHi[rend={small-caps}]{xix}\teiHi[rend={sup}]{e} siècle, ne tient pas compte des émaux des armoiries représentées et ne doit, de ce fait, pas être prise en compte. }}. La généalogie familiale commence avec Antoine de Castelnau et Catherine de Chauvigny\teiNote[n={21},place={foot},xmlid={ftn52-7}]{\teiP[xmlid={p-367}]{Leur contrat de mariage a été dressé le 16 août 1436 à Issoudun. Voir Gaspard Thaumas de La Thaumassière, \teiHi[rend={italic}]{Histoire de Berry}, Bourges-Paris, F. Toubeau-Morel, 1689, p. 527.}} : aux armes des Castelnau sont associées celles de son épouse, qui porte d’argent à la fasce de fusées de gueules au lambel de six pendants d’azur\teiNote[n={22},place={foot},xmlid={ftn53-7}]{\teiP[xmlid={p-368}]{Voir notamment la base Sigilla pour les armes des Chauvigny de Brosse, http://www.sigilla.org/armoiries/chauvigny-25103.}}. À la génération suivante, ce sont les armes de Marie de Culan qui sont associées à celle de son mari Jean II : d’azur semé de molettes d’or, chargé d’un lion du même surmonté d’un lambel du même\teiNote[n={23},place={foot},xmlid={ftn54-7}]{\teiP[xmlid={p-369}]{Ces armoiries se retrouvent également sur la clé de voûte de la chapelle sud de la collégiale, portées par deux lions. }}. Enfin, les armoiries de Jacques, baron de Castelnau en titre lors de l’achèvement des travaux de la collégiale en 1514 et de son épouse Françoise de La Tour-Turenne terminent la lignée\teiNote[n={24},place={foot},xmlid={ftn55-7}]{\teiP[xmlid={p-370}]{Françoise de La Tour-Turenne portait un écartelé : en 1 et 3, d’azur semé de fleur de lys d’or à la tour d’argent maçonnée de sable ; en 2 et 4 bandé d’or et de gueules. }}. Comme pour Guy de Castelnau, l’écu est supporté par deux aigles qui tiennent entre leurs becs un phylactère, dont la partie basse s’enroule sur un bâton écôté. L’ensemble est posé sur un fond échiqueté. Ce lignage ainsi figuré par l’héraldique sur les clés de voûtes de l’église, éléments architectoniques qui assurent l’équilibre et la stabilité de l’édifice, illustre la vocation même de la collégiale, qui est d’être le panthéon dynastique des Castelnau-Calmont, famille solide et prestigieuse.}
              \teiP[xmlid={p9-9}]{Les vitraux qui ornent encore les baies du chœur de la collégiale complètent l’héraldique familiale. Jean II de Castelnau est figuré à genoux, revêtu d’une riche armure damasquinée et d’un tabard à ses armes. Ces dernières figurent également sur le prie-Dieu devant lequel il est agenouillé. Son saint patron, Jean-Baptiste, placé derrière lui, le désigne du doigt (fig. 121). Face à lui, son épouse Marie de Culan est présentée également par sa sainte patronne, Marie Jacobé, sœur légendaire de la Vierge, associée à ses deux enfants, saint Jean l’évangéliste et saint Jacques le Majeur. Les armoiries des Culan, initialement figurées sur le prie-Dieu, ont aujourd’hui disparu. Il manque également, déposées en 1949 lors d’une restauration des verrières, une représentation d’un chevalier en armure dorée, certainement Jacques de Castelnau, et son épouse Françoise de La Tour, dont les armes devaient également accompagner leur images\teiNote[n={25},place={foot},xmlid={ftn56-7}]{\teiP[xmlid={p-371}]{Voir Françoise Gatouillat, \teiHi[rend={italic}]{op. cit.}}}.}
              \teiP[xmlid={p10-9}]{Le décor héraldique se poursuit sur les stalles de la collégiale\teiNote[n={26},place={foot},xmlid={ftn57-7}]{\teiP[xmlid={p-372}]{Sur cet ensemble, voir Nelly Blaya et Nicolas Bru, \teiHi[rend={italic}]{Le mobilier des églises du Moyen Âge dans le Lot}, Toulouse, Conseil général Midi-Pyrénées, 2014, p. 42 et 43.}}. Cet ensemble remarquable en bois occupe encore la dernière travée de la nef avant le chevet (fig. 122). Deux grands corps se font face, chacun comportant treize sièges disposés sur deux rangées. Deux retours en équerre, à l’ouest, associés à une porte aujourd’hui disparue, isolaient à l’origine le chœur liturgique réservé aux chanoines de l’espace ouvert aux simples fidèles. Les appuie-mains et les miséricordes des stalles comportent un riche décor sculpté de personnages du quotidien ou d’animaux et les dorsaux des réseaux de remplages flamboyants aveugles. Les deux derniers panneaux des dorsaux situés près de l’autel sont toutefois ornés au nord par les armes de Guy de Castelnau-Calmont, timbrées de la crosse et de la mitre et soutenu par deux anges ; au sud des armoiries du baron de Castelnau, surmontées d’un casque à cimier avec la tête d’un lion, appuyées sur le même animal et supportées par un homme et une femme sauvages (fig. 123). Là encore, la présence appuyée des armes de Guy de Castelnau renforce l’hypothèse que ce prélat a suivi et accompagné les travaux de la collégiale, y compris la commande et la réalisation des sièges des chanoines.}
              \teiP[xmlid={p11-9}]{D’autres armoiries sont présentes sur les stalles. Les parties supérieures des jouées placées au plus près de l’autel comportent, inscrites dans deux motifs circulaires, au nord les armes des Castelnau-Calmont et les armes mi-parti Castelnau-Culan (pour Jean II et sa femme) (fig. 124) ; au sud, le blason des Castelnau et celui mi-parti Castelnau-Rochefort (pour Jean III et son épouse Charlotte de Rochefort\teiNote[n={27},place={foot},xmlid={ftn58-7}]{\teiP[xmlid={p-373}]{Charlotte de Rochefort, fille de Guy de Rochefort-Pluvant, chancelier de France, a épousé vers 1507 Jean III de Castelnau-Caumont. Elle portait comme armoiries d’azur semé de billettes d’or au chef d’argent chargé d’un lion léopardé de gueules. }}). L’association des armoiries du fondateur de la collégiale et de Jean III de Castelnau nous permet ainsi de proposer que les sièges des chanoines ont été réalisés après le décès de Jacques de Castelnau, survenu en 1514 ou avant fin avril 1515\teiNote[n={28},place={foot},xmlid={ftn59-7}]{\teiP[xmlid={p-374}]{Voir Henri de Lavaur de Sainte-Fortunade, « Documents sur la baronnie de Castelnau de Bretenoux », \teiHi[rend={italic}]{Bulletin de la Société scientifique, historique et archéologique de la Corrèze}, vol. 34, n\teiHi[rend={sup}]{o} 1, 1912, p. 353-403, ici p. 356.}} et avant celui de Guy de Castelnau en 1523. Un détail supplémentaire pourrait affiner cette fourchette de datation déjà étroite. Le revers de l’un des motifs circulaires avec les armes de Jean II est en effet sculpté d’un dauphin (fig. 125). Or, cet animal héraldique ne figure pas dans le bestiaire traditionnel des Castelnau-Calmont. Le 15 février 1518 est né François de France, Dauphin de France et duc de Bretagne, fils du roi François I\teiHi[rend={sup}]{er}. Le dauphin présent dans la collégiale de Prudhomat, vu les liens importants noués entre les Castelnau-Calmont et les familles alliées avec la Cour de France, pourrait être un hommage direct à la naissance de l’héritier du trône royal\teiNote[n={29},place={foot},xmlid={ftn60-7}]{\teiP[xmlid={p-375}]{De nombreuses fleurs de lys ponctuent également les colonnettes des dorsaux des stalles. }}. Ainsi, la réalisation des stalles serait plutôt à situer entre 1518 et 1523.}
              \teiP[xmlid={p12-9}]{Sur les jouées des sièges canoniaux prennent également place des statuettes qui figurent des saintes et des saints. Leur examen attentif montre que le programme iconographique retenu est directement lié à la famille du patron du chapitre collégial. En effet, l’ensemble des saints protecteurs de la famille des barons de Castelnau-Calmont et de leurs épouses est représenté : Jean le Baptiste et Jean l’évangéliste, Catherine, Marie, Jacques, Guy\teiNote[n={30},place={foot},xmlid={ftn61-7}]{\teiP[xmlid={p-376}]{Il s’agit très certainement de la statuette d’un évêque, dans laquelle nous proposons d’y voir celle de saint Guy, évêque d’Acqui. }}, Marguerite et Anne\teiNote[n={31},place={foot},xmlid={ftn62-7}]{\teiP[xmlid={p-377}]{Ces deux dernières saintes renvoient à Marie et Marguerite de Castelnau, filles de Jean II et de Marie de Culan.}}. Deux saints toutefois ne participent pas de ce paradis familial. Il s’agit d’un saint pape, reconnaissable à sa tiare (fig. 126) et saint Jérôme en habit de cardinal, avec un lion à ses pieds (fig. 127). Ces deux images sont placées respectivement sur les panneaux des jouées nord et sud des stalles situées au plus près de l’autel. Elles occupent ainsi un emplacement privilégié et paraissent protéger les chanoines et la collégiale. Le pape n’est autre que Jules II, figuré sous les traits de son saint patron Jules I\teiHi[rend={sup}]{er}, qui a, nous l’avons vu, confirmé l’érection de la collégiale Saint-Jean-Baptiste de Prudhomat et ainsi le vœu émis par Jean II de Castelnau. Quant au saint cardinal, il faut y voir une allégorie de François-Guillaume de Clermont-Lodève, cardinal depuis 1503, archevêque d’Auch (1507-1511), ambassadeur du roi Louis XII auprès de Jules II. Ce cousin direct des barons de Castelnau-Calmont\teiNote[n={32},place={foot},xmlid={ftn63-7}]{\teiP[xmlid={p-378}]{Son frère Pierre a épousé Marguerite de La Tour à Prudhomat en 1514 et est devenu après le décès de Jean III en 1530 l’héritier des biens et titres de la maison de Castelnau-Caumont. }} a très certainement dû appuyer auprès du Souverain Pontife la demande de reconnaissance en collégiale des chapellenies érigées par Jean II à Prudhomat, d’où sa représentation symbolique sur les stalles. Outre ces deux statues, les colonnettes des dorsaux des stalles de Prudhomat portent également des feuilles de chêne – allusion aux armes parlantes de la famille de Jules II, les della Rovere – et des hermines, meuble héraldique des Clermont-Lodève (fig. 128).}
              \teiP[xmlid={p13-9}]{La collégiale Saint-Jean-Baptiste-Saint-Louis de Prudhomat est un rare exemple dans le Midi de la France d’un modeste chapitre de chanoines séculiers fondé au début du \teiHi[rend={small-caps}]{xvi}\teiHi[rend={sup}]{e} siècle par une riche famille aristocratique qui conserve encore à la fois son décor héraldique mais également ses stalles richement ornées. C’est à travers la sculpture de ces dernières que le commanditaire de cet ensemble, Guy de Castelnau, évêque de Périgueux, a rendu hommage aux deux bienfaiteurs du chapitre canonial voulu par son frère, le pape Jules II et le cardinal de Clermont-Lodève, mis à l’honneur dans le chœur liturgique. Il s’agit très certainement d’un des seuls exemplaires conservés \teiHi[rend={italic}]{in situ} à ce jour d’un mobilier religieux qui relie directement, via l’héraldique et la sculpture, un pape et une communauté canoniale reconnue par ce dernier.}
              \begin{teiFigure}[xmlid={figure01-8}]

                \teiGraphic[url={../icono/br/Ch09\_Moureau/Fig123.JPG}]
                \teiHead[xmlid={figure-120-cle-de-voute-aux-armes-de-guy-de-castelnau-calmont-eveque-de-perigueux-001}]{Figure 120. Clé de voûte aux armes de Guy de Castelnau-Calmont, évêque de Périgueux.}
              
\end{teiFigure}
              \begin{teiFigure}[xmlid={figure02-8}]

                \teiGraphic[url={../icono/br/Ch09\_Moureau/Fig124.jpg}]
                \teiHead[xmlid={figure-121-jean-ii-de-castelnau-calmont-et-son-saint-patron-verriere-du-choeur-vers-1510-1514-001}]{Figure 121. Jean II de Castelnau-Calmont et son saint patron, verrière du chœur, vers 1510-1514.}
              
\end{teiFigure}
              \begin{teiFigure}[xmlid={figure03-8}]

                \teiGraphic[url={../icono/br/Ch09\_Moureau/Fig125.jpg}]
                \teiHead[xmlid={figure-122-vue-generale-des-stalles-de-la-collegiale-de-prudhomat-001}]{Figure 122. Vue générale des stalles de la collégiale de Prudhomat.}
              
\end{teiFigure}
              \begin{teiFigure}[xmlid={figure04-8}]

                \teiGraphic[url={../icono/br/Ch09\_Moureau/Fig126.jpg}]
                \teiHead[xmlid={figure-123-armoiries-du-baron-de-castelnau-calmont-stalles-du-choeur-001}]{Figure 123. Armoiries du baron de Castelnau-Calmont, stalles du chœur.}
              
\end{teiFigure}
              \begin{teiFigure}[xmlid={figure05-8}]

                \teiGraphic[url={../icono/br/Ch09\_Moureau/Fig127.jpg}]
                \teiHead[xmlid={figure-124-armoiries-de-jean-ii-de-castelnau-calmont-et-de-marie-de-culan-stalles-du-choeur-001}]{Figure 124. Armoiries de Jean II de Castelnau-Calmont et de Marie de Culan, stalles du chœur.}
              
\end{teiFigure}
              \begin{teiFigure}[xmlid={figure06-8}]

                \teiGraphic[url={../icono/br/Ch09\_Moureau/Fig128.jpg}]
                \teiHead[xmlid={figure-125-detail-du-dauphin-stalles-du-choeur-001}]{Figure 125. Détail du dauphin, stalles du chœur.}
              
\end{teiFigure}
              \begin{teiFigure}[xmlid={figure07-7}]

                \teiGraphic[url={../icono/br/Ch09\_Moureau/Fig129.jpg}]
                \teiHead[xmlid={figure-126-saint-jules-i-er-allegorie-du-pape-jules-ii-stalles-du-choeur-001}]{Figure 126. Saint Jules I\teiHi[rend={sup}]{er}, allégorie du pape Jules II, stalles du chœur.}
              
\end{teiFigure}
              \begin{teiFigure}[xmlid={figure08-6}]

                \teiGraphic[url={../icono/br/Ch09\_Moureau/Fig130.JPG}]
                \teiHead[xmlid={figure-127-saint-jerome-allegorie-du-cardinal-francois-guillaume-de-clermont-lodeve-stalles-du-choeur-001}]{Figure 127. Saint Jérôme, allégorie du cardinal François Guillaume de Clermont-Lodève, stalles du chœur.}
              
\end{teiFigure}
              \begin{teiFigure}[xmlid={figure09-6}]

                \teiGraphic[url={../icono/br/Ch09\_Moureau/Fig131.jpg}]
                \teiHead[xmlid={figure-128-detail-des-hermines-et-des-feuilles-de-chene-sculptees-sur-une-colonnette-des-stalles-du-choeur-001}]{Figure 128. Détail des hermines et des feuilles de chêne sculptées sur une colonnette des stalles du chœur.}
              
\end{teiFigure}
            
\end{teiDiv}
          
\end{teiElement}
        
\end{teiElement}
        \begin{teiElement}[name={group},type={chapter},data-page-title={Légitimer la légation avignonnaise},data-page-subtitle={Héraldique et politique de l’Altera Roma (xvie-xviie siècles)},data-page-authors={Camille Larraz},data-include-href={Ch10\_Legitimer\_legation\_avignonnaise.xml},xmlid={chapter-007}]

          \begin{teiElement}[name={front}]

            \begin{teiDiv}[type={titlePage},xmlid={titlepage-012}]

              \teiP[rend={title-main},xmlid={p-379}]{Légitimer la légation avignonnaise}
              \teiP[rend={title-sub},xmlid={p-380}]{Héraldique et politique de l’Altera Roma (\teiHi[rend={small-caps}]{xvi}\teiHi[rend={sup}]{e}-\teiHi[rend={small-caps}]{xvii}\teiHi[rend={sup}]{e} siècles)}
              \teiP[rend={author-aut},xmlid={p-381}]{Camille Larraz}
            
\end{teiDiv}
          
\end{teiElement}
          \begin{teiElement}[name={body},xmlid={body-012}]

            \teiP[xmlid={p1-12}]{À partir du \teiHi[rend={small-caps}]{xiv}\teiHi[rend={sup}]{e} siècle, le séjour des souverains pontifes donne une nouvelle dimension à Avignon, qui commence alors à définir son identité culturelle, politique et religieuse en tant qu’\teiHi[rend={italic}]{Altera Roma}, réaffirmant ainsi son statut de cité sainte. Même après leur départ en 1418, la ville demeure la propriété du Saint-Siège et est gouvernée par un légat exerçant l’autorité du pape. Si le concept de l’\teiHi[rend={italic}]{Altera Roma} est déjà bien étudié\teiNote[n={1},place={foot},xmlid={ftn1-10}]{\teiP[xmlid={p-382}]{Notamment par Olivier Rouchon, « Rituels publics, souveraineté et identité citadine. Les cérémonies d’entrée en Avignon (\teiHi[rend={small-caps}]{xvi}\teiHi[rend={sup}]{e}-\teiHi[rend={small-caps}]{xvii}\teiHi[rend={sup}]{e} siècles) », \teiHi[rend={italic}]{Cahiers de la Méditerranée}, n\teiHi[rend={sup}]{o} 77, 2008, p. 39-59, https://journals.openedition.org/cdlm/4362.}}, cet article se propose de l’examiner à travers le prisme de l’histoire de l’art, premièrement par le rôle des peintres, lors de somptueuses entrées au cours desquelles est mise en jeu une abondance héraldique, et secondement par le rôle de la gravure, qui favorise une large diffusion de ces idées. Ces quelques exemples nous permettront ainsi d’interroger l’importance du visuel, en particulier des armoiries et autres dispositifs héraldiques, pour exprimer et légitimer le concept d’\teiHi[rend={italic}]{Altera Roma} entre la seconde moitié du \teiHi[rend={small-caps}]{xvi}\teiHi[rend={sup}]{e} siècle et la seconde moitié du \teiHi[rend={small-caps}]{xvii}\teiHi[rend={sup}]{e} siècle.}
            \begin{teiDiv}[type={section1},xmlid={div01-9}]

              \teiHead[xmlid={lentree-farnese-de-1553-001}]{L’entrée Farnèse de 1553}
              \teiP[xmlid={p2-11}]{Les somptueuses entrées à Avignon sont des événements particulièrement propices pour exprimer son identité urbaine et sa loyauté envers le Saint-Siège, en particulier à travers le faste visuel qui est déployé à cette occasion. L’entrée Farnèse de 1553 étant bien documentée par les sources municipales, les archives permettent de comprendre comment une profusion héraldique et décorative est déployée à cette occasion, afin de légitimer la prise de pouvoir du légat et son appropriation de l’espace urbain. Les comptes de la ville soulignent notamment l’importance du métier de peintre, ainsi que des matériaux et pigments qui leur sont alloués pour créer les décors et éléments symboliques visant à recréer une pompe à la romaine. Le 16 mars 1553, le cardinal légat Alexandre Farnèse fait ainsi son entrée solennelle dans la ville d’Avignon, alors qu’il revient de Paris, où l’usage veut qu’il rencontre le roi à la cour. Vice-chancelier de l’Église et archevêque d’Avignon depuis 1535, après la mort du cardinal Hippolyte de Médicis, il est également resté dans les mémoires pour avoir exercé la fonction de légat en Avignon à partir de 1541. Alexandre Farnèse n’ayant lui-même jamais séjourné sur le long terme à Avignon, un vice-légat qui y demeurait de façon permanente faisait office de représentant de la diplomatie pontificale et de sa juridiction spirituelle sur les provinces ecclésiastiques du comtat Venaissin (Vienne, Embrun, Arles, Aix et parfois Narbonne\teiNote[n={2},place={foot},xmlid={ftn37-6}]{\teiP[xmlid={p-383}]{Bernard Barbiche et Ségolène de Dainville-Barbiche, « Les légats \teiHi[rend={italic}]{a latere} en France et leurs facultés aux \teiHi[rend={small-caps}]{xvi}\teiHi[rend={sup}]{e} et \teiHi[rend={small-caps}]{xvii}\teiHi[rend={sup}]{e} siècles », \teiHi[rend={italic}]{Archivum Historiae Pontificae}, n\teiHi[rend={sup}]{o} 23, 1985, p. 93-165, ici p. 94.}}). Bien que les cérémonies d’entrée soient assez rares au \teiHi[rend={small-caps}]{xvi}\teiHi[rend={sup}]{e} siècle, elles sont supervisées par le conseil de ville et ses trois consuls. En outre, les magistrats municipaux veillent à ce qu’un livret imprimé relatant l’événement soit rédigé afin de le mémoriser et de le célébrer, mais aussi d’expliquer les subtilités des allégories élaborées et de rendre immuable un moment éphémère par sa nature même\teiNote[n={3},place={foot},xmlid={ftn38-6}]{\teiP[xmlid={p-384}]{Yves Pauwels, « Le thème de l’arc de triomphe dans l’architecture urbaine à la Renaissance, entre pouvoir politique et pouvoir religieux », dans Patrick Boucheron et Jean-Philippe Genet (dir.), \teiHi[rend={italic}]{Marquer la ville. Signes, traces, empreintes du pouvoir (}\teiHi[rend={small-caps italic}]{xiii}\teiHi[rend={italic sup}]{e}\teiHi[rend={italic}]{-}\teiHi[rend={small-caps italic}]{xvi}\teiHi[rend={italic sup}]{e}\teiHi[rend={italic}]{ siècle)}, Paris-Rome, Éditions de la Sorbonne-École française de Rome, 2014, p. 181-188.}}. Dans le cas d’Avignon, l’entrée Farnèse est la première à être retranscrite sous la forme imprimée, par Honoré Henry, source essentielle pour cette étude\teiNote[n={4},place={foot},xmlid={ftn39-6}]{\teiP[xmlid={p-385}]{La page de couverture présente les armoiries de la ville d’Avignon.}}. Les entrées du cardinal Farnèse dans les villes d’Avignon et de Carpentras, respectivement en mars et en avril 1553, ont notamment été étudiées par Olivier Rouchon et Richard Cooper, qui en proposent de fines analyses historiques et iconographiques, permettant de mieux saisir les enjeux identitaires de ces rituels\teiNote[n={5},place={foot},xmlid={ftn40-6}]{\teiP[xmlid={p-386}]{Olivier Rouchon, art. cité et Richard Cooper, « Legate’s Luxury: the Entries of Cardinal Alessandro Farnese to Avignon and Carpentras, 1553 », dans Nicolas Russell et Hélène Visentin (dir.), \teiHi[rend={italic}]{French Ceremonial Entries in the Sixteenth century: Event, Image, Text}, Toronto, Center for Renaissance and Reformation Studies, 2007, p. 133-161.}}. Ces travaux ne s’attardent toutefois pas sur le rôle des peintres et de l’héraldique, qui nous semble essentiel pour comprendre les mécanismes de pouvoir en jeu.}
              \teiP[xmlid={p3-11}]{Formellement, l’entrée consiste en un accueil formalisé de celui qu’on veut honorer, d’abord à l’église de Saint-Véran, juste en-dehors des remparts, puis à la porte Saint-Lazare, avant que l’hôte ne soit mené par un cortège à travers un itinéraire précis dans la ville. Près de la porte est installée une tribune aux harangues depuis laquelle étaient déclamées des paroles rituelles. À la différence d’une entrée royale, l’entrée d’un légat prévoyait l’aménagement d’une chapelle temporaire sous forme de dais, appelé « poisle ». Ce dais fait d’ailleurs partie des éléments codifiés les plus importants de la cérémonie. Honoré Henry, dans son livret de 1553 relatant la cérémonie, nous le décrit précisément : « aiant au dessus de luy un poisle de damas cramoisi rouge à franges d’or, \& à sex bastons semés de fleurs de lis que les trois Consulz, assesseurs \& deux gentilhommes de la Ville portoient à pied\teiNote[n={6},place={foot},xmlid={ftn41-6}]{\teiP[xmlid={p-387}]{Honoré Henry, \teiHi[rend={italic}]{La magnifique entrée du reverendissime et très illustre seigneur, Monseigneur Alexandre, Cardinal de Farnez, Vichancelier du Sainct siege apostolique, et Legat de la ville et Cité d’Avignon, faicte en icelle le }\teiHi[rend={small-caps italic}]{xvi}\teiHi[rend={italic}]{ mars 1553}, Avignon, s. n., 1553, p. C iii.}} ». Le rouge cramoisi et l’or semblent ainsi faire allusion à la Rome antique, alors que l’usage du poisle est lui aussi fortement symbolique. En effet, ce sont les trois consuls de la ville, le viguier (un magistrat aux fonctions analogues à celles d’un prévôt), et d’autres membres de l’élite locale qui doivent saisir les six montants lorsqu’il est utilisé comme instrument d’honneur. Il est intéressant de souligner les fleurs de lys ornant les montants de bois : la présence des armoiries Farnèse et d’Avignon est un véritable \teiHi[rend={italic}]{leitmotiv} qui traverse l’ensemble de la ville, qui demande la protection du légat, auquel elle délègue son pouvoir. Nous retrouvons l’usage du poisle armorié dans d’autres villes du royaume de France lors d’entrées de légats. Par exemple, à Rouen en 1596, « \lbrack{}s\rbrack{}uivant les lettres du Roi, l’on fera Entrée au sieur légat \lbrack{}Alexandre de Médicis, futur Léon XI\rbrack{}, auquel sera presenté poisle de damas, aux armaryes dudit 3\teiHi[rend={sup}]{r}, porté par les quatre quarteniers\teiNote[n={7},place={foot},xmlid={ftn42-6}]{\teiP[xmlid={p-388}]{\teiHi[rend={italic}]{Entrée à Rouen du roi Henri IV en 1596, précédée d’une introduction par J. Félix et de notes par Ch. de Robillard de Beaurepaire}, Rouen, Imprimerie de E. Cagniard, 1887, p. XLIII.}} ». De même, les entrées de légats à Paris suivent un cérémonial précis qui comprend également l’emploi d’un damas blanc peint aux armes du légat et de la ville, « porté à tour de rôle par les échevins, les bourgeois et les marchands des six corps de la ville\teiNote[n={8},place={foot},xmlid={ftn43-7}]{\teiP[xmlid={p-389}]{Bernard Barbiche et Ségolène de Dainville-Barbiche, art. cité, p. 106.}} ».}
              \teiP[xmlid={p4-11}]{Les rues que le cortège traverse sont aussi ornées de tentures aux couleurs et décorations significatives. À nouveau, Honoré Henry nous dit à ce propos : « ledict Seigneur, avec toutes les compaignies disposées en bon \& convenant ordre, vindrent en une Eglise nommée Sainct Veran \lbrack{}...\rbrack{}. Puis apres vint au devant de ladicte eglise en une Loge tout au tout \& au dedans tendue de riche tappisserie, avec arcade \& appuys sur le devant, le tout enrichi à force festons, chapeaux de triumphe, \& armoiries dudict Seigneur environnez d’or cliquant \lbrack{}...\rbrack{}\teiNote[n={9},place={foot},xmlid={ftn44-7}]{\teiP[xmlid={p-390}]{Honoré Henry, \teiHi[rend={italic}]{op. cit.}, p. B ii. Saint Véran est la première étape du parcours, juste en-dehors de la ville. La représentation du comtat Venaissin dans la galerie des cartes du Vatican par Antonio Danti (1581) montre bien le lieu hors des remparts.}}. » Comme nous le verrons plus loin, l’or fait partie des matériaux les plus importants alloués aux peintres, un statut qui se reflète visuellement lors du faste de la cérémonie. Ces décors éphémères sont réalisés par des artisans locaux sous la direction d’un architecte ayant conçu le programme, en général d’après les nombreux traités représentant les antiques, tels que celui de Sebastiano Serlio sur les Joyeuses entrées en 1549. Nous ignorons qui a conçu le programme avignonnais, mais il a été supposé qu’il s’agit des mêmes personnes chargées de l’entrée du cardinal à Carpentras en avril 1553, dont le livret écrit par Antoine Blégier, imprimé en italien et en français, présente des gravures de la main de Bernard Salomon, à la différence de la version avignonnaise. Cet imprimé peut ainsi nous servir de point de comparaison pour tenter de reconstituer les décors réalisés à l’occasion de l’entrée du légat à Avignon. On y trouve notamment des architectures éphémères comme des arcs de triomphe ou des tableaux vivants permettant de façonner l’espace urbain et de recréer par là une pompe antique digne des empereurs romains (fig. 129 et 130).}
              \teiP[xmlid={p5-11}]{Les décors signalent visuellement la souveraineté pontificale à travers la figure du légat. Symboliquement, la procession sert donc à concrétiser par l’image la prise de possession consentie du légat, qui endosse son rôle de dirigeant et de protecteur, tout en réaffirmant la souveraineté du pape dont il est le représentant. Le faste visuel de la cérémonie devient acte politique : il lui permet de se faire connaître physiquement et publiquement aux citoyens et de manifester son pouvoir par des signes visibles. Comme le note Richard Cooper, un élément créé pour l’occasion est intéressant à plus d’un égard : l’emblème réalisé chaque année à Noël et utilisé pour l’année à venir. Pour 1553, probablement en prévision de l’entrée de mars, il présente la tiare papale, les clés de Saint-Pierre, un chapeau de cardinal dont deux colombes avec une branche d’olivier descendent, deux cornes d’abondance, deux palmes et deux mains jointes avec la devise « \teiHi[rend={italic}]{Sic omnia prospera}\teiNote[n={10},place={foot},xmlid={ftn45-7}]{\teiP[xmlid={p-391}]{Honoré Henry, \teiHi[rend={italic}]{op. cit.,} p. B iiiiv à B v ; et Richard Cooper, art. cité, p. 138 et 139.}} ». Cet emblème, représentant l’unité et la paix apportées par la papauté et le cardinal légat, est combiné avec les armes de Farnèse et des six fleurs de lys, pratique qui permet d’intégrer facilement l’alliance de la France et des Farnèse puisque les fleurs font naturellement allusion au royaume. Nous pouvons donc y voir une représentation métaphorique du statut d’Avignon, celle-ci étant autant loyale à Rome et au pape qu’au royaume de France, dont elle est régnicole depuis 1536. La profusion héraldique lors de la cérémonie d’entrée joue donc un rôle fondamental pour « une appropriation symbolique temporaire de l’espace urbain\teiNote[n={11},place={foot},xmlid={ftn46-7}]{\teiP[xmlid={p-392}]{Olivier Rouchon, art. cité. }} ».}
            
\end{teiDiv}
            \begin{teiDiv}[type={section1},xmlid={div02-9}]

              \teiHead[xmlid={peintres-et-heraldique-001}]{Peintres et héraldique}
              \teiP[xmlid={p6-11}]{L’abondance héraldique présente lors de la cérémonie fait partie des nombreux éléments de chantier confiés aux peintres de la ville, engagés sur plusieurs mois pour réaliser ces travaux considérables, entre les mois de novembre 1552 et mars 1553. À la tête de l’équipe des artisans, composée d’une douzaine de peintres locaux et peut-être étrangers, se trouve Simon de Mailly, dit de Châlons, premier peintre de la ville\teiNote[n={12},place={foot},xmlid={ftn47-7}]{\teiP[xmlid={p-393}]{Voir le manuscrit de Pierre Pansier, \teiHi[rend={italic}]{Les peintres d’Avignon au }\teiHi[rend={small-caps italic}]{xvi}\teiHi[rend={italic sup}]{e}\teiHi[rend={italic}]{ siècle. Biographies et documents}, Avignon, BM, ms. 5712 ; et Camille Larraz, \teiHi[rend={italic}]{Simon de Châlons}, Cinisello Balsamo, Silvana Editoriale, 2022.}}. Les documents d’archives révèlent quantité d’informations sur les matériaux qui sont alloués aux peintres (fig. 131), notamment des pigments, presque exclusivement confiés à Simon de Châlons, probablement en raison de leur valeur. Nous trouvons une grande variété de pigments fournis à Simon de Mailly : des bleus (florée et azur), des verts (verdet, vert de vessie), des rouges (ocre rouge, cinabre, rosette, rouge de minium, dit « mini »), des jaunes (ocre jaune, orpiment, massicot, le stil de grain, dit « estut de grum »), des blancs (céruse, céruse menue), et du noir broyé\teiNote[n={13},place={foot},xmlid={ftn48-8}]{\teiP[xmlid={p-394}]{Pierre Pansier, \teiHi[rend={italic}]{Les peintres d’Avignon…},\teiHi[rend={italic}]{ op. cit.}, fol. 365, 391, 516 et 517.}}. On lui octroie par la suite des ingrédients essentiels à la préparation desdits pigments, comme de l’huile de noix, du vernis liquide, de la laque fine, de la colle claire, de l’aigue ardent, du papier « d’estrasse\teiHi[rend={italic}]{ »} et du coton battu (pour appliquer les feuilles d’or). Quantitativement, le blanc et le noir sont les pigments les plus employés, suivis de l’ocre jaune, de l’ocre rouge et de l’azur, de l’orpiment, du cinabre et des verts, de la rouzette et du massicot. Il n’est pas étonnant de voir le blanc et le noir utilisés majoritairement (un très grand nombre de textes poétiques et allégoriques ornaient les décors), tout comme le jaune, rouge et bleu pour les armoiries d’Alexandre Farnèse (d’or à six fleurs de lys d’azur) et de la Ville d’Avignon (de gueules à trois clefs d’or posées en fasce), elles aussi reproduites de très nombreuses fois.}
              \teiP[xmlid={p7-11}]{Les quantités d’or et d’argent nécessaires étant vraisemblablement trop importantes pour l’ensemble des décorations, de l’oripeau (« auripel ») est employé comme substitut. Ces feuilles battues de laiton ou de cuivre avaient l’apparence de l’or ou du bronze antique et permettaient de rendre un effet doré à moindre coût. Ces substances ont peut-être également été utilisées pour des armoiries, réalisées en de très nombreux exemplaires ainsi que d’autres éléments de décors intégrés à l’ensemble. Le fait qu’« une grosse auripel » et « deux douzènes auripel » soient données respectivement aux peintres Antoine Duplan et Claude Roc, et que la ligne suivante mentionne l’impressionnant nombre de 1 125 pignons à « dourer\teiNote[n={14},place={foot},xmlid={ftn49-8}]{\teiP[xmlid={p-395}]{\teiHi[rend={italic}]{Ibid.}}} », laisse imaginer que ce matériau est employé à cet effet. Un mandat individuel engage précisément Antoine Duplan pour dorer 1 125 « pignes » « pour l’entrée de Mgr le légat », pour le prix de 4 livres et 10 sous\teiNote[n={15},place={foot},xmlid={ftn50-8}]{\teiP[xmlid={p-396}]{Avignon, AM, CC468 ; Avignon, BM, ms. 5712, fol. 515. }}. Certains autres peintres reçoivent en effet des mandats précis, par exemple pour la peinture de « trois douzaine de petites armes », « unes grandes armes de toile » et 50 fleurs de lis sur « gros papier » pour Louis de Champclaus\teiNote[n={16},place={foot},xmlid={ftn51-8}]{\teiP[xmlid={p-397}]{Avignon, AM, CC468.}}, alors que Jean de Linga, ou Langue, réalise 1 250 petites armes, ainsi que deux douzaines de grandes armes pour la porte Saint-Lazare\teiNote[n={17},place={foot},xmlid={ftn52-8}]{\teiP[xmlid={p-398}]{\teiHi[rend={italic}]{Ibid}.}}, là où commence l’entrée. Il est difficile de déterminer si ces 1 250 armes sont destinées à l’ensemble du parcours ou simplement à quelques étapes, au nombre de huit. Gonet Chanis reçoit également un mandat pour 4 grands écussons armoriés « de notre saint père le pape de l’eglise, de monseigneur le legat et de la ville\teiNote[n={18},place={foot},xmlid={ftn53-8}]{\teiP[xmlid={p-399}]{\teiHi[rend={italic}]{Ibid}.}} ». Malheureusement, nous ignorons à nouveau où sont destinées ces armes, qui devaient être autant accrochées sur l’ensemble du parcours qu’à des endroits spécifiques où étaient habituellement disposés les emblèmes du pape et de la ville. Les gravures de la main de Bernard Salomon dans le livret de Carpentras nous montrent, par exemple, un bâtiment recouvert sur presque chaque brique, ou pierre, de fleurs de lys, surmonté de deux étendards montrant à nouveau l’emblème floral de la France ainsi que les armoiries de la ville (fig. 132). De même, un arc de triomphe laisse apparaître les armes du pape, du cardinal Farnèse et de la couronne de France, dans une mise en scène qui a dû être reproduite à Avignon si l’on en croit les descriptions d’Honoré Henry. Si les concepteurs des deux entrées se révèlent être les mêmes, il y a de fortes chances que les décors soient relativement similaires dans leur conception, avec la présence de figures mythologiques ou allégoriques parsemant les étapes, également ornées de blasons et de devises.}
              \teiP[xmlid={p8-11}]{Comme nous avons déjà pu le comprendre avec les extraits cités de Honoré Henry, la confrontation des archives avec son texte et les gravures de Carpentras confirme cette abondance décorative et héraldique : en-dehors des nombreuses mentions d’armoiries présentes sur les décors, l’auteur insiste presque à chaque page sur les tissus portés par les participants de la cérémonie, les différentes compagnies, en particulier sur des tissus chers comme le velours et le satin, toujours de couleurs significatives, surtout de rouge ou d’or. L’importance des costumes se retrouve à nouveau dans les archives, où les artisans sont payés pour réaliser notamment les manches doublées des vêtements : « 2 pans satin rouge fin balhat à mestre Anthoine de la Bancaso pour fere \lbrack{}...\rbrack{} manches »\teiNote[n={19},place={foot},xmlid={ftn54-8}]{\teiP[xmlid={p-400}]{Avignon, AM, CC468.}}. Nous pouvons encore citer Honoré Henry, qui décrit entre autres les costumes de la compagnie des veloutiers : « Lequel \lbrack{}gouiat\rbrack{} estoit suyvi de sa bande, quatre à quatre en chasque rang, tous en habits de velours jaulne : des quelz la meilleure part estoit reluysante en mainst beaux boutons, \& fers d’or. L’Enseigne trop plus brave en conference des autres, aiant son mourrion en teste, \& maniant son enseigne qui estoit de taffetas jaulne, semée de fleurs de list d’asur, armoiries dudict Seigneur\teiNote[n={20},place={foot},xmlid={ftn55-8}]{\teiP[xmlid={p-401}]{Honoré Henry, \teiHi[rend={italic}]{op. cit.}, p. C ii-C iiv.}}. » Ou encore les trompettes de la ville : « Apres eulx venoient les trompettes de la ville, vestus de satin rouge, avec la baniere de leur trompe de taffetas de pareille couleur, semés de trois clefs d’or, les armoiries de la Ville\teiNote[n={21},place={foot},xmlid={ftn56-8}]{\teiP[xmlid={p-402}]{\teiHi[rend={italic}]{Ibid.}, p. B iiii.}}. » Les vêtements des participants font donc écho aux décors de la ville ; ce sont des éléments tout aussi codifiés et importants pour visualiser leur consentement et loyauté envers le pape.}
              \teiP[xmlid={p9-10}]{De façon générale, chaque étape du parcours, ornementée de grands décors éphémères à l’antique, est accompagnée par l’omniprésence d’armoiries, dont le très grand nombre s’accorde avec les documents d’archives qui voient plusieurs peintres réaliser des blasons, parfois en quantité impressionnante. Par ailleurs, le texte d’Honoré Henry fait fréquemment référence à cette abondance héraldique. Par exemple, il insiste, à propos du portail de la ville, sur l’impossibilité de rendre précisément compte de l’ensemble des décors et des armes : « Je laisse à dire (à cause de briefvesté) les armoiries de nostre Sainct Pere, dudict seigneur Legat, Vicelegat, Gouverneur, \& de la Ville, ensemble l’Embleme des Consulz, \& devise, le tous assis sus un front d’espic au devant du portal, avec tymbre \& festons enrichis : chose certes belle à veoir pour la variété des interpositions des metaulx \& couleurs\teiNote[n={22},place={foot},xmlid={ftn57-8}]{\teiP[xmlid={p-403}]{\teiHi[rend={italic}]{Ibid.}, p. C iii.}}. » Le mot « metaulx » est particulièrement intéressant dans le contexte qui nous occupe, car nous l’avons vu dans les documents d’archives, l’or et l’argent apparaissent régulièrement, ainsi que l’oripeau. Dans le même ordre d’idée, d’autres extraits soulignent le grand nombre d’armoiries exposées parmi les décors éphémères : « Au desus duquel estoit autre couroir à clair voie à la façon Ionique, composé de bien vint \& cinq, ou trente arcs enrichis de festons \& armoiries dudict Seigneur, \& de la Ville\teiNote[n={23},place={foot},xmlid={ftn58-8}]{\teiP[xmlid={p-404}]{\teiHi[rend={italic}]{Ibid.}, p. C iiii.}} », ou encore, l’auteur remarque aux Augustins « une infinité de factures, enrichies des couleurs dudict Seigneur \& de la Ville. \lbrack{}...\rbrack{} Et en l’arc du milieu estoit un grand nombre de fleurs de lis d’asur en champ d’or. \lbrack{}...\rbrack{} Semblablement les deux arcs, costoians celuy du milieu, estoient renouvellez par dedans de clefz d’or en champs de gueules\teiNote[n={24},place={foot},xmlid={ftn59-8}]{\teiP[xmlid={p-405}]{\teiHi[rend={italic}]{Ibid.}, p. D iiv.}}. » La rencontre dans les comptes de la ville de paiements pour des milliers de blasons semble ainsi confirmée par le rapport visuel d’Honoré Henry. Vis-à-vis des documents, il convient également de remarquer la concordance iconographique entre son texte et les archives conservées. En effet, nous retrouvons dans les deux cas une distinction importante entre les petites et les grandes armes, notamment au palais des Papes, dont l’entrée « estoit ornée de festons \& entrelassements magnifiques, bayes, \& tymbres remplis par le dedans de petites armoiries \& par dehors de quatre grands escussons de luy, \& messeigneurs ses trois freres\teiNote[n={25},place={foot},xmlid={ftn60-8}]{\teiP[xmlid={p-406}]{\teiHi[rend={italic}]{Ibid.}, p. E ii.}} ». Bien qu’il soit difficile de savoir à quelle taille ces blasons correspondaient, il semble notable que l’on privilégie des armoiries de grande taille pour les lieux les plus significatifs, en l’occurrence le palais.}
            
\end{teiDiv}
            \begin{teiDiv}[type={section1},xmlid={div03-7}]

              \teiHead[xmlid={avignon-une-altera-roma-001}]{Avignon, une « Altera Roma »}
              \teiP[xmlid={p10-10}]{Comme le texte d’Honoré Henry nous le confirme, les armes du légat Alexandre Farnèse apparaissent presque systématiquement avec les armes de la ville d’Avignon, permettant à la cité pontificale de réaffirmer son obéissance à Rome. En se plaçant sous la protection de son nouveau dirigeant, la ville espère en retirer non seulement sa bienveillance, mais surtout des bénéfices concrets. Concept éminemment politique et religieux, le principe d’\teiHi[rend={italic}]{Altera Roma} est adopté très tôt par les papes avignonnais. Par exemple, une chaire en ogive à Notre-Dame-des-Doms qui servait de trône pontifical depuis 1307 était soutenue par trois anges et complétée de l’inscription « \teiHi[rend={italic}]{SEDES SVMMORVM PONTIFCVM QVI AB AN. M.CCC.VII. PER PLVS QVAM LXX. ANN. AVENIONE ALTERA ROMA DEGENTES ORBI CHRISTIANO PRAEFERVNT}\teiNote[n={26},place={foot},xmlid={ftn61-8}]{\teiP[xmlid={p-407}]{Joseph Guérin, \teiHi[rend={italic}]{Panorama d’Avignon, de Vaucluse, du Mont-Ventoux et du Col-Longet suivi de quelques vues des Alpes françaises}, Avignon, Imprimerie de Guichard aîné, 1829, p. 86.}} »\teiHi[rend={italic}]{.} Dans la même idée, Clément VI renomma de nombreux sites de la ville afin de refléter ses origines romaines. En 1344, il donna le nom de \teiHi[rend={italic}]{Roma} à la chapelle Saint-Michel, située au-dessus de la tour de la Garde-robe du palais, et il construisit sa propre chapelle de Saint-Pierre pour les couronnements des papes\teiNote[n={27},place={foot},xmlid={ftn62-8}]{\teiP[xmlid={p-408}]{Joëlle Rollo-Koster, \teiHi[rend={italic}]{Avignon and its papacy 1309- 1417: Popes, Institutions and Society}, Londres, Rowman \& Littlefield, 2015, p. 85.}}. Cette tradition se poursuit toujours au \teiHi[rend={small-caps}]{xvi}\teiHi[rend={sup}]{e} siècle. Pierre Pansier relève qu’un ouvrage imprimé en 1500 à Avignon, dont nous ne conservons pas d’exemplaire, portait la mention : « \teiHi[rend={italic}]{Impressit Dominicus Anselmus Avenionensis, Avenione altera Roma ultimo kalendas martii 1500}\teiNote[n={28},place={foot},xmlid={ftn63-8}]{\teiP[xmlid={p-409}]{Pierre Pansier « Les débuts de l’imprimerie à Avignon (\teiHi[rend={small-caps}]{xv}\teiHi[rend={sup}]{e} et \teiHi[rend={small-caps}]{xvi}\teiHi[rend={sup}]{e} siècles) », \teiHi[rend={italic}]{Mémoires de l’Académie de Vaucluse}, vol. XIX, n\teiHi[rend={sup}]{o} 2, 1919, p. 153-179, ici p. 166. Il s’agit d’un traité d’un jurisconsulte italien, Odofredo de Bénévent. }} ». Cette politique identitaire se poursuit encore au \teiHi[rend={small-caps}]{xvii}\teiHi[rend={sup}]{e} siècle et l’on trouve de nombreuses références visuelles au statut particulier de la ville, comme en atteste le frontispice du \teiHi[rend={italic}]{Recueil de statuts et de privilèges de la ville d’Avignon }(Avignon, bibliothèque municipale, ms. 2834, \teiHi[rend={small-caps}]{xiv}\teiHi[rend={sup}]{e}-\teiHi[rend={small-caps}]{xv}\teiHi[rend={sup}]{e} siècles et \teiHi[rend={small-caps}]{xvii}\teiHi[rend={sup}]{e} siècle), où l’expression SPQA apparaît sur des étendards rouges à la place du traditionnel SPQR\teiNote[n={29},place={foot},xmlid={ftn64-7}]{\teiP[xmlid={p-410}]{Voir Victor de Baumefort, \teiHi[rend={italic}]{Cession de la ville et de l’état d’Avignon au pape Clément VI par Jeanne I}\teiHi[rend={italic sup}]{re}\teiHi[rend={italic}]{, reine de Naples}, Apt, Typographie et lithographie J.-S. Jean, 1873, p. 59 et 60.}} (fig. 133). Deux gerfauts, emblèmes de la cité, entourent les anciennes armoiries de la ville (toutefois sans le dôme crucigère sur la tour du milieu) complétées du texte « \teiHi[rend={italic}]{SENATUS POPULUSQUE AVENIONENSIS} ». Des devises sont également inscrites dans des cartouches : « \teiHi[rend={italic}]{FELI\lbrack{}CIUM\rbrack{} TEMPORUM / REPARATIO », « UNGUIBUS ET ROSTRO / ALISQUE ARMATUS IN / HOSTEM} » avec les lettres « \teiHi[rend={italic}]{R.A} », soit la devise de la ville avec l’expression \teiHi[rend={italic}]{respublica avenionensis}, et « \teiHi[rend={italic}]{SINGULIS VARIUS / UTILIS OMNIBUS} » (différent pour chacun, avantageux pour tous »). L’expression « \teiHi[rend={italic}]{Felicium temporum reparatio} » semble faire allusion à des pièces de monnaie à l’effigie de différents empereurs romains à partir du règne de Constantin\teiNote[n={30},place={foot},xmlid={ftn65-7}]{\teiP[xmlid={p-411}]{Voir Harold Mattingly, « “Fel. Temp. Reparatio” », \teiHi[rend={italic}]{The Numismatic Chronicle and Journal of the Royal Numismatic Society}, vol. 5, n\teiHi[rend={sup}]{o} 13, 1933, p. 182-202.}}.}
              \teiP[xmlid={p11-10}]{Cette conscience civique s’étend bien au-delà du \teiHi[rend={small-caps}]{xvi}\teiHi[rend={sup}]{e} siècle, comme en attestent deux gravures, peu étudiées. La première est un frontispice (fig. 134), gravé pour un ouvrage non déterminé, conservé à la Bibliothèque nationale de France, qui fut ré-employé en 1648 pour l’\teiHi[rend={italic}]{Avenio Christiania} de Joseph-Marie Suarès, bibliothécaire pontifical. Il a cependant été réalisé en 1598 au plus tôt, car il porte non seulement la signature de Jean Bœuf « \teiHi[rend={italic}]{Joannes Bouis Avenionensis fec.} » en bas, un graveur buriniste actif à Avignon entre l’extrême fin du \teiHi[rend={small-caps}]{xvi}\teiHi[rend={sup}]{e} siècle et 1629, date de sa mort\teiNote[n={31},place={foot},xmlid={ftn66-6}]{\teiP[xmlid={p-412}]{Camille Larraz, « La réception tardive de l’art bellifontain dans le Comtat Venaissin : le cas du graveur Jean Bœuf (\teiHi[rend={small-caps}]{xvi}\teiHi[rend={sup}]{e}-\teiHi[rend={small-caps}]{xvii}\teiHi[rend={sup}]{e} siècles) », \teiHi[rend={italic}]{Nouvelles de l’estampe}, n\teiHi[rend={sup}]{o} 268, décembre 2022, \teiRef[target={https://doi.org/10.4000/estampe.3641}]{https://doi.org/10.4000/estampe.3641}.}}, mais aussi les armoiries du pape Clément VIII (1592-1605), du cardinal Ottavio Acquaviva (1591-1612), légat d’Avignon, et de l’archevêque Jean François Bordini\teiNote[n={32},place={foot},xmlid={ftn67-6}]{\teiP[xmlid={p-413}]{Marcel Chossat, \teiHi[rend={italic}]{Les Jésuites et leurs œuvres à Avignon : 1553-1768}, Avignon, François Seguin, 1896, p. 116 et 117.}} (1598-1609). Nous trouvons également les anciennes et nouvelles armoiries de la ville d’Avignon, une référence évidente au passé de la ville avant le séjour des papes, car elles furent modifiées par Clément VI. Ces blasons ornent un portique aux colonnes torsadées d’ordre corinthien agrémenté de nombreuses figures saintes et éléments symboliques. Face à face se situent des statues de saint Ruf et de sainte Marthe, qui sont identifiables grâce à des inscriptions latines. Patronne de la ville, Marthe est considérée comme la première apôtresse d’Avignon, qu’elle avait évangélisée, faisant ainsi parfaitement pendant à Ruf, premier évêque de la cité. En haut de l’attique se trouve la figure de la ville d’Avignon, « \teiHi[rend={italic}]{secu\lbrack{}n\rbrack{}da sede sap\lbrack{}ientiae\rbrack{}} », associée à la figure de la Vierge sur un trône, symbolisant la sagesse, tenant une tiare papale et encadrée par la Foi et la Vérité qui surmontent elles-mêmes la devise de la ville. Enfin, le trône de la figure d’Avignon est placé sous le nom de Jésus. Commandé par les jésuites, ainsi que nous le confirme l’inscription abrégée « \teiHi[rend={italic}]{SCHOL. RHETOR. AVEN. SOC. IESU. P. S. C. P. A. D.} », il est vraisemblable que le frontispice ait été commandé dans le cadre de l’entrée du cardinal et archevêque Bordini à Avignon en 1598\teiNote[n={33},place={foot},xmlid={ftn68-6}]{\teiP[xmlid={p-414}]{\teiHi[rend={italic}]{Ibid}, p. 117.}}. En effet, l’ordre était régulièrement sollicité par le conseil de ville depuis la fin du \teiHi[rend={small-caps}]{xvi}\teiHi[rend={sup}]{e} siècle pour organiser ces somptueuses fêtes\teiNote[n={34},place={foot},xmlid={ftn69-6}]{\teiP[xmlid={p-415}]{Olivier Rouchon, art. cité.}}. Le frontispice aurait donc pu servir de projet au livret qui aurait dû être imprimé pour relater l’événement. À nouveau, nous constatons qu’Avignon se réclame d’ascendance romaine et papale par un discours visuel qui honore ce lien avec Rome, autant par les armoiries de ses dirigeants, c’est-à-dire le pape, l’archevêque et le légat, que par les figures saintes et allégoriques qui font allusion au passé et au présent chrétien de la ville, iconographie renforcée par les anciennes et nouvelles armes d’Avignon.}
              \teiP[xmlid={p12-10}]{La seconde gravure (fig. 135), illustrant le \teiHi[rend={italic}]{Bullarium civitatis Avenionensis seu Bullae ac constitutiones apostolicae summorum pontificum et diplomata regum}, paru à Lyon en 1657, représente précisément cette vision d’\teiHi[rend={italic}]{Altera Roma}. L’ouvrage recueille les statuts et privilèges, privilèges qui se trouvaient dans les bulles des Papes, dont les originaux risquaient davantage la destruction\teiNote[n={35},place={foot},xmlid={ftn70-4}]{\teiP[xmlid={p-416}]{Pierre Charpenne, \teiHi[rend={italic}]{Histoire de réunions temporaires d’Avignon et du Comtat Venaissin à la France}, Paris, Calmann Lévy, 1886, p. 64.}}. Le conseil municipal d’Avignon demanda ainsi au pape de les confirmer afin de les faire imprimer, dans ce recueil de bulles et autres documents concernant les privilèges de la ville. On observe sur la gravure des personnifications allégoriques féminines des concepts d’\teiHi[rend={italic}]{Altera Roma} et d’\teiHi[rend={italic}]{Avinio Cavarum}, c’est-à-dire de Rome et d’Avignon (cette dernière était considérée comme la ville de la Narbonnaise). Installées sur des piédestaux, elles se font face et entourent une représentation de la ville avec le palais des Papes trônant au centre, du Rhône et du pont Saint-Bénezet. Cette vue est encadrée des armoiries de la ville et de celles du pape Alexandre VII (1655-1667). On trouve également des anges portant des phylactères aux textes suivants : « \teiHi[rend={italic}]{amica clavium veritas} » et « \teiHi[rend={italic}]{ungibus et rostro fidelitas} », soit « la vérité est amie des clés » et « fidélité de bec et ongles » (« \teiHi[rend={italic}]{ungibus et rostro} » est la devise d’Avignon, d’où les deux gerfauts entourant son blason). La figure de l’\teiHi[rend={italic}]{Altera Roma} est nue, portant un simple voilage sur son avant-bras gauche retombant pudiquement sur le bas du corps. Elle désigne de sa main droite son collier en forme de cœur, et sa main gauche tient un soleil, souvent porté par la figure de la vérité dans les livres d’emblèmes, faisant ainsi probablement référence à la devise « \teiHi[rend={italic}]{amica clavium veritas} ». Quant à la représentation de l’\teiHi[rend={italic}]{Avinio Cavarum}, elle est vêtue d’une robe évoquant un modèle antique, avec à ses pieds un aigle. Elle arbore un chien autour du cou, symbole de fidélité qui doit faire allusion à la devise de la cité, et porte une clé dans sa main droite, probablement celle de la ville. Avignon est donc ici mise sur un pied d’égalité avec Rome, justifiant sa demande au pape de faire imprimer les statuts et privilèges qu’il lui accorde. C’est donc grâce à l’image, et surtout le discours héraldique et allégorique, qu’Avignon peut tisser ce lien immuable avec la cité pontificale, dont elle est en quelque sorte le substitut.}
            
\end{teiDiv}
            \begin{teiDiv}[type={section1},xmlid={div04-5}]

              \teiHead[xmlid={conclusion-001}]{Conclusion}
              \teiP[xmlid={p13-10}]{Cette étude vise à démontrer que la question de la légitimation de la légation avignonnaise passe en grande partie par l’image, notamment par l’héraldique, mais aussi par les devises et les figures emblématiques, déployées lors de cérémonies grandioses, et par l’imprimé, dont on connaît bien le pouvoir de diffusion. Le pouvoir pontifical, malgré la distance géographique et politique, se construit à travers le concept de l’\teiHi[rend={italic}]{Altera Roma} qui vise à faire d’Avignon un autre centre névralgique de la papauté. La communication liée à cette représentation repose sur des outils visuels destinés à promouvoir la légation. Il s’agit donc de créer un lien de familiarité dans l’imaginaire collectif avignonnais, un objectif qui devient probablement d’autant plus important dans les années 1550 et à partir de la Contre-Réforme pour confirmer le statut d’Avignon comme défenseur du catholicisme face à un Languedoc et un Dauphiné protestants. Cette volonté affichée d’apparenter Avignon à Rome se remarque également à travers les arts avignonnais de façon générale, surtout à partir des années 1525, avec des peintres français qui séjournent à Rome avant de s’établir à Avignon comme Henri Guigues, qui importe de façon précoce les modèles de Raphaël en Provence, ou Simon de Châlons, qui imite le Sanzio et Michel-Ange tout au long de sa carrière avignonnaise\teiNote[n={36},place={foot},xmlid={ftn71-4}]{\teiP[xmlid={p-417}]{Camille Larraz, \teiHi[rend={italic}]{op. cit}.}}. La pérennité et le succès que rencontrent les modèles romains à Avignon doivent sans doute à la politique de l’\teiHi[rend={italic}]{Altera Roma}, qui veut elle aussi faire de ses peintres un « autre Raphaël ».}
              \begin{teiFigure}[xmlid={figure01-9}]

                \teiGraphic[url={../icono/br/Ch10\_Larraz/Fig132.jpg}]
                \teiHead[xmlid={figure-129-bernard-salomon-dans-la-magnificque-et-triumphante-entree-de-carpentras-faicte-a-tresillustre-trespuisssant-prince-alexandre-farnes-cardinal-legat-davignon-vicechancelier-du-s-siege-apostolique-par-antoine-blegier-avignon-par-mace-bonhomme-aux-changes-1553-p-da-verso-collection-particuliere-camille-larraz-001}]{Figure 129. Bernard Salomon, dans \teiHi[rend={italic}]{La Magnificque et triumphante entrée de Carpentras, faicte à tresillustre \& trespuisssant prince ALEXANDRE FARNES, Cardinal, Legat d’Avignon, Vicechancelier du S. siege Apostolique} par Antoine Blégier, Avignon, Par Macé Bonhomme, aux Changes, 1553, p. Da verso, collection particulière, Camille Larraz.}
              
\end{teiFigure}
              \begin{teiFigure}[xmlid={figure02-9}]

                \teiGraphic[url={../icono/br/Ch10\_Larraz/Fig133.jpg}]
                \teiHead[xmlid={figure-130-bernard-salomon-dans-la-magnificque-et-triumphante-entree-de-carpentras-par-antoine-blegier-avignon-1553-p-d-collection-particuliere-camille-larraz-001}]{Figure 130. Bernard Salomon, dans \teiHi[rend={italic}]{La Magnificque et triumphante entrée de Carpentras… }par Antoine Blégier, Avignon, 1553, p. D, collection particulière, Camille Larraz.}
              
\end{teiFigure}
              \begin{teiFigure}[xmlid={figure03-9}]

                \teiGraphic[url={../icono/br/Ch10\_Larraz/Fig134.jpg}]
                \teiHead[xmlid={figure-131-avignon-archives-municipales-cc-468-comptes-de-gaspard-droin-depenses-extraordinaires-non-numerote-collection-particuliere-camille-larraz-001}]{Figure 131. Avignon, Archives municipales, CC 468, Comptes de Gaspard Droin, dépenses extraordinaires, non numéroté, collection particulière, Camille Larraz.}
              
\end{teiFigure}
              \begin{teiFigure}[xmlid={figure04-9}]

                \teiGraphic[url={../icono/br/Ch10\_Larraz/Fig135.jpg}]
                \teiHead[xmlid={figure-132-bernard-salomon-dans-la-magnificque-et-triumphante-entree-de-carpentras-par-antoine-blegier-avignon-1553-p-b-ii-collection-particuliere-camille-larraz-001}]{Figure 132. Bernard Salomon, dans \teiHi[rend={italic}]{La Magnificque et triumphante entrée de Carpentras… }par Antoine Blégier, Avignon, 1553, p. B ii, collection particulière, Camille Larraz.}
              
\end{teiFigure}
              \begin{teiFigure}[xmlid={figure05-9}]

                \teiGraphic[url={../icono/br/Ch10\_Larraz/Fig136.png}]
                \teiHead[xmlid={figure-133-anonyme-recueil-de-statuts-et-de-privileges-de-la-ville-davignon-avignon-bm-ms-2884-001}]{Figure 133. Anonyme, \teiHi[rend={italic}]{Recueil de statuts et de privilèges de la ville d’Avignon}, Avignon, BM, ms. 2884.}
              
\end{teiFigure}
              \begin{teiFigure}[xmlid={figure06-9}]

                \teiGraphic[url={../icono/br/Ch10\_Larraz/Fig137.jpg}]
                \teiHead[xmlid={figure-134-jean-boeuf-frontispice-vers-1598-paris-bnf-latin-8971-001}]{Figure 134. Jean Bœuf, \teiHi[rend={italic}]{Frontispice}, vers 1598. Paris, BNF, Latin 8971.}
              
\end{teiFigure}
              \begin{teiFigure}[xmlid={figure07-8}]

                \teiGraphic[url={../icono/br/Ch10\_Larraz/Fig138.jpg}]
                \teiHead[xmlid={figure-135-anonyme-page-de-titre-du-bullarium-civitatis-avenionensis-seu-bullae-lyon-joannis-amati-candy-1657-paris-bnf-departement-philosophie-histoire-sciences-de-lhomme-fol-lk7-647-001}]{Figure 135. Anonyme, page de titre du \teiHi[rend={italic}]{Bullarium civitatis Avenionensis seu Bullae}, Lyon, Joannis Amati Candy, 1657. Paris, BNF, département Philosophie, histoire, sciences de l’homme, FOL-LK7-647.}
              
\end{teiFigure}
            
\end{teiDiv}
          
\end{teiElement}
        
\end{teiElement}
        \begin{teiElement}[name={group},type={chapter},data-page-title={Sainte-Prisca-sur-l’Aventin},data-page-subtitle={Notes d’héraldique pontificale et cardinalice},data-page-authors={Rocco Ronzani},data-include-href={Ch11\_Sainte\_Prisca\_sur\_Aventin.xml},xmlid={chapter-008}]

          \begin{teiElement}[name={front}]

            \begin{teiDiv}[type={titlePage},xmlid={titlepage-013}]

              \teiP[rend={title-main},xmlid={p-418}]{Sainte-Prisca-sur-l’Aventin}
              \teiP[rend={title-sub},xmlid={p-419}]{Notes d’héraldique pontificale et cardinalice}
              \teiP[rend={author-aut},xmlid={p-420}]{Rocco Ronzani}
            
\end{teiDiv}
          
\end{teiElement}
          \begin{teiElement}[name={body},xmlid={body-013}]

            \begin{teiDiv}[type={section1},xmlid={div01-10}]

              \teiHead[xmlid={remarques-sur-le-monument-001}]{Remarques sur le monument}
              \teiP[xmlid={p1-13}]{Dans cette contribution, il n’est pas possible de résumer l’histoire pluriséculaire de l’antique édifice chrétien de l’Urbe situé sur le mont Aventin. Des études approfondies lui ont été consacrées, il mériterait toutefois d’être à nouveau examiné sous de nombreux aspects : archéologique et architectural, mais aussi historique et artistique\teiNote[n={1},place={foot},xmlid={ftn1-11}]{\teiP[xmlid={p-421}]{Isidoro Carini, \teiHi[rend={italic}]{Sul titolo presbiterale di Santa Prisca}, Palerme, Tip. Camillo Tamburello, 1885 ; Johann Peter Kirsch, \teiHi[rend={italic}]{Die römischen Titelkirchen im Altertum}, Paderborn, Druck ind Verlag von Ferdinand Schöningh, 1918, p. 101-104 ; Francesco Lanzoni, « I titoli presbiterali di Roma antica nella storia e nella leggenda », \teiHi[rend={italic}]{Rivista di archeologia cristiana},n\teiHi[rend={sup}]{o} 2, 1925, p. 247-250 ; Geremia Sangiorgi, \teiHi[rend={italic}]{S. Prisca e il suo mitreo}, Rome, Marietti, « Le chiese di Roma illustrate, 101 », 1968 ; Mariano Armellini et Carlo Cecchelli, \teiHi[rend={italic}]{Le chiese di Roma dal secolo IV al XIX}, Rome, RORE, 1942, t. 2 ; \teiHi[rend={italic}]{Monasticon Italiae}, Cesena, Badia di Santa Maria del Monte, 1981, t. 1 : \teiHi[rend={italic}]{Roma e Lazio}, p. 73 et 74 ; Mario Bevilacqua et Daniela Gallavotti Cavallero (dir.), \teiHi[rend={italic}]{L’Aventino dal Rinascimento a oggi. Arte e Architettura}, Rome, Artemide, 2010 ; Rocco Ronzani, « Nota storico-iconografica sulla committenza del cardinale Benedetto Giustiniani in Santa Prisca all’Aventino », \teiHi[rend={italic}]{Studi Romani}, n. s. III, 1-4, 2021, p. 21-32.}}.}
              \teiP[xmlid={p2-12}]{Si, de fait, sa préexistence classique a été amplement étudiée, spécialement après la découverte d’un mithréum du \teiHi[rend={small-caps}]{iii}\teiHi[rend={sup}]{e} siècle, situé dans sa partie souterraine, il reste beaucoup à rechercher sur l’histoire du monument et des institutions religieuses qui au cours des siècles se sont succédées à la garde du lieu sacré et des structures édifiées pour l’abriter\teiNote[n={2},place={foot},xmlid={ftn20}]{\teiP[xmlid={p-422}]{Je me réfère en premier lieu au haut Moyen Âge, \teiHi[rend={italic}]{monasterium sancti Donati qui ponitur iuxta titulum sanctae Priscae,} évoqué pour la première et unique fois dans le \teiHi[rend={italic}]{Liber Pontificalis,} à la biographie de Léon III (795‑816), mais aussi à la présence bénédictine prolongée (italienne et française : les bénédictins de Vendôme) et enfin aux autres nombreuses communautés religieuses qui ont gardé le lieu jusqu’à présent. Voir \teiHi[rend={italic}]{Liber Pontificalis}, Louis Duchesne (éd.), Paris, E. de Boccard, 1956, t. 2, p. 20.}}.}
              \teiP[xmlid={p3-12}]{En outre, il paraît nécessaire de revoir le complexe dossier hagiographique de la sainte titulaire qui, comme on le suppose, à une fondatrice homonyme du \teiHi[rend={italic}]{titulus} tardo-antique\teiNote[n={3},place={foot},xmlid={ftn21}]{\teiP[xmlid={p-423}]{Il s’agissait peut-être d’une \teiHi[rend={italic}]{Aquilia Prisca}, nom bien attesté dans l’épigraphie, d’où la confusion avec le couple apostolique qui survint ensuite. Voir Giovanni Battista de Rossi, « Priscilla e gli Acilii Glabrioni. § VIII. Aquila e Prisca e gli Acilii Glabrioni », \teiHi[rend={italic}]{Bullettino di archeologia cristiana}, n\teiHi[rend={sup}]{o} 6, n\teiHi[rend={sup}]{o} 3-4, 1888-1889, p. 128-133, ici p. 131 et 132.}} superpose la figure d’une martyre vénérée sur les lieux – sa chronologie oscillant entre \teiHi[rend={small-caps}]{i}\teiHi[rend={sup}]{er} et \teiHi[rend={small-caps}]{iii}\teiHi[rend={sup}]{e} siècle – dont les reliques furent vraisemblablement transférées depuis la Via Ostiensis et à laquelle fut consacrée une tardive\teiHi[rend={italic}]{ Passio}\teiNote[n={4},place={foot},xmlid={ftn22}]{\teiP[xmlid={p-424}]{Voir \teiHi[rend={italic}]{Corinthiens}, 16:19 ; \teiHi[rend={italic}]{Romains}, 16:3-5 ; \teiHi[rend={italic}]{Actes}, 18:1-15, 18-19 et 26 ; \teiHi[rend={italic}]{Bibliotheca Hagiographica Latina, antiquae et mediae aetatis}, Bruxelles, Société des Hollandistes, 1900-1901, n\teiHi[rend={sup}]{o} 6926, p. 1008 ;\teiHi[rend={italic}]{ Acta Sanctorum Ianuarii}, Antverpiae, Apud Joannem Meursium, 1643, t. 2, p. 184-187 ; Iohannes Baptista de Rossi et Ludovicus Duchesne, \teiHi[rend={italic}]{Acta Sanctorum Novembris}, Bruxelles, 1894, vol. 2, part. 1 : \teiHi[rend={italic}]{Martyrologium Hieronymianum} ; Angelo Di Berardino, « Prisca, santa », dans \teiHi[rend={italic}]{Nuovo dizionario patristico e di antichità cristiane}, Gênes-Milan, Marietti, 2008, vol. 3, col. 4331-4332.}}. Ensuite, certainement entre \teiHi[rend={small-caps}]{viii}\teiHi[rend={sup}]{e} et \teiHi[rend={small-caps}]{ix}\teiHi[rend={sup}]{e} siècle, l’église de l’Aventin fut aussi liée aux saints Aquila et Prisca (ou Priscilla), le couple d’époque apostolique cité dans les lettres de saint Paul et dans les \teiHi[rend={italic}]{Actes des apôtres}. Une antique inscription métrique située sur la porte du sanctuaire commençait ainsi : « \teiHi[rend={italic}]{Haec domus est Aquilae seu Priscae uirginis almae} » \lbrack{}Cette demeure est à Aquila ou à Prisca, la douce vierge\teiNote[n={5},place={foot},xmlid={ftn23}]{\teiP[xmlid={p-425}]{Voir Giovanni Battista de Rossi, « Della casa d’Aquila e Prisca all’Aventino », \teiHi[rend={italic}]{Bullettino di archeologia cristiana}, vol. 3, 1867, p. 44-47.}}\rbrack{}.}
              \teiP[xmlid={p4-12}]{De tels développements hagiographiques ont contribué de manière décisive à faire de l’église de l’Aventin une sorte de « mémorial apostolique » : saint Paul y aurait peut-être demeuré, en alternance avec la \teiHi[rend={italic}]{domus} de la Regola, et certainement saint Pierre y aurait exercé son ministère, célébrant l’eucharistie pour une des plus antiques communautés chrétiennes de Rome, baptisant en outre Prisca, la martyre, qui, dans cette version de la légende hagiographique se trouve par conséquent rétrodatée au \teiHi[rend={small-caps}]{i}\teiHi[rend={sup}]{er} siècle.}
              \teiP[xmlid={p5-12}]{La valorisation de la mémoire apostolique de l’Aventin est aussi liée à l’existence plus évidente de témoignages héraldiques présents dans l’édifice sacré comme dans sa crypte, réalisée et décorée de fresques pour permettre aux pèlerins de rejoindre le sous-sol qui, aux yeux des ecclésiastiques députés à la conservation du lieu de culte, conservait vraisemblablement l’un des vestiges chrétiens les plus anciens de Rome. Ce n’est pas un hasard si l’épisode du baptême de Prisca, administré par le prince des apôtres, fut un des motifs les plus valorisés à l’occasion du Jubilé de 1600\teiNote[n={6},place={foot},xmlid={ftn24}]{\teiP[xmlid={p-426}]{C’est dans ce contexte qu’il faut replacer la valorisation des fonts baptismaux de l’église, réalisés dans un beau chapiteau dorique, peut-être au \teiHi[rend={small-caps}]{ii}\teiHi[rend={sup}]{e} siècle, sur lequel au Moyen Âge fut gravée l’inscription suivante en commémoration du ministère de Pierre sur l’Aventin : \teiHi[rend={italic}]{« bactismu(m) s(an)c(t)i Petri} ». Aujourd’hui, les fonts baptismaux sont situés dans une chapelle latérale de l’église d’après l’aménagement conçu par l’ingénieur Arturo Hoerner, procurateur du chapitre de Saint-Pierre, et par le sculpteur carrarais Antonio Biggi. Voir Mariano Armellini, \teiHi[rend={italic}]{Le chiese di Roma, dal secolo IV al XIX}, Rome, Tipografia vaticana, 1891, p. 578 ; Alfonso Camillo De Romanis, \teiHi[rend={italic}]{Per un nuovo Battistero a Santa Prisca. Nova et vetera}, Cité du Vatican, Tipografia poliglotta vaticana, 1948.}}. Le commanditaire des travaux jubilaires fut Benedetto Giustiniani\teiNote[n={7},place={foot},xmlid={ftn25}]{\teiP[xmlid={p-427}]{Voir Simona Feci et Luca Bortolotti, « Giustiniani, Benedetto », dans \teiHi[rend={italic}]{DBI}, Rome, Istituto della Enciclopedia italiana, 2001, vol. 57, p. 366-377.}}, cardinal du titre presbytéral de 1599 à 1611 (fig. 136). L’intervention de Giustiniani s’insère dans le cadre de la redécouverte et de la promotion des restes paléochrétiens de l’Urbe, centre de la catholicité, afin de légitimer la primauté de l’Église romaine et de démontrer la continuité comme la cohérence de son patrimoine liturgique et cultuel, doctrinal et disciplinaire, le message évangélique transmis depuis les premières générations chrétiennes assumant ainsi un rôle pertinent dans la polémique antiprotestante. D’analogues initiatives furent déployées non seulement par la monumentalisation des reliques du christianisme romain des origines, mais aussi et surtout dans le domaine de la recherche archéologique, dans la codification liturgique, dans les nombreuses éditions des pères de l’Église et dans les grandes synthèses d’érudition historique, parmi lesquelles il faut en particulier signaler les \teiHi[rend={italic}]{Annales} de Cesare Baronio (1538-1607), publiées entre 1588 et 1607, dans les années de la restauration de Sainte-Prisca, restauration à laquelle le cardinal oratorien, lié à Giustiniani, ne dut pas être étranger. À Baronio est aussi due la relecture du dossier hagiographique de Prisca qui, initialement identifiée avec la martyre du \teiHi[rend={small-caps}]{iii}\teiHi[rend={sup}]{e} siècle – elle aurait été exécutée sous le règne de Claude le Gothique (268-270) –, fut rétrodatée par cet historien de l’Église au \teiHi[rend={small-caps}]{i}\teiHi[rend={sup}]{er} siècle, fixant sa mort à la treizième et ultime année du règne de Claude I (41-54), selon toute évidence pour valoriser la tradition de son baptême par saint Pierre.}
              \teiP[xmlid={p6-12}]{C’est le mérite de Daniela Gallavotti Cavallero d’avoir illustré le rôle central de la figure de saint Pierre dans le cycle de fresques commandé par Giustiniani à Sainte-Prisca\teiNote[n={8},place={foot},xmlid={ftn26}]{\teiP[xmlid={p-428}]{Daniela Gallavotti Cavallero, « L’iconografia di santa Prisca e l’immagine di Pietro dalla basilica Vaticana al titulus della santa sull’Aventino nei secoli XVI e XVII », dans Mario Bevilacqua et Daniela Gallavotti Cavallero (dir.), \teiHi[rend={italic}]{op. cit.}, p. 57-71.}}. En fait, au détriment de la sainte éponyme de l’édifice sacré, ces fresques sont principalement centrées sur la figure et sur le ministère du premier « pape\teiHi[rend={italic}]{ }», aussi bien le cycle de la \teiHi[rend={italic}]{confessio }que celui du presbytère de l’église, sans oublier le grand retable du maître autel dû à Domenico Cresti, dit le Passignano\teiNote[n={9},place={foot},xmlid={ftn27}]{\teiP[xmlid={p-429}]{Simonetta Prosperi Valenti Rodinò, « Cresti, Domenico, detto il Passignano », dans \teiHi[rend={italic}]{DBI}, \teiHi[rend={italic}]{op. cit}., 1984, vol. 30, p. 741-749. Le peintre est mentionné par Giovan Battista Marino dans l’\teiHi[rend={italic}]{Adone} et sa patrie, Passignano, est le siège d’une importante abbaye vallombrosane. }} (1559-1638), qui représente le baptême de Prisca par le prince des apôtres (fig. 137). En complément des références iconographiques déjà relevées par la savante, je voudrais attirer l’attention sur l’abside où la fresque qui représente actuellement saint Augustin (fig. 138), fruit d’un maladroit repeint situable entre le \teiHi[rend={small-caps}]{xix}\teiHi[rend={sup}]{e} siècle et le début du \teiHi[rend={small-caps}]{xx}\teiHi[rend={sup}]{e} siècle, supplante certainement l’image originale d’un saint Pierre en chaire auquel renvoient clairement les symboles pontificaux tenus par des anges : la tiare, la férule pontificale et les clefs pétriniennes (fig. 138 à 140).}
              \teiP[xmlid={p7-12}]{La théorie des anges et des saints des fresques de la nef centrale de l’église – où réapparaît la figure de Pierre, accompagné de Paul et des apôtres André et Jean – a été étudiée par Gallavotti Cavallero, mais il reste à élucider la présence de deux saints identifiés par la chercheuse comme saint Benoît et saint Antoine abbé. Le premier (fig. 141) aurait été « choisi en tant que saint éponyme du commanditaire », le cardinal Benedetto Giustiniani ; pour l’autre, Gallavotti Cavallero ne trouve pas de motif évident\teiNote[n={10},place={foot},xmlid={ftn28}]{\teiP[xmlid={p-430}]{Daniela Gallavotti Cavallero, « L’iconografia di santa Prisca e l’immagine di Pietro…\teiHi[rend={italic}]{ »}, art. cité, p. 64. Il n’y a pas de doute que le premier saint représente le patriarche de la vie monastique occidentale Benoît de Nurcie ; son iconographie est tout à fait traditionnelle : le saint corrige l’indiscipline monastique de sa verge, inspirée de l’hymne \teiHi[rend={italic}]{Signifer invictissime} – attribué au moins en partie à saint Pierre Damien – qui résume en vers les épisodes saillants de la \teiHi[rend={italic}]{vita Benedicti} au second livre des \teiHi[rend={italic}]{Dialogues} de Grégoire le Grand. La strophe de référence dans l’hymne est la suivante : « \teiHi[rend={italic}]{Frater quem tunc nequissimum / vagum raptabat spiritum / dum tua virga ceditur / stabilitati redditur.} »}}. En réalité, le second saint (fig. 142) n’est pas identifiable comme Antoine, mais plutôt comme le réformateur bénédictin Jean Gualbert (995-1073), saint initiateur de la réforme de Vallombreuse. La figure est en fait représentée suivant une iconographie tout à fait traditionnelle : avec l’habit monastique des vallombrosains, tenant de sa main le bâton en forme de tau dont les deux branches s’achèvent par de caractéristiques protomes léonins, qui est aussi la férule héraldique représentée dans les armes de la congrégation de Vallombreuse\teiNote[n={11},place={foot},xmlid={ftn29}]{\teiP[xmlid={p-431}]{Sur les Vallombrosains voir Francesco Salvestrini, « Il patrimonio fondiario del monastero di Vallombrosa fra XIII e XVI secolo: presenza e utilizzazione del bosco », dans Simonetta Cavaciocchi (dir), \teiHi[rend={italic}]{L’uomo e la foresta, secc. XIII-XVIII}, Prato, Istituto Internazionale di Storia Economica Francesco Datini, 1996, p. 1057-1068 ; « L’apporto dei Vallombrosani e dei Camaldolesi all’edificazione della marina toscana (seconda metà del XVII-anni ’20 del XVIII secolo) », \teiHi[rend={italic}]{Archivio Storico Italiano}, vol. 156, n\teiHi[rend={sup}]{o} 2, 1998, p. 307-329 ; \teiHi[rend={italic}]{Santa Maria di Vallombrosa. Patrimonio e vita economica di un grande monastero medievale}, Florence, Olschki, « Biblioteca Storica Toscana, I, 33 », 1998 ; « La storiografia sul movimento e sull’Ordine monastico di Vallombrosa », \teiHi[rend={italic}]{Quaderni Medievali,} n\teiHi[rend={sup}]{o} 53, 2002, p. 294-323 ; « L’esperienza di Vallombrosa nella documentazione archivistica (secoli XI-XVI) », dans Giancarlo Andenna et Renata Salvarani (dir), \teiHi[rend={italic}]{La memoria dei chiostri, Atti delle prime Giornate di studi medievali, Laboratorio di storia monastica dell’Italia settentrionale}, Castiglione delle Stiviere, 11-13 octobre 2001, Brescia, CESIMB, 2002, p. 215-230 ; « L’art et la magnificence “contre” le pouvoir du prince. L’abbé général Biagio Milanesi, l’ordre monastique de Vallombreuse, Laurent de Médicis et le pape Léon X », dans Élisabeth Crouzet-Pavan et Jean-Claude Maire Vigueur (dir), \teiHi[rend={italic}]{L’art au service du prince. Paradigme italien, expériences européennes (vers 1250-vers 1550)}, actes du colloque, Paris, Université de Paris-Sorbonne, Centre Roland-Mousnier, 17-19 octobre 2013, Rome, Viella, « Italia comunale e signorile, 8 », 2015, p. 253-279 ; Francesco Salvestrini et Rosagemma Ciliberti, \teiHi[rend={italic}]{I Vallombrosani nel Piemonte medievale e moderno. Ospizi e monasteri intorno alla strada di Francia}, Rome, Viella, 2014 ; Sofia Boesch Gajano, « \teiHi[rend={italic}]{Storia e tradizione vallombrosane} », dans Antonella Degl’Innocenti (dir.), \teiHi[rend={italic}]{Vallombrosa. Memorie agiografiche e culto delle reliquie}, Rome, Viella, 2012, p. 15-115. Il faut aussi remarquer que le peintre Domenico Cresti, auteur du retable de l’église romaine, avait toujours été actif dans les chantiers toscans des Vallombrosains et que dans la dernière décennie du \teiHi[rend={small-caps}]{xvi}\teiHi[rend={sup}]{e }siècle il avait peint à fresque le cycle gualbertien dans la chapelle du saint à la Sainte-Trinité de Florence (1593-1594). Par ailleurs, jusqu’à l’époque contemporaine l’autel latéral de la nef droite de l’église était dédié au saint vallombrosain, remplacé seulement dans la première moitié du \teiHi[rend={small-caps}]{xx}\teiHi[rend={sup}]{e} siècle par la sainte agostinienne Rita de Cascia avec un retable déjà vénéré dans l’église située au pied du Capitole, démolie puis reconstruite près de la place Campitelli.}} (fig. 143).}
              \teiP[xmlid={p8-12}]{Les figures de Benoît et de Jean Gualbert dans les cycles des fresques de Sainte-Prisca renvoient au milieu bénédictin vallombrosain et à un projet précis, longuement caressé par le cardinal Giustiniani, jusqu’à présent ignoré des études sur le monument de l’Aventin, celui d’installer à Sainte-Prisca une communauté de vallombrosains, ordre duquel le cardinal était aussi le protecteur près le Saint-Siège. Comme le montre le destin postérieur du complexe religieux et la documentation d’archives\teiNote[n={12},place={foot},xmlid={ftn30}]{\teiP[xmlid={p-432}]{Comme me l’a aimablement indiqué Francesco Salvestrini, des informations sur le destin du complexe figurent dans les pièces D.V.3 et D.V.4, \teiHi[rend={italic}]{Scritture diverse pertinenti alla nostra congregazione}, de l’archive de l’abbaye de Vallombreuse. Par ailleurs, Salvestrini a aussi observé que, comme cela se produira au \teiHi[rend={small-caps}]{xviii}\teiHi[rend={sup}]{e} siècle à San Vigilio de Sienne, les Vallombrosains soulignaient immédiatement l’acquisition ou la récupération d’une fondation en commandant des œuvres d’art représentant leur père fondateur ou d’autres saints de leur obédience.}}, la tentative d’installation monastique n’eut pas de suite. Après l’échec de ce premier projet, entre 1606 et 1607, Giustiniani appela à résider à Sainte-Prisca les augustins de la congrégation de l’observance de Lombardie, présents à Rome depuis le \teiHi[rend={small-caps}]{xv}\teiHi[rend={sup}]{e} siècle à Sainte-Marie-du-Peuple et qui se maintinrent à la tête du complexe sans interruption jusqu’à l’époque napoléonienne, et jusqu’aujourd’hui non sans diverses péripéties\teiNote[n={13},place={foot},xmlid={ftn31}]{\teiP[xmlid={p-433}]{Voir Carlos Alonso, « La chiesa e il convento di santa Prisca a Roma. Vicende storiche, 1535-1724 », \teiHi[rend={italic}]{Analecta Augustiniana, }n\teiHi[rend={sup}]{o} 58, 1995, p. 215-263 ; Salvatore Enrico Anselmi, « Novità documentarie su S. Prisca in età barocca. L’attività di Carlo Lambardi e di Carlo Fontana », dans Mario Bevilacqua et Daniela Gallavotti Cavallero (dir.), \teiHi[rend={italic}]{op. cit}., p. 73-77. }}.}
              \teiP[xmlid={p9-11}]{Après s’être engagé dans la restauration de l’édifice sacré, il était tout à fait naturel que Giustiniani pensât à assurer pour les fidèles et surtout pour les pèlerins, le soin du culte divin par une présence stable, monastique ou religieuse, qui mieux que le clergé séculier pût garantir la continuité et le décorum du lieu saint.}
            
\end{teiDiv}
            \begin{teiDiv}[type={section1},xmlid={div02-10}]

              \teiHead[xmlid={lheraldique-cardinalice-et-pontificale-du-monument-001}]{L’héraldique cardinalice et pontificale du monument}
              \teiP[xmlid={p10-11}]{L’héraldique du cardinal Giustiniani est largement représentée dans l’église comme dans la crypte et le décor exploite constamment les meubles de son blason comme ceux des pontifes qui le voulurent membre du Sacré Collège et cardinal prêtre de la basilique aventinienne. Par ailleurs, ces éléments héraldiques, auxquels on n’a pas accordé toute l’importance qu’ils méritaient, permettent de dater pour ainsi dire à l’année près les fresques commandées par le cardinal.}
              \teiP[xmlid={p11-11}]{Il est connu que l’ambitieuse politique économique de Sixte Quint, pape de 1585 à 1590, trouva chez les banquiers génois Giustiniani un interlocuteur privilégié et, pour des raisons évidentes, le pape voulut les récompenser en favorisant la brillante carrière du fils ecclésiastique de la famille, le créant finalement cardinal le 16 décembre 1586. Ce lien étroit entre Benedetto Giustiniani et le pontife issu des Marches est exprimé dans le décor de la crypte de Sainte-Prisca où paraissent tantôt les armes complètes du pape (fig. 144, en haut, état très dégradé) tantôt leurs meubles séparés (le lion, fig. 145). Cela rappelle beaucoup les décors picturaux qu’il commanda pour les architectures de Domenico Fontana, dus à Cesare Nebbia et Giovanni Guerra, du monumental salon Sixtin de la Bibliothèque vaticane au complexe du Latran, sans oublier les édifices sacrés : chapelle Sixtine de Sainte-Marie-Majeure ou église Saint-Jérôme-des-Esclavons. Dans ces fresques sont sans cesse répétés les meubles héraldiques du pape : le lion, les poires (par allusion au nom de Peretti, armes parlantes), le \teiHi[rend={italic}]{trimonzio} \lbrack{}triple mont\rbrack{} et l’étoile ou comète\teiNote[n={14},place={foot},xmlid={ftn32}]{\teiP[xmlid={p-434}]{Marcello Fagiolo et Maria Luisa Madonna (dir.), \teiHi[rend={italic}]{Sisto V}, Rome, Istituto poligrafico e Zecca dello Stato, Libreria dello Stato, 1992 ; Maria Luisa Madonna (dir.), \teiHi[rend={italic}]{Roma di Sisto V. Le arti e la cultura}, Rome, De Luca, 1993 ; Silvano Giordano, « Sisto V », dans \teiHi[rend={italic}]{Enciclopedia dei Papi}, Rome, Treccani, 3 vol., t. 3, 2000, p. 202-222, ici p. 213 et 214 ; Massimo Ceresa (dir.), \teiHi[rend={italic}]{La Biblioteca Vaticana tra riforma cattolica, crescita delle collezioni e nuovo edificio (1535-1590)}, Cité du Vatican, BAV, « Storia della Biblioteca Apostolica Vaticana, 2 », 2012. Pour les aspects héraldiques, on se reportera à Yvan Loskoutoff, \teiHi[rend={italic}]{Un art de la Réforme catholique, La symbolique du pape Sixte-Quint et des Peretti-Montalto (1566-1655)}, Paris, Honoré Champion, 2011.}}. Les armes de Sixte Quint Peretti peuvent en effet se blasonner ainsi : « D’azur au lion d’or tenant une branche de poirier chargée de trois poires au naturel, à la bande voûtée de gueules brochant sur le tout, et chargée d’une comète d’or, et d’une montagne de trois coupeaux d’argent\teiNote[n={15},place={foot},xmlid={ftn33-2}]{\teiP[xmlid={p-435}]{Jean-Baptiste Rietstap, \teiHi[rend={italic}]{Armorial général}, Gouda, G. B. van Goor, 1861, « Peretti », p. 803. Giovanni Battista di Crollalanza donne pour sa part une description où il perd le caractère parlant de l’écu : « \teiHi[rend={italic}]{D’azzurro, al leone tenente colla branca anteriore destra un ramo di tre pomi, fogliato, il tutto d’oro; colla banda di rosso attraversante, caricata verso il capo da una cometa d’oro, e verso la punta da un monte di tre cime d’argento; il tutto nel verso della banda} », \teiHi[rend={italic}]{Dizionario storico-blasonico delle famiglie nobili e notabili italiane estinte e fiorenti}, Pise, Presso la direzione del giornale araldico, 1886-1890, 3 vol., t 2, « Peretti », p. 313.}}. »}
              \teiP[xmlid={p12-11}]{Comme c’était l’usage, les armes du pape furent intégrées au chef \lbrack{}tiers supérieur\rbrack{} de celles du cardinal qu’il avait créé (fig. 146 à 148). Celles-ci sont présentes dans les fresques, les inscriptions et dans les éléments décoratifs de marbre du maître autel et de l’église. On peut les décrire de la manière suivante : « Coupé : au 1 d’or à l’aigle naissante de sable ; au 2 de gueules au château d’argent\teiNote[n={16},place={foot},xmlid={ftn34-3}]{\teiP[xmlid={p-436}]{Jean-Baptiste Rietstap, \teiHi[rend={italic}]{Armorial général}, \teiHi[rend={italic}]{op. cit.}, « Giustiniani », p. 429. Chez Giovanni Battista Crollalanza, le blasonnement diffère dans la mesure où l’écu n’est pas coupé mais comporte un chef et où le château est décrit plus précisément : « \teiHi[rend={italic}]{Di rosso, al castello d’argento sormontato da tre torri, quella di mezzo più elevata, aperto e finestrato del campo ; col capo d’oro, all’aquila nascente di nero, coronata del campo} », \teiHi[rend={italic}]{Dizionario storico-blasonico, op. cit}., t. 1, « Giustiniani ».}}. » Comme pour Sixte Quint, outre les armoiries complètes, les meubles furent aussi utilisés séparément, le château (fig. 149) ou l’aigle (fig. 150 et 151).}
              \teiP[xmlid={p13-11}]{En ce qui regarde la datation du cycle de saint Pierre et de la sainte martyre, il est important de considérer la date de 1599, où le cardinal Giustiniani fut nommé titulaire de Sainte-Prisca par Clément VIII Aldobrandini (1592-1605), dont les armes figurent dans les fresques de la crypte (fig. 152), et la date de la mort du même pontife en 1605 : c’est dans ces années que peut se situer la réalisation des fresques, sans doute entre 1599 et 1600, date de l’achèvement des travaux et de l’inauguration ainsi que de la pose de la pierre commémorative du presbytère. Les armes de Clément VIII peuvent se décrire ainsi : « D’azur à la bande brétessée et contre-brétessée d’or, accostée de six étoiles du même, rangées en orle\teiNote[n={17},place={foot},xmlid={ftn35-3}]{\teiP[xmlid={p-437}]{Jean-Baptiste Rietstap, \teiHi[rend={italic}]{Armorial général}, \teiHi[rend={italic}]{op. cit}., « Aldobrandini », p. 58. Di Crollalanza donne : « \teiHi[rend={italic}]{D’azzurro, alla banda doppio-merlata d’oro, accompagnata da sei stelle di otto raggi dello stesso} », \teiHi[rend={italic}]{Dizionario storico-blasonico}, \teiHi[rend={italic}]{op. cit}., t. 1, p. 26, « Aldobrandini ». }}. »}
            
\end{teiDiv}
            \begin{teiDiv}[type={section1},xmlid={div03-8}]

              \teiHead[xmlid={autres-temoignages-heraldiques-interessants-presents-dans-leglise-001}]{Autres témoignages héraldiques intéressants présents dans l’église}
              \teiP[xmlid={p14-8}]{Pour conclure et compléter l’étude héraldique, on doit signaler dans l’abside les armes au bœuf passant de Calixte III Borgia († 1458), qui rappellent les restaurations effectuées sous son pontificat (fig. 153), comme les armes de Clément XII Corsini († 1740), aux trois bandes et à la fasce brochant sur le tout, qui sont aussi liées à une phase de restauration de l’église (fig. 141). Enfin, on remarquera les armes cardinalices puis papales d’Angelo Giuseppe Roncalli, saint Jean XXIII († 1963), qui fut cardinal titulaire de Sainte-Prisca de 1953 à 1958\teiNote[n={18},place={foot},xmlid={ftn36-5}]{\teiP[xmlid={p-438}]{Roncalli, d’humbles origines, n’avait pas d’armes familiales. Quand il accéda au siège titulaire d’Areopoli, avec le titre d’archevêque, puis à celui de Mesembria, puis au cardinalat, puis au pontificat, il porta un écu peut-être inspiré de celui d’autres familles du même nom ayant vécu dans la région bergamasque, attesté dans les fresques de l’église Saint-Augustin à Bergame, ou, pour certains, de celui des Maytini di Sotto il Monte. Roncalli, dans une de ses notes autobiographiques, raconte qu’à cet écu il voulut ajouter deux lys et le mot \teiHi[rend={italic}]{Oboedientia et pax}. Le pontife lui-même, historien de formation, donne plusieurs fois l’explication de ce \teiHi[rend={italic}]{motto} dans son \teiHi[rend={italic}]{Journal d’une âme} : « Motto de mes armes sont les mots \teiHi[rend={italic}]{Oboedientia et pax}, que le père Cesare Baronio prononçait chaque jour en baisant à Saint-Pierre le pied de l’apôtre. Ces mots sont un peu mon histoire et ma vie. Oh, qu’ils soient la glorification de mon pauvre nom dans les siècles » (cité par Mario Benigni et Goffredo Zanchi, \teiHi[rend={italic}]{Giovanni XXIII}, Cinisello Balsamo, Edizioni San Paolo, 2000, p. 174. Voir Angelo Roncalli, \teiHi[rend={italic}]{Il Cardinale Cesare Baronio}, Rome, Edizioni di storia e letteratura, 1961, p. 46 et 47 ; Camillo Fumagalli, « Dello stemma di Papa Giovanni XXIII. La genealogia dei Roncalli, Ronco e Roncaglia », \teiHi[rend={italic}]{Atti dell'Ateneo di Scienze, Lettere ed Arti di Bergamo}, n\teiHi[rend={sup}]{o} 34, 1968-1969, p. 11-22 ; Jacques Martin, \teiHi[rend={italic}]{Heraldry in the Vatican / L’araldica in Vaticano}, Gerrards Cross, Van Duren, 1987, p. 235-245 ; Pietro Messa, « Il motto di Giovanni XXIII: obbedienza e pace », \teiHi[rend={italic}]{Forma Sororum}, n\teiHi[rend={sup}]{o} 50, 2013, p. 250-253). L’écu peut se décrire ainsi : « De gueules à la fasce d’argent, à la tour au naturel sur le tout, accostée de deux lys d’argent, au chef de saint Marc (d’argent au lion au naturel passant ailé et nimbé, tenant de sa patte droite antérieure un livre ouvert où est inscrite la légende \teiHi[rend={small-caps italic}]{pax tibi marce evangelista meus}). » La présence dans l’écu de ce que l’on appelle le chef patriarcal de saint Marc rappelle, comme pour Pie X († 1914) et pour Jean-Paul I († 1978), la période où le cardinal Roncalli fut patriarche de Venise. Il s’agit d’une transposition de l’emblème de la Sérénissime République. Le lion est l’animal symbolique de l’évangéliste Marc qui lors de ses voyages aurait séjourné dans une île de la lagune où il aurait eu la vision d’un ange et entendu le salut : « \teiHi[rend={italic}]{Pax tibi Marce, evangelista meus} » \lbrack{}Paix à toi, Marc, mon évangéliste\rbrack{}. Dans les armes de la Sérénissime, le lion était d’or en champ de gueules, rappelant la pourpre impériale, et il tenait un sabre. Adopté par les patriarches, l’emblème devint religieux et non plus politique, le champ passa du gueules à l’argent et le lion au naturel. En outre, pour manifester le caractère pacifique de la religion chrétienne, le lion perdit son sabre. }} (fig. 154). Par ailleurs, les écus de l’ordre augustin ne manquent pas, la gestion du complexe de l’Aventin lui ayant été confiée en 1606 et ayant duré jusqu’à ce jour. Mais c’est là un autre chapitre que nous ne pouvons approfondir, comme celui non moins intéressant de l’héraldique sur les tissus et parements, souvent négligés par ceux mêmes qui en ont la tutelle. Ils sont pourtant d’importants témoins du destin des monuments sacrés pour lesquels l’héraldique est souvent l’unique clef de lecture et d’accès à une histoire complexe. Montrons la chape portant l’écu prélatice de Camillo Alfonso de Romanis (1885-1950), historien de l’ordre augustin, curé de Sainte-Prisca puis vicaire général de la Cité du Vatican et préfet du sanctuaire apostolique (1937-1950) (fig. 155). Les armes figurées sur ce vêtement liturgique dont il fit don à Sainte-Prisca sont : « \lbrack{}p\rbrack{}arti au 1 d’argent à l’emblème des augustins, au 2 d’azur au lion d’argent à la bande de gueules\teiNote[n={19},place={foot},xmlid={ftn37-7}]{\teiP[xmlid={p-439}]{Voir Davide Aurelio Perini, \teiHi[rend={italic}]{Bibliographia Augustiniana}, Florence, Tip. Sordomuti, 1929-1938, 4 vol., t. 2, p. 29-30 et t. 4, \teiHi[rend={italic}]{Additamenta}, p. 74-75 ; Rocco Ronzani, « Note di sigillografia dell’ordine di sant’Agostino », \teiHi[rend={italic}]{Subsidia Augustiniana Italica} II, 1, Rome, Centro culturale agostiniano, 2009, p. 22-25.}} ».}
              \begin{teiFigure}[xmlid={figure01-10}]

                \teiGraphic[url={../icono/br/Ch11\_Ronzani/Fig139.jpeg}]
                \teiHead[xmlid={figure-136-portrait-du-cardinal-benedetto-giustiniani-michel-natalis-dapres-un-dessin-de-giovanni-citosibio-guidi-gravure-238-mm-170-mm-rome-1636-1647-dans-galleria-giustiniana-del-marchese-vincenzo-giustiniani-rome-partie-2-pl-5-p-169-recueil-de-gravures-des-sculptures-antiques-de-la-collection-de-vincenzo-giustiniani-avec-une-galerie-de-portraits-de-famille-bnf-001}]{Figure 136. Portrait du cardinal Benedetto Giustiniani. Michel Natalis, d’après un dessin de Giovanni Citosibio Guidi, gravure (238 mm × 170 mm), Rome, 1636-1647, dans \teiHi[rend={italic}]{Galleria Giustiniana del Marchese Vincenzo Giustiniani}, Rome, partie 2, pl. 5, p. 169 ; recueil de gravures des sculptures antiques de la collection de Vincenzo Giustiniani, avec une galerie de portraits de famille, BNF.}
              
\end{teiFigure}
              \begin{teiFigure}[xmlid={figure02-10}]

                \teiGraphic[url={../icono/br/Ch11\_Ronzani/Fig140.jpg}]
                \teiHead[xmlid={figure-137-rome-eglise-sainte-prisca-maitre-autel-domenico-cresti-bapteme-de-sainte-prisca-001}]{Figure 137. Rome, église Sainte-Prisca, maître autel, Domenico Cresti,\teiHi[rend={italic}]{ Baptême de sainte Prisca}.}
              
\end{teiFigure}
              \begin{teiFigure}[xmlid={figure03-10}]

                \teiGraphic[url={../icono/br/Ch11\_Ronzani/Fig141.jpg}]
                \teiHead[xmlid={figure-138-rome-eglise-sainte-prisca-abside-anastasio-fontebuoni-et-collaborateurs-image-de-saint-pierre-repeint-en-saint-augustin-au-xix-e-siecle-001}]{Figure 138. Rome, église Sainte-Prisca, abside, Anastasio Fontebuoni et collaborateurs, image de saint Pierre, repeint en saint Augustin au \teiHi[rend={small-caps}]{xix}\teiHi[rend={sup}]{e} siècle.}
              
\end{teiFigure}
              \begin{teiFigure}[xmlid={figure04-10}]

                \teiGraphic[url={../icono/br/Ch11\_Ronzani/Fig142.jpg}]
                \teiHead[xmlid={figure-139-rome-eglise-sainte-prisca-abside-anastasio-fontebuoni-et-collaborateurs-ange-portant-la-ferule-pontificale-001}]{Figure 139. Rome, église Sainte-Prisca, abside, Anastasio Fontebuoni et collaborateurs, \teiHi[rend={italic}]{Ange portant la férule pontificale}.}
              
\end{teiFigure}
              \begin{teiFigure}[xmlid={figure05-10}]

                \teiGraphic[url={../icono/br/Ch11\_Ronzani/Fig143.jpg}]
                \teiHead[xmlid={figure-140-rome-eglise-sainte-prisca-abside-anastasio-fontebuoni-et-collaborateurs-ange-portant-les-clefs-de-saint-pierre-001}]{Figure 140. Rome, église Sainte-Prisca, abside, Anastasio Fontebuoni et collaborateurs, \teiHi[rend={italic}]{Ange portant les clefs de saint Pierre}.}
              
\end{teiFigure}
              \begin{teiFigure}[xmlid={figure06-10}]

                \teiGraphic[url={../icono/br/Ch11\_Ronzani/Fig144.jpg}]
                \teiHead[xmlid={figure-141-rome-eglise-sainte-prisca-fresques-de-la-nef-centrale-anastasio-fontebuoni-et-collaborateurs-saint-benoit-de-nursie-001}]{Figure 141. Rome, église Sainte-Prisca, fresques de la nef centrale, Anastasio Fontebuoni et collaborateurs, \teiHi[rend={italic}]{Saint Benoît de Nursie}.}
              
\end{teiFigure}
              \begin{teiFigure}[xmlid={figure07-9}]

                \teiGraphic[url={../icono/br/Ch11\_Ronzani/Fig145.jpg}]
                \teiHead[xmlid={figure-142-rome-eglise-sainte-prisca-fresques-de-la-nef-centrale-anastasio-fontebuoni-et-collaborateurs-saint-jean-gualbert-001}]{Figure 142. Rome, église Sainte-Prisca, fresques de la nef centrale, Anastasio Fontebuoni et collaborateurs, \teiHi[rend={italic}]{Saint Jean Gualbert}.}
              
\end{teiFigure}
              \begin{teiFigure}[xmlid={figure08-7}]

                \teiGraphic[url={../icono/br/Ch11\_Ronzani/Fig146.jpg}]
                \teiHead[xmlid={figure-143-armes-de-la-congregation-de-vallombreuse-miniature-bibliotheque-nationale-de-florence-conventi-soppressi-c-v-1912-fol-2-8-mai-1513-001}]{Figure 143. Armes de la congrégation de Vallombreuse, miniature, Bibliothèque nationale de Florence, Conventi soppressi, C.V.1912, fol. 2 (8 mai 1513).}
              
\end{teiFigure}
              \begin{teiFigure}[xmlid={figure09-7}]

                \teiGraphic[url={../icono/br/Ch11\_Ronzani/Fig147.jpg}]
                \teiHead[xmlid={figure-144-rome-eglise-sainte-prisca-fresques-de-la-crypte-anastasio-fontebuoni-et-collaborateurs-armes-de-sixte-quint-peretti-et-au-dessous-du-cardinal-giustiniani-001}]{Figure 144. Rome, église Sainte-Prisca, fresques de la crypte, Anastasio Fontebuoni et collaborateurs, armes de Sixte Quint Peretti et, au-dessous, du cardinal Giustiniani.}
              
\end{teiFigure}
              \begin{teiFigure}[xmlid={figure10-6}]

                \teiGraphic[url={../icono/br/Ch11\_Ronzani/Fig148.jpg}]
                \teiHead[xmlid={figure-145-rome-eglise-sainte-prisca-fresques-de-la-crypte-anastasio-fontebuoni-et-collaborateurs-lion-heraldique-peretti-couronne-par-un-ange-001}]{Figure 145. Rome, église Sainte-Prisca, fresques de la crypte, Anastasio Fontebuoni et collaborateurs, lion héraldique Peretti couronné par un ange.}
              
\end{teiFigure}
              \begin{teiFigure}[xmlid={figure11-6}]

                \teiGraphic[url={../icono/br/Ch11\_Ronzani/Fig149.jpg}]
                \teiHead[xmlid={figure-146-rome-eglise-sainte-prisca-decor-de-marbre-du-maitre-autel-armes-du-cardinal-giustiniani-001}]{Figure 146. Rome, église Sainte-Prisca, décor de marbre du maître autel, armes du cardinal Giustiniani.}
              
\end{teiFigure}
              \begin{teiFigure}[xmlid={figure12-6}]

                \teiGraphic[url={../icono/br/Ch11\_Ronzani/Fig150.jpg}]
                \teiHead[xmlid={figure-147-rome-eglise-sainte-prisca-inscription-commemorative-encadree-des-armes-du-cardinal-giustiniani-et-de-leurs-meubles-aigle-eployee-et-chateau-001}]{Figure 147. Rome, église Sainte-Prisca, inscription commémorative encadrée des armes du cardinal Giustiniani et de leurs meubles (aigle éployée et château).}
              
\end{teiFigure}
              \begin{teiFigure}[xmlid={figure13-5}]

                \teiGraphic[url={../icono/br/Ch11\_Ronzani/Fig151.jpg}]
                \teiHead[xmlid={figure-148-rome-eglise-sainte-prisca-fresques-de-la-crypte-anastasio-fontebuoni-et-collaborateurs-armes-giustiniani-lhumidite-a-fait-perdre-laigle-qui-avait-peut-etre-ete-peinte-a-tempera-001}]{Figure 148. Rome, église Sainte-Prisca, fresques de la crypte, Anastasio Fontebuoni et collaborateurs, armes Giustiniani (l’humidité a fait perdre l’aigle qui avait peut-être été peinte \teiHi[rend={italic}]{a tempera}).}
              
\end{teiFigure}
              \begin{teiFigure}[xmlid={figure14-5}]

                \teiGraphic[url={../icono/br/Ch11\_Ronzani/Fig152.jpg}]
                \teiHead[xmlid={figure-149-rome-eglise-sainte-prisca-frise-de-lescalier-dacces-a-la-crypte-detail-heraldique-giustiniani-chateau-001}]{Figure 149. Rome, église Sainte-Prisca, frise de l’escalier d’accès à la crypte, détail héraldique Giustiniani (château).}
              
\end{teiFigure}
              \begin{teiFigure}[xmlid={figure15-5}]

                \teiGraphic[url={../icono/br/Ch11\_Ronzani/Fig153.jpg}]
                \teiHead[xmlid={figure-150-rome-eglise-sainte-prisca-frise-de-lescalier-dacces-a-la-crypte-detail-heraldique-giustiniani-aigle-001}]{Figure 150. Rome, église Sainte-Prisca, frise de l’escalier d’accès à la crypte, détail héraldique Giustiniani (aigle).}
              
\end{teiFigure}
              \begin{teiFigure}[xmlid={figure16-5}]

                \teiGraphic[url={../icono/br/Ch11\_Ronzani/Fig154.jpg}]
                \teiHead[xmlid={figure-151-rome-eglise-sainte-prisca-crypte-meuble-des-armes-giustiniani-aigle-sur-le-benitier-001}]{Figure 151. Rome, église Sainte-Prisca, crypte, meuble des armes Giustiniani (aigle) sur le bénitier.}
              
\end{teiFigure}
              \begin{teiFigure}[xmlid={figure17-4}]

                \teiGraphic[url={../icono/br/Ch11\_Ronzani/Fig155.jpg}]
                \teiHead[xmlid={figure-152-rome-eglise-sainte-prisca-crypte-armes-de-clement-viii-aldobrandini-001}]{Figure 152. Rome, église Sainte-Prisca, crypte, armes de Clément VIII Aldobrandini.}
              
\end{teiFigure}
              \begin{teiFigure}[xmlid={figure18-4}]

                \teiGraphic[url={../icono/br/Ch11\_Ronzani/Fig156.jpg}]
                \teiHead[xmlid={figure-153-rome-eglise-sainte-prisca-abside-inscription-surmontee-des-armes-de-callixte-iii-borgia-001}]{Figure 153. Rome, église Sainte-Prisca, abside, inscription surmontée des armes de Callixte III Borgia.}
              
\end{teiFigure}
              \begin{teiFigure}[xmlid={figure19-3}]

                \teiGraphic[url={../icono/br/Ch11\_Ronzani/Fig157.jpg}]
                \teiHead[xmlid={figure-154-rome-eglise-sainte-prisca-presbytere-inscription-commemorative-avec-les-armes-cardinalices-et-papales-dangelo-giuseppe-roncalli-saint-jean-xxiii-001}]{Figure 154. Rome, église Sainte-Prisca, presbytère, inscription commémorative avec les armes cardinalices et papales d’Angelo Giuseppe Roncalli, saint Jean XXIII.}
              
\end{teiFigure}
              \begin{teiFigure}[xmlid={figure20-3}]

                \teiGraphic[url={../icono/br/Ch11\_Ronzani/Fig158.jpg}]
                \teiHead[xmlid={figure-155-rome-eglise-sainte-prisca-sacristie-chape-portant-lecu-prelatice-de-camillo-alfonso-de-romanis-1885-1950-historien-de-lordre-augustin-cure-de-sainte-prisca-puis-vicaire-general-de-la-cite-du-vatican-et-prefet-du-sacraire-apostolique-1937-1950-001}]{Figure 155. Rome, église Sainte-Prisca, sacristie, chape portant l’écu prélatice de Camillo Alfonso de Romanis (1885-1950), historien de l’ordre augustin, curé de Sainte-Prisca puis vicaire général de la Cité du Vatican et préfet du sacraire apostolique (1937-1950).}
              
\end{teiFigure}
            
\end{teiDiv}
          
\end{teiElement}
        
\end{teiElement}
        \begin{teiElement}[name={group},type={chapter},data-page-title={Héraldique et représentation à Saint-Pierre de Rome},data-page-subtitle={Entre le durable et l’éphémère},data-page-authors={Martine Boiteux},data-include-href={Ch12\_Heraldique\_et\_representation.xml},xmlid={chapter-009}]

          \begin{teiElement}[name={front}]

            \begin{teiDiv}[type={titlePage},xmlid={titlepage-014}]

              \teiP[rend={title-main},xmlid={p-440}]{Héraldique et représentation à Saint-Pierre de Rome}
              \teiP[rend={title-sub},xmlid={p-441}]{Entre le durable et l’éphémère}
              \teiP[rend={author-aut},xmlid={p-442}]{Martine Boiteux}
            
\end{teiDiv}
          
\end{teiElement}
          \begin{teiElement}[name={body},xmlid={body-014}]

            \teiP[xmlid={p1-14}]{À partir d’un corpus visuel d’images commémoratives et mémorielles, sont présentés les blasons et emblèmes sur les monuments de la basilique Saint-Pierre, lieu de mémoire, ou leur absence, ainsi que sur les décors éphémères des cérémonies, et leurs représentations sur des gravures, en éclairant leurs significations par les textes écrits. Les monuments et apparats éphémères sont représentés sur des images à double signification : à la fois projet d’architecte et objet mémoriel pour la fabrication et la diffusion de l’événement éphémère, communication et propagande au service d’une stratégie politique. Les images sont de deux types : iconiques et narratives ; faites avant l’événement et distribuées, elles portent l’interprétation officielle.}
            \teiP[xmlid={p2-13}]{Les emblèmes sont utilisés comme instrument de communication ; l’héraldique sert à identifier le commanditaire du monument durable ou éphémère. Indispensables sur les monuments et sur les architectures éphémères des fêtes nombreuses à l’époque baroque\teiNote[n={1},place={foot},xmlid={ftn1-12}]{\teiP[xmlid={p-443}]{Pierre Francastel, « La Contre-Réforme et les arts en Italie à la fin du \teiHi[rend={small-caps}]{xvi}\teiHi[rend={sup}]{e }siècle », dans Henri Bédarida (dir.), \teiHi[rend={italic}]{À travers l’art italien du xv}\teiHi[rend={italic sup}]{e }\teiHi[rend={italic}]{au }\teiHi[rend={small-caps italic}]{xx}\teiHi[rend={italic sup}]{e}\teiHi[rend={italic}]{ siècle}, Paris, Boivin-Société d’études italiennes, 1949, p. 63-113.}}, ils sont à la fois ornements et signes identitaires\teiNote[n={2},place={foot},xmlid={ftn29-2}]{\teiP[xmlid={p-444}]{Pierre Bourdieu, \teiHi[rend={italic}]{La Distinction : critique sociale du jugement}, Paris, Éditions de Minuit, 1979.}}. Le contexte culturel pousse à la théâtralisation dans une Rome \teiHi[rend={italic}]{caput mundi }et théâtre du monde, où tout est représentation comme l’écrivait l’ambassadeur Wicquefort\teiNote[n={3},place={foot},xmlid={ftn30-2}]{\teiP[xmlid={p-445}]{Abraham de Wicquefort, \teiHi[rend={italic}]{Mémoire touchant les ambassadeurs et les ministres publics}, La Haye, Jean et Daniel Steucker, 1677, p. 229 ; \teiHi[rend={italic}]{L’ambassadeur et ses fonctions}, La Haye, Maurice Georges Veneur, 1682, 2 vol.}}\teiHi[rend={italic}]{.} Le pape, souverain absolu au double pouvoir, temporel et spirituel sur Rome et la Chrétienté\teiNote[n={4},place={foot},xmlid={ftn31-2}]{\teiP[xmlid={p-446}]{Paolo Prodi, \teiHi[rend={italic}]{Il sovrano pontefice : un corpo e due anime : la monarchia papale nella prima età moderna}, Bologne, Il Mulino, 1982.}}, surveille la magnificence de la basilique et des rituels qui s’y déroulent\teiNote[n={5},place={foot},xmlid={ftn32-2}]{\teiP[xmlid={p-447}]{Martine Boiteux, « Il cerimoniale. La necessità della magnificenza », dans Francesco Buranelli (dir.), \teiHi[rend={italic}]{Il Vaticano Barocco, Arte, architettura e cerimoniale,} Milan, Jaca Book, 2014, p. 10-27.}}. Antonio Pinelli montre la profusion des emblèmes\teiNote[n={6},place={foot},xmlid={ftn33-3}]{\teiP[xmlid={p-448}]{Antonio Pinelli, « \teiHi[rend={italic}]{La philosophie des images}. Emblemi e imprese fra manierismo e barocco », \teiHi[rend={italic}]{Ricerche di storia dell’arte}, n\teiHi[rend={sup}]{o} 1-2, 1976, p. 3-28 ; Carlo Ginzburg, \teiHi[rend={italic}]{Miti, emblemi}, \teiHi[rend={italic}]{spie. Morfologia e storia}, Turin, Einaudi, 1986.}} et rappelle les nombreux traités publiés à partir du \teiHi[rend={small-caps}]{xvi}\teiHi[rend={sup}]{e }siècle\teiNote[n={7},place={foot},xmlid={ftn34-4}]{\teiP[xmlid={p-449}]{Paolo Giovio, \teiHi[rend={italic}]{Dialogo dell’imprese militari et amorose}, In Roma, Appresso Antonio Barre, 1555.}} montrant l’utilisation de l’\teiHi[rend={italic}]{impresa} comme figure de rhétorique\teiNote[n={8},place={foot},xmlid={ftn35-4}]{\teiP[xmlid={p-450}]{Giulio Carlo Argan, « La rettorica et l’arte barocca », dans Enrico Castelli (dir.), \teiHi[rend={italic}]{Retorica e Barocco. Atti del 3 congresso internazionale di studi umanistici}, Rome, Fratelli Bocca, 1955, p. 10-14, republié dans \teiHi[rend={italic}]{Studi e Note. Da Bramante a Canova}, Rome, Bulzoni, 1970.}}. Robert Klein\teiNote[n={9},place={foot},xmlid={ftn36-6}]{\teiP[xmlid={p-451}]{Robert Klein, « La théorie de l’expression figurée dans les traités italiens sur les \teiHi[rend={italic}]{Imprese}, 1555-1612 », \teiHi[rend={italic}]{Bibliothèque d’humanisme et Renaissance}, vol. 19, n\teiHi[rend={sup}]{o} 2, n\teiHi[rend={sup}]{o} 2, 1957, p. 320-341 ; \teiHi[rend={italic}]{La forme et l’intelligible}. \teiHi[rend={italic}]{Écrits sur la Renaissance et l’art moderne}, Paris, Gallimard, 1970, p. 125-150.}} souligne sa valeur de distinction sociale ainsi que sa fonction de talisman voulu et consciemment utilisé à l’époque de la Contre-Réforme catholique. Héraldique, pouvoir et efficacité symbolique\teiNote[n={10},place={foot},xmlid={ftn37-8}]{\teiP[xmlid={p-452}]{Claude Lévi-Strauss, « L’efficacité symbolique », \teiHi[rend={italic}]{Revue d’histoire des religions}, vol. 135, n\teiHi[rend={sup}]{o} 1, 1949, p. 5-27, repris dans \teiHi[rend={italic}]{Anthropologie structurale}, Paris, Plon, 1958, chap. X ; Martine Boiteux, « Linguaggio figurativo ed efficacia rituale », dans Francesca Cantù (dir.), \teiHi[rend={italic}]{I linguaggi del potere. Politica e religione nell’età barocca}, Rome, Viella, 2007, p. 39-79.}} sont liés et interagissent.}
            \teiP[xmlid={p3-13}]{Sont ici proposés quelques exemples d’utilisation, de typologie et d’emplacement des armoiries, blasons, devises et emblèmes dans la basilique Saint-Pierre sur des supports variés. Comment s’expriment les décorations et la propagande ? La présence ou l’absence des emblèmes du pape sur les monuments et sur les apparats éphémères ? Quels sont les signes identitaires et mémoriels ?}
            \begin{teiDiv}[type={section1},xmlid={div01-11}]

              \teiHead[xmlid={la-basilique-un-monument-marque-par-ses-constructeurs-et-ses-utilisateurs-001}]{La basilique : un monument marqué par ses constructeurs et ses utilisateurs}
              \teiP[xmlid={p4-13}]{Un seul exemple pour le monument est présenté ici : le baldaquin Barberini (fig. 156). De nombreuses études lui ayant été consacrées, l’analyse est donc limitée à quelques remarques en relation avec l’héraldique. Il est construit de 1624 à 1633 par le jeune Gian Lorenzo Bernini, avec la participation de son rival Borromini. Ce baldaquin est conçu comme une machine processionnelle et des expérimentations de forme sont faites : ainsi, les dessins et gravures réalisés en 1630 pour le catafalque du frère du pape, Carlo Barberini. Il est bien connu que le marbre et le bronze sont pris au Panthéon pour sa construction. Les abeilles sont déjà présentes sur les architectures éphémères. C’est l’image de l’Église de la Contre-Réforme avec la présence dominante des emblèmes d’Urbain VIII Barberini : abeilles, soleil, laurier et emblèmes du pontificat fondus avec les symboles de la gloire de Pierre et de la religion catholique romaine\teiNote[n={11},place={foot},xmlid={ftn38-7}]{\teiP[xmlid={p-453}]{Maria Grazia D’Amelio, \teiHi[rend={italic}]{L’obelisco marmoreo del foro italico a Roma. Storia, immagini e note tecniche}, Rome, Palombi \& Partner, 2009 ; \teiHi[rend={italic}]{Giovan Lorenzo Bernini e l’oro per il baldacchino di San Pietro (1624-1633)}, Rome, Argos, 2021.}}. Sur les bases au milieu du blason sont sculptés des bas-reliefs avec des visage douloureux mystérieux. Les colonnes salomoniques sont couvertes de laurier grimpant. La lumière est un élément constitutif de l’architecture du monument. L’intention du monument est la célébration : le pape Barberini, commanditaire, très présent dans la réalisation du baldaquin, veut démontrer la légitimité de la succession apostolique et le primat du pape qui en procède\teiNote[n={12},place={foot},xmlid={ftn39-7}]{\teiP[xmlid={p-454}]{Maria Grazia D’Amelio, \teiHi[rend={italic}]{Giovan Lorenzo Bernini…},\teiHi[rend={italic}]{ op. cit}.}} : sur la coupole, la présence de \teiHi[rend={italic}]{putti} avec les emblèmes du pouvoir, tiare et clefs, en témoignent. L’importance du symbole de la remise des clefs est soulignée sur une peinture de Perugino (1581-1582) dans la chapelle Sixtine ; et également par le relief sur la façade de la basilique sculpté par Ambrogio Buonvicino dans les années 1614-1632 (fig. 157).}
            
\end{teiDiv}
            \begin{teiDiv}[type={section1},xmlid={div02-11}]

              \teiHead[xmlid={les-personnes-001}]{Les personnes}
              \begin{teiDiv}[type={section2},xmlid={div03-9}]

                \teiHead[xmlid={les-tombeaux-des-papes-001}]{Les tombeaux des papes}
                \teiP[xmlid={p5-13}]{Les tombeaux de certains papes se trouvent dans les cryptes et leurs monuments funéraires dans la basilique. Ils se composent toujours de la statue, rarement du blason et souvent de sculptures allégoriques. C’est le cas à plusieurs reprises : en 1492, pour Innocent VIII Cibo par Antonio Pollaiolo, le monument le plus ancien avec les clefs sur la base (fig. 158) ; en 1545, pour Paul III Farnèse par Guglielmo della Porta, sans blason, avec des statues allégoriques dont deux se trouvent au palais Farnèse par manque de place dans Saint-Pierre ; en 1585, pour Grégoire XIII Boncompagni par Camillo Rusconi, avec des statues allégoriques, inscription et le dragon, animal du blason du pape, sous l’urne ; pour Urbain VIII Barberini, terminé après sa mort en 1644, sans blason ni abeille ; en 1667, pour Alexandre VII Chigi par Le Bernin, avec statues ; pour Benoît XIV par Pietro Bracci, 1764-1769 ; en 1769, pour Clément XIII Rezzonico par Antonio Canova, avec la tiare et les lions qui gardent la porte.}
              
\end{teiDiv}
              \begin{teiDiv}[type={section2},xmlid={div04-6}]

                \teiHead[xmlid={les-tombeaux-de-personnalites-001}]{Les tombeaux de personnalités }
                \teiP[xmlid={p6-13}]{Trois femmes seulement sont présentes dans la basilique Saint-Pierre.}
                \teiP[xmlid={p7-13}]{Mathilde de Canossa ou de Toscane (1046-1115) est une souveraine forte et un soutien fidèle de l’Église ; en 1632, Urbain VIII fait transférer sa dépouille au château Saint Ange, puis, en 1644, à Saint-Pierre, où un monument funèbre est réalisé par Le Bernin (fig. 159).}
                \teiP[xmlid={p8-13}]{Christine de Suède, la grande convertie, arrivée à Rome en 1655 après son abdication, est accueillie avec les honneurs dus à son rang de reine de sang royal ; elle mène une vie mondaine, religieuse et intellectuelle active. Elle est célébrée avec faste lors de sa guérison en 1689 à San Salvatore in Lauro. Lors de ses grandioses funérailles d’État en 1689 à la Chiesa Nuova, église des Oratoriens, voulues par le pape malgré ses volontés testamentaires\teiNote[n={13},place={foot},xmlid={ftn40-7}]{\teiP[xmlid={p-455}]{Martine Boiteux, « Funérailles féminines dans la Rome baroque », dans Bernard Dompnier (dir.), \teiHi[rend={italic}]{Les cérémonies extraordinaires du catholicisme baroque}, Clermont-Ferrand, PUBP, 2009, p. 389-421.}}, les signes emblématiques de sa famille et de sa dignité royale sont présents : sceptre, couronne et manteau royal posés sur son corps exposé et dominé par le dais sacral. La gravure du transport à Saint-Pierre, où se déroulent les secondes funérailles, porte le blason des Wasa, famille royale suédoise\teiNote[n={14},place={foot},xmlid={ftn41-7}]{\teiP[xmlid={p-456}]{Visible sur : \teiRef[target={https://collections.vam.ac.uk/item/O649881/pompae-funebres-habitae-in-funere-print-robert-van-auden/}]{https://collections.vam.ac.uk/item/O649881/pompae-funebres-habitae-in-funere-print-robert-van-auden/} (consulté le 03 mars 2025).}}. C’est une reine qui est célébrée dans la basilique par une célébration éphémère et par un tombeau. Construit sur un projet de l’architecte Carlo Fontana, ce monument intègre son portrait modelé par Giovanni Giardini sur un médaillon en bronze dominant la couronne royale. Christine est identifiée par l’inscription.}
                \teiP[xmlid={p9-12}]{Maria Clementina Sobieska, morte en 1735, a des funérailles d’État sur le modèle de celles de Christine de Suède, parce qu’elle aussi est de sang royal et mariée au prétendant catholique au trône d’Angleterre, Jacques III Stuart\teiNote[n={15},place={foot},xmlid={ftn42-7}]{\teiP[xmlid={p-457}]{Martine Boiteux, « Triomphes romains des Sobieski, 1683-1735 », dans \teiHi[rend={italic}]{I Sobieski a Roma. La famiglia reale polacca nella Città Eterna}, Varsovie, Muzeum Pałacu Króla Jana III w Wiłanwie, 2018, p. 48-51.}}. Cette célébration a lieu dans sa paroisse des Santi Apostoli puis dans la basilique Saint-Pierre. Son tombeau de style encore baroque comporte des sculptures de Pietro Bracci exécutées sur dessin de Filippo Barigioni (1739-1742). Son portrait, en mosaïque de Pietro Paolo Cristofari à partir d’une peinture d’Ignaz Stern, fait scandale dans Saint-Pierre car trop mondain et pas assez religieux ; il est retiré, mais le pape le fait remettre en place. Son mari, Jacques Stuart, toujours prétendant catholique au trône britannique, meurt en 1766 et a droit aux mêmes honneurs funèbres. Un tombeau par Canova en 1819 l’accueille avec leurs deux fils et est installé en face de celui de sa femme ; de style néoclassique, il évoque les tombeaux étrusques, la porte est gardée par deux anges ; il est surmonté du blason familial et de la couronne royale.}
                \teiP[xmlid={p10-12}]{Les cérémonies de funérailles pour des personnalités dans la basilique Saint-Pierre ne suscitent pas d’images mais seulement la réalisation de tombeaux, monuments durables sans blason pour les trois femmes.}
              
\end{teiDiv}
            
\end{teiDiv}
            \begin{teiDiv}[type={section1},xmlid={div05-3}]

              \teiHead[xmlid={la-basilique-lieu-du-ceremonial-ceremonies-ephemeres-et-representations-001}]{La basilique lieu du cérémonial : cérémonies éphémères et représentations}
              \begin{teiDiv}[type={section2},xmlid={div06-3}]

                \teiHead[xmlid={les-funerailles-001}]{Les funérailles}
                \teiP[xmlid={p11-12}]{Le cortège funèbre du pape se déroule du palais du Quirinal, où il réside, à la basilique Saint-Pierre ; il est fait de manière anonyme et sans blason ni emblèmes, seulement une croix et le corps exposé : il n’est plus pape, sans pouvoir, ni individu vivant. Après l’exposition au Quirinal, il est exposé à Saint-Pierre, toujours sans emblème.}
                \teiP[xmlid={p12-12}]{En temps de \teiHi[rend={italic}]{Sede vacante}, vacance du Siège Apostolique, les plans du conclave apparaissent à la fin \teiHi[rend={small-caps}]{xvi}\teiHi[rend={sup}]{e }siècle, après le Concile de Trente et dans le contexte de réaffirmation du pouvoir pontifical. La vie est régulée par des cérémoniels écrits\teiNote[n={16},place={foot},xmlid={ftn43-8}]{\teiP[xmlid={p-458}]{Marc Dykmans, \teiHi[rend={italic}]{Le cérémonial pontifical de la fin du Moyen Âge à la Renaissance}, Rome-Bruxelles, Institut historique belge de Rome, 1977-1985, 4 vol. }} et figurés\teiNote[n={17},place={foot},xmlid={ftn44-8}]{\teiP[xmlid={p-459}]{Martine Boiteux, « Rituels funéraires pontificaux », \teiHi[rend={italic}]{Studi romani}, vol. 60, n\teiHi[rend={sup}]{o} 1-4, 2012, p. 31-56.}} ; ces images sérielles de natures variées, normatives et mémorielles, sont toujours au nom du pape défunt, excepté en 1585 pour Sixte Quint à peine élu : la norme n’est pas encore fixée. L’image iconique souligne les lieux avec la figuration du Saint-Esprit et de l’emblème pontifical, exprimant la sacralisation du pouvoir (fig. 160) ; l’image narrative souligne les actions rituelles spécifiques à cette période, avec parfois le portrait du pape défunt ou avec des vignettes, comme sur une bande dessinée, depuis la mort jusqu’à l’adoration du nouvel élu. Ces plans sont le plus souvent sans emblème personnalisé, avec seulement ceux du pontife (tiare et clefs croisées), quelquefois avec le portrait du défunt. Le blason et le portrait du cardinal dédicataire peuvent être représentés, comme en 1691 pour le cardinal de Médicis, au décès d’Alexandre VIII Ottoboni (fig. 161). Ces plans sont souvent réutilisés et actualisés en changeant les inscriptions et les noms. Les images sont identifiées par la présence de la colombe du Saint-Esprit en position dominante et l’emblème pontifical, c’est-à-dire les clefs entrecroisées et la tiare ou triple couronne.}
                \teiP[xmlid={p13-12}]{Les cérémonies funèbres sont somptueuses, de même que les catafalques des papes présents sur des gravures autonomes et sur les plans des conclaves, avec la présence d’emblèmes. Ils restent dans Saint-Pierre le temps des \teiHi[rend={italic}]{Novendiales}, neuf jours de deuil jusqu’à l’entrée en conclave des cardinaux pour l’élection du nouveau pape.}
                \teiP[xmlid={p14-9}]{Les premiers catafalques pontificaux apparaissent un an après le décès à la basilique Sainte-Marie-Majeure, avec les emblèmes familiaux sur les apparats éphémères et le blason des Peretti sur la coupole pour Sixte Quint en 1591, puis pour Paul V Borghèse en 1622, où l’architecture est de Venturini, les sculptures du Bernin et la gravure de Greuter, avec des blasons sur les médaillons du décor de la nef.}
                \teiP[xmlid={p15-8}]{Deux modèles existent : l’un en forme de tholos ou de tiare, l’autre incluant l’urne entourée d’obélisques. Tous les catafalques portent la tiare sans les clefs (l’homme est alors rendu à son statut individuel et familial) et le blason familial, sur le piédestal comme sur le baldaquin de Saint-Pierre pour les abeilles Barberini\teiNote[n={18},place={foot},xmlid={ftn45-8}]{\teiP[xmlid={p-460}]{Martine Boiteux, « Les Barberini, Rome et la France : fête et politique », dans Lorenza Mochi Onori, Sebastian Schutze et Francesco Solinas (dir.), \teiHi[rend={italic}]{I Barberini e la cultura europea del Seicento}, Rome, De Luca, 2007, p. 345-360.}}. Le modèle avec les obélisques l’emportera, magnifié par le dessin du Bernin pour Alexandre VII Chigi en 1667 où il opère une métaphore conceptuelle\teiNote[n={19},place={foot},xmlid={ftn46-8}]{\teiP[xmlid={p-461}]{Antonio Pinelli, art. cité.}} : l’emblématique intégrée est fondamentale dans le projet\teiHi[rend={italic}]{.}}
                \teiP[xmlid={p16-8}]{Les catafalques se ressemblent parce que le modèle architectural est souvent repris avec peu de variantes, sinon le renouvellement des emblèmes et des médaillons narratifs.}
              
\end{teiDiv}
              \begin{teiDiv}[type={section2},xmlid={div07-3}]

                \teiHead[xmlid={les-rituels-dinvestiture-001}]{Les rituels d’investiture}
                \teiP[xmlid={p17-8}]{Le nouvel élu choisit son blason ; le plus souvent le blason familial est alors surmonté de la tiare, triple couronne signe du pouvoir, et des clefs de Saint-Pierre, signe de continuité.}
                \teiP[xmlid={p18-8}]{Le couronnement a lieu à l’extérieur de la basilique et suscite rarement une représentation, sauf la fresque de Giovanni Guerra et Cesare Nebbia du \teiHi[rend={italic}]{Salone Sistino} à la Bibliothèque vaticane pour Sixte Quint, en 1585. Pour le couronnement d’Innocent X en octobre 1644 le cardinal Francesco Barberini, toujours vice-chancelier et ex-cardinal neveu d’Urbain VIII, pour la décoration de la porte de la basilique Saint-Pierre, prête deux draps de soie brodée représentant deux épisodes de la vie de Déborah, dont le nom en hébreu signifie abeille ; l’allusion au pape-poète précédent, Urbain VIII Barberini, est transparente. Par ce rappel du pontificat précédent et de sa propre famille, Francesco Barberini utilise l’emblème familial symboliquement et de manière évidente pour les participants cultivés de la cérémonie\teiNote[n={20},place={foot},xmlid={ftn47-8}]{\teiP[xmlid={p-462}]{Pascal-François Bertrand, « L’uso dell’arazzo a Roma: l’esempio del Cardinale Francesco Barberini », dans Lorenza Mochi Onori, Sebastian Schutze et Francesco Solinas (dir.), \teiHi[rend={italic}]{op. cit.}, p. 421-430, ici p. 429.}}.}
                \teiP[xmlid={p19-8}]{Le \teiHi[rend={italic}]{Possesso} transporte hors du Vatican mais il fait partie du cérémonial pontifical\teiNote[n={21},place={foot},xmlid={ftn48-9}]{\teiP[xmlid={p-463}]{Francesco Cancellieri, \teiHi[rend={italic}]{Storia de'solenni possessi ,} Rome, Luigi Lazzarini stampatore della R. C. A., 1802 ; Gaetano Moroni, \teiHi[rend={italic}]{Dizionario di erudizione storico-ecclesiastica}, In Venezia, dalla Tipografia Emiliana, 1840-1861, 103 vol. ; Martine Boiteux, « \teiHi[rend={italic}]{Il Possesso}. La presa di potere del Sovrano Pontefice sulla città di Roma », dans Francesco Buranelli (dir.),\teiHi[rend={italic}]{ Habemus Papam. Le elezioni pontificie da San Pietro a Benedetto XIV}, Rome, De Luca, 2006, p. 131-140.}}. Il marque la prise de possession de la ville et de l’évêché de Rome par une grande procession qui traverse l’espace urbain d’ouest en est, de Saint-Pierre à Saint-Jean-de-Latran, avec une étape au Capitole. Au \teiHi[rend={small-caps}]{xviii}\teiHi[rend={sup}]{e} siècle, il part parfois du Quirinal. Pour tous les papes, sont publiées des gravures sérielles et normatives répétitives identifiées par le blason et souvent le portrait du nouvel élu, excepté pour Sixte Quint dont le \teiHi[rend={italic}]{Possesso} est représenté sur une peinture comme le couronnement. Ainsi, en 1605, Léon XI Médicis est figuré avec son blason et son portrait en compagnie des trois autres papes Médicis qui l’ont précédé (fig. 162). Les arcs de triomphe éphémères qui rythment le parcours sont actualisés et identifiés par les emblèmes, blasons et inscriptions : offerts par la commune et situés en haut de l’escalier du Capitole, ils font voisiner les blasons du pape et de la municipalité, comme le montre en 1605 la gravure de Tempesta. Offerts par le duc de Parm\teiHi[rend={strikethrough}]{e}, un Farnèse, puis son héritier le roi de Naples, un Bourbon, les emblèmes identifient les arcs élevés dans le forum, notamment les lys communs aux familles ; souvent presque identiques avec comme seules variantes les inscriptions, portraits ou emblèmes en position sommitale.}
              
\end{teiDiv}
              \begin{teiDiv}[type={section2},xmlid={div08-3}]

                \teiHead[xmlid={lexercice-du-pouvoir-quelques-exemples-dutilisation-des-emblemes-001}]{L’exercice du pouvoir : quelques exemples d’utilisation des emblèmes}
                \teiP[xmlid={p20-8}]{Lors de la Semaine sainte, pour la cérémonie de la \teiHi[rend={italic}]{Croce Luminosa}\teiNote[n={22},place={foot},xmlid={ftn49-9}]{\teiP[xmlid={p-464}]{Martine Boiteux, « La Croce luminosa, un rito della Settimana Santa nella basilica di San Pietro », dans Roberto Regoli et Ilaria Fiumi Sermattei (dir\teiHi[rend={italic}]{.), La religione dei nuovi tempi. Il riformismo spirituale nell’età di Leone XII}, Ancone, Consiglio regionale delle Marche, 2020, p. 281-316.}}, on constate une absence de blason et de tout signe identitaire du commanditaire de la cérémonie, comme le montre un tableau de Desprez représentant la participation à la cérémonie du roi Gustave III de Suède. C’est un rituel religieux et liturgique du pouvoir pontifical, il n’est donc pas personnalisé.}
                \teiP[xmlid={p21-8}]{Les béatifications donnent lieu à de somptueuses cérémonies à partir de celle de François de Sales en 1662 ; le cérémonial est identique à la canonisation qui reprend le même apparat et les mêmes signes et symboles. Lors des canonisations l’emblèmes de chaque saint est figuré sur les étendards comme signe identitaire\teiNote[n={23},place={foot},xmlid={ftn50-9}]{\teiP[xmlid={p-465}]{Martine Boiteux, « La cerimonia della canonizzazione. Teatro riti, stendardi e immagini », dans Alessandra Bartolomei Romagnoli (dir.), \teiHi[rend={italic}]{La canonizzazione di Santa Francesca Romana, Santità, cultura e istituzioni a Roma tra Medioevo e età moderna}, Florence, Edizioni del Galluzzo per la Fondazione Ezio Franceschini, 2013, p. 99-121.}}. La façade de la basilique est décorée avec les images des saints et de leurs emblèmes et aussi avec les blasons des papes. En 1610, pour la canonisation de Carlo Borromeo, la façade est totalement recouverte d’un apparat éphémère ; la gloire de Milan éclate opposée à celle de Rome, manifeste lors la récente canonisation en 1608 de Francesca Romana patronne de la ville. Sur la gravure sont mis en valeur le blason du dédicataire, le cardinal Federico Borromeo neveu du canonisé, ainsi que de nombreux portraits des archevêques et évêques de Milan, prédécesseurs du canonisé, soulignant l’absence du blason du pape.}
                \teiP[xmlid={p22-8}]{Le théâtre où se déroule la célébration est représenté sur les gravures de canonisation. Ainsi, en 1622, Greuter figure le blason du pape célébrant Grégoire XV Ludovisi et les emblèmes identitaires des cinq nouveaux saints sur les vignettes du pourtour ainsi que sur les étendards suspendus dans la basilique comme signes identitaires (fig. 163) : pour Ignace et François Xavier est présent l’emblème jésuite ; pour Philippe Neri son image iconique avec sa vision de la Madone et la présence du lys ; pour Isidore le bâton et le surgissement de l’eau ; pour Thérèse d’Avila sa vision et le livre. Les étendards sont suspendus à des couronnes largement éclairées comme symboles de la gloire céleste.}
                \teiP[xmlid={p23-8}]{Sur les images de la présentation de la \teiHi[rend={italic}]{Chinea} du \teiHi[rend={small-caps}]{xvii}\teiHi[rend={sup}]{e} siècle dans la basilique et du \teiHi[rend={small-caps}]{xviii}\teiHi[rend={sup}]{e} siècle dans la \teiHi[rend={italic}]{Sala Regia} du Vatican, ainsi que sur celle de l’arrivée du cortège à Saint-Pierre, aucun emblème ni blason n’est figuré.}
              
\end{teiDiv}
              \begin{teiDiv}[type={section2},xmlid={div09-3}]

                \teiHead[xmlid={laccueil-des-personnalites-001}]{L’accueil des personnalités}
                \teiP[xmlid={p24-6}]{Christine de Suède\teiNote[n={24},place={foot},xmlid={ftn51-9}]{\teiP[xmlid={p-466}]{Carlo Festini, \teiHi[rend={italic}]{I trionfi della magnificenza pontificia celebrati per lo passaggio nelle città e luoghi dello Stato ecclesiatico e in Roma per lo ricevimento della Maestà della regina di Svezia, }In Roma, Nella Stamperia della Rev. Camera Apost., 1656. La bibliographie concernant Christine de Suède est infinie ; parmi bien d’autres titres en rapport direct avec notre sujet : Galeazzo Gualdo Priorato, \teiHi[rend={italic}]{Historia della Sua Reale Maestà Cristina Alessandra, Regina di Svezia}, In Roma, Nella Stamperia della Rev. Camera Apost., 1656 ; Per Bjurström, \teiHi[rend={italic}]{Feast and Theatre in Queen Christina’s Rome}, Stockholm, s. n., 1966 ; Georgina Masson, \teiHi[rend={italic}]{Queen Christina}, Londres, Secker \& Warburg, 1968 ; Cesare d’Onofrio, \teiHi[rend={italic}]{Roma val bene un’abiura}, Rome, Fratelli Palombi, 1976 ; Daniela Pizzagalli, \teiHi[rend={italic}]{La regina di Roma. Vita e misteri di Cristina di Svezia nell’Italia barocca}, Milan, Rizzoli, 2002.}}, convertie au catholicisme, est accueillie comme une reine par le pape Clément IX Rospigliosi, malgré son abdication. À son arrivée, il la loge au Vatican et lui offre un dîner en 1668 ; certes ils sont à deux tables séparées comme le prévoit le protocole ; le blason de Clément IX Rospigliosi et les \teiHi[rend={italic}]{Trionfi} symboliques qui décorent la table sont représentés sur l’image qui montre l’emblème papal au fond de la pièce. Lorsque sa résidence est transférée, sur volonté du pape, au palais Farnèse\teiNote[n={25},place={foot},xmlid={ftn52-9}]{\teiP[xmlid={p-467}]{Martine Boiteux, « Le palais Farnèse : la représentation et l’identité », \teiHi[rend={italic}]{Mélanges de l’École française de Rome}, vol. 122, n\teiHi[rend={sup}]{o} 2, 2010, p. 311-338.}}, son blason et celui de sa famille sont présents mais pas celui du pape sur la façade éphémère plaquée sur le palais\teiNote[n={26},place={foot},xmlid={ftn53-9}]{\teiP[xmlid={p-468}]{Michel Pastoureau, dans « L’emblématique des Farnèse », \teiHi[rend={italic}]{Le palais Farnèse}, Rome, École française de Rome, 1981, 2 vol., t. 1, p. 431-455, a montré la large utilisation des emblèmes dans le palais par les Farnèse, comme le faisaient les grandes familles romaines dans leur palais familial.}}.}
                \teiP[xmlid={p25-6}]{Les ambassadeurs rivalisent de faste à l’occasion des cortèges pour venir remettre leurs lettres de créance, et lors les audiences protocolaires. Ainsi une publication de Luca Antonio Chracas décrit le cortège du marquis de Fontès, ambassadeur portugais se rendant à l’audience du pape Clément XI le 8 juillet 1716\teiNote[n={27},place={foot},xmlid={ftn54-9}]{\teiP[xmlid={p-469}]{\teiHi[rend={italic}]{Distinto raguaglio del sontuoso treno delle carrozze con cui andò all’udienza di Sua Santità, il di 8 luglio 1716, l’illustrissimo ed eccellentissimo signore Don Rodrigo Annes de Saa, Almeida e Meneses, marchese di Fontes, conte di Pennaghiano}, In Roma, Nella Stamperia di Gio. Francesco Chracas, 1716.}}, jour de la fête de sainte Élisabeth reine du Portugal. Il souligne l’opulence des carrosses dorés ornés d’éléments héraldiques, mythologiques et symboliques, comme en témoigne un dessin de Vieira\teiNote[n={28},place={foot},xmlid={ftn55-9}]{\teiP[xmlid={p-470}]{Dario Beccarini, « Riflessioni e nuove proposte su Francisco Vieira de Matos, detto Vieira Lusitano, “eccellente disegnatore \lbrack{}…\rbrack{} felice e franco nell’inventare” », dans Elisa Debenedetti (dir.), \teiHi[rend={italic}]{Studi sul Settecento romano}, Rome, Edizioni Quazar, 2021, p. 207-220, ici p. 208.}} : le texte décrit dans le détail le carrosse des océans, qui évoque l’extension de l’empire portugais, pour l’ambassadeur, et quatre autres pour sa suite, véritables chars de triomphe, ainsi que dix carrosses d’escorte. Sont également énumérés les ornements dorés magnifiques surmontés par le blason de l’ambassadeur, les nombreux éléments allégoriques et symboliques comme le dragon symbole du Portugal, le soleil pour l’immortalité, le serpent qui se mord la queue pour l’éternité, le lion pour l’intrépidité.}
              
\end{teiDiv}
            
\end{teiDiv}
            \begin{teiDiv}[type={section1},xmlid={div10-3}]

              \teiHead[xmlid={conclusion-002}]{Conclusion}
              \teiP[xmlid={p26-6}]{Dans la basilique Saint-Pierre, les emblèmes servent la glorification, l’autoreprésentation du pouvoir pontifical et l’identité de ses représentants ; chaque pape laisse la marque de son passage et de ce qu’il accepte d’accueillir, personnes et cérémonies.}
              \teiP[xmlid={p27-6}]{Sur les monuments, sur les apparats éphémères et les gravures qui les représentent, les emblèmes servent comme ornements et comme signes identitaires. Blasons, armes et emblèmes sont nombreux. Une plus grande complexité héraldique est laissée aux créations durables tandis que les architectures éphémères, et les images qui en sont faites, donnent un image héraldique simplifiée.}
              \begin{teiFigure}[xmlid={figure01-11}]

                \teiGraphic[url={../icono/br/Ch12\_Boiteux/Fig159.jpg}]
                \teiHead[xmlid={figure-156-le-baldaquin-de-la-basilique-saint-pierre-001}]{Figure 156. Le baldaquin de la basilique Saint-Pierre.}
              
\end{teiFigure}
              \begin{teiFigure}[xmlid={figure02-11}]

                \teiGraphic[url={../icono/br/Ch12\_Boiteux/Fig160.jpg}]
                \teiHead[xmlid={figure-157-la-remise-des-clefs-bas-relief-de-ambrogio-buonvicino-1614-1632-facade-de-la-basilique-saint-pierre-001}]{Figure 157. La remise des clefs, bas-relief de Ambrogio Buonvicino, 1614-1632, façade de la basilique Saint-Pierre.}
              
\end{teiFigure}
              \begin{teiFigure}[xmlid={figure03-11}]

                \teiGraphic[url={../icono/br/Ch12\_Boiteux/Fig161.JPG}]
                \teiHead[xmlid={figure-158-monument-funeraire-du-pape-innocent-viii-cibo-1492-antonio-pollaiolo-basilique-saint-pierre-001}]{Figure 158. Monument funéraire du pape Innocent VIII Cibo, 1492, Antonio Pollaiolo, basilique Saint-Pierre.}
              
\end{teiFigure}
              \begin{teiFigure}[xmlid={figure04-11}]

                \teiGraphic[url={../icono/br/Ch12\_Boiteux/Fig162.jpg}]
                \teiHead[xmlid={figure-159-monument-funeraire-de-mathilde-de-canossa-1644-le-bernin-basilique-saint-pierre-001}]{Figure 159. Monument funéraire de Mathilde de Canossa, 1644, Le Bernin, basilique Saint-Pierre.}
              
\end{teiFigure}
              \begin{teiFigure}[xmlid={figure05-11}]

                \teiGraphic[url={../icono/br/Ch12\_Boiteux/Fig164.jpg}]
                \teiHead[xmlid={figure-160-plan-de-sede-vacante-au-deces-durbain-viii-barberini-1644-gravure-rome-istituto-centrale-per-la-grafica-001}]{Figure 160. Plan de \teiHi[rend={italic}]{Sede Vacante} au décès d’Urbain VIII Barberini, 1644, gravure, Rome, Istituto centrale per la grafica.}
              
\end{teiFigure}
              \begin{teiFigure}[xmlid={figure06-11}]

                \teiGraphic[url={../icono/br/Ch12\_Boiteux/Fig165.jpg}]
                \teiHead[xmlid={figure-161-plan-de-sede-vacante-au-deces-dalexandre-viii-ottoboni-avec-le-blason-du-cardinal-medicis-dedicataire-1691-gravure-001}]{Figure 161. Plan de \teiHi[rend={italic}]{Sede Vacante} au décès d’Alexandre VIII Ottoboni, avec le blason du cardinal Médicis dédicataire, 1691, gravure.}
              
\end{teiFigure}
              \begin{teiFigure}[xmlid={figure07-10}]

                \teiGraphic[url={../icono/br/Ch12\_Boiteux/Fig166.jpg}]
                \teiHead[xmlid={figure-162-cortege-du-possesso-de-leon-xi-medicis-1605-gravure-001}]{Figure 162. Cortège du \teiHi[rend={italic}]{Possesso} de Léon XI Médicis, 1605, gravure.}
              
\end{teiFigure}
              \begin{teiFigure}[xmlid={figure08-8}]

                \teiGraphic[url={../icono/br/Ch12\_Boiteux/Fig167.jpg}]
                \teiHead[xmlid={figure-163-theatre-de-la-canonisation-des-cinq-saints-1622-matteo-greuter-gravure-001}]{Figure 163. Théâtre de la canonisation des cinq saints, 1622, Matteo Greuter, gravure.}
              
\end{teiFigure}
            
\end{teiDiv}
          
\end{teiElement}
        
\end{teiElement}
        \begin{teiElement}[name={group},type={chapter},data-page-title={Apothéoses armoriales dans la Rome baroque},data-page-subtitle={De Clément VIII Aldobrandini à Clément X Altieri},data-page-authors={Frédéric Cousinié},data-include-href={Ch13\_Apotheoses\_armoriales.xml},xmlid={chapter-010}]

          \begin{teiElement}[name={front}]

            \begin{teiDiv}[type={titlePage},xmlid={titlepage-015}]

              \teiP[rend={title-main},xmlid={p-471}]{Apothéoses armoriales dans la Rome baroque}
              \teiP[rend={title-sub},xmlid={p-472}]{De Clément VIII Aldobrandini à Clément X Altieri}
              \teiP[rend={author-aut},xmlid={p-473}]{Frédéric Cousinié}
            
\end{teiDiv}
            \begin{teiElement}[name={epigraph}]

              \begin{teiCit}

                \begin{teiQuote}
« Cette salle Clémentine, immense, semblait sans bornes à cette heure, dans la clarté crépusculaire des lampes. La décoration si riche, les sculptures, les peintures, les dorures, se noyait, n’était plus qu’une vague apparition fauve, des murs de rêve où dormaient des reflets de joyaux et de pierreries. Et, d’ailleurs, pas un meuble, le dallage sans fin, une solitude élargie, se perdant au fond des demi-ténèbres. \lbrack{}…\rbrack{} Aujourd’hui encore, elle était la salle accessible à tous et que tous devaient traverser, simplement une salle des gardes, pleine toujours d’un tumulte de pas, d’allées et de venues sans nombre. Mais quelle mort pesante, dès que la nuit l’avait envahie, et comme elle était désespérée et lasse d’avoir vu défiler tant de choses et tant d’êtres ! »
\end{teiQuote}
                \begin{teiBibl}[xmlid={bibl-008}]
Émile Zola, \teiHi[rend={italic}]{Les Trois Villes. Rome}, Paris, G. Charpentier et E. Fasquelle, 1896, chap. XIV, p. 600.
\end{teiBibl}
              
\end{teiCit}
            
\end{teiElement}
          
\end{teiElement}
          \begin{teiElement}[name={body},xmlid={body-015}]

            \teiP[xmlid={p1-15}]{La célébration des grandes dynasties princières de l’époque moderne a bien souvent pris la forme de grandes compositions monumentales ascensionnelles. Le héros familial – ancêtre mythique ou glorieux prédécesseur – y est emporté dans un halo lumineux et élevé vers le ciel : pensons à la flagorneuse composition de Luca Giordano pour les Médicis du palais Medicis-Riccardi à Florence (1684-1686) ou, du même peintre, à sa non moins obséquieuse \teiHi[rend={italic}]{Gloria de la Monarquía de España} qui magnifie la voûte du grand escalier de l’Escorial. Le modèle reconduit celui, en vigueur pour les empereurs romains, de l’apothéose (étymologiquement un « devenir-dieu » : une déification) et il redouble les propres apothéoses ou glorifications des saints catholiques qui se multiplièrent à l’âge baroque. Dans plusieurs cas néanmoins, la représentation figurative d’un individu singulier a été éclipsée par celle du blason familial lui-même élevé « en gloire » vers le ciel selon une forme inédite dans la longue histoire de cet objet. L’apothéose individuelle s’y conjugue ou se soumet à l’apothéose collective d’une dynastie\teiNote[n={1},place={foot},xmlid={ftn1-13}]{\teiP[xmlid={p-474}]{Sur la glorification des pouvoirs religieux et civils à l’époque moderne, nous nous permettons de renvoyer à Frédéric Cousinié, Annick Lemoine, Michèle Virol (dir.), \teiHi[rend={italic}]{Soleils baroques. La Gloire de Dieu et des Princes en représentation dans l’époque moderne}, Rome-Paris, Académie de France à Rome-Somogy, 2018.}}.}
            \teiP[xmlid={p2-14}]{Avant même que le motif peint ou sculpté de la gloire lumineuse ne se banalise au \teiHi[rend={small-caps}]{xvii}\teiHi[rend={sup}]{e} siècle, ce dispositif était déjà fréquemment employé par les graveurs. Nous le trouvons par exemple, dès 1612, gravé par le bolonais Oliviero Gatti (1579-1649), dans les armes de Maffeo Barberini (devenu cardinal-prêtre depuis 1606), dont la lourde masse en suspension qui se détache sur un soleil rayonnant est flanquée d’Apollon et d’une autre allégorie féminine. Le même artiste est également l’auteur en 1615, d’après une composition d’Annibal Castelli, d’une réalisation semblable pour la dynastie des Farnèse\teiNote[n={2},place={foot},xmlid={ftn44-9}]{\teiP[xmlid={p-475}]{John T. Spike (éd.), \teiHi[rend={italic}]{The Illustrated Bartsch}, New York, Abaris Book, 1981, vol. 41, Oliver Gatti, et 41 (15), 1612 ; 50 (20), 1615.}}. Giovanni Luigi Valesio (\teiHi[rend={italic}]{circa} 1583-1633), autre graveur bolonais, affectionnait également ce motif auquel il recourt par exemple pour les armes du cardinal Ludovisi entourées des quatre éléments et surmontant le char solaire ; pour celles du cardinal Spinola dont les armes rayonnantes et en lévitation au-dessus de Rome se superposent sur une grande tapisserie tendue par deux figures allégoriques à l’arrière-plan ; ou pour celles encore, destinées à un frontispice d’ouvrage, du cardinal Scipione Borghese, supportées par \teiHi[rend={italic}]{la Renommée}\teiNote[n={3},place={foot},xmlid={ftn45-9}]{\teiP[xmlid={p-476}]{Veronika Birke (éd.), \teiHi[rend={italic}]{The Illustrated Bartsch}, New York, Abaris Book, 1982, vol. 40, part. 1 et 1987, part. 2, 91 (234) ; 95 (235) et 100 (238).}}. Dans tous ces cas l’écu, doté de son cartouche, conserve sa forte présence opaque, la lumière venant constituer un nouveau fond immatériel qui illumine \teiHi[rend={italic}]{de retro} le motif.}
            \teiP[xmlid={p3-14}]{Pour la grande sculpture monumentale romaine, la formule, déjà utilisée par la célébration du trigramme jésuite, a trouvé sa plus remarquable incarnation dans l’étonnante voûte de la chapelle Spada à Santa Maria in Vallicella (conçue par Camillo Arcucci puis Giuseppe Brusati Arcucci et Carlo Rainaldi), dont les stucs sont dus à Giovanni Francesco De Rossi (1666-1679), l’un des collaborateurs du Bernin pour la décoration intérieure de Santa Maria del Popolo. L’œuvre, que j’ai pu ailleurs étudier en détail\teiNote[n={4},place={foot},xmlid={ftn46-9}]{\teiP[xmlid={p-477}]{Frédéric Cousinié, \teiHi[rend={italic}]{Gloriae. Figurabilité du divin, esthétique de la lumière et dématérialisation de l’œuvre d’art à l’âge Baroque}, Rennes, PUR, 2018, « La chapelle Spada/Borromée de S. Maria in Vallicella », p. 181-206.}}, associe de façon exceptionnelle le motif rayonnant de la gloire, omniprésent à Rome sur nombre d’autels, à l’écu, détaché par un ange, de saint Charles Borromée. Sa présence joue encore, moyennant toute une série de relations (supports de l’architecture, géométrie) avec les armes des fondateurs de la chapelle, disposées au centre du pavement de marbre.}
            \teiP[xmlid={p4-14}]{Dans le champ de la peinture, la conjonction de l’élévation ou de la suspension/ lévitation et de la glorification armoriale s’applique à plusieurs exemples exceptionnels du \teiHi[rend={small-caps}]{xvii}\teiHi[rend={sup}]{e} siècle. Leur analyse se propose d’approcher la poétique inédite qui les organise et de préciser leurs implications interprétatives. Nous constaterons notamment que l’expression de la \teiHi[rend={italic}]{potestas} pontificale ne se donne pas dans le seul déchiffrement de ses signes mais dans l’intégration sensible et subjective des dispositifs scéniques et cérémoniaux qui lui sont associés.}
            \begin{teiDiv}[type={section1},xmlid={div01-12}]

              \teiHead[xmlid={seul-en-scene-lapotheose-ludovisi-circa-1623-001}]{Seul en scène : l’apothéose Ludovisi (\teiHi[rend={italic}]{circa} 1623)}
              \teiP[xmlid={p5-14}]{Le premier exemple, le moins connu, tardivement attribué et peu accessible, est celui de la composition représentant \teiHi[rend={italic}]{L’Apothéose des armes Ludovisi} (\teiHi[rend={italic}]{L’Apoteosi dello stemma dei Ludovisi}), peinte vers 1623 sur la voûte d’une des salles dite « dell’armadio de’ libri » ou « della Stemma » (petite bibliothèque et salle de musique) du Casino Ludovisi (fig. 164). Le lieu est célèbre pour la fameuse \teiHi[rend={italic}]{Aurore} du Guerchin et nous savons qu’il constitue le dernier vestige de la villa du cardinal Ludovico Ludovisi (1595-1632), l’un des deux neveux d’Alessandro Ludovisi, archevêque de Bologne devenu pape sous le nom de Grégoire XV (1621-1623). L’exemple est important, car il serait dû à l’un des graveurs bolonais déjà évoqués qui s’illustrèrent par leur prolixité et leur inventivité en matière héraldique : Giovanni Luigi Valesio, auquel l’œuvre a été attribuée par Armando Schiavo sur la base des mentions apportées par Giovanni Baglione en 1642 et Fioravante Martinelli dans sa description de Rome en 1660\teiNote[n={5},place={foot},xmlid={ftn47-9}]{\teiP[xmlid={p-478}]{Voir Armando Schiavo, \teiHi[rend={italic}]{Villa Ludovisi e Palazzo Margherita}, Rome, Roma Amor, 1981, p. 128 et 129, citant Baglione, p. 354 : « \teiHi[rend={italic}]{Dipinse in quel Palagio alcune stanze con diversi capricci di puttini in fresco coloriti} » ; Carolyn Harwood Wood, \teiHi[rend={italic}]{The Indian Summer of Bolognese Painting: Gregory XV (1621-1623) and Ludovisi art patronage in Rome}, thèse, The University of North Carolina, 1988, chap. II, « The Casino Ludovisi », p. 94-99, qui relève bien p. 96 que les armes ne sont pas juste « \teiHi[rend={italic}]{hovering but ascending} », ; Carla Benocci, \teiHi[rend={italic}]{Villa Ludovisi}, Rome, Istituto poligrafico e Zecca dello Stato-Libreria dello Stato, 2010, chap. IV, p. 156, sur le salaire payé à Valesio comme « \teiHi[rend={italic}]{guardaroba} » et pour diverses peintures, et p. 172-174, sur un possible programme alchimique et spirituel qui organiserait le décor du Casino ; en dernier lieu, Laura Bartoni, « L’allestimento delle collezioni nella villa pinciana: spazi, criteri, funzioni », \teiHi[rend={italic}]{Storia del’arte – « Guercino nel Casino Ludovisi (1621-1623) »}, vol. 157, n\teiHi[rend={sup}]{o} 1, 2022, p. 99-117, notamment p. 101 et 102 sur le décor de la salle « \teiHi[rend={italic}]{dell’armadio de’libri} » et, dans le même volume, Caterina Volpi, « La collezione Ludovisi come ″Scuola del mondo″. Il ruolo di Agucchi, Marino, Valesio, Guercino e Domenichino », \teiHi[rend={italic}]{ibid., }p. 118-133. Voir également, sur l’artiste, Kenichi Takahashi, \teiHi[rend={italic}]{Giovanni Luigi Valesio. Ritratto de « l’Instabile academico incaminato »,} Bologne, Clueb, 2007 et Giovanna Perini Folesani, « Pittura e poesia, arti sorelle nella Bologna rinascimentale e barocca », \teiHi[rend={italic}]{Letteratura e arte}, n\teiHi[rend={sup}]{o} 18, 2020, p. 67-80. Sur l’œuvre gravée, Veronika Birke (éd.), \teiHi[rend={italic}]{op. cit.}, part. 1; et \teiHi[rend={italic}]{ibid}., part. 2.}}. L’artiste, actif dans de nombreux domaines (musique, ballet, miniature, calligraphie, poésie), fut également théoricien, auteur d’un traité sur le dessin et d’un discours sur la nouvelle rhétorique, publié en défense du poète Giovan Battista Marino. Passé en 1621 au service des Ludovisi à Rome, il fut aussi « gardien » des jardins de Porta Pinciana, avant d’être chargé de l’apparat funèbre célébrant à Bologne l’anniversaire de la mort du pape le 25 juillet 1624.}
              \teiP[xmlid={p6-14}]{La voûte rectangulaire de la salle \teiHi[rend={italic}]{Dell’armadio de’libri} est occupée par une balustrade feinte alternant balustres et piédestaux couverts de riches marbres feints, sur laquelle prennent appui de nombreux angelots ou \teiHi[rend={italic}]{putti} représentés dans des attitudes variées et animées. Leur fantaisie est proche de celle qui caractérise les angelots supportant l’écu papal déjà en suspension au sommet de la chambre de la Signature de Raphaël (1508) (fig. 165), et elle peut encore rappeler les \teiHi[rend={italic}]{putti} manipulant les lettres du nom de Jules III (Giovanni Maria del Monte) dans les appartements du pontife au Belvédère\teiNote[n={6},place={foot},xmlid={ftn48-10}]{\teiP[xmlid={p-479}]{Ce dernier exemple, aimablement signalé par Yvan Loskoutoff, est attribué à Prospero Fontana et à son atelier : voir Denis Ribouillault, « Jeux de mots et d’images : les frises peintes de l’appartement de Jules III au Vatican (vers 1550-1553) », dans Antonella Fenech Kroke, Annick Lemoine (dir.), \teiHi[rend={italic}]{Frises peintes. Les décors des villas et palais au Cinquecento}, Paris-Rome, Somogy-Académie de France à Rome, 2016, p. 103-129 qui signale également d’autres exemples. Je remercie Elli Doulkaridou pour ses suggestions lors de nos échanges durant le colloque.}} (1550-1553). Elle anticipe également les attitudes des \teiHi[rend={italic}]{putti} jouant avec des guirlandes que l’on retrouve plus tard sur la frise intérieure de l’église de Sant’Ignazio (\teiHi[rend={italic}]{circa} 1649-1650), dont Ludovico Ludovisi fut précisément le fondateur, conjuguant ainsi une dimension ludique au jeu « sérieux » des contenus héraldiques et idéologiques de ces programmes\teiNote[n={7},place={foot},xmlid={ftn49-10}]{\teiP[xmlid={p-480}]{Voir Jennifer Montagu, \teiHi[rend={italic}]{Alessandro Algardi}, New Haven-Londres, Yale University Press, 1985, 2 vol., t. I, p. 110-113 et cat. A. 202, p. 456 et 457 avec référence aux archives : Archivio della Compagnia Di Gesù, Fondo Gesuitico, Fabbrica di Sant’Ignazio, vol. 1343, p. 339-354 et vol. 1344, p. 232 ; Jacopo Curzietti, \teiHi[rend={italic}]{Antonio Raggi: scultore ticinese nella Roma barocca}, Canterano, Aracne editrice, 2020, cat. 3 (Œuvres rejetées), p. 296-297. L’Algarde, au revers de la façade, a donné les modèles pour les figures de la \teiHi[rend={italic}]{Religion} et de la \teiHi[rend={italic}]{Magnificence} qui encadrent le grand cadre de stuc avec les inscriptions rappelant le rôle des deux frères Ludovisi dans la construction de l’église. La frise en stuc avec des \teiHi[rend={italic}]{putti} supportant les écus du cardinal Ludovico Ludovisi, de son frère le prince Nicolò (associé à celui de sa femme), a été réalisée d’après un projet de l’Algarde vers 1649-1650, par une série de sculpteurs sous la direction de Stefano Castelli. Montagu attribue les sculptures à Raggi ou Giovan Battista Morelli, Curzietti rejetant la première attribution.}}. Au Casino Ludovisi, un long phylactère ondulant unit ces figures et célèbre comme il se doit le propriétaire des lieux. On peut y lire, en lettres détachées qui obligent à faire le tour de la composition, le nom du pape régnant et celui de « Ludovicus Ludovisi », suivi de son titre de cardinal camerlingue de la sainte Église : « \teiHi[rend={small-caps italic}]{Ludovicus Ludovisius S.R.E. Cardinalis Camerarius SS. Domini Nostri Gregorij PP. XV Nepos}. » Essentiel et original est le motif central isolé dont le rouge (« de gueules ») éclate sur le fond ou le « champ » bleu intense (d’« azur ») d’un \teiHi[rend={italic}]{cielo aperto} que l’on retrouvera dans les autres exemples\teiNote[n={8},place={foot},xmlid={ftn50-10}]{\teiP[xmlid={p-481}]{La prégnance du motif tient à son isolement dans le champ et à l’opposition radicale, par superposition de deux des couleurs primaires : rouge et bleu quasi giottesque. Ce jeu entre primaires se retrouve dans l’œuvre de Pietro da Cortona pour la salle de Mars du palais Pitti (étendu au jaune) et, de façon plus modérée, au palais Barberini où le peintre donne leur présence aux formes par un usage plus marqué des valeurs (clair-obscur). Ces contrastes sont en revanche absents tant de la réalisation de Canuti que de celle des frères Alberti, fondées sur une harmonisation tonale (avec usage de couleurs primaires atténuées et de couleurs secondaires : orange, violet, vert) et un faible contraste des valeurs.}}. Insistons sur le fait que l’intervalle ou l’espace vide et marginal qui n’avait pour fonctions que de marquer l’indépendance de l’écu à l’égard de son support (les écus « détachés »), ou sa mobilité et l’intervalle temporel entre deux instants (l’un des projets du Bernin pour l’écu en suspension du tombeau d’Alexandre VII qui échappe à l’allégorie de la Mort), se déploie désormais dans la profondeur d’un fond dominant – vraie reconquête picturale et pariétale du ciel et du champ – qui prend désormais une valeur transcendantale : il indique le lieu céleste visé par l’écu. Anges et \teiHi[rend={italic}]{putti} agglutinés y supportent ou sont supportés par le vaste écu familial dépourvu de cartouche – insistant ainsi sur le décrochement aérien du motif mis en gloire – et surmonté du chapeau cardinalice en suspension dont virevoltent les cordons. Une épaisse guirlande de feuilles et de fruits exprime l’abondance et la fertilité associées à l’action du cardinal honoré, tout en faisant écho aux trophées ou festons végétaux fleuris que portent encore les \teiHi[rend={italic}]{putti} de la balustrade.}
              \teiP[xmlid={p7-14}]{Nous pourrions retrouver certaines des formules déjà employées par Valesio dans ses propres gravures célébrant le même cardinal. Celle où une cohorte de \teiHi[rend={italic}]{putti} tenait déjà en suspension aérienne l’écu surmonté du chapeau et présenté à une figure agenouillée ; ou celle, plus proche de la formule commune du \teiHi[rend={italic}]{Stemma in arrivo}, où \teiHi[rend={italic}]{putti} et allégorie de la \teiHi[rend={italic}]{Renommée} présentent l’armoirie cardinalice à deux héros antiques auxquels le cardinal s’identifiait : Jules César et Fabius Cunctator (Quintus Fabius Maximus Verrucosus), consul et héros de la guerre contre Hannibal loué pour sa prudence (le Temporisateur : \teiHi[rend={italic}]{cunctator}\teiNote[n={9},place={foot},xmlid={ftn51-10}]{\teiP[xmlid={p-482}]{Veronica Birke (éd), \teiHi[rend={italic}]{The Illustrated Bartsch}, vol. 40, \teiHi[rend={italic}]{op. cit}., respectivement 73 (228) et 89 (23).}}). Dans le Casino Ludovisi, le motif est en revanche splendidement isolé de tout contexte narratif ou allégorique et il célèbre, par métaphore et métonymie, autant l’événement ponctuel de « l’élévation » à la pourpre du \teiHi[rend={italic}]{cardinal-nipote}, le 15 février 1621 que, plus largement, l’apothéose céleste ou la gloire mondaine de la dynastie Ludovisi évoqué par le blason familial.}
              \teiP[xmlid={p8-14}]{Si l’amplification monumentale de l’écu, la suspension, l’élévation et la mise en gloire sont les principes déterminants auxquels recourt Valesio, ils vont de pair avec une forme de léger pivotement du motif qui implique un dédoublement spéculaire : la création d’au moins deux points de vue privilégiés pour le destinataire. L’écu est ainsi affecté d’une inclinaison axiale qui contraste avec la présentation frontale dominante des nombreuses armoiries pontificales ou cardinalices fixées à l’aplomb de la plupart des appartements princiers ou des lourds et statiques plafonds à caissons des églises romaines (Santo Crisogono, Santa Maria in Trastevere, Santi Cosma e Damiano, plus tard Santo Clemente, San Giovanni in Laterano, etc.). Ce choix oblige au déplacement du visiteur, instaurant une « relation de mobilité réciproque\teiNote[n={10},place={foot},xmlid={ftn52-10}]{\teiP[xmlid={p-483}]{Gérard Le Don, \teiHi[rend={italic}]{Voir la sculpture. Essai sur le dispositif visuel}, Rennes, PUR, « Aesthetica », 2018, p. 48.}} », tout en participant d’une forme de déséquilibre qui anime avec vivacité toute la composition, conforme en cela au surnom d’« \teiHi[rend={italic}]{Instabile} academico incaminato » qu’avait reçu l’artiste à Bologne.}
            
\end{teiDiv}
            \begin{teiDiv}[type={section1},xmlid={div02-12}]

              \teiHead[xmlid={dematerialisation-inchoativite-contextualisation-les-apotheoses-barberini-1632-1636-et-medici-1644-1646-001}]{Dématérialisation, inchoativité, contextualisation : les Apothéoses Barberini (1632-1636) et Médici (1644-1646)}
              \teiP[xmlid={p9-13}]{Les historiens de l’art ont relevé l’importance qu’a sans doute eu la réalisation novatrice de Valesio pour un second ensemble constitué par deux des œuvres les plus célèbres et spectaculaires de Pietro da Cortona : la fresque s’étendant sur la voûte du grand salon du palais Barberini (1632-1638) à Rome et celle qui occupe, à Florence, la voûte d’une des pièces du palais Pitti, la dite « Anticamera de’ cortigiani » ou salle de Mars (1644-1646). Sans revenir sur le complexe ensemble d’allégories et de scènes narratives qui, absent de la composition de Valesio, occupe les marges de ces deux compositions, je m’arrêterai sur le seul motif central\teiNote[n={11},place={foot},xmlid={ftn53-10}]{\teiP[xmlid={p-484}]{Voir, respectivement, les études fondatrices de Walter Vitsthum, « A comment on the Iconography of Pietro da Cortona’s Barberini Ceiling », \teiHi[rend={italic}]{The Burlington Magazine}, vol. 103, n\teiHi[rend={sup}]{o} 703, octobre 1961, p. 427-433, repris dans Susan M. Dixon (dir.), \teiHi[rend={italic}]{Italian Baroque Art. Blackwell Anthologies in Art History}, Malden-Oxford, Carlton, Blackwell Publishing, 2008, p. 201-208 ; John Beldon Scott, \teiHi[rend={italic}]{Images of Nepotism. The Painted Ceilings of Palazzo Barberini,} Princeton, Princeton University Press, 1991, t. 3, chap. X-XII ; Sebastian Schütze, « Urbano VIII e il concetto di Palazzo Barberini. Alla ricerca di un primato culturale di rinascimentale memoria », dans Christoph Luitpold Frommel et Sebastian Schütze (dir.), \teiHi[rend={italic}]{Pietro da Cortona}, Milan, Electa, 1998, p. 86-97 ; et Malcolm B. Campbell, \teiHi[rend={italic}]{Pietro da Cortona at the Pitti Palace. A Study of the Planetary Rooms and Related Projects}, Princeton, Princeton University Press, 1977. La lecture iconographique du complexe romain est fondée pour l’essentiel sur Francesco Bracciolini, \teiHi[rend={italic}]{L’elettione di Urbano Papa VIII}, \lbrack{}Rome\rbrack{}, s. n., \lbrack{}1628\rbrack{}, les \teiHi[rend={italic}]{Vite de Pittori} de Giovanni Battista Passeri, la description de Mattia Rosichino, \teiHi[rend={italic}]{Dichiarazione delle Pitture della Sala de’ Signori Barberini in Rome}, In Roma, Appresso Vitale Mascardi, 1640 et celle de Girolamo Teti, \teiHi[rend={italic}]{Aedes Barberinae ad Quirinalem}, Romæ, Mascardus, 1642.}}.}
              \teiP[xmlid={p10-13}]{À Rome (fig. 166), les armes des Barberini, réduites aux seules abeilles devenues indépendantes de l’écu désormais évoqué \teiHi[rend={italic}]{in absentia} par son seul contour et ses « meubles », sont entourées d’une couronne de lauriers (autre emblème familial) que soutiennent les trois vertus théologales (sans attributs) de la \teiHi[rend={italic}]{Foi}, de l’\teiHi[rend={italic}]{Espérance} et de la \teiHi[rend={italic}]{Charité} (où Giovanni Battista Passeri voulait voir « \teiHi[rend={italic}]{le tre muse principali, Urania, Calliope e Clio} »). Au sommet, elles sont complétées de l’allégorie de la \teiHi[rend={italic}]{Religion} tenant les immenses clefs de saint Pierre, d’une personnification de Rome qui couronne de la tiare papale le groupe central, et d’un \teiHi[rend={italic}]{putto} sur la gauche portant une autre couronne évoquant les qualités poétiques attribuées au pape-poète qu’était Urbain VIII. En contrebas, obéissant à la \teiHi[rend={italic}]{Divine Providence} occupant la partie inférieure de la fresque et qui est le moteur central de l’action, une allégorie de l’\teiHi[rend={italic}]{Immortalité} (ou de la Gloire/\teiHi[rend={italic}]{Fama}), tenant une couronne d’étoiles, est censée rejoindre le groupe principal. Dans ce qui était décrit par Giovanni Baglione ou plus tard Filippo Titi comme un « triomphe de la gloire » (\teiHi[rend={italic}]{il Trionfo della Gloria}) célébrant le règne du pontife orné des plus hautes vertus et accédant à l’immortalité, nous assistons donc à un triple, voire à un quadruple couronnement surclassant la simple élévation cardinalice de Ludovico Ludovisi.}
              \teiP[xmlid={p11-13}]{Nous retrouvons dans cette composition complexe, outre la monumentalisation générale, plusieurs des procédés associés fréquemment aux armoiries du baroque romain : l’adjonction d’éléments signifiants complémentaires autour des meubles familiaux (les allégories), l’évidement de la surface de l’écu (où les « meubles » que sont les abeilles se détachent directement sur le ciel et se rassemblent à l’intérieur des branches du laurier), la dimension inchoative de la scène où l’écu, loin d’être donné dans une présence éternelle, paraît se former sous nos yeux par agglomération de ses divers éléments (abeilles, lauriers, couronnes\teiNote[n={12},place={foot},xmlid={ftn54-10}]{\teiP[xmlid={p-485}]{Voir Irving Lavin, « Bernini’s Bumbling Barberini Bees », dans Irving Lavin, \teiHi[rend={italic}]{Visible Spirit. The Art of Gianlorenzo Bernini}, Londres, The Pinder Press, 2009, t. 2 ; p. 955-1017, ici p. 1004-1005, qui met en relation cette assemblée spontanée d’abeilles avec les circonstances fabuleuses de l’élection du nouveau pape où un essaim d’abeilles aurait pénétré le Conclave des cardinaux. Le motif de l’autoconstitution de l’écu apparaît à plusieurs reprises dans les gravures de Valesio, Oliviero Gatti ou Giovanni Battista Coriolano.}}), l’élévation armoriale et le soutien de ce qui s’est substitué aux limites de l’écu (les rameaux de laurier), la suspension et bien sûr la mise en gloire qu’avait déjà expérimentées Valesio. Ajoutons – nouvelle amplification des recherches de l’artiste bolonais – l’animation générale impliquée par le quadruple couronnement – avec un effet de suspens temporel en trois temps correspondant à la plus ou moins grande proximité des couronnes en instance d’être apposées sur le motif central – ; la multiplication et la diffusion des meubles ou des emblèmes des Barberini (on retrouve ailleurs les abeilles ou le soleil) ; et enfin l’incorporation spatiale de cette partie de la fresque dans un vaste et complexe ensemble narratif, symbolique et emblématique qui se déploie latéralement en diverses scènes.}
              \teiP[xmlid={p12-13}]{À Florence (fig. 167), programme qui nous éloigne un instant du contexte pontifical mais qui est fondé sur une logique analogue, la salle de Mars du palais des Offices avait été décorée par Pietro da Cortona entre 1644 et 1646 pour le grand-duc Ferdinand II de Medici (1610-1670). Elle s’insère au sein d’un ensemble de cinq pièces consacrées aux principales divinités antiques. Celles-ci sont mises en relation avec l’itinéraire ascendant d’un prince idéal (donné en modèle au jeune fils du duc, futur Cosme III), depuis ses exercices vertueux dans la salle de Vénus jusqu’à son ascension céleste vers la \teiHi[rend={italic}]{Gloire} et l’\teiHi[rend={italic}]{Éternité} dans la pièce placée sous l’invocation de Saturne. La voûte de la salle de Mars est occupée, sur les côtés, par divers épisodes héroïques où le prince combattant est associé à divers dieux et héros antiques (Mars, Hercule, les Dioscures). Au centre, le motif principal est de nouveau l’armoirie familiale des Medici avec les cinq boules (\teiHi[rend={italic}]{palle}) rouges et une bleue décorée par des lis d’or. L’écu, dans ce cas maintenu dans son intégrité unitaire, se confond avec son cartouche et s’appuie encore sur un motif ondulant évoquant une grande coquille soutenue par trois \teiHi[rend={italic}]{putti}. La monumentalité des armoiries, qui avait été annulée au palais Barberini par leur dématérialisation, est ici intensifiée. L’ensemble se démarque encore d’un grand cercle d’armes agglutinées formant trophée, tandis que l’écu reçoit une couronne portée par d’autres \teiHi[rend={italic}]{putti,} reprenant le motif déjà employé à Rome par le peintre. Comme l’a relevé Malcolm B. Campbell\teiNote[n={13},place={foot},xmlid={ftn55-10}]{\teiP[xmlid={p-486}]{Malcom B. Campbell, \teiHi[rend={italic}]{op. cit}., p. 153 et 154.}}, le nom inscrit à l’intérieur de la couronne fait de la célébration de la dynastie des Medici une apothéose également \teiHi[rend={italic}]{personnelle} du duc Ferdinand. L’ensemble, que l’on peut rapprocher de la fresque d’Angelo Michele Colonna et Agostino Mitelli célébrant Alexandre le Grand (1641), autre modèle de Ferdinand II de Médicis, au palais Pitti, ne sera pas sans incidences sur les grandes réalisations florentines d’Antonio Domenico Gabbiani (1652-1726) : la fresque de \teiHi[rend={italic}]{L’Apothéose de la famille Corsini} au palais Corsini al Parione (fig. 168) et celle de \teiHi[rend={italic}]{L’Apothéose de Cosme l’Ancien} (ou \teiHi[rend={italic}]{Cosme l’Ancien conduit par la gloire au trône de Jupiter}) à la villa de Poggio a Caiano (1698), où la présence céleste du héros familial est associée à celle de l’écu Medici. La formule trouvera l’un de ses ultimes accomplissements dans la fresque d’Ermenegildo Costantini (1731-1791) pour le \teiHi[rend={italic}]{Triomphe de la Maison Borghese et des arts} (1768) dans la salle d’audience du palais Borghese à Rome\teiNote[n={14},place={foot},xmlid={ftn56-9}]{\teiP[xmlid={p-487}]{Autre exemple tardif plus modeste : la fresque centrale du \teiHi[rend={italic}]{pòrtego} de Costantino Cedini (1773-1783), au palais Corner della Regina à Venise. Pour un parcours sélectif et synthétique dans ces grands cycles, voir Steffi Roettgen, \teiHi[rend={italic}]{Fresques italiennes du Baroque aux Lumières (1600-1797)}, Jeanne Étoré-Lortholary (trad.), Paris, Citadelle \& Mazenod, 2007. En France, relevons le projet de Charles Le Brun pour les armes de Fouquet puis du roi dans le grand salon du château de Vaux-Le-Vicomte.}}.}
              \teiP[xmlid={p13-13}]{La proximité est grande avec l’autre réalisation romaine de Pietro da Cortona ; la principale différence, outre le maintien du volume de l’écu que présentait également Valesio, étant le mouvement ascensionnel dynamique de la réalisation florentine qui s’oppose à l’effet de suspension immobilisée de Rome, où ce sont les figures allégoriques qui prennent davantage en charge l’animation. Le modèle est plutôt à rechercher dans l’autre réalisation de Pietro da Cortona pour l’une des petites galeries du palais Barberini (1632), où était également employé le motif du \teiHi[rend={italic}]{Stemma in arrivo} associé au couronnement de l’écu (fig. 169). Une autre différence, qui nous rapproche de la solution de Valesio ou de celle de la petite galerie du palais Barberini, est le pivotement angulaire de l’armoirie, qui se distingue de la position, plus orthogonale par rapport au sol, de la réalisation du grand salon du même palais. S’il est donné à la contemplation du spectateur présent dans la salle du palais Pitti, l’écu des Medici est ici d’abord dirigé, au sein même de la composition, vers le prince héroïque qui est représenté sur la partie gauche de la composition.}
            
\end{teiDiv}
            \begin{teiDiv}[type={section1},xmlid={div03-10}]

              \teiHead[xmlid={lecu-familial-et-lancetre-fondateur-lapotheose-altieri-circa-1675-1676-001}]{L’écu familial et l’ancêtre fondateur : l’apothéose Altieri (\teiHi[rend={italic}]{circa} 1675-1676)}
              \teiP[xmlid={p14-10}]{À ces deux célèbres exemples il faut ajouter, troisième ensemble exceptionnel\teiNote[n={15},place={foot},xmlid={ftn57-9}]{\teiP[xmlid={p-488}]{Nous pourrions ajouter la réalisation, plus modeste mais fondée sur le même principe suspensif, de la fresque du salon du palais du prince Camillo Pamphili à Valmontone (Gaspard Dughet et Guglielmo Cortese, 1658 ; le motif étant sans doute postérieur) où l’écu assemble les meubles du Prince et de son épouse Olimpia Aldobrandini : voir Marie-Nicolas Boisclair, \teiHi[rend={italic}]{Gaspard Dughet, 1615-1675}, Paris, Arthéna, 1986, cat. 171-175, p. 226 ; ou celle de Francesco Cozza pour la bibliothèque (1667-1673) du palais Doria-Pamphilj (place Navone) : l’écu familial, tenu en suspension par un ange dans ciel, n’y a qu’une place modeste, mais l’œuvre est intéressante pour la dissémination des meubles familiaux dans le reste de la fresque, qui évoque la réalisation antérieure de la Sala Clementina.}}, le cas des décors plus tardifs du palais Altieri à Rome. Conçu autour de 1633 par Lorenzo Altieri à proximité du Gesù, puis étendu par l’architecte Giovanni Antonio De’Rossi entre 1650-1655 pour le cardinal Giambattista Altieri, le palais fut complété entre 1670 et 1676 à l’occasion de l’élection d’Emilio Altieri (frère de Giambattista) comme pontife sous le nom de Clément X (1670-1676). Le palais était alors occupé par le cardinal Paluzzo Paluzzi-Albertoni, \teiHi[rend={italic}]{cardinal-nipote} adopté par le pape, ainsi que par son frère Angelo et son neveu Gaspare Paluzzi-Albertoni\teiNote[n={16},place={foot},xmlid={ftn58-9}]{\teiP[xmlid={p-489}]{Le premier était doté des titres de « \teiHi[rend={italic}]{Vicario di Roma} », « \teiHi[rend={italic}]{Camerlengo di Santa Chiesa} » et « \teiHi[rend={italic}]{Prefetto di Propaganda Fide} », le second était « \teiHi[rend={italic}]{Generale di Santa Romana Chiesa} ».}}, époux depuis 1667 de Laura Caterina Altieri, seule descendante de la famille.}
              \teiP[xmlid={p15-9}]{Le palais est célèbre pour abriter la fresque de la Sala della Clemenza peinte par Carlo Maratta (1674-1675), tardive et quelque peu plate réponse « classique » au chef-d’œuvre conçu par Pietro da Cortona pour les Barberini. L’usage des éléments du blason papal et familial y est largement présent, quoique relativement discret et associé à plusieurs figures porteuses dont les armoiries restent tributaires : à gauche, la \teiHi[rend={italic}]{Prudenza}-Minerva s’y appuie sur l’écu portant les étoiles des Altieri tandis que le personnage vêtu à l’antique présent sur la droite de la composition, évoquant Gaspare Paluzzi-Albertoni, porte le « \teiHi[rend={italic}]{gonfalone} » papal avec les clés de saint Pierre ; au sommet de la composition, deux figures de la \teiHi[rend={italic}]{Fama} supportent encore la tiare et, de nouveau, les clés.}
              \teiP[xmlid={p16-9}]{À proximité, dans l’antichambre du palais qui était commune aux deux appartements familiaux, le plafond est occupé par une vaste fresque réalisée par le bolonais Domenico Maria Canuti, célèbre pour ses grands décors au palais Pepoli de Bologne (\teiHi[rend={italic}]{Apothéose d’Hercule}, 1669-1671), puis à la galerie Colonna (1672-1673) et à l’église romaine des Santi Domenico e Sisto (\teiHi[rend={italic}]{Apothéose de saint Dominique}, 1673-1675). Au palais Altieri, la fresque représente en son centre \teiHi[rend={italic}]{L’Apothéose de Romulus }(\teiHi[rend={italic}]{circa} 1675-1676) (fig. 170). Lieu commun de nombreux décors romains (\teiHi[rend={italic}]{L’Apothéose d’Énée} de Pietro da Cortona au palais Pamphilj, 1651-1654), elle exalte la continuité de la Rome antique et moderne, chrétienne et païenne. Canuti, associé à Enrico Haffner pour la \teiHi[rend={italic}]{quadratura}, renoue avec le faste et l’illusionnisme de Pietro da Cortona ou de Valesio. Ce n’est cependant pas le blason familial qui est en suspension et en gloire, car ses « meubles » – six étoiles « en champ d’azur » extraites de l’écu – sont soutenus et élevés par des \teiHi[rend={italic}]{putti} qui restent appuyés latéralement sur l’architecture feinte de la composition et sont subordonnés au motif central (fig. 171). L\teiHi[rend={italic}]{’élévation} des meubles dans une zone intermédiaire entre terre et ciel anticipe et accompagne donc, mais ne se confond pas, avec l’\teiHi[rend={italic}]{apothéose} du héros dynastique. L’apothéose familiale de la dynastie romaine associée à Clément X passe ainsi avant tout par l’ascension céleste d’un ancêtre mythique du peuple romain auquel pouvait prétendre s’identifier l’antique famille des Altieri, s’autorisant en particulier de la symbolique stellaire qui était un élément commun aux Altieri et aux prestigieux fondateurs romains (Romulus et Remus sont en effet désignés par Ovide comme « \teiHi[rend={italic}]{stelle congiunte} », \teiHi[rend={italic}]{Fastes}, chant II, 459 et 471-472\teiNote[n={17},place={foot},xmlid={ftn59-9}]{\teiP[xmlid={p-490}]{Le lion des Paluzzi degli Albertoni, associé à des devises, est encore présent à quatre reprises.}}).}
              \teiP[xmlid={p17-9}]{L’extraction, la diffusion dans le décor et la monumentalisation des meubles de l’écu s’accompagnent par ailleurs de leur association avec des éléments complémentaires qui, tout comme dans le cas de Pietro da Cortona, en enrichissent les connotations : divers types de couronnes (or, roseaux, chêne, herbe, laurier) entourent les différentes étoiles, des figures historiques sont en contrebas (au centre des grands côtés : \teiHi[rend={italic}]{Remus et Romulus allaités par la louve} ; \teiHi[rend={italic}]{Romulus apparaissant à Julius Proculus}) et des allégories prennent place également à leurs bases (les quatre éléments dans les angles de la composition\teiNote[n={18},place={foot},xmlid={ftn60-9}]{\teiP[xmlid={p-491}]{Sur cet ensemble, voir Angela Cipriani, « Un programma belloriano », dans Franco Borsi, Gabriele Morolli, Maurizio Caperna \teiHi[rend={italic}]{et al}., \teiHi[rend={italic}]{Palazzo Altieri}, Rome, Editalia, 1991, p. 177-208 ; Eduard A. Safarik, \teiHi[rend={italic}]{L’Apoteosi di Romolo di Domenico Maria Canuti in Palazzo Altieri}, Rome, Finnat Euramerica, 1995 ; texte repris dans « L’Apoteosi di Romolo », dans Daniela Di Castro, Gabriele Reina et Eduard A. Safarik, \teiHi[rend={italic}]{Roma Palazzo Altieri. Le stanze al piano nobile dei cardinali Giovanni Battista e Paluzzo Altieri}, Milan, Franco Maria Ricci, 1999, p. 24-29 ; Sybille Ebert-Schifferer, « La romanità in termini bolognesi: Domenico Maria Canuti a palazzo Altieri », dans Sabine Frommel (dir.), \teiHi[rend={italic}]{Crocevia e capitale della migrazione artistica: forestieri a Bologna e bolognesi nel mondo (secolo XVII)}, Bologne, Bononia University Press, 2012, p. 327-345. La fresque, disposée dans un espace commun aux deux appartements de Paluzzo Paluzzi-Albertoni et d’Angelo et Gaspare Paluzzi-Albertoni, est décrite dans un manuscrit attribué à Giovanni Pietro Bellori, conservé à la bibliothèque du Vatican (Barb. lat. 4342, fol. 82r-86r, publié en annexe de l’étude de Angela Cipriani citée plus haut) : « \teiHi[rend={italic}]{Romolo portato in cielo da Marte, avanti la maestà di Giove per deificarsi, là dove è Venere in atto supplicante per Romolo }\lbrack{}…\rbrack{}\teiHi[rend={italic}]{ vi sono ancora alcuni amorini che portano in cielo le insegne romane }\lbrack{}dans l’axe principal, au-dessus de l’étoile couronnée des Altieri\rbrack{}\teiHi[rend={italic}]{ e le rappresentano al sommo Giove }\lbrack{}…\rbrack{}\teiHi[rend={italic}]{ vi sono alcune cose que alludono a Romolo et altre alla casa Altieri }\lbrack{}…\rbrack{} ». La fresque fait écho à celle de Pietro da Cortona pour le palais Barberini et à celle, consacrée à Énée, peinte au palais Doria Pamphilj. Comme le rappelle Sybille Ebert-Schifferer, la famille Altieri pouvait, à la différence des Barberini ou des Pamphili, se revendiquer d’une ancienne origine romaine et faire valoir les liens établis notamment par Marcantonio Altieri au \teiHi[rend={small-caps}]{xvi}\teiHi[rend={sup}]{e} siècle avec l’administration du Capitole. Le programme, attribué par Eduard A. Safarik à l’érudit cardinal Paluzzo, aurait été conçu, selon Sybille Ebert-Schifferer, avec l’aide de Bellori et du cardinal Camillo Massimi, grand connaisseur d’art, également d’antique famille romaine, chargé par le pape des antiquités romaines et peut-être à l’origine du choix du peintre.}}).}
            
\end{teiDiv}
            \begin{teiDiv}[type={section1},xmlid={div04-7}]

              \teiHead[xmlid={extension-scenique-dissemination-et-transformisme-lapotheose-aldobrandini-1596-1602-001}]{Extension scénique, dissémination et transformisme : l’apothéose Aldobrandini (1596-1602)}
              \teiP[xmlid={p18-9}]{L’originalité des réalisations du \teiHi[rend={small-caps}]{xvii}\teiHi[rend={sup}]{e} siècle ne doit pas dissimuler ce qu’elles doivent à une œuvre inaugurale sur laquelle nous nous devons d’insister : la salle Clémentine, située au second étage des appartements privés du pape, entre le débouché de la Scala Regia et la salle du Consistoire. Elle était destinée notamment, au sein de la nouvelle résidence papale initiée par Sixte Quint, aux réceptions des délégations étrangères, aux réunions du collège des cardinaux mais également, semble-t-il, à l’exposition du corps du pape défunt comme ce fut encore le cas pour Jean-Paul I et Jean-Paul II\teiNote[n={19},place={foot},xmlid={ftn61-9}]{\teiP[xmlid={p-492}]{Voir Giovanni Pietro Chattard, \teiHi[rend={italic}]{Nuova descrizione del Vaticano : O Si a Della Sacrosanta Basilica Di S. Pietro Data in Lvce}, Rome, Dalle stampe del Mainardi, 1762-1767,3. vol., t. II, chap. XX : « \teiHi[rend={italic}]{Appartamento Nobile pontificio nel’Palazzo di Clemente VIII detto di Sisto }», p. 153-160 et p. 174, pour la description où la salle est donnée comme « \teiHi[rend={italic}]{ordinaria residenza della Guardia Svizzera} », la fresque du Baptême est donnée comme celle de Constantin par saint Silvestre et non comme celui de saint Clément (p.158) ; voir de même Agostino Taja, \teiHi[rend={italic}]{Descrizione del Palazzo Apostolico Vaticano}, In Roma, Appresso Niccolò e Marco Pagliarini, Marco Pagliarini, 1750, p. 493 et 494 (reprenant Baglioni). L’exposition du corps du défunt avait lieu d’abord dans ses appartements privés, puis dans les appartements officiels (salle du Trône, salle du Perroquet ou salle Clémentine), puis dans la chapelle Sixtine et enfin dans la basilique Saint-Pierre : voir Antonio Margheriti Mastino\teiHi[rend={italic}]{, La Mort du pape}, Monaco, Liamar, 2015, p. 191-195.}}. Sa décoration est due essentiellement à deux artistes toscans, Cherubino Alberti (1553-1615) et son frère Giovanni (1588-1601), que leur virtuosité en matière de perspective avait rendu célèbres. Les travaux initiés en 1596 se prolongèrent jusqu’en 1602 et furent commandés par Clément VIII Aldobrandini (1592-1605\teiNote[n={20},place={foot},xmlid={ftn62-9}]{\teiP[xmlid={p-493}]{Voir toujours la riche monographie de Ludwig Freiherr von Pastor, \teiHi[rend={italic}]{The History of the Popes: From the Close of the Middle Ages, Drawn From the Secret Archives of the Vatican and Other Original Sources}, R. F. Kerr (éd.), Londres, Paul Kegan Trench, Trubner \& Co., 1933, vol. 33 : \teiHi[rend={italic}]{Clement VIII (1592-1605)}. Sur le décor : Stefania Macioce, \teiHi[rend={italic}]{« Undique Splendent ». Aspetti della pittura sacra nella Roma di Clemente VIII Aldobrandini (1592-1605)}, Rome, De Luca, 1990, chap. II, p. 74 et suiv. ; « Per un’introduzione al pontificato di Clemente VIII Aldobrandini (1592-1605): politica e committenza », \teiHi[rend={italic}]{Bollettino della Deputazione di Storia Patria per l’Umbria}, vol. 111, n\teiHi[rend={sup}]{o} 1-2, 2014, p. 709-726 ; Christopher L. C. Ewart Witcombe, « Cesare Ripa and the Sala Clementina », \teiHi[rend={italic}]{Journal of the Warburg and Courtauld Institutes}, n\teiHi[rend={sup}]{o} 55, 1992, p. 277-282 ; Laura Carlevaris, « La Sala Clementina: La costruzione pittorica delle pareti dallo schema compositivo alla griglia prospettica », \teiHi[rend={italic}]{Bollettino-Monumenti, Musei e Gallerie Pontificie}, n\teiHi[rend={sup}]{o} 21, 2001, p. 320-361 ; Anna Laura Carlevaris, Laura De Carlo et Daniel Di Marzio, « La sala Clementina in Vaticano tra manierismo e barocco », dans Fauzia Farneti et Deanna Lenzi (dir.), \teiHi[rend={italic}]{L’architettura dell’inganno. Quadraturismo e grande decorazione nella pittura di età barocca}, Florence, Alinea, 2004, p. 133-148 ; Elisabeth Oy-Marra, \teiHi[rend={italic}]{Profane Repräsentationskunst in Rom von Clemens VIII Aldobrandini bis Alexander VII Chigi. Studien zur Funktion und Semantik römischer Deckenfresken im höfischen Kontext}, Munich-Berlin, Deutscher Kunstverlag, 2005, chap. I : « Der Papst als spirituelles Haupt der Kirche : Die Fresken der Sala Clementina im Vatikan », p. 15-46. Antérieurement : Morton C. Abromson, « Clement VIII’s Patronage of the Brothers Alberti », \teiHi[rend={italic}]{The Art Bulletin}, vol. LX, n\teiHi[rend={sup}]{o} 3, septembre 1978, p. 531-547 ; \teiHi[rend={italic}]{Painting in Rome during the Reign of Clement VIII (1592-1607), a Documentary Study}, these de doctorat, Columbia University, 1976 ; ainsi que Kristina Herrmann Fiore, « Giovanni Albertis Kunst und Wissenschaft der Quadratur. Eine Allegorie in der Sala Clementina des Vatikan », \teiHi[rend={italic}]{Mitteilungen des Kunsthistorisches Institutes in Florens}, n\teiHi[rend={sup}]{o} 22, 1978, p. 61-84 ; Christopher L. C. Ewart Witcombe, « A new fresco by Cherubino Alberti in the Vatican », \teiHi[rend={italic}]{Notes in the History of Art}, 1984, vol. 4, n\teiHi[rend={sup}]{o} 1, p. 12-16. Plus largement et récemment : Clare Robertson, \teiHi[rend={italic}]{The city and the visual arts under Clement VIII}, New Haven, Yale University Press, 2015 ; Guido Cornini, Anna Maria De Strobel et Maria Serlupi Crescenzi, « L’appartamento papale di rappresentanza », dans Carlo Pietrangeli (dir.), \teiHi[rend={italic}]{l Palazzo apostolico vaticano}, Florence, Nardini Editori, 1992, p. 169-185 ; la contribution de Lucia Simonato dans Francesco Buranelli (dir.), \teiHi[rend={italic}]{Vaticano Barocco. Arte, architettura e cerimoniale} \lbrack{}Cité du Vatican, Libreria Editrice Vaticana, 2012\rbrack{}, Milan, Jaca Book, 2014, chap. IV : « Il Palazzo Apostolico Vaticano » ; Romeo Panciroli, \teiHi[rend={italic}]{The audience suite of the papal apartments}, Cité du Vatican, Libreria Editrice Vaticana, 2004.}}) afin de célébrer le saint homonyme du pape régnant, le martyr saint Clément, saint du premier siècle et troisième successeur de saint Pierre.}
              \teiP[xmlid={p19-9}]{Deux grandes compositions occupent les parois des extrémités de la salle. Un vaste paysage, dû à Paul Bril, illustre \teiHi[rend={italic}]{Le Martyre} du saint (fig. 172) et fait face à une scène de \teiHi[rend={italic}]{Baptême}\teiNote[n={21},place={foot},xmlid={ftn63-9}]{\teiP[xmlid={p-494}]{La scène a été décrite à plusieurs reprises comme le baptême de Constantin : il s’agit plus vraisemblablement soit du propre baptême de Clément, soit, et le seul mentionné par Jacques de Voragine (\teiHi[rend={italic}]{La Légende dorée}, chap. CLXVII, II, 652), du baptême par Clément de l’ex-païen Sisinnius, « ami de l’empereur » Domitien, auquel Clément avait rendu la vue et l’ouïe, qu’il convertit puis baptisa avec 313 « personnes de sa famille », Sisinnius ayant par la suite converti nobles et familiers de l’empereur Nerva, le martyre du saint ayant eu lieu sous Trajan.}} (fig. 173), opposant les deux termes extrêmes de la vie chrétienne : l’entrée dans l’Église et le martyre du saint qui ouvre sur son apothéose céleste. Sur les parois latérales (fig. 174), se déploient nombre d’anges ou de \teiHi[rend={italic}]{putti} ainsi que toute une série d’allégories des vertus cardinales et théologales, étroitement inspirées de l’édition de 1593 (non encore illustrée) de l’\teiHi[rend={italic}]{Iconologia} de Cesare Ripa, traité dont nous connaissons la fortune au \teiHi[rend={small-caps}]{xvii}\teiHi[rend={sup}]{e} siècle. D’autres allégories (dont la \teiHi[rend={italic}]{Religion} et la \teiHi[rend={italic}]{Charité} sur l’axe médian, la \teiHi[rend={italic}]{Clémence} et la \teiHi[rend={italic}]{Justice} sur l’autre axe, l’\teiHi[rend={italic}]{Abondance} et la \teiHi[rend={italic}]{Bénignité} ou \teiHi[rend={italic}]{Bienfaisance}), se dressent encore sur la balustrade feinte du registre supérieur.}
              \teiP[xmlid={p20-9}]{La voûte (fig. 175) est occupée, au centre, par saint Clément agenouillé sur des nuées, dressé vers la Trinité lumineuse, et environné d’anges dont l’un soutient l’ancre avec laquelle il avait été précipité au sein des flots de la mer Noire. Aux deux extrémités, un autre ensemble d’anges soutient dans un savant désordre les emblèmes institutionnels de la papauté. D’un côté, les clefs et l’\teiHi[rend={italic}]{ombrellino}, insignes de l’Église romaine et de sa puissance temporelle conférée éternellement à l’Église, même en l’absence du pape\teiNote[n={22},place={foot},xmlid={ftn64-8}]{\teiP[xmlid={p-495}]{Bruno Bernard Heim\teiHi[rend={italic}]{, Coutumes et droits héraldiques de l’Église} \lbrack{}1949\rbrack{}, Paris, Beauchesne, 2000, p. 67 rappelle que cet insigne, à la mort du pape, passe, tout comme les clefs, dans le timbre du cardinal camerlingue chargé d’administrer les droits temporels de l’Église pendant la vacance du Saint-Siège au nom du Sacré Collège. Les clefs en tant qu’attributs de l’Église deviennent celles de la papauté qui se confond avec l’Église tout entière : voir Édouard Bouyé, « Les armoiries pontificales à la fin du \teiHi[rend={small-caps}]{xiii}\teiHi[rend={sup}]{e} siècle. Construction d’une campagne de communication », \teiHi[rend={italic}]{Médiévales}, n\teiHi[rend={sup}]{o} 44, 2003, p. 177 et 178, http://medievales.revues.org/938 ; de même Martine Boiteux, « La vacance du siège pontifical. De la mort et des funérailles à l’investiture du Pape : les rites de l’époque moderne », dans José Pedro Paiva (dir.), \teiHi[rend={italic}]{Religious Ceremonials and Images}. \teiHi[rend={italic}]{Power and social meaning (1400-1750)}, Coimbra, Palimage Editores, 2002, p. 103-141, ici p. 110 : la \teiHi[rend={italic}]{potestas} papale, lors d’un interrègne, est assurée par un « État éphémère » doté de son blason particulier « par l’assemblage du baldaquin au-dessus des clefs entrecroisées ».}}. Ils sont placés au-dessus du siège pontifical et de la vive mais provisoire présence de celui qui incarne par sa propre personne cette puissance. À l’opposé, associée à l’étoile (mais non à la bande dentée ou \teiHi[rend={italic}]{rastrelli} des Aldobrandini), s’élève la tiare rayonnante sur le monde (\teiHi[rend={italic}]{signum imperii} depuis Innocent III), que reçoit le nouveau pape lors de son couronnement en substitution de la mitre. Cette scène décisive de la vie du nouveau pontife est par exemple représentée par Pietro Bernini au sommet du tombeau du pontife dans la chapelle Pauline de Santa Maria Maggiore\teiNote[n={23},place={foot},xmlid={ftn65-8}]{\teiP[xmlid={p-496}]{Martine Boiteux, art. cité, p. 130 et 131 sur, une fois l’élection faite et antérieurement au \teiHi[rend={italic}]{Possesso}, le couronnement public du pape qui s’effectue devant la façade de Saint-Pierre.}}. Modèle fondateur mythique de la papauté, signes universels de la puissance de l’Église et incarnation individuelle de ce pouvoir sont ainsi conjoints.}
              \teiP[xmlid={p21-9}]{Le thème central de la voûte est celui de l’apothéose d’un saint déjà présent « dans l’Empyrée, heureux possesseur de sa Gloire\teiNote[n={24},place={foot},xmlid={ftn66-7}]{\teiP[xmlid={p-497}]{Voir la description inédite de Michelangelo Lualdi Romano, ms. Bib. Corsininiana cod. 275 : « Memorie istoriche, e curiose del Tempio e Palazzo Vaticano », vol. III, p. 23-25, cité par Kristina Hermann-Fiore, « La salle clémentine au Palais apostolique du Vatican », dans Jean-Marc Olivesi (dir.), \teiHi[rend={italic}]{Les cieux en gloire. Paradis en trompe-l’œil pour la Rome baroque}, catalogue d’exposition, Ajaccio, Musée Fesch, 2002, p. 101-115.}} ». Quand, dans nombre de cas, l’élévation et la suspension de l’armoirie ne semblaient pas avoir de terme ultime, sinon le fait même d’être élevé et en suspens, le point visé est bien le dernier ciel de l’empyrée céleste, explicitant la dimension eschatologique qui est aussi vraisemblablement celle des grandes compositions postérieures de Valesio, Pietro da Cortona ou Canuti. L’exaltation du martyr Clément sert, par homonymie, à la gloire du pape régnant comme le faisait également, à l’égard d’Urbain VIII, le tableau attribué à Charles Mellin et présent dans les collections de Francesco Barberini : \teiHi[rend={italic}]{L’Apothéose de saint Urbain }(vers 1630, Rome, Galleria Nazionale di Arte Antica, Palazzo Barberini), pape entre 222 et 230. Le ciel est le lieu auquel peut également prétendre un Clément VIII dont toutes les vertus suprêmes, rassemblées sous la Providence divine, semblent attester l’exemplarité chrétienne et anticiper la propre future apothéose : l’exposition semi-publique des corps des papes défunts en ce lieu, sous l’apothéose même de saint Clément, présentait un suggestif raccourci de l’itinéraire espéré.}
              \teiP[xmlid={p22-9}]{Par extension, le martyr Clément offre le modèle d’une gloire céleste non seulement à Clément VIII, mais à toute la dynastie des Aldobrandini, comme l’attestent, sur les parois, l’élévation et la suspension des emblèmes pontificaux et familiaux du pape et, sur la voûte, l’élévation angélique et le couronnement de l’étoile par la tiare. Dans un autre registre historique, \teiHi[rend={italic}]{L’Apothéose de saint Clément }pourrait jouer, relativement cette fois aux Aldobrandini, le rôle qu’eut plus tard \teiHi[rend={italic}]{L’Apothéose de Romulus} à l’égard de Clément X Altieri et de sa lignée dans l’antichambre du palais familial décorée par Canuti, où nous retrouvons le même meuble héraldique stellaire. Le thème fut repris au \teiHi[rend={small-caps}]{xviii}\teiHi[rend={sup}]{e} siècle sous Clément XI Albani (1700-1721) qui célébra à son tour, avec un large déploiement héraldique, \teiHi[rend={italic}]{L’Apothéose de saint Clément} au plafond de l’église Santo Clemente qui en conserve les reliques\teiNote[n={25},place={foot},xmlid={ftn67-7}]{\teiP[xmlid={p-498}]{C’est \teiHi[rend={italic}]{L’Apothéose de saint Clément} peinte par Giuseppe Chiari en 1715.}}. L’association est celle que l’on trouve encore dans les grands plafonds à caissons des églises romaines déjà signalées, où l’exaltation des armes pontificales (Borghese, Barberini, etc.) s’accorde à plusieurs reprises aux scènes d’apothéose de saints éminents (Crisogono, Cosme et Damien) ou à celles, également ascensionnelles, de l’Assomption et de l’Ascension. Tel est par exemple le cas à Santa Maria in Trastevere, où le plafond à caissons du Dominiquin intègre une \teiHi[rend={italic}]{Assomption} répondant à l’écu dédoublé du cardinal Pietro Aldobrandini, qui en fut le commanditaire en 1615. De la même façon, \teiHi[rend={italic}]{L’Ascension du Christ} due au Cavalier d’Arpin, qui domine l’autel du transept gauche de San Giovanni in Laterano, fait écho à l’écu monumental de Clément VIII sur le plafond (1597-1601\teiNote[n={26},place={foot},xmlid={ftn68-7}]{\teiP[xmlid={p-499}]{Rappelons que la sacristie des chanoines de la basilique (dite également \teiHi[rend={italic}]{sala Clementina}) honore également saint Clément par un décor de nouveau confié à Giovanni et Cherubino Alberti, associés à Agostino Ciampelli (\teiHi[rend={italic}]{circa} 1600).}}).}
            
\end{teiDiv}
            \begin{teiDiv}[type={section1},xmlid={div05-4}]

              \teiHead[xmlid={poetique-001}]{Poétique}
              \teiP[xmlid={p23-9}]{Nous retrouvons dans ce dispositif complexe, outre l’élévation et la mise en gloire des armes pontificales (clefs, \teiHi[rend={italic}]{ombrellino}, tiare, étoile des Adobrandini), l’ensemble quasi complet des nombreuses opérations formelles et rhétoriques qui ont organisé par la suite l’usage des armoiries dans les grands décors du \teiHi[rend={small-caps}]{xvii}\teiHi[rend={sup}]{e} siècle romain :}
              \begin{teiList}[rendition={\#list-ndash},type={unordered},xmlid={list1}]

                \teiItem{la monumentalisation, par amplification, des emblèmes de la papauté et des meubles de l’écu personnel et familial ;}
                \teiItem{le soutien et l’élévation angélique de ces éléments où, variante du \teiHi[rend={italic}]{Stemma in arrivo} (dans l’espace terrestre), se substitue celui du \teiHi[rend={italic}]{Stemma in partenza} (vers le ciel)\teiHi[rend={italic}]{. }Les anges \teiHi[rend={italic}]{reggi-stemma }associés au commanditaire et à son statut font ici échos aux anges psychopompes (\teiHi[rend={italic}]{reggi-anima}) porteurs de l’âme de saint Clément ;}
                \teiItem{l’extraction et la dissémination des meubles dans un contexte inédit, devenus ainsi indépendants de l’armoirie et de son écu. Cherubino Alberti en a repris le principe en 1604 pour la plus modeste échelle de la voûte de la chapelle Aldobrandini à Santa Maria sopra Minerva ;}
                \teiItem{une forme généralisée d’animation et d’activation des objets qui trouve, sur la voûte, son expression ultime dans le couronnement de l’étoile Aldobrandini par la tiare papale ;}
                \teiItem{l’incorporation spatiale des éléments armoriaux au sein d’un nouveau et élargi contexte scénique, narratif et théologique ;}
                \teiItem{le motif de « l’armoirie transformiste » où les armes familiales sont soumises à d’inédites manipulations.}
              
\end{teiList}
              \teiP[xmlid={p24-7}]{Insistons sur deux opérations essentielles dans ce décor : la tension entre dissémination et intégration unitaire des meubles ou des autres éléments de l’armoirie ; la dynamique « transformiste » ou « métamorphique » de l’armoirie.}
              \teiP[xmlid={p25-7}]{Relevons tout d’abord que emblèmes et meubles héraldiques, dispersés principalement sur les parties hautes, tendent à se rassembler ailleurs : au centre du pavement de la salle, où est présent l’écu familial complet timbré de la tiare et des clefs en sautoir, et sur le registre inférieur des parois latérales de marbre entre les fenêtres et les portes (l’écu avec tiare et clefs, ou bien une seule étoile surmontée de la tiare et des clefs). L’écu complet est également disposé sous une forme monumentale au niveau supérieur de l’architecture fictive de la composition, dans un espace intermédiaire entre ciel et terre qui évoque la solution reprise plus tard par Canuti au palais Altieri. Dans l’un des angles de la voûte (côté trône) deux figures – qui évoquent les \teiHi[rend={italic}]{ignudi} de Michel-Ange – soutiennent l’écu Aldobrandini reconstitué (avec étoiles et bande dentée), un \teiHi[rend={italic}]{putto} ou un ange élevant la tiare en gloire accompagnée des habituelles clefs en sautoir. Le même motif est présent sur l’angle diagonal opposé (côté entrée) (fig. 176) mais accompagné cette fois d’une reprise extériorisée, amplifiée et transformée de la bande dentée devenue un surprenant fragment de cylindre : il forme une sorte de couronne glissée sous l’entablement du fronton qu’il semble supporter et surélever. Le meuble inscrit dans l’écu est ainsi devenu un élément autonome associé à une fonction à la fois référentielle (l’objet est le support de l’inscription avec le nom du pape régnant), structurelle ou tectonique (il contribue à porter l’édifice), et symbolique : soutenir ce qui peut sans doute s’assimiler à l’\teiHi[rend={italic}]{Ecclesia}, situation qui était également celle, sous Jules III del Monte, du \teiHi[rend={italic}]{trimonzio} placé sous l’édifice symbolique de l’Église dans certaines représentations\teiNote[n={27},place={foot},xmlid={ftn69-7}]{\teiP[xmlid={p-500}]{Sur ce point, voir l’article d’Yvan Loskoutoff dans le présent volume, « L’héraldique de Jules III Ciocchi del Monte (1550-1555) dans l’ornement pour le livre », fig. 16.}}.}
              \teiP[xmlid={p26-7}]{La fragmentation/dispersion des meubles armoriaux est ailleurs portée à son comble. Si l’étoile est presque partout présente, de nombreuses figures, anges ou génies placés en amortissement sur la balustrade feinte ou dans les ouvertures des parois latérales, brandissent encore l’étoile ou la bande dentée dans de nouveaux contextes signifiants. D’autres angelots, face-à-face sur les parois latérales, portent également en suspension la tiare ou les clefs de saint Pierre. La figure rhétorique employée est celle de la synecdoque, où un élément isolé renvoie à une unité complète. Étoile, bande dentée, tiare et clefs dispersées ou associées de façon partielle évoquent le blason que le spectateur est amené à reconstituer imaginairement en recomposant les éléments épars : l’écu familial monumental au centre du pavement avec la tiare et les clefs pourrait constituer à cet égard le lieu matriciel où se rassemblent, ou bien d’où se sont dispersées, les unités éparses ailleurs dans l’espace\teiNote[n={28},place={foot},xmlid={ftn70-5}]{\teiP[xmlid={p-501}]{Dans l’espace multidirectionnel de la salle, l’ordre de lecture et d’interprétation des lieux n’est plus celui auquel oblige la linéarité et la surface du texte et permet des associations ouvertes mais néanmoins réglées et contraignantes.}}. Diffuse ou ramassée sous différentes configurations, la présence pontificale est ainsi démultipliée et donnée comme une puissance agissante qui investit toute la salle et exige l’implication perceptive et cognitive des destinataires dont la propre activité relationnelle doit se déployer dans l’ensemble de l’espace.}
              \teiP[xmlid={p27-7}]{Acmé de cette puissance, la seconde opération décisive qu’il faut relever dans ce grand décor est celle de l’armoirie que nous disons transformiste, dont les meubles sont appelés non plus seulement à se \teiHi[rend={italic}]{rassembler} sous la forme unitaire de l’écu et de ses ornements obligés, mais à se \teiHi[rend={italic}]{réagencer et métamorphoser} en une autre forme. Virtuose est ici le motif, présent sur une des parois latérales (à droite du trône pontifical) (fig. 177), où les meubles de l’écu – étoiles et « \teiHi[rend={italic}]{banda dentata} » – sont recomposés pour former de nouvelles et inédites figures. Deux personnages, l’un vêtu d’une peau animale qui ferait référence au \teiHi[rend={italic}]{Labeur} (\teiHi[rend={italic}]{Labor}) manuel du peintre (tenant les pinceaux d’une main et un livre au-dessous) et l’autre, avec règle et compas qui évoquent la \teiHi[rend={italic}]{Perspective} selon Ripa\teiNote[n={29},place={foot},xmlid={ftn71-5}]{\teiP[xmlid={p-502}]{Jean Baudoin, \teiHi[rend={italic}]{Iconologie ou les principales choses qui peuvent tomber dans la pensée touchant les Vices et les Vertus}..., Paris, Aux amateurs de livres, 1643, t. I, p. 148 et figure p. 146 rapportant la chaîne d’or « où pend un œil en lieu de joyau », la règle, équerre, fil à plomb, miroir (absent ici), le livre mais tenu par l’autre figure.}}, se trouvent à proximité de deux sphères issues du réagencement des meubles Aldobrandini. Ces sphères font allusion aux sciences de la perspective et des mathématiques maîtrisées par nos artistes\teiNote[n={30},place={foot},xmlid={ftn72-3}]{\teiP[xmlid={p-503}]{\teiHi[rend={italic}]{Ibid}., p. 109, figure p. 107, avec présence de la sphère et tracé d’une étoile.}}. L’ensemble constitue également une forme d’autoportrait ou de signature des peintres, conjuguant pratique et théorie de l’art tout en témoignant encore de leurs investissements spirituels\teiNote[n={31},place={foot},xmlid={ftn73-3}]{\teiP[xmlid={p-504}]{Le personnage en contrebas pourrait être un autoportrait de Giovanni Alberti, associant dans ses vêtements les attributs de son protecteur personnel, saint Jean-Baptiste, et de saint Luc (la tête cornée), le chien renvoyant à la Fidélité.}}.}
              \teiP[xmlid={p28-6}]{Les sphères renvoient en outre aux globes armillaires de la science astronomique. Stefania Macioce, à la suite de Kristina Herrmann Fiore, y voit une allégorie sacrée où la Vérité, découverte grâce à l’association du travail et de la science perspective qui scrute le ciel, accède à une connaissance supérieure. Celle-ci est éclairée par la lumière céleste et placée sous la protection pontificale à laquelle ferait allusion la tiare présente au sommet de la composition, associée au \teiHi[rend={italic}]{motto} « \teiHi[rend={small-caps italic}]{Undique Sple\lbrack{}n\rbrack{}dent} » \lbrack{}il resplendit partout\rbrack{}. Sans doute faut-il établir une relation entre le ciel physique auquel s’attachent les efforts contraints de la science humaine et celui, métaphysique (l’empyrée), auquel seule l’âme sauvée accède au niveau de la voûte, grâce aux autres qualités et vertus acquises.}
              \teiP[xmlid={p29-6}]{Le meuble héraldique renaît ainsi et se déploie dans l’univers de l’emblématique, de la devise\teiNote[n={32},place={foot},xmlid={ftn74-3}]{\teiP[xmlid={p-505}]{Les nombreux cas, attestés dès le \teiHi[rend={italic}]{Dialogo delle imprese militari e amorose }de Paolo Giovio en 1551, des devises issues de la reprise de meubles armoriaux.}} et de l’allégorie, assumant une fonction spirituelle et également poétique ou métalinguistique, où il se prend lui-même pour objet\teiNote[n={33},place={foot},xmlid={ftn75-3}]{\teiP[xmlid={p-506}]{Sur ces échanges voir par exemple Anne Rolet, « Aux sources de l’emblème : blasons et devises », \teiHi[rend={italic}]{Littérature}, vol. 145, n\teiHi[rend={sup}]{o} 1, 2007, p. 53-78.}}.}
            
\end{teiDiv}
            \begin{teiDiv}[type={section1},xmlid={div06-4}]

              \teiHead[xmlid={de-la-scene-au-milieu-de-linterpretation-a-la-subjectivation-001}]{De la scène au milieu, de l’interprétation à la subjectivation}
              \teiP[xmlid={p30-5}]{L’entreprise monumentale de la salle Clémentine constitue sans doute à Rome le plus ambitieux exemple de déploiement armorial de l’époque moderne. L’espace architectural est entièrement investi et absorbé par le pouvoir pontifical qui en marque chaque lieu de façon différenciée, impliquant une organisation à la fois axialisée, polarisée et hiérarchisée de tout l’espace, susceptible d’effets multiples sur le spectateur.}
              \teiP[xmlid={p31-5}]{Le point focal déterminant (et le point de fuite de l’architecture de la paroi principale) est bien entendu le trône papal. Il est disposé en contrebas et au centre du portique feint représenté sur la paroi arrière. Le déploiement de la fresque établit un véritable \teiHi[rend={italic}]{frons scenae} théâtral qui, simultanément, encadre, détache et permet d’exalter la présence pontificale en s’étalant largement sur les angles des parois latérales pour mettre en évidence toute la partie antérieure de la salle. Le siège papal se situe également au croisement de l’axe horizontal (communautaire) du volume rectangulaire et de l’axe vertical (transcendant) dont la continuité est rendue sensible par les percements successifs du double portique juxtaposé et la superposition de motifs éminents : \teiHi[rend={italic}]{ombrellino} et clefs / Saint-Esprit rayonnant (dédoublé sur l’axe transversal pour illuminer également la Justice) / allégorie de la \teiHi[rend={italic}]{Religion} (qui reprend aussi les attributs de l’\teiHi[rend={italic}]{Autorité}\teiNote[n={34},place={foot},xmlid={ftn76-3}]{\teiP[xmlid={p-507}]{Jean Baudoin, \teiHi[rend={italic}]{op. cit}., p. 24 et figure p. 22 : les clefs renvoyant à la puissance « notamment spirituelle » issue de Dieu, le sceptre étant « une marque de la puissance temporelle », les Livres (ici Évangiles et Tables de la Loi) étant « un signe expres de l’Authorité des Escritures ».}}) / baptême du futur martyr recevant l’eau et sur lequel des anges déversent des fleurs / trône papal. L’autre point focal et complémentaire est précisément le motif central de l’écu familial formé par le dallage du sol. Il se trouve à l’aplomb exact de l’\teiHi[rend={italic}]{Apothéose de saint Clément}, ultime destin eschatologique visé également par Clément VIII au terme de sa vie terrestre, et aspiration commune des membres de la dynastie Aldobrandini\teiNote[n={35},place={foot},xmlid={ftn77-3}]{\teiP[xmlid={p-508}]{La disposition en ce lieu, si tel était le cas, du corps du pape défunt irait dans le sens de cette démonstration symbolique. }} et de tous les chrétiens.}
              \teiP[xmlid={p32-5}]{Comment comprendre l’espace ainsi produit ? Maîtrisé par les artistes, « l’art de peindre la scène » (\teiHi[rend={italic}]{scaeno-graphia}) est ce qui accomplit la mise en forme et la « mise en sens » de l’espace architectural\teiNote[n={36},place={foot},xmlid={ftn78-3}]{\teiP[xmlid={p-509}]{Laurent Hablot, « Le décor emblématique chez les princes de la fin du Moyen Âge : un outil pour construire et qualifier l’espace », dans Thomas Lienhard (dir.), \teiHi[rend={italic}]{Construction de l’espace au Moyen Âge : pratiques et représentations}, actes du 37\teiHi[rend={sup}]{e} congrès de la Société des historiens médiévistes de l’enseignement supérieur public, Paris, Publications de la Sorbonne, 2007, p. 147-165, article de référence pour notre sujet.}}, de ses acteurs et plus généralement de ses actants au sein desquels s’insèrent les objets et notamment, aux côtés des nombreuses inscriptions, les signes héraldiques, emblématiques et symboliques. En résulte, plus largement, une « mise en scène » qu’il faut entendre, dans le contexte de la salle Clémentine, au sens premier du terme. C’est-à-dire comme l’établissement d’une \teiHi[rend={italic}]{skènè} où se manifestent dans leur actualité phénoménale les dieux et ceux – êtres et choses – qui prétendent les représenter sous différentes formes cumulatives d’instanciation présentielle : corps réel et singulier du pape sur son trône, corps modèle d’identification projective idéale qu’est saint Clément, corps collectif dynastique qu’évoquent les armoiries familiales, corps statutaire trans-individuel que signifient les emblèmes pontificaux.}
              \teiP[xmlid={p33-5}]{En cette scène, les êtres représentés et les personnes réelles, les acteurs et les spectateurs devenus à leur tour acteurs, mais aussi les choses représentant ces êtres, paraissent se confondre dans un même espace. L’objet héraldique joue ici, de même que les autres images, un rôle déterminant. Transformé en chose animée et agissante, notamment par la mobilité et les actions dont il est alors fréquemment investi, il tend à devenir ce que nous pourrions désigner comme un « presqu’être ». Insigne (\teiHi[rend={italic}]{insignis}) ou ensemble d’insignes, il ne se contente pas, comme le fait un simple signe (\teiHi[rend={italic}]{signum}) ou un signal, de renvoyer en termes sémiotiques à son référent et de \teiHi[rend={italic}]{représenter} l’individu ou une communauté. Mais il peut tendre effectivement, comme certaines images du pouvoir politique ou religieux\teiNote[n={37},place={foot},xmlid={ftn79-3}]{\teiP[xmlid={p-510}]{Démonstration classique de Louis Marin, \teiHi[rend={italic}]{Le portrait du roi} \lbrack{}1981\rbrack{}, Paris, Éditions de Minuit, 2001, « Introduction », p. 13 : « ses signes \lbrack{}les signes iconiques de la présence royale\rbrack{} sont la réalité royale, l’être et la substance du prince, cette croyance est nécessairement exigée par les signes eux-mêmes… ».}}, à \teiHi[rend={italic}]{s’identifier} à ceux qu’il représente, voire même à les \teiHi[rend={italic}]{investir }de la puissance et de la dignité qu’ils revendiquaient. Sans doute est-ce cette qualité présentielle et vitaliste qui explique notamment, évoquée par Michel Pastoureau, la « terreur héraldique » des révolutionnaires à l’égard des « signes » dits de « féodalité », en grande partie analogue à l’acharnement des iconoclastes du \teiHi[rend={small-caps}]{xvi}\teiHi[rend={sup}]{e} siècle vis-à-vis d’images également suspectées d’entretenir la confusion idolâtre entre représentation et divinité.}
              \teiP[xmlid={p34-4}]{Une autre caractéristique éminente du lieu commun où s’articulent ces différents actants est qu’il s’agit d’un espace englobant. Bien que polarisé, hiérarchisé et déterminant une relation frontale vis-à-vis du siège papal, il est presque infiniment dilaté dans toutes ses directions grâce à l’usage de la perspective. Un tel espace ne se réduit pas pour autant à ses coordonnées géométriques abstraites ni à un illusionnisme spatial qui en feraient un simple contenant pour divers contenus. Il est une structure relationnelle, un « milieu » (le \teiHi[rend={italic}]{medium }aristotélicien), un « environnement », peut-être une « ambiance » (\teiHi[rend={italic}]{Stimmung}) ou une « atmosphère » dont la description qu’en livre Émile Zola et que nous citons en exemple est un précieux bien que tardif témoignage\teiNote[n={38},place={foot},xmlid={ftn80-4}]{\teiP[xmlid={p-511}]{Émile Zola, \teiHi[rend={italic}]{op. cit}, p. 600.}}. Ce milieu est celui où importent autant les coordonnées géométriques et les signes qui investissent l’espace (relevant des sphères théologiques, politiques et sociales), que les couleurs (fondues en « une vague apparition fauve »), les reflets, la lumière et le clair-obscur (les « reflets de joyaux et de pierreries », les « demi-ténèbres »), faisant de l’espace interne de la salle un lieu indéterminé en accord avec l’ouverture céleste de la voûte. Le lieu est également indissociable des actions qui s’y déroulent et actualisent ses potentialités\teiNote[n={39},place={foot},xmlid={ftn81-4}]{\teiP[xmlid={p-512}]{Voir par exemple la catégorisation de Gernot Böhme, \teiHi[rend={italic}]{Aisthetique. Pour une esthétique de l’expérience sensible }\lbrack{}2001\rbrack{}\teiHi[rend={italic}]{, }Martin Kaltenecker et Franck Lemonde (trad.), Dijon, Les Presses du réel, 2020, p. 112, qui distingue cinq « classes de caractères atmosphériques » : les caractères sociaux, les synesthésies, les humeurs, les caractères communicationnels, les impressions de mouvement.}}. L’ensemble des actants s’inscrit en effet au sein de la représentation d’un récit, ponctué d’étapes et d’épreuves intégratrices (conversion et baptême, vie vertueuse, martyre, apothéose) qui forment autant d’événements se rapportant à l’économie du salut. Le récit, distribué sur tous les axes de l’espace-milieu de la salle Clémentine, est celui que les spectateurs-acteurs étaient eux-mêmes amenés à partager et à revivre pour leur propre salvation, par leurs positions, leurs perceptions, leurs trajectoires et leurs projections imaginaires. En leur absence, que constate le héros d’Émile Zola qui traverse nuitamment des salles désertées, le lieu redevient « une solitude élargie » énigmatique dépourvue de signification.}
              \teiP[xmlid={p35-4}]{Une telle conception de l’espace n’est pas sans conséquences sur les usagers comme sur nos propres pratiques interprétatives.}
              \teiP[xmlid={p36-4}]{Au-delà des fonctions référentielles et cognitives de ce décor et du déchiffrement érudit presqu’infini qu’il peut communément susciter (lecture iconographique, sémiotique analytique ou synthétique, interprétations exégétiques et herméneutiques), la mise en forme et la mise en scène génèrent des effets peut-être plus essentiels encore sur le spectateur/acteur. Effets à la fois phatiques (contacts), expressifs (sensations synesthésiques, affects, sentiments, \teiHi[rend={italic}]{pathos}), injonctifs, conatifs et pragmatiques (obligations, résolutions). Les \teiHi[rend={italic}]{dispositifs} formels (l’antique \teiHi[rend={italic}]{dispositio}) impliquent notamment des \teiHi[rend={italic}]{dispositions }corporelles et intérieures, c’est-à-dire des « tonalités affectives » au sens de Gernot Böhme\teiNote[n={40},place={foot},xmlid={ftn82-4}]{\teiP[xmlid={p-513}]{Gernot Böhme, \teiHi[rend={italic}]{op. cit}., chap. V. La tonalité renvoie au terme traditionnel en allemand de \teiHi[rend={italic}]{Stimmung.}}} qui adviennent lors des « expériences d’insertion » que vivent les usagers de ces espaces\teiNote[n={41},place={foot},xmlid={ftn83-4}]{\teiP[xmlid={p-514}]{\teiHi[rend={italic}]{Ibid}., p. 57.}}. L’attention méditative, l’admiration, la surprise, le saisissement, la fascination, le désir d’imitation, l’exaltation, l’impression d’une intégration immersive dans un univers ascensionnel dont le chrétien est supposé être partie intégrante au sein de l’\teiHi[rend={italic}]{Ecclesia}, tous ces sentiments peuvent céder la place à des impressions dysphoriques : sentiments de subordination, voire d’assujettissement, d’écrasement ou d’égarement au sein d’un univers saturé d’objets, de signes et d’obligations cérémonielles. Ces dernières peuvent dès lors être conçues non seulement comme des « marques de la protestation de foy, que nous faisons à Dieu », mais aussi comme un « tesmoignage de notre franche \teiHi[rend={italic}]{servitude}\teiNote[n={42},place={foot},xmlid={ftn84-3}]{\teiP[xmlid={p-515}]{Claude Vavre, \teiHi[rend={italic}]{La cour de Rome la saincte. Ou traicté des ceremonies et coustumes qui s’observent dans la ville de Rome}, Paris, Nicolas Buon, 1623, p. 134, sur les cérémonies conçues comme « marques de la protestation de foy, que nous faisons à Dieu, \& tesmoignage de notre franche servitude, au meilleur maistre de tout le monde, \& de l’ordre des creatures inferieures à leur Createur superieur » : à l’inacessibilité d’un Dieu invisible répond « l’exterieur des Ceremonies sacrees » visibles.}} » à l’égard du Créateur et de ses représentants. Le lieu peut ainsi susciter, à l’image du héros d’Émile Zola dont la traversée déceptive des appartements privés du pape va signer l’échec de son ambition réformatrice, rejet et sentiment de discordance (la « discrépance » au sens de nouveau de Gernot Böhme).}
              \teiP[xmlid={p37-4}]{L’espace-milieu et ses objets – où les éléments armoriaux jouent un rôle déterminant dans la salle Clémentine comme dans la plupart des grands décors baroques postérieurs – impliquent ainsi l’assignation à des positions obligées et l’incorporation de postures de soumission et de vénération où le sujet percevant est largement \teiHi[rend={italic}]{construit} par l’objet perçu : ce que Michel Foucault désignait comme un processus de subjectivation qui peut être une forme d’aliénation. L’ordre spirituel évoqué – celui du ciel et de ses représentants ici-bas, où la hiérarchie céleste verticale qui culmine avec saint Clément trouve son répondant dans la hiérarchie terrestre horizontale dominée par le corps du pape régnant – oblige à l’intégration d’un ordre politique et social fondé sur l’exaltation d’un modèle de monarchie universelle d’institution divine, qu’attestent les anciens cérémoniaux curiaux\teiNote[n={43},place={foot},xmlid={ftn85-3}]{\teiP[xmlid={p-516}]{Rappelons que Clément VIII est à l’origine du premier cérémonial publié à la suite du Concile de Trente en 1600 : le \teiHi[rend={italic}]{Cæremoniale episcoporum jussu Clementis VIII }(édition parisienne en 1633, avec gravures). Toute une littérature livre la hiérarchie et les rites de la cour romaine, voir notamment Claude Vavre, \teiHi[rend={italic}]{op. cit.} ; Frédéric Tantouche, \teiHi[rend={italic}]{Traicté des coustumes et ceremonies qui s’observent à Rome, \& ailleurs. Tant par nostre S. Père que par Messeigneurs les Eminentissimes \& Reverendissimes Cardinaux…}, Paris, Fleury Bouriquant, 1631, notamment p. 88 sur les « Audiences du Pape » avec indication des différentes hiérarchies à respecter ; Francesco Sestini da Bibbiena\teiHi[rend={italic}]{, Il maestro di camera \lbrack{}…\rbrack{} Di nuove ricorretto secondo il Ceremoniale Romano}, Viterbo, Per Il Diotallevi, 1639 (sur les cardinaux) ; Girolamo Lunadoro, \teiHi[rend={italic}]{Relatione della corte di Roma, e de’riti da osservarsi in essa, e de’suoi magistrati, \& offitij, con la loro distinta giurisdittione}, Bracciano, Andrea Fei, 1650, p. 86-88 sur la réception des ambassadeurs et des « \teiHi[rend={italic}]{Dame all’Audienza} » (mais sans indication des lieux concernés) ; Gregorio Leti, \teiHi[rend={italic}]{Itinerario della corte di Roma ò vero teatro historico, cronologico, e Politico}…, Valenza, Pietro Francesco Guerini, 1675, 3 vol., t. 2, p. 481-484, sur les réceptions publiques du pape (mais toujours sans indications de lieu) ; ainsi que Bernard Picart, \teiHi[rend={italic}]{Cérémonies et coutumes religieuses de tous les peuples du monde}, Amsterdam, J. F. Bernard, 1723, 8 vol., t. II, « Supplément à ce qui concerne la Hiérarchie de l’Église ; où l’on donne un détail abrégé du spirituel de la Cour du Pape », p. 165-197 ; pour la période médiévale voir la somme de Michel Andrieu, \teiHi[rend={italic}]{Le pontifical romain au Moyen Âge}, Cité du Vatican, BAV, 1938-1941. Le \teiHi[rend={italic}]{Diario di Roma} mentionne la salle Clémentine à plusieurs occasions au \teiHi[rend={small-caps}]{xix}\teiHi[rend={sup}]{e} siècle pour la cérémonie du lavement des « \teiHi[rend={italic}]{piedi ai tredici Sacerdoti di diverses nazioni} » (29 mars 1823, p. 3, de même le 17 avril 1824, p. 3, etc.), mais cette cérémonie, concernant les « pauvres », se faisait dans la salle Ducale antérieurement (voir Frédéric Tantouche, \teiHi[rend={italic}]{op. cit}., p. 64). L’étude de référence sur les relations de l’architecture palatiale (mais sans intégrer les résidences papales) et du cérémonial est celle de Patricia Waddy, \teiHi[rend={italic}]{Seventeenth-Century Roman Palaces. Use and the Art of the Plan}, Cambridge-Londres, The MIT Press, 1990 ; sur le palais apostolique, mais avant notre période, voir Tristan Weddigen, Sible de Blaauw et Bram Kempers (éd.), \teiHi[rend={italic}]{Functions and Decorations : Art and Ritual at the Vatican Palace in the Middle Ages and the Renaissance}, Cité du Vatican-Turnhout, BAV-Brepols, 2003, notamment la contribution de Bram Kempers, « Ritual and its images. Paris de Grassis, Raphaël and the “signatures” in the Vatican Stanze », p. 71-93. Si manquent les publications, une recherche systématique sur les usages de cette salle serait à mener dans les témoignages anciens et, au Vatican, dans l’\teiHi[rend={italic}]{Archivio delle celebrazioni liturgiche del Sommo Pontefice}. Je remercie Martine Boiteux pour ces conseils.}}. Les usages actuels en démontrent encore la force contraignante : pape en position centrale et élevée, membres éminents de la Curie à ses côtés, cardinaux disposés en rangs latéralement ou représentants extérieurs placés en face et en retrait du pape.}
              \teiP[xmlid={p38-4}]{L’ordre symbolique auquel participe de façon décisive l’économie spatiale des armoiries détermine, en retour, une « mise en ordre » sociale et statutaire qui, si elle n’est pas immuable, reste toujours intensément prégnante.}
              \begin{teiFigure}[xmlid={figure01-12}]

                \teiGraphic[url={../icono/br/Ch13\_Cousinie/Fig168.jpg}]
                \teiHead[xmlid={figure-164-giovanni-luigi-valesio-lapotheose-des-armes-ludovisi-c-1623-fresque-rome-casino-ludovisi-voute-de-la-salle-dellarmadio-de-libri-ou-della-stemma-001}]{Figure 164. Giovanni Luigi Valesio, \teiHi[rend={italic}]{L’Apothéose des armes Ludovisi, }c. 1623, fresque, Rome, Casino Ludovisi, voûte de la salle « \teiHi[rend={italic}]{dell’armadio de’ libri} » ou « \teiHi[rend={italic}]{della Stemma} ».}
              
\end{teiFigure}
              \begin{teiFigure}[xmlid={figure02-12}]

                \teiGraphic[url={../icono/br/Ch13\_Cousinie/Fig169.jpg}]
                \teiHead[xmlid={figure-165-raphael-angelots-supportant-lecu-papal-1508-fresque-cite-du-vatican-voute-de-la-chambre-de-la-signature-001}]{Figure 165. Raphaël, \teiHi[rend={italic}]{Angelots supportant l’écu papal}, 1508, fresque, Cité du Vatican, voûte de la chambre de la Signature.}
              
\end{teiFigure}
              \begin{teiFigure}[xmlid={figure03-12}]

                \teiGraphic[url={../icono/br/Ch13\_Cousinie/Fig170.jpg}]
                \teiHead[xmlid={figure-166-pietro-da-cortona-voute-du-grand-salon-du-palais-barberini-1632-1638-fresque-rome-detail-001}]{Figure 166. Pietro da Cortona, \teiHi[rend={italic}]{Voûte du grand salon du palais Barberini}, 1632-1638, fresque, Rome (détail).}
              
\end{teiFigure}
              \begin{teiFigure}[xmlid={figure04-12}]

                \teiGraphic[url={../icono/br/Ch13\_Cousinie/Fig171.jpg}]
                \teiHead[xmlid={figure-167-pietro-da-cortona-voute-de-l-anticamera-de-cortigiani-ou-salle-de-mars-1644-1646-fresque-florence-palais-pitti-001}]{Figure 167. Pietro da Cortona, \teiHi[rend={italic}]{Voûte de l’« Anticamera de’ cortigiani » ou salle de Mars}, 1644-1646, fresque, Florence, palais Pitti.}
              
\end{teiFigure}
              \begin{teiFigure}[xmlid={figure05-12}]

                \teiGraphic[url={../icono/br/Ch13\_Cousinie/Fig172.jpg}]
                \teiHead[xmlid={figure-168-antonio-domenico-gabbiani-lapotheose-de-la-famille-corsini-c-1696-huile-sur-toile-modele-pour-la-fresque-du-salon-du-palais-corsini-al-parione-florence-001}]{Figure 168. Antonio Domenico Gabbiani, \teiHi[rend={italic}]{L’Apothéose de la famille Corsini}, c. 1696, huile sur toile, modèle pour la fresque du salon du palais Corsini al Parione, Florence.}
              
\end{teiFigure}
              \begin{teiFigure}[xmlid={figure06-12}]

                \teiGraphic[url={../icono/br/Ch13\_Cousinie/Fig173.jpg}]
                \teiHead[xmlid={figure-169-pietro-da-cortona-decor-de-la-voute-de-la-petite-galerie-du-palais-barberini-1632-fresque-rome-001}]{Figure 169. Pietro da Cortona, \teiHi[rend={italic}]{Décor de la voûte de la petite galerie du palais Barberini}, 1632, fresque, Rome.}
              
\end{teiFigure}
              \begin{teiFigure}[xmlid={figure07-11}]

                \teiGraphic[url={../icono/br/Ch13\_Cousinie/Fig174.JPG}]
                \teiHead[xmlid={figure-170-domenico-maria-canuti-lapotheose-de-romulus-c-1675-1676-fresque-rome-palais-altieri-voute-de-lantichambre-001}]{Figure 170. Domenico Maria Canuti, \teiHi[rend={italic}]{L’Apothéose de Romulus}, c. 1675-1676, fresque, Rome, palais Altieri, voûte de l’antichambre.}
              
\end{teiFigure}
              \begin{teiFigure}[xmlid={figure08-9}]

                \teiGraphic[url={../icono/br/Ch13\_Cousinie/Fig175.JPG}]
                \teiHead[xmlid={figure-171-domenico-maria-canuti-lapotheose-de-romulus-c-1675-1676-fresque-detail-dun-des-angles-de-la-voute-avec-letoile-des-altieri-001}]{Figure 171. Domenico Maria Canuti, \teiHi[rend={italic}]{L’Apothéose de Romulus}, c. 1675-1676, fresque, détail d’un des angles de la voûte avec l’étoile des Altieri.}
              
\end{teiFigure}
              \begin{teiFigure}[xmlid={figure09-8}]

                \teiGraphic[url={../icono/br/Ch13\_Cousinie/Fig176.jpg}]
                \teiHead[xmlid={figure-172-paul-bril-le-martyre-de-saint-clement-c-1696-1602-fresque-cite-du-vatican-salle-clementine-001}]{Figure 172. Paul Bril, \teiHi[rend={italic}]{Le Martyre de saint Clément}, c. 1696-1602, fresque, Cité du Vatican, salle Clémentine.}
              
\end{teiFigure}
              \begin{teiFigure}[xmlid={figure10-7}]

                \teiGraphic[url={../icono/br/Ch13\_Cousinie/Fig177.JPG}]
                \teiHead[xmlid={figure-173-cherubino-et-giovanni-alberti-la-sala-clementina-c-1696-1602-cite-du-vatican-001}]{Figure 173. Cherubino et Giovanni Alberti, \teiHi[rend={italic}]{La Sala Clementina}, c. 1696-1602, Cité du Vatican.}
              
\end{teiFigure}
              \begin{teiFigure}[xmlid={figure11-7}]

                \teiGraphic[url={../icono/br/Ch13\_Cousinie/Fig178.jpg}]
                \teiHead[xmlid={figure-174-cherubino-et-giovanni-alberti-anges-putti-et-allegories-c-1696-1602-fresque-cite-du-vatican-salle-clementine-detail-des-parois-laterales-001}]{Figure 174. Cherubino et Giovanni Alberti, \teiHi[rend={italic}]{Anges, putti et allégories}, c. 1696-1602, fresque, Cité du Vatican, salle Clémentine, détail des parois latérales.}
              
\end{teiFigure}
              \begin{teiFigure}[xmlid={figure12-7}]

                \teiGraphic[url={../icono/br/Ch13\_Cousinie/Fig179.JPG}]
                \teiHead[xmlid={figure-175-cherubino-et-giovanni-alberti-lapotheose-de-saint-clement-c-1696-1602-fresque-cite-du-vatican-salle-clementine-voute-001}]{Figure 175. Cherubino et Giovanni Alberti, \teiHi[rend={italic}]{L’Apothéose de saint Clément}, c. 1696-1602, fresque, Cité du Vatican, salle Clémentine, voûte.}
              
\end{teiFigure}
              \begin{teiFigure}[xmlid={figure13-6}]

                \teiGraphic[url={../icono/br/Ch13\_Cousinie/Fig180.JPG}]
                \teiHead[xmlid={figure-176-cherubino-et-giovanni-alberti-lecu-aldobrandini-c-1696-1602-fresque-cite-du-vatican-salle-clementine-detail-dun-des-angles-de-la-voute-001}]{Figure 176. Cherubino et Giovanni Alberti, \teiHi[rend={italic}]{L’écu Aldobrandini}, c. 1696-1602, fresque, Cité du Vatican, salle Clémentine, détail d’un des angles de la voûte.}
              
\end{teiFigure}
              \begin{teiFigure}[xmlid={figure14-6}]

                \teiGraphic[url={../icono/br/Ch13\_Cousinie/Fig181.jpg}]
                \teiHead[xmlid={figure-177-cherubino-et-giovanni-alberti-allegorie-c-1696-1602-fresque-cite-du-vatican-salle-clementine-detail-des-parois-laterales-001}]{Figure 177. Cherubino et Giovanni Alberti, \teiHi[rend={italic}]{Allégorie, }c. 1696-1602, fresque, Cité du Vatican, salle Clémentine, détail des parois latérales.}
              
\end{teiFigure}
            
\end{teiDiv}
          
\end{teiElement}
        
\end{teiElement}
      
\end{teiElement}
      \begin{teiElement}[name={group},type={bibliography},data-page-title={Sources et bibliographie},data-include-href={Ch14\_Source\_bibliographie.xml},xmlid={bibliography-001}]

        \begin{teiElement}[name={front}]

          \begin{teiDiv}[type={titlePage},xmlid={titlepage-016}]

            \teiP[rend={title-main},xmlid={p-517}]{Sources et bibliographie}
          
\end{teiDiv}
        
\end{teiElement}
        \begin{teiElement}[name={body},xmlid={body-016}]

          \begin{teiDiv}[type={section1},xmlid={div01-13}]

            \teiHead[xmlid={manuscrits-et-documents-darchives-001}]{Manuscrits et documents d’archives}
            \begin{teiDiv}[type={section2},xmlid={div02-13}]

              \teiHead[xmlid={avignon-001}]{Avignon}
              \begin{teiBibl}[xmlid={bibl-009}]
Archives municipales : \teiHi[rend={small-caps}]{CC468} (Comptes de Gaspard Droin, dépenses extraordinaires, non numéroté).
\end{teiBibl}
              \begin{teiBibl}[xmlid={bibl-010}]
Bibliothèque municipale : ms. 2384 (Recueil de statuts et de privilèges de la ville d’Avignon) ; ms. 5712 (Pierre Pansier, \teiHi[rend={italic}]{Les peintres d’Avignon au }\teiHi[rend={small-caps italic}]{xvi}\teiHi[rend={italic sup}]{e}\teiHi[rend={italic}]{ siècle. Biographies et documents}).
\end{teiBibl}
            
\end{teiDiv}
            \begin{teiDiv}[type={section2},xmlid={div03-11}]

              \teiHead[xmlid={bologne-001}]{Bologne}
              \begin{teiBibl}[xmlid={bibl-011}]
Archivio di Stato : \teiHi[rend={italic}]{Insignia degli Anziani consoli}, vol. \teiHi[rend={small-caps}]{VIII}.
\end{teiBibl}
            
\end{teiDiv}
            \begin{teiDiv}[type={section2},xmlid={div04-8}]

              \teiHead[xmlid={cahors-001}]{Cahors}
              \begin{teiBibl}[xmlid={bibl-012}]
Archives départementales du Lot : F 365 (donation du monastère de Beaulieu-sur-Dordogne par Hugues de Castelnau, 1076).
\end{teiBibl}
            
\end{teiDiv}
            \begin{teiDiv}[type={section2},xmlid={div05-5}]

              \teiHead[xmlid={cite-du-vatican-001}]{Cité du Vatican}
              \begin{teiBibl}[xmlid={bibl-013}]
Archivio Apostolico Vaticano : Armadio \teiHi[rend={small-caps}]{XXXIX}.59, fol. 267 et 444 (concessions héraldiques de Jules \teiHi[rend={small-caps}]{III}) ; Misc. Arm. \teiHi[rend={small-caps}]{II} 150 (relation sur les familles romaines, début \teiHi[rend={small-caps}]{xvii}\teiHi[rend={sup}]{e} siècle).
\end{teiBibl}
              \begin{teiBibl}[xmlid={bibl-014}]
BAV : Barb. gr. 265 (ms. aux armes de Marcello Cervini, cardinal bibliothécaire) ; Barb. lat. 4342, fol. 82-86 (description de fresque attribuée à Pietro Bellori) ; Cappella Sistina 67 (\teiHi[rend={italic}]{Officium ad Tertiam}) ; Cappella Sistina 154 (missel) ; Cappella Sistina 213 (lectionnaire) ; Chigiani A.\teiHi[rend={small-caps}]{II}.29, fol. 45 (lettre de Chigi à Malvezzi, 1644) ; Ottoboniani latini 3105 (album Pier Leone Ghezzi) ; Vat. lat. 317 (saint Jérôme, \teiHi[rend={italic}]{Exposition sur les Psaumes}) ; Vat. lat. 3561, fol. 8 (discours d’Andrea de Monte) ; Vat. lat. 3790 (Léon \teiHi[rend={small-caps}]{IX}, \teiHi[rend={italic}]{Lettres}) ; Vat. lat. 3841 (Hildebert de Tours, \teiHi[rend={italic}]{Lettres}) ; Vat. lat. 3965, fol. 29r-v (paiement du cardinal bibliothécaire Marcello Cervini) ; Vat. lat. 3967-3969 (\teiHi[rend={italic}]{Registres des dépenses} de la Bibliothèque Vaticane) ; Vat. lat. 8253 (Gigli, recueil d’épitaphes) ; Vat. lat. 8254 \teiHi[rend={small-caps}]{II} (sur Gauges de’ Gozze) ; Vat. lat. 11258 (papiers de Virgilio Spada) ; Vat. lat. 13295 (dessins de Filippo Juvarra).
\end{teiBibl}
            
\end{teiDiv}
            \begin{teiDiv}[type={section2},xmlid={div06-5}]

              \teiHead[xmlid={florence-001}]{Florence}
              \begin{teiBibl}[xmlid={bibl-015}]
Biblioteca Medicea Laurenziana : Plutei 41.3 (\teiHi[rend={italic}]{Fior di virtù) }; 76.63 (\teiHi[rend={italic}]{Fior di virtù}) ; Tempiano, n\teiHi[rend={sup}]{o} 3 (vues de Sienne et de Florence).
\end{teiBibl}
              \begin{teiBibl}[xmlid={bibl-016}]
Biblioteca Riccardiana : ms. 1761 (\teiHi[rend={italic}]{Fior di virtù}) ; ms. 1763 ; ms. 1774 (\teiHi[rend={italic}]{Fior di virtù}).
\end{teiBibl}
            
\end{teiDiv}
            \begin{teiDiv}[type={section2},xmlid={div07-4}]

              \teiHead[xmlid={harvard-001}]{Harvard}
              \begin{teiBibl}[xmlid={bibl-017}]
Houghton Library : ms. Typ 136 (pontifical).
\end{teiBibl}
            
\end{teiDiv}
            \begin{teiDiv}[type={section2},xmlid={div08-4}]

              \teiHead[xmlid={milan-001}]{Milan}
              \begin{teiBibl}[xmlid={bibl-018}]
Bibliothèque Ambrosienne : ms. F 227 inf. (Giacomo Grimaldi, \teiHi[rend={italic}]{Instrumenta translationum}).
\end{teiBibl}
            
\end{teiDiv}
            \begin{teiDiv}[type={section2},xmlid={div09-4}]

              \teiHead[xmlid={new-york-001}]{New York}
              \begin{teiBibl}[xmlid={bibl-019}]
New York Public Library : ms. Spencer 50 (Ésope des Médicis).
\end{teiBibl}
              \begin{teiBibl}[xmlid={bibl-020}]
Pierpont Morgan Library : ms. H.6 (\teiHi[rend={italic}]{Praeparatio ad Missam Pontificalem}).
\end{teiBibl}
              \begin{teiBibl}[xmlid={bibl-021}]
The Metropolitan Museum of Art : \teiHi[rend={italic}]{Coat of arms of Clement XIII (Rezzonico)}, Accession Number 52.570.155.
\end{teiBibl}
            
\end{teiDiv}
            \begin{teiDiv}[type={section2},xmlid={div10-4}]

              \teiHead[xmlid={paris-001}]{Paris}
              \begin{teiBibl}[xmlid={bibl-022}]
Archives nationales : ms. Fr. 23154, fol. 273 (bulle de fondation de la collégiale Saint-Louis de Castelnau-Bretenoux)
\end{teiBibl}
              \begin{teiBibl}[xmlid={bibl-023}]
BNF : Fr. 9542 (lettre de Gozze à Peiresc) ; It. 364 (Gauges de’ Gozze, \teiHi[rend={italic}]{Dell’antichità delle armi, o insegne delle famiglie}) ; Latin 8971 (Jean Bœuf, \teiHi[rend={italic}]{Frontispice}, vers 1598).
\end{teiBibl}
            
\end{teiDiv}
            \begin{teiDiv}[type={section2},xmlid={div11-3}]

              \teiHead[xmlid={rome-002}]{Rome}
              \begin{teiBibl}[xmlid={bibl-024}]
Archivio della Compagnia di Gesù : Fondo Gesuitico, Fabbrica di Sant’Ignazio, vol. 1343-1344 (décor de l’église Saint-Ignace).
\end{teiBibl}
              \begin{teiBibl}[xmlid={bibl-025}]
Archivio di Stato di Roma : Cartari Febei, reg. 164, f. 122 (blasonnement des armes de Jules \teiHi[rend={small-caps}]{III}).
\end{teiBibl}
              \begin{teiBibl}[xmlid={bibl-026}]
ADP : Fonds Archiviolo, n\teiHi[rend={sup}]{o} 145 (blasonnement des armes Pamphili) ; scaff. 93, busta 54, interno 2 (\teiHi[rend={italic}]{Riflessioni intorno alla chiarezza dell’Eccellentissima Casa Panphilia dall’anno 850 fino al 1300}) ; 86.2, 88.34.1 et 88.34.2 (paiements).
\end{teiBibl}
              \begin{teiBibl}[xmlid={bibl-027}]
Biblioteca Angelica : ms. 928 (Mariano Cavense, \teiHi[rend={italic}]{Triumphus Montium}, 1551).
\end{teiBibl}
              \begin{teiBibl}[xmlid={bibl-028}]
Biblioteca Casanatense : ms. 1335.
\end{teiBibl}
              \begin{teiBibl}[xmlid={bibl-029}]
Biblioteca Corsiniana : ms. 275 (Michelangelo Lualdi Romano, \teiHi[rend={italic}]{Memorie istoriche, e curiose del Tempio e Palazzo Vaticano}).
\end{teiBibl}
            
\end{teiDiv}
            \begin{teiDiv}[type={section2},xmlid={div12-3}]

              \teiHead[xmlid={tolede-001}]{Tolède}
              \begin{teiBibl}[xmlid={bibl-030}]
Bibliothèque capitulaire : Évangéliaire de Jules \teiHi[rend={small-caps}]{III} ; Missel du cardinal Girolamo Dandini.
\end{teiBibl}
            
\end{teiDiv}
            \begin{teiDiv}[type={section2},xmlid={div13-2}]

              \teiHead[xmlid={vallombreuse-001}]{Vallombreuse}
              \begin{teiBibl}[xmlid={bibl-031}]
Archives de l’abbaye : pièces D.V.3 et D.V.4 (Scritture diverse pertinenti alla nostra congregazione).
\end{teiBibl}
            
\end{teiDiv}
            \begin{teiDiv}[type={section2},xmlid={div14-2}]

              \teiHead[xmlid={viterbe-001}]{Viterbe}
              \begin{teiBibl}[xmlid={bibl-032}]
Archivio comunale : Pergamene n\teiHi[rend={sup}]{o} 11 (concession héraldique impériale).
\end{teiBibl}
            
\end{teiDiv}
          
\end{teiDiv}
          \begin{teiDiv}[type={section1},xmlid={div15-2}]

            \teiHead[xmlid={livres-001}]{Livres}
            \begin{teiDiv}[type={section2},xmlid={div16}]

              \teiHead[xmlid={livres-anciens-editions-anterieures-a-1800-001}]{Livres anciens (éditions antérieures à 1800)}
              \begin{teiBibl}[xmlid={bibl-033}]
\teiHi[rend={italic}]{Acta Sanctorum Ianuarii}, Antverpiae, Apud Joannem Meursium, 1643, t. 2.
\end{teiBibl}
              \begin{teiBibl}[xmlid={bibl-034}]
\teiHi[rend={small-caps}]{Agrippa} Camillo, \teiHi[rend={italic}]{Trattato di Scientia d’Arme}, In Roma, Per Antonio Blado stampadore apostolico, 1553.
\end{teiBibl}
              \begin{teiBibl}[xmlid={bibl-035}]
\teiHi[rend={small-caps}]{Albani} Giovanni Girolamo, \teiHi[rend={italic}]{Erudita atque luculenta disputatio}, Romae, Excudebant Valerius et Aloysius Dorici fratres Brixienses, 1553.
\end{teiBibl}
              \begin{teiBibl}[xmlid={bibl-036}]
\teiHi[rend={italic}]{Annales mediolanenses ab anno 1230 usque ad annum 1402}, dans Lodovico Antonio Muratori (éd.), \teiHi[rend={italic}]{Rerum Italicarum Scriptores}, Mediolani, Typographia societatis palatinae in regiae curia, 1730, t. 16.
\end{teiBibl}
              \begin{teiBibl}[xmlid={bibl-037}]
\teiHi[rend={small-caps}]{Archinto} Filippo, \teiHi[rend={italic}]{De noua christiani orbis pace oratio}, Romae, Apud Antonium Bladum Asulanum, 1544.
\end{teiBibl}
              \begin{teiBibl}[xmlid={bibl-038}]
\teiHi[rend={small-caps}]{Armanni} Vincenzo, \teiHi[rend={italic}]{Delle Lettere del Signor V. A., colla vita dell’autore per Carlo Cartari}, Roma, Iacomo Dragondelli, 1663, t. I ; Macerata, Giuseppe Piccini, 1674, t. \teiHi[rend={small-caps}]{II} et \teiHi[rend={small-caps}]{III}.
\end{teiBibl}
              \begin{teiBibl}[xmlid={bibl-039}]
\teiHi[rend={small-caps}]{Baudoin} Jean, \teiHi[rend={italic}]{Iconologie ou les principales choses qui peuvent tomber dans la pensée touchant les Vices et les Vertus}..., Paris, Aux amateurs de livres, 1643.
\end{teiBibl}
              \begin{teiBibl}[xmlid={bibl-040}]
\teiHi[rend={small-caps}]{Bernal di Gioya} Antonio,\teiHi[rend={italic}]{ Copiosissimo discorso della fontana, e guglia eretta in Piazza Navona, per ordine della Santità di Nostro Signore Innocentio X dal Signor Cavalier Bernini : con una abondante dichiaratione delli quattro fiumi...}, Roma, Tiberius, 1651.
\end{teiBibl}
              \begin{teiBibl}[xmlid={bibl-041}]
\teiHi[rend={small-caps}]{Blégier} Antoine, \teiHi[rend={italic}]{La Magnificque et triumphante entrée de Carpentras, faicte à tresillustre \& trespuisssant prince ALEXANDRE FARNES, Cardinal, Legat d’Avignon, Vicechancelier du S. siege Apostolique}, Avignon, Par Macé Bonhomme, aux Changes, 1553.
\end{teiBibl}
              \begin{teiBibl}[xmlid={bibl-042}]
\teiHi[rend={small-caps}]{Bracciolini} Francesco, \teiHi[rend={italic}]{L’elettione di Urbano Papa VIII}, \lbrack{}Rome\rbrack{}, s. n., \lbrack{}1628\rbrack{}.
\end{teiBibl}
              \begin{teiBibl}[xmlid={bibl-043}]
\teiHi[rend={small-caps}]{Brusoni} Girolamo, \teiHi[rend={italic}]{Degli allori d’Eurota, poesie di diuersi all’eccellentiss. sig. principe D. Camillo Pamphilio, raccolte dal caualier Girolamo Brusoni, e dedicate all’eccellentissima signora principessa donna Olimpia Aldobrandina Pamphilij}, In Venetia, Per il Valuasense, 1662-1663, 2 vol.
\end{teiBibl}
              \begin{teiBibl}[xmlid={bibl-044}]
\teiHi[rend={bold}]{—}, \teiHi[rend={italic}]{Le glorie Pamphilie, oda di Girolamo Brusoni}, s. l., s. n., \lbrack{}après 1656\rbrack{}.
\end{teiBibl}
              \begin{teiBibl}[xmlid={bibl-045}]
\teiHi[rend={italic}]{Bulla confirmationis priuilegiorum militum sancti Pauli}, \lbrack{}Rome\rbrack{}, s. n., \lbrack{}v. 1550\rbrack{}.
\end{teiBibl}
              \begin{teiBibl}[xmlid={bibl-046}]
\teiHi[rend={italic}]{Bulla indulgentiarum et facultatum pro erectione Ecclesiae, \& Monasterii S. Tholemaei in Vrbae Nepesina}, \lbrack{}Rome, Antonio Blado, 1543\rbrack{}.
\end{teiBibl}
              \begin{teiBibl}[xmlid={bibl-047}]
\teiHi[rend={italic}]{Bullae diuersorum pontificum incipiente a Ioanne XXII usque ad sanctiss. d.n.d. Iulium papam III,} Romae, Apud d. Hieronymam de Cartulariis Perusinam, 1550 mense Iunii.
\end{teiBibl}
              \begin{teiBibl}[xmlid={bibl-048}]
\teiHi[rend={italic}]{Bullarium civitatis Avenionensis seu Bullae…}, Lugduni, Ex Typographia Ioannis Amati Candy, 1657.
\end{teiBibl}
              \begin{teiBibl}[xmlid={bibl-049}]
\teiHi[rend={italic}]{Cæremoniale episcoporum jussu Clementis VIII… Reformatum}, Parisiis, Societatis typographicae libr. off. eccl., 1633.
\end{teiBibl}
              \begin{teiBibl}[xmlid={bibl-050}]
\teiHi[rend={small-caps}]{Caferri} Angelo, \teiHi[rend={italic}]{Synthema Vetustatis}, Romae, Ex tipographia Iacobi Dragondelli, 1670.
\end{teiBibl}
              \begin{teiBibl}[xmlid={bibl-051}]
\teiHi[rend={small-caps}]{Casali} Giovanni Battista, \teiHi[rend={italic}]{De profanis, et sacris veteribus ritibus}, Romae, Ex typographia Andreae Phaei, 1645.
\end{teiBibl}
              \begin{teiBibl}[xmlid={bibl-052}]
\teiHi[rend={small-caps}]{Castiglione} Sabba da, \teiHi[rend={italic}]{Ricordi di monsignor Sabba da Castiglione caualier Gierosolimitano, di nuouo corretti, et ristampati. Con vna tauola copiosissima nuouamente aggiunta. Et appresso breuemente e descritta la vita dell’auttore}, In Venetia, Per Paolo Gerardo, 1560.
\end{teiBibl}
              \begin{teiBibl}[xmlid={bibl-053}]
\teiHi[rend={small-caps}]{Catharino} Ambrosio, \teiHi[rend={italic}]{Enarrationes r.p.f. Ambrosii Catharini Politi Senensis archiepiscopi Compsani in quinque priora capita libri Geneseos. Adduntur plerique alij tractatus et quaestiones…}, Romae, Apud Antonium Bladum Camerae Apostolicae typographum, 1552.
\end{teiBibl}
              \begin{teiBibl}[xmlid={bibl-054}]
\teiHi[rend={small-caps}]{Cellonese} Andrea, \teiHi[rend={italic}]{Specchio simbolico delle armi gentilizie}, Naples, Giovanni Francesco Paci, 1667.
\end{teiBibl}
              \begin{teiBibl}[xmlid={bibl-055}]
\teiHi[rend={small-caps}]{Chasseneuz} Barthélémy de, \teiHi[rend={italic}]{Catalogus gloriae mundi, in quo doctissime simul \& copiosissime de dignitatibus, honoribus, prærogatiuis, \& excellentia spirituum, hominum, animantium, rerumque cæterarum omnium, quæ cælo, mari, terra, inferno que ipso continentur, ita differitur, ut nihil . præteritum sit.}, Lugduni, Per honoratum virum Dionysium de Harsy, 1529.
\end{teiBibl}
              \begin{teiBibl}[xmlid={bibl-056}]
\teiHi[rend={small-caps}]{Chattard} Giovanni Pietro, \teiHi[rend={italic}]{Nuova descrizione del Vaticano : O Si a Della Sacrosanta Basilica Di S. Pietro Data in Lvce}, In Roma, Dalle stampe del Mainardi, 1762-1767, 3 vol.
\end{teiBibl}
              \begin{teiBibl}[xmlid={bibl-057}]
\teiHi[rend={small-caps}]{Chracas} Luca Antonio, \teiHi[rend={italic}]{Distinto raguaglio del sontuoso treno delle carrozze con cui andò all’udienza di Sua Santità, il di 8 luglio 1716, l’illustrissimo ed eccellentissimo signore Don Rodrigo Annes de Saa, Almeida e Meneses, marchese di Fontes, conte di Pennaghiano}, In Roma, Nella Stamperia di Gio. Francesco Chracas, 1716.
\end{teiBibl}
              \begin{teiBibl}[xmlid={bibl-058}]
\teiHi[rend={small-caps}]{Commeno} Andrea Angelo, \teiHi[rend={italic}]{Genealogia d’imperatori romani et constantinopolitani}, In Roma, per Valerio Dorico, \& Luigi fratelli bresciani, 1551.
\end{teiBibl}
              \begin{teiBibl}[xmlid={bibl-059}]
\teiHi[rend={italic}]{Contra Hebreos retinentes libros in quibus aliquid contra fidem catholicam notetur vel scribatur}, Romae, Apud Antonium Bladum impressorem Cameralem, \lbrack{}v. 1554\rbrack{}.
\end{teiBibl}
              \begin{teiBibl}[xmlid={bibl-060}]
\teiHi[rend={italic}]{Copia delle lettere del Serenissimo Re d’Inghilterra, \& del Reuerendissimo Card. Polo Legato della S. Sede Apostolica alla Santità di N. S. Iulio Papa III. sopra la reduttione di quel Regno alla vnione della Santa Madre Chiesa ; \& obedienza della Sede Apostolica}, \lbrack{}Rome, Valerio e Luigi Dorico, v. 1554\rbrack{}.
\end{teiBibl}
              \begin{teiBibl}[xmlid={bibl-061}]
\teiHi[rend={small-caps}]{Crescenzi} Giovanni Pietro, \teiHi[rend={italic}]{Il nobile romano, o sia trattato di nobiltà}, Bologna, Per gli eredi d’Antonio Pisarri, 1693.
\end{teiBibl}
              \begin{teiBibl}[xmlid={bibl-062}]
\teiHi[rend={small-caps}]{Crespi} Luigi, \teiHi[rend={italic}]{Felsina pittrice}, In Roma, Nella Stamperia di Marco Pagliarini, 1769.
\end{teiBibl}
              \begin{teiBibl}[xmlid={bibl-063}]
\teiHi[rend={small-caps}]{Dal Gambaro Sclarici} Lorenzo, \teiHi[rend={italic}]{De concilio liber I}, Bononiae, Apud Anselmum Giaccarelum, 1551.
\end{teiBibl}
              \begin{teiBibl}[xmlid={bibl-064}]
\teiHi[rend={small-caps}]{Danza} Prospero, \teiHi[rend={italic}]{La elettion del beatissimo \& summo pontifice Iulio tertio}, \lbrack{}Venise, stampata per Proxpero Danza, 1550\rbrack{}.
\end{teiBibl}
              \begin{teiBibl}[xmlid={bibl-065}]
\teiHi[rend={small-caps}]{Davidico} Lorenzo, \teiHi[rend={italic}]{Anotomia delli vitii}, In Firenze, s. n., 1550.
\end{teiBibl}
              \begin{teiBibl}[xmlid={bibl-066}]
\teiHi[rend={bold}]{—}, \teiHi[rend={italic}]{Il vittorioso trionfo di Maria V. Contra Luterani}, In Firenze, s. n., 1550.
\end{teiBibl}
              \begin{teiBibl}[xmlid={bibl-067}]
\teiHi[rend={italic}]{Delle antichità delle armi gentilizie, trattato di Celso Cittadini, colle annotazioni di Giovan Girolamo Carli}, Lucca, Per Salvatore e Giandomenico Marescandoli, 1741.
\end{teiBibl}
              \begin{teiBibl}[xmlid={bibl-068}]
\teiHi[rend={italic}]{Effigies cardinalium nunc viventium}, Rome, Giovanni Giacomo de Rossi, 1660-1676.
\end{teiBibl}
              \begin{teiBibl}[xmlid={bibl-069}]
\teiHi[rend={small-caps}]{Eustathius Thessalonicensis}, \teiHi[rend={italic}]{Parekbolai eis ten Homerou Iliada kai Odysseian meta euporotatou kai panu ophelimou pinakos}, Romae, Apud Antonium Bladum impressorem cameralem, 1550.
\end{teiBibl}
              \begin{teiBibl}[xmlid={bibl-070}]
\teiHi[rend={small-caps}]{Festini} Carlo, \teiHi[rend={italic}]{I trionfi della magnificenza pontificia celebrati per lo passaggio nelle città e luoghi dello Stato ecclesiatico e in Roma per lo ricevimento della Maestà della regina di Svezia}, In Roma, Nella Stamperia della Rev. Camera Apost., 1656.
\end{teiBibl}
              \begin{teiBibl}[xmlid={bibl-071}]
\teiHi[rend={small-caps}]{Foglietta} Uberto, \teiHi[rend={italic}]{In festo omnium Beatorum}, \lbrack{}Rome, Antonio Blado, v. 1553\rbrack{}.
\end{teiBibl}
              \begin{teiBibl}[xmlid={bibl-072}]
\teiHi[rend={small-caps}]{Fortunio} Agostino, \teiHi[rend={italic}]{Cronichetta del Monte San Savino di Toscana}, In Firenze, Nella Stamperia di Bartolomeo Sermartelli, 1583.
\end{teiBibl}
              \begin{teiBibl}[xmlid={bibl-073}]
\teiHi[rend={italic}]{Galleria Giustiniana del marchese Vincenzo Giustiniani}, \lbrack{}Rome\rbrack{}, s. n., \lbrack{}v. 1631\rbrack{}, 2 vol.
\end{teiBibl}
              \begin{teiBibl}[xmlid={bibl-074}]
\teiHi[rend={small-caps}]{Galletti} Pier Luigi, \teiHi[rend={italic}]{Inscriptiones Romanae infimi aevi Romae exstantes}, Romae, Typis Generosi Salomoni Bibliopolae, 1760.
\end{teiBibl}
              \begin{teiBibl}[xmlid={bibl-075}]
\teiHi[rend={small-caps}]{Georgijevic} Bartolomej, \teiHi[rend={italic}]{Libellus vere Christiana lectione dignus diuersas res Turcharum}, \lbrack{}Romae, apud Anthonium Bladum, 1552\rbrack{}.
\end{teiBibl}
              \begin{teiBibl}[xmlid={bibl-076}]
\teiHi[rend={small-caps}]{Giannotti Rangoni} Tommaso, \teiHi[rend={italic}]{De vita hominis ultra CXX. annos protrahenda}, Venetiis, s. n., 1553.
\end{teiBibl}
              \begin{teiBibl}[xmlid={bibl-077}]
\teiHi[rend={small-caps}]{Giovio} Paolo,\teiHi[rend={italic}]{ Dialogo dell’imprese militari et amorose di Monsignor Paolo Giouio Vescouo di Nucera}, In Roma, Appresso Antonio Barre, 1555.
\end{teiBibl}
              \begin{teiBibl}[xmlid={bibl-078}]
\teiHi[rend={bold}]{—}, \teiHi[rend={italic}]{Ragionamento di Mons. Paolo Giouio sopra i motti, \& disegni d’arme, \& d’amore, che communemente chiamano imprese, con vn discorso di Girolamo Ruscelli intorno allo stesso soggetto}, In Venetia, Appresso Giordano Ziletti, 1556.
\end{teiBibl}
              \begin{teiBibl}[xmlid={bibl-079}]
\teiHi[rend={small-caps}]{Gozze} Gauges de’, \teiHi[rend={italic}]{Se dalle armi, o insegne, che parlano, overo se da’ corpi delle armi, che rappresentano i cognomi, si possi argomentare ignobiltà in quella fameglia, che le usa \lbrack{}...\rbrack{}}, In Roma, Per Vitale Mascardi, 1637.
\end{teiBibl}
              \begin{teiBibl}[xmlid={bibl-080}]
\teiHi[rend={small-caps}]{Gualdo Priorato} Galeazzo, \teiHi[rend={italic}]{Historia della Sua Reale Maestà Cristina Alessandra, Regina di Svezia}, In Roma, Nella Stamperia della Rev. Camera Apost., 1656.
\end{teiBibl}
              \begin{teiBibl}[xmlid={bibl-081}]
\teiHi[rend={small-caps}]{Henry} Honoré, \teiHi[rend={italic}]{La magnifique entrée du reverendissime et très illustre seigneur, Monseigneur Alexandre, Cardinal de Farnez, Vichancelier du Sainct siege apostolique, et Legat de la ville et Cité d’Avignon, faicte en icelle le xvi mars 1553}, Avignon, s. n., 1553.
\end{teiBibl}
              \begin{teiBibl}[xmlid={bibl-082}]
\teiHi[rend={small-caps}]{Juvarra} Filippo,\teiHi[rend={italic}]{ Raccolta Di Varie Targhe di Roma Fatte Da Professori Primarj In Roma, Disegnate, ed intagliate}, Roma, Per Antonio de’ Rossi, 1711.
\end{teiBibl}
              \begin{teiBibl}[xmlid={bibl-083}]
\teiHi[rend={small-caps}]{Laguna} Andrès, \teiHi[rend={italic}]{Rimedio delle podagre}, Roma, Ad istantia di Gio. Maria Scotto, 1552.
\end{teiBibl}
              \begin{teiBibl}[xmlid={bibl-084}]
\teiHi[rend={small-caps}]{Le Laboureur} Jean, \teiHi[rend={italic}]{Relation du voyage de la Royne de Pologne et du retour de Mme la mareschalle de Guébriant, ambassadrice extraordinaire... par la Hongrie, l’Autriche, Styrie, Carinthie, le Frioul et l’Italie, avec un discours historique de toutes les villes et estats par où elle a passé et un traitté particulier du royaume de Pologne...}, Paris, Chez la veuve Jean Camusat et Pierre Le Petit, 1647.
\end{teiBibl}
              \begin{teiBibl}[xmlid={bibl-085}]
\teiHi[rend={small-caps}]{Leti} Gregorio, \teiHi[rend={italic}]{Itinerario della corte di Roma ò vero teatro historico, cronologico, e Politico}…, Valenza, Pietro Francesco Guerini, 1675, 3 vol.
\end{teiBibl}
              \begin{teiBibl}[xmlid={bibl-086}]
\teiHi[rend={small-caps}]{Lualdi} Michelangelo, \teiHi[rend={italic}]{Descrittione della fontana Pamphilia, dove fu già il cerchio Agonale}, In Roma, Nella stamperia di Francesco Moneta, 1651.
\end{teiBibl}
              \begin{teiBibl}[xmlid={bibl-087}]
\teiHi[rend={small-caps}]{Lunadoro} Girolamo, \teiHi[rend={italic}]{Relatione della corte di Roma, e de’riti da osservarsi in essa, e de’suoi magistrati, \& offitij, con la loro distinta giurisdittione}, Bracciano, Andrea Fei, 1650.
\end{teiBibl}
              \begin{teiBibl}[xmlid={bibl-088}]
\teiHi[rend={small-caps}]{Maggi} Bartolomeo, \teiHi[rend={italic}]{De vulnerum sclopetorum et bombardarum curatione}, Bononiae, Per Bartholomeum Bonardum, 1552.
\end{teiBibl}
              \begin{teiBibl}[xmlid={bibl-089}]
\teiHi[rend={small-caps}]{Magnus} Iohannes, \teiHi[rend={italic}]{Historia de omnibus Gothorum Sueonumque regibus}, Romae, \lbrack{}Giovanni Maria Viotti\rbrack{}, 1554.
\end{teiBibl}
              \begin{teiBibl}[xmlid={bibl-090}]
\teiHi[rend={small-caps}]{Marescalchi} Vincenzo Maria, \teiHi[rend={italic}]{Aulam Vidoniam picturis exornatam Alexandri VII Pont. Opt. Max. Nomini et Numini consecratam Vincentius Maria Marescalchius hisce carminibus adumbrabat}, Bologna, Typis Ferronianis, 1665.
\end{teiBibl}
              \begin{teiBibl}[xmlid={bibl-091}]
\teiHi[rend={small-caps}]{Massa} Antonio, \teiHi[rend={italic}]{De exercitatione iurisperitorum}, \lbrack{}Impressum Romae, Apud Valerium et Aloysium Doricos fratres Brixienses, v. 1550\rbrack{}.
\end{teiBibl}
              \begin{teiBibl}[xmlid={bibl-092}]
\teiHi[rend={small-caps}]{Modesti} Publio Francesco, \teiHi[rend={italic}]{Christiana pietas}, Rimini, Erasmo Virginio, 1552.
\end{teiBibl}
              \begin{teiBibl}[xmlid={bibl-093}]
\teiHi[rend={small-caps}]{Monte} Andrea de, \teiHi[rend={italic}]{Oratio habita apud S. D. N. Iulium diuina prouidentia Papam Tertium per Andream de Monte doctorem, quondam Haebreum iam vero sacro baptismate initiatum, ex hebreo in latinum versa, in qua suadetur summum Pontificem ex illustri familia montis \& illius insignia in sacris literis fuisse praenunciata}, Rome, Antonio Blado, s. d.
\end{teiBibl}
              \begin{teiBibl}[xmlid={bibl-094}]
\teiHi[rend={small-caps}]{Morales} Cristobal de,\teiHi[rend={italic}]{ Cristophori Moralis Hyspalensis Missarum liber primus}, \lbrack{}Impressum Rome, Per Valerium Doricum et Ludouicum fratres, 1544\rbrack{}.
\end{teiBibl}
              \begin{teiBibl}[xmlid={bibl-095}]
\teiHi[rend={bold}]{—}, \teiHi[rend={italic}]{Christophori Moralis Hyspalensis missarum liber secundus}, \lbrack{}Impressum Rome, Per Valerium Doricum et Ludouicum fratres, 1544\rbrack{}.
\end{teiBibl}
              \begin{teiBibl}[xmlid={bibl-096}]
\teiHi[rend={small-caps}]{Nacchiante} Giacomo, \teiHi[rend={italic}]{Enarrationes piae, doctae, et catholicae : in epistolam Pauli ad Ephesios}, Venetiis, 1553.
\end{teiBibl}
              \begin{teiBibl}[xmlid={bibl-097}]
\teiHi[rend={italic}]{Nouveau Theatre de l’Italie }\lbrack{}Pierre Mortier, 1704-1705\rbrack{}, Amsterdam, Rutgert Christophle Alberts, 1724.
\end{teiBibl}
              \begin{teiBibl}[xmlid={bibl-098}]
\teiHi[rend={italic}]{Novae regvlae cancellariae apostolicae}, Rome, s. n., 1552.
\end{teiBibl}
              \begin{teiBibl}[xmlid={bibl-099}]
\teiHi[rend={small-caps}]{Nucula} Orazio, \teiHi[rend={italic}]{Commentariorum de bello Aphrodisiensi libri quinque}, \lbrack{}Romae, Apud Valerium et Ludouicum fratres Brixienses, 1552\rbrack{}.
\end{teiBibl}
              \begin{teiBibl}[xmlid={bibl-100}]
\teiHi[rend={italic}]{Ossequii poetici nell’ammirarsi in Bologna l’Heroica risoluzione della sereniss. Principessa Caterina Farnese di monacarsi nel Ven. Monastero delle Madri carmelitane scalze di Santa Teresa in Parma}, In Bologna, Per Giacomo Monti, 1661.
\end{teiBibl}
              \begin{teiBibl}[xmlid={bibl-101}]
\teiHi[rend={small-caps}]{Paleotti} Gabriele, \teiHi[rend={italic}]{De nothis spuriisque filiis liber}, Bononiae, Apud Anselmum Giaccarellum, 1550.
\end{teiBibl}
              \begin{teiBibl}[xmlid={bibl-102}]
\teiHi[rend={small-caps}]{Palestrina} Giovanni Pierluigi da, \teiHi[rend={italic}]{Missarum liber primus}, \lbrack{}Impressum Romae, Apud Valerium Doricum \& Aloysium fratres, 1554\rbrack{}.
\end{teiBibl}
              \begin{teiBibl}[xmlid={bibl-103}]
\teiHi[rend={small-caps}]{Passeri} Giovanni Battista, \teiHi[rend={italic}]{Vite de’ pittori}, In Roma, presso Natale Barbiellini, 1772.
\end{teiBibl}
              \begin{teiBibl}[xmlid={bibl-104}]
\teiHi[rend={small-caps}]{Picart} Bernard, \teiHi[rend={italic}]{Cérémonies et coutumes religieuses de tous les peuples du monde}, Amsterdam, J.-F. Bernard, 1723-1743, 8 vol.
\end{teiBibl}
              \begin{teiBibl}[xmlid={bibl-105}]
\teiHi[rend={small-caps}]{Pietrasanta} Silvestro, \teiHi[rend={italic}]{Tesserae gentilitiae}, Romae, Typis haeredum Francisci Corbelletti, 1638.
\end{teiBibl}
              \begin{teiBibl}[xmlid={bibl-106}]
\teiHi[rend={small-caps}]{Piranesi} Giovanni Battista, \teiHi[rend={italic}]{Le antichità romane}, In Roma, Nella stamperia di Angelo Rotilj, 1756, 4 vol.
\end{teiBibl}
              \begin{teiBibl}[xmlid={bibl-107}]
\teiHi[rend={bold}]{—}, \teiHi[rend={italic}]{Lettere di giustificazione scritte a Milord Charlemont e a’ di lui agenti di Roma, dal signor Piranesi socio della real società degli antiquari di Londra intorno la dedica della sua opera delle antichità rom. fatta allo stesso signore ed vltimamente soppressa}, In Roma, s. n., 1757.
\end{teiBibl}
              \begin{teiBibl}[xmlid={bibl-108}]
\teiHi[rend={bold}]{—}, \teiHi[rend={italic}]{Della magnificenza ed architettura de’ Romani}, \lbrack{}Rome\rbrack{}, s. n., 1761.
\end{teiBibl}
              \begin{teiBibl}[xmlid={bibl-109}]
\teiHi[rend={bold}]{—}, \teiHi[rend={italic}]{Le rovine del castello dell’acqua Giulia}, \lbrack{}Rome\rbrack{}, s. n., 1761.
\end{teiBibl}
              \begin{teiBibl}[xmlid={bibl-110}]
\teiHi[rend={bold}]{—}, \teiHi[rend={italic}]{Lapides Capitolini sive fasti consularesque Romanorum}, \lbrack{}Rome\rbrack{}, s. n., 1762.
\end{teiBibl}
              \begin{teiBibl}[xmlid={bibl-111}]
\teiHi[rend={bold}]{—}, \lbrack{}\teiHi[rend={italic}]{Vedute di Roma}, Rome, s. n., 1745-1763\rbrack{}.
\end{teiBibl}
              \begin{teiBibl}[xmlid={bibl-112}]
\teiHi[rend={bold}]{—}, \teiHi[rend={italic}]{Antichità d’Albano e di Castel Gandolfo}, In Roma, s. n., 1764.
\end{teiBibl}
              \begin{teiBibl}[xmlid={bibl-113}]
\teiHi[rend={bold}]{—}, \teiHi[rend={italic}]{Diverse maniere d’adornare i cammini}, Roma, Generoso Salomoni, 1769.
\end{teiBibl}
              \begin{teiBibl}[xmlid={bibl-114}]
\teiHi[rend={bold}]{—}, \teiHi[rend={italic}]{Vasi, candelabri, cippi, sarcofagi, tripodi, lucerne ed ornamenti antichi}, s. l., s. n., 1778.
\end{teiBibl}
              \begin{teiBibl}[xmlid={bibl-115}]
\teiHi[rend={small-caps}]{Poiana} Giovanni Battista, \teiHi[rend={italic}]{De iobileo et indulgentiis}, \lbrack{}Romae, Apud Valerium, et Aloysium Doricos fratres, 1550\rbrack{}.
\end{teiBibl}
              \begin{teiBibl}[xmlid={bibl-116}]
\teiHi[rend={italic}]{Regulae tres S. D. N. D. Iulii diuina prouidentia Papae. III}, \lbrack{}Rome, Antonio Blado, v. 1550\rbrack{}.
\end{teiBibl}
              \begin{teiBibl}[xmlid={bibl-117}]
\teiHi[rend={small-caps}]{Ripa} Cesare, \teiHi[rend={italic}]{Iconologia}, In Roma, Appresso Lepido Facji, 1603.
\end{teiBibl}
              \begin{teiBibl}[xmlid={bibl-118}]
\teiHi[rend={small-caps}]{Rosichino} Mattia, \teiHi[rend={italic}]{Dichiarazione delle Pitture della Sala de’ Signori Barberini in Rome}, In Roma, Appresso Vitale Mascardi, 1640.
\end{teiBibl}
              \begin{teiBibl}[xmlid={bibl-119}]
\teiHi[rend={small-caps}]{Ruscelli} Girolamo, \teiHi[rend={italic}]{Le imprese illustri con espositioni, et discorsi del S.or Ieronimo Ruscelli}, In Venetia, \lbrack{}Damiano Zenaro\rbrack{}, 1566.
\end{teiBibl}
              \begin{teiBibl}[xmlid={bibl-120}]
\teiHi[rend={small-caps}]{Sestini Da Bibbiena} Francesco,\teiHi[rend={italic}]{ Il maestro di camera \lbrack{}…\rbrack{} Di nuove ricorretto secondo il Ceremoniale Romano}, Viterbo, Per Il Diotallevi, 1639.
\end{teiBibl}
              \begin{teiBibl}[xmlid={bibl-121}]
\teiHi[rend={small-caps}]{Sigonio} Silvestro, \teiHi[rend={italic}]{De concordia servanda}, Romae, Apud Antonium Bladum, Stephanus Nicolini Sabiensis imprimendum curauit, 1550.
\end{teiBibl}
              \begin{teiBibl}[xmlid={bibl-122}]
\teiHi[rend={small-caps}]{Taja} Agostino, \teiHi[rend={italic}]{Descrizione del Palazzo Apostolico Vaticano}, In Roma, Appresso Niccolò e Marco Pagliarini, 1750.
\end{teiBibl}
              \begin{teiBibl}[xmlid={bibl-123}]
\teiHi[rend={small-caps}]{Tantouche} Frédéric, \teiHi[rend={italic}]{Traicté des coustumes et ceremonies qui s’observent à Rome, \& ailleurs. Tant par nostre S. Père que par Messeigneurs les Eminentissimes \& Reverendissimes Cardinaux…}, Paris, Fleury Bouriquant, 1631.
\end{teiBibl}
              \begin{teiBibl}[xmlid={bibl-124}]
\teiHi[rend={small-caps}]{Teti} Girolamo, \teiHi[rend={italic}]{Aedes Barberinae ad Quirinalem}, Romæ, excudebat Mascardus, 1642.
\end{teiBibl}
              \begin{teiBibl}[xmlid={bibl-125}]
\teiHi[rend={small-caps}]{Thaumas De La Thaumassière} Gaspard, \teiHi[rend={italic}]{Histoire de Berry}, Bourges-Paris, F. Toubeau-Morel, 1689.
\end{teiBibl}
              \begin{teiBibl}[xmlid={bibl-126}]
\teiHi[rend={italic}]{Theatrum Civitatum et admirandorum Italiae}, vol. 1, \teiHi[rend={italic}]{Stato della Chiesa}, Amsterdam, s. n., 1663.
\end{teiBibl}
              \begin{teiBibl}[xmlid={bibl-127}]
\teiHi[rend={small-caps}]{Tomasi} Tommaso, \teiHi[rend={italic}]{Della esaltazione di papa Innocentio X. Lettera panegirica}, In Roma, Nella stamperia di Bernardino Tani, 1644.
\end{teiBibl}
              \begin{teiBibl}[xmlid={bibl-128}]
\teiHi[rend={small-caps}]{Ughelli} Ferdinando, \teiHi[rend={italic}]{Italia Sacra}, Romae, Apud Bernardinum Tanum, 1647.
\end{teiBibl}
              \begin{teiBibl}[xmlid={bibl-129}]
\teiHi[rend={small-caps}]{Vavre} Claude, \teiHi[rend={italic}]{La cour de Rome la saincte. Ou traicté des ceremonies et coustumes qui s’observent dans la ville de Rome}, Paris, Nicolas Buon, 1623.
\end{teiBibl}
              \begin{teiBibl}[xmlid={bibl-130}]
\teiHi[rend={italic}]{Vita del b. Lodolfo Pamphilij vescovo di Gubbio, e abate fondatore della Congreg. Dell’Avellana, detta della Colomba descritta dal p.d. Agostino Romano Fiori monaco camaldolese}, Rome, Per Antonio de’ Rossi nella strada del Seminario romano, vicino alla Rotonda, 1722.
\end{teiBibl}
              \begin{teiBibl}[xmlid={bibl-131}]
\teiHi[rend={small-caps}]{Wicquefort} Abraham de, \teiHi[rend={italic}]{Mémoire touchant les ambassadeurs et les ministres publics}, La Haye, Jean et Daniel Steucker, 1677.
\end{teiBibl}
              \begin{teiBibl}[xmlid={bibl-132}]
—, \teiHi[rend={italic}]{L’ambassadeur et ses fonctions}, La Haye, Maurice Georges Veneur, 1682, 2 vol.
\end{teiBibl}
              \begin{teiBibl}[xmlid={bibl-133}]
\teiHi[rend={small-caps}]{Zanotti} Giovan Pietro, \teiHi[rend={italic}]{Storia dell’Accademia Clementina di Bologna}, \teiHi[rend={italic}]{Giuseppe Rolli}, In Bologna, Per Lelio Dalla Volpe, 1739, vol. I.
\end{teiBibl}
            
\end{teiDiv}
            \begin{teiDiv}[type={section2},xmlid={div17}]

              \teiHead[xmlid={livres-modernes-editions-posterieures-a-1800-001}]{Livres modernes (éditions postérieures à 1800)}
              \begin{teiBibl}[xmlid={bibl-134}]
\teiHi[rend={small-caps}]{Alexander} Jonathan J. G. (dir.), \teiHi[rend={italic}]{The painted page: Italian Renaissance Book Illumination, 1450-1550}, catalogue d’exposition, Munich-New York, Prestel, 1994.
\end{teiBibl}
              \begin{teiBibl}[xmlid={bibl-135}]
\teiHi[rend={small-caps}]{Amayden} Dirk \lbrack{}Teodoro\rbrack{}, \teiHi[rend={italic}]{La Storia delle famiglie romane, con note ed aggiunte del Comm. Carlo Augusto Bertini} \lbrack{}Rome, Collegio Araldico, 1910\rbrack{}, rééd. \teiHi[rend={italic}]{facsimile} : Bologne, Forni, 1967 ; Edizioni Romane Colosseum, 1987, 2 vol.
\end{teiBibl}
              \begin{teiBibl}[xmlid={bibl-136}]
\teiHi[rend={small-caps}]{Andrieu} Michel, \teiHi[rend={italic}]{Le pontifical romain au Moyen Âge}, Cité du Vatican, BAV, 1938-1941.
\end{teiBibl}
              \begin{teiBibl}[xmlid={bibl-137}]
\teiHi[rend={small-caps}]{Armellini} Mariano, \teiHi[rend={italic}]{Le chiese di Roma, dal secolo IV al XIX}, Rome, Tipografia vaticana, 1891.
\end{teiBibl}
              \begin{teiBibl}[xmlid={bibl-138}]
— et \teiHi[rend={small-caps}]{Cecchelli} Carlo, \teiHi[rend={italic}]{Le chiese di Roma dal secolo IV al XIX}, Rome, \teiHi[rend={small-caps}]{RORE}, 1942, vol. 2.
\end{teiBibl}
              \begin{teiBibl}[xmlid={bibl-139}]
\teiHi[rend={small-caps}]{Aurigemma} Maria Giulia, \teiHi[rend={italic}]{Palazzo Firenze in Campo Marzio}, Rome, Istituto poligrafico e zecca dello Stato-Libreria dello Stato, 2007.
\end{teiBibl}
              \begin{teiBibl}[xmlid={bibl-140}]
\teiHi[rend={small-caps}]{Barrau} Hippolyte de, \teiHi[rend={italic}]{Documents historiques et généalogiques sur les familles et les hommes remarquables du Rouergue dans les temps anciens et modernes}, Rodez, Impr. de N. Ratery, 1853-1860, 4 vol.
\end{teiBibl}
              \begin{teiBibl}[xmlid={bibl-141}]
\teiHi[rend={small-caps}]{Bascapé} Giacomo Carlo et \teiHi[rend={small-caps}]{Dal Piazzo} Marcello, \teiHi[rend={italic}]{Insegne e simboli. Araldica pubblica e privata medievale e moderna}, Rome, Ministero per i beni culturali e ambientali, 1983.
\end{teiBibl}
              \begin{teiBibl}[xmlid={bibl-142}]
\teiHi[rend={small-caps}]{Baumefort} Victor de, \teiHi[rend={italic}]{Cession de la ville et de l’État d’Avignon au pape Clément VI par Jeanne I}\teiHi[rend={italic sup}]{re},\teiHi[rend={italic}]{ reine de Naples}, Apt, Typographie et lithographie J.-S. Jean, 1873.
\end{teiBibl}
              \begin{teiBibl}[xmlid={bibl-143}]
\teiHi[rend={small-caps}]{Bayard} Marc (dir.), \teiHi[rend={italic}]{Rome-Paris 1640. Transferts culturels et renaissance d’un centre artistique}, actes du colloque, Rome, villa Médicis, 17-19 avril 2008, Paris, Somogy, 2010.
\end{teiBibl}
              \begin{teiBibl}[xmlid={bibl-144}]
\teiHi[rend={small-caps}]{Beldon Scott} John, \teiHi[rend={italic}]{Images of Nepotism. The Painted Ceilings of Palazzo Barberini}, Princeton, Princeton University Press, 1991, t. 3.
\end{teiBibl}
              \begin{teiBibl}[xmlid={bibl-145}]
\teiHi[rend={small-caps}]{Benigni} Mario et \teiHi[rend={small-caps}]{zanchi} Goffredo, \teiHi[rend={italic}]{Giovanni XXIII}, Cinisello Balsamo, Edizioni San Paolo, 2000.
\end{teiBibl}
              \begin{teiBibl}[xmlid={bibl-146}]
\teiHi[rend={small-caps}]{Benocci} Carla, \teiHi[rend={italic}]{Villa Ludovisi}, Rome, Istituto poligrafico e Zecca dello Stato-Libreria dello Stato, 2010.
\end{teiBibl}
              \begin{teiBibl}[xmlid={bibl-147}]
\teiHi[rend={small-caps}]{Bevilacqua} Mario (éd.), \teiHi[rend={italic}]{Piranesi, Taccuini di Modena}, Rome, Artemide, 2008, 2  vol.
\end{teiBibl}
              \begin{teiBibl}[xmlid={bibl-148}]
— et \teiHi[rend={small-caps}]{Gallavotti Cavallero} Daniela (dir.), \teiHi[rend={italic}]{L’Aventino dal Rinascimento a oggi. Arte e Architettura}, Rome, Artemide, 2010.
\end{teiBibl}
              \begin{teiBibl}[xmlid={bibl-149}]
\teiHi[rend={italic}]{Bibliotheca Hagiographica Latina, antiquae et mediae aetatis}, Bruxelles, Société des Bollandistes, 1900-1901.
\end{teiBibl}
              \begin{teiBibl}[xmlid={bibl-150}]
\teiHi[rend={small-caps}]{Birke} Veronika (éd.), \teiHi[rend={italic}]{The Illustrated Bartsch}, New York, Abaris Book, 1982, vol. 40, part. 1 et 1987, part. 2.
\end{teiBibl}
              \begin{teiBibl}[xmlid={bibl-151}]
\teiHi[rend={small-caps}]{Bizzocchi} Roberto, \teiHi[rend={italic}]{Genealogie incredibili. Scritti di storia nell’Europa moderna}, Bologne, Il Mulino, 1995.
\end{teiBibl}
              \begin{teiBibl}[xmlid={bibl-152}]
\teiHi[rend={small-caps}]{Bjurström} Per, \teiHi[rend={italic}]{Feast and Theatre in Queen Christina’s Rome}, Stockholm, s. n., 1966.
\end{teiBibl}
              \begin{teiBibl}[xmlid={bibl-153}]
\teiHi[rend={small-caps}]{Blaya} Nelly et \teiHi[rend={small-caps}]{Bru} Nicolas, \teiHi[rend={italic}]{Le mobilier des églises du Moyen Âge dans le Lot}, Toulouse, Conseil général Midi-Pyrénées, 2014.
\end{teiBibl}
              \begin{teiBibl}[xmlid={bibl-154}]
\teiHi[rend={small-caps}]{Böhme} Gernot, \teiHi[rend={italic}]{Aisthetique. Pour une esthétique de l’expérience sensible }\lbrack{}2001\rbrack{}, Martin Kaltenecker et Franck Lemonde (trad.), Dijon, Les Presses du réel, 2020.
\end{teiBibl}
              \begin{teiBibl}[xmlid={bibl-155}]
\teiHi[rend={small-caps}]{Bohn} Babette, \teiHi[rend={italic}]{Ludovico Carracci and the art of drawing}, Turnhout, Brepols, 2004.
\end{teiBibl}
              \begin{teiBibl}[xmlid={bibl-156}]
\teiHi[rend={small-caps}]{Boisclair} Marie-Nicolas, \teiHi[rend={italic}]{Gaspard Dughet, 1615-1675}, Paris, Arthéna, 1986.
\end{teiBibl}
              \begin{teiBibl}[xmlid={bibl-157}]
\teiHi[rend={small-caps}]{Borello} Benedetta, \teiHi[rend={italic}]{L’apprentissage de Rome à la Renaissance, officiers à l’ombre de la Curie }\teiHi[rend={small-caps italic}]{xv}\teiHi[rend={italic}]{-}\teiHi[rend={small-caps italic}]{xvii}\teiHi[rend={italic sup}]{e}\teiHi[rend={italic}]{ siècles}, Turnhout, Brepols, 2021.
\end{teiBibl}
              \begin{teiBibl}[xmlid={bibl-158}]
\teiHi[rend={small-caps}]{Botana} Federico, \teiHi[rend={italic}]{Learning through Images in the Italian Renaissance: Illustrated Manuscripts and Education in Quattrocento Florence}, Cambridge, Cambridge University Press, 2020, https://doi.org/10.1017/9781108867313.004
\end{teiBibl}
              \begin{teiBibl}[xmlid={bibl-159}]
\teiHi[rend={small-caps}]{Bottino} Camilla (dir.), \teiHi[rend={italic}]{Il palazzo comunale di Bologna. Storia, architettura e restauri}, Bologne, Compositori, 1999.
\end{teiBibl}
              \begin{teiBibl}[xmlid={bibl-160}]
\teiHi[rend={small-caps}]{Bourdieu} Pierre, \teiHi[rend={italic}]{La Distinction. Critique sociale du jugement}, Paris, Éditions de Minuit, 1979.
\end{teiBibl}
              \begin{teiBibl}[xmlid={bibl-161}]
\teiHi[rend={small-caps}]{Bruel} Alexandre, \teiHi[rend={italic}]{Inventaire d’une partie des titres de famille et documents historiques de la maison de La Tour d’Auvergne conservés dans les papiers Bouillon pour faire suite aux inventaires rédigés par Baluze (2}\teiHi[rend={italic sup}]{e}\teiHi[rend={italic}]{ pt.)}, Paris, Archives nationales, 1905.
\end{teiBibl}
              \begin{teiBibl}[xmlid={bibl-162}]
\teiHi[rend={small-caps}]{Bursellis} Hieronymus de, \teiHi[rend={italic}]{Cronica gestorum ac factorum memorabilium civitatis Bononie}, A. Sorbelli (éd.), Città di Castello, S. Lapi, 1912.
\end{teiBibl}
              \begin{teiBibl}[xmlid={bibl-163}]
\teiHi[rend={small-caps}]{Campbell} Malcolm B., \teiHi[rend={italic}]{Pietro da Cortona at the Pitti Palace. A Study of the Planetary Rooms and Related Projects}, Princeton, Princeton University Press, 1977.
\end{teiBibl}
              \begin{teiBibl}[xmlid={bibl-164}]
\teiHi[rend={small-caps}]{Cancellieri} Francesco, \teiHi[rend={italic}]{Storia de’solenni possessi de’Sommi Pontefici, detti anticamente processi}, Rome, Presso Luigi Lazzarini stampatore della R. C. A., 1802.
\end{teiBibl}
              \begin{teiBibl}[xmlid={bibl-165}]
\teiHi[rend={small-caps}]{Carini} Isidoro, \teiHi[rend={italic}]{Sul titolo presbiterale di Santa Prisca}, Palerme, Tip. Camillo Tamburello, 1885.
\end{teiBibl}
              \begin{teiBibl}[xmlid={bibl-166}]
\teiHi[rend={small-caps}]{Ceresa} Massimo (dir.), \teiHi[rend={italic}]{La Biblioteca Vaticana tra riforma cattolica, crescita delle collezioni e nuovo edificio (1535-1590)}, Cité du Vatican, BAV, « Storia della Biblioteca Apostolica Vaticana, 2 », 2012.
\end{teiBibl}
              \begin{teiBibl}[xmlid={bibl-167}]
\teiHi[rend={small-caps}]{Charpenne} Pierre, \teiHi[rend={italic}]{Histoire de réunions temporaires d’Avignon et du Comtat Venaissin à la France}, Paris, Calmann-Lévy, 1886.
\end{teiBibl}
              \begin{teiBibl}[xmlid={bibl-168}]
\teiHi[rend={small-caps}]{Chastel} André, \teiHi[rend={italic}]{La grottesque}, Paris, Le Promeneur, 1988.
\end{teiBibl}
              \begin{teiBibl}[xmlid={bibl-169}]
\teiHi[rend={small-caps}]{Chossat} Marcel, \teiHi[rend={italic}]{Les Jésuites et leurs œuvres à Avignon : 1553-1768}, Avignon, François Seguin, 1896.
\end{teiBibl}
              \begin{teiBibl}[xmlid={bibl-170}]
\teiHi[rend={italic}]{Chronique métrique de Godefroy Paris}, J.-A. Buchon (éd.), Paris, Verdière, 1827.
\end{teiBibl}
              \begin{teiBibl}[xmlid={bibl-171}]
\teiHi[rend={italic}]{Clemente XIII Rezzonico, Un papa veneto nella Roma di metà Settecento}, catalogue d’exposition, Padoue, Museo Diocesiano, 12 décembre 2008-15 mars 2009, Milan, Silvana Editoriale, 2008.
\end{teiBibl}
              \begin{teiBibl}[xmlid={bibl-172}]
\teiHi[rend={italic}]{Cox}\teiHi[rend={small-caps italic}]{-}\teiHi[rend={italic}]{Rearick} Janet, \teiHi[rend={italic}]{Dynasty and Destiny in Medici Art. Pontormo, Leo X and the two Cosimos}, Princeton, Princeton University Press, 1984.
\end{teiBibl}
              \begin{teiBibl}[xmlid={bibl-173}]
\teiHi[rend={small-caps}]{Cousinié} Frédéric, \teiHi[rend={italic}]{Gloriae. Figurabilité du divin, esthétique de la lumière et dématérialisation de l’œuvre d’art à l’âge Baroque}, Rennes, \teiHi[rend={small-caps}]{PUR}, 2018.
\end{teiBibl}
              \begin{teiBibl}[xmlid={bibl-174}]
—, \teiHi[rend={small-caps}]{Lemoine} Annick, \teiHi[rend={small-caps}]{Virol} Michèle (dir.), \teiHi[rend={italic}]{Soleils baroques. La Gloire de Dieu et des Princes en représentation dans l’époque moderne}, Rome-Paris, Académie de France à Rome-Somogy, 2018.
\end{teiBibl}
              \begin{teiBibl}[xmlid={bibl-175}]
\teiHi[rend={small-caps}]{Cummings} Anthony M., \teiHi[rend={italic}]{The politicized muse: music for Medici festivals, 1512-1537}, Princeton, Princeton University Press, 1992.
\end{teiBibl}
              \begin{teiBibl}[xmlid={bibl-176}]
—, \teiHi[rend={italic}]{The Lion’s Ear: Pope Leo X, the Renaissance Papacy, and Music}, Ann Arbor, University of Michigan Press, 2012.
\end{teiBibl}
              \begin{teiBibl}[xmlid={bibl-177}]
\teiHi[rend={small-caps}]{Curzietti} Jacopo, \teiHi[rend={italic}]{Antonio Raggi: scultore ticinese nella Roma barocca}, Canterano, Aracne editrice, 2020.
\end{teiBibl}
              \begin{teiBibl}[xmlid={bibl-178}]
\teiHi[rend={small-caps}]{D’Amelio} Maria Grazia, \teiHi[rend={italic}]{L’obelisco marmoreo del foro italico a Roma. Storia, immagini e note tecniche}, Rome, Palombi \& Partner, 2009.
\end{teiBibl}
              \begin{teiBibl}[xmlid={bibl-179}]
—, \teiHi[rend={italic}]{Giovan Lorenzo Bernini e l’oro per il baldacchino di San Pietro (1624-1633)}, Rome, Argos, 2021.
\end{teiBibl}
              \begin{teiBibl}[xmlid={bibl-180}]
\teiHi[rend={small-caps}]{De Benedictis} Angela, \teiHi[rend={italic}]{Repubblica per contratto. Bologna, una città europea nello Stato della Chiesa}, Bologne, Il Mulino, 1995.
\end{teiBibl}
              \begin{teiBibl}[xmlid={bibl-181}]
\teiHi[rend={small-caps}]{Deloche} Maximin, \teiHi[rend={italic}]{Cartulaire de l’abbaye de Beaulieu (en Limousin)}, Paris, Imprimerie impériale, 1859.
\end{teiBibl}
              \begin{teiBibl}[xmlid={bibl-182}]
\teiHi[rend={small-caps}]{Desachy} Matthieu (dir.), \teiHi[rend={italic}]{L’héraldique et le livre}, catalogue d’exposition, Montauban-Albi-Toulouse, Paris, Somogy, 2002.
\end{teiBibl}
              \begin{teiBibl}[xmlid={bibl-183}]
\teiHi[rend={small-caps}]{Desplat} Christian et \teiHi[rend={small-caps}]{Mironneau} Paul (dir.), \teiHi[rend={italic}]{Les entrées : gloire et déclin d’un cérémonial}, actes du colloque de Pau, Biarritz, Société Henri \teiHi[rend={small-caps}]{IV}, 1997.
\end{teiBibl}
              \begin{teiBibl}[xmlid={bibl-184}]
\teiHi[rend={small-caps}]{Di Crollalanza} Giovanni Battista, \teiHi[rend={italic}]{Dizionario storico-blasonico delle famiglie nobili e notabili italiane estinte e fiorenti}, Pise, Presso la direzione del giornale araldico, 1886-1890, 3 vol.
\end{teiBibl}
              \begin{teiBibl}[xmlid={bibl-185}]
\teiHi[rend={small-caps}]{Di Crollalanza} Goffredo, \teiHi[rend={italic}]{Gli emblemi dei guelfi e dei ghibellini} \lbrack{}San Casciano, F. Cappelli, 1878\rbrack{}, Milan, Orsini De Marzo, 2010.
\end{teiBibl}
              \begin{teiBibl}[xmlid={bibl-186}]
\teiHi[rend={small-caps}]{Donato} Maria Monica et \teiHi[rend={small-caps}]{Parenti} Daniela (dir.), \teiHi[rend={italic}]{Dal Giglio al David, Arte civica a Firenze fra Medioevo e Rinascimento}, catalogue d’exposition, Florence, Giunti, 2013.
\end{teiBibl}
              \begin{teiBibl}[xmlid={bibl-187}]
\teiHi[rend={small-caps}]{Dykmans} Marc, \teiHi[rend={italic}]{Le cérémonial pontifical de la fin du Moyen Âge à la Renaissance}, Rome-Bruxelles, Institut historique belge de Rome, 1977-1985, 4 vol.
\end{teiBibl}
              \begin{teiBibl}[xmlid={bibl-188}]
\teiHi[rend={small-caps}]{Emiliani} Andrea, \teiHi[rend={italic}]{Leggi, bandi e provvedimenti per la tutela dei beni artistici e culturali negli antichi stati italiani 1571-1860} \lbrack{}Alfa, 1978\rbrack{}, Bologne, Nuova Alfa, 1996.
\end{teiBibl}
              \begin{teiBibl}[xmlid={bibl-189}]
\teiHi[rend={italic}]{Entrée à Rouen du roi Henri IV en 1596, précédée d’une introduction par J. Félix et de notes par Ch. de Robillard de Beaurepaire}, Rouen, Imprimerie de E. Cagniard, 1887.
\end{teiBibl}
              \begin{teiBibl}[xmlid={bibl-190}]
\teiHi[rend={small-caps}]{Fagiolo} Marcello et \teiHi[rend={small-caps}]{Madonna} Maria Luisa (dir.), \teiHi[rend={italic}]{Sisto V}, Rome, Istituto poligrafico e Zecca dello Stato, Libreria dello Stato, 1992.
\end{teiBibl}
              \begin{teiBibl}[xmlid={bibl-191}]
\teiHi[rend={small-caps}]{Fagiolo Dell’arco} Maurizio, \teiHi[rend={italic}]{Corpus delle feste a Roma}, vol.\teiHi[rend={italic}]{ }1,\teiHi[rend={italic}]{ La festa barocca}, Rome, De Luca, 1997.
\end{teiBibl}
              \begin{teiBibl}[xmlid={bibl-192}]
\teiHi[rend={small-caps}]{Favini} Vieri et \teiHi[rend={small-caps}]{Savorelli} Alessandro, \teiHi[rend={italic}]{Segni di Toscana. Identità e territorio attraverso l’araldica dei comuni: storia e invenzione grafica (secoli }\teiHi[rend={small-caps italic}]{XIII}\teiHi[rend={italic}]{-}\teiHi[rend={small-caps italic}]{XVII}\teiHi[rend={italic}]{)}, Florence, Le lettere, 2006.
\end{teiBibl}
              \begin{teiBibl}[xmlid={bibl-193}]
\teiHi[rend={small-caps}]{Favreau} Robert (éd.), \teiHi[rend={italic}]{Poitiers sous le règne de Louis XI, de 1471 à 1482. Registre de délibérations du corps de ville}, n\teiHi[rend={sup}]{o} 7, Poitiers, Société des Antiquaires de l’Ouest, 2015.
\end{teiBibl}
              \begin{teiBibl}[xmlid={bibl-194}]
\teiHi[rend={small-caps}]{Federici} Fabrizio et \teiHi[rend={small-caps}]{Garms} Jörg, « “Tombs of illustrious Italians at Rome”. L’album di disegni \teiHi[rend={small-caps}]{RCIN} 970334 della Royal Library di Windsor », \teiHi[rend={italic}]{Bollettino d’Arte}, volume spécial, 2011.
\end{teiBibl}
              \begin{teiBibl}[xmlid={bibl-195}]
\teiHi[rend={small-caps}]{Feinblatt} Ebria, \teiHi[rend={italic}]{Seventeenth-century Bolognese ceiling decorators}, Santa Barbara, Fithian Press, 1992.
\end{teiBibl}
              \begin{teiBibl}[xmlid={bibl-196}]
\teiHi[rend={small-caps}]{Ferrari} Matteo (dir.), \teiHi[rend={italic}]{L’arme segreta, Araldica e storia dell’arte nel Medioevo (secoli XIII-XV)}, Alessandro Savorelli (intr.), Florence, Le Lettere, 2015.
\end{teiBibl}
              \begin{teiBibl}[xmlid={bibl-197}]
\teiHi[rend={small-caps}]{—}, \teiHi[rend={italic}]{La politica in figure. Temi, funzioni, attori della communicazione visiva nei Comuni lombardi (XII-XIV secolo)}, Rome, Viella, 2022.
\end{teiBibl}
              \begin{teiBibl}[xmlid={bibl-198}]
\teiHi[rend={small-caps}]{Filippini} Guido et \teiHi[rend={small-caps}]{Zucchini} Francesco,\teiHi[rend={italic}]{ Miniatori e pittori a Bologna. Documenti del secolo XV}, Rome, Accademia nazionale dei Lincei, 1968.
\end{teiBibl}
              \begin{teiBibl}[xmlid={bibl-199}]
\teiHi[rend={small-caps}]{Focillon} Henri, \teiHi[rend={italic}]{G. B. Piranesi}, Paris, Librairie Renouard-Henri Laurens, 1928.
\end{teiBibl}
              \begin{teiBibl}[xmlid={bibl-200}]
\teiHi[rend={small-caps}]{Forcella} Vincenzo, \teiHi[rend={italic}]{Iscrizioni delle chiese e d’altri edificii di Roma dal secolo XI fino ai nostri giorni}, Rome, Tipografia delle scienze matematiche e fisiche, 1869-1884, 14 vol.
\end{teiBibl}
              \begin{teiBibl}[xmlid={bibl-201}]
\teiHi[rend={small-caps}]{Gaeta Bertelà} Giovanna (éd.), avec la collaboration de \teiHi[rend={small-caps}]{Stefano} Ferrara, \teiHi[rend={italic}]{Incisori bolognesi ed emiliani del sec. XVII}, Bologne, Associazione per le arti Francesco Francia, \lbrack{}1973\rbrack{}.
\end{teiBibl}
              \begin{teiBibl}[xmlid={bibl-202}]
\teiHi[rend={small-caps}]{Galbreath }Donald Lindsay, \teiHi[rend={italic}]{Papal Heraldry}, Cambridge, W. Heffer, 1930.
\end{teiBibl}
              \begin{teiBibl}[xmlid={bibl-203}]
\teiHi[rend={small-caps}]{Garms} Jörg, \teiHi[rend={small-caps}]{Juffinger} Roswitha et \teiHi[rend={small-caps}]{Ward}-\teiHi[rend={small-caps}]{Perkins} Bryan (dir.),\teiHi[rend={italic}]{ Die mittelalterlichen Grabmäler in Rom und Latium vom 13. bis zum 15. Jahrhundert}: I, \teiHi[rend={italic}]{Die Grabplatten und Tafeln}, Rome-Vienne, Österreichische Akademie der Wissenschaften, 1981-1994, 2 vol.
\end{teiBibl}
              \begin{teiBibl}[xmlid={bibl-204}]
\teiHi[rend={small-caps}]{Gauchat }Patrice (dir.), \teiHi[rend={italic}]{Hierarchia Catholica Medii et Recentioris Aevi}, Regensburg, Sumptibus et typis Librariae Monasterii, 1935, vol. V, \teiHi[rend={italic}]{1592-1667}.
\end{teiBibl}
              \begin{teiBibl}[xmlid={bibl-205}]
\teiHi[rend={small-caps}]{Ginzburg }Carlo, \teiHi[rend={italic}]{Miti, emblemi}, \teiHi[rend={italic}]{spie. Morfologia e storia}, Turin, Einaudi, 1986.
\end{teiBibl}
              \begin{teiBibl}[xmlid={bibl-206}]
\teiHi[rend={small-caps}]{Giordani} Gaetano, \teiHi[rend={italic}]{Pitture della sala Farnese}, Bologne, Tipografia Guidi all’Ancora, 1845.
\end{teiBibl}
              \begin{teiBibl}[xmlid={bibl-207}]
\teiHi[rend={small-caps}]{Gnoli} Umberto, \teiHi[rend={italic}]{Pittori e miniatori nell’Umbria}, Spolète, Argentieri, 1923.
\end{teiBibl}
              \begin{teiBibl}[xmlid={bibl-208}]
\teiHi[rend={small-caps}]{Gordon} Dillian, \teiHi[rend={italic}]{The Wilton Diptych}, Londres, The National Gallery, 2015.
\end{teiBibl}
              \begin{teiBibl}[xmlid={bibl-209}]
—, \teiHi[rend={small-caps}]{Monnas} Lisa et \teiHi[rend={small-caps}]{Elam} Caroline (dir.), \teiHi[rend={italic}]{The regal image of Richard II and the Wilton Diptych}, Londres, Harvey Miller, 1997.
\end{teiBibl}
              \begin{teiBibl}[xmlid={bibl-210}]
\teiHi[rend={small-caps}]{Gouzou} Joseph-Simon, \teiHi[rend={italic}]{Bretenoux en Haut-Quercy}, Villefranche-de-Rouergue, Impr. Salingardes, 1955.
\end{teiBibl}
              \begin{teiBibl}[xmlid={bibl-211}]
\teiHi[rend={italic}]{Grande Dizionario della lingua italiana}, Turin, \teiHi[rend={small-caps}]{UTET}, 1972, t. \teiHi[rend={small-caps}]{VI}.
\end{teiBibl}
              \begin{teiBibl}[xmlid={bibl-212}]
\teiHi[rend={small-caps}]{Gualdi} Francesco, \teiHi[rend={italic}]{Memorie sepolcrali}, Fabrizio Federici (éd.), Cité du Vatican, BAV, à paraître.
\end{teiBibl}
              \begin{teiBibl}[xmlid={bibl-213}]
\teiHi[rend={small-caps}]{Guenée} Bernard, \teiHi[rend={small-caps}]{Lehoux} Françoise, \teiHi[rend={italic}]{Les entrées royales de 1328 à 1515}, Paris, \teiHi[rend={small-caps}]{CNRS}, 1968. 
\end{teiBibl}
              \begin{teiBibl}[xmlid={bibl-214}]
\teiHi[rend={small-caps}]{Guérin} Joseph, \teiHi[rend={italic}]{Panorama d’Avignon, de Vaucluse, du Mont-Ventoux et du Col-Longet suivi de quelques vues des Alpes françaises}, Avignon, Imprimerie de Guichard aîné, 1829.
\end{teiBibl}
              \begin{teiBibl}[xmlid={bibl-215}]
\teiHi[rend={small-caps}]{Heim} Bruno Bernard,\teiHi[rend={italic}]{ Coutumes et droits héraldiques de l’Église} \lbrack{}1949\rbrack{}, Paris, Beauchesne, 2000.
\end{teiBibl}
              \begin{teiBibl}[xmlid={bibl-216}]
\teiHi[rend={small-caps}]{Herklotz} Ingo, \teiHi[rend={italic}]{Cassiano Dal Pozzo und die Archäologie des 17. Jahrhunderts}, Munich, Hirmer, 1999.
\end{teiBibl}
              \begin{teiBibl}[xmlid={bibl-217}]
\teiHi[rend={small-caps}]{Hubert} Hans W., \teiHi[rend={italic}]{Der Palazzo Comunale von Bologna} Cologne, Weimar, 1993.
\end{teiBibl}
              \begin{teiBibl}[xmlid={bibl-218}]
\teiHi[rend={small-caps}]{Jatta} Barbara (dir.), \teiHi[rend={italic}]{Piranesi e l’Aventino}, Milan, Electa, 1998.
\end{teiBibl}
              \begin{teiBibl}[xmlid={bibl-219}]
\teiHi[rend={small-caps}]{Jong} Jan de, \teiHi[rend={italic}]{Tombs in Early Modern Rome (1400-1600). Monuments of Mourning, Memory and Meditation}, Leiden, Brill, 2023.
\end{teiBibl}
              \begin{teiBibl}[xmlid={bibl-220}]
\teiHi[rend={small-caps}]{Juillet} Jacques, \teiHi[rend={italic}]{Les 38 barons de Castelnau : seconds barons chrétiens du royaume et leur baronnie, leur château, leurs églises}, s. l., Imprimerie Fabrègue, 1971.
\end{teiBibl}
              \begin{teiBibl}[xmlid={bibl-221}]
\teiHi[rend={small-caps}]{Kirsch} Johann Peter, \teiHi[rend={italic}]{Die römischen Titelkirchen im Altertum}, Paderborn, Druck und Verlag von Ferdinand Schöningh, 1918.
\end{teiBibl}
              \begin{teiBibl}[xmlid={bibl-222}]
\teiHi[rend={small-caps}]{Klein} Robert, \teiHi[rend={italic}]{La forme et l’intelligible. Écrits sur la Renaissance et l’art moderne}, Paris, Gallimard, 1970.
\end{teiBibl}
              \begin{teiBibl}[xmlid={bibl-223}]
\teiHi[rend={small-caps}]{Kybal} Vlastimil et \teiHi[rend={small-caps}]{Incisa Della Rocchetta} Giovanni (éd.), \teiHi[rend={italic}]{La nunziatura di Fabio Chigi}, Rome, Tip. Ist. Grafico Tiberino, 1943, vol. 1.
\end{teiBibl}
              \begin{teiBibl}[xmlid={bibl-224}]
\teiHi[rend={small-caps}]{Larraz} Camille, \teiHi[rend={italic}]{Simon de Châlons}, Cinisello Balsamo, Silvana Editoriale, 2022.
\end{teiBibl}
              \begin{teiBibl}[xmlid={bibl-225}]
\teiHi[rend={small-caps}]{Laurentiis} Elena de et \teiHi[rend={small-caps}]{Talamo} Emilia Anna, \teiHi[rend={italic}]{The Lost Manuscripts from the Sistine Chapel, An Epic Journey from Rome to Toledo}, catalogue d’exposition, Dallas, Meadows Museum-Southern Methodist University, 2010.
\end{teiBibl}
              \begin{teiBibl}[xmlid={bibl-226}]
\teiHi[rend={small-caps}]{Le Don} Gérard, \teiHi[rend={italic}]{Voir la sculpture. Essai sur le dispositif visuel}, Rennes, \teiHi[rend={small-caps}]{PUR}, « Aesthetica », 2018.
\end{teiBibl}
              \begin{teiBibl}[xmlid={bibl-227}]
\teiHi[rend={italic}]{Legature papali da Eugenio IV a Paolo VI}, catalogue d’exposition, Cité du Vatican, BAV, 1977.
\end{teiBibl}
              \begin{teiBibl}[xmlid={bibl-228}]
\teiHi[rend={small-caps}]{Lenzuni} Anna (dir.), \teiHi[rend={italic}]{All’ombra del lauro. Documenti librari della cultura in età laurenziana}, catalogue d’exposition, Cinisello Balsamo, Silvana Editoriale, 1992.
\end{teiBibl}
              \begin{teiBibl}[xmlid={bibl-229}]
\teiHi[rend={small-caps}]{Leone} Stéphanie, \teiHi[rend={italic}]{The Palazzo Pamphilj in Piazza Navona: Constructing Identity in Early Modern Rome}, Londres, H. Miller, 2008.
\end{teiBibl}
              \begin{teiBibl}[xmlid={bibl-230}]
\teiHi[rend={small-caps}]{—} (dir.),\teiHi[rend={italic}]{ The Pamphilj and the Arts: Patronage and Consumption in Baroque Rome}, Chestnut Hill-Chicago, McMullen Museum of Art-Boston College-University of Chicago Press, 2011.
\end{teiBibl}
              \begin{teiBibl}[xmlid={bibl-231}]
\teiHi[rend={small-caps}]{Le Pogam} Pierre-Yves, \teiHi[rend={italic}]{De la « cité de Dieu » au « palais du pape » : les résidences pontificales dans la seconde moitié du }\teiHi[rend={small-caps italic}]{xiii}\teiHi[rend={italic sup}]{e}\teiHi[rend={italic}]{ siècle (1254-1304)}, Rome, École française de Rome, 2005.
\end{teiBibl}
              \begin{teiBibl}[xmlid={bibl-232}]
\teiHi[rend={small-caps}]{Lévi}-\teiHi[rend={small-caps}]{Strauss} Claude, \teiHi[rend={italic}]{Anthropologie structurale}, Paris, Plon, 1958.
\end{teiBibl}
              \begin{teiBibl}[xmlid={bibl-233}]
\teiHi[rend={italic}]{Liber Pontificalis}, Louis Duchesne (éd.), Paris, E. de Boccard, 1956, t. 2.
\end{teiBibl}
              \begin{teiBibl}[xmlid={bibl-234}]
\teiHi[rend={small-caps}]{Loskoutoff} Yvan, \teiHi[rend={italic}]{L’armorial de Calliope, L’œuvre du Père Le Moyne S. J. (1602-1671) : littérature, héraldique, spiritualité}, Tübingen, Gunter Narr Verlag, 2000.
\end{teiBibl}
              \begin{teiBibl}[xmlid={bibl-235}]
—, \teiHi[rend={italic}]{Rome des Césars Rome des Papes. La propagande du cardinal Mazarin}, Paris, Honoré Champion, 2007.
\end{teiBibl}
              \begin{teiBibl}[xmlid={bibl-236}]
—, \teiHi[rend={italic}]{Un art de la Réforme catholique. La symbolique du pape Sixte-Quint et des Peretti-Montalto (1566-1655)}, Paris, Honoré Champion, 2011.
\end{teiBibl}
              \begin{teiBibl}[xmlid={bibl-237}]
—, \teiHi[rend={italic}]{Un art de la Réforme catholique}, vol. 2,\teiHi[rend={italic}]{ La symbolique du pape Grégoire XIII (1572-1585) et des Boncompagni}, Paris, Honoré Champion, 2018.
\end{teiBibl}
              \begin{teiBibl}[xmlid={bibl-238}]
\teiHi[rend={small-caps}]{—} (dir.), \teiHi[rend={italic}]{Héraldique et papauté. Moyen Âge-Temps modernes}, Mont-Saint-Aignan, PURH, 2020.
\end{teiBibl}
              \begin{teiBibl}[xmlid={bibl-239}]
\teiHi[rend={small-caps}]{Lucarelli} Oderigi, \teiHi[rend={italic}]{Memorie e guida storica di Gubbio}, Città di Castello, S. Lappi, 1888.
\end{teiBibl}
              \begin{teiBibl}[xmlid={bibl-240}]
\teiHi[rend={small-caps}]{Macioce} Stefania, « \teiHi[rend={italic}]{Undique Splendent ». Aspetti della pittura sacra nella Roma di Clemente VIII Aldobrandini (1592-1605)}, Rome, De Luca, 1990.
\end{teiBibl}
              \begin{teiBibl}[xmlid={bibl-241}]
\teiHi[rend={small-caps}]{Madonna} Maria Luisa (dir.), \teiHi[rend={italic}]{Roma di Sisto V. Le arti e la cultura}, Rome, De Luca, 1993.
\end{teiBibl}
              \begin{teiBibl}[xmlid={bibl-242}]
\teiHi[rend={small-caps}]{Manfredi} Tommaso, \teiHi[rend={italic}]{Filippo Juvarra. Gli anni giovanili}, Rome, Argos, 2010.
\end{teiBibl}
              \begin{teiBibl}[xmlid={bibl-243}]
\teiHi[rend={small-caps}]{Margheriti Mastino} Antonio,\teiHi[rend={italic}]{ La Mort du pape}, Monaco, Liamar, 2015.
\end{teiBibl}
              \begin{teiBibl}[xmlid={bibl-244}]
\teiHi[rend={small-caps}]{Marin} Louis, \teiHi[rend={italic}]{Le portrait du roi} \lbrack{}1981\rbrack{}, Paris, Éditions de Minuit, 2001.
\end{teiBibl}
              \begin{teiBibl}[xmlid={bibl-245}]
\teiHi[rend={small-caps}]{Marinis} Tammaro de, \teiHi[rend={italic}]{La legatura artistica in Italia nei secoli XV e XVI, Notizie ed elenchi}, Florence, Fratelli Alinari, 1960, 3 vol.
\end{teiBibl}
              \begin{teiBibl}[xmlid={bibl-246}]
\teiHi[rend={small-caps}]{Martin} Jacques, \teiHi[rend={italic}]{Heraldry in the Vatican / L’araldica in Vaticano}, Gerrards Cross, Van Duren, 1987.
\end{teiBibl}
              \begin{teiBibl}[xmlid={bibl-247}]
\teiHi[rend={small-caps}]{Marucci} Valerio, \teiHi[rend={small-caps}]{Mazzo} Antonio, \teiHi[rend={small-caps}]{Romano} Angela et \teiHi[rend={small-caps}]{Aquilecchia} Giovanni (éd.), \teiHi[rend={italic}]{Pasquinate romane del Cinquecento}, Rome, Salerno Editrice, 1983, 2 vol.
\end{teiBibl}
              \begin{teiBibl}[xmlid={bibl-248}]
\teiHi[rend={small-caps}]{Masson} Georgina, \teiHi[rend={italic}]{Queen Christina}, Londres, Secker \& Warburg, 1968.
\end{teiBibl}
              \begin{teiBibl}[xmlid={bibl-249}]
\teiHi[rend={small-caps}]{Micalizzi} Paolo, \teiHi[rend={italic}]{Gubbio. Storia dell’architettura e della città} \lbrack{}2\teiHi[rend={sup}]{e} éd.\rbrack{}, Gubbio, L’arte grafica, 2009.
\end{teiBibl}
              \begin{teiBibl}[xmlid={bibl-250}]
\teiHi[rend={small-caps}]{Millon} Henry A., \teiHi[rend={italic}]{Filippo Juvarra. Drawings from the roman period 1704-1714}, Rome, Edizioni dell’Elefante, 1984, vol. 1.
\end{teiBibl}
              \begin{teiBibl}[xmlid={bibl-251}]
—, \teiHi[rend={small-caps}]{Griseri} Andreina, \teiHi[rend={small-caps}]{Mc Phee} Sarah et \teiHi[rend={small-caps}]{Viale Ferrero} Mercedes (dir.), \teiHi[rend={italic}]{Filippo Juvarra. Drawings from the roman period 1704-1714}, Rome, Edizioni dell’Elefante, 1999, vol. 2.
\end{teiBibl}
              \begin{teiBibl}[xmlid={bibl-252}]
\teiHi[rend={small-caps}]{Mochi Onori} Lorenza, \teiHi[rend={small-caps}]{Schutze} Sebastian et \teiHi[rend={small-caps}]{Solinas} Francesco (dir.), \teiHi[rend={italic}]{I Barberini e la cultura europea del Seicento}, Rome, De Luca, 2007.
\end{teiBibl}
              \begin{teiBibl}[xmlid={bibl-253}]
\teiHi[rend={small-caps}]{Modesti} Adolfo, \teiHi[rend={italic}]{Corpus numismatum omnium romanorum pontificum}, Rome, Adolfo Modesti, 2002-2018, 6 vol.
\end{teiBibl}
              \begin{teiBibl}[xmlid={bibl-254}]
\teiHi[rend={italic}]{Monasticon Italiae}, Cesena, Badia di Santa Maria del Monte, 1981, t. 1 : \teiHi[rend={italic}]{Roma e Lazio}.
\end{teiBibl}
              \begin{teiBibl}[xmlid={bibl-255}]
\teiHi[rend={small-caps}]{Monstrelet} Enguerrand de, \teiHi[rend={italic}]{Chronique}, Louis Douët d’Arcq (éd.), Paris, V\teiHi[rend={sup}]{e} Jules Renouard, 1857-1862, 6 vol.
\end{teiBibl}
              \begin{teiBibl}[xmlid={bibl-256}]
\teiHi[rend={small-caps}]{Montagu} Jennifer, \teiHi[rend={italic}]{Alessandro Algardi}, New Haven-Londres, Yale University Press, 1985, 2 vol.
\end{teiBibl}
              \begin{teiBibl}[xmlid={bibl-257}]
\teiHi[rend={small-caps}]{Moretti} Italo et \teiHi[rend={small-caps}]{Stopani} Renato, \teiHi[rend={italic}]{Architettura romanica religiosa a Gubbio}, Florence, Salimbeni, 1973.
\end{teiBibl}
              \begin{teiBibl}[xmlid={bibl-258}]
\teiHi[rend={small-caps}]{Moroni} Gaetano, \teiHi[rend={italic}]{Dizionario di erudizione storico-ecclesiastica da S. Pietro sino ai nostri giorni}, In Venezia, dalla Tipografia Emiliana, 1840-1878, 103 vol.
\end{teiBibl}
              \begin{teiBibl}[xmlid={bibl-259}]
\teiHi[rend={small-caps}]{Myers} Mary, \teiHi[rend={italic}]{Architectural and ornament drawings}: \teiHi[rend={italic}]{Juvarra, Vanvitelli, The Bibiena family \& other Italian draughtsmen}, New York, The Metropolitan Museum of Art, 1975.
\end{teiBibl}
              \begin{teiBibl}[xmlid={bibl-260}]
\teiHi[rend={small-caps}]{Nardelli} Giuseppe Maria, \teiHi[rend={italic}]{La pianta di Gubbio di Ignazio Cassetta stampata da Joannis Blaeu e Pierre Mortier: rappresentazione e simbolismi}, Gubbio, Rotary Club, 2001.
\end{teiBibl}
              \begin{teiBibl}[xmlid={bibl-261}]
\teiHi[rend={small-caps}]{Neppi} Lionello, \teiHi[rend={italic}]{Schede del catalogo di opere d’arte mobili (Santa Maria in Vallicella)}, Rome, Ministero della Pubblica Istruzione, 1957.
\end{teiBibl}
              \begin{teiBibl}[xmlid={bibl-262}]
\teiHi[rend={small-caps}]{Noblet} Julien, \teiHi[rend={italic}]{En perpétuelle mémoire, Collégiales castrales et Saintes-Chapelles à vocation funéraire en France (1450-1560)}, Rennes, PUR, 2009.
\end{teiBibl}
              \begin{teiBibl}[xmlid={bibl-263}]
\teiHi[rend={small-caps}]{Nova} Alessandro, \teiHi[rend={italic}]{The Artistic Patronage of Pope Julius III (1550-1555), Profane Imagery and Buildings for the De Monte Family in Rome}, New York-Londres, Garland Publishing, 1988.
\end{teiBibl}
              \begin{teiBibl}[xmlid={bibl-264}]
\teiHi[rend={small-caps}]{Onofrio} Cesare d’, \teiHi[rend={italic}]{Roma val bene un’abiura}, Rome, Fratelli Palombi, 1976.
\end{teiBibl}
              \begin{teiBibl}[xmlid={bibl-265}]
\teiHi[rend={small-caps}]{Oy}-\teiHi[rend={small-caps}]{Marra} Elisabeth, \teiHi[rend={italic}]{Profane Repräsentationskunst in Rom von Clemens VIII Aldobrandini bis Alexander VII Chigi. Studien zur Funktion und Semantik römischer Deckenfresken im höfischen Kontext}, Munich-Berlin, Deutscher Kunstverlag, 2005.
\end{teiBibl}
              \begin{teiBibl}[xmlid={bibl-266}]
\teiHi[rend={small-caps}]{Palliolo} Paolo, \teiHi[rend={italic}]{Le feste pel conferimento del patriziato romano a Giuliano e Lorenzo de’ Medici}. Bologne, G. Romagnoli, 1885.
\end{teiBibl}
              \begin{teiBibl}[xmlid={bibl-267}]
\teiHi[rend={small-caps}]{Panciroli} Romeo, \teiHi[rend={italic}]{The audience suite of the papal apartments}, Cité du Vatican, Libreria Editrice Vaticana, 2004.
\end{teiBibl}
              \begin{teiBibl}[xmlid={bibl-268}]
\teiHi[rend={small-caps}]{Panza} Pierluigi, \teiHi[rend={italic}]{Piranesi architetto}, Milan, Guerini Studio, 1998.
\end{teiBibl}
              \begin{teiBibl}[xmlid={bibl-269}]
\teiHi[rend={small-caps}]{Paravicini Bagliani} Agostino, \teiHi[rend={italic}]{Il bestiario del papa}, Turin, Einaudi, 2016.
\end{teiBibl}
              \begin{teiBibl}[xmlid={bibl-270}]
\teiHi[rend={small-caps}]{Pasini Frassoni} Ferruccio, \teiHi[rend={italic}]{Essai d’armorial des papes}, Rome, Collège héraldique, 1906.
\end{teiBibl}
              \begin{teiBibl}[xmlid={bibl-271}]
\teiHi[rend={small-caps}]{Pastoureau} Michel, \teiHi[rend={italic}]{Typologie des sources du Moyen Âge occidental}, Turnhout, Brepols, 1976, fasc. 20 : \teiHi[rend={italic}]{Les armoiries}.
\end{teiBibl}
              \begin{teiBibl}[xmlid={bibl-272}]
—, \teiHi[rend={italic}]{L’ours. L’histoire d’un roi déchu}, Paris, Seuil, 2016.
\end{teiBibl}
              \begin{teiBibl}[xmlid={bibl-273}]
\teiHi[rend={small-caps}]{Pecchiai} Pio, \teiHi[rend={italic}]{Un assassinio politico a Roma nel Cinquecento}, Rome, Biblioteca d’Arte Editrice, 1956.
\end{teiBibl}
              \begin{teiBibl}[xmlid={bibl-274}]
\teiHi[rend={bold}]{—}, \teiHi[rend={italic}]{I Barberini}, Rome, Biblioteca d’Arte Editrice, 1959.
\end{teiBibl}
              \begin{teiBibl}[xmlid={bibl-275}]
\teiHi[rend={small-caps}]{Perini} Davide Aurelio, \teiHi[rend={italic}]{Bibliographia Augustiniana}, Florence, Tip. Sordomuti, 1929-1938, 4 vol.
\end{teiBibl}
              \begin{teiBibl}[xmlid={bibl-276}]
\teiHi[rend={italic}]{Piranesi architetto}, catalogue d’exposition, American Academy in Rome, 15 mai-5 juillet 1992, Rome, Edizioni dell’Elefante, 1992.
\end{teiBibl}
              \begin{teiBibl}[xmlid={bibl-277}]
\teiHi[rend={small-caps}]{Pizzagalli} Daniela, \teiHi[rend={italic}]{La regina di Roma. Vita e misteri di Cristina di Svezia nell’Italia barocca}, Milan, Rizzoli, 2002.
\end{teiBibl}
              \begin{teiBibl}[xmlid={bibl-278}]
\teiHi[rend={small-caps}]{Plessi} Giuseppe (dir.), \teiHi[rend={italic}]{Le Insignia degli Anziani del comune dal 1530 al 1796: catalogo-inventario}, Bologne, ASBO, 1954.
\end{teiBibl}
              \begin{teiBibl}[xmlid={bibl-279}]
\teiHi[rend={small-caps}]{Polo} Anna, \teiHi[rend={italic}]{Las }Annotationi \teiHi[rend={italic}]{de Gauges de’ Gozze (1631) en relación con la obra didáctica de Franciosini}, Padoue, \teiHi[rend={small-caps}]{CLEUP}, 2021.
\end{teiBibl}
              \begin{teiBibl}[xmlid={bibl-280}]
\teiHi[rend={small-caps}]{Poos} Adam, \teiHi[rend={italic}]{San Lorenzo in Damaso a Roma}, \teiHi[rend={italic}]{Guida alla basilica}, Rome, \teiHi[rend={small-caps}]{GB} Editoria, 2015.
\end{teiBibl}
              \begin{teiBibl}[xmlid={bibl-281}]
\teiHi[rend={small-caps}]{Poulbrière} Jean-Baptiste, \teiHi[rend={italic}]{Notice historique et archéologique sur Castelnau-de-Bretenoux, Lot}, Tulle, Imprimerie E. Crauffon, 1873.
\end{teiBibl}
              \begin{teiBibl}[xmlid={bibl-282}]
\teiHi[rend={small-caps}]{Poutrin} Isabelle, \teiHi[rend={italic}]{Les convertis du pape. Une famille de banquiers juifs à Rome au }\teiHi[rend={small-caps italic}]{xvi}\teiHi[rend={italic sup}]{e}\teiHi[rend={italic}]{ siècle}, Paris, Seuil, 2023.
\end{teiBibl}
              \begin{teiBibl}[xmlid={bibl-283}]
\teiHi[rend={small-caps}]{Prodi} Paolo, \teiHi[rend={italic}]{Il sovrano pontefice: un corpo e due anime: la monarchia papale nella prima età moderna}, Bologne, Il Mulino, 1982.
\end{teiBibl}
              \begin{teiBibl}[xmlid={bibl-284}]
\teiHi[rend={small-caps}]{Rietstap} Jean-Baptiste, \teiHi[rend={italic}]{Armorial général}, Gouda, G. B. van Goor, 1861.
\end{teiBibl}
              \begin{teiBibl}[xmlid={bibl-285}]
\teiHi[rend={bold}]{—}, \teiHi[rend={italic}]{Armorial général}, Londres, Heraldry Today, 1988, 2 vol.
\end{teiBibl}
              \begin{teiBibl}[xmlid={bibl-286}]
\teiHi[rend={small-caps}]{Rivosecchi} Valerio, \teiHi[rend={italic}]{Esotismo in Roma Barocca. Studi sul Padre Kircher}, Rome, Bulzoni, 1982.
\end{teiBibl}
              \begin{teiBibl}[xmlid={bibl-287}]
\teiHi[rend={small-caps}]{Robertson} Clare, \teiHi[rend={italic}]{The city and the visual arts under Clement VIII}, New Haven, Yale University Press, 2015.
\end{teiBibl}
              \begin{teiBibl}[xmlid={bibl-288}]
\teiHi[rend={small-caps}]{Roca De Amicis} Augusto, \teiHi[rend={italic}]{L’opera di Borromini in San Giovanni in Laterano: gli anni della fabbrica (1646,1650} ), Rome, Dedalo, 1996.
\end{teiBibl}
              \begin{teiBibl}[xmlid={bibl-289}]
\teiHi[rend={small-caps}]{Roettgen} Steffi, \teiHi[rend={italic}]{Fresques italiennes du Baroque aux Lumières (1600-1797)}, Jeanne Étoré-Lortholary (trad.), Paris, Citadelle \& Mazenod, 2007.
\end{teiBibl}
              \begin{teiBibl}[xmlid={bibl-290}]
\teiHi[rend={small-caps}]{Rollo}-\teiHi[rend={small-caps}]{Koster} Joëlle, \teiHi[rend={italic}]{Avignon and its papacy 1309-1417: Popes, Institutions and Society}, Londres, Rowman \& Littlefield, 2015.
\end{teiBibl}
              \begin{teiBibl}[xmlid={bibl-291}]
\teiHi[rend={small-caps}]{Romanis} Alfonso Camillo De, \teiHi[rend={italic}]{Per un nuovo Battistero a Santa Prisca. Nova et vetera}, Cité du Vatican, Tipografia poliglotta vaticana, 1948.
\end{teiBibl}
              \begin{teiBibl}[xmlid={bibl-292}]
\teiHi[rend={small-caps}]{Roncalli} Angelo, \teiHi[rend={italic}]{Il Cardinale Cesare Baronio}, Rome, Edizioni di storia e letteratura, 1961.
\end{teiBibl}
              \begin{teiBibl}[xmlid={bibl-293}]
\teiHi[rend={small-caps}]{Roscoe} William, \teiHi[rend={italic}]{Life and Pontificate of Leo X}, Londres, Henry G. Bohn, 1846, 2 vol.
\end{teiBibl}
              \begin{teiBibl}[xmlid={bibl-294}]
\teiHi[rend={small-caps}]{Rossi} Iohannes Baptista de et \teiHi[rend={small-caps}]{Duchesne} Ludovicus, \teiHi[rend={italic}]{Acta Sanctorum Novembris.} Bruxelles, 1894, vol. 2, part. 1 : \teiHi[rend={italic}]{Martyrologium Hieronymianum}.
\end{teiBibl}
              \begin{teiBibl}[xmlid={bibl-295}]
\teiHi[rend={small-caps}]{Safarik} Eduard A., \teiHi[rend={italic}]{L’Apoteosi di Romolo di Domenico Maria Canuti in Palazzo Altieri}, Rome, Finnat Euramerica, 1995.
\end{teiBibl}
              \begin{teiBibl}[xmlid={bibl-296}]
\teiHi[rend={small-caps}]{Salvestrini} Francesco, \teiHi[rend={italic}]{Santa Maria di Vallombrosa. Patrimonio e vita economica di un grande monastero medievale}, Florence, Olschki, « Biblioteca Storica Toscana, I, 33 », 1998.
\end{teiBibl}
              \begin{teiBibl}[xmlid={bibl-297}]
\teiHi[rend={small-caps}]{Salvestrini} Francesco et \teiHi[rend={small-caps}]{Ciliberti} Rosagemma, \teiHi[rend={italic}]{I Vallombrosani nel Piemonte medievale e moderno. Ospizi e monasteri intorno alla strada di Francia}, Rome, Viella, 2014.
\end{teiBibl}
              \begin{teiBibl}[xmlid={bibl-298}]
\teiHi[rend={small-caps}]{Sangiorgi} Geremia, \teiHi[rend={italic}]{S. Prisca e il suo mitreo}, Rome, Marietti, « Le chiese di Roma illustrate, 101 », 1968.
\end{teiBibl}
              \begin{teiBibl}[xmlid={bibl-299}]
\teiHi[rend={small-caps}]{Schiavo} Armando, \teiHi[rend={italic}]{Villa Ludovisi e Palazzo Margherita}, Rome, Roma Amor, 1981.
\end{teiBibl}
              \begin{teiBibl}[xmlid={bibl-300}]
\teiHi[rend={small-caps}]{Simonetta} Giuseppe, \teiHi[rend={small-caps}]{Gigli} Laura et \teiHi[rend={small-caps}]{Marchetti} Gabriella, \teiHi[rend={italic}]{Sant’Angese in Agone a Piazza Navona. Bellezza, Proporzione, Armonia nelle fabbriche Pamphili}, Rome, Gangemi, 2003.
\end{teiBibl}
              \begin{teiBibl}[xmlid={bibl-301}]
\teiHi[rend={small-caps}]{—}, \teiHi[rend={small-caps}]{—} et \teiHi[rend={small-caps}]{—}, \teiHi[rend={italic}]{Sant’Agnese in Agone a Piazza Navona, Immagine, Luce, Ordine, Suono nelle fabbriche Pamphili}, Rome, Gangemi, 2005.
\end{teiBibl}
              \begin{teiBibl}[xmlid={bibl-302}]
\teiHi[rend={small-caps}]{Spike} John T. (éd.), \teiHi[rend={italic}]{The Illustrated Bartsch}, New York, Abaris Book, 1981, vol. 41.
\end{teiBibl}
              \begin{teiBibl}[xmlid={bibl-303}]
\teiHi[rend={small-caps}]{Stagni} Simonetta, \teiHi[rend={italic}]{Domenico Maria Canuti}, Rimini, Luisè, 1988.
\end{teiBibl}
              \begin{teiBibl}[xmlid={bibl-304}]
\teiHi[rend={small-caps}]{Takahashi} Kenichi, \teiHi[rend={italic}]{Giovanni Luigi Valesio. Ritratto de « l’Instabile academico incaminato »}, Bologne, Clueb, 2007.
\end{teiBibl}
              \begin{teiBibl}[xmlid={bibl-305}]
\teiHi[rend={small-caps}]{Talamo} Emilia Anna, \teiHi[rend={italic}]{Codices cantorum, miniature e disegni nei codici della Cappella Sistina}, Florence, Officine del Novecento, 1997.
\end{teiBibl}
              \begin{teiBibl}[xmlid={bibl-306}]
\teiHi[rend={small-caps}]{Thiry} Steven, \teiHi[rend={italic}]{Matter(s) of State. Heraldic Display and Discourse in the Early Modern Monarchy (c. 1480-1650)}, Ostfildern, Thorbecke Verlag, « Heraldic Studies, 2 », 2018.
\end{teiBibl}
              \begin{teiBibl}[xmlid={bibl-307}]
\teiHi[rend={small-caps}]{Todini} Filippo, \teiHi[rend={italic}]{La pittura umbra dal Duecento al primo Cinquecento}, Milan, Longanesi, 1989.
\end{teiBibl}
              \begin{teiBibl}[xmlid={bibl-308}]
\teiHi[rend={small-caps}]{Vaillancourt} Daniel et \teiHi[rend={small-caps}]{Wagner} Marie-France (dir.), \teiHi[rend={italic}]{Dix-septième siècle}, numéro spécial :\teiHi[rend={italic}]{ Les entrées royales}, n\teiHi[rend={sup}]{o} 3, 2001.
\end{teiBibl}
              \begin{teiBibl}[xmlid={bibl-309}]
\teiHi[rend={small-caps}]{Villani} Giovanni, \teiHi[rend={italic}]{Nuova Cronica}, Giuseppe Porta (éd.), Parme, Guanda, 1990-1991, livre \teiHi[rend={small-caps}]{VII}, vol. 2.
\end{teiBibl}
              \begin{teiBibl}[xmlid={bibl-310}]
\teiHi[rend={small-caps}]{Von Pastor} Ludwig Freiherr, \teiHi[rend={italic}]{The History of the Popes: From the Close of the Middle Ages, Drawn From the Secret Archives of the Vatican and Other Original Sources}, R. F. Kerr (éd.), Londres, Kegan Paul, Trench, Trubner \& Co., 1910, vol. 8; 1933, vol. 23.
\end{teiBibl}
              \begin{teiBibl}[xmlid={bibl-311}]
\teiHi[rend={small-caps}]{Waddy} Patricia, \teiHi[rend={italic}]{Seventeenth-Century Roman Palaces. Use and the Art of the Plan}, Cambridge-Londres, The \teiHi[rend={small-caps}]{MIT} Press, 1990.
\end{teiBibl}
              \begin{teiBibl}[xmlid={bibl-312}]
\teiHi[rend={small-caps}]{Weber} Christoph Friedrich, \teiHi[rend={italic}]{Zeichen der Ordnung und des Aufruhrs. Heraldische Symbolik in italienischen Stadtkommunen des Mittelalters}, Cologne, Böhlau, 2011.
\end{teiBibl}
              \begin{teiBibl}[xmlid={bibl-313}]
\teiHi[rend={small-caps}]{Weddigen} Tristam, De \teiHi[rend={small-caps}]{Blaauw} Sible et \teiHi[rend={small-caps}]{Kempers} Bram (éd.), \teiHi[rend={italic}]{Functions and Decorations: Art and Ritual at the Vatican Palace in the Middle Ages and the Renaissance}, Cité du Vatican-Turnhout, BAV-Brepols, 2003.
\end{teiBibl}
              \begin{teiBibl}[xmlid={bibl-314}]
\teiHi[rend={small-caps}]{Wilton}-\teiHi[rend={small-caps}]{Ely} John, \teiHi[rend={italic}]{Piranesi as Architect and Designer}, New York-New Haven-Londres, The Pierpont Morgan Library-Yale University Press, 1993.
\end{teiBibl}
              \begin{teiBibl}[xmlid={bibl-315}]
\teiHi[rend={small-caps}]{Yorke} Carolyn et \teiHi[rend={small-caps}]{Minor} Heather Hyde, \teiHi[rend={italic}]{Piranese unbound}, Princeton et Oxford, Princeton University Press, 2020.
\end{teiBibl}
              \begin{teiBibl}[xmlid={bibl-316}]
\teiHi[rend={small-caps}]{Zola} Émile, \teiHi[rend={italic}]{Les trois villes. Rome}, Paris, G. Charpentier et E. Fasquelle, 1896.
\end{teiBibl}
              \begin{teiBibl}[xmlid={bibl-317}]
\teiHi[rend={small-caps}]{Zuccari} Alessandro, \teiHi[rend={small-caps}]{Benocci} Carla et \teiHi[rend={small-caps}]{Bagolan} Maria Cristina (dir.)\teiHi[rend={italic}]{, La Spagna sul Gianicolo}, Rome, Eurografica, 2004, 3 vol.
\end{teiBibl}
              \begin{teiBibl}[xmlid={bibl-318}]
\teiHi[rend={small-caps}]{Zwart} Hub, \teiHi[rend={italic}]{Ethical Consensus and the Truth of Laughter: The Structure of Moral Transformations}, Kampen, Kok Pharos, 1996.
\end{teiBibl}
            
\end{teiDiv}
          
\end{teiDiv}
          \begin{teiDiv}[type={section1},xmlid={div18}]

            \teiHead[xmlid={articles-001}]{Articles}
            \begin{teiDiv}[type={section2},xmlid={div19}]

              \teiHead[xmlid={articles-douvrages-001}]{Articles d’ouvrages}
              \begin{teiBibl}[xmlid={bibl-319}]
\teiHi[rend={small-caps}]{Alteri} Giancarlo, « Monete e medaglie del pontificato di Clemente \teiHi[rend={small-caps}]{XIII} (1758-1769) », dans Barbara Jatta (dir.), \teiHi[rend={italic}]{Piranesi e l’Aventino}, Milan, Electa, 1998, p. 204.
\end{teiBibl}
              \begin{teiBibl}[xmlid={bibl-320}]
\teiHi[rend={small-caps}]{Andretta} Stefano, « Farnese, Girolamo », dans \teiHi[rend={italic}]{DBI}, Rome, Istituto della Enciclopedia italiana, 1995, vol. 45, p. 95-98.
\end{teiBibl}
              \begin{teiBibl}[xmlid={bibl-321}]
\teiHi[rend={small-caps}]{Anselmi} Salvatore Enrico, « Novità documentarie su S. Prisca in età barocca. L’attività di Carlo Lambardi e di Carlo Fontana », dans Mario Bevilacqua et Daniela Gallavotti Cavallero (dir.), \teiHi[rend={italic}]{L’Aventino dal Rinascimento a oggi. Arte e Architettura}, Rome, Artemide, 2010, p. 73-77.
\end{teiBibl}
              \begin{teiBibl}[xmlid={bibl-322}]
\teiHi[rend={small-caps}]{Argan} Giulio Carlo, « La rettorica et l’arte barocca », dans Enrico Castelli (dir.), \teiHi[rend={italic}]{Retorica e Barocco. Atti del 3 congresso internazionale di studi umanistici}, Rome, Fratelli Bocca, 1955, p. 10-14, republié dans \teiHi[rend={italic}]{Studi e Note. Da Bramante a Canova}, Rome, Bulzoni, 1970.
\end{teiBibl}
              \begin{teiBibl}[xmlid={bibl-323}]
\teiHi[rend={small-caps}]{Arrighi} Vanna et \teiHi[rend={small-caps}]{Insabato} Elisabetta, « Tra storia e mito, la ricostruzione del passato familiare nella nobiltà toscana dei secoli \teiHi[rend={small-caps}]{XVI}-\teiHi[rend={small-caps}]{XVIII} », dans \teiHi[rend={italic}]{L’identità genealogica e araldica. Fonti, metodologie, interdisciplinarietà}, actes du colloque, Turin, Archive d’État, 21-26 septembre 1988, Rome, Ministero per i beni e le attività culturali, 2000, p. 1099-1121.
\end{teiBibl}
              \begin{teiBibl}[xmlid={bibl-324}]
\teiHi[rend={small-caps}]{Barry} Fabio, « “Onward Christian Soldier” : Piranesi at Santa Maria del Priorato », dans Mario Bevilacqua et Daniela Gallavotti Cavallero (dir.), \teiHi[rend={italic}]{L’Aventino dal Rinascimento a oggi, Arte e architettura}, Rome, Editoriale Artemide, 2010, p. 140-175.
\end{teiBibl}
              \begin{teiBibl}[xmlid={bibl-325}]
\teiHi[rend={small-caps}]{Beccarini} Dario, « Riflessioni e nuove proposte su Francisco Vieira de Matos, detto Vieira Lusitano, “eccellente disegnatore \lbrack{}…\rbrack{} felice e franco nell’inventare” », dans Elisa Debenedetti (dir.), \teiHi[rend={italic}]{Studi sul Settecento romano}, Rome, Edizioni Quazar, 2021, p. 207-220.
\end{teiBibl}
              \begin{teiBibl}[xmlid={bibl-326}]
\teiHi[rend={small-caps}]{Bernardini} Carla, « Bologna 1660. Una committenza Farnese per papa Alessandro \teiHi[rend={small-caps}]{VII} », dans Andrea Bacchi et Luca M. Barbero (dir.), \teiHi[rend={italic}]{Studi in onore di Stefano Tumidei}, Venise-Bologne, Fondazione Cini-Fondazione Zeri, 2016, p. 279-287.
\end{teiBibl}
              \begin{teiBibl}[xmlid={bibl-327}]
\teiHi[rend={small-caps}]{Bertrand} Pascal-François, « L’uso dell’arazzo a Roma: l’esempio del Cardinale Francesco Barberini », dans Lorenza Mochi Onori, Sebastian Schutze et Francesco Solinas (dir.), \teiHi[rend={italic}]{I Barberini e la cultura europea del Seicento}, Rome, De Luca, 2007, p. 421-430.
\end{teiBibl}
              \begin{teiBibl}[xmlid={bibl-328}]
\teiHi[rend={small-caps}]{Bevilacqua} Mario, « The Young Piranesi: The Itineraries of his formation », dans Mario Bevilacqua, Heather Hyde Minor et Fabio Barri (dir.), \teiHi[rend={italic}]{The Serpent and the Stylus, Essays on G. B. Piranesi}, Ann Arbor, The University of Michigan Press, « Supplements to the Memoirs of the American Academy in Rome », 2006, p. 13-53.
\end{teiBibl}
              \begin{teiBibl}[xmlid={bibl-329}]
\teiHi[rend={small-caps}]{Blackburn} Bonnie, « Messages in Miniature: Pictorial Programme and Theological Implications in the Alamire Choirbooks », dans \teiHi[rend={italic}]{The Burgundian-Habsburg Court Complex of Music Manuscripts (1500-1535) and the Workshop of Petrus Alamire: Colloquium Proceedings, Leuven, 25-28 November 1999}, Louvain-Neerpelt, Alamire Foundation, 2003, p. 161-184.
\end{teiBibl}
              \begin{teiBibl}[xmlid={bibl-330}]
\teiHi[rend={small-caps}]{Boesch Gajano} Sofia, « Storia e tradizione vallombrosane », dans Antonella Degl’Innocenti (dir.), \teiHi[rend={italic}]{Vallombrosa. Memorie agiografiche e culto delle reliquie}, Rome, Viella, 2012, p. 15-115.
\end{teiBibl}
              \begin{teiBibl}[xmlid={bibl-331}]
\teiHi[rend={small-caps}]{Boiteux} Martine, « La vacance du siège pontifical. De la mort et des funérailles à l’investiture du Pape : les rites de l’époque moderne », dans José Pedro Paiva (dir.), \teiHi[rend={italic}]{Religious Ceremonials and Images}. \teiHi[rend={italic}]{Power and social meaning (1400-1750)}, Coimbra, Palimage Editores, 2002, p. 103-141.
\end{teiBibl}
              \begin{teiBibl}[xmlid={bibl-332}]
\teiHi[rend={bold}]{—}, « \teiHi[rend={italic}]{Il Possesso}. La presa di potere del Sovrano Pontefice sulla città di Roma », dans Francesco Buranelli (dir.),\teiHi[rend={italic}]{ Habemus Papam. Le elezioni pontificie da San Pietro a Benedetto XIV}, Rome, De Luca, 2006, p. 131-140.
\end{teiBibl}
              \begin{teiBibl}[xmlid={bibl-333}]
\teiHi[rend={bold}]{—}, « Linguaggio figurativo ed efficacia rituale », dans Francesca Cantù (dir.), \teiHi[rend={italic}]{I linguaggi del potere. Politica e religione nell’età barocca}, Rome, Viella, 2007, p. 39-79.
\end{teiBibl}
              \begin{teiBibl}[xmlid={bibl-334}]
\teiHi[rend={bold}]{—}, « Les Barberini, Rome et la France : fête et politique », dans Lorenza Mochi Onori, Sebastian Schutze et Francesco Solinas (dir.), \teiHi[rend={italic}]{I Barberini e la cultura europea del Seicento}, Rome, De Luca, 2007, p. 345-360.
\end{teiBibl}
              \begin{teiBibl}[xmlid={bibl-335}]
\teiHi[rend={bold}]{—}, « Funérailles féminines dans la Rome baroque », dans Bernard Dompnier (dir.), \teiHi[rend={italic}]{Les cérémonies extraordinaires du catholicisme baroque}, Clermont-Ferrand, PUBP, 2009, p. 389-421.
\end{teiBibl}
              \begin{teiBibl}[xmlid={bibl-336}]
\teiHi[rend={bold}]{—}, « La cerimonia della canonizzazione. Teatro riti, stendardi e immagini », dans Alessandra Bartolomei Romagnoli (dir.), \teiHi[rend={italic}]{La canonizzazione di Santa Francesca Romana, Santità, cultura e istituzioni a Roma tra Medioevo e età moderna}, Florence, Edizioni del Galluzzo per la Fondazione Ezio Franceschini, 2013, p. 99-121.
\end{teiBibl}
              \begin{teiBibl}[xmlid={bibl-337}]
\teiHi[rend={bold}]{—}, « Il cerimoniale. La necessità della magnificenza », dans Francesco Buranelli (dir.), \teiHi[rend={italic}]{Il Vaticano Barocco, Arte, architettura e cerimoniale}, Milan, Jaca Book, 2014, p. 10-27.
\end{teiBibl}
              \begin{teiBibl}[xmlid={bibl-338}]
\teiHi[rend={bold}]{—}, « Triomphes romains des Sobieski, 1683-1735 », dans \teiHi[rend={italic}]{I Sobieski a Roma. La famiglia reale polacca nella Città Eterna}, Varsovie, Muzeum Pałacu Króla Jana \teiHi[rend={small-caps}]{III} w Wiłanwie, 2018, p. 48-51.
\end{teiBibl}
              \begin{teiBibl}[xmlid={bibl-339}]
\teiHi[rend={bold}]{—}, « La Croce luminosa, un rito della Settimana Santa nella basilica di San Pietro », dans Roberto Regoli et Ilaria Fiumi Sermattei (dir\teiHi[rend={italic}]{.}),\teiHi[rend={italic}]{ La religione dei nuovi tempi. Il riformismo spirituale nell’età di Leone XII}, Ancone, Consiglio regionale delle Marche, 2020, p. 281-316.
\end{teiBibl}
              \begin{teiBibl}[xmlid={bibl-340}]
\teiHi[rend={small-caps}]{Borello} Benedetta, « Strategie di insediamento in città: i Pamphilj a Roma nel primo Cinquecento », dans Maria Antonietta Visceglia (dir.), \teiHi[rend={italic}]{La nobiltà romana in età moderna}, Rome, Carocci, 2001, p. 31-61.
\end{teiBibl}
              \begin{teiBibl}[xmlid={bibl-341}]
\teiHi[rend={bold}]{—}, « Pamphili, Girolamo », dans \teiHi[rend={italic}]{DBI}, Rome, Istituto della Enciclopedia italiana, 2014, vol. 80, p. 670-671.
\end{teiBibl}
              \begin{teiBibl}[xmlid={bibl-342}]
\teiHi[rend={small-caps}]{Bouyé} Édouard, « Les clefs de saint Pierre, sur la terre comme au ciel », dans Catalina Girbea, Christine Manigand, Jérôme Grévy, Laurent Hablot, Denise Turrel et Martin Aurell (dir.), \teiHi[rend={italic}]{Signes et couleurs des identités politiques}, Rennes, \teiHi[rend={small-caps}]{PUR}, 2008, p. 275-312.
\end{teiBibl}
              \begin{teiBibl}[xmlid={bibl-343}]
\teiHi[rend={small-caps}]{Brunelli} Giampiero, « Vidoni, Pietro senior », dans \teiHi[rend={italic}]{DBI}, Rome, Istituto della Enciclopedia italiana, 2020, vol. 99, p. 204-207.
\end{teiBibl}
              \begin{teiBibl}[xmlid={bibl-344}]
\teiHi[rend={small-caps}]{Buscaroli Fabbri} Beatrice, « Il Cardinal Farnese e la sua sala. Un ciclo di affreschi per la famiglia e la città », dans Camilla Bottino (dir.), \teiHi[rend={italic}]{Il palazzo comunale di Bologna. Storia, architettura e restauri}, Bologne, Compositori, 1999, p. 99-110.
\end{teiBibl}
              \begin{teiBibl}[xmlid={bibl-345}]
\teiHi[rend={small-caps}]{Coldagelli} Umberto, « Armanni, Vincenzo », dans\teiHi[rend={italic}]{ DBI}, Rome, Istituto della Enciclopedia italiana, 1962, vol. 4, p. 222-224.
\end{teiBibl}
              \begin{teiBibl}[xmlid={bibl-346}]
\teiHi[rend={small-caps}]{Cacciamani} Giuseppe, \teiHi[rend={italic}]{Dizionario degli Istituti di Perfezione}, Milan, Paoline, 1974, vol. 1.
\end{teiBibl}
              \begin{teiBibl}[xmlid={bibl-347}]
\teiHi[rend={small-caps}]{Carlevaris} Anna Laura, D\teiHi[rend={small-caps}]{e Carlo} Laura et \teiHi[rend={small-caps}]{Di Marzio} Daniel, « La sala Clementina in Vaticano tra manierismo e barocco », dans Fauzia Farneti et Deanna Lenzi (dir.), \teiHi[rend={italic}]{L’architettura dell’inganno. Quadraturismo e grande decorazione nella pittura di età barocca}, Florence, Alinea, 2004, p. 133-148.
\end{teiBibl}
              \begin{teiBibl}[xmlid={bibl-348}]
\teiHi[rend={small-caps}]{Cattaneo} Maria Vittoria et \teiHi[rend={small-caps}]{Gianasso} Elena, « Il Corpus Juvarrianum della Biblioteca Nazionale Universitaria di Torino », dans Franca Porticelli, Costanza Roggero, Chiara Devoti et Gustavo Mola Di Nomaglio (dir.), \teiHi[rend={italic}]{Filippo Juvarra regista di corti e capitali dalla Sicilia al Piemonte all’Europa}, Turin, Centro Studi Piemontesi, 2020, p. 197-210.
\end{teiBibl}
              \begin{teiBibl}[xmlid={bibl-349}]
« Chauvigny de Brosse, armes des », dans la base Sigilla, http://www.sigilla.org/armoiries/chauvigny-25103.
\end{teiBibl}
              \begin{teiBibl}[xmlid={bibl-350}]
\teiHi[rend={small-caps}]{Cipriani} Angela, « Un programma belloriano », dans Franco Borsi, Gabriele Morolli, Maurizio Caperna \teiHi[rend={italic}]{et al}., \teiHi[rend={italic}]{Palazzo Altieri}, Gianfranco Spagnesi (introduction), Rome, Editalia, 1991, p. 177-208.
\end{teiBibl}
              \begin{teiBibl}[xmlid={bibl-351}]
\teiHi[rend={small-caps}]{Cooper} Richard, « Legate’s Luxury : the Entries of Cardinal Alessandro Farnese to Avignon and Carpentras, 1553 », dans Nicolas Russell et Hélène Visentin (dir.), \teiHi[rend={italic}]{French Ceremonial Entries in the Sixteenth century : Event, Image, Text}, Toronto, Center for Renaissance and Reformation Studies, 2007, p. 133-161.
\end{teiBibl}
              \begin{teiBibl}[xmlid={bibl-352}]
\teiHi[rend={small-caps}]{Cornini} Guido, \teiHi[rend={small-caps}]{Strobel} Anna Maria de et \teiHi[rend={small-caps}]{Serlupi Crescenzi} Maria, « L’appartamento papale di rappresentanza », dans Carlo Pietrangeli (dir.), \teiHi[rend={italic}]{Il Palazzo apostolico vaticano}, Florence, Nardini Editori, 1992, p. 169-185.
\end{teiBibl}
              \begin{teiBibl}[xmlid={bibl-353}]
\teiHi[rend={small-caps}]{Critien} John Edward, « Un manoscritto dell’Archivio Magistrale del Sovrano Militare Ordine di Malta di Roma », dans Barbara Jatta (dir.), \teiHi[rend={italic}]{Piranesi e l’Aventino}, Milan, Electa, 1998, p. 79-85.
\end{teiBibl}
              \begin{teiBibl}[xmlid={bibl-354}]
\teiHi[rend={small-caps}]{D’Amelio} Maria Grazia et \teiHi[rend={small-caps}]{Marder} Tod A., « La fontana dei Quattro Fiumi a piazza Navona, iconologia e costruzione », dans Jean-François Bernard (dir.), \teiHi[rend={italic}]{Piazza Navona, ou Place Navone, la plus belle \& \lbrack{}et\rbrack{} la plus grande}, Rome, École française de Rome, « Collection de l’École Française de Rome, 493 », 2014, p. 399-419.
\end{teiBibl}
              \begin{teiBibl}[xmlid={bibl-355}]
\teiHi[rend={small-caps}]{De Benedictis} Angela, « Tra “affratellamento comunitario” e “commensualità rituale”. La “condotta di vita” dei cittadini di governo tra ’400 e ’500 », dans Maurizio Ricci (dir.), \teiHi[rend={italic}]{L’architettura a Bologna nel Rinascimento (1460-1550): centro o periferia?}, Bologne, Minerva, 2001, p. 17-27.
\end{teiBibl}
              \begin{teiBibl}[xmlid={bibl-356}]
\teiHi[rend={small-caps}]{De Caro} Gaspare, « Brusoni, Girolamo », dans\teiHi[rend={italic}]{ Dizionario biografico degli italiani}, Rome, Istituto della Enciclopedia italiana, 1972, vol. 14, p. 712-720.
\end{teiBibl}
              \begin{teiBibl}[xmlid={bibl-357}]
\teiHi[rend={small-caps}]{Di Berardino} Angelo, « Prisca, santa », dans \teiHi[rend={italic}]{Nuovo dizionario patristico e di antichità cristiane}, Gênes-Milan, Marietti, 2008, vol. 3, col. 4331-4332.
\end{teiBibl}
              \begin{teiBibl}[xmlid={bibl-358}]
\teiHi[rend={small-caps}]{Di Monte} Michele, « La cultura antica e lo specchio della storia », dans Maurizia Cicconi, Flaminia Gennari Santori et Sebastian Schütze (dir.), \teiHi[rend={italic}]{L’immagine sovrana. Urbano VIII e i Barberini}, catalogue d’exposition, Rome, Officina libraria, 2023, p. 265-269.
\end{teiBibl}
              \begin{teiBibl}[xmlid={bibl-359}]
\teiHi[rend={small-caps}]{Donato} Maria Monica, « “Ogni cosa è pieno d’arme”. Uno sguardo dall’esterno », dans Matteo Ferrari (dir.), \teiHi[rend={italic}]{L’Arme segreta. Araldica e storia dell’arte nel Medioevo (secoli XIII-XIV)}, Alessandro Savorelli (introduction), Florence, Le Lettere, 2015, p. 267-283.
\end{teiBibl}
              \begin{teiBibl}[xmlid={bibl-360}]
\teiHi[rend={small-caps}]{Doulkaridou}-\teiHi[rend={small-caps}]{Ramantani} Elli, « Jules de Médicis et Pompeo Colonna. Enjeux héraldiques dans les missels de deux cardinaux antagonistes », dans Yvan Loskoutoff (dir.), \teiHi[rend={italic}]{Héraldique \& Papauté, Moyen Âge-Temps modernes}, Mont-Saint-Aignan, PURH, 2020, p. 197-233.
\end{teiBibl}
              \begin{teiBibl}[xmlid={bibl-361}]
\teiHi[rend={small-caps}]{Ebert}-\teiHi[rend={small-caps}]{Schifferer} Sybille, « La romanità in termini bolognesi: Domenico Maria Canuti a palazzo Altieri », dans Sabine Frommel (dir.), \teiHi[rend={italic}]{Crocevia e capitale della migrazione artistica: forestieri a Bologna e bolognesi nel mondo (secolo XVII)}, Bologne, Bologna University Press, 2012, p. 327-345.
\end{teiBibl}
              \begin{teiBibl}[xmlid={bibl-362}]
\teiHi[rend={small-caps}]{Ermini} Gian-Paolo, « La campana del palazzo del popolo di Orvieto (1316) », dans Matteo Ferrari (dir.), \teiHi[rend={italic}]{L’arme segreta, Araldica e storia dell’arte nel Medioevo (secoli XIII-XV)}, Alessandro Savorelli (introduction), Florence, Le Lettere, 2015, p. 109-125.
\end{teiBibl}
              \begin{teiBibl}[xmlid={bibl-363}]
\teiHi[rend={bold}]{—}, « La décoration picturale de la grande salle du palais de la Commune d’Orvieto », dans Torsten Hiltman et Miguel de Seixas (dir.), \teiHi[rend={italic}]{Heraldry in Medieval and Early Modern State Rooms}, Ostfildern, Thorbecke Verlag, « Heraldic Studies, 3 », 2020, p. 167-171.
\end{teiBibl}
              \begin{teiBibl}[xmlid={bibl-364}]
\teiHi[rend={small-caps}]{Esposito} Anna, « Il rione Parione durante il Pontificato Sistino: analisi di un’area campione. Osservazione sulla popolazione rionale », dans Massimo Miglio, Francesca Niutta, Diego Quaglioni et Concetta Ranieri (dir.), \teiHi[rend={italic}]{Un pontificato e una città. Sisto IV (1471-1484)}, actes de colloque, Rome, 3-7 décembre 1984, Cité du Vatican, Scuola Vaticana di Paleografia, Diplomatica e Archivistica, « Littera antiqua, 5 », 1986, p. 651-666.
\end{teiBibl}
              \begin{teiBibl}[xmlid={bibl-365}]
\teiHi[rend={small-caps}]{Fagiolo} Marcello, « I geroglifici, gli emblemi e l’araldica: note sul ragionamento simbolico di Borromini », dans Richard Boesel et Christoph L. Frommel (dir.), \teiHi[rend={italic}]{Borromini e l’universo barocco}, Milan, Electa, 1999, p. 95-105.
\end{teiBibl}
              \begin{teiBibl}[xmlid={bibl-366}]
\teiHi[rend={small-caps}]{Feci} Simona et \teiHi[rend={small-caps}]{Bortolotti} Luca, « Giustiniani, Benedetto », dans \teiHi[rend={italic}]{DBI}, Rome, Istituto della Enciclopedia italiana, 2001, vol. 57, p. 366-377.
\end{teiBibl}
              \begin{teiBibl}[xmlid={bibl-367}]
\teiHi[rend={small-caps}]{Federici} Fabrizio, « “Anticomoderno”: significati ed usi del termine nella letteratura artistica tra Cinque e Settecento », dans Maria Monica Donato et Massimo Ferretti (dir.), \teiHi[rend={italic}]{« Conosco un ottimo storico dell’arte… ». Per Enrico Castelnuovo. Scritti di allievi e amici pisani}, Pise, Edizioni della Normale, 2012, p. 291-296.
\end{teiBibl}
              \begin{teiBibl}[xmlid={bibl-368}]
\teiHi[rend={bold}]{—}, « Un partisan de la France en quête de protection : ce “bon vieux” chevalier Francesco Gualdi et le cardinal Mazarin », dans Yvan Loskoutoff et Patrick Michel (dir.), \teiHi[rend={italic}]{Mazarin, Rome et l’Italie}, Mont-Saint-Aignan, PURH, 2022, t. 2, p. 199-238.
\end{teiBibl}
              \begin{teiBibl}[xmlid={bibl-369}]
\teiHi[rend={bold}]{—}, « Baroque Returns: the Donations and Reuses of Francesco Gualdi », \teiHi[rend={italic}]{Contested Holdings: Museum Collections in Political}, dans Eva-Maria Troelenberg, Felicity Bodenstein, Damiana Otoiu (dir.), \teiHi[rend={italic}]{Contested Holdings. Museum Collections in Political, Epistemic and Artistic Processes of Return}, actes de colloque, Florence, Kunsthistorisches Institut in Florenz, 21-22 octobre 2016, New York-Oxford, Berghahn, 2022, p. 197-219.
\end{teiBibl}
              \begin{teiBibl}[xmlid={bibl-370}]
\teiHi[rend={small-caps}]{Ferrari} Matteo, « Stemmi esposti. Presenze araldiche nei broletti lombardi \teiHi[rend={italic}]{»}, dans Matteo Ferrari (dir.), \teiHi[rend={italic}]{L’arme segreta, Araldica e storia dell’arte nel Medioevo (secoli XIII-XV)}, Alessandro Savorelli (introduction), Florence, Le Lettere, 2015, p. 91-107.
\end{teiBibl}
              \begin{teiBibl}[xmlid={bibl-371}]
\teiHi[rend={bold}]{—}, « Héraldique et “mise en signe” de l’espace urbain en Poitou au Moyen Âge », \teiHi[rend={italic}]{L’héraldique dans la ville}, Ostfildern, Thorbecke Verlag, « Heraldic Studies, 6 », 2025.
\end{teiBibl}
              \begin{teiBibl}[xmlid={bibl-372}]
\teiHi[rend={small-caps}]{—}, \teiHi[rend={small-caps}]{Rao} Riccardo et \teiHi[rend={small-caps}]{Terenzi} Pierluigi, « Rappresentazioni del potere angioino nell’Italia comunale: sovrani, ufficiali, città », dans Thierry Pécout (dir.), \teiHi[rend={italic}]{Les officiers et la chose publique dans les territoires angevins (}\teiHi[rend={small-caps italic}]{xiii}\teiHi[rend={italic sup}]{e}\teiHi[rend={italic}]{-}\teiHi[rend={small-caps italic}]{xvi}\teiHi[rend={italic sup}]{e}\teiHi[rend={italic}]{ siècle) : vers une culture politique}, Rome, École française de Rome, « École française de Rome, 518, 4 », 2020, p. 279-319.
\end{teiBibl}
              \begin{teiBibl}[xmlid={bibl-373}]
\teiHi[rend={small-caps}]{Forni} Giorgio, « Nicolò Machiavelli », dans Ezio Raimondi (dir.), \teiHi[rend={italic}]{La letteratura italiana. Dalle origini al Cinquecento}, Milan, Bruno Mondadori, 2007, p. 346.
\end{teiBibl}
              \begin{teiBibl}[xmlid={bibl-374}]
\teiHi[rend={small-caps}]{Francastel} Pierre, « La Contre-Réforme et les arts en Italie à la fin du \teiHi[rend={small-caps}]{xvi}\teiHi[rend={sup}]{e} siècle », dans Henri Bédarida (dir.), \teiHi[rend={italic}]{À travers l’art italien du }\teiHi[rend={small-caps italic}]{xv}\teiHi[rend={italic sup}]{e }\teiHi[rend={italic}]{au }\teiHi[rend={small-caps italic}]{xx}\teiHi[rend={italic sup}]{e}\teiHi[rend={italic}]{ siècle}, Paris, Boivin-Société d’études italiennes, 1949, p. 63-113.
\end{teiBibl}
              \begin{teiBibl}[xmlid={bibl-375}]
\teiHi[rend={small-caps}]{Fumi Cambi Gado} Francesca, « Stemmi ed emblemi nella decorazione degli edifici\teiHi[rend={italic}]{ »}, dans Amerigo Restucci (dir.),\teiHi[rend={italic}]{ L’architettura civile in Toscana. Il Medioevo}, Sienne, Monte dei Paschi di Siena, 1995, p. 401-441.
\end{teiBibl}
              \begin{teiBibl}[xmlid={bibl-376}]
\teiHi[rend={small-caps}]{Gallavotti Cavallero} Daniela, « L’iconografia di santa Prisca e l’immagine di Pietro dalla basilica Vaticana al titulus della santa sull’Aventino nei secoli \teiHi[rend={small-caps}]{XVI} e \teiHi[rend={small-caps}]{XVII} », dans Mario Bevilacqua et Daniela Gallavotti Cavallero (dir.), \teiHi[rend={italic}]{L’Aventino dal Rinascimento a oggi. Arte e Architettura}, Rome, Artemide, 2010, p. 57-71.
\end{teiBibl}
              \begin{teiBibl}[xmlid={bibl-377}]
\teiHi[rend={small-caps}]{Gatouillat} Françoise, « Prudhomat. Église Saint-Louis, ancienne collégiale Saint-Jean-Baptiste », dans Michel Hérold (dir.), \teiHi[rend={italic}]{Corpus vitrearum. Les vitraux du Midi de la France}, Rennes, PUR, 2020, p. 238-240.
\end{teiBibl}
              \begin{teiBibl}[xmlid={bibl-378}]
\teiHi[rend={small-caps}]{Giacomelli} Alfeo, « La storia di Bologna dal 1650 al 1796: un racconto e una cronologia », dans Adriano Prosperi (dir.), \teiHi[rend={italic}]{Bologna nell’età moderna (secoli XVI-XVIII). }Vol. I.\teiHi[rend={italic}]{ Istituzioni, forme del potere, economia e società}, Bologne, Bologna University Press, 2009, p. 61-198.
\end{teiBibl}
              \begin{teiBibl}[xmlid={bibl-379}]
\teiHi[rend={small-caps}]{Giordano} Silvano, « Sisto V », dans \teiHi[rend={italic}]{Enciclopedia dei Papi}, Rome, Treccani, 2000, 3 vol., t. 3, p. 202-222.
\end{teiBibl}
              \begin{teiBibl}[xmlid={bibl-380}]
\teiHi[rend={small-caps}]{Hablot} Laurent, « Le décor emblématique chez les princes de la fin du Moyen Âge : un outil pour construire et qualifier l’espace », dans Thomas Lienhard (dir.), \teiHi[rend={italic}]{Construction de l’espace au Moyen Âge : pratiques et représentations}, actes du 37\teiHi[rend={sup}]{e} congrès de la Société des historiens médiévistes de l’enseignement supérieur public, Paris, Publications de la Sorbonne, 2007, p. 147-165.
\end{teiBibl}
              \begin{teiBibl}[xmlid={bibl-381}]
\teiHi[rend={bold}]{—}, « “\teiHi[rend={italic}]{Sens dessoubz dessus}.” Le blason de la trahison », dans Myriam Soria et Maïté Billoré (dir.), \teiHi[rend={italic}]{La trahison au Moyen Âge, De la monstruosité au crime politique (}\teiHi[rend={small-caps italic}]{v}\teiHi[rend={italic sup}]{e}\teiHi[rend={italic}]{-}\teiHi[rend={small-caps italic}]{xv}\teiHi[rend={italic sup}]{e}\teiHi[rend={italic}]{ siècle)}, Rennes, PUR, 2010, p. 331-347.
\end{teiBibl}
              \begin{teiBibl}[xmlid={bibl-382}]
\teiHi[rend={bold}]{—}, « Le bris des armes. L’iconoclasme héraldique dans la société médiévale », dans Marc Gil, Pascale Charron et Ambre Vilain (dir.), \teiHi[rend={italic}]{La pensée du regard. Études d’histoire de l’art du Moyen Âge offertes à Christian Heck}, Turnhout, Brepols, 2016, p. 181-191.
\end{teiBibl}
              \begin{teiBibl}[xmlid={bibl-383}]
\teiHi[rend={bold}]{—}, « Sacralisation of the royal coats of arms in Europe in the Middle Ages », dans Monserrat Herrero, Jaume Aurell, Angela Concetta Micelli Stout (dir.), \teiHi[rend={italic}]{Political theology in Medieval and Early Modern Europe. Discourses, Rites and Representations}, Turnhout, Brepols, 2017, p. 313-336.
\end{teiBibl}
              \begin{teiBibl}[xmlid={bibl-384}]
\teiHi[rend={small-caps}]{Hayez} Jérôme, « Un segno fra altri segni. Forme, significati e usi della marca mercantile verso il 1400 », dans Elena Cecchi Aste (dir.), \teiHi[rend={italic}]{Di mio nome e segno. « Marche » di mercanti nel carteggio Datini (secc. XIV-XV)}, Prato, Istituto di studi postali Aldo Cecchi, « Quaderni di storia postale, 30 », 2010 p. \teiHi[rend={small-caps}]{IX}-\teiHi[rend={small-caps}]{XLVI}.
\end{teiBibl}
              \begin{teiBibl}[xmlid={bibl-385}]
\teiHi[rend={small-caps}]{Hermann}-\teiHi[rend={small-caps}]{Fiore} Kristina, « La salle clémentine au Palais apostolique du Vatican », dans Jean-Marc Olivesi (dir.), \teiHi[rend={italic}]{Les cieux en gloire. Paradis en trompe-l’œil pour la Rome baroque}, catalogue d’exposition, Ajaccio, Musée Fesch, 2002, p. 101-115.
\end{teiBibl}
              \begin{teiBibl}[xmlid={bibl-386}]
\teiHi[rend={small-caps}]{Hubert} Hans W., « La nascita e lo sviluppo architettoncio del Palazzo del Comune di Bologna fra potere comunale e potere papale », dans Camilla Bottino (dir.), \teiHi[rend={italic}]{Il palazzo comunale di Bologna. Storia, architettura e restauri}, Bologne, Compositori, 1999, p. 64-87.
\end{teiBibl}
              \begin{teiBibl}[xmlid={bibl-387}]
\teiHi[rend={small-caps}]{Ladner} Gerhart Burian, « Vegetation Symbolism and the Concept of Renaissance », dans Millard Meiss (dir.), \teiHi[rend={italic}]{De artibus opuscula XL. Essays in Honor of Erwin Panofsky}, New York, New York University Press, 1961, vol. I, p. 303-322. 
\end{teiBibl}
              \begin{teiBibl}[xmlid={bibl-388}]
\teiHi[rend={small-caps}]{Laurentiis} Elena de, « \teiHi[rend={italic}]{Preparatio ad missam pontificalem }Vincent Raymond di Lodève (illuminator) », dans Jeffrey F. Hamburger, William P. Stoneman, Anne-Mare Eze et Nancy Netzer (dir.), \teiHi[rend={italic}]{Beyond words: illuminated manuscripts in Boston collections}, catalogue d’exposition, Chestnut Hill, McMullen Museum of Art-Boston College, 2016.
\end{teiBibl}
              \begin{teiBibl}[xmlid={bibl-389}]
\teiHi[rend={small-caps}]{Lavin} Irving, « Giambologna’s Neptune at the crossroads », dans Irving Lavin, \teiHi[rend={italic}]{Past-present essays on historicism in art from Donatello to Picasso}, Berkeley, University of California Press, 1993, p. 63-83.
\end{teiBibl}
              \begin{teiBibl}[xmlid={bibl-390}]
\teiHi[rend={bold}]{—}, « Bernini’s Bumbling Barberini Bees », dans Irving Lavin, \teiHi[rend={italic}]{Visible Spirit. The Art of Gianlorenzo Bernini}, Londres, The Pinder Press, 2009, t. 2, p. 955-1017.
\end{teiBibl}
              \begin{teiBibl}[xmlid={bibl-391}]
\teiHi[rend={small-caps}]{Leone} Stéphanie, « La costruzione di palazzo Pamphilj », dans \teiHi[rend={italic}]{Palazzo Pamphilj: ambasciata del Brasile a Roma}, Turin, U. Allemandi, 2016, p. 15-67. 
\end{teiBibl}
              \begin{teiBibl}[xmlid={bibl-392}]
\teiHi[rend={small-caps}]{Llorens} Juan Maria, « Miniaturas de Vincent Raymond en los manuscritos musicales de la Capilla Sixtina », dans \teiHi[rend={italic}]{Miscelànea Homenaje a Mons. Higinio Anglès}, Barcelone, Consejo Superior de Investigaciones científicas, 1958, 2. vol., t. 1, p. 475-494.
\end{teiBibl}
              \begin{teiBibl}[xmlid={bibl-393}]
\teiHi[rend={small-caps}]{Loskoutoff} Yvan, « Augures héraldiques de la papauté », dans Andreas Gottsmann, Pierantonio Piatti et Andreas E. Rehberg (dir.), \teiHi[rend={italic}]{Incorrupta monumenta Ecclesiam defendunt, Studi offerti a Sergio Pagano, Prefetto dell’Archivio Segreto Vaticano}, Cité du Vatican, Archivio Segreto Vaticano, 2018, 3 vol., t. 2, p. 976-994.
\end{teiBibl}
              \begin{teiBibl}[xmlid={bibl-394}]
\teiHi[rend={bold}]{—}, « Concessions héraldiques à la fin de la Renaissance d’après les recueils de brefs de saint Pie V (1566-1572), Grégoire \teiHi[rend={small-caps}]{XIII} (1572-1585) et Sixte-Quint (1585-1590) conservés aux Archives du Vatican », dans Yvan Loskoutoff (dir.), \teiHi[rend={italic}]{Héraldique et papauté. Moyen Âge-Temps modernes}, Mont-Saint-Aignan, PURH, 2020, p. 179-196.
\end{teiBibl}
              \begin{teiBibl}[xmlid={bibl-395}]
\teiHi[rend={small-caps}]{Minor} Heather Hyde, « Engraved in Porphyry, Printed on Paper: Piranesi and Lord Charlemont », dans Mario Bevilacqua, Heather Hyde Minor et Fabio Barri (dir.), \teiHi[rend={italic}]{The Serpent and the Stylus, Essays on G. B. Piranesi}, Ann Arbor, The University of Michigan Press, « Supplements to the Memoirs of the American Academy in Rome », 2006, p. 123-147.
\end{teiBibl}
              \begin{teiBibl}[xmlid={bibl-396}]
\teiHi[rend={small-caps}]{Moretti} Simona, « Manno di Bandino», dans \teiHi[rend={italic}]{DBI}, Rome, Istituto della Enciclopedia italiana, 2007, vol. 69, p. 109-110.
\end{teiBibl}
              \begin{teiBibl}[xmlid={bibl-397}]
\teiHi[rend={small-caps}]{Paravicini Bagliani} Agostino, « Boniface \teiHi[rend={small-caps}]{VIII} en images, Vision d’église et mémoire de soi », dans Dominic Olariu (dir.), \teiHi[rend={italic}]{Le portrait individuel. Réflexions autour d’une forme de représentation, }\teiHi[rend={small-caps italic}]{xiii}\teiHi[rend={italic sup}]{e}\teiHi[rend={italic}]{-}\teiHi[rend={small-caps italic}]{xv}\teiHi[rend={italic sup}]{e}\teiHi[rend={italic}]{ siècles}, Berne, Peter Lang, 2009, p. 65-81.
\end{teiBibl}
              \begin{teiBibl}[xmlid={bibl-398}]
\teiHi[rend={small-caps}]{Pastoureau} Michel, « L’emblématique des Farnèse », dans \teiHi[rend={italic}]{Le palais Farnèse}, Rome, École française de Rome, 1981, 2. vol., t. 1, p. 431-455.
\end{teiBibl}
              \begin{teiBibl}[xmlid={bibl-399}]
\teiHi[rend={small-caps}]{Pauwels} Yves, « Le thème de l’arc de triomphe dans l’architecture urbaine à la Renaissance, entre pouvoir politique et pouvoir religieux », dans Patrick Boucheron et Jean-Philippe Genet (dir.), \teiHi[rend={italic}]{Marquer la ville. Signes, traces, empreintes du pouvoir (}\teiHi[rend={small-caps italic}]{xiii}\teiHi[rend={italic sup}]{e}\teiHi[rend={italic}]{-}\teiHi[rend={small-caps italic}]{xvi}\teiHi[rend={italic sup}]{e}\teiHi[rend={italic}]{ siècle)}, Paris-Rome, Éditions de la Sorbonne-École française de Rome, 2014, p. 181-188.
\end{teiBibl}
              \begin{teiBibl}[xmlid={bibl-400}]
\teiHi[rend={small-caps}]{Petitmengin} Pierre, avec la collaboration de \teiHi[rend={small-caps}]{Fohlen} Jeannine, « I manoscritti latini della Vaticana. Uso, acquisizioni, classificazioni », dans Massimo Ceresa (dir.), \teiHi[rend={italic}]{Storia della Biblioteca apostolica vaticana, }vol. II :\teiHi[rend={italic}]{ La Biblioteca vaticana tra riforma cattolica, crescita delle collezioni e nuovo edificio (1535-1590)}, Cité du Vatican, BAV, 2012, p. 43-90.
\end{teiBibl}
              \begin{teiBibl}[xmlid={bibl-401}]
\teiHi[rend={small-caps}]{Pini} Raffaella, « La statua di Bonifacio \teiHi[rend={small-caps}]{VIII}, Manno da Siena e gli orefici a Bologna », dans Maria Andaloro (dir.), \teiHi[rend={italic}]{Le culture di Bonifacio VIII}, Rome, Istituto Storico Italiano per il Medioevo, 2006, p. 231-240.
\end{teiBibl}
              \begin{teiBibl}[xmlid={bibl-402}]
\teiHi[rend={small-caps}]{Popoff} Michel, « Le “capo dello scudo” dans l’héraldique florentine, \teiHi[rend={small-caps}]{xiii}\teiHi[rend={sup}]{e}\teiHi[rend={small-caps}]{-xvi}\teiHi[rend={sup}]{e }siècles », dans \teiHi[rend={italic}]{Brisures, augmentations et changements d'armoiries}, Bruxelles, Académie internationale d’héraldique, 1988, p. 251-255.
\end{teiBibl}
              \begin{teiBibl}[xmlid={bibl-403}]
\teiHi[rend={small-caps}]{Portoghesi} Paolo, « Borromini e Innocenzo X: architettura e politica pontificia », dans Alessandro Zuccari et Stefania Macioce (dir.), \teiHi[rend={italic}]{Innocenzo X Pamphili. Arte e potere a Roma nell’età barocca}, actes de colloque, Rome, novembre 1990 \lbrack{}2\teiHi[rend={sup}]{e} éd.\rbrack{}, Rome, Logart Press, 2001, p. 27-34.
\end{teiBibl}
              \begin{teiBibl}[xmlid={bibl-404}]
\teiHi[rend={small-caps}]{Prosperi Valenti Rodinò} Simonetta, « Cresti, Domenico, detto il Passignano », dans \teiHi[rend={italic}]{DBI}, Rome, Istituto della Enciclopedia italiana, 1984, vol. 30, p. 741-749.
\end{teiBibl}
              \begin{teiBibl}[xmlid={bibl-405}]
\teiHi[rend={small-caps}]{Raspe} Martin, « La Corsia Nuova dell’Ospedale del \teiHi[rend={small-caps}]{SS}. Salvatore ad Sancta Sanctorum e i suoi architetti Giacomo e Giovanni Battista Mola », dans Philine Helas et Patrizia Tosini (dir.), \teiHi[rend={italic}]{Tra Campidoglio e curia: l’Ospedale del SS. Salvatore ad Sancta Sanctorum tra Medioevo ed Età Moderna}, Cinisello Balsamo, Silvana, « Studi della Bibliotheca Hertziana, 11 », 2017, p. 137-147.
\end{teiBibl}
              \begin{teiBibl}[xmlid={bibl-406}]
\teiHi[rend={small-caps}]{Rehberg} Andreas, « Collecting and Drawing against Oblivion. Panvinio, Ceccarelli and Chacón and their Search for the Genealogical-Heraldic Identity of the Families of Rome », dans Noëlle-Laetitia Perret et Hans-Joachim Schmidt (dir.), \teiHi[rend={italic}]{Memories Lost in the Middle Ages. Collective Forgetting as an Alternative Procedure of Social Cohesion}, Turnhout, Brepols, 2023, p. 295-324.
\end{teiBibl}
              \begin{teiBibl}[xmlid={bibl-407}]
\teiHi[rend={small-caps}]{Ribouillault} Denis, « Jeux de mots et d’images : les frises peintes de l’appartement de Jules \teiHi[rend={small-caps}]{III} au Vatican (vers 1550-1553) », dans Antonella Fenech Kroke et Annick Lemoine (dir.), \teiHi[rend={italic}]{Frises peintes. Les décors des villas et palais au Cinquecento}, Paris-Rome, Somogy-Académie de France à Rome, 2016, p. 103-129.
\end{teiBibl}
              \begin{teiBibl}[xmlid={bibl-408}]
\teiHi[rend={small-caps}]{Rice} Louise, « Branding, Barberini-style », dans Maurizia Cicconi, Flaminia Gennari Santori, Sebastian Schütze (dir.), \teiHi[rend={italic}]{L’immagine sovrana. Urbano VIII e i Barberini}, catalogue d’exposition, Rome, Officina libraria, 2023, p. 245-250.
\end{teiBibl}
              \begin{teiBibl}[xmlid={bibl-409}]
\teiHi[rend={small-caps}]{Roli Guidetti} Clara, « Caccioli, Giovanni Battista », dans \teiHi[rend={italic}]{DBI}, Rome, Istituto della Enciclopedia italiana, 1973, vol. 16, p. 39-40.
\end{teiBibl}
              \begin{teiBibl}[xmlid={bibl-410}]
\teiHi[rend={small-caps}]{Rothstein} Bret, « The Rule of Metaphor and the Play of the Viewer in the Hours of Mary of Burgundy », dans Reindert Falkenburg, Walter S. Melion et Todd M. Richardson (dir.), \teiHi[rend={italic}]{Image and Imagination of the Religious Self in Late Medieval and Early Modern Europe}, Turnhout, Brepols, 2007, p. 237-275.
\end{teiBibl}
              \begin{teiBibl}[xmlid={bibl-411}]
\teiHi[rend={small-caps}]{Safarik} Eduard A., « L’Apoteosi di Romolo », dans \teiHi[rend={italic}]{Roma Palazzo Altieri. Le stanze al piano nobile dei cardinali Giovanni Battista e Paluzzo Altieri}, Milan, Franco Maria Ricci, 1999, p. 24-29. 
\end{teiBibl}
              \begin{teiBibl}[xmlid={bibl-412}]
\teiHi[rend={small-caps}]{Saffiotti Dale} Maria Francesca, « Raymond de Lodève, Vincent », dans \teiHi[rend={italic}]{Dizionario biografico dei miniatori italiani}, Milan, Sylvestre Bonnard, 2004, p. 899-902.
\end{teiBibl}
              \begin{teiBibl}[xmlid={bibl-413}]
\teiHi[rend={small-caps}]{Salvestrini} Francesco, « Il patrimonio fondiario del monastero di Vallombrosa fra \teiHi[rend={small-caps}]{XIII} e \teiHi[rend={small-caps}]{XVI} secolo: presenza e utilizzazione del bosco », dans Simonetta Cavaciocchi (dir.), \teiHi[rend={italic}]{L’uomo e la foresta, secc. XIII-XVIII}, Prato, Istituto Internazionale di Storia Economica Francesco Datini, 1996, p. 1057-1068.
\end{teiBibl}
              \begin{teiBibl}[xmlid={bibl-414}]
\teiHi[rend={bold}]{—}, « L’esperienza di Vallombrosa nella documentazione archivistica (secoli \teiHi[rend={small-caps}]{XI}-\teiHi[rend={small-caps}]{XVI}) », dans Giancarlo Andenna, Renata Salvarani (dir), \teiHi[rend={italic}]{La memoria dei chiostri: atti delle prime Giornate di studi medievali, Laboratorio di storia monastica dell’Italia settentrionale}, Castiglione delle Stiviere, 11-13 octobre 2001, Brescia, \teiHi[rend={small-caps}]{CESIMB}, 2002, p. 215-230.
\end{teiBibl}
              \begin{teiBibl}[xmlid={bibl-415}]
\teiHi[rend={bold}]{—}, « L’art et la magnificence “contre” le pouvoir du prince. L’abbé général Biagio Milanesi, l’ordre monastique de Vallombreuse, Laurent de Médicis et le pape Léon X », dans Élisabeth Crouzet-Pavan, Jean-Claude Maire Vigueur (dir), \teiHi[rend={italic}]{L’art au service du prince. Paradigme italien, expériences européennes (vers 1250-vers 1550)}, actes du colloque, Paris, Université de Paris-Sorbonne, Centre Roland-Mousnier, 17-19 octobre 2013, Rome, Viella, « Italia comunale e signorile, 8 », 2015, p. 253-279.
\end{teiBibl}
              \begin{teiBibl}[xmlid={bibl-416}]
\teiHi[rend={small-caps}]{Savorelli} Alessandro, « L’araldica per la storia: una fonte ausiliaria ? », dans Maria Pia Paoli (dir.), \teiHi[rend={italic}]{Nel laboratorio della storia: una guida alle fonti dell’età moderna}, Rome, Carocci, 2013, p. 289-315.
\end{teiBibl}
              \begin{teiBibl}[xmlid={bibl-417}]
\teiHi[rend={bold}]{—}, « Segni e simboli araldici nell’arte fiorentina dal Medioevo al Rinascimento », dans Maria Monica Donato et Daniela Parenti (dir.), \teiHi[rend={italic}]{Dal giglio al David. Arte civica a Firenze fra Medioevo e Rinascimento}, Florence-Milan, Giunti, 2013, p. 73-77.
\end{teiBibl}
              \begin{teiBibl}[xmlid={bibl-418}]
\teiHi[rend={small-caps}]{Schütze} Sebastian, « Urbano \teiHi[rend={small-caps}]{VIII} e il concetto di Palazzo Barberini. Alla ricerca di un primato culturale di rinascimentale memoria », dans Christoph Luitpold Frommel et Sebastian Schütze (dir.), \teiHi[rend={italic}]{Pietro da Cortona}, Milan, Electa, 1998 p. 86-97.
\end{teiBibl}
              \begin{teiBibl}[xmlid={bibl-419}]
\teiHi[rend={small-caps}]{Simonato} Lucia, « Il Palazzo Apostolico Vaticano », dans Francesco Buranelli (dir.), \teiHi[rend={italic}]{Vaticano Barocco. Arte, architettura e cerimoniale} \lbrack{}Cité du Vatican, Libreria Editrice Vaticana, 2012\rbrack{}, Milan, Jaca Book, 2014, chap. \teiHi[rend={small-caps}]{IV}.
\end{teiBibl}
              \begin{teiBibl}[xmlid={bibl-420}]
\teiHi[rend={small-caps}]{Von Flüe} Carlo, « Davidico, Lorenzo », \teiHi[rend={italic}]{DBI}, Rome, Istituto della Enciclopedia italiana, vol. 33, 1987, p. 157-160.
\end{teiBibl}
              \begin{teiBibl}[xmlid={bibl-421}]
\teiHi[rend={small-caps}]{Zuccari} Alessandro, « Borromini tra religiosità borromaica e cultura scientifica », dans Alessandro Zuccari et Stefania Macioce (dir.), \teiHi[rend={italic}]{Innocenzo X Pamphili. Arte e potere a Roma nell’età barocca} \lbrack{}2\teiHi[rend={sup}]{e} éd.\rbrack{}, actes de colloque, Rome, novembre 1990, Rome, Logart Press, 2001, p. 35-74.
\end{teiBibl}
            
\end{teiDiv}
            \begin{teiDiv}[type={section2},xmlid={div20}]

              \teiHead[xmlid={articles-de-periodiques-001}]{Articles de périodiques}
              \begin{teiBibl}[xmlid={bibl-422}]
\teiHi[rend={small-caps}]{Abromson} Morton C., « Clement \teiHi[rend={small-caps}]{VIII’s} Patronage of the Brothers Alberti », \teiHi[rend={italic}]{The Art Bulletin}, vol. \teiHi[rend={small-caps}]{LX}, n\teiHi[rend={sup}]{o} 3, septembre 1978, p. 531-547.
\end{teiBibl}
              \begin{teiBibl}[xmlid={bibl-423}]
\teiHi[rend={small-caps}]{Alonso} Carlos, « La chiesa e il convento di santa Prisca a Roma. Vicende storiche, 1535-1724 », \teiHi[rend={italic}]{Analecta Augustiniana}, n\teiHi[rend={sup}]{o} 58, 1995, p. 215-263.
\end{teiBibl}
              \begin{teiBibl}[xmlid={bibl-424}]
\teiHi[rend={small-caps}]{As}-\teiHi[rend={small-caps}]{Vijvers} Anne Margreet W., « The Grotesque Initials in the Alamire Choirbooks », \teiHi[rend={italic}]{Journal of the Alamire Foundation}, vol. 11, n\teiHi[rend={sup}]{o} 1‑2, 2019, p. 13‑46.
\end{teiBibl}
              \begin{teiBibl}[xmlid={bibl-425}]
\teiHi[rend={small-caps}]{Barbiche} Bernard et \teiHi[rend={small-caps}]{Dainville}-\teiHi[rend={small-caps}]{Barbiche} Ségolène de, « Les légats \teiHi[rend={italic}]{a latere} en France et leurs facultés aux \teiHi[rend={small-caps}]{xvi}\teiHi[rend={sup}]{e} et \teiHi[rend={small-caps}]{xvii}\teiHi[rend={sup}]{e} siècles », \teiHi[rend={italic}]{Archivum Historiae Pontificae}, n\teiHi[rend={sup}]{o} 23, 1985, p. 93-165.
\end{teiBibl}
              \begin{teiBibl}[xmlid={bibl-426}]
\teiHi[rend={small-caps}]{Bartoni} Laura, « L’allestimento delle collezioni nella villa pinciana: spazi, criteri, funzioni », \teiHi[rend={italic}]{Storia dell’arte – « Guercino nel Casino Ludovisi (1621-1623) »}, vol. 157, n\teiHi[rend={sup}]{o} 1, 2022, p. 99-117.
\end{teiBibl}
              \begin{teiBibl}[xmlid={bibl-427}]
\teiHi[rend={small-caps}]{Bernardini} Carla, « Per una storia dell’Appartamento del Legato, sede delle Collezioni Comunali d’Arte di Palazzo d’Accursio. Una ricerca in prospettiva museografica », \teiHi[rend={italic}]{Schede Umanistiche}, n\teiHi[rend={sup}]{o} 3, 1989, p. 81-92.
\end{teiBibl}
              \begin{teiBibl}[xmlid={bibl-428}]
\teiHi[rend={small-caps}]{Boislile} Arthur de, « Rapport de M. de Boislile sur une communication de M. Destandau », \teiHi[rend={italic}]{Bulletin historique et philologique du Comité des travaux historiques et scientifiques}, 1905, p. 356-359.
\end{teiBibl}
              \begin{teiBibl}[xmlid={bibl-429}]
\teiHi[rend={small-caps}]{Boiteux} Martine, « Le palais Farnèse : la représentation et l’identité », \teiHi[rend={italic}]{Mélanges de l’École française de Rome}, vol. 122, n\teiHi[rend={sup}]{o} 2, 2010, p. 311-338.
\end{teiBibl}
              \begin{teiBibl}[xmlid={bibl-430}]
\teiHi[rend={bold}]{—}, « Rituels funéraires pontificaux », \teiHi[rend={italic}]{Studi romani}, vol. 60, n\teiHi[rend={sup}]{o} 1-4, 2012, p. 31-56.
\end{teiBibl}
              \begin{teiBibl}[xmlid={bibl-431}]
\teiHi[rend={small-caps}]{Borello} Benedetta, « Alleanze matrimoniali e mobilità sociale e geografica. Il caso dei Pamphilj (\teiHi[rend={small-caps}]{XV}-\teiHi[rend={small-caps}]{XVII} secolo) », \teiHi[rend={italic}]{Mélanges de l’École française de Rome}, vol. 115, n\teiHi[rend={sup}]{o} 1, 2003, p. 345-366.
\end{teiBibl}
              \begin{teiBibl}[xmlid={bibl-432}]
\teiHi[rend={small-caps}]{Borghetti} Vincenzo, « Music and the Representation of Princely Power in the Fifteenth and Sixteenth Century », \teiHi[rend={italic}]{Acta Musicologica}, vol. 80, n\teiHi[rend={sup}]{o} 2, 2008, p. 179‑214.
\end{teiBibl}
              \begin{teiBibl}[xmlid={bibl-433}]
\teiHi[rend={small-caps}]{Bouyé} Édouard, « Les armoiries pontificales à la fin du \teiHi[rend={small-caps}]{xiii}\teiHi[rend={sup}]{e} siècle : construction d’une campagne de communication », \teiHi[rend={italic}]{Médiévales}, n\teiHi[rend={sup}]{o} 44, 2003, p. 177 et 178, http://medievales.revues.org/938.
\end{teiBibl}
              \begin{teiBibl}[xmlid={bibl-434}]
\teiHi[rend={small-caps}]{Bühler} Curt F., « Studies in the Early Editions of the “Fiore Di Virtù” », \teiHi[rend={italic}]{The Papers of the Bibliographical Society of America}, vol. 49, n\teiHi[rend={sup}]{o} 4, 1955, p. 315-339.
\end{teiBibl}
              \begin{teiBibl}[xmlid={bibl-435}]
\teiHi[rend={small-caps}]{Busiri Vici} Andrea, « L’erudito della corte russa del Settecento, Ivan Ivanovitch Schouvaloff, ed i suoi rapporti con Roma », \teiHi[rend={italic}]{L’Urbe}, n\teiHi[rend={sup}]{o} 39, 1976, p. 39-52.
\end{teiBibl}
              \begin{teiBibl}[xmlid={bibl-436}]
\teiHi[rend={small-caps}]{Campbell} Tom, « Pope Leo X’s Consistorial “letto de paramento” and the Boughton House Cartoons », \teiHi[rend={italic}]{The Burlington Magazine}, vol. 138, n\teiHi[rend={sup}]{o} 1120, 1996, p. 436-445.
\end{teiBibl}
              \begin{teiBibl}[xmlid={bibl-437}]
\teiHi[rend={small-caps}]{Carlevaris} Laura, « La Sala Clementina: La costruzione pittorica delle pareti dallo schema compositivo alla griglia prospettica », \teiHi[rend={italic}]{Bollettino-Monumenti, Musei e Gallerie Pontificie}, n\teiHi[rend={sup}]{o} 21, 2001, p. 320-361.
\end{teiBibl}
              \begin{teiBibl}[xmlid={bibl-438}]
\teiHi[rend={small-caps}]{Cavazzini} Patrizia, « Towards a Chronology of Agostino Tassi », \teiHi[rend={italic}]{The Burlington Magazine}, n\teiHi[rend={sup}]{o} 144, 2002, p. 396-408\teiHi[rend={italic}]{.}
\end{teiBibl}
              \begin{teiBibl}[xmlid={bibl-439}]
\teiHi[rend={small-caps}]{Cavicchioli} Sonia, « Il cardinale Pietro Vidoni e l’omaggio ad Alessandro \teiHi[rend={small-caps}]{VII} Chigi nell’appartamento del legato a Bologna (1665). Una decorazione politica », \teiHi[rend={italic}]{Studi di Storia dell’Arte}, n\teiHi[rend={sup}]{o} 29, 2018, p. 167-178.
\end{teiBibl}
              \begin{teiBibl}[xmlid={bibl-440}]
\teiHi[rend={small-caps}]{Coupe} C., « An Old Picture », \teiHi[rend={italic}]{The Month}, n\teiHi[rend={sup}]{o} 84, 1895, p. 229–242.
\end{teiBibl}
              \begin{teiBibl}[xmlid={bibl-441}]
\teiHi[rend={small-caps}]{De Gregori} Luigi, « Piazza Navona prima d’Innocenzo X », \teiHi[rend={italic}]{Roma}, n\teiHi[rend={sup}]{o} 4, 1926, p. 14-25 et 97-116.
\end{teiBibl}
              \begin{teiBibl}[xmlid={bibl-442}]
\teiHi[rend={small-caps}]{Dang} Benjamin, « Le bouffon, un héraut du vide », \teiHi[rend={italic}]{Les chantiers de la création. Revue pluridisciplinaire en lettres, langues, arts et civilisations}, n\teiHi[rend={sup}]{o} 13, 12 juillet 2021, \teiRef[target={https://doi.org/10.4000/lcc.4119}]{https://doi.org/10.4000/lcc.4119}.
\end{teiBibl}
              \begin{teiBibl}[xmlid={bibl-443}]
\teiHi[rend={small-caps}]{Deonna} Waldemar. « L’abeille et le roi », \teiHi[rend={italic}]{Revue belge d’archéologie et d’histoire de l’art}, vol. 25, 1956, p. 105‑131.
\end{teiBibl}
              \begin{teiBibl}[xmlid={bibl-444}]
\teiHi[rend={small-caps}]{Elam} Caroline, « Art and Diplomacy in Renaissance Florence », \teiHi[rend={italic}]{RSA Journal}, vol. 136, n\teiHi[rend={sup}]{o} 5387, 1988, p. 813‑826.
\end{teiBibl}
              \begin{teiBibl}[xmlid={bibl-445}]
\teiHi[rend={small-caps}]{Ewart Witcombe} Christopher L. C., « A new fresco by Cherubino Alberti in the Vatican », \teiHi[rend={italic}]{Notes in the History of Art}, vol. 4, n\teiHi[rend={sup}]{o} 1, 1984, p. 12-16.
\end{teiBibl}
              \begin{teiBibl}[xmlid={bibl-446}]
\teiHi[rend={bold}]{—}, « Cesare Ripa and the Sala Clementina », \teiHi[rend={italic}]{Journal of the Warburg and Courtauld Institutes}, n\teiHi[rend={sup}]{o} 55, 1992, p. 277-282.
\end{teiBibl}
              \begin{teiBibl}[xmlid={bibl-447}]
\teiHi[rend={small-caps}]{Fagiolo} Marcello, « Borromini in Laterano. Il Nuovo Tempio per il Concilio universale », \teiHi[rend={italic}]{L’arte}, n\teiHi[rend={sup}]{o} 13, 1971, p. 1-46.
\end{teiBibl}
              \begin{teiBibl}[xmlid={bibl-448}]
\teiHi[rend={small-caps}]{Federici} Fabrizio, « Il trattato \teiHi[rend={italic}]{Delle memorie sepolcrali} del cavalier Francesco Gualdi: un collezionista del Seicento e le testimonianze figurative medievali », \teiHi[rend={italic}]{Prospettiva}, n\teiHi[rend={sup}]{o} 110-111, avril-juin 2003, p. 149-159.
\end{teiBibl}
              \begin{teiBibl}[xmlid={bibl-449}]
\teiHi[rend={bold}]{—}, « Battaglie per la tutela nella Roma barocca: Francesco Gualdi e la difesa delle “memorie antiche” », \teiHi[rend={italic}]{Studi Romani}, vol. \teiHi[rend={small-caps}]{LXII}, n\teiHi[rend={sup}]{o} 1-4, 2014, p. 149-172.
\end{teiBibl}
              \begin{teiBibl}[xmlid={bibl-450}]
\teiHi[rend={small-caps}]{Feinblatt} Ebria, « Giovan Battista Cavalcaselle: a Drawing identified », \teiHi[rend={italic}]{Master Drawings}, vol. V, n\teiHi[rend={sup}]{o} 2, 1967, p. 569-576.
\end{teiBibl}
              \begin{teiBibl}[xmlid={bibl-451}]
\teiHi[rend={small-caps}]{Franzoni} Claudio et \teiHi[rend={small-caps}]{Tempesta} Alessandra, « Il museo di Francesco Gualdi nella Roma del Seicento tra raccolta privata ed esibizione pubblica », \teiHi[rend={italic}]{Bollettino d’Arte}, vol. \teiHi[rend={small-caps}]{LXXVII}, n\teiHi[rend={sup}]{o} 73, 1992, p. 1-42.
\end{teiBibl}
              \begin{teiBibl}[xmlid={bibl-452}]
\teiHi[rend={small-caps}]{Fumagalli} Camillo, « Dello stemma di Papa Giovanni \teiHi[rend={small-caps}]{XXIII}. La genealogia dei Roncalli, Ronco e Roncaglia », \teiHi[rend={italic}]{Atti dell’Ateneo di Scienze, Lettere ed Arti di Bergamo}, n\teiHi[rend={sup}]{o} 34, 1968-1969, p. 11-22.
\end{teiBibl}
              \begin{teiBibl}[xmlid={bibl-453}]
\teiHi[rend={small-caps}]{Fusconi} Giulia, « Da Bartoli a Piranesi: spigolature dai Codici Ottoboniani Latini della Raccolta Ghezzi », \teiHi[rend={italic}]{Xenia Antiqua}, n\teiHi[rend={sup}]{o} 3, 1994, p. 145-172.
\end{teiBibl}
              \begin{teiBibl}[xmlid={bibl-454}]
\teiHi[rend={small-caps}]{Gentile} Luisa Clotilde, « La diffusion du capo dell'Impero dans l’Italie du nord », \teiHi[rend={italic}]{Armas e troféus}, n\teiHi[rend={sup}]{o} 20, 2018, p. 109-138.
\end{teiBibl}
              \begin{teiBibl}[xmlid={bibl-455}]
\teiHi[rend={small-caps}]{Gnoli} Umberto, « \teiHi[rend={italic}]{Nobiltà romana e marchi mercantili nel Medioevo »}, \teiHi[rend={italic}]{L’Urbe}, vol. \teiHi[rend={small-caps}]{VIII}, n\teiHi[rend={sup}]{o} 11-12, 1943, p. 15-32.
\end{teiBibl}
              \begin{teiBibl}[xmlid={bibl-456}]
\teiHi[rend={small-caps}]{Harvey} J. H., « The Wilton Diptych: A Re-examination », \teiHi[rend={italic}]{Archaeologia}, n\teiHi[rend={sup}]{o} 98, 1961, p. 1-28.
\end{teiBibl}
              \begin{teiBibl}[xmlid={bibl-457}]
\teiHi[rend={small-caps}]{Hermann Fiore} Kristina, « Giovanni Albertis Kunst und Wissenschaft der Quadratur. Eine Allegorie in der Sala Clementina des Vatikan », \teiHi[rend={italic}]{Mitteilungen des Kunsthistorisches Institutes in Florens}, n\teiHi[rend={sup}]{o} 22, 1978, p. 61-84.
\end{teiBibl}
              \begin{teiBibl}[xmlid={bibl-458}]
\teiHi[rend={small-caps}]{Klein} Robert, « La théorie de l’expression figurée dans les traités italiens sur les \teiHi[rend={italic}]{Imprese}, 1555-1612 », \teiHi[rend={italic}]{Bibliothèque d’humanisme et Renaissance}, vol. 19, n\teiHi[rend={sup}]{o} 2, 1957, p. 320-341.
\end{teiBibl}
              \begin{teiBibl}[xmlid={bibl-459}]
\teiHi[rend={small-caps}]{Kruse} Jeremy, « Hunting, Magnificence and the Court of Leo X », \teiHi[rend={italic}]{Renaissance Studies}, vol. 7, n\teiHi[rend={sup}]{o} 3, 1993, p. 243‑257.
\end{teiBibl}
              \begin{teiBibl}[xmlid={bibl-460}]
\teiHi[rend={small-caps}]{Lanzoni} Francesco, « I titoli presbiterali di Roma antica nella storia e nella leggenda », \teiHi[rend={italic}]{Rivista di archeologia cristiana}, n\teiHi[rend={sup}]{o} 2, 1925, p. 247-250.
\end{teiBibl}
              \begin{teiBibl}[xmlid={bibl-461}]
\teiHi[rend={small-caps}]{Larraz} Camille, « La réception tardive de l’art bellifontain dans le Comtat Venaissin : le cas du graveur Jean Bœuf (\teiHi[rend={small-caps}]{xvi}\teiHi[rend={sup}]{e}-\teiHi[rend={small-caps}]{xvii}\teiHi[rend={sup}]{e} siècles) », \teiHi[rend={italic}]{Nouvelles de l’estampe}, n\teiHi[rend={sup}]{o} 268, décembre 2022, \teiRef[target={https://doi.org/10.4000/estampe.3641}]{https://doi.org/10.4000/estampe.3641}.
\end{teiBibl}
              \begin{teiBibl}[xmlid={bibl-462}]
\teiHi[rend={small-caps}]{Lartigaut} Jean, « Sur la date de construction de la collégiale de Castelnau-Bretenoux », \teiHi[rend={italic}]{Bulletin de la Société des études du Lot}, vol. XCVIII, 1977, p. 167-168.
\end{teiBibl}
              \begin{teiBibl}[xmlid={bibl-463}]
\teiHi[rend={small-caps}]{Lavaur De Sainte}-\teiHi[rend={small-caps}]{Fortunade} Henri de, « Documents sur la baronnie de Castelnau de Bretenoux », \teiHi[rend={italic}]{Bulletin de la Société scientifique},\teiHi[rend={italic}]{ historique et archéologique de la Corrèze}, vol. 34, n\teiHi[rend={sup}]{o} 1, 1912, p. 353-403.
\end{teiBibl}
              \begin{teiBibl}[xmlid={bibl-464}]
\teiHi[rend={small-caps}]{Leone} Stéphanie, « Cardinal Pamphilj Builds a Palace: Self-Representation and Familial Ambition in Seventeenth-Century Rome », \teiHi[rend={italic}]{Journal of the Society of Architectural Historians}, vol. 63, n\teiHi[rend={sup}]{o} 4, 2004, p. 440-471.
\end{teiBibl}
              \begin{teiBibl}[xmlid={bibl-465}]
\teiHi[rend={small-caps}]{Lévi}-\teiHi[rend={small-caps}]{Strauss} Claude, « L’efficacité symbolique », \teiHi[rend={italic}]{Revue d’histoire des religions}, vol. 135, n\teiHi[rend={sup}]{o} 1, 1949, p. 5-27.
\end{teiBibl}
              \begin{teiBibl}[xmlid={bibl-466}]
\teiHi[rend={small-caps}]{Loskoutoff} Yvan, « Le symbolisme des \teiHi[rend={italic}]{palle} médicéennes à la villa Madama », \teiHi[rend={italic}]{Journal des savants}, n\teiHi[rend={sup}]{o} 2, 2001, p. 351-391.
\end{teiBibl}
              \begin{teiBibl}[xmlid={bibl-467}]
\teiHi[rend={bold}]{—}, « Hercule et Atlas à Rome puis à Modène, Un présent diplomatique du cardinal Francesco degli Albizzi au cardinal Mazarin », \teiHi[rend={italic}]{Dix-septième siècle}, n\teiHi[rend={sup}]{o} 241, octobre-décembre 2008, p. 677-708.
\end{teiBibl}
              \begin{teiBibl}[xmlid={bibl-468}]
\teiHi[rend={bold}]{—}, « Regrets et pasquinades, Le cycle de Jules \teiHi[rend={small-caps}]{III} dans le recueil de Du Bellay », \teiHi[rend={italic}]{Revue d’histoire littéraire de la France}, vol. 123, n\teiHi[rend={sup}]{o} 1, 2023, p. 5-24.
\end{teiBibl}
              \begin{teiBibl}[xmlid={bibl-469}]
\teiHi[rend={small-caps}]{Macioce} Stefania, « Per un’introduzione al pontificato di Clemente \teiHi[rend={small-caps}]{VIII} Aldobrandini (1592-1605): politica e committenza », \teiHi[rend={italic}]{Bollettino della Deputazione di Storia Patria per l’Umbria}, vol. 111, n\teiHi[rend={sup}]{o} 1-2, 2014, p. 709-726 et p. 1412-1413.
\end{teiBibl}
              \begin{teiBibl}[xmlid={bibl-470}]
\teiHi[rend={small-caps}]{Mattingly} Harold, « “Fel. Temp. Reparatio” », \teiHi[rend={italic}]{The Numismatic Chronicle and Journal of the Royal Numismatic Society}, vol. 5, n\teiHi[rend={sup}]{o} 13, 1933, p. 182-202.
\end{teiBibl}
              \begin{teiBibl}[xmlid={bibl-471}]
\teiHi[rend={small-caps}]{Mc Phee} Sarah, « A new sketch-book by Filippo Juvarra », \teiHi[rend={italic}]{The Burlington Magazine}, n\teiHi[rend={sup}]{o} 135, 1993, p. 346-350.
\end{teiBibl}
              \begin{teiBibl}[xmlid={bibl-472}]
\teiHi[rend={bold}]{—}, « Filippo Juvarra: a new volume of drawings from the Roman years », \teiHi[rend={italic}]{Studi piemontesi}, n\teiHi[rend={sup}]{o} 22, 1993, p. 377-383.
\end{teiBibl}
              \begin{teiBibl}[xmlid={bibl-473}]
\teiHi[rend={small-caps}]{Medvedkova} Olga, «\teiRef[target={https://fr.wikipedia.org/wiki/Š}]{ Š}uvalov à Rome (1765-1774), Histoire d’une dédicace », \teiHi[rend={italic}]{Cahiers du monde russe et soviétique}, vol. 52, n\teiHi[rend={sup}]{o} 1, 2011, p. 45-73.
\end{teiBibl}
              \begin{teiBibl}[xmlid={bibl-474}]
\teiHi[rend={small-caps}]{Messa} Pietro, « Il motto di Giovanni XXIII: obbedienza e pace », \teiHi[rend={italic}]{Forma Sororum}, n\teiHi[rend={sup}]{o} 50, 2013, p. 250-253.
\end{teiBibl}
              \begin{teiBibl}[xmlid={bibl-475}]
\teiHi[rend={small-caps}]{Moureau} Emmanuel, « Les collégiales du Quercy et la Curie pontificale du \teiHi[rend={small-caps}]{xiv}\teiHi[rend={sup}]{e} au \teiHi[rend={small-caps}]{xvi}\teiHi[rend={sup}]{e} siècle », \teiHi[rend={italic}]{Bulletin de la société archéologique et historique de Tarn-et-Garonne}, vol. 145, 2020, p. 19-48.
\end{teiBibl}
              \begin{teiBibl}[xmlid={bibl-476}]
\teiHi[rend={small-caps}]{Nassiet} Michel, « Signes de parenté signes de seigneurie : un système idéologique », \teiHi[rend={italic}]{Mémoires de la Société d’histoire et d’archéologie de Bretagne}, n\teiHi[rend={sup}]{o} 67, 1991, p. 175-232.
\end{teiBibl}
              \begin{teiBibl}[xmlid={bibl-477}]
\teiHi[rend={small-caps}]{Nova} Alessandro, « Bartolomeo Ammannati e Prospero Fontana a Palazzo Firenze. Architettura e emblemi per Giulio \teiHi[rend={small-caps}]{III} Del Monte », \teiHi[rend={italic}]{Ricerche di Storia dell’arte}, n\teiHi[rend={sup}]{o} 21, 1983, p. 53-76.
\end{teiBibl}
              \begin{teiBibl}[xmlid={bibl-478}]
\teiHi[rend={small-caps}]{O’Bryan} Robin, « The Medici Pope, Curative Puns, and the Panacean Dwarf in the Sala di Costantino », \teiHi[rend={italic}]{Secac Review}, vol. 16, n\teiHi[rend={sup}]{o} 5, p. 590-607.
\end{teiBibl}
              \begin{teiBibl}[xmlid={bibl-479}]
\teiHi[rend={small-caps}]{Pansier} Pierre, « Les débuts de l’imprimerie à Avignon (\teiHi[rend={small-caps}]{xv}\teiHi[rend={sup}]{e} et \teiHi[rend={small-caps}]{xvi}\teiHi[rend={sup}]{e} siècles) », \teiHi[rend={italic}]{Mémoires de l’Académie de Vaucluse}, vol. XIX, n\teiHi[rend={sup}]{o} 2, 1919, p. 153-179.
\end{teiBibl}
              \begin{teiBibl}[xmlid={bibl-480}]
\teiHi[rend={small-caps}]{Panza} Pierluigi, « L’araldica e Giovan Battista Piranesi », \teiHi[rend={italic}]{Rivista del Collegio araldico}, 2017, p. 142-154.
\end{teiBibl}
              \begin{teiBibl}[xmlid={bibl-481}]
\teiHi[rend={small-caps}]{Pastoureau} Michel, « Aux origines des hachures héraldiques (\teiHi[rend={small-caps}]{xv}\teiHi[rend={sup}]{e}-\teiHi[rend={small-caps}]{xvii}\teiHi[rend={sup}]{e} siècle) », \teiHi[rend={italic}]{Revue française d’héraldique et de sigillographie}, n\teiHi[rend={sup}]{o} 65, 1995, p. 21-31.
\end{teiBibl}
              \begin{teiBibl}[xmlid={bibl-482}]
\teiHi[rend={small-caps}]{Perini Folesani} Giovanna, « Pittura e poesia, arti sorelle nella Bologna rinascimentale e barocca », \teiHi[rend={italic}]{Letteratura e arte}, n\teiHi[rend={sup}]{o} 18, 2020, p. 67-80.
\end{teiBibl}
              \begin{teiBibl}[xmlid={bibl-483}]
\teiHi[rend={small-caps}]{Pinelli} Antonio, « \teiHi[rend={italic}]{La philosophie des images}. Emblemi e imprese fra manierismo e barocco », \teiHi[rend={italic}]{Ricerche di storia dell’arte}, n\teiHi[rend={sup}]{o} 1-2, 1976, p. 3-28.
\end{teiBibl}
              \begin{teiBibl}[xmlid={bibl-484}]
\teiHi[rend={small-caps}]{Preimesberger} Rudolf, « Obeliscus Pamphilius. Beitrage zur Vorgeschichte und Ikonographie des Vierstromebrunnes auf Piazza Navona », \teiHi[rend={italic}]{Münchner Jahrbuch der Bildenden Kunst}, n\teiHi[rend={sup}]{o} 25, 1974, p. 77-162.
\end{teiBibl}
              \begin{teiBibl}[xmlid={bibl-485}]
\teiHi[rend={bold}]{—}, « “Pontifex Romanus per Aeneam Praesignatus”, Die Galleria Pamphilj und ihre Fresken », \teiHi[rend={italic}]{Romisches Jahrbuch fur Kunstgeschichte}, n\teiHi[rend={sup}]{o} 16, 1976, p. 221-287.
\end{teiBibl}
              \begin{teiBibl}[xmlid={bibl-486}]
\teiHi[rend={small-caps}]{Ravaioli} Davide, « Pitture murali di Francesco Lola sulla facciata del Palazzo Comunale », \teiHi[rend={italic}]{Arte a Bologna}, n\teiHi[rend={sup}]{o} 6, 2007, p. 133-136.
\end{teiBibl}
              \begin{teiBibl}[xmlid={bibl-487}]
\teiHi[rend={small-caps}]{Righini} Davide, « Iconografie celebrative. Gli emblemi e i miti raffigurati nelle sale di rappresentanza del Palazzo Pubblico di Bologna (secoli \teiHi[rend={small-caps}]{XV}-\teiHi[rend={small-caps}]{XVIII}) », \teiHi[rend={italic}]{Arte a Bologna. Bollettino dei Musei Civici d’Arte Antica}, n\teiHi[rend={sup}]{o} 7-8, 2010-2011, p. 128-140.
\end{teiBibl}
              \begin{teiBibl}[xmlid={bibl-488}]
\teiHi[rend={small-caps}]{Roland} Martin et \teiHi[rend={small-caps}]{Zajic} Andreas, « Les chartes médiévales enluminées dans les pays d’Europe centrale », \teiHi[rend={italic}]{Bibliothèque de l’École des chartes}, n\teiHi[rend={sup}]{o} 169, 2011, p. 151-253.
\end{teiBibl}
              \begin{teiBibl}[xmlid={bibl-489}]
\teiHi[rend={small-caps}]{Rolet} Anne, « Aux sources de l’emblème : blasons et devises », \teiHi[rend={italic}]{Littérature}, vol. 145, n\teiHi[rend={sup}]{o} 1, 2007, p. 53-78.
\end{teiBibl}
              \begin{teiBibl}[xmlid={bibl-490}]
\teiHi[rend={small-caps}]{Romani} Valentino, « Per la storia dell’editoria italiana del Cinquecento: le edizioni romane “In aedibus sanctae Brigidae (1553-1557)” », \teiHi[rend={italic}]{Rara Volumina}, n\teiHi[rend={sup}]{o} 1, 1998, p. 23-36.
\end{teiBibl}
              \begin{teiBibl}[xmlid={bibl-491}]
\teiHi[rend={small-caps}]{Ronzani} Rocco, « Note di sigillografia dell’ordine di sant’Agostino », \teiHi[rend={italic}]{Subsidia Augustiniana Italica} \teiHi[rend={small-caps}]{II}, n\teiHi[rend={sup}]{o} 1, 2009, p. 22-25.
\end{teiBibl}
              \begin{teiBibl}[xmlid={bibl-492}]
\teiHi[rend={bold}]{—}, « Nota storico-iconografica sulla committenza del cardinale Benedetto Giustiniani in Santa Prisca all’Aventino », \teiHi[rend={italic}]{Studi Romani}, à paraître.
\end{teiBibl}
              \begin{teiBibl}[xmlid={bibl-493}]
\teiHi[rend={small-caps}]{Rossi} Giovanni Battista de, « Della casa d’Aquila e Prisca all’Aventino », \teiHi[rend={italic}]{Bullettino di archeologia cristiana}, vol. 3, 1867, p. 44-47.
\end{teiBibl}
              \begin{teiBibl}[xmlid={bibl-494}]
\teiHi[rend={bold}]{—}, « Priscilla e gli Acilii Glabrioni. § \teiHi[rend={small-caps}]{VIII}, Aquila e Prisca e gli Acilii Glabrioni », \teiHi[rend={italic}]{Bullettino di archeologia cristiana}, n\teiHi[rend={sup}]{o} 6, 1888-1889, p. 128-133.
\end{teiBibl}
              \begin{teiBibl}[xmlid={bibl-495}]
\teiHi[rend={small-caps}]{Rouchon} Olivier, « Rituels publics, souveraineté et identité citadine. Les cérémonies d’entrée en Avignon (\teiHi[rend={small-caps}]{xvi}\teiHi[rend={sup}]{e}-\teiHi[rend={small-caps}]{xvii}\teiHi[rend={sup}]{e} siècles) », \teiHi[rend={italic}]{Cahiers de la Méditerranée}, n\teiHi[rend={sup}]{o} 77, 2008, p. 39-59, https://journals.openedition.org/cdlm/4362.
\end{teiBibl}
              \begin{teiBibl}[xmlid={bibl-496}]
\teiHi[rend={small-caps}]{Salvestrini} Francesco, « L’apporto dei Vallombrosani e dei Camaldolesi all’edificazione della marina toscana (seconda metà del XVII-anni ’20 del XVIII secolo) », \teiHi[rend={italic}]{Archivio Storico Italiano}, vol. 156, n\teiHi[rend={sup}]{o} 2, 1998, p. 307-329.
\end{teiBibl}
              \begin{teiBibl}[xmlid={bibl-497}]
—, « La storiografia sul movimento e sull’Ordine monastico di Vallombrosa », \teiHi[rend={italic}]{Quaderni Medievali}, n\teiHi[rend={sup}]{o} 53, 2002, p. 294-323.
\end{teiBibl}
              \begin{teiBibl}[xmlid={bibl-498}]
\teiHi[rend={small-caps}]{Savorelli} Alessandro, « “Les fleurs de lys y sont partout”. Origine et signification du chef d’Anjou en Italie (\teiHi[rend={small-caps}]{xiii}-\teiHi[rend={small-caps}]{xiv}\teiHi[rend={sup}]{e} siècles) », \teiHi[rend={italic}]{Armas e troféus}, vol. 20, 2018, p. 93-107.
\end{teiBibl}
              \begin{teiBibl}[xmlid={bibl-499}]
\teiHi[rend={small-caps}]{Schwarzenberg} Erkinger, « Glovis, impresa di Giuliano de’ Medici », \teiHi[rend={italic}]{Mitteilungen des Kunsthistorischen Institutes in Florenz}, n\teiHi[rend={sup}]{o} 39, 1995, p. 140-166.
\end{teiBibl}
              \begin{teiBibl}[xmlid={bibl-500}]
\teiHi[rend={small-caps}]{Sickel} Lothar, « L’antica cappella Casali in Sant’Agostino a Roma in un disegno di Taddeo Kuntz: una testimonianza per un affresco scomparso del Pintoricchio? », \teiHi[rend={italic}]{Bollettino d’Arte}, n\teiHi[rend={sup}]{o} 43, 2020, p. 59-78.
\end{teiBibl}
              \begin{teiBibl}[xmlid={bibl-501}]
\teiHi[rend={small-caps}]{Tacconi} Marica, « Appropriating the Instruments of Worship: The 1512 Medici Restoration and the Florentine Cathedral », \teiHi[rend={italic}]{Renaissance Quarterly}, vol. 56, n\teiHi[rend={sup}]{o} 2, 2003, p. 333-376.
\end{teiBibl}
              \begin{teiBibl}[xmlid={bibl-502}]
\teiHi[rend={small-caps}]{Van Der Heide} Klaas, « How Many Paths Must a Choirbook Tread Before it Reaches the Pope? », \teiHi[rend={italic}]{Journal of the Alamire Foundation}, vol. 11, n\teiHi[rend={sup}]{o} 1‑2, 2019, p. 47‑70.
\end{teiBibl}
              \begin{teiBibl}[xmlid={bibl-503}]
\teiHi[rend={small-caps}]{Visceglia} Maria Antonietta, « Giubilei tra pace e guerre (1625-1650) », \teiHi[rend={italic}]{Roma moderna e contemporanea}, vol. V, n\teiHi[rend={sup}]{o} 2-3, 1997, p. 431-474.
\end{teiBibl}
              \begin{teiBibl}[xmlid={bibl-504}]
\teiHi[rend={small-caps}]{Vitsthum} Walter, « A comment on the Iconography of Pietro da Cortona’s Barberini Ceiling », \teiHi[rend={italic}]{The Burlington Magazine}, vol. 103, n\teiHi[rend={sup}]{o} 703, octobre 1961, p. 427-433, repris dans Susan M. Dixon (dir.), \teiHi[rend={italic}]{Italian Baroque Art. Blackwell Anthologies in Art History}, Malden-Oxford, Carlton, Blackwell Publishing, 2008, p. 201-208.
\end{teiBibl}
              \begin{teiBibl}[xmlid={bibl-505}]
\teiHi[rend={small-caps}]{Volpi} Caterina, « La collezione Ludovisi come “Scuola del mondo”. Il ruolo di Agucchi, Marino, Valesio, Guercino e Domenichino », \teiHi[rend={italic}]{Storia del’arte – « Guercino nel Casino Ludovisi (1621-1623) »}, vol. 157, n\teiHi[rend={sup}]{o} 1, 2022, p. 118-133.
\end{teiBibl}
              \begin{teiBibl}[xmlid={bibl-506}]
\teiHi[rend={small-caps}]{Weber} Christof, « Papstgeschichte und Genealogie », \teiHi[rend={italic}]{Römische Quartalschrift}, n\teiHi[rend={sup}]{o} 84, 1989, p. 331-400.
\end{teiBibl}
              \begin{teiBibl}[xmlid={bibl-507}]
\teiHi[rend={small-caps}]{Zega} Raffaella, « La “Raccolta di varie targhe di Roma...” nell’attività di Filippo Juvarra », \teiHi[rend={italic}]{Archivum Arcis}, n\teiHi[rend={sup}]{o} 2, 1989, p. 10-51.
\end{teiBibl}
            
\end{teiDiv}
            \begin{teiDiv}[type={section2},xmlid={div21}]

              \teiHead[xmlid={theses-et-memoires-001}]{Thèses et mémoires}
              \begin{teiBibl}[xmlid={bibl-508}]
\teiHi[rend={small-caps}]{Abromson} Morton C., \teiHi[rend={italic}]{Painting in Rome during the Reign of Clement VIII (1592-1607), a Documentary Study}, thèse de doctorat, Columbia University, 1976.
\end{teiBibl}
              \begin{teiBibl}[xmlid={bibl-509}]
\teiHi[rend={small-caps}]{Beneš} Mirka, \teiHi[rend={italic}]{The Villa Pamphilj (1630,1670): Family, Gardens, and Land in Papal Rome}, thèse de doctorat, Université de Yale, 1989.
\end{teiBibl}
              \begin{teiBibl}[xmlid={bibl-510}]
\teiHi[rend={small-caps}]{Broucke} Paul-François, \teiHi[rend={italic}]{Les prééminences héraldiques de la cathédrale de Quimper}, mémoire de master I, Université de Bretagne Occidentale, 2010.
\end{teiBibl}
              \begin{teiBibl}[xmlid={bibl-511}]
\teiHi[rend={bold}]{—}, \teiHi[rend={italic}]{La cathédrale de Saint-Pol de Léon : le chantier des parties orientales du }\teiHi[rend={small-caps italic}]{xi}\teiHi[rend={italic sup}]{e}\teiHi[rend={italic}]{ au }\teiHi[rend={small-caps italic}]{xvi}\teiHi[rend={italic sup}]{e}\teiHi[rend={italic}]{ siècle}, mémoire de master \teiHi[rend={small-caps}]{II}, Université de Bretagne Occidentale, 2012.
\end{teiBibl}
              \begin{teiBibl}[xmlid={bibl-512}]
\teiHi[rend={small-caps}]{Hablot} Laurent, \teiHi[rend={italic}]{Affinités héraldiques}. \teiHi[rend={italic}]{Concessions, augmentations et partages d’armoiries en Europe au Moyen Âge}, mémoire d’habilitation à diriger des recherches, Michel Pastoureau (dir.), École pratique des hautes études, 2015.
\end{teiBibl}
              \begin{teiBibl}[xmlid={bibl-513}]
\teiHi[rend={small-caps}]{Harwood Wood} Carolyn, \teiHi[rend={italic}]{The Indian Summer of Bolognese Painting: Gregory XV (1621-23) and Ludovisi art patronage in Rome}, thèse de doctorat, The University of North Carolina, 1988.
\end{teiBibl}
              \begin{teiBibl}[xmlid={bibl-514}]
\teiHi[rend={small-caps}]{Ziller Camenietzki} Carlos, \teiHi[rend={italic}]{L’harmonie du monde au }\teiHi[rend={small-caps italic}]{xvii}\teiHi[rend={italic sup}]{e}\teiHi[rend={italic}]{ siècle. Essai sur la pensée scientifique d’Athanasius Kircher}, thèse de doctorat, Maurice Clavelin (dir.), Université de Paris \teiHi[rend={small-caps}]{IV}, 1995.
\end{teiBibl}
            
\end{teiDiv}
          
\end{teiDiv}
        
\end{teiElement}
      
\end{teiElement}
      \begin{teiElement}[name={group},type={loi},data-page-title={Table des figures},data-include-href={Ch15\_Table\_des\_figures.xml},xmlid={loi-001}]

        \begin{teiElement}[name={front}]

          \begin{teiDiv}[type={titlePage},xmlid={titlepage-017}]

            \teiP[rend={title-main},xmlid={p-518}]{Table des figures}
          
\end{teiDiv}
        
\end{teiElement}
        \begin{teiElement}[name={body},xmlid={body-017}]

          \teiP[xmlid={p1-16}]{Figure 1. \teiHi[rend={italic}]{Antiphonaire de Léon X}, Rome, ms. Capp.Sist.10 : Amico Aspertini,n fol. 66v (détail). © Biblioteca Apostolica Vaticana, 2025.}
          \teiP[xmlid={p2-15}]{Figure 2. \teiHi[rend={italic}]{Praeparatio ad missam pontificalem}, New York, The Pierpont Morgan Library, ms. H.6 : Attavante degli Attavanti, fol. Iv-IIr. Gift of the Heineman. The Morgan Library \& Museum, New York.}
          \teiP[xmlid={p3-15}]{Figure 3. \teiHi[rend={italic}]{Praeparatio ad missam pontificalem}, New York, The Pierpont Morgan Library, ms. H.6 : Attavante degli Attavanti, fol. 7v. Gift of the Heineman. The Morgan Library \& Museum, New York.}
          \teiP[xmlid={p4-15}]{Figure 4. \teiHi[rend={italic}]{Fiore di Virtù}, Florence, Biblioteca Riccardiana, ms. Ricc. 1774 : Mariano del Buono, fol. 18r. Su concessione del Ministero della Cultura Biblioteca Riccardiana.}
          \teiP[xmlid={p5-15}]{Figure 5. \teiHi[rend={italic}]{Fior di Virtù} Florence, Biblioteca Riccardiana, ms. Ricc.1763 : Giovanni di Gregorio d’Antonio Ghinghi (copiste), fol. 13r. Su concessione del Ministero della Cultura Biblioteca Riccardiana.}
          \teiP[xmlid={p6-15}]{Figure 6. \teiHi[rend={italic}]{Motets}, Jena, ms. 2, fol. 2r (détail). Thüringer Universitäts- und Landesbibliothek.}
          \teiP[xmlid={p7-15}]{Figure 7. \teiHi[rend={italic}]{Motets}, atelier de Petrus Alamire, Rome, ms. Capp.Sist.36, fol. 66v (détails). © Biblioteca Apostolica Vaticana, 2025.}
          \teiP[xmlid={p8-15}]{Figure 8. Ferdinando Ruano, \teiHi[rend={italic}]{Registre des dépenses}, Vat. lat. 3968, fol. 1r. © Biblioteca Apostolica Vaticana, 2025.}
          \teiP[xmlid={p9-14}]{Figure 9. Ferdinando Ruano, \teiHi[rend={italic}]{Lettres de Léon IX}, 1551, Vat. lat. 3790, fol. 3v. © Biblioteca Apostolica Vaticana, 2025.}
          \teiP[xmlid={p10-14}]{Figure 10. Ferdinando Ruano, \teiHi[rend={italic}]{Lettres d’Hildebert, évêque du Mans}, 1551, Vat. lat. 3841, fol. 1r. © Biblioteca Apostolica Vaticana, 2025.}
          \teiP[xmlid={p11-14}]{Figure 11. Ferdinando Ruano, saint Jérôme, \teiHi[rend={italic}]{Exposition sur les Psaumes}, 1554, Vat. lat. 317, fol. 2v. © Biblioteca Apostolica Vaticana, 2025.}
          \teiP[xmlid={p12-14}]{Figure 12. Missel, ms. Capp.Sist.154, fol. 3v (détail). © Biblioteca Apostolica Vaticana, 2025.}
          \teiP[xmlid={p13-14}]{Figure 13. Missel, ms. Capp.Sist.154, fol. 4r (détail). © Biblioteca Apostolica Vaticana, 2025.}
          \teiP[xmlid={p14-11}]{Figure 14. \teiHi[rend={italic}]{Officium ad Tertiam}, ms. Capp.Sist.67, fol. 4v (détail). © Biblioteca Apostolica Vaticana, 2025.}
          \teiP[xmlid={p15-10}]{Figure 15. \teiHi[rend={italic}]{Officium ad Tertiam}, ms. Capp.Sist.67, fol. 5r (détail). © Biblioteca Apostolica Vaticana, 2025.}
          \teiP[xmlid={p16-10}]{Figure 16.\teiHi[rend={italic}]{ Officium ad Tertiam}, ms. Capp.Sist.67, fol. 7v (détail). © Biblioteca Apostolica Vaticana, 2025.}
          \teiP[xmlid={p17-10}]{Figure 17.\teiHi[rend={italic}]{ Officium ad Tertiam}, ms. Capp.Sist.67, fol. 6v (détail). © Biblioteca Apostolica Vaticana, 2025.}
          \teiP[xmlid={p18-10}]{Figure 18. Vincent Raymond et Apollonio de’ Bonfratelli, Lectionnaire, ms. Capp.Sist.213, folio de titre. © Biblioteca Apostolica Vaticana, 2025.}
          \teiP[xmlid={p19-10}]{Figure 19. Prospero Fontana et atelier, devise de Jules III, portique de la villa Giulia. Photographie Yvan Loskoutoff.}
          \teiP[xmlid={p20-10}]{Figure 20. Vincent Raymond (attribué à), \teiHi[rend={italic}]{Évangéliaire de Jules III}, premier folio. © Bibliothèque capitulaire de Tolède.}
          \teiP[xmlid={p21-10}]{Figure 21. Missel du cardinal Girolamo Dandini, premier folio. © Bibliothèque capitulaire de Tolède.}
          \teiP[xmlid={p22-10}]{Figure 22. Vignette de titre pour Antonio Massa,\teiHi[rend={italic}]{ De exercitatione iurisperitorum}, Rome, 1550. Gallica.bnf.fr/BNF, 8-J-2363.}
          \teiP[xmlid={p23-10}]{Figure 23. Vignette de titre pour Giovanni Girolamo Albani, \teiHi[rend={italic}]{Erudita atque luculenta disputatio, }Rome, 1553. Biblioteca General Histórica de la Universidad de Salamanca (Espagne).}
          \teiP[xmlid={p24-8}]{Figure 24. Vignette de titre pour Giovanni Battista Poiana, \teiHi[rend={italic}]{De iobileo et indulgentiis}, Rome, 1550. Su concessione della Biblioteca Casanatense, Rome, MiC. Photographie Yvan Loskoutoff.}
          \teiP[xmlid={p25-8}]{Figure 25. Vignette de titre pour les \teiHi[rend={italic}]{Regulae tres S. D. N. D. Iulii diuina prouidentia Papae. III, }Rome, 1550. Biblioteca Nazionale Centrale « Vittorio Emanuele II » di Roma.}
          \teiP[xmlid={p26-8}]{Figure 26. Vignette de titre pour Silvestro Sigonio,\teiHi[rend={italic}]{ De concordia servanda,} Rome, 1550. Biblioteca Nazionale Centrale « Vittorio Emanuele II » di Roma.}
          \teiP[xmlid={p27-8}]{Figure 27. Vignette de titre pour \teiHi[rend={italic}]{Contra hebreos}…, Rome, 1554. Su concessione della Biblioteca Casanatense, Roma, MiC. Photographie Yvan Loskoutoff.}
          \teiP[xmlid={p28-7}]{Figure 28. Vignette pour le privilège du commentaire d’Eustathius sur l’\teiHi[rend={italic}]{Iliade, }vol. IV, Rome, 1550. Bibliothèque Sainte-Geneviève, FOL Y 7 INV 11 RES.}
          \teiP[xmlid={p29-7}]{Figure 29. Vignette de titre pour Uberto Foglietta, \teiHi[rend={italic}]{In festo omnium Beatorum, }Rome, 1553. Université de Turin.}
          \teiP[xmlid={p30-6}]{Figure 30. Vignette de titre pour les \teiHi[rend={italic}]{Novae regvlae cancellariae apostolicae, }Rome, 1552. Biblioteca Nazionale Centrale « Vittorio Emanuele II » di Roma.}
          \teiP[xmlid={p31-6}]{Figure 31. Lettrine pour Ambrosio Catharino, \teiHi[rend={italic}]{Enarrationes in quinque priora capita libri geneseos}, Rome, 1552. National Central Library of Rome, domaine public, PDM 1.0.}
          \teiP[xmlid={p32-6}]{Figure 32. Vignette pour \teiHi[rend={italic}]{Bullae diuersorum pontificum}, Rome, 1550. Biblioteca Nazionale Centrale « Vittorio Emanuele II » di Roma. Photographie Yvan Loskoutoff.}
          \teiP[xmlid={p33-6}]{Figure 33. Vignette pour Bartolomej Georgijevic, \teiHi[rend={italic}]{Libellus vere Christiana lectione…, }Rome, 1552. Biblioteca Nazionale Centrale « Vittorio Emanuele II » di Roma.}
          \teiP[xmlid={p34-5}]{Figure 34. Ioannes Magnus, \teiHi[rend={italic}]{Historia de omnibus Gothorum Sueonumque regibus,} Rome, 1554. Médiathèque d’Orléans, Rés.E2460, CESR, http://www.bvh.univ-tours.fr.}
          \teiP[xmlid={p35-5}]{Figure 35. Page de titre pour Palestrina, \teiHi[rend={italic}]{Missarum liber primus}, Rome, Valerio \& Luigi Dorico, 1554. Museo internazionale e biblioteca della musica, Bologne.}
          \teiP[xmlid={p36-5}]{Figure 36. Lettrine pour Palestrina, \teiHi[rend={italic}]{Missarum liber primus}, Rome, Valerio \& Luigi Dorico, 1554. Bibliothèque du conservatoire de Sainte Cécile, Rome. Photographie Yvan Loskoutoff.}
          \teiP[xmlid={p37-5}]{Figure 37. Vignette de titre pour Prospero Danza, \teiHi[rend={italic}]{La elettion del beatissimo \& summo pontifice Iulio tertio}, Venise, 1550, R.I.VI.340(int.8). © Biblioteca Apostolica Vaticana, 2025.}
          \teiP[xmlid={p38-5}]{Figure 38. Vignette de titre pour Jacopo Naccchiante, \teiHi[rend={italic}]{Enarrationes piae, doctae, et catholicae : in epistolam Pauli ad Ephesios, }Venetiis, 1553. By permission of Biblioteca civica Bertoliana, Vicenza.}
          \teiP[xmlid={p39-4}]{Figure 39. Vignette de titre pour Tommaso Giannoti Rangoni, \teiHi[rend={italic}]{De vita hominis ultra CXX. annos protrahenda, }Venise, 1553. BNF, RES RT-49. Photographie Yvan Loskoutoff.}
          \teiP[xmlid={p40-4}]{Figure 40. Vignette de titre pour Lorenzo Davidico, \teiHi[rend={italic}]{Anotomia delli vitii, }Florence, 1550, Stamp.De.Luca.V.40255. Su concessione della Biblioteca Casanatense, Rome, MiC. Photographie Yvan Loskoutoff.}
          \teiP[xmlid={p41-3}]{Figure 41. Vignette de titre pour Lorenzo Dal Gambaro Sclarici, \teiHi[rend={italic}]{De concilio liber I, }Bologne, apud Anselmum Giaccarelum, 1551. Su concessione della Biblioteca Casanatense, Rome, MiC. Photographie Yvan Loskoutoff.}
          \teiP[xmlid={p42-2}]{Figure 42. Vignette de titre pour Bartolomeo Maggi, \teiHi[rend={italic}]{De vulnerum sclopetorum et bombardarum curatione, }Bologne, 1552, Biblioteca Nazionale Centrale « Vittorio Emanuele II » di Roma. Photographie Yvan Loskoutoff.}
          \teiP[xmlid={p43-2}]{Figure 43. Gauges de’ Gozze, \teiHi[rend={italic}]{Se dalle armi, o insegne, che parlano }\lbrack{}...\rbrack{}, Roma, Vitale Mascardi, 1637, page de titre, Biblioteca Nazionale Centrale « Vittorio Emanuele II » di Roma. Photographie Yvan Loskoutoff.}
          \teiP[xmlid={p44-2}]{Figure 44. Silvestro Pietrasanta, \teiHi[rend={italic}]{Tesserae gentilitiae}, Romae, Typis haeredum Francisci Corbelletti, 1638, page de titre gravée par Johann Friedrich Greuter sur un dessin de Giovan Francesco Romanelli. © Getty Research Institute, Los Angeles (92-B22642).}
          \teiP[xmlid={p45-2}]{Figure 45. Silvestro Pietrasanta, \teiHi[rend={italic}]{Tesserae gentilitiae}, p. 661. © Getty Research Institute, Los Angeles (92-B22642).}
          \teiP[xmlid={p46-2}]{Figure 46. Silvestro Pietrasanta, \teiHi[rend={italic}]{Tesserae gentilitiae}, p. 677, détail des effigies de Richard II et de la reine Anne sur le tableau perdu du Collège anglais. © Getty Research Institute, Los Angeles (92-B22642).}
          \teiP[xmlid={p47-2}]{Figure 47. Xylographie pour les \teiHi[rend={italic}]{Memorie sepolcrali} figurant la pierre tombale perdue d’Angelotto Normanni († 1378), autrefois dans l’église de Sant’Antonio à Camaldoli. Collection particulière, Fabrizio Federici.}
          \teiP[xmlid={p48-2}]{Figure 48. Écu d’Everso II dell’Anguillara, première moitié du \teiHi[rend={small-caps}]{xv}\teiHi[rend={sup}]{e} siècle, Rome, Hôpital du Latran. Photographie Fabrizio Federici.}
          \teiP[xmlid={p49-2}]{Figure 49. Armes de la famille Casali, 1429, Rome, Palais Casali. Photographie Fabrizio Federici.}
          \teiP[xmlid={p50-2}]{Figure 50. Xylographie représentant les armes de la famille Casali, d’après G. B. Casali, \teiHi[rend={italic}]{De profanis, et sacris veteribus ritibus}, Romae, Ex typographia Andreae Phaei, 1645. Ville de Grenoble, bibliothèque municipale, C.1969.}
          \teiP[xmlid={p51-2}]{Figure 51. Xylographie pour le traité des \teiHi[rend={italic}]{Memorie sepolcrali} figurant l’écu Gualdi sur le flanc nord du Palais sénatorial. Collection particulière, Fabrizio Federici.}
          \teiP[xmlid={p52-2}]{Figure 52. Armes d’un sénateur de la famille Gualdi, Galeotto ou Francesco, première moitié du \teiHi[rend={small-caps}]{xvi}\teiHi[rend={sup}]{e} siècle, Rome, Palais sénatorial. Photographie Fabrizio Federici.}
          \teiP[xmlid={p53-2}]{Figure 53. Pierre tombale de Giuliano Porcari († 1282), Rome, San Giovanni della Pigna. Photographie Fabrizio Federici.}
          \teiP[xmlid={p54-2}]{Figure 54. Xylographie pour le traité des \teiHi[rend={italic}]{Memorie sepolcrali} figurant la pierre tombale de Pietro Muti († 1390) dans l’église de San Pantaleo. Collection particulière, Fabrizio Federici.}
          \teiP[xmlid={p55-2}]{Figure 55. Xylographie pour le traité des \teiHi[rend={italic}]{Memorie sepolcrali} figurant la pierre tombale de Lello Maddaleni († 1390) dans l’église de Santa Maria sopra Minerva. Collection particulière, Fabrizio Federici.}
          \teiP[xmlid={p56-2}]{Figure 56. Pierre tombale de Lello Maddaleni, Rome, Santa Maria sopra Minerva. Photographie Fabrizio Federici.}
          \teiP[xmlid={p57}]{Figure 57. Armes de la famille Cavalieri à l’entrée de la Tour du Papito, Rome, Largo di Torre Argentina. Photographie Fabrizio Federici.}
          \teiP[xmlid={p58}]{Figure 58. Armes de Taddeo di Cecco da Barberino, 2\teiHi[rend={sup}]{e} moitié du \teiHi[rend={small-caps}]{xiv}\teiHi[rend={sup}]{e} siècle, Barberino di Val d’Elsa, chapelle de San Bartolo, sacristie. Photographie Gruppo Archeologico Achu APS.}
          \teiP[xmlid={p59}]{Figure 59. Giovanni Luigi Valesio, \teiHi[rend={italic}]{Stemma del cardinale Innocenzo X Panfili e figure allegoriche, }gravure. Bologne, Pinacoteca Nazionale, Gabinetto Disegni e Stampe, 83 Picc. Coll. II. Attribution 4.0 International (CC BY 4.0).}
          \teiP[xmlid={p60}]{Figure 60. Rome, Santa Maria in Vallicella, Paolo Maruscelli et collaborateurs, mémorial du cardinal Girolamo Pamphili, 1634. Photographie Daniela del Pesco.}
          \teiP[xmlid={p61}]{Figure 61. Rome, palais Pamphili, salle des Paysages, armes papales d’Innocent X Pamphili. © Daniela del Pesco. Photographie Daniela del Pesco.}
          \teiP[xmlid={p62}]{Figure 62. Gubbio, palais Pamphili, façade sur la place. Photographie Daniela del Pesco.}
          \teiP[xmlid={p63}]{Figure 63. Gubbio, palais Pamphili, salle à l’étage noble, emblème de la colombe Pamphili inscrit dans la cheminée. Photographie Daniela del Pesco.}
          \teiP[xmlid={p64}]{Figure 64. Gubbio, palais Pamphili, salon à l’étage noble, détail de la décoration à fresque avec l’emblème de la colombe Pamphili. Photographie Daniela del Pesco.}
          \teiP[xmlid={p65}]{Figure 65. Gubbio, ancien cimetière près de l’église San Secondo, chapelle Pamphili, Jacopo Bedi, \teiHi[rend={italic}]{Martyre de saint Sébastien}, 1458, fresque (détail). Photographie Daniela del Pesco.}
          \teiP[xmlid={p66}]{Figure 66. Gubbio, ancien cimetière près de l’église San Secondo, chapelle Pamphili, pierre tombale avec l’écu des Pamphili. Photographie Daniela del Pesco.}
          \teiP[xmlid={p67}]{Figure 67. Gubbio, palais des Consuls, architrave du portail d’entrée avec emblèmes de la ville (1332-1336). Photographie Daniela del Pesco.}
          \teiP[xmlid={p68}]{Figure 68. \teiHi[rend={italic}]{Gubbio Città Regia Antichiss\lbrack{}ima\rbrack{} dell'Umbria}, gravure, d’après un dessin d’Ignazio Cassetta, Johannes Blaeu, Amsterdam, 1663, reproduit dans Johannes Blaeu, \teiHi[rend={italic}]{Nouveau théâtre d’Italie ou description exacte de ses villes, palais, églises, etc. Tome Quatrième}, À Amsterdam, Par le soin de Pierre Mortier Libraire, 1704. Détail montrant la maison des Pamphili, indiquée par la lettre Z. Institut national d’histoire de l’art (France), Attribution 4.0 International (CC BY 4.0).}
          \teiP[xmlid={p69}]{Figure 69. Armoiries « \teiHi[rend={italic}]{Romae Panfiliorum} », dans Silvestro Pietrasanta, \teiHi[rend={italic}]{Tesserae gentilitiae,} Romae, Typis haered. Francisci Corbelletti, 1638, p. 438. © Getty Research Institute, Los Angeles ((92-B22642).}
          \teiP[xmlid={p70}]{Figure 70. Armoiries de la famille Pamphili (sans attributs cardinalices ou papaux), Rome, palais Pamphili, escalier de l’aile vers la place Pasquino. Photographie Daniela del Pesco.}
          \teiP[xmlid={p71}]{Figure 71. Rome, palais Pamphili, loggia de la cour centrale, angle nord-est. Photographie Daniela del Pesco.}
          \teiP[xmlid={p72}]{Figure 72. Rome, palais Pamphili, loggia de la cour centrale (détail). Photographie Daniela del Pesco.}
          \teiP[xmlid={p73}]{Figure 73. Portrait d’Innocent X, dans Athanasius Kircher, \teiHi[rend={italic}]{Obeliscus Pamphilius}, gravure. © Emblem Collection of the University of Illinois.}
          \teiP[xmlid={p74}]{Figure 74. Jacopo Antonio Fancelli d’après un dessin de Gian Lorenzo Bernini, armes d’Innocent X Pamphili, marbre, Fontaine des Quatre Fleuve (armoiries côté nord). Photographie Daniela del Pesco.}
          \teiP[xmlid={p75}]{Figure 75. Antonio Raggi, d’après un dessin de Gian Lorenzo Bernini, armes d’Innocenzo X Pamphili, marbre, Fontaine des Quatre Fleuves (armoiries côté sud). Jastrow, domaine public, via Wikimedia Commons.}
          \teiP[xmlid={p76}]{Figure 76. Basilique de Saint-Jean-de-Latran, intérieur. Photographie Daniela del Pesco.}
          \teiP[xmlid={p77}]{Figure 77. Basilique de Saint-Jean-de-Latran, décor en stuc avec l’écu du pape Pamphili, double porte entre la basilique et le palais du Latran, d’après un dessin de Francesco Borromini. Photographie Daniela del Pesco.}
          \teiP[xmlid={p78}]{Figure 78. Francesco Borromini, étude pour la porte double entre le palais du Latran et la basilique, dessin à la mine de plomb (détail). Vienne, Albertina, Az Rom 416.}
          \teiP[xmlid={p79}]{Figure 79. Francesco Borromini, étude pour la verrière de la fenêtre de la façade orientale de Saint-Jean-de-Latran, dessin à la mine de plomb. Vienne, Albertina, Az Rom 454.}
          \teiP[xmlid={p80}]{Figure 80. Armes d’Innocent X Pamphili à la fontaine des Quatre Fleuves, côté nord, dans Filippo Juvarra, \teiHi[rend={italic}]{Raccolta Di Varie Targhe di Roma Fatte Da Professori Primarj In Roma, Disegnate, ed intagliate,} Antonio De Rossi ed., a Roma, 1711, pl. 15, gravure. @ Getty research institute, Los Angeles (1379-314).}
          \teiP[xmlid={p81}]{Figure 81. Armes d’Innocent X Pamphili, à la fontaine des Quatre Fleuves, côté sud, dans Filippo Juvarra\teiHi[rend={italic}]{, Raccolta Di Varie Targhe di Roma Fatte Da Professori Primarj In Roma, Disegnate, ed intagliate}, Antonio De Rossi ed., a Roma, 1711, pl. 18, gravure. @ Getty research institute, Los Angeles (1379-314).}
          \teiP[xmlid={p82}]{Figure 82. Armes d’Innocent X Pamphili à Saint-Jean-de-Latran, arche centrale de la nef principale, in Filippo Juvarra, \teiHi[rend={italic}]{Raccolta Di Varie Targhe di Roma Fatte Da Professori Primarj In Roma, Disegnate, ed intagliate}, Antonio De Rossi ed., a Roma, 1711, pl. 30, gravure. @ Getty research institute, Los Angeles (1379-314).}
          \teiP[xmlid={p83}]{Figure 83. Filippo Juvarra, étude pour la pl. 30 de la \teiHi[rend={italic}]{Raccolta Di Varie Targhe }(Écu Pamphili, Saint-Jean-de-Latran, arche centrale de la nef principale), dessin, MM 117, plume, crayon, encre marron, aquarelle. New York, The Metropolitan Museum of Art, Rogers Fund, 1969, 69.655.}
          \teiP[xmlid={p84}]{Figure 84. \teiHi[rend={italic}]{Écus}, Bologne, Porte Galliera, façade vers la ville. Sailko, licence CC BY-SA 3.0, via Wikimedia Commons.}
          \teiP[xmlid={p85}]{Figure 85. Anonyme, \teiHi[rend={italic}]{Miniature du primier bimestre 1658}, Bologne, Archives d’État, \teiHi[rend={italic}]{Anziani Consoli, Insignia}, vol. VIII, p. 28. Photographie Sonia Cavicchioli.}
          \teiP[xmlid={p86}]{Figure 86. Giovanni Andrea Abanti, \teiHi[rend={italic}]{Miniature du cinquième bimestre 1658}, Bologne, Archives d’État, \teiHi[rend={italic}]{Anziani Consoli, Insignia}, vol. VIII, p. 38. Photographie Sonia Cavicchioli.}
          \teiP[xmlid={p87}]{Figure 87. Giovanni Andrea Abanti, \teiHi[rend={italic}]{Miniature du sixième bimestre 1660}, Bologne, Archives d’État, \teiHi[rend={italic}]{Anziani Consoli, Insignia}, vol. VIII, p. 46. Photographie Sonia Cavicchioli.}
          \teiP[xmlid={p88}]{Figure 88. Carlo Cignani e Emilio Taruffi, \teiHi[rend={italic}]{Entrée de Paul III Farnèse à Bologne en 1543}, Bologne, Palais Public, sala Regia. Archivio fotografico Musei Civici d’Arte Antica di Bologna, avec l'aimable autorisation de Silvia Battistini, directrice, que nous remercions.}
          \teiP[xmlid={p89}]{Figure 89. Albert Clouwet, sur un dessin de Pietro Martire Neri, \teiHi[rend={italic}]{Le cardinal Pietro Vidoni}, in \teiHi[rend={italic}]{Effigies cardinalium nunc viventium, }Rome, Giovanni Giacomo de Rossi, 1660-1676. © The Trustees of the British Museum.}
          \teiP[xmlid={p90}]{Figure 90. Giovanni Andrea Abanti, \teiHi[rend={italic}]{Miniature du troisième bimestre 1662}, Bologne, Archive d’État, \teiHi[rend={italic}]{Anziani Consoli, Insignia}, vol. VIII, p. 56. Photographie Sonia Cavicchioli.}
          \teiP[xmlid={p91}]{Figure 91. Domenico Santi, dit le Mengazzino et Giovan Battista Caccioli, \teiHi[rend={italic}]{Galleria Vidoniana}, Appartement du légat, Bologne, Palais communal, vue d’ensemble. Di Waltre manni - Opera propria, licence CC BY-SA 4.0, via Wikimedia Commons.}
          \teiP[xmlid={p92}]{Figure 92. Domenico Fontana, frontispice des \teiHi[rend={italic}]{Ossequii poetici nell’ammirarsi in Bologna l’Heroica risoluzione della sereniss. Principessa Caterina Farnese di monacarsi nel Ven. Monastero delle Madri carmelitane scalze di Santa Teresa in Parma}, In Bologna, Per Giacomo Monti, 1661, Bologne, Bibliothèque communale de l’Archiginnasio.}
          \teiP[xmlid={p93}]{Figure 93. Giovan Battista Caccioli et Domenico Santi, dit le Mengazzino, \teiHi[rend={italic}]{Jupiter place l’étoile chigienne sur la tiare présentée par la Justice et la Paix}, Bologne, Palais communal, Galleria Vidoniana, détail de la voûte. Archivio fotografico Musei Civici d’Arte Antica di Bologna, avec l'aimable autorisation de Silvia Battistini, directrice, que nous remercions.}
          \teiP[xmlid={p94}]{Figure 94. Giovan Battista Caccioli et Domenico Santi, dit le Mengazzino, \teiHi[rend={italic}]{La Prudence, inspirée par l’étoile chigienne tient la Fortune en lisière}, Bologne, Palais communal, Galleria Vidoniana, détail de la voûte. Archivio fotografico Musei Civici d’Arte Antica di Bologna, avec l'aimable autorisation de Silvia Battistini, directrice, que nous remercions.}
          \teiP[xmlid={p95}]{Figure 95. Giovan Battista Caccioli et Domenico Santi, dit le Mengazzino, \teiHi[rend={italic}]{Le Pouvoir domine les puissances infernales}, Bologne, Palais communal, Galleria Vidoniana, détail de la voûte. Archivio fotografico Musei Civici d’Arte Antica di Bologna, avec l'aimable autorisation de Silvia Battistini, directrice, que nous remercions.}
          \teiP[xmlid={p96}]{Figure 96. Giovanni Battista Piranesi, \teiHi[rend={italic}]{Le antichità romane}, 1756, détail du frontispice, 1\teiHi[rend={sup}]{er} état. Florence, proprietà Stato, Accademia di Belle Arti di Firenze, no 0900452291-8, licence CC-BY 4.0.}
          \teiP[xmlid={p97}]{Figure 97. Giovanni Battista Piranesi, \teiHi[rend={italic}]{Le antichità romane}, 1756, détail du frontispice, 2\teiHi[rend={sup}]{e} état, Stampe.VI.5. © Biblioteca Apostolica Vaticana, 2025.}
          \teiP[xmlid={p98}]{Figure 98. Giovanni Battista Piranesi, « Frontispiece: Urbis Aeternae Vestigia (The Footsteps of the Eternal City) », \teiHi[rend={italic}]{Le antichità romane}, 1756, etching, 515 × 760 mm, détail du frontispice, 3\teiHi[rend={sup}]{e} état., inv. BdH boek 2 a 3 (PK). Collection Museum Boijmans Van Beuningen, Rotterdam. Gift Dr J. C. J. Bierens de Haan / photographie : Rik Klein Gotink.}
          \teiP[xmlid={p99}]{Figure 99. Giovanni Battista Piranesi, \teiHi[rend={italic}]{Vasi, candelabri, …}, 1778, frontispice du vol. 2. © Getty Research Institute, Los Angeles (90-B18452).}
          \teiP[xmlid={p100}]{Figure 100. Dessin d’un rhyton, album Ghezzi, Ott. lat. 3105. © Biblioteca Apostolica Vaticana, 2025.}
          \teiP[xmlid={p101}]{Figure 101. Giovanni Battista Piranesi, détail du tombeau de Cecilia Metella, \teiHi[rend={italic}]{Vedute di Roma}. Collection particulière, Yvan Loskoutoff.}
          \teiP[xmlid={p102}]{Figure 102. Giovanni Battista Piranesi, dessin des armes de Clément XIII Rezzonico. The Met Museum, 52.570.155. © The Elisha Whittelsey Collection, The Elisha Whittelsey Fund, 1952.}
          \teiP[xmlid={p103}]{Figure 103. Giovanni Battista Piranesi, \teiHi[rend={italic}]{De Romanorum magnificentia et architectura}, 1761, frontispice. Domaine public, via Wikimedia Commons.}
          \teiP[xmlid={p104}]{Figure 104. Giovanni Battista Piranesi, \teiHi[rend={italic}]{Le rovine del castello dell’acqua Giulia, }1761, cul-de-lampe, Yale University Art Gallery, The Arthur Ross Collection, 2012.159.21.}
          \teiP[xmlid={p105}]{Figure 105. Giovanni Battista Piranesi, \teiHi[rend={italic}]{Lapides Capitolini sive fasti consularesque Romanorum, }Romae, Typis Generosi Salomoni, 1762, frontispice. Livres rares et collections spécialisées, Bibliothèques McGill.}
          \teiP[xmlid={p106}]{Figure 106. Giovanni Battista Piranesi, \teiHi[rend={italic}]{Antichità d’Albano e di Castel Gandolfo}, Roma, s. l., 1764, frontispice. © Getty Research Institut, Los Angeles (90-B18472).}
          \teiP[xmlid={p107}]{Figure 107. Giovanni Battista Piranesi, \teiHi[rend={italic}]{Diverse maniere d’adornare i cammini,} 1769, frontispice. Gallica.bnf.fr/BNF, Montpellier, médiathèque centrale Émile Zola, RTG008.}
          \teiP[xmlid={p108}]{Figure 108. Giovanni Battista Piranesi, portail du prieuré de Malte, Rome (détail). Photographie Yvan Loskoutoff.}
          \teiP[xmlid={p109}]{Figure 109. Giovanni Battista Piranesi, portail du prieuré de Malte, Rome (détail). Photographie Yvan Loskoutoff.}
          \teiP[xmlid={p110}]{Figure 110. Giovanni Battista Piranesi, façade de Santa Maria del Priorato, Rome (détail). © Archi Diap. Dipartimento di Architettura. Progetto Sapienza Università di Roma, Paolo Marcoaldi.}
          \teiP[xmlid={p111}]{Figure 111. Giovanni Battista Piranesi, extérieur de l’abside du chœur de Santa Maria del Priorato, Rome (détail). Photographie Yvan Loskoutoff.}
          \teiP[xmlid={p112}]{Figure 112. Giovanni Battista Piranesi, armes de Giovan Battista Rezzonico, chœur de Santa Maria del Priorato, Rome. Photographie Yvan Loskoutoff.}
          \teiP[xmlid={p113}]{Figure 113. Giovanni Battista Piranesi, composition héraldique en l’honneur de Giovan Battista Rezzonico, Santa Maria del Priorato, Rome. Photographie Yvan Loskoutoff.}
          \teiP[xmlid={p114}]{Figure 114. Giacomo Grimaldi, \teiHi[rend={italic}]{Instrumenta translationum}, Milan, bibliothèque Ambrosienne, ms. F 227 inf., copie de la scène de la bénédiction du cycle de fresques de la loggia du Latran. Domaine public, via Wikimedia Commons.}
          \teiP[xmlid={p115}]{Figure 115. La prééminence des insignes pontificaux sur le Palazzo Vechio. Photographie Matteo Ferrari.}
          \teiP[xmlid={p116}]{Figure 116. Architrave du palais du Peuple de Gubbio. Photographie Laurent Hablot.}
          \teiP[xmlid={p117}]{Figure 117. Relief de pierre de la cité d’Ascoli Piceno. Photographie Andrea Di Benedetto.}
          \teiP[xmlid={p118}]{Figure 118. Allégorie héraldique de la ville de Florence, vers 1360-1370, par Chini Galileo (fin du \teiHi[rend={small-caps}]{xix}\teiHi[rend={sup}]{e} siècle). Florence, Palazzo dell’Arte della Seta, 0900749643. Licence CC-BY 4.0.}
          \teiP[xmlid={p119}]{Figure 119. Loggia de Viterbe. Photographie Jean-Pierre Dalbéra.}
          \teiP[xmlid={p120}]{Figure 120. Clé de voûte aux armes de Guy de Castelnau-Calmont, évêque de Périgueux. Photographie Emmanuel Moureau.}
          \teiP[xmlid={p121}]{Figure 121. Jean II de Castelnau-Calmont et son saint patron, verrière du chœur, vers 1510-1514. Photographie Emmanuel Moureau.}
          \teiP[xmlid={p122}]{Figure 122. Vue générale des stalles de la collégiale de Prudhomat. Photographie Emmanuel Moureau.}
          \teiP[xmlid={p123}]{Figure 123. Armoiries du baron de Castelnau-Calmont, stalles du chœur. Photographie Emmanuel Moureau.}
          \teiP[xmlid={p124}]{Figure 124. Armoiries de Jean II de Castelnau-Calmont et de Marie de Culan, stalles du chœur. Photographie Emmanuel Moureau.}
          \teiP[xmlid={p125}]{Figure 125. Détail du dauphin, stalles du chœur. Photographie Emmanuel Moureau.}
          \teiP[xmlid={p126}]{Figure 126. Saint Jules I\teiHi[rend={sup}]{er}, allégorie du pape Jules II, stalles du chœur. Photographie Emmanuel Moureau.}
          \teiP[xmlid={p127}]{Figure 127. Saint Jérôme, allégorie du cardinal François Guillaume de Clermont-Lodève, stalles du chœur. Photographie Emmanuel Moureau.}
          \teiP[xmlid={p128}]{Figure 128. Détail des hermines et des feuilles de chêne sculptées sur une colonnette des stalles du chœur. Photographie Emmanuel Moureau.}
          \teiP[xmlid={p129}]{Figure 129. Bernard Salomon, dans \teiHi[rend={italic}]{La Magnificque et triumphante entrée de Carpentras, faicte à tresillustre \& trespuisssant prince ALEXANDRE FARNES, Cardinal, Legat d’Avignon, Vicechancelier du S. siege Apostolique} par Antoine Blégier, Avignon, Par Macé Bonhomme, aux Changes, 1553, p. Da verso. Collection particulière, Camille Larraz.}
          \teiP[xmlid={p130}]{Figure 130. Bernard Salomon, dans \teiHi[rend={italic}]{La Magnificque et triumphante entrée de Carpentras… }par Antoine Blégier, Avignon, 1553, p. D. Collection particulière, Camille Larraz.}
          \teiP[xmlid={p131}]{Figure 131. Avignon, Archives municipales, CC 468, Comptes de Gaspard Droin, dépenses extraordinaires, non numéroté. Collection particulière, Camille Larraz.}
          \teiP[xmlid={p132}]{Figure 132. Bernard Salomon, dans \teiHi[rend={italic}]{La Magnificque et triumphante entrée de Carpentras… }par Antoine Blégier, Avignon, 1553, p. B ii. Collection particulière, Camille Larraz.}
          \teiP[xmlid={p133}]{Figure 133. Anonyme, \teiHi[rend={italic}]{Recueil de statuts et de privilèges de la ville d’Avignon}. Avignon, Bibliothèque municipale, ms. 2834.}
          \teiP[xmlid={p134}]{Figure 134. Jean Bœuf, \teiHi[rend={italic}]{Frontispice}, vers 1598. Paris, BNF, Latin 8971.}
          \teiP[xmlid={p135}]{Figure 135. Anonyme, page de titre du \teiHi[rend={italic}]{Bullarium civitatis Avenionensis seu Bullae}, Lyon, Joannis Amati Candy, 1657. Paris, BNF, département Philosophie, histoire, sciences de l’homme, FOL-LK7-647.}
          \teiP[xmlid={p136}]{Figure 136. Portrait du cardinal Benedetto Giustiniani. Michel Natalis, d’après un dessin de Giovanni Citosibio Guidi, gravure (238 mm × 170 mm), Rome, 1636-1647, dans \teiHi[rend={italic}]{Galleria Giustiniana del Marchese Vincenzo Giustiniani}, Rome, vol. 2. BNF, département Réserve des livres rares, Rés.J.485 © Gallica.bnf.fr/BNF.}
          \teiP[xmlid={p137}]{Figure 137. Rome, église Sainte-Prisca, maître autel, Domenico Cresti,\teiHi[rend={italic}]{ Baptême de sainte Prisca}. Photographie Rocco Ronzani.}
          \teiP[xmlid={p138}]{Figure 138. Rome, église Sainte-Prisca, abside, Anastasio Fontebuoni et collaborateurs, image de saint Pierre, repeint en saint Augustin au \teiHi[rend={small-caps}]{xix}\teiHi[rend={sup}]{e} siècle. Photographie Rocco Ronzani.}
          \teiP[xmlid={p139}]{Figure 139. Rome, église Sainte-Prisca, abside, Anastasio Fontebuoni et collaborateurs, \teiHi[rend={italic}]{Ange portant la férule pontificale}. Photographie Rocco Ronzani.}
          \teiP[xmlid={p140}]{Figure 140. Rome, église Sainte-Prisca, abside, Anastasio Fontebuoni et collaborateurs, \teiHi[rend={italic}]{Ange portant les clefs de saint Pierre}. Photographie Rocco Ronzani.}
          \teiP[xmlid={p141}]{Figure 141. Rome, église Sainte-Prisca, fresques de la nef centrale, Anastasio Fontebuoni et collaborateurs, \teiHi[rend={italic}]{Saint Benoît de Nursie}. Photographie Rocco Ronzani.}
          \teiP[xmlid={p142}]{Figure 142. Rome, église Sainte-Prisca, fresques de la nef centrale, Anastasio Fontebuoni et collaborateurs, \teiHi[rend={italic}]{Saint Jean Gualbert}. Photographie Rocco Ronzani.}
          \teiP[xmlid={p143}]{Figure 143. Armes de la congrégation de Vallombreuse, miniature, Bibliothèque nationale de Florence, Conventi soppressi, C.V.1912, f. 2r (8 mai 1513).}
          \teiP[xmlid={p144}]{Figure 144. Rome, église Sainte-Prisca, fresques de la crypte, Anastasio Fontebuoni et collaborateurs, armes de Sixte Quint Peretti et, au-dessous, du cardinal Giustiniani. Photographie Rocco Ronzani.}
          \teiP[xmlid={p145}]{Figure 145. Rome, église Sainte-Prisca, fresques de la crypte, Anastasio Fontebuoni et collaborateurs, lion héraldique Peretti couronné par un ange. Photographie Rocco Ronzani.}
          \teiP[xmlid={p146}]{Figure 146. Rome, église Sainte-Prisca, décor de marbre du maître autel, armes du cardinal Giustiniani. Photographie Rocco Ronzani.}
          \teiP[xmlid={p147}]{Figure 147. Rome, église Sainte-Prisca, inscription commémorative encadrée des armes du cardinal Giustiniani et de leurs meubles (aigle éployée et château). Photographie Rocco Ronzani.}
          \teiP[xmlid={p148}]{Figure 148. Rome, église Sainte-Prisca, fresques de la crypte, Anastasio Fontebuoni et collaborateurs, armes Giustiniani (l’humidité a fait perdre l’aigle qui avait peut-être été peinte \teiHi[rend={italic}]{a tempera}). Photographie Rocco Ronzani.}
          \teiP[xmlid={p149}]{Figure 149. Rome, église Sainte-Prisca, frise de l’escalier d’accès à la crypte, détail héraldique Giustiniani (château). Photographie Rocco Ronzani.}
          \teiP[xmlid={p150}]{Figure 150. Rome, église Sainte-Prisca, frise de l’escalier d’accès à la crypte, détail héraldique Giustiniani (aigle). Photographie Rocco Ronzani.}
          \teiP[xmlid={p151}]{Figure 151. Rome, église Sainte-Prisca, crypte, meuble des armes Giustiniani (aigle) sur le bénitier. Photographie Rocco Ronzani.}
          \teiP[xmlid={p152}]{Figure 152. Rome, église Sainte-Prisca, crypte, armes de Clément VIII Aldobrandini. Photographie Rocco Ronzani.}
          \teiP[xmlid={p153}]{Figure 153. Rome, église Sainte-Prisca, abside, inscription surmontée des armes de Callixte III Borgia. Photographie Rocco Ronzani.}
          \teiP[xmlid={p154}]{Figure 154. Rome, église Sainte-Prisca, presbytère, inscription commémorative avec les armes cardinalices et papales d’Angelo Giuseppe Roncalli, saint Jean XXIII. Photographie Rocco Ronzani.}
          \teiP[xmlid={p155}]{Figure 155. Rome, église Sainte-Prisca, sacristie, chape portant l’écu prélatice de Camillo Alfonso de Romanis (1885-1950), historien de l’ordre augustin, curé de Sainte-Prisca puis vicaire général de la Cité du Vatican et préfet du sacraire apostolique (1937-1950). Photographie Rocco Ronzani.}
          \teiP[xmlid={p156}]{Figure 156. Le baldaquin de la basilique Saint-Pierre, photographie. Jebulon, licence CC0 1.0 Universel, via Wikimedia Commons.}
          \teiP[xmlid={p157}]{Figure 157. La remise des clefs, bas-relief de Ambrogio Buonvicino, 1614-1632, façade de la basilique Saint-Pierre. Vincenzo.paglione, domaine public, via Wikimedia Commons.}
          \teiP[xmlid={p158}]{Figure 158. Monument funéraire du pape Innocent VIII Cibo, 1492, Antonio Pollaiolo, basilique Saint-Pierre. Torvindus, licence CC BY-SA 3.0, via Wikimedia Commons.}
          \teiP[xmlid={p159}]{Figure 159. Monument funéraire de Mathilde de Canossa, 1644, Le Bernin, basilique Saint-Pierre. © St Peters basilica.}
          \teiP[xmlid={p160}]{Figure 160. Plan de \teiHi[rend={italic}]{Sede Vacante} au décès d’Urbain VIII Barberini, 1644. Issu de Robert Van Audenaerd, \teiHi[rend={italic}]{Nuova Et Esatta Pianta Del Conclave Con Le Funzioni}…, 1647, Domenico De Rossi, Rome. Fondo Corsini, volume 58N19, S-FC122177. Rome, Istituto centrale per la grafica, per gentile concessione del Ministero della cultura, proprietà Accademia Naxionale dei Lincei.}
          \teiP[xmlid={p161}]{Figure 161. Plan de \teiHi[rend={italic}]{Sede Vacante} au décès d’Alexandre VIII Ottoboni, avec le blason du cardinal Médicis dédicataire, anonyme, 1691, gravure. Domaine public.}
          \teiP[xmlid={p162}]{Figure 162. Cortège du \teiHi[rend={italic}]{Possesso} de Léon XI Médicis, \lbrack{}Rome\rbrack{}, Cum priuilegio et superior permissu Io. Antonij de Paulis for. Romae, 1605, gravure. © Getty Research Institute, Los Angeles (2563-899).}
          \teiP[xmlid={p163}]{Figure 163. Théâtre de la canonisation des cinq saints, 1622, Matteo Greuter, gravure. Domaine public, via Wikimedia Commons.}
          \teiP[xmlid={p164}]{Figure 164. Giovanni Luigi Valesio, \teiHi[rend={italic}]{L’Apothéose des armes Ludovisi, }c. 1623, fresque, Rome, Casino Ludovisi, voûte de la salle « \teiHi[rend={italic}]{dell’armadio de’ libri} » ou « \teiHi[rend={italic}]{della Stemma} ». Photographie Frédéric Cousinié.}
          \teiP[xmlid={p165}]{Figure 165. Raphaël, \teiHi[rend={italic}]{Angelots supportant l’écu papal}, 1508, fresque, Cité du Vatican, voûte de la chambre de la Signature. Photographie Frédéric Cousinié.}
          \teiP[xmlid={p166}]{Figure 166. Pietro da Cortona, \teiHi[rend={italic}]{Voûte du grand salon du palais Barberini}, 1632-1638, fresque, Rome (détail). Photographie Frédéric Cousinié.}
          \teiP[xmlid={p167}]{Figure 167. Pietro da Cortona, \teiHi[rend={italic}]{Voûte de l’« Anticamera de’ cortigiani » ou salle de Mars}, 1644-1646, fresque, Florence, palais Pitti. Photographie Frédéric Cousinié.}
          \teiP[xmlid={p168}]{Figure 168. Antonio Domenico Gabbiani, \teiHi[rend={italic}]{L’Apothéose de la famille Corsini}, c. 1696, huile sur toile, modèle pour la fresque du salon du palais Corsini al Parione, Florence. Photographie Frédéric Cousinié.}
          \teiP[xmlid={p169}]{Figure 169. Pietro da Cortona, \teiHi[rend={italic}]{Décor de la voûte de la petite galerie du palais Barberini}, 1632, fresque, Rome. Photographie Frédéric Cousinié.}
          \teiP[xmlid={p170}]{Figure 170. Domenico Maria Canuti, \teiHi[rend={italic}]{L’Apothéose de Romulus}, c. 1675-1676, fresque, Rome, palais Altieri, voûte de l’antichambre. Photographie Frédéric Cousinié.}
          \teiP[xmlid={p171}]{Figure 171. Domenico Maria Canuti, \teiHi[rend={italic}]{L’Apothéose de Romulus}, c. 1675-1676, fresque, détail d’un des angles de la voûte avec l’étoile des Altieri. Photographie Frédéric Cousinié.}
          \teiP[xmlid={p172}]{Figure 172. Paul Bril, \teiHi[rend={italic}]{Le Martyre de saint Clément}, c. 1696-1602, fresque, Cité du Vatican, salle Clémentine. Photographie Frédéric Cousinié.}
          \teiP[xmlid={p173}]{Figure 173. Cherubino et Giovanni Alberti, \teiHi[rend={italic}]{La Sala Clementina}, c. 1696-1602, Cité du Vatican. Photographie Frédéric Cousinié.}
          \teiP[xmlid={p174}]{Figure 174. Cherubino et Giovanni Alberti, \teiHi[rend={italic}]{Anges, putti et allégories}, c. 1696-1602, fresque, Cité du Vatican, salle Clémentine, détail des parois latérales. Photographie Frédéric Cousinié.}
          \teiP[xmlid={p175}]{Figure 175. Cherubino et Giovanni Alberti, \teiHi[rend={italic}]{L’Apothéose de saint Clément}, c. 1696-1602, fresque, Cité du Vatican, salle Clémentine, voûte. Photographie Frédéric Cousinié.}
          \teiP[xmlid={p176}]{Figure 176. Cherubino et Giovanni Alberti, \teiHi[rend={italic}]{L’écu Aldobrandini}, c. 1696-1602, fresque, Cité du Vatican, salle Clémentine, détail d’un des angles de la voûte. Photographie Frédéric Cousinié.}
          \teiP[xmlid={p177}]{Figure 177. Cherubino et Giovanni Alberti, \teiHi[rend={italic}]{Allégorie, }c. 1696-1602, fresque, Cité du Vatican, salle Clémentine, détail des parois latérales. Photographie Frédéric Cousinié.}
        
\end{teiElement}
      
\end{teiElement}
      \begin{teiElement}[name={group},type={chapter},data-page-title={Les auteurs},data-include-href={Ch16\_Auteurs.xml},xmlid={chapter-011}]

        \begin{teiElement}[name={front}]

          \begin{teiDiv}[type={titlePage},xmlid={titlepage-018}]

            \teiP[rend={title-main},xmlid={p-519}]{Les auteurs}
          
\end{teiDiv}
        
\end{teiElement}
        \begin{teiElement}[name={body},xmlid={body-018}]

          \teiP[xmlid={p1-17}]{Martine \teiHi[rend={small-caps}]{Boiteux} est agrégée d’histoire-géographie, ancienne membre de l’École française de Rome, professeure à l’École des Hautes Études en Sciences Sociales. Elle étudie l’histoire, l’histoire de l’art et l’anthropologie culturelle de la Rome pontificale (\teiHi[rend={small-caps}]{xv}-\teiHi[rend={small-caps}]{xix}\teiHi[rend={sup}]{e} siècle), avec des travaux sur les résidences, les cérémonies et rituels sociaux, politiques et religieux, leur déroulement dans l’espace urbain, les apparats et les architectures éphémères, les fêtes populaires et aristocratiques, l’histoire urbaine et l’utilisation des fêtes, et enfin la fonction de la représentation. Les sources utilisées conjuguent documents écrits et figurés pour une approche globale du patrimoine physique, du patrimoine immatériel, du pouvoir des images. L’étude de l’iconographie, spécialement des saintes mystiques, a été abordée. Un autre champ de recherche est la communauté juive. Elle a publié environ 150 articles et co-dirigé quatre ouvrages. Elle a également créé et dirigé deux revues : \teiHi[rend={italic}]{Temps Libre }de 1981 à 1986 et \teiHi[rend={italic}]{Epure }de 1983 à 1998.}
          \teiP[xmlid={p2-16}]{Sonia \teiHi[rend={small-caps}]{Cavicchioli }est docteure et professeure d’histoire de l’art moderne à l’université de Bologne. Elle s’intéresse, entre autres, à l’art et à la décoration intérieure des \teiHi[rend={small-caps}]{xvi}\teiHi[rend={sup}]{e }et \teiHi[rend={small-caps}]{xvii}\teiHi[rend={sup}]{e} siècles en Émilie, ainsi qu’au mécénat des Estensi, seigneurs de Ferrare et de Modène, à l’époque moderne. Parmi ses publications récentes figurent la monographie \teiHi[rend={italic}]{Dipingere pensieri. Francesco IV d’Austria-Este e la decorazione del Palazzo Ducale di Modena (1814-1846)} (EUT, 2022), et le recueil d’études \teiHi[rend={italic}]{Giuseppe Maria Soli (1747-1822). Un architetto modenese fra Antico Regime e Restaurazione}, sous la direction de Sonia Cavicchioli et de Vincenzo Vandelli (Académie nationale des sciences, lettres et arts, 2024).}
          \teiP[xmlid={p3-16}]{Frédéric \teiHi[rend={small-caps}]{Cousinié} est professeur d’histoire et théorie de l’art à l’université de Rouen Normandie. Auteur d’une douzaine d’ouvrages sur la période moderne, entre France et Italie, il a notamment publié : \teiHi[rend={italic}]{Paysage du paysage. Nicolas Poussin, Claude Gellée Le Lorrain, Sébastien Bourdon }(Les presses du réel, 2022) ; \teiHi[rend={italic}]{Tensivité des images. Surgissement, Révélation, Extase, Apothéose dans la France du }\teiHi[rend={ small-caps italic}]{xvii}\teiHi[rend={italic sup}]{e}\teiHi[rend={italic}]{ siècle} (Mare \& Martin, 2023) ; et \teiHi[rend={italic}]{L’espace imaginal : méditer dans l’image au }\teiHi[rend={ small-caps italic}]{xvii}\teiHi[rend={italic sup}]{e}\teiHi[rend={italic}]{ siècle} (1:1, « ars », 2024).}
          \teiP[xmlid={p4-16}]{Daniela \teiHi[rend={small-caps}]{del Pesco} est professeur émérite en histoire de l’art et de l’architecture modernes à l’université de Rome III. Ses domaines de recherche principaux concernent l’architecture baroque à Rome et l’œuvre du Bernin, l’architecture et l’urbanisme des villes dans le royaume espagnol de l’Italie du Sud ainsi que l’art et les rapports diplomatiques entre l’Italie et la France au \teiHi[rend={small-caps}]{xvii}\teiHi[rend={sup}]{e} siècle.}
          \teiP[xmlid={p5-16}]{Elli \teiHi[rend={small-caps}]{Doulkaridou-Ramantani} est docteur en histoire de l’art moderne de l’université Paris I Panthéon-Sorbonne (2019). Elle a réalisé une thèse intitulée \teiHi[rend={italic}]{La Renaissance enluminée de Rome. Systèmes décoratifs dans les manuscrits du }\teiHi[rend={ small-caps italic}]{xvi}\teiHi[rend={italic sup}]{e}\teiHi[rend={italic}]{ siècle}, sous la direction de Philippe Morel. Elle est l’auteur de plusieurs articles publiés dans des revues à comité de lecture ou des actes de colloques, lesquels s’intéressent aux fonctions du décor dans les manuscrits enluminés de la Renaissance italienne. Elle a notamment coédité l’ouvrage collectif \teiHi[rend={italic}]{Jeux et enjeux du cadre dans les systèmes décoratifs de la première modernité 1500-1700} (PUR, 2019). }
          \teiP[xmlid={p6-16}]{Fabrizio \teiHi[rend={small-caps}]{Federici} a étudié l’histoire de l’art à l’université et l’École normale supérieure de Pise, où il a obtenu son doctorat en soutenant une thèse sur le collectionneur du \teiHi[rend={small-caps}]{xvii}\teiHi[rend={sup}]{e} siècle Francesco Gualdi. Ses intérêts portent sur l'histoire sociale de l’art, l’art à Rome et en Toscane au \teiHi[rend={small-caps}]{xvii}\teiHi[rend={sup}]{e} siècle, l’histoire de l’érudition et de l’antiquarianisme, la réception de l’art du Moyen Âge. De 2016 à 2018, il a été post-doctorant à la Bibliotheca Hertziana de Rome. Il est actuellement chercheur à l’université de Florence.}
          \teiP[xmlid={p7-16}]{Laurent \teiHi[rend={small-caps}]{Hablot} est directeur d’études à l’École pratique des hautes études-PSL, titulaire de la Conférence d’emblématique occidentale, membre de l’équipe SAPRAT et du conseil de Biblissima +, chercheur associé à l’IRHT. Il travaille sur les formes et les fonctions des systèmes de signes en Europe au Moyen Âge et à l'époque moderne et coordonne les bases de données en ligne ARMMA-\teiHi[rend={italic}]{Armorial monumental du Moyen Âge}, SIGILLA-\teiHi[rend={italic}]{base numérique des sceaux conservés en France} et DEVISE-\teiHi[rend={italic}]{emblématique européenne à la fin du Moyen Âge}.}
          \teiP[xmlid={p8-16}]{Camille \teiHi[rend={small-caps}]{Larraz}, après un doctorat en histoire de l’art soutenu à l’université de Genève (2021) au sein du projet « Peindre en France à la Renaissance », a intégré le programme « La peinture produite en France au \teiHi[rend={small-caps}]{xvi}\teiHi[rend={sup}]{e} siècle » à l’Institut national d’histoire de l’art à Paris, en collaboration avec le musée du Louvre. De retour à Genève depuis 2024, elle contribue notamment à différents ouvrages collectifs tant sur la peinture française de la Renaissance que sur l’art genevois des \teiHi[rend={small-caps}]{xviii}\teiHi[rend={sup}]{e} et \teiHi[rend={small-caps}]{xix}\teiHi[rend={sup}]{e} siècles.}
          \teiP[xmlid={p9-15}]{Yvan \teiHi[rend={small-caps}]{Loskoutoff}, agrégé de lettres modernes et docteur de Sorbonne Université, est professeur de littérature française des \teiHi[rend={small-caps}]{xvi}\teiHi[rend={sup}]{e} et \teiHi[rend={small-caps}]{xvii}\teiHi[rend={sup}]{e} siècles à l’université du Havre (laboratoire GRIC). Il est membre de l’Académie des jeux floraux et de la Société des antiquaires de Londres. Ses études portent sur littérature et spiritualité, littérature et symbolique. Il a publié cinq livres (sur la mode des contes de fées, l’héraldique chez le P. Le Moyne, la propagande du cardinal Mazarin, la symbolique des papes Sixte Quint et Grégoire XIII), dix actes de colloques (sur héraldique et numismatique, héraldique et papauté, Mazarin Rome et l’Italie, et l’histoire métallique de Louis XIV) et une centaine d’articles.}
          \teiP[xmlid={p10-15}]{Emmanuel \teiHi[rend={small-caps}]{Moureau}, docteur en histoire de l’art médiéval, chargé de cours à l’université de Toulouse-Jean-Jaurès, participe au programme ANR Col\&Mon sur les collégiales séculières en France de 816 à 1561. Il a publié une vingtaine d’articles sur ce sujet, et deux livres : \teiHi[rend={italic}]{Vivre en ville au temps des papes d’Avignon. Montauban 1317-1378} (La Louve, 2019), \teiHi[rend={italic}]{Un marchand au Moyen Âge. Regards sur la vie quotidienne au }\teiHi[rend={ small-caps italic}]{xiv}\teiHi[rend={italic sup}]{e}\teiHi[rend={italic}]{ siècle : les comptes de Barthélémy Bonis 1345-1365} (La Louve, 2012). Il est également conservateur des antiquités et objets d’art de Tarn-et-Garonne.}
          \teiP[xmlid={p11-15}]{Le Père Rocco \teiHi[rend={small-caps}]{Ronzani}, religieux de l’ordre de saint Augustin, a étudié à l’université pontificale grégorienne et s’est spécialisé avec un doctorat en philologie à l’Institut patristique Augustinianum de Rome. Il est archiviste-paléographe de l’École vaticane de paléographie. En 2024, il a été nommé préfet des Archives apostoliques du Vatican, membre du Comité des sciences historiques et de l’Académie pontificale des sciences.}
        
\end{teiElement}
      
\end{teiElement}
    
\end{teiElement}
    \begin{teiElement}[name={back}]

\end{teiElement}
  
\end{teiElement}

\end{teiElement}
\end{lateiDocument}

\cleardoublepage
\pagestyle{plain}
\titleformat{\chapter}[display]{\PURHTitleFont\bfseries\fontsize{16pt}{19pt}\selectfont\centering}{}{10pt}{\MakeUppercase}
\tableofcontents

\end{document}
